\documentclass[11pt,a4paper,twoside,openright]{book}

% ============================================
% BANCO ZK - Volumen I: Introducción Teórica
% Versión MODULAR - Archivos separados por capítulo
% ============================================

% --- Codificación y fuentes ---
\usepackage[utf8]{inputenc}
\usepackage[T1]{fontenc}
\usepackage{lmodern}
\usepackage{microtype}

% --- Geometría ---
\usepackage[a4paper, top=2.5cm, bottom=2.5cm, inner=3cm, outer=2.5cm]{geometry}

% --- Matemáticas ---
\usepackage{amsmath, amssymb, amsthm}
\usepackage{mathtools}
\usepackage{bm}

% --- TikZ ---
\usepackage{tikz}
\usetikzlibrary{arrows.meta, shapes, positioning, calc, decorations.pathreplacing, patterns, fit}

% --- Tablas ---
\usepackage{booktabs}
\usepackage{array}
\usepackage{tabularx}
\usepackage{longtable}

% --- Hyperref ---
\usepackage[colorlinks=true, linkcolor=blue!60!black, citecolor=green!50!black, urlcolor=blue!70!black]{hyperref}

% --- Encabezados ---
\usepackage{fancyhdr}
\pagestyle{fancy}
\fancyhf{}
\fancyhead[LE]{\small\textit{\leftmark}}
\fancyhead[RO]{\small\textit{\rightmark}}
\fancyfoot[LE,RO]{\thepage}
\renewcommand{\headrulewidth}{0.4pt}

% --- Espaciado ---
\usepackage{setspace}
\onehalfspacing

% --- Teoremas ---
\theoremstyle{definition}
\newtheorem{definition}{Definición}[chapter]
\newtheorem{example}{Ejemplo}[chapter]

\theoremstyle{plain}
\newtheorem{theorem}{Teorema}[chapter]
\newtheorem{proposition}[theorem]{Proposición}

\theoremstyle{remark}
\newtheorem{remark}{Observación}[chapter]

% --- Comandos matemáticos ---
\newcommand{\F}{\mathbb{F}}
\newcommand{\Fp}{\mathbb{F}_p}
\newcommand{\Z}{\mathbb{Z}}
\newcommand{\N}{\mathbb{N}}
\newcommand{\G}{\mathbb{G}}
\newcommand{\Hash}{\mathcal{H}}

% --- Cajas tcolorbox ---
\usepackage{tcolorbox}
\tcbuselibrary{skins, breakable}

\newtcolorbox{keyidea}{
    enhanced, colback=yellow!5!white, colframe=orange!70!black,
    title={\textbf{Idea Clave}}, breakable
}

\newtcolorbox{bancozk}{
    enhanced, colback=blue!5!white, colframe=blue!70!black,
    title={\textbf{En Banco ZK}}, breakable
}

\newtcolorbox{mathnote}{
    enhanced, colback=green!3!white, colframe=green!50!black,
    title={\textbf{Matemáticas}}, breakable
}

\newtcolorbox{intuicion}{
    enhanced, colback=blue!3!white, colframe=blue!50!black,
    title={\textbf{Intuición}}, breakable
}

\newtcolorbox{advertencia}{
    enhanced, colback=red!5!white, colframe=red!70!black,
    title={\textbf{Advertencia}}, breakable
}

% ============================================
% INICIO DEL DOCUMENTO - VOLUMEN 1
% ============================================
\begin{document}
% ============================================
% PORTADA
% ============================================
\begin{titlepage}
    \centering
    \vspace*{1cm}
    
    {\Large\scshape Serie Banco ZK}\\[0.3cm]
    {\small 3 Vol\'umenes $\cdot$ 1000+ p\'aginas $\cdot$ 20 Cap\'itulos}\\[0.5cm]
    \rule{\textwidth}{1.5pt}\\[0.5cm]
    
    {\Huge\bfseries Volumen I}\\[0.5cm]
    {\LARGE Fundamentos Te\'oricos}\\[0.3cm]
    {\large Cap\'itulos 1--5: Criptograf\'ia, Blockchain, Zero-Knowledge}\\[1.5cm]
    
    \begin{tikzpicture}
        \node[rectangle, draw, rounded corners, minimum width=12cm, minimum height=5cm, fill=zkblue!5] {
            \begin{minipage}{11cm}
            \centering
            \textbf{Contenido del Volumen 1:}\\[0.5cm]
            \small
            \textbf{Parte I:} Fundamentos Criptogr\'aficos\\
            Grupos $\bullet$ Curvas El\'ipticas $\bullet$ Hash $\bullet$ Firmas Digitales\\[0.3cm]
            \textbf{Parte II:} Blockchain y Escalabilidad\\
            Consenso $\bullet$ L1/L2/L3 $\bullet$ Rollups $\bullet$ Sidechains\\[0.3cm]
            \textbf{Parte III:} Zero-Knowledge Proofs\\
            Definici\'on $\bullet$ Propiedades $\bullet$ Protocolos Interactivos $\bullet$ SNARKs\\[0.3cm]
            \textbf{Parte IV:} Circuitos ZK\\
            Circuitos Aritm\'eticos $\bullet$ R1CS $\bullet$ QAP $\bullet$ Groth16
            \end{minipage}
        };
    \end{tikzpicture}
    
    \vfill
    
    {\Large Walter Milstein}\\[0.5cm]
    {\large 2025}
\end{titlepage}
% ============================================
% P\'AGINA DE COPYRIGHT
% ============================================
\thispagestyle{empty}
\null\vfill
\begin{center}
\textbf{Banco ZK: Criptograf\'ia y Privacidad en Blockchain}\\
\textbf{Volumen 1: Fundamentos Te\'oricos}\\[1cm]
\textcopyright\ 2025 Walter Milstein\\
Todos los derechos reservados.\\[0.5cm]
Primera Edici\'on: 2025\\[1cm]
\textbf{Serie completa:}\\
Volumen 1: Fundamentos Te\'oricos (Caps. 1--5)\\
Volumen 2: Aplicaciones y Privacidad (Caps. 6--13)\\
Volumen 3: Herramientas Avanzadas y Ecosistema (Caps. 14--20)\\[1cm]
El c\'odigo incluido es de dominio p\'ublico y puede ser utilizado libremente.\\
Este libro es para uso personal y educativo.
\end{center}
\vfill
\cleardoublepage
% --- BLOQUES ---
\thispagestyle{empty}
\null\vfill
\begin{center}
{\Large\bfseries Los 20 Bloques Temáticos}\\[1cm]

\begin{tabular}{cl}
\toprule
\textbf{Bloque} & \textbf{Tema} \\
\midrule
\textbf{1} & \textbf{Fundamentos Criptográficos} \\
\textbf{2} & \textbf{Blockchain Base (L1)} \\
\textbf{3} & \textbf{Escalabilidad (L2/L3)} \\
\textbf{4} & \textbf{Zero-Knowledge Proofs (Core)} \\
\textbf{5} & \textbf{Circuitos ZK} \\
\midrule
6 & Compromisos y Estructuras \\
7 & Infraestructura ZK \\
8 & Privacidad Avanzada \\
9 & DeFi y Aplicaciones \\
10 & Avanzado/Investigación \\
11 & Protocolos de Anonimización \\
12 & Monedas Privadas (Privacy Coins) \\
13 & Técnicas de Unlinkability \\
14 & Identidad y Credenciales Anónimas \\
15 & Compliance vs Privacidad \\
16 & Metadata y Network Privacy \\
17 & Private Smart Contracts \\
18 & Criptografía Avanzada para Privacidad \\
19 & Ataques y Defensa \\
20 & Herramientas y Ecosistema \\
\bottomrule
\end{tabular}

\vspace{1cm}
\textit{En negrita: capítulos incluidos en este volumen}
\end{center}
\vfill
\cleardoublepage

% --- ÍNDICE ---
\tableofcontents
\cleardoublepage

% ============================================
\mainmatter
% ============================================

% ============================================
% BLOQUE 1: FUNDAMENTOS CRIPTOGRAFICOS
% ============================================
\chapter{Fundamentos Criptograficos}
\label{ch:bloque1}

\begin{center}
\textit{Bloque 1 de 20}
\end{center}

% ----- HASH FUNCTIONS -----
\section{Hash Functions}
% ========================  SHA 

%===========================================
% SUBSECCIÓN: SHA-256 - VERSIÓN EXPANDIDA COMPLETA
%===========================================

\subsection{SHA-256}

\begin{definition}[SHA-256]
Función hash de la familia SHA-2 que produce salida de 256 bits.
\end{definition}

Características:
\begin{itemize}
    \item Construcción Merkle-Damgård
    \item Bloques de 512 bits
    \item 64 rondas por bloque
    \item Operaciones: rotaciones, shifts, XOR, sumas mod $2^{32}$
\end{itemize}

Usado en: Bitcoin (PoW, merkle roots), certificados SSL, integridad de archivos.

%-------------------------------------------
\subsubsection{Contexto Histórico}

\begin{intuicion}
SHA significa \textit{Secure Hash Algorithm}. La familia SHA fue desarrollada por la NSA y publicada por NIST. SHA-256 es parte de SHA-2, publicado en 2001 como sucesor de SHA-1 (que hoy se considera inseguro tras el ataque de colisión de Google en 2017).
\end{intuicion}

SHA-256 se convirtió en el estándar de facto para criptomonedas cuando Satoshi Nakamoto lo eligió para Bitcoin en 2008. Esta elección no fue accidental: SHA-256 ofrecía el mejor balance entre seguridad probada, eficiencia en hardware, y resistencia a ataques conocidos.

\begin{clave}
256 bits de salida significa $2^{256}$ posibles hashes. Este número es astronómicamente grande:
\[ 2^{256} \approx 1.16 \times 10^{77} \]
Para contexto: hay aproximadamente $10^{80}$ átomos en el universo observable. Un ataque de fuerza bruta es físicamente imposible.
\end{clave}

% ============================================
% SECCIÓN 1.2: ÁLGEBRA PARA CRIPTO
% Insertar después del contenido de SHA-256
% ============================================

\section{Álgebra para Cripto (Campos y Polinomios)}

La criptografía moderna, especialmente Zero-Knowledge Proofs, se construye sobre estructuras algebraicas precisas. Esta sección establece los fundamentos matemáticos que aparecerán constantemente en SNARKs, STARKs, KZG, y sistemas como PLONK.

\begin{intuicion}
¿Por qué álgebra abstracta en criptografía? Porque necesitamos operaciones que:
\begin{enumerate}
    \item Sean \textbf{eficientes} de calcular
    \item Tengan \textbf{estructura predecible} (asociatividad, inversos)
    \item Presenten \textbf{problemas difíciles} para atacantes
\end{enumerate}
Los campos finitos y grupos cíclicos satisfacen perfectamente estos requisitos.
\end{intuicion}

% ============================================
\subsection{Campos Finitos $\mathbb{F}_p$}
% ============================================

\begin{definition}[Campo Finito Primo]
Un \textbf{campo finito} $\mathbb{F}_p$ (donde $p$ es primo) es el conjunto $\{0, 1, 2, \ldots, p-1\}$ con dos operaciones:
\begin{itemize}
    \item \textbf{Suma}: $a + b \mod p$
    \item \textbf{Multiplicación}: $a \cdot b \mod p$
\end{itemize}
que satisfacen los axiomas de campo (asociatividad, conmutatividad, distributividad, existencia de neutros e inversos).
\end{definition}

\begin{example}[Campo $\mathbb{F}_7$]
En $\mathbb{F}_7 = \{0, 1, 2, 3, 4, 5, 6\}$:
\begin{align*}
    3 + 5 &= 8 \mod 7 = 1 \\
    3 \cdot 5 &= 15 \mod 7 = 1 \\
    4 + 6 &= 10 \mod 7 = 3
\end{align*}
Observa que $3 \cdot 5 = 1$, lo que significa que $5$ es el \textbf{inverso multiplicativo} de $3$.
\end{example}

\subsubsection{Inversos Multiplicativos}

\begin{definition}[Inverso Multiplicativo]
Para $a \in \mathbb{F}_p$ con $a \neq 0$, el \textbf{inverso multiplicativo} $a^{-1}$ es el único elemento tal que:
\[ a \cdot a^{-1} \equiv 1 \pmod{p} \]
\end{definition}

\begin{mathnote}
El inverso existe para todo $a \neq 0$ porque $p$ es primo (por el teorema de Bézout: $\gcd(a, p) = 1$ implica que existen $x, y$ con $ax + py = 1$, entonces $ax \equiv 1 \pmod{p}$).
\end{mathnote}

Hay dos métodos principales para calcular inversos:

\textbf{Método 1: Algoritmo de Euclides Extendido}
\begin{verbatim}
def extended_gcd(a, b):
    if a == 0:
        return b, 0, 1
    gcd, x1, y1 = extended_gcd(b % a, a)
    x = y1 - (b // a) * x1
    y = x1
    return gcd, x, y

def mod_inverse(a, p):
    _, x, _ = extended_gcd(a % p, p)
    return x % p
\end{verbatim}

\textbf{Método 2: Pequeño Teorema de Fermat}
\begin{theorem}[Pequeño Teorema de Fermat]
Si $p$ es primo y $a \not\equiv 0 \pmod{p}$, entonces:
\[ a^{p-1} \equiv 1 \pmod{p} \]
Por lo tanto:
\[ a^{-1} \equiv a^{p-2} \pmod{p} \]
\end{theorem}

\begin{example}[Calculando $3^{-1}$ en $\mathbb{F}_7$]
Usando Fermat: $3^{-1} = 3^{7-2} = 3^5 \mod 7$
\begin{align*}
    3^2 &= 9 \mod 7 = 2 \\
    3^4 &= 2^2 = 4 \\
    3^5 &= 3^4 \cdot 3 = 4 \cdot 3 = 12 \mod 7 = 5
\end{align*}
Verificación: $3 \cdot 5 = 15 \mod 7 = 1$ \checkmark
\end{example}

\subsubsection{Orden de Elementos}

\begin{definition}[Orden de un Elemento]
El \textbf{orden} de $a \in \mathbb{F}_p^*$ (donde $\mathbb{F}_p^* = \mathbb{F}_p \setminus \{0\}$) es el menor entero positivo $k$ tal que:
\[ a^k \equiv 1 \pmod{p} \]
Se denota $\text{ord}(a) = k$.
\end{definition}

\begin{theorem}[Lagrange]
El orden de cualquier elemento divide a $|\mathbb{F}_p^*| = p - 1$.
\end{theorem}

\begin{definition}[Generador / Raíz Primitiva]
Un elemento $g \in \mathbb{F}_p^*$ es un \textbf{generador} (o raíz primitiva) si $\text{ord}(g) = p - 1$. Esto significa que:
\[ \mathbb{F}_p^* = \{g^0, g^1, g^2, \ldots, g^{p-2}\} = \{1, g, g^2, \ldots, g^{p-2}\} \]
\end{definition}

\begin{example}[Generador en $\mathbb{F}_7$]
Probemos si $g = 3$ es generador:
\begin{align*}
    3^1 &= 3 \\
    3^2 &= 9 \mod 7 = 2 \\
    3^3 &= 6 \\
    3^4 &= 18 \mod 7 = 4 \\
    3^5 &= 12 \mod 7 = 5 \\
    3^6 &= 15 \mod 7 = 1
\end{align*}
Las potencias generan $\{1, 2, 3, 4, 5, 6\} = \mathbb{F}_7^*$. Por tanto, $3$ es generador con orden $6 = p-1$.
\end{example}

\begin{bancozk}
En Banco ZK, los campos finitos que usamos tienen primos especiales. Por ejemplo, el campo de BN254:
\[ p = 21888242871839275222246405745257275088696311157297823662689037894645226208583 \]
Este primo de 254 bits fue elegido para que las curvas elípticas sobre $\mathbb{F}_p$ tengan propiedades específicas para pairings.
\end{bancozk}

% ============================================
\subsection{Extensiones de Campo $\mathbb{F}_{p^k}$}
% ============================================

\begin{intuicion}
Así como extendemos $\mathbb{R}$ a $\mathbb{C}$ añadiendo $i = \sqrt{-1}$, podemos extender $\mathbb{F}_p$ añadiendo raíces de polinomios irreducibles. Esto nos da campos más grandes con estructura rica, esenciales para pairings.
\end{intuicion}

\begin{definition}[Extensión de Campo]
Sea $f(x) \in \mathbb{F}_p[x]$ un polinomio irreducible de grado $k$. La \textbf{extensión} $\mathbb{F}_{p^k}$ se construye como:
\[ \mathbb{F}_{p^k} = \mathbb{F}_p[x] / (f(x)) \]
Los elementos son polinomios de grado $< k$ con coeficientes en $\mathbb{F}_p$.
\end{definition}

\subsubsection{Extensión Cuadrática $\mathbb{F}_{p^2}$}

\begin{example}[Construcción de $\mathbb{F}_{p^2}$]
Si $-1$ no es un cuadrado en $\mathbb{F}_p$ (i.e., $p \equiv 3 \pmod 4$), podemos usar $f(x) = x^2 + 1$:
\[ \mathbb{F}_{p^2} = \mathbb{F}_p[i] \quad \text{donde } i^2 = -1 \]
Los elementos tienen la forma $a + bi$ con $a, b \in \mathbb{F}_p$.
\end{example}

\textbf{Aritmética en $\mathbb{F}_{p^2}$:}
\begin{align*}
    (a + bi) + (c + di) &= (a+c) + (b+d)i \\
    (a + bi) \cdot (c + di) &= (ac - bd) + (ad + bc)i
\end{align*}

\begin{mathnote}
La multiplicación requiere 4 multiplicaciones en $\mathbb{F}_p$. Con el truco de Karatsuba se reduce a 3:
\begin{align*}
    ac - bd &= ac - bd \\
    ad + bc &= (a+b)(c+d) - ac - bd
\end{align*}
Calculamos $ac$, $bd$, $(a+b)(c+d)$: solo 3 multiplicaciones.
\end{mathnote}

\subsubsection{Extensión $\mathbb{F}_{p^{12}}$ para Pairings}

\begin{intuicion}
Los pairings en curvas BN y BLS requieren $\mathbb{F}_{p^{12}}$. ¿Por qué 12? Es el \textit{embedding degree}: el menor $k$ tal que $r | p^k - 1$, donde $r$ es el orden del subgrupo de la curva.
\end{intuicion}

La construcción típica usa una \textbf{torre de extensiones}:
\[ \mathbb{F}_p \to \mathbb{F}_{p^2} \to \mathbb{F}_{p^6} \to \mathbb{F}_{p^{12}} \]

\begin{example}[Torre para BN254]
\begin{align*}
    \mathbb{F}_{p^2} &= \mathbb{F}_p[u] / (u^2 + 1) \\
    \mathbb{F}_{p^6} &= \mathbb{F}_{p^2}[v] / (v^3 - (9 + u)) \\
    \mathbb{F}_{p^{12}} &= \mathbb{F}_{p^6}[w] / (w^2 - v)
\end{align*}
Un elemento de $\mathbb{F}_{p^{12}}$ se representa con 12 elementos de $\mathbb{F}_p$.
\end{example}

\begin{advertencia}
La aritmética en $\mathbb{F}_{p^{12}}$ es costosa: una multiplicación requiere $\sim 54$ multiplicaciones en $\mathbb{F}_p$ (con optimizaciones). Por eso los pairings son operaciones ``caras'' comparadas con operaciones en curvas.
\end{advertencia}

\begin{keyidea}
\textbf{Resumen de extensiones:}
\begin{itemize}
    \item $|\mathbb{F}_{p^k}| = p^k$ elementos
    \item $\mathbb{F}_{p^k}^*$ es cíclico de orden $p^k - 1$
    \item Toda extensión finita de $\mathbb{F}_p$ tiene la forma $\mathbb{F}_{p^k}$
    \item Pairings mapean $\mathbb{G}_1 \times \mathbb{G}_2 \to \mathbb{G}_T \subset \mathbb{F}_{p^{12}}^*$
\end{itemize}
\end{keyidea}

% ============================================
\subsection{Polinomios: Interpolación y Evaluación}
% ============================================

Los polinomios son la herramienta fundamental en ZK-SNARKs y STARKs. Permiten codificar computaciones completas en objetos matemáticos compactos.

\begin{definition}[Polinomio sobre $\mathbb{F}_p$]
Un polinomio de grado $\leq n$ sobre $\mathbb{F}_p$ es:
\[ f(x) = a_0 + a_1 x + a_2 x^2 + \cdots + a_n x^n, \quad a_i \in \mathbb{F}_p \]
El conjunto de todos estos polinomios se denota $\mathbb{F}_p[x]_{\leq n}$.
\end{definition}

\subsubsection{Evaluación de Polinomios}

\begin{definition}[Evaluación]
Dado $f(x) \in \mathbb{F}_p[x]$ y un punto $\alpha \in \mathbb{F}_p$, la \textbf{evaluación} $f(\alpha)$ sustituye $x$ por $\alpha$.
\end{definition}

\textbf{Método de Horner} (evaluación eficiente en $O(n)$):
\[ f(x) = a_0 + x(a_1 + x(a_2 + \cdots + x(a_{n-1} + x \cdot a_n))) \]

\begin{verbatim}
def horner_eval(coeffs, x, p):
    """Evalúa polinomio en x módulo p usando Horner"""
    result = 0
    for coeff in reversed(coeffs):
        result = (result * x + coeff) % p
    return result
\end{verbatim}

\subsubsection{Interpolación de Lagrange}

\begin{theorem}[Interpolación de Lagrange]
Dados $n+1$ puntos distintos $(x_0, y_0), (x_1, y_1), \ldots, (x_n, y_n)$ con $x_i \in \mathbb{F}_p$ distintos, existe un \textbf{único} polinomio $f(x)$ de grado $\leq n$ tal que $f(x_i) = y_i$ para todo $i$.
\end{theorem}

\begin{definition}[Polinomios Base de Lagrange]
Para un conjunto de puntos $\{x_0, \ldots, x_n\}$, el $i$-ésimo \textbf{polinomio base de Lagrange} es:
\[ L_i(x) = \prod_{j \neq i} \frac{x - x_j}{x_i - x_j} \]
Satisface: $L_i(x_j) = \begin{cases} 1 & \text{si } i = j \\ 0 & \text{si } i \neq j \end{cases}$
\end{definition}

El polinomio interpolante es:
\[ f(x) = \sum_{i=0}^{n} y_i \cdot L_i(x) \]

\begin{example}[Interpolación en $\mathbb{F}_{11}$]
Interpolar los puntos $(1, 3), (2, 7), (4, 5)$:

Bases de Lagrange:
\begin{align*}
    L_0(x) &= \frac{(x-2)(x-4)}{(1-2)(1-4)} = \frac{(x-2)(x-4)}{(-1)(-3)} = \frac{(x-2)(x-4)}{3} \\
    L_1(x) &= \frac{(x-1)(x-4)}{(2-1)(2-4)} = \frac{(x-1)(x-4)}{-2} \\
    L_2(x) &= \frac{(x-1)(x-2)}{(4-1)(4-2)} = \frac{(x-1)(x-2)}{6}
\end{align*}

En $\mathbb{F}_{11}$: $3^{-1} = 4$, $(-2)^{-1} = 9^{-1} = 5$, $6^{-1} = 2$

\[ f(x) = 3 \cdot 4(x-2)(x-4) + 7 \cdot 5(x-1)(x-4) + 5 \cdot 2(x-1)(x-2) \]

Expandiendo y simplificando módulo 11: $f(x) = 2x^2 + 9x + 3$

Verificación: $f(1) = 2 + 9 + 3 = 14 \equiv 3$, $f(2) = 8 + 18 + 3 = 29 \equiv 7$, $f(4) = 32 + 36 + 3 = 71 \equiv 5$ \checkmark
\end{example}

\begin{keyidea}
\textbf{Dualidad evaluación-coeficientes:}
\begin{itemize}
    \item Un polinomio de grado $n$ se representa con $n+1$ coeficientes
    \item Alternativamente, con $n+1$ evaluaciones en puntos distintos
    \item Lagrange convierte evaluaciones $\to$ coeficientes
    \item Horner convierte coeficientes $\to$ evaluación
\end{itemize}
Esta dualidad es fundamental en SNARKs: el prover trabaja con coeficientes, el verifier con evaluaciones.
\end{keyidea}

\subsubsection{Polinomio de Desvanecimiento (Vanishing Polynomial)}

\begin{definition}[Vanishing Polynomial]
Para un conjunto $H = \{h_1, \ldots, h_n\} \subset \mathbb{F}_p$, el \textbf{polinomio de desvanecimiento} es:
\[ Z_H(x) = \prod_{i=1}^{n} (x - h_i) \]
Satisface: $Z_H(h) = 0$ para todo $h \in H$.
\end{definition}

\begin{bancozk}
En PLONK, $H$ es el dominio de evaluación (típicamente raíces de la unidad). Si $|H| = n$ y $\omega$ es raíz $n$-ésima primitiva:
\[ Z_H(x) = x^n - 1 \]
¡Un polinomio de grado $n$ con solo 2 coeficientes no nulos! Esto hace que dividir por $Z_H$ sea extremadamente eficiente.
\end{bancozk}

% ============================================
\subsection{FFT/NTT: Transformada Rápida}
% ============================================

\begin{intuicion}
La interpolación de Lagrange ingenua cuesta $O(n^2)$. Para circuitos con millones de gates, esto es prohibitivo. La FFT/NTT reduce el costo a $O(n \log n)$, haciendo posibles los SNARKs prácticos.
\end{intuicion}

\subsubsection{Raíces de la Unidad}

\begin{definition}[Raíz $n$-ésima de la Unidad]
En $\mathbb{F}_p$, una \textbf{raíz $n$-ésima primitiva de la unidad} es un elemento $\omega$ tal que:
\begin{enumerate}
    \item $\omega^n = 1$
    \item $\omega^k \neq 1$ para $0 < k < n$
\end{enumerate}
Las $n$ raíces $n$-ésimas son $\{1, \omega, \omega^2, \ldots, \omega^{n-1}\}$.
\end{definition}

\begin{theorem}[Existencia de Raíces]
$\mathbb{F}_p$ contiene una raíz $n$-ésima primitiva si y solo si $n | (p-1)$.
\end{theorem}

\begin{example}[Raíces en $\mathbb{F}_{17}$]
Como $17 - 1 = 16 = 2^4$, existen raíces $n$-ésimas para $n \in \{1, 2, 4, 8, 16\}$.

Para $n = 4$: buscamos $\omega$ con $\omega^4 = 1$ pero $\omega^2 \neq 1$.
\begin{itemize}
    \item $2^4 = 16 \equiv -1 \pmod{17}$, así que $2^8 \equiv 1$. Entonces $\omega = 2^2 = 4$ tiene orden 4.
\end{itemize}
Verificación: $4^1 = 4$, $4^2 = 16 \equiv -1$, $4^3 = -4 \equiv 13$, $4^4 = 1$ \checkmark

Las 4 raíces cuartas son: $\{1, 4, 16, 13\} = \{4^0, 4^1, 4^2, 4^3\}$.
\end{example}

\subsubsection{Number Theoretic Transform (NTT)}

La NTT es la versión de la FFT sobre campos finitos.

\begin{definition}[NTT]
Sea $\omega$ raíz $n$-ésima primitiva en $\mathbb{F}_p$. La \textbf{NTT} de un vector $(a_0, \ldots, a_{n-1})$ es:
\[ \hat{a}_k = \sum_{j=0}^{n-1} a_j \cdot \omega^{jk} \mod p \]
Equivalentemente, si $f(x) = \sum_{j=0}^{n-1} a_j x^j$, entonces:
\[ \text{NTT}(a) = (f(1), f(\omega), f(\omega^2), \ldots, f(\omega^{n-1})) \]
\end{definition}

\begin{keyidea}
\textbf{La NTT convierte coeficientes en evaluaciones sobre raíces de la unidad.}
\begin{itemize}
    \item NTT: coeficientes $\to$ evaluaciones en $\{1, \omega, \omega^2, \ldots\}$
    \item INTT: evaluaciones $\to$ coeficientes
\end{itemize}
\end{keyidea}

\subsubsection{Algoritmo Cooley-Tukey}

La eficiencia viene de la estructura recursiva. Si $n = 2m$:

\begin{align*}
    f(x) &= f_{\text{even}}(x^2) + x \cdot f_{\text{odd}}(x^2)
\end{align*}

donde $f_{\text{even}}$ tiene los coeficientes de índice par y $f_{\text{odd}}$ los impares.

\begin{verbatim}
def ntt(a, omega, p):
    """NTT recursiva - O(n log n)"""
    n = len(a)
    if n == 1:
        return a
    
    # Dividir en pares e impares
    a_even = ntt(a[0::2], pow(omega, 2, p), p)
    a_odd = ntt(a[1::2], pow(omega, 2, p), p)
    
    # Combinar (butterfly)
    result = [0] * n
    omega_k = 1
    for k in range(n // 2):
        result[k] = (a_even[k] + omega_k * a_odd[k]) % p
        result[k + n//2] = (a_even[k] - omega_k * a_odd[k]) % p
        omega_k = (omega_k * omega) % p
    
    return result
\end{verbatim}

\begin{mathnote}
\textbf{Complejidad:}
\begin{itemize}
    \item Evaluación ingenua en $n$ puntos: $O(n^2)$
    \item NTT: $O(n \log n)$
    \item Para $n = 2^{20}$ (un millón de gates): 
    \begin{itemize}
        \item Ingenuo: $10^{12}$ operaciones
        \item NTT: $20 \times 10^6$ operaciones
        \item ¡Speedup de $50,000\times$!
    \end{itemize}
\end{itemize}
\end{mathnote}

\subsubsection{Multiplicación de Polinomios con NTT}

\begin{theorem}[Convolución vía NTT]
Para multiplicar dos polinomios $f, g$ de grado $< n$:
\begin{enumerate}
    \item Pad a longitud $2n$
    \item $\hat{f} = \text{NTT}(f)$, $\hat{g} = \text{NTT}(g)$
    \item $\hat{h}_k = \hat{f}_k \cdot \hat{g}_k$ (multiplicación punto a punto)
    \item $h = \text{INTT}(\hat{h})$
\end{enumerate}
Costo total: $O(n \log n)$ vs $O(n^2)$ ingenuo.
\end{theorem}

\begin{bancozk}
\textbf{¿Por qué importa en SNARKs/STARKs?}

\textbf{PLONK}: Necesita multiplicar y dividir polinomios de grado $\sim n$ (número de gates). Sin NTT, generar una prueba para un circuito de 1 millón de gates tomaría horas. Con NTT, toma segundos.

\textbf{STARKs}: Usan FRI (Fast Reed-Solomon IOP) que requiere evaluar polinomios en dominios exponencialmente más grandes. La NTT hace esto factible.

\textbf{KZG Commitments}: El trusted setup genera evaluaciones de $\tau^i$ que se combinan con NTT para commits eficientes.
\end{bancozk}

\begin{advertencia}
Para usar NTT necesitas que $n | (p-1)$ donde $n$ es potencia de 2. Los primos ``amigables con SNARK'' se eligen específicamente para esto. Ejemplos:
\begin{itemize}
    \item BN254: $p - 1 = 2^{28} \cdot \ldots$ (soporta NTT hasta $2^{28}$)
    \item Goldilocks: $p = 2^{64} - 2^{32} + 1$, $p - 1 = 2^{32} \cdot \ldots$
    \item BLS12-381: $p - 1 = 2^{32} \cdot \ldots$
\end{itemize}
\end{advertencia}

% ============================================
\subsection{Grupos Cíclicos y Supuestos de Seguridad}
% ============================================

Esta sección formaliza la notación de grupos que usaremos constantemente y establece los problemas computacionalmente difíciles que fundamentan la seguridad criptográfica.

\subsubsection{Grupos Cíclicos}

\begin{definition}[Grupo Cíclico]
Un grupo $\mathbb{G}$ es \textbf{cíclico} si existe un elemento $g \in \mathbb{G}$ (llamado \textbf{generador}) tal que:
\[ \mathbb{G} = \langle g \rangle = \{g^0, g^1, g^2, \ldots, g^{n-1}\} \]
donde $n = |\mathbb{G}|$ es el \textbf{orden} del grupo.
\end{definition}

\begin{remark}[Notación]
Usamos dos notaciones según el contexto:
\begin{itemize}
    \item \textbf{Multiplicativa}: $g^x$ (típica en $\mathbb{F}_p^*$, grupos abstractos)
    \item \textbf{Aditiva}: $[x]G = xG = G + G + \cdots + G$ ($x$ veces) (típica en curvas elípticas)
\end{itemize}
Son equivalentes: $g^x \leftrightarrow [x]G$.
\end{remark}

\begin{example}[Grupos Cíclicos Comunes]
\begin{enumerate}
    \item $\mathbb{Z}_n = (\{0, 1, \ldots, n-1\}, +)$: grupo aditivo, generador $1$
    \item $\mathbb{F}_p^* = (\{1, 2, \ldots, p-1\}, \cdot)$: grupo multiplicativo, generador = raíz primitiva
    \item $E(\mathbb{F}_p)$: puntos de curva elíptica (bajo ciertas condiciones)
\end{enumerate}
\end{example}

\subsubsection{El Problema del Logaritmo Discreto (DLOG)}

\begin{definition}[Problema DLOG]
Sea $\mathbb{G} = \langle g \rangle$ un grupo cíclico de orden $n$. Dado $h \in \mathbb{G}$, el \textbf{problema del logaritmo discreto} es encontrar $x \in \mathbb{Z}_n$ tal que:
\[ g^x = h \]
El valor $x$ se llama el \textbf{logaritmo discreto} de $h$ en base $g$, denotado $x = \log_g h$.
\end{definition}

\begin{intuicion}
Calcular $g^x$ dado $g$ y $x$ es fácil (exponenciación rápida: $O(\log x)$ multiplicaciones). Pero encontrar $x$ dado $g$ y $g^x$ es difícil en ciertos grupos. Esta \textit{asimetría} es la base de la criptografía de clave pública.
\end{intuicion}

\begin{definition}[Supuesto DLOG]
El \textbf{supuesto DLOG} afirma que para todo adversario probabilístico de tiempo polinomial $\mathcal{A}$:
\[ \Pr\left[ \mathcal{A}(g, g^x) = x \;:\; x \leftarrow \mathbb{Z}_n \right] \leq \mathsf{negl}(\lambda) \]
donde $\lambda$ es el parámetro de seguridad y $\mathsf{negl}(\lambda)$ es una función negligible.
\end{definition}

\begin{advertencia}
DLOG \textbf{no} es difícil en todos los grupos:
\begin{itemize}
    \item En $(\mathbb{Z}_n, +)$: $\log_1 x = x$ (trivial)
    \item En $\mathbb{F}_p^*$ con $p$ pequeño: algoritmo de Pohlig-Hellman
    \item En $\mathbb{F}_p^*$ con $p$ grande pero $(p-1)$ suave: también vulnerable
\end{itemize}
Necesitamos grupos donde el mejor algoritmo conocido sea exponencial.
\end{advertencia}

\subsubsection{Supuestos Relacionados}

\begin{definition}[Computational Diffie-Hellman (CDH)]
Dados $g$, $g^a$, $g^b$ (con $a, b$ aleatorios), es difícil computar $g^{ab}$.
\end{definition}

\begin{definition}[Decisional Diffie-Hellman (DDH)]
Es difícil distinguir $(g, g^a, g^b, g^{ab})$ de $(g, g^a, g^b, g^c)$ donde $c$ es aleatorio.
\end{definition}

\begin{theorem}[Relación entre supuestos]
\[ \text{DLOG} \implies \text{CDH} \implies \text{DDH} \]
(Si puedes resolver DLOG, puedes resolver CDH; si puedes resolver CDH, puedes resolver DDH.)
\end{theorem}

\begin{definition}[Knowledge of Exponent Assumption (KEA)]
Si un algoritmo $\mathcal{A}$ produce $(C, \hat{C})$ tales que $\hat{C} = C^\alpha$ (donde $\alpha$ es desconocido para $\mathcal{A}$), entonces $\mathcal{A}$ ``conoce'' $c$ tal que $C = g^c$.

Formalmente: existe un extractor $\mathcal{E}$ que, con acceso a $\mathcal{A}$, puede extraer $c$.
\end{definition}

\begin{keyidea}
\textbf{KEA es fundamental para SNARKs:} permite argumentar que el prover ``conoce'' el witness, no solo que produjo una prueba válida. Groth16 y otros sistemas usan variantes de KEA.
\end{keyidea}

\subsubsection{Tabla de Grupos y Seguridad}

\begin{center}
\begin{tabular}{lccl}
\toprule
\textbf{Grupo} & \textbf{DLOG} & \textbf{DDH} & \textbf{Mejor ataque conocido} \\
\midrule
$\mathbb{F}_p^*$ & Difícil & Difícil & Index calculus: $L_p[1/3]$ \\
$E(\mathbb{F}_p)$ (curva genérica) & Difícil & Difícil & Pollard-rho: $O(\sqrt{n})$ \\
$E(\mathbb{F}_p)$ (con pairing) & Difícil & \textbf{Fácil} & Pairing $\to \mathbb{F}_{p^k}^*$ \\
\bottomrule
\end{tabular}
\end{center}

\begin{mathnote}
En curvas con pairing, DDH es fácil porque:
\[ e(g^a, g^b) = e(g, g)^{ab} \]
Podemos verificar si $g^c = g^{ab}$ chequeando $e(g, g^c) \stackrel{?}{=} e(g^a, g^b)$.

Esto no rompe DLOG ni CDH, pero significa que debemos usar otros supuestos (como co-CDH o q-SDH) para seguridad en sistemas basados en pairings.
\end{mathnote}

\subsubsection{Notación Estándar}

Resumimos la notación que usaremos en el resto del libro:

\begin{center}
\begin{tabular}{ll}
\toprule
\textbf{Símbolo} & \textbf{Significado} \\
\midrule
$\mathbb{G}$ & Grupo cíclico genérico \\
$g, h$ & Generadores del grupo \\
$n$ u $|\mathbb{G}|$ & Orden del grupo \\
$\mathbb{Z}_n$ & Exponentes (enteros módulo $n$) \\
$g^x$ o $[x]g$ & Exponenciación / multiplicación escalar \\
$\mathbb{G}_1, \mathbb{G}_2$ & Grupos fuente de pairing \\
$\mathbb{G}_T$ & Grupo objetivo de pairing \\
$e: \mathbb{G}_1 \times \mathbb{G}_2 \to \mathbb{G}_T$ & Pairing bilineal \\
$\mathsf{negl}(\lambda)$ & Función negligible en parámetro de seguridad \\
$x \leftarrow S$ & $x$ muestreado uniformemente de $S$ \\
$x \leftarrow \mathcal{A}(\cdot)$ & $x$ es output del algoritmo $\mathcal{A}$ \\
\bottomrule
\end{tabular}
\end{center}

\begin{bancozk}
Con esta base algebraica, estamos preparados para entender:
\begin{itemize}
    \item \textbf{Pedersen Commitments}: $C = g^v h^r$ (hiding por DLOG)
    \item \textbf{KZG Commitments}: evaluaciones en $\tau$ secreto, verificación con pairing
    \item \textbf{PLONK}: polinomios sobre $\mathbb{F}_p$, NTT para eficiencia, KZG para commits
    \item \textbf{STARKs}: campos grandes, FRI usando NTT extensiva
    \item \textbf{Groth16}: pairings para verificación en 3 operaciones
\end{itemize}
Todo se construye sobre los conceptos de esta sección.
\end{bancozk}

% ============================================
% FIN DE SECCIÓN 1.2
% ============================================

%-------------------------------------------
\subsubsection{La Construcción Merkle-Damgård}

SHA-256 utiliza la construcción \textbf{Merkle-Damgård}, propuesta independientemente por Ralph Merkle e Ivan Damgård en 1989. Este descubrimiento fue fundamental porque resolvió cómo construir una función hash segura para mensajes de cualquier longitud a partir de una función de compresión que solo procesa bloques de tamaño fijo.

\begin{definition}[Función de Compresión]
Una \textbf{función de compresión} es una función determinística:
\[ f: \{0,1\}^n \times \{0,1\}^b \to \{0,1\}^n \]
que toma dos entradas:
\begin{itemize}
    \item Un estado de $n$ bits (llamado \textit{chaining value})
    \item Un bloque de mensaje de $b$ bits
\end{itemize}
y produce una salida de $n$ bits (el nuevo estado).
\end{definition}

\begin{definition}[Construcción Merkle-Damgård]
Sea $f: \{0,1\}^n \times \{0,1\}^b \to \{0,1\}^n$ una función de compresión resistente a colisiones. La construcción MD define una función hash $H: \{0,1\}^* \to \{0,1\}^n$ como sigue:

\textbf{Entrada:} Mensaje $M$ de longitud arbitraria

\textbf{Proceso:}
\begin{enumerate}
    \item \textbf{Padding:} Extender $M$ a longitud múltiplo de $b$ bits:
    \[ M' = M \,||\, 1 \,||\, 0^k \,||\, \langle |M| \rangle_{64} \]
    donde $k$ es el mínimo tal que $|M| + 1 + k + 64 \equiv 0 \pmod{b}$
    
    \item \textbf{División:} Particionar $M'$ en bloques de $b$ bits:
    \[ M' = M_1 \,||\, M_2 \,||\, \cdots \,||\, M_L \]
    
    \item \textbf{Iteración:} Comenzando con un valor inicial fijo $H_0 = IV$:
    \begin{align}
        H_1 &= f(H_0, M_1) = f(IV, M_1) \\
        H_2 &= f(H_1, M_2) \\
        &\vdots \\
        H_L &= f(H_{L-1}, M_L)
    \end{align}
    
    \item \textbf{Salida:} $H(M) = H_L$
\end{enumerate}
\end{definition}

\begin{center}
\begin{tikzpicture}[
    block/.style={draw, minimum width=1.2cm, minimum height=1cm, align=center},
    arrow/.style={->, thick},
    node distance=1.8cm
]
    % IV
    \node (iv) {$IV$};
    
    % Funciones de compresión
    \node[block, right=of iv] (f1) {$f$};
    \node[block, right=of f1] (f2) {$f$};
    \node[right=of f2] (dots) {$\cdots$};
    \node[block, right=of dots] (fL) {$f$};
    
    % Salida
    \node[right=of fL] (out) {$H(M)$};
    
    % Bloques de mensaje (arriba)
    \node[above=0.8cm of f1] (m1) {$M_1$};
    \node[above=0.8cm of f2] (m2) {$M_2$};
    \node[above=0.8cm of fL] (mL) {$M_L$};
    
    % Flechas horizontales
    \draw[arrow] (iv) -- (f1);
    \draw[arrow] (f1) -- node[above] {\small $H_1$} (f2);
    \draw[arrow] (f2) -- (dots);
    \draw[arrow] (dots) -- (fL);
    \draw[arrow] (fL) -- (out);
    
    % Flechas de mensajes
    \draw[arrow] (m1) -- (f1);
    \draw[arrow] (m2) -- (f2);
    \draw[arrow] (mL) -- (fL);
\end{tikzpicture}
\end{center}

\begin{intuicion}
\textbf{Analogía de la trituradora de papel:}

Imagina una trituradora de documentos especial:
\begin{enumerate}
    \item Comienzas con una ``semilla'' específica en la máquina ($IV$)
    \item Introduces la primera página del documento ($M_1$)
    \item La máquina mezcla la semilla con la página y produce un nuevo estado interno ($H_1$)
    \item Introduces la segunda página ($M_2$), se mezcla con $H_1$ produciendo $H_2$
    \item Continúas hasta procesar todas las páginas
    \item El estado final de la máquina es tu ``hash''
\end{enumerate}

Lo crucial es que cada página afecta todas las siguientes (efecto cascada), no puedes retroceder, y cambiar cualquier página cambia completamente el resultado final.
\end{intuicion}

\begin{ejemplo}[Merkle-Damgård con mensaje ``Hola'']
Usaremos parámetros simplificados para ilustrar:

\textbf{Paso 1: Mensaje original}
\begin{verbatim}
"Hola" en ASCII:
H = 0x48, o = 0x6F, l = 0x6C, a = 0x61
Mensaje M = 32 bits
\end{verbatim}

\textbf{Paso 2: Calcular padding (para bloques de 512 bits)}
\begin{align}
\text{Longitud} &= 32 \text{ bits} \\
\text{Padding} &= 1 \,||\, 0^{423} \,||\, \langle 32 \rangle_{64}
\end{align}

\textbf{Paso 3: Aplicar función de compresión}
\begin{align}
H_1 &= f(IV, M_1) = f(\texttt{0x6a09e667...}, \texttt{[bloque con padding]})
\end{align}

Como solo hay un bloque: $H(\text{``Hola''}) = H_1$
\end{ejemplo}

%-------------------------------------------
\subsubsection{Parámetros Específicos de SHA-256}

\begin{center}
\begin{tabular}{ll}
\toprule
\textbf{Parámetro} & \textbf{Valor en SHA-256} \\
\midrule
Tamaño de salida & 256 bits \\
Tamaño de bloque & 512 bits \\
Tamaño de palabra & 32 bits \\
Número de rondas & 64 \\
Tamaño del estado interno & 256 bits (8 palabras de 32 bits) \\
Campo de longitud & 64 bits \\
\bottomrule
\end{tabular}
\end{center}

%-------------------------------------------
\subsubsection{Paso 1: Padding del Mensaje}

\begin{definition}[Padding SHA-256]
Dado un mensaje $M$ de longitud $L$ bits:
\begin{enumerate}
    \item Añadir un bit '1' al final: $M \leftarrow M || 1$
    \item Añadir $k$ bits '0', donde $k$ es el mínimo tal que $L + 1 + k \equiv 448 \pmod{512}$
    \item Añadir la longitud original $L$ como entero de 64 bits (big-endian)
\end{enumerate}
\end{definition}

\begin{ejemplo}[Padding del mensaje ``abc'']
El mensaje ``abc'' en ASCII es:
\begin{verbatim}
a = 0x61 = 01100001
b = 0x62 = 01100010  
c = 0x63 = 01100011
\end{verbatim}

Mensaje original: 24 bits. Padding:
\begin{enumerate}
    \item Añadir bit '1': ahora 25 bits
    \item $25 + k \equiv 448 \pmod{512}$, entonces $k = 423$ ceros
    \item Añadir longitud $L = 24$ como 64 bits: \texttt{0x0000000000000018}
\end{enumerate}

Resultado (512 bits = 1 bloque):
\begin{verbatim}
61626380 00000000 00000000 00000000
00000000 00000000 00000000 00000000
00000000 00000000 00000000 00000000
00000000 00000000 00000000 00000018
\end{verbatim}
\end{ejemplo}

%-------------------------------------------
\subsubsection{Paso 2: Vector de Inicialización}

\begin{definition}[IV de SHA-256]
El estado inicial consiste en 8 palabras de 32 bits, derivadas de las partes fraccionarias de las raíces cuadradas de los primeros 8 primos:
\begin{align}
H_0^{(0)} &= \texttt{0x6a09e667} \leftarrow \lfloor 2^{32} \times \text{frac}(\sqrt{2}) \rfloor \\
H_1^{(0)} &= \texttt{0xbb67ae85} \leftarrow \lfloor 2^{32} \times \text{frac}(\sqrt{3}) \rfloor \\
H_2^{(0)} &= \texttt{0x3c6ef372} \leftarrow \lfloor 2^{32} \times \text{frac}(\sqrt{5}) \rfloor \\
H_3^{(0)} &= \texttt{0xa54ff53a} \leftarrow \lfloor 2^{32} \times \text{frac}(\sqrt{7}) \rfloor \\
H_4^{(0)} &= \texttt{0x510e527f} \leftarrow \lfloor 2^{32} \times \text{frac}(\sqrt{11}) \rfloor \\
H_5^{(0)} &= \texttt{0x9b05688c} \leftarrow \lfloor 2^{32} \times \text{frac}(\sqrt{13}) \rfloor \\
H_6^{(0)} &= \texttt{0x1f83d9ab} \leftarrow \lfloor 2^{32} \times \text{frac}(\sqrt{17}) \rfloor \\
H_7^{(0)} &= \texttt{0x5be0cd19} \leftarrow \lfloor 2^{32} \times \text{frac}(\sqrt{19}) \rfloor
\end{align}
\end{definition}

\begin{intuicion}
¿Por qué raíces cuadradas de primos? Es una forma de generar constantes ``nothing-up-my-sleeve'' (nada bajo la manga). Cualquiera puede verificar que estas constantes no fueron elegidas maliciosamente --- simplemente se derivan de matemáticas públicas.
\end{intuicion}

%-------------------------------------------
\subsubsection{Paso 3: Constantes de Ronda}

SHA-256 usa 64 constantes $K_0, \ldots, K_{63}$, derivadas de las partes fraccionarias de las raíces cúbicas de los primeros 64 primos:
\[ K_i = \lfloor 2^{32} \times \text{frac}(\sqrt[3]{p_i}) \rfloor \]

\begin{verbatim}
K[0..7]   = 428a2f98 71374491 b5c0fbcf e9b5dba5 
            3956c25b 59f111f1 923f82a4 ab1c5ed5
K[8..15]  = d807aa98 12835b01 243185be 550c7dc3 
            72be5d74 80deb1fe 9bdc06a7 c19bf174
    ... (continúa hasta K[63])
\end{verbatim}

%-------------------------------------------
\subsubsection{Paso 4: Expansión del Mensaje (Message Schedule)}

Cada bloque de 512 bits se divide en 16 palabras de 32 bits ($W_0, \ldots, W_{15}$), y luego se expanden a 64 palabras:

\begin{definition}[Message Schedule]
\begin{align}
W_t &= M_t^{(i)} \quad \text{para } t = 0, \ldots, 15 \\
W_t &= \sigma_1(W_{t-2}) + W_{t-7} + \sigma_0(W_{t-15}) + W_{t-16} \quad \text{para } t = 16, \ldots, 63
\end{align}

Donde las funciones $\sigma$ son:
\begin{align}
\sigma_0(x) &= \text{ROTR}^7(x) \oplus \text{ROTR}^{18}(x) \oplus \text{SHR}^3(x) \\
\sigma_1(x) &= \text{ROTR}^{17}(x) \oplus \text{ROTR}^{19}(x) \oplus \text{SHR}^{10}(x)
\end{align}
\end{definition}

\begin{intuicion}
La expansión de mensaje asegura que cada bit del mensaje original influya en muchas rondas diferentes, amplificando el efecto avalancha.
\end{intuicion}

%-------------------------------------------
\subsubsection{Paso 5: Función de Compresión (64 Rondas)}

\begin{definition}[Funciones Auxiliares]
\begin{align}
Ch(x, y, z) &= (x \land y) \oplus (\neg x \land z) \quad \text{(``choice'')} \\
Maj(x, y, z) &= (x \land y) \oplus (x \land z) \oplus (y \land z) \quad \text{(``majority'')} \\
\Sigma_0(x) &= \text{ROTR}^2(x) \oplus \text{ROTR}^{13}(x) \oplus \text{ROTR}^{22}(x) \\
\Sigma_1(x) &= \text{ROTR}^6(x) \oplus \text{ROTR}^{11}(x) \oplus \text{ROTR}^{25}(x)
\end{align}
\end{definition}

\begin{intuicion}
\begin{itemize}
    \item $Ch(x,y,z)$: ``Si $x$, entonces $y$, sino $z$'' --- cada bit de $x$ elige entre el bit correspondiente de $y$ o $z$
    \item $Maj(x,y,z)$: Voto mayoritario bit a bit --- el resultado es 1 si al menos 2 de los 3 bits son 1
\end{itemize}
\end{intuicion}

\begin{definition}[Una Ronda de SHA-256]
Dado el estado $(a, b, c, d, e, f, g, h)$, en la ronda $t$:
\begin{align}
T_1 &= h + \Sigma_1(e) + Ch(e, f, g) + K_t + W_t \\
T_2 &= \Sigma_0(a) + Maj(a, b, c) \\
h &\leftarrow g, \quad g \leftarrow f, \quad f \leftarrow e, \quad e \leftarrow d + T_1 \\
d &\leftarrow c, \quad c \leftarrow b, \quad b \leftarrow a, \quad a \leftarrow T_1 + T_2
\end{align}
\end{definition}

%-------------------------------------------
\subsubsection{Ejemplo Completo: SHA-256(``abc'')}

\begin{ejemplo}[Cálculo paso a paso]
\textbf{Input:} ``abc'' (3 bytes)

\textbf{Paso 1 - Padding:}
\begin{verbatim}
61626380 00000000 00000000 00000000
00000000 00000000 00000000 00000000
00000000 00000000 00000000 00000000
00000000 00000000 00000000 00000018
\end{verbatim}

\textbf{Paso 2 - Estado inicial:}
\begin{verbatim}
a = 6a09e667    e = 510e527f
b = bb67ae85    f = 9b05688c
c = 3c6ef372    g = 1f83d9ab
d = a54ff53a    h = 5be0cd19
\end{verbatim}

\textbf{Paso 3 - Primera ronda (t=0):}
\begin{align}
W_0 &= \texttt{61626380} \\
K_0 &= \texttt{428a2f98} \\
\Sigma_1(e) &= \texttt{3a6fe667} \\
Ch(e,f,g) &= \texttt{1f85c98c} \\
T_1 &= h + \Sigma_1(e) + Ch(e,f,g) + K_0 + W_0 = \texttt{5d6aebb1}
\end{align}

(Continúa por 64 rondas...)

\textbf{Resultado final:}
\begin{verbatim}
SHA-256("abc") = ba7816bf8f01cfea414140de5dae2223
                 b00361a396177a9cb410ff61f20015ad
\end{verbatim}
\end{ejemplo}

%-------------------------------------------
\subsubsection{Efecto Avalancha}

\begin{clave}
El \textbf{efecto avalancha} significa que cambiar un solo bit del input cambia aproximadamente el 50\% de los bits del output.
\end{clave}

\begin{ejemplo}[Demostración del efecto avalancha]
\begin{verbatim}
SHA-256("abc") = ba7816bf8f01cfea414140de5dae2223...
SHA-256("abd") = 6ca13d52ca70c883e0f0bb101e425a89...
\end{verbatim}

Diferencia: 130 de 256 bits cambiaron (50.8\%). Solo cambió 'c' por 'd'.

\begin{verbatim}
SHA-256("abc") = ba7816bf...
SHA-256("Abc") = 06d90109... (solo mayúscula)
\end{verbatim}

Diferencia: 125 bits (48.8\%)
\end{ejemplo}

%-------------------------------------------
\subsubsection{Propiedades de Seguridad}

\begin{center}
\begin{tabular}{lcc}
\toprule
\textbf{Propiedad} & \textbf{Seguridad Ideal} & \textbf{Mejor Ataque Conocido} \\
\midrule
Preimagen & $2^{256}$ & $2^{254}$ (teórico, no práctico) \\
Segunda preimagen & $2^{256}$ & $2^{256}$ \\
Colisión & $2^{128}$ & $2^{128}$ (birthday attack) \\
\bottomrule
\end{tabular}
\end{center}

\begin{intuicion}
La seguridad contra colisiones es $2^{128}$ (no $2^{256}$) debido a la \textbf{paradoja del cumpleaños}: en un grupo de 23 personas, hay >50\% de probabilidad de que dos compartan cumpleaños. Análogamente, después de $\sqrt{2^{256}} = 2^{128}$ hashes, hay buena probabilidad de encontrar una colisión.
\end{intuicion}

%-------------------------------------------
\subsubsection{Vulnerabilidad: Ataque de Extensión de Longitud}

\begin{advertencia}
La construcción Merkle-Damgård tiene una vulnerabilidad inherente: el \textbf{ataque de extensión de longitud}.
\end{advertencia}

\begin{definition}[Length Extension Attack]
Conociendo $H(M)$ y la longitud de $M$ (pero NO el contenido de $M$), un atacante puede calcular:
\[ H(M || \text{padding} || M') \]
para cualquier sufijo $M'$ elegido por el atacante.
\end{definition}

\begin{intuicion}
El hash $H(M)$ es exactamente el estado interno después de procesar $M$. Si conocemos este estado y la longitud, podemos ``continuar'' el proceso de hashing con bloques adicionales, sin conocer $M$.
\end{intuicion}

\begin{ejemplo}[Explotación práctica]
\textbf{Escenario:} Un servidor web autentica requests usando:
\begin{verbatim}
token = SHA-256(secret_key || user_data)
\end{verbatim}

\textbf{Situación del atacante:}
\begin{itemize}
    \item Conoce un token válido para \texttt{user\_data = "user=alice\&role=guest"}
    \item Sabe que \texttt{secret\_key} tiene 16 bytes
    \item Quiere crear token para \texttt{"...\&role=admin"}
\end{itemize}

\textbf{Ataque:}
\begin{enumerate}
    \item Calcula el padding que SHA-256 añadió internamente
    \item Construye: \texttt{user\_data || padding || "\&role=admin"}
    \item Usa el token original como estado inicial (en lugar de $IV$)
    \item Calcula el nuevo hash para "\&role=admin"
\end{enumerate}

¡El servidor aceptará este token como válido, sin conocer \texttt{secret\_key}!
\end{ejemplo}

\textbf{Mitigaciones:}
\begin{center}
\begin{tabular}{lp{9cm}}
\toprule
\textbf{Mitigación} & \textbf{Descripción} \\
\midrule
HMAC & $\text{HMAC}(K, M) = H((K \oplus \text{opad}) \,||\, H((K \oplus \text{ipad}) \,||\, M))$ \\
SHA-3 & Construcción sponge, inmune por diseño \\
BLAKE2/3 & Incluye contador de bytes, previene extensión \\
Truncamiento & SHA-224 es SHA-256 truncado con IV diferente \\
\bottomrule
\end{tabular}
\end{center}

%-------------------------------------------
\subsubsection{Operaciones Bitwise}

\begin{definition}[Operaciones en SHA-256]
\begin{itemize}
    \item \textbf{ROTR}$^n(x)$: Rotación circular a la derecha por $n$ bits
    \[ \text{ROTR}^n(x) = (x \gg n) \lor (x \ll (32-n)) \]
    
    \item \textbf{SHR}$^n(x)$: Shift lógico a la derecha (rellena con ceros)
    \[ \text{SHR}^n(x) = x \gg n \]
    
    \item \textbf{XOR} ($\oplus$): OR exclusivo bit a bit
    
    \item \textbf{Suma mod $2^{32}$}: Aritmética modular, descartando overflow
\end{itemize}
\end{definition}

\begin{ejemplo}[Rotación vs Shift]
Sea $x = \texttt{0x80000001}$ (bit más significativo y menos significativo en 1)

En binario: \texttt{10000000 00000000 00000000 00000001}

\begin{align}
\text{ROTR}^1(x) &= \texttt{0xC0000000} \quad \text{(el bit 1 ``envuelve'' a la izquierda)} \\
\text{SHR}^1(x) &= \texttt{0x40000000} \quad \text{(el bit 1 se pierde)}
\end{align}
\end{ejemplo}

%-------------------------------------------
\subsubsection{Aplicaciones en Criptomonedas}

\paragraph{Bitcoin: Proof of Work}

\begin{definition}[Mining de Bitcoin]
Los mineros buscan un \texttt{nonce} tal que:
\[ \text{SHA-256}(\text{SHA-256}(\text{block\_header} || \text{nonce})) < \text{target} \]
\end{definition}

\begin{intuicion}
Bitcoin usa SHA-256 \textbf{doble} (SHA-256d) principalmente como defensa contra ataques de extensión de longitud.
\end{intuicion}

\begin{ejemplo}[Dificultad de Bitcoin]
En 2024, el target requiere aproximadamente 79 ceros iniciales en binario:
\begin{itemize}
    \item Probabilidad de éxito por intento: $\approx 2^{-79}$
    \item Hashrate de la red: $\approx 500$ EH/s ($5 \times 10^{20}$ hashes/segundo)
    \item Tiempo promedio para encontrar bloque: 10 minutos
\end{itemize}
\end{ejemplo}

\paragraph{Árboles de Merkle en Bitcoin}

Las transacciones se organizan en un árbol binario donde cada nodo es:
\[ H_{ij} = \text{SHA-256d}(H_i || H_j) \]

Esto permite \textbf{SPV (Simplified Payment Verification)}: verificar una transacción con solo $\log_2(n)$ hashes.

\paragraph{Direcciones de Bitcoin}

\begin{verbatim}
public_key (65 bytes)
    ↓ SHA-256
hash_1 (32 bytes)
    ↓ RIPEMD-160
hash_2 (20 bytes)
    ↓ + version + checksum (SHA-256d)
    ↓ Base58Check
Dirección Bitcoin (25-34 caracteres)
\end{verbatim}

%-------------------------------------------
\subsubsection{Comparación con Otras Funciones Hash}

\begin{center}
\begin{tabular}{lccccc}
\toprule
\textbf{Función} & \textbf{Output} & \textbf{Bloques} & \textbf{Construcción} & \textbf{Estado} \\
\midrule
MD5 & 128 bits & 512 bits & Merkle-Damgård & \textcolor{red}{ROTO} \\
SHA-1 & 160 bits & 512 bits & Merkle-Damgård & \textcolor{red}{ROTO} \\
SHA-256 & 256 bits & 512 bits & Merkle-Damgård & Seguro \\
SHA-512 & 512 bits & 1024 bits & Merkle-Damgård & Seguro \\
SHA-3-256 & 256 bits & 1088 bits & Sponge & Seguro \\
BLAKE3 & 256 bits & 64 bytes & Merkle tree & Seguro \\
\bottomrule
\end{tabular}
\end{center}

%-------------------------------------------
\subsubsection{Rendimiento}

\begin{center}
\begin{tabular}{lcc}
\toprule
\textbf{Plataforma} & \textbf{Velocidad Típica} & \textbf{Notas} \\
\midrule
CPU moderno (software) & 200-500 MB/s & Con SHA-NI: 2+ GB/s \\
GPU (RTX 3090) & $\sim$1.5 GH/s & Optimizado para mining \\
ASIC (Antminer S19) & $\sim$110 TH/s & Hardware especializado \\
\bottomrule
\end{tabular}
\end{center}

%-------------------------------------------
\subsubsection{Aplicación en Banco ZK}

\begin{bancozk}
\textbf{¿Por qué Banco ZK NO usa SHA-256 internamente?}

Aunque SHA-256 es excelente para uso general, tiene un problema crítico para pruebas de conocimiento cero: su estructura de operaciones bitwise (rotaciones, XOR) es extremadamente costosa de representar en circuitos aritméticos.

\begin{center}
\begin{tabular}{lc}
\toprule
\textbf{Función Hash} & \textbf{Restricciones en Circuito ZK} \\
\midrule
SHA-256 & $\sim$25,000 \\
Keccak & $\sim$150,000 \\
Poseidon & $\sim$250 \\
\bottomrule
\end{tabular}
\end{center}

\textbf{Decisión de diseño en Banco ZK:}
\begin{itemize}
    \item \textbf{Dentro del circuito ZK:} Poseidon (250 restricciones)
    \item \textbf{Interfaz con L1/externos:} SHA-256 cuando sea necesario
    \item \textbf{Nunca:} Construcciones vulnerables a length extension para autenticación
\end{itemize}

La reducción de 25,000 a 250 restricciones (100x) significa que las pruebas se generan mucho más rápido y son más baratas de verificar on-chain.
\end{bancozk}

%-------------------------------------------
\subsubsection{Implementación de Referencia (Pseudocódigo)}

\begin{verbatim}
function sha256(message):
    # Paso 1: Padding
    padded = pad(message)
    blocks = split_into_512bit_blocks(padded)
    
    # Paso 2: Inicializar estado
    H = [0x6a09e667, 0xbb67ae85, 0x3c6ef372, 0xa54ff53a,
         0x510e527f, 0x9b05688c, 0x1f83d9ab, 0x5be0cd19]
    
    # Paso 3: Procesar cada bloque
    for block in blocks:
        W = expand_message(block)  # 16 → 64 palabras
        a,b,c,d,e,f,g,h = H
        
        for t in range(64):
            T1 = h + Σ1(e) + Ch(e,f,g) + K[t] + W[t]
            T2 = Σ0(a) + Maj(a,b,c)
            h,g,f,e,d,c,b,a = g,f,e,(d+T1),c,b,a,(T1+T2)
        
        H = [x + y for x,y in zip(H, [a,b,c,d,e,f,g,h])]
    
    return concatenate(H)
\end{verbatim}

%-------------------------------------------
\subsubsection{Ejercicios Propuestos}

\begin{enumerate}
    \item Calcular manualmente el padding para el mensaje ``Hello'' (5 bytes).
    
    \item Verificar que SHA-256(``'') (string vacío) = \texttt{e3b0c442...}
    
    \item Implementar las funciones $Ch$, $Maj$, $\Sigma_0$, $\Sigma_1$ en tu lenguaje favorito.
    
    \item Demostrar el efecto avalancha: hashear ``test0'', ``test1'', ..., ``test9'' y contar bits diferentes entre hashes consecutivos.
    
    \item ¿Por qué Bitcoin usa SHA-256 doble en lugar de simple?
    
    \item Calcular el padding completo para ``ZK'' (2 bytes) usando parámetros de SHA-256.
    
    \item ¿Cuántos bloques de 512 bits se necesitan para hashear 447 bits? ¿Y para 448 bits? ¿Por qué la diferencia?
\end{enumerate}

%  FIN ==============================================================================
% ALE 
%===========================================
% SUBSECCIÓN: KECCAK (SHA-3) - VERSIÓN EXPANDIDA
% Usa tus entornos: definition, example, keyidea, bancozk, 
%                   mathnote, intuicion, advertencia, remark
%===========================================

\subsection{Keccak (SHA-3)}

\begin{definition}[Keccak]
Función hash basada en construcción esponja, estandarizada como SHA-3.
\end{definition}

Arquitectura esponja:
\begin{itemize}
    \item Estado de 1600 bits
    \item Parámetros: rate $r$ y capacity $c$ ($r + c = 1600$)
    \item Fase de absorción: XOR input con estado
    \item Fase de exprimido: extraer output
    \item Permutación Keccak-$f[1600]$: 24 rondas
\end{itemize}

Ventajas sobre SHA-2:
\begin{itemize}
    \item Resistente a ataques de extensión de longitud
    \item Diseño completamente diferente (diversidad)
    \item Parámetros flexibles (SHAKE128, SHAKE256)
\end{itemize}

Usado en: Ethereum (direcciones, state root), aplicaciones post-SHA-2.

%-------------------------------------------
\subsubsection{Contexto Histórico: La Competencia SHA-3}

\begin{intuicion}
En 2007, NIST (National Institute of Standards and Technology) anunció una competencia pública para seleccionar un nuevo estándar de hash llamado SHA-3. La motivación no era que SHA-2 estuviera roto, sino que SHA-2 compartía diseño similar con SHA-1 (que mostraba debilidades). NIST quería un ``plan B'' con arquitectura completamente diferente.
\end{intuicion}

La competencia SHA-3 fue una de las más importantes en la historia de la criptografía:

\begin{center}
\begin{tabular}{ll}
\toprule
\textbf{Fase} & \textbf{Detalle} \\
\midrule
2007 & Convocatoria abierta \\
2008 & 64 candidatos presentados \\
2009 & Ronda 1: 51 candidatos aceptados, 14 pasan a Ronda 2 \\
2010 & Ronda 2: 5 finalistas seleccionados \\
2012 & \textbf{Keccak} declarado ganador \\
2015 & Publicación oficial de SHA-3 (FIPS 202) \\
\bottomrule
\end{tabular}
\end{center}

\textbf{Los 5 finalistas fueron:}
\begin{itemize}
    \item \textbf{Keccak} (Bertoni, Daemen, Peeters, Van Assche) --- \textit{Ganador}
    \item BLAKE (Aumasson et al.) --- Muy rápido en software
    \item Grøstl (Knudsen et al.) --- Basado en AES
    \item JH (Wu) --- Diseño elegante
    \item Skein (Schneier et al.) --- Basado en Threefish
\end{itemize}

\begin{remark}
Joan Daemen, uno de los creadores de Keccak, también co-diseñó AES (Rijndael). Su experiencia en diseño de primitivas criptográficas fue crucial para la seguridad de Keccak.
\end{remark}

%-------------------------------------------
\subsubsection{La Construcción Sponge}

A diferencia de Merkle-Damgård (usada en SHA-2), Keccak usa la \textbf{construcción sponge} (esponja), una arquitectura fundamentalmente diferente.

\begin{definition}[Construcción Sponge]
Una construcción sponge es una función criptográfica que:
\begin{enumerate}
    \item Mantiene un \textbf{estado interno} de $b$ bits, dividido en:
    \begin{itemize}
        \item \textbf{Rate} ($r$ bits): parte que interactúa con input/output
        \item \textbf{Capacity} ($c$ bits): parte que \textit{nunca} se expone directamente
    \end{itemize}
    \item Opera en dos fases:
    \begin{itemize}
        \item \textbf{Absorción}: el mensaje se ``absorbe'' en el estado
        \item \textbf{Exprimido}: el hash se ``exprime'' del estado
    \end{itemize}
    \item Usa una \textbf{permutación} $f$ para mezclar el estado
\end{enumerate}
\end{definition}

\begin{center}
\begin{tikzpicture}[
    block/.style={draw, minimum width=1.5cm, minimum height=1cm},
    state/.style={draw, minimum width=3cm, minimum height=1.5cm},
    arrow/.style={->, thick}
]
    % Título
    \node at (0, 3) {\textbf{Construcción Sponge}};
    
    % Estado inicial
    \node[state, fill=gray!10] (s0) at (0, 1.5) {Estado = $0^b$};
    \node[below=0.1cm of s0, font=\small] {$r$ bits | $c$ bits};
    
    % Fase de absorción
    \node at (-4, 0) {\textbf{Absorción}};
    
    \node[block] (xor1) at (-2, 0) {$\oplus$};
    \node[above=0.3cm of xor1, font=\small] {$M_1$};
    \node[block, fill=blue!20] (f1) at (0, 0) {$f$};
    
    \node[block] (xor2) at (2, 0) {$\oplus$};
    \node[above=0.3cm of xor2, font=\small] {$M_2$};
    \node[block, fill=blue!20] (f2) at (4, 0) {$f$};
    
    \node at (5.5, 0) {$\cdots$};
    
    % Flechas absorción
    \draw[arrow] (s0) -- (xor1);
    \draw[arrow] (xor1) -- (f1);
    \draw[arrow] (f1) -- (xor2);
    \draw[arrow] (xor2) -- (f2);
    
    % Fase de exprimido
    \node at (-4, -1.5) {\textbf{Exprimido}};
    
    \node[block, fill=blue!20] (f3) at (0, -1.5) {$f$};
    \node[block, fill=blue!20] (f4) at (3, -1.5) {$f$};
    
    \node[below=0.5cm of f3, font=\small] {$Z_1$ ($r$ bits)};
    \node[below=0.5cm of f4, font=\small] {$Z_2$ ($r$ bits)};
    
    \draw[arrow] (-2, -1.5) -- (f3);
    \draw[arrow] (f3) -- (f4);
    \draw[arrow] (f3) -- ++(0, -1);
    \draw[arrow] (f4) -- ++(0, -1);
    
\end{tikzpicture}
\end{center}

\begin{keyidea}
La seguridad de la construcción sponge depende de la \textbf{capacity} $c$. Esta parte del estado \textbf{nunca} se expone ni se modifica directamente por el input. Un atacante tendría que adivinar $2^{c/2}$ estados posibles para montar un ataque de colisión.

En SHA-3-256: $c = 512$ bits $\Rightarrow$ seguridad de $2^{256}$ contra colisiones.
\end{keyidea}

\begin{intuicion}
\textbf{Analogía de la esponja real:}

Imagina una esponja de cocina:
\begin{enumerate}
    \item \textbf{Absorción}: Sumerges la esponja en agua (tu mensaje). El agua entra gradualmente, mezclándose con lo que ya había dentro.
    \item \textbf{Exprimido}: Aprietas la esponja y sale agua. Pero no sale \textit{toda} el agua --- parte queda atrapada en el interior (la capacity).
\end{enumerate}

La parte que ``queda atrapada'' es lo que hace segura a la construcción: aunque veas lo que sale (el hash), no puedes saber el estado interno completo.
\end{intuicion}

%-------------------------------------------
\subsubsection{Parámetros de SHA-3}

\begin{center}
\begin{tabular}{lccccc}
\toprule
\textbf{Variante} & \textbf{Output} & \textbf{Rate $r$} & \textbf{Capacity $c$} & \textbf{Seguridad} \\
\midrule
SHA-3-224 & 224 bits & 1152 & 448 & 112 bits \\
SHA-3-256 & 256 bits & 1088 & 512 & 128 bits \\
SHA-3-384 & 384 bits & 832 & 768 & 192 bits \\
SHA-3-512 & 512 bits & 576 & 1024 & 256 bits \\
\midrule
SHAKE128 & Variable & 1344 & 256 & 128 bits \\
SHAKE256 & Variable & 1088 & 512 & 256 bits \\
\bottomrule
\end{tabular}
\end{center}

\begin{mathnote}
La relación fundamental es:
\[ r + c = b = 1600 \text{ bits} \]

El \textbf{nivel de seguridad} contra ataques genéricos es:
\begin{itemize}
    \item Preimagen: $\min(2^{output}, 2^c)$
    \item Colisión: $\min(2^{output/2}, 2^{c/2})$
\end{itemize}

Por ejemplo, SHA-3-256 tiene $c = 512$, por lo que la seguridad contra colisiones es $\min(2^{128}, 2^{256}) = 2^{128}$.
\end{mathnote}

%-------------------------------------------
\subsubsection{El Estado Interno: Estructura 3D}

\begin{definition}[Estado de Keccak]
El estado de 1600 bits se organiza como un arreglo tridimensional:
\[ A[x][y][z] \quad \text{donde } x, y \in \{0,1,2,3,4\}, \quad z \in \{0,1,\ldots,63\} \]

Esto forma:
\begin{itemize}
    \item $5 \times 5 = 25$ \textbf{lanes} (carriles) de 64 bits cada uno
    \item $25 \times 64 = 1600$ bits totales
    \item Visualización: un cubo de $5 \times 5 \times 64$
\end{itemize}
\end{definition}

\begin{center}
\begin{tikzpicture}[scale=0.6]
    % Dibujar cuadrícula 5x5
    \foreach \x in {0,1,2,3,4} {
        \foreach \y in {0,1,2,3,4} {
            \draw (\x, \y) rectangle (\x+1, \y+1);
        }
    }
    
    % Etiquetas
    \node at (2.5, -0.5) {$x$};
    \node at (-0.5, 2.5) {$y$};
    \node at (6, 2.5) {Cada celda = 1 lane (64 bits)};
    \node at (2.5, 5.7) {\textbf{Vista superior del estado}};
    
    % Marcar una celda
    \fill[blue!30] (2,2) rectangle (3,3);
    \node at (2.5, 2.5) {\small $A[2][2]$};
\end{tikzpicture}
\end{center}

%-------------------------------------------
\subsubsection{La Permutación Keccak-$f$: 24 Rondas}

La permutación $f$ que mezcla el estado consiste en 24 rondas, cada una aplicando 5 operaciones en secuencia:

\begin{definition}[Ronda de Keccak-$f$]
Cada ronda $R_i$ consiste en:
\[ R_i = \iota \circ \chi \circ \pi \circ \rho \circ \theta \]

Aplicadas en este orden:
\begin{enumerate}
    \item $\theta$ (theta) --- Paridad y difusión lineal
    \item $\rho$ (rho) --- Rotación de bits dentro de cada lane
    \item $\pi$ (pi) --- Permutación de posiciones de lanes
    \item $\chi$ (chi) --- Transformación no lineal
    \item $\iota$ (iota) --- Adición de constante de ronda
\end{enumerate}
\end{definition}

\paragraph{Operación $\theta$ (Theta): Difusión}

\begin{definition}[Theta]
\begin{align}
C[x] &= A[x][0] \oplus A[x][1] \oplus A[x][2] \oplus A[x][3] \oplus A[x][4] \\
D[x] &= C[x-1] \oplus \text{ROT}(C[x+1], 1) \\
A'[x][y] &= A[x][y] \oplus D[x]
\end{align}

donde los índices son módulo 5.
\end{definition}

\begin{intuicion}
$\theta$ calcula la ``paridad'' de cada columna y la propaga a las columnas vecinas. Esto asegura que cambiar un solo bit afecte a múltiples posiciones --- el inicio del efecto avalancha.
\end{intuicion}

\paragraph{Operación $\rho$ (Rho): Rotación}

\begin{definition}[Rho]
Cada lane $A[x][y]$ se rota por un número diferente de bits:
\[ A'[x][y][z] = A[x][y][z - r[x][y] \mod 64] \]

Los offsets de rotación $r[x][y]$ son:
\begin{center}
\begin{tabular}{c|ccccc}
$r[x][y]$ & $y=0$ & $y=1$ & $y=2$ & $y=3$ & $y=4$ \\
\hline
$x=0$ & 0 & 36 & 3 & 41 & 18 \\
$x=1$ & 1 & 44 & 10 & 45 & 2 \\
$x=2$ & 62 & 6 & 43 & 15 & 61 \\
$x=3$ & 28 & 55 & 25 & 21 & 56 \\
$x=4$ & 27 & 20 & 39 & 8 & 14 \\
\end{tabular}
\end{center}
\end{definition}

\paragraph{Operación $\pi$ (Pi): Permutación}

\begin{definition}[Pi]
Los lanes se reordenan según:
\[ A'[x][y] = A[x + 3y \mod 5][x] \]

Esta es una permutación de las 25 posiciones $(x, y)$ en el plano.
\end{definition}

\paragraph{Operación $\chi$ (Chi): No Linealidad}

\begin{definition}[Chi]
\[ A'[x][y][z] = A[x][y][z] \oplus ((\neg A[x+1][y][z]) \land A[x+2][y][z]) \]
\end{definition}

\begin{keyidea}
$\chi$ es la \textbf{única operación no lineal} en Keccak. Sin ella, todo sería lineal (sumas y XORs) y fácilmente invertible. La no linealidad es esencial para la seguridad criptográfica --- impide que un atacante resuelva ecuaciones simples para revertir el hash.
\end{keyidea}

\paragraph{Operación $\iota$ (Iota): Constante de Ronda}

\begin{definition}[Iota]
\[ A'[0][0] = A[0][0] \oplus RC_i \]

donde $RC_i$ es una constante diferente para cada ronda $i \in \{0, \ldots, 23\}$.
\end{definition}

\begin{intuicion}
Sin $\iota$, todas las rondas serían idénticas. Un atacante podría explotar esta simetría. Las constantes $RC_i$ ``rompen la simetría'' y aseguran que cada ronda sea única.
\end{intuicion}

%-------------------------------------------
\subsubsection{Algoritmo Completo}

\begin{example}[Procesamiento de un mensaje]
\textbf{Input:} Mensaje $M$ de longitud arbitraria

\textbf{Paso 1: Padding}

Keccak usa padding simple: $M' = M \,||\, 0x06 \,||\, 0x00^* \,||\, 0x80$

(El padding de SHA-3 es ligeramente diferente: añade bits específicos para distinguir SHA-3 de SHAKE)

\textbf{Paso 2: Inicializar estado}
\[ S = 0^{1600} \quad \text{(1600 bits en cero)} \]

\textbf{Paso 3: Absorción}

Dividir $M'$ en bloques de $r$ bits: $M' = M_1 || M_2 || \cdots || M_k$

\begin{verbatim}
for i = 1 to k:
    S[0..r-1] = S[0..r-1] XOR M_i    # XOR con rate
    S = Keccak-f(S)                   # Aplicar permutación
\end{verbatim}

\textbf{Paso 4: Exprimido}

\begin{verbatim}
Z = ""
while len(Z) < output_length:
    Z = Z || S[0..r-1]    # Extraer rate
    S = Keccak-f(S)       # Aplicar permutación
return Z[0..output_length-1]
\end{verbatim}
\end{example}

%-------------------------------------------
\subsubsection{¿Por Qué es Inmune al Ataque de Extensión de Longitud?}

\begin{advertencia}
Recordemos que SHA-256 (Merkle-Damgård) es vulnerable: conociendo $H(M)$ y $|M|$, se puede calcular $H(M || pad || M')$ sin conocer $M$.
\end{advertencia}

\begin{keyidea}
Keccak es \textbf{inmune} a este ataque porque el hash NO revela el estado interno completo.

En Merkle-Damgård:
\[ H(M) = \text{estado completo después de procesar } M \]

En Sponge:
\[ H(M) = \text{solo los primeros } r \text{ bits del estado} \]

La capacity ($c$ bits) \textbf{nunca se revela}. Sin conocer esos $c$ bits, el atacante no puede continuar el cálculo.
\end{keyidea}

\begin{example}[Por qué falla el ataque]
Supongamos SHA-3-256 ($r = 1088$, $c = 512$):

\begin{enumerate}
    \item El atacante conoce $H(M)$ = 256 bits
    \item El estado interno tiene 1600 bits
    \item El atacante conoce como máximo 256 bits (el hash)
    \item Faltan $1600 - 256 = 1344$ bits desconocidos
    \item Para continuar el cálculo, necesitaría adivinar al menos los 512 bits de capacity
    \item Complejidad: $2^{512}$ intentos --- completamente impracticable
\end{enumerate}
\end{example}

%-------------------------------------------
\subsubsection{SHAKE: Funciones de Salida Extensible (XOF)}

\begin{definition}[XOF - Extendable Output Function]
Una XOF es una función hash que puede producir salida de longitud arbitraria:
\[ \text{SHAKE}n(M, d) : \{0,1\}^* \times \mathbb{N} \to \{0,1\}^d \]

donde $d$ es la longitud de salida deseada.
\end{definition}

\begin{example}[Usos de SHAKE]
\begin{verbatim}
SHAKE128("password", 256)  -> 256 bits de output
SHAKE128("password", 512)  -> 512 bits de output
SHAKE128("password", 1024) -> 1024 bits de output
SHAKE256("key_material", 32768) -> 32KB de output
\end{verbatim}

Todos con el mismo mensaje pero diferente longitud de salida.
\end{example}

\begin{keyidea}
SHAKE es posible gracias a la construcción sponge. Durante la fase de exprimido, puedes seguir ``exprimiendo'' tantos bloques como necesites. En Merkle-Damgård esto no es posible porque el output es exactamente el estado final, de tamaño fijo.
\end{keyidea}

\textbf{Aplicaciones de XOF:}
\begin{itemize}
    \item \textbf{Derivación de claves}: Generar múltiples claves desde un secreto master
    \item \textbf{Padding criptográfico}: Generar máscaras de longitud variable
    \item \textbf{Generación de números pseudoaleatorios}: CSPRNG basado en SHAKE
    \item \textbf{Firmas post-cuánticas}: SPHINCS+, Dilithium usan SHAKE internamente
\end{itemize}

%-------------------------------------------
\subsubsection{Comparación SHA-2 vs SHA-3}

\begin{center}
\begin{tabular}{lcc}
\toprule
\textbf{Característica} & \textbf{SHA-2 (SHA-256)} & \textbf{SHA-3 (Keccak)} \\
\midrule
Construcción & Merkle-Damgård & Sponge \\
Operaciones & AND, OR, XOR, ROT, ADD & XOR, ROT, AND, NOT \\
Estado interno & 256 bits & 1600 bits \\
Rondas & 64 & 24 \\
Extension attack & Vulnerable & Inmune \\
Salida variable & No & Sí (SHAKE) \\
Velocidad (software) & Más rápido & Más lento \\
Velocidad (hardware) & Similar & Muy eficiente \\
Madurez & Más analizado & Más nuevo \\
\bottomrule
\end{tabular}
\end{center}

\begin{remark}
SHA-3 es generalmente más lento que SHA-256 en CPUs convencionales, pero puede ser muy eficiente en hardware dedicado debido a su estructura regular. En ASICs, Keccak puede superar a SHA-256.
\end{remark}

%-------------------------------------------
\subsubsection{Keccak en Ethereum}

\begin{intuicion}
Ethereum usa una versión de Keccak que es \textbf{ligeramente diferente} del SHA-3 estandarizado. Esto se debe a que Ethereum adoptó Keccak antes de que NIST finalizara SHA-3 con cambios menores en el padding.
\end{intuicion}

\begin{definition}[Keccak-256 vs SHA-3-256]
\begin{itemize}
    \item \textbf{Keccak-256} (Ethereum): Padding $0x01$
    \item \textbf{SHA-3-256} (NIST): Padding $0x06$
\end{itemize}

Para el mismo input, producen hashes \textbf{diferentes}.
\end{definition}

\begin{example}[Diferencia práctica]
\begin{verbatim}
Keccak-256("") = c5d2460186f7233c927e7db2dcc703c0
                 e500b653ca82273b7bfad8045d85a470
                 
SHA-3-256("")  = a7ffc6f8bf1ed76651c14756a061d662
                 f580ff4de43b49fa82d80a4b80f8434a
\end{verbatim}

¡Hashes completamente diferentes para el string vacío!
\end{example}

\textbf{Usos de Keccak-256 en Ethereum:}
\begin{itemize}
    \item \textbf{Direcciones}: $\text{address} = \text{Keccak-256}(\text{pubkey})[12..31]$ (últimos 20 bytes)
    \item \textbf{State root}: Raíz del Merkle Patricia Trie
    \item \textbf{Transaction hash}: Identificador único de transacciones
    \item \textbf{Function selectors}: $\text{selector} = \text{Keccak-256}(\text{signature})[0..3]$ (primeros 4 bytes)
\end{itemize}

\begin{example}[Selector de función en Solidity]
\begin{verbatim}
function transfer(address to, uint256 amount)

signature = "transfer(address,uint256)"
selector = Keccak-256(signature)[0:4] = 0xa9059cbb
\end{verbatim}

Por eso en los calldata de transacciones ERC-20 siempre ves \texttt{0xa9059cbb...}
\end{example}

%-------------------------------------------
\subsubsection{Rendimiento}

\begin{center}
\begin{tabular}{lccc}
\toprule
\textbf{Plataforma} & \textbf{SHA-256} & \textbf{SHA-3-256} & \textbf{Ratio} \\
\midrule
Intel x86-64 (software) & 500 MB/s & 300 MB/s & SHA-2 1.7x \\
Intel con SHA-NI & 2000 MB/s & 300 MB/s & SHA-2 6.7x \\
ARM Cortex-A72 & 200 MB/s & 350 MB/s & SHA-3 1.75x \\
FPGA (optimizado) & Variable & Muy eficiente & SHA-3 gana \\
Circuito ZK (R1CS) & 25,000 & 150,000 & SHA-2 6x \\
\bottomrule
\end{tabular}
\end{center}

\begin{advertencia}
En \textbf{circuitos ZK}, tanto SHA-256 como Keccak son extremadamente costosos:
\begin{itemize}
    \item SHA-256: $\sim$25,000 restricciones
    \item Keccak-256: $\sim$150,000 restricciones
\end{itemize}

Ambos usan operaciones bitwise (XOR, AND, rotaciones) que son muy ineficientes en aritmética de campos finitos. Por esto se desarrollaron hashes ``ZK-friendly'' como Poseidon.
\end{advertencia}

%-------------------------------------------
\subsubsection{Seguridad: Estado del Arte}

\begin{center}
\begin{tabular}{lcc}
\toprule
\textbf{Ataque} & \textbf{SHA-3-256 Completo} & \textbf{Mejor Ataque Conocido} \\
\midrule
Preimagen & $2^{256}$ & $2^{256}$ (no hay ataque) \\
Colisión & $2^{128}$ & $2^{128}$ (no hay ataque) \\
Distinguisher & --- & 8 rondas (de 24) \\
\bottomrule
\end{tabular}
\end{center}

\begin{keyidea}
SHA-3 tiene un margen de seguridad muy amplio. Los mejores ataques conocidos solo afectan versiones reducidas (8 de 24 rondas). El diseño con 24 rondas tiene un factor de seguridad de 3x sobre los mejores ataques.
\end{keyidea}

%-------------------------------------------
\subsubsection{Aplicación en Banco ZK}

\begin{bancozk}
\textbf{Rol de Keccak en Banco ZK:}

\textbf{1. Interoperabilidad con Ethereum:}
Para verificar datos de Ethereum dentro de circuitos ZK (como probar que una transacción existe), necesitamos verificar hashes Keccak. Esto es costoso ($\sim$150,000 restricciones) pero necesario para bridges y verificación cross-chain.

\textbf{2. Decisión de diseño:}
\begin{center}
\begin{tabular}{lcc}
\toprule
\textbf{Uso} & \textbf{Hash} & \textbf{Restricciones} \\
\midrule
Estado interno & Poseidon & $\sim$250 \\
Nullifiers & Poseidon & $\sim$250 \\
Merkle trees & Poseidon & $\sim$250 \\
Verificar Ethereum & Keccak & $\sim$150,000 \\
Compatibilidad Bitcoin & SHA-256 & $\sim$25,000 \\
\bottomrule
\end{tabular}
\end{center}

\textbf{3. Optimización:}
Cuando es posible, movemos la verificación de Keccak/SHA-256 fuera del circuito. Por ejemplo, verificamos el hash en la capa de contrato inteligente y solo pasamos el resultado al circuito ZK como input público.

\textbf{4. Futuro:}
Proyectos como ``Keccak circuits optimizados'' y ``precompiles de hash en L2'' están reduciendo estos costos. Algunos ZK-rollups ya ofrecen precompiles que reducen el costo de Keccak a $\sim$50,000 restricciones.
\end{bancozk}

%-------------------------------------------
\subsubsection{Ejercicios}

\begin{enumerate}
    \item Calcular manualmente una ronda de $\theta$ para un estado simplificado de $5 \times 5$ bits (en lugar de $5 \times 5 \times 64$).
    
    \item Si SHA-3-256 tiene $c = 512$ bits de capacity, ¿cuál es la complejidad de un ataque de colisión? ¿Y de preimagen?
    
    \item Explicar por qué SHAKE puede producir salida de cualquier longitud mientras que SHA-256 no puede.
    
    \item Un sistema usa: \texttt{token = Keccak(secret || data)}. ¿Es vulnerable a ataques de extensión de longitud? Justificar.
    
    \item Dado que Ethereum usa Keccak-256 (no SHA-3-256), ¿qué sucedería si un desarrollador usa una librería de SHA-3 estándar para verificar direcciones Ethereum?
    
    \item Calcular el function selector para \texttt{approve(address,uint256)} en Solidity.
    
    \item Si Keccak cuesta 150,000 restricciones en un circuito ZK y Poseidon cuesta 250, ¿cuántas operaciones Poseidon equivalen a una operación Keccak? ¿Qué implicaciones tiene esto para el diseño de protocolos ZK?
\end{enumerate}
% ALE FIN


% empieza aca poseidon 
%===========================================
% SUBSECCIÓN: POSEIDON - VERSIÓN EXPANDIDA
% Usa tus entornos: definition, example, keyidea, bancozk, 
%                   mathnote, intuicion, advertencia, remark
%===========================================

\subsection{Poseidon}

\begin{definition}[Poseidon]
Función hash algebraica diseñada para eficiencia en circuitos ZK.
\end{definition}

Arquitectura:
\begin{itemize}
    \item Opera sobre campos finitos $\Fp$ nativamente
    \item S-box: $x \mapsto x^\alpha$ (típicamente $\alpha = 5$)
    \item Capa lineal: multiplicación por matriz MDS
    \item Rondas parciales + rondas completas
\end{itemize}

\begin{mathnote}
Costo en constraints:
\begin{itemize}
    \item Poseidon: $\sim$250 constraints
    \item SHA-256: $\sim$25,000 constraints
    \item Keccak: $\sim$150,000 constraints
\end{itemize}
Por esto Poseidon es el estándar en ZK.
\end{mathnote}

Otros hashes algebraicos:
\begin{itemize}
    \item \textbf{MiMC:} $x \mapsto x^3$, simple pero lento
    \item \textbf{Rescue:} Alternancia de S-boxes
    \item \textbf{Griffin:} Optimizado para STARKs
    \item \textbf{Anemoi:} Nuevo, muy eficiente
\end{itemize}

%-------------------------------------------
\subsubsection{El Problema: ¿Por Qué Necesitamos Hashes Nuevos?}

\begin{intuicion}
Las pruebas de conocimiento cero funcionan convirtiendo computaciones en ecuaciones polinomiales sobre campos finitos. Cada operación del programa original se traduce en ``restricciones'' (constraints) que el prover debe satisfacer.

El problema es que SHA-256 y Keccak usan operaciones \textbf{bitwise}: AND, OR, XOR, rotaciones. Estas operaciones son baratas en CPUs pero \textbf{extremadamente costosas} en aritmética de campos finitos.
\end{intuicion}

\begin{example}[Costo de operaciones bitwise en ZK]
Para representar XOR de dos valores de 256 bits en un circuito aritmético:

\textbf{Paso 1:} Descomponer cada número en 256 bits individuales
\[ x = \sum_{i=0}^{255} b_i \cdot 2^i \quad \text{donde } b_i \in \{0, 1\} \]
Esto requiere 256 constraints para verificar que cada $b_i$ es binario: $b_i \cdot (1 - b_i) = 0$

\textbf{Paso 2:} Calcular XOR bit a bit
\[ (a \oplus b)_i = a_i + b_i - 2 \cdot a_i \cdot b_i \]
256 constraints más.

\textbf{Paso 3:} Recomponer el resultado
Otros 256 constraints.

\textbf{Total para UN XOR de 256 bits:} $\sim$768 constraints

SHA-256 tiene cientos de operaciones XOR, AND, rotaciones... Por eso cuesta $\sim$25,000 constraints.
\end{example}

\begin{keyidea}
La solución es diseñar funciones hash que usen operaciones \textbf{nativas} en campos finitos:
\begin{itemize}
    \item Suma: $a + b \mod p$ → 0 constraints (gratis)
    \item Multiplicación: $a \times b \mod p$ → 1 constraint
    \item Potenciación: $x^5$ → 4 multiplicaciones → 4 constraints
\end{itemize}

Poseidon está diseñado específicamente para minimizar el número de multiplicaciones.
\end{keyidea}

%-------------------------------------------
\subsubsection{Historia y Contexto}

Poseidon fue presentado en 2019 por un equipo de investigadores:
\begin{itemize}
    \item Lorenzo Grassi (Universidad de Graz)
    \item Dmitry Khovratovich (Ethereum Foundation)
    \item Christian Rechberger (TU Graz)
    \item Arnab Roy (Fujitsu Labs)
    \item Markus Schofnegger (TU Graz)
\end{itemize}

\begin{remark}
El nombre ``Poseidon'' viene del dios griego del mar. Los autores usaron una temática de dioses griegos para sus hashes: \textbf{Poseidon} (hash general), \textbf{Neptune} (para compromisos), \textbf{Starkad} y \textbf{Poseidon$^\pi$} (variantes optimizadas).
\end{remark}

La evolución de hashes ZK-friendly:
\begin{center}
\begin{tabular}{lccc}
\toprule
\textbf{Año} & \textbf{Hash} & \textbf{Constraints (aprox.)} & \textbf{Innovación} \\
\midrule
2016 & MiMC & $\sim$700 & Primera propuesta algebraica \\
2019 & Poseidon & $\sim$250 & Rondas parciales \\
2019 & Rescue & $\sim$400 & S-boxes alternadas \\
2022 & Griffin & $\sim$150 & Optimizado para STARKs \\
2022 & Anemoi & $\sim$120 & Estado del arte \\
\bottomrule
\end{tabular}
\end{center}

%-------------------------------------------
\subsubsection{Arquitectura de Poseidon}

\begin{definition}[Estructura de Poseidon]
Poseidon opera sobre un estado de $t$ elementos de campo $\Fp$:
\[ \text{state} = (s_0, s_1, \ldots, s_{t-1}) \in \Fp^t \]

Cada ronda aplica tres operaciones:
\begin{enumerate}
    \item \textbf{AddRoundConstants (ARC):} $s_i \leftarrow s_i + c_i$
    \item \textbf{S-box:} $s_i \leftarrow s_i^\alpha$ (o solo a $s_0$ en rondas parciales)
    \item \textbf{MDS:} $\text{state} \leftarrow M \cdot \text{state}$
\end{enumerate}

La estructura de rondas es:
\[ R_F/2 \text{ rondas completas} \to R_P \text{ rondas parciales} \to R_F/2 \text{ rondas completas} \]
\end{definition}

\begin{center}
\begin{tikzpicture}[
    roundbox/.style={draw, rounded corners, minimum width=2cm, minimum height=0.8cm},
    arrow/.style={->, thick}
]
    % Input
    \node (input) at (0, 0) {Input};
    
    % Rondas completas (inicio)
    \node[roundbox, fill=blue!20] (rf1) at (3, 0) {$R_F/2$ Full};
    
    % Rondas parciales
    \node[roundbox, fill=green!20] (rp) at (6.5, 0) {$R_P$ Partial};
    
    % Rondas completas (fin)
    \node[roundbox, fill=blue!20] (rf2) at (10, 0) {$R_F/2$ Full};
    
    % Output
    \node (output) at (13, 0) {Hash};
    
    % Flechas
    \draw[arrow] (input) -- (rf1);
    \draw[arrow] (rf1) -- (rp);
    \draw[arrow] (rp) -- (rf2);
    \draw[arrow] (rf2) -- (output);
    
    % Leyenda
    \node[font=\small] at (3, -1) {S-box a todos};
    \node[font=\small] at (6.5, -1) {S-box solo a $s_0$};
    \node[font=\small] at (10, -1) {S-box a todos};
\end{tikzpicture}
\end{center}

%-------------------------------------------
\subsubsection{La S-box: $x \mapsto x^\alpha$}

\begin{definition}[S-box de Poseidon]
La S-box es la función no lineal:
\[ S(x) = x^\alpha \]

donde $\alpha$ es el menor entero positivo tal que $\gcd(\alpha, p-1) = 1$.

Valores típicos:
\begin{itemize}
    \item $\alpha = 5$ para la mayoría de campos primos usados en SNARKs
    \item $\alpha = 3$ no sirve cuando $3 | (p-1)$ (común en curvas elípticas)
    \item $\alpha = 7$ usado en algunos campos específicos
\end{itemize}
\end{definition}

\begin{mathnote}
¿Por qué necesitamos $\gcd(\alpha, p-1) = 1$?

Para que la S-box sea una \textbf{permutación} (biyectiva), necesitamos que $x \mapsto x^\alpha$ sea invertible. Esto ocurre si y solo si $\alpha$ es coprimo con $p-1$.

La inversa es:
\[ S^{-1}(y) = y^{\alpha^{-1} \mod (p-1)} \]

Por ejemplo, si $\alpha = 5$ y $p - 1 = 4 \cdot k$ para algún $k$, entonces $\gcd(5, p-1)$ podría no ser 1. Hay que verificar para cada campo.
\end{mathnote}

\begin{example}[Cálculo de S-box]
Sea $p = 17$ (campo $\F_{17}$) y $\alpha = 5$.

Verificación: $p - 1 = 16$, y $\gcd(5, 16) = 1$ ✓

Calculamos $S(x) = x^5 \mod 17$:
\begin{center}
\begin{tabular}{c|cccccccccccccccc}
$x$ & 1 & 2 & 3 & 4 & 5 & 6 & 7 & 8 & 9 & 10 & 11 & 12 & 13 & 14 & 15 & 16 \\
\hline
$x^5$ & 1 & 15 & 5 & 4 & 14 & 7 & 11 & 9 & 16 & 10 & 6 & 3 & 13 & 2 & 12 & 8 \\
\end{tabular}
\end{center}

Observamos que es una permutación: cada valor de 1 a 16 aparece exactamente una vez.
\end{example}

\begin{keyidea}
El costo de $x^5$ en un circuito R1CS es solo \textbf{4 constraints}:
\begin{align}
x^2 &= x \cdot x \quad (1 \text{ constraint}) \\
x^4 &= x^2 \cdot x^2 \quad (1 \text{ constraint}) \\
x^5 &= x^4 \cdot x \quad (1 \text{ constraint})
\end{align}

Más un constraint auxiliar para el routing del resultado. Comparado con los $\sim$768 constraints de un XOR, esto es 200x más eficiente.
\end{keyidea}

%-------------------------------------------
\subsubsection{La Matriz MDS}

\begin{definition}[Matriz MDS]
Una matriz $M \in \F_p^{t \times t}$ es \textbf{MDS} (Maximum Distance Separable) si todas sus submatrices cuadradas son invertibles.

Esto garantiza \textbf{difusión óptima}: cambiar un elemento del estado afecta a todos los demás después de una multiplicación por $M$.
\end{definition}

\begin{intuicion}
La matriz MDS es como un ``mezclador perfecto''. Imagina $t$ líquidos de colores en $t$ vasos. Multiplicar por una matriz MDS es como mezclarlos de tal forma que cada vaso final contiene partes de \textbf{todos} los líquidos originales.

Sin esta mezcla, los bits de entrada podrían ``viajar'' por el estado sin interactuar, debilitando la seguridad.
\end{intuicion}

\begin{example}[Matriz MDS para $t = 3$]
Una construcción común es la matriz de Cauchy. Para $t = 3$:

\[
M = \begin{pmatrix}
2 & 1 & 1 \\
1 & 2 & 1 \\
1 & 1 & 2
\end{pmatrix}
\]

Verificación de que es MDS (todas las submatrices $k \times k$ son invertibles):
\begin{itemize}
    \item Submatrices $1 \times 1$: $\det(2) = 2 \neq 0$ ✓
    \item Submatrices $2 \times 2$: $\det\begin{pmatrix} 2 & 1 \\ 1 & 2 \end{pmatrix} = 3 \neq 0$ ✓
    \item Matriz completa: $\det(M) = 4 \neq 0$ ✓
\end{itemize}
\end{example}

\begin{mathnote}
En la práctica, Poseidon usa matrices MDS construidas a partir de:
\begin{itemize}
    \item \textbf{Matrices de Cauchy}: $M_{ij} = \frac{1}{x_i - y_j}$ para secuencias $x, y$ disjuntas
    \item \textbf{Matrices circulantes}: Eficientes para implementación
\end{itemize}

La elección depende del campo y del tamaño del estado. Las matrices se precomputan y son parte de la especificación de cada instancia de Poseidon.
\end{mathnote}

%-------------------------------------------
\subsubsection{Rondas Completas vs Parciales}

\begin{definition}[Tipos de Rondas]
\textbf{Ronda Completa (Full Round):}
\begin{enumerate}
    \item Sumar constantes: $s_i \leftarrow s_i + c_i$ para todo $i$
    \item Aplicar S-box: $s_i \leftarrow s_i^5$ para \textbf{todo} $i$
    \item Multiplicar por MDS: $\text{state} \leftarrow M \cdot \text{state}$
\end{enumerate}

\textbf{Ronda Parcial (Partial Round):}
\begin{enumerate}
    \item Sumar constantes: $s_i \leftarrow s_i + c_i$ para todo $i$
    \item Aplicar S-box: $s_0 \leftarrow s_0^5$ \textbf{solo al primer elemento}
    \item Multiplicar por MDS: $\text{state} \leftarrow M \cdot \text{state}$
\end{enumerate}
\end{definition}

\begin{keyidea}
Esta es la innovación clave de Poseidon. Las rondas parciales ahorran $(t-1)$ S-boxes por ronda, manteniendo la seguridad gracias a la difusión de la matriz MDS.

\textbf{Costo comparativo para $t = 3$:}
\begin{itemize}
    \item Ronda completa: $3 \times 4 = 12$ constraints (3 S-boxes)
    \item Ronda parcial: $1 \times 4 = 4$ constraints (1 S-box)
\end{itemize}

Ahorro: 67\% por ronda parcial.
\end{keyidea}

\begin{example}[Parámetros típicos para seguridad de 128 bits]
Para campo BN254 (usado en muchos SNARKs) con $t = 3$:
\begin{center}
\begin{tabular}{lccc}
\toprule
\textbf{Parámetro} & \textbf{Símbolo} & \textbf{Valor} & \textbf{Constraints} \\
\midrule
Rondas completas & $R_F$ & 8 & $8 \times 3 \times 4 = 96$ \\
Rondas parciales & $R_P$ & 57 & $57 \times 1 \times 4 = 228$ \\
Constantes de ronda & --- & $65 \times 3 = 195$ & 0 (precomputadas) \\
Multiplicación MDS & --- & 65 matrices & $\sim$18 \\
\midrule
\textbf{Total} & & & $\sim$342 \\
\bottomrule
\end{tabular}
\end{center}

En la práctica, con optimizaciones, se llega a $\sim$250 constraints.
\end{example}

%-------------------------------------------
\subsubsection{Algoritmo Completo}

\begin{example}[Poseidon Hash paso a paso]
\textbf{Parámetros:}
\begin{itemize}
    \item Campo: $\F_p$ con $p = 101$ (ejemplo pequeño)
    \item Estado: $t = 3$
    \item S-box: $\alpha = 5$
    \item Rondas: $R_F = 2$, $R_P = 1$ (simplificado)
    \item Matriz MDS: $M = \begin{pmatrix} 2 & 1 & 1 \\ 1 & 2 & 1 \\ 1 & 1 & 2 \end{pmatrix}$
\end{itemize}

\textbf{Input:} Mensaje $m = (15, 27)$

\textbf{Paso 1: Inicializar estado}
\[ \text{state} = (0, 15, 27) \]
(El primer elemento es el ``capacity'', los otros son el rate)

\textbf{Paso 2: Ronda completa 1}
\begin{align}
\text{ARC}: \quad &(0 + 3, 15 + 7, 27 + 11) = (3, 22, 38) \\
\text{S-box}: \quad &(3^5, 22^5, 38^5) \mod 101 = (42, 79, 85) \\
\text{MDS}: \quad &M \cdot (42, 79, 85)^T = (89, 67, 52)
\end{align}

\textbf{Paso 3: Ronda parcial}
\begin{align}
\text{ARC}: \quad &(89 + 13, 67 + 17, 52 + 19) = (1, 84, 71) \\
\text{S-box (solo } s_0): \quad &(1^5, 84, 71) = (1, 84, 71) \\
\text{MDS}: \quad &M \cdot (1, 84, 71)^T = (56, 68, 57)
\end{align}

\textbf{Paso 4: Ronda completa 2}
\begin{align}
\text{ARC}: \quad &(56 + 23, 68 + 29, 57 + 31) = (79, 97, 88) \\
\text{S-box}: \quad &(79^5, 97^5, 88^5) \mod 101 = (31, 12, 44) \\
\text{MDS}: \quad &M \cdot (31, 12, 44)^T = (18, 99, 32)
\end{align}

\textbf{Salida:} $H(15, 27) = 18$ (primer elemento del estado final)
\end{example}

%-------------------------------------------
\subsubsection{Modos de Operación}

\begin{definition}[Construcción Sponge en Poseidon]
Poseidon típicamente usa una construcción sponge con:
\begin{itemize}
    \item \textbf{Capacity}: Primer(os) elemento(s) del estado (nunca absorbemos/exprimimos aquí)
    \item \textbf{Rate}: Resto del estado (interactúa con input/output)
\end{itemize}

Para $t = 3$ con capacity 1 y rate 2:
\begin{verbatim}
state = [capacity | rate_0 | rate_1]
             ↑          ↑        ↑
         nunca      absorbe/   absorbe/
         expuesto   exprime    exprime
\end{verbatim}
\end{definition}

\textbf{Modos comunes:}

\begin{itemize}
    \item \textbf{Hash}: Absorber mensaje, exprimir un elemento
    \item \textbf{Merkle tree}: $H(left, right) \to root$
    \item \textbf{Commitment}: $C = H(value, randomness)$
    \item \textbf{PRF}: $k = H(secret, nonce)$
\end{itemize}

%-------------------------------------------
\subsubsection{Seguridad}

\begin{keyidea}
Poseidon está diseñado para resistir tres tipos principales de ataques:
\begin{enumerate}
    \item \textbf{Ataques algebraicos}: Resolver el sistema de ecuaciones polinomiales
    \item \textbf{Ataques diferenciales}: Explotar diferencias que se propagan de forma predecible
    \item \textbf{Ataques de interpolación}: Aproximar la función con polinomios de bajo grado
\end{enumerate}
\end{keyidea}

\begin{definition}[Márgenes de Seguridad]
Para Poseidon con 128 bits de seguridad:
\begin{itemize}
    \item \textbf{Ataques algebraicos}: Requieren resolver sistemas de grado $5^{R}$ donde $R$ es el número de rondas con S-boxes activas
    \item \textbf{Ataques diferenciales}: La matriz MDS asegura que diferencias se propaguen a todos los elementos
    \item \textbf{Interpolación}: El grado del polinomio después de todas las rondas excede $2^{128}$
\end{itemize}
\end{definition}

\begin{advertencia}
Poseidon es relativamente nuevo (2019) comparado con SHA-256 (2001). Aunque ha sido bien analizado, algunos criptógrafos recomiendan usar parámetros conservadores (más rondas) en aplicaciones críticas.

En 2023, se publicaron ataques de interpolación mejorados que no rompen Poseidon con parámetros estándar, pero redujeron el margen de seguridad. Los diseñadores respondieron recomendando parámetros actualizados.
\end{advertencia}

%-------------------------------------------
\subsubsection{Otros Hashes Algebraicos}

\paragraph{MiMC (2016)}

\begin{definition}[MiMC]
MiMC es el primer hash algebraico ampliamente usado:
\[ F(x) = (x + k + c_i)^3 \]

aplicado iterativamente por $r$ rondas (típicamente $r \approx 82$ para 128 bits de seguridad).
\end{definition}

\begin{center}
\begin{tabular}{lcc}
\toprule
& \textbf{MiMC} & \textbf{Poseidon} \\
\midrule
S-box & $x^3$ & $x^5$ \\
Rondas & $\sim$82 & $\sim$65 (8 full + 57 partial) \\
Constraints & $\sim$700 & $\sim$250 \\
Estructura & Una sola rama & Estado de $t$ elementos \\
\bottomrule
\end{tabular}
\end{center}

\begin{remark}
MiMC tiene estructura más simple pero es menos eficiente porque no tiene el concepto de rondas parciales.
\end{remark}

\paragraph{Rescue (2019)}

\begin{definition}[Rescue]
Rescue alterna entre dos S-boxes:
\begin{itemize}
    \item $x \mapsto x^\alpha$ (forward)
    \item $x \mapsto x^{1/\alpha}$ (backward, es decir $x^{\alpha^{-1} \mod (p-1)}$)
\end{itemize}

Esto proporciona seguridad simétrica: es igualmente difícil atacar en dirección forward o backward.
\end{definition}

\begin{mathnote}
La inversión $x^{1/\alpha}$ es costosa de calcular pero verifica que el prover realmente hizo el trabajo. Esto es útil en esquemas donde queremos que el prover trabaje más que el verifier (asimetría favorable).
\end{mathnote}

\paragraph{Griffin (2022)}

\begin{definition}[Griffin]
Griffin está optimizado para pruebas STARKs (que usan AIR en lugar de R1CS):
\begin{itemize}
    \item Usa una combinación de operaciones lineales y no lineales
    \item Aprovecha que STARKs tienen estructura de constraints diferente
    \item Especialmente eficiente para campos pequeños (ej: $p = 2^{64} - 2^{32} + 1$)
\end{itemize}
\end{definition}

\paragraph{Anemoi (2022)}

\begin{definition}[Anemoi]
Anemoi es el estado del arte en eficiencia:
\begin{itemize}
    \item S-box: Combinación de $x \mapsto x^{1/5}$ y operaciones cuadráticas
    \item Estructura de ``mariposa'' (butterfly)
    \item Optimizado tanto para R1CS como para AIR
\end{itemize}

Costo: $\sim$120 constraints para seguridad de 128 bits.
\end{definition}

\begin{center}
\begin{tabular}{lccc}
\toprule
\textbf{Hash} & \textbf{Constraints} & \textbf{Rondas} & \textbf{Madurez} \\
\midrule
MiMC & $\sim$700 & 82 & Alta (2016) \\
Rescue & $\sim$400 & 10 & Media (2019) \\
Poseidon & $\sim$250 & 65 & Alta (2019) \\
Griffin & $\sim$150 & 12 & Baja (2022) \\
Anemoi & $\sim$120 & 10 & Baja (2022) \\
\bottomrule
\end{tabular}
\end{center}

\begin{intuicion}
El campo de hashes ZK-friendly evoluciona rápidamente. La tendencia es hacia diseños más eficientes, pero con el trade-off de menor tiempo de análisis criptográfico. Poseidon representa un buen balance: significativamente mejor que MiMC, con 5+ años de escrutinio por la comunidad.
\end{intuicion}

%-------------------------------------------
\subsubsection{Implementación Práctica}

\begin{example}[Poseidon en Circom]
\begin{verbatim}
pragma circom 2.0.0;

include "poseidon.circom";

template HashExample() {
    signal input a;
    signal input b;
    signal output hash;
    
    component poseidon = Poseidon(2);  // 2 inputs
    poseidon.inputs[0] <== a;
    poseidon.inputs[1] <== b;
    hash <== poseidon.out;
}

component main = HashExample();
\end{verbatim}

Este circuito produce $\sim$250 constraints y calcula $H(a, b)$.
\end{example}

\begin{example}[Poseidon en Rust (Arkworks)]
\begin{verbatim}
use ark_crypto_primitives::sponge::poseidon::{
    PoseidonConfig, PoseidonSponge
};
use ark_crypto_primitives::sponge::CryptographicSponge;
use ark_bn254::Fr;  // Campo del BN254

fn hash_two_elements(a: Fr, b: Fr) -> Fr {
    let params = PoseidonConfig::new(...);  // Parámetros precomputados
    let mut sponge = PoseidonSponge::new(&params);
    
    sponge.absorb(&[a, b]);
    sponge.squeeze_field_elements(1)[0]
}
\end{verbatim}
\end{example}

%-------------------------------------------
\subsubsection{Parámetros para Diferentes Curvas}

\begin{center}
\begin{tabular}{lcccc}
\toprule
\textbf{Curva/Campo} & \textbf{$t$} & \textbf{$R_F$} & \textbf{$R_P$} & \textbf{Uso} \\
\midrule
BN254 & 3 & 8 & 57 & Ethereum, zkSync \\
BLS12-381 & 3 & 8 & 57 & Filecoin, Zcash \\
Pallas & 3 & 8 & 56 & Zcash Orchard \\
Vesta & 3 & 8 & 56 & Zcash Orchard \\
Goldilocks ($p = 2^{64} - 2^{32} + 1$) & 12 & 8 & 22 & Plonky2 \\
\bottomrule
\end{tabular}
\end{center}

\begin{remark}
Cada instancia de Poseidon requiere parámetros específicos: constantes de ronda y matriz MDS. Estos se generan de forma determinista a partir de una semilla, permitiendo verificación independiente.
\end{remark}

%-------------------------------------------
\subsubsection{Aplicación en Banco ZK}

\begin{bancozk}
\textbf{Poseidon es el hash principal de Banco ZK.} Lo usamos en:

\textbf{1. Árboles de Merkle del estado}
\begin{itemize}
    \item Profundidad: 20 niveles ($2^{20} \approx 1M$ notas)
    \item Cada nivel: 1 hash Poseidon
    \item Costo total del path: $20 \times 250 = 5,000$ constraints
    \item Comparación con SHA-256: $20 \times 25,000 = 500,000$ constraints (100x más)
\end{itemize}

\textbf{2. Compromisos a notas}
\begin{align}
\text{note\_commitment} &= \text{Poseidon}(\text{value}, \text{owner\_pubkey}, \text{randomness})
\end{align}

\textbf{3. Derivación de nullifiers}
\begin{align}
\text{nullifier} &= \text{Poseidon}(\text{note\_commitment}, \text{spending\_key})
\end{align}

\textbf{4. Por qué elegimos Poseidon sobre alternativas:}
\begin{center}
\begin{tabular}{lcc}
\toprule
\textbf{Criterio} & \textbf{Poseidon} & \textbf{Anemoi} \\
\midrule
Constraints & 250 & 120 \\
Años de análisis & 5+ & 2 \\
Librerías maduras & Muchas & Pocas \\
Auditorías & Múltiples & Limitadas \\
\midrule
\textbf{Decisión} & \checkmark Elegido & Para futuro \\
\bottomrule
\end{tabular}
\end{center}

En Banco ZK priorizamos seguridad probada sobre eficiencia marginal. Cuando Anemoi tenga más escrutinio, consideraremos migrarlo.

\textbf{5. Optimizaciones implementadas:}
\begin{itemize}
    \item Precomputación de constantes de ronda
    \item Matrices MDS optimizadas para multiplicación rápida
    \item Caché de estados intermedios para pruebas batch
\end{itemize}
\end{bancozk}

%-------------------------------------------
\subsubsection{Comparación Final: El Costo Real}

\begin{example}[Impacto en tiempo de prueba]
Para un circuito de transferencia privada típico:

\begin{center}
\begin{tabular}{lccc}
\toprule
\textbf{Componente} & \textbf{Con SHA-256} & \textbf{Con Poseidon} & \textbf{Ahorro} \\
\midrule
Merkle path (20 niveles) & 500,000 & 5,000 & 99\% \\
2 nullifiers & 50,000 & 500 & 99\% \\
2 commitments & 50,000 & 500 & 99\% \\
Lógica adicional & 10,000 & 10,000 & 0\% \\
\midrule
\textbf{Total} & 610,000 & 16,000 & 97\% \\
\bottomrule
\end{tabular}
\end{center}

\textbf{Tiempo de generación de prueba} (estimado en laptop):
\begin{itemize}
    \item Con SHA-256: $\sim$30 segundos
    \item Con Poseidon: $\sim$1 segundo
\end{itemize}

Esta diferencia de 30x hace que las pruebas sean prácticas para usuarios móviles.
\end{example}
% hasta aca poseidon 
% ----- CURVAS ELIPTICAS -----
\section{Curvas Elipticas}

% ALE DESDE Grupo
%===========================================
% SUBSECCIÓN: CURVAS ELÍPTICAS - DEFINICIÓN Y LEY DE GRUPO
% Usa tus entornos: definition, example, keyidea, bancozk, 
%                   mathnote, intuicion, advertencia, remark
%===========================================

\subsection{Definición y Ley de Grupo}

\begin{definition}[Curva elíptica]
Sobre un campo $K$, una curva elíptica en forma de Weierstrass es:
\[
E: y^2 = x^3 + ax + b \quad \text{con } 4a^3 + 27b^2 \neq 0
\]
\end{definition}

Los puntos de $E$ forman un grupo abeliano con:
\begin{itemize}
    \item Elemento neutro: punto en el infinito $\mathcal{O}$
    \item Suma geométrica: línea por $P, Q$ intersecta en $R$, reflejar
    \item Inverso: $-P = (x, -y)$
\end{itemize}

%-------------------------------------------
\subsubsection{Contexto Histórico}

\begin{intuicion}
Las curvas elípticas tienen una historia fascinante que abarca más de 2000 años. Originalmente estudiadas en el contexto de calcular la longitud del arco de una elipse (de ahí el nombre ``elípticas''), hoy son fundamentales para la criptografía moderna.
\end{intuicion}

\textbf{Línea temporal:}
\begin{center}
\begin{tabular}{ll}
\toprule
\textbf{Época} & \textbf{Desarrollo} \\
\midrule
Siglo III a.C. & Diofanto estudia ecuaciones similares \\
Siglo XVII & Fermat, Newton estudian curvas cúbicas \\
1901 & Poincaré descubre la estructura de grupo \\
1985 & Koblitz y Miller proponen ECC independientemente \\
2009 & Bitcoin usa secp256k1 \\
2016 & Ethereum adopta la misma curva \\
\bottomrule
\end{tabular}
\end{center}

\begin{remark}
Neal Koblitz y Victor Miller, trabajando independientemente, propusieron usar curvas elípticas para criptografía en 1985. Su insight fue que el problema del logaritmo discreto en curvas elípticas (ECDLP) era más difícil que en grupos multiplicativos, permitiendo claves más cortas con la misma seguridad.
\end{remark}

%-------------------------------------------
\subsubsection{La Ecuación de Weierstrass}

\begin{definition}[Forma de Weierstrass Generalizada]
La forma más general de una curva elíptica es:
\[
E: y^2 + a_1 xy + a_3 y = x^3 + a_2 x^2 + a_4 x + a_6
\]

Para campos de característica $\neq 2, 3$, esto se simplifica a la \textbf{forma corta de Weierstrass}:
\[
E: y^2 = x^3 + ax + b
\]
\end{definition}

\begin{mathnote}
La condición $\Delta = 4a^3 + 27b^2 \neq 0$ es el \textbf{discriminante}. Garantiza que la curva es ``suave'' (no tiene puntos singulares como cúspides o nodos).

Geométricamente:
\begin{itemize}
    \item Si $\Delta = 0$: la curva tiene un punto donde se ``cruza'' consigo misma o tiene una ``punta''
    \item Si $\Delta \neq 0$: la curva es suave, lo que permite definir tangentes en todos los puntos
\end{itemize}
\end{mathnote}

\begin{example}[Curvas válidas e inválidas]
\textbf{Curva válida:} $E: y^2 = x^3 - 3x + 3$
\begin{align}
a &= -3, \quad b = 3 \\
\Delta &= 4(-3)^3 + 27(3)^2 = 4(-27) + 27(9) = -108 + 243 = 135 \neq 0 \quad \checkmark
\end{align}

\textbf{Curva inválida (singular):} $E: y^2 = x^3$
\begin{align}
a &= 0, \quad b = 0 \\
\Delta &= 4(0)^3 + 27(0)^2 = 0 \quad \text{(tiene cúspide en el origen)}
\end{align}

\textbf{Curva inválida (nodo):} $E: y^2 = x^3 + x^2$
\begin{align}
a &= 0, \quad b = 0 \quad \text{(forma } y^2 = x^2(x+1) \text{)}
\end{align}
El origen es un punto donde la curva se cruza consigo misma.
\end{example}

%-------------------------------------------
\subsubsection{Visualización Geométrica}

\begin{intuicion}
Sobre los números reales $\mathbb{R}$, una curva elíptica tiene una forma característica simétrica respecto al eje $x$. Dependiendo de los coeficientes, puede tener una o dos componentes conexas.
\end{intuicion}

\begin{center}
\begin{tikzpicture}[scale=0.8]
    % Ejes
    \draw[->] (-3, 0) -- (3.5, 0) node[right] {$x$};
    \draw[->] (0, -3) -- (0, 3) node[above] {$y$};
    
    % Curva y² = x³ - x (dos componentes)
    % Componente izquierda (aproximada)
    \draw[thick, blue, domain=-1:0, samples=50] plot (\x, {sqrt(\x*\x*\x - \x)});
    \draw[thick, blue, domain=-1:0, samples=50] plot (\x, {-sqrt(\x*\x*\x - \x)});
    
    % Componente derecha
    \draw[thick, blue, domain=1:2.5, samples=50] plot (\x, {sqrt(\x*\x*\x - \x)});
    \draw[thick, blue, domain=1:2.5, samples=50] plot (\x, {-sqrt(\x*\x*\x - \x)});
    
    % Etiqueta
    \node at (2, -2.5) {$y^2 = x^3 - x$};
\end{tikzpicture}
\hspace{1cm}
\begin{tikzpicture}[scale=0.8]
    % Ejes
    \draw[->] (-2, 0) -- (3.5, 0) node[right] {$x$};
    \draw[->] (0, -3) -- (0, 3) node[above] {$y$};
    
    % Curva y² = x³ - x + 1 (una componente)
    \draw[thick, red, domain=-1.3:2.3, samples=100] plot (\x, {sqrt(\x*\x*\x - \x + 1)});
    \draw[thick, red, domain=-1.3:2.3, samples=100] plot (\x, {-sqrt(\x*\x*\x - \x + 1)});
    
    % Etiqueta
    \node at (2, -2.5) {$y^2 = x^3 - x + 1$};
\end{tikzpicture}
\end{center}

\begin{remark}
En criptografía \textbf{no} trabajamos sobre $\mathbb{R}$, sino sobre campos finitos $\mathbb{F}_p$. En estos campos, la ``curva'' es realmente un conjunto discreto de puntos que satisfacen la ecuación módulo $p$. La visualización geométrica sobre $\mathbb{R}$ es útil para intuición, pero los cálculos reales son en aritmética modular.
\end{remark}

%-------------------------------------------
\subsubsection{El Punto en el Infinito}

\begin{definition}[Punto en el Infinito]
El punto $\mathcal{O}$ (punto en el infinito) es un punto ``especial'' que no tiene coordenadas finitas. Se puede pensar como el punto donde todas las líneas verticales se encuentran.

Formalmente, $\mathcal{O}$ es el elemento neutro del grupo:
\[
P + \mathcal{O} = \mathcal{O} + P = P \quad \text{para todo punto } P
\]
\end{definition}

\begin{intuicion}
Imagina que extiendes el plano hacia arriba y hacia abajo al infinito. Todas las líneas verticales, que parecen paralelas, eventualmente se ``encuentran'' en un único punto en el infinito $\mathcal{O}$.

En coordenadas proyectivas, $\mathcal{O} = [0 : 1 : 0]$. En la práctica, lo representamos con un flag especial o coordenadas $(0, 0)$ con una marca.
\end{intuicion}

%-------------------------------------------
\subsubsection{La Ley de Grupo: Suma de Puntos}

\begin{keyidea}
Los puntos de una curva elíptica forman un \textbf{grupo abeliano} bajo una operación de ``suma'' definida geométricamente. Esta estructura algebraica es lo que hace a las curvas elípticas útiles para criptografía.
\end{keyidea}

\begin{definition}[Suma de Puntos - Definición Geométrica]
Dados dos puntos $P$ y $Q$ en la curva $E$:
\begin{enumerate}
    \item Trazar la línea $\ell$ que pasa por $P$ y $Q$
    \item Esta línea intersecta la curva en un tercer punto $R'$
    \item Reflejar $R'$ sobre el eje $x$ para obtener $R = P + Q$
\end{enumerate}

Casos especiales:
\begin{itemize}
    \item Si $P = Q$: usar la línea tangente a la curva en $P$
    \item Si $P = -Q$ (misma $x$, opuesta $y$): la línea es vertical, $P + Q = \mathcal{O}$
\end{itemize}
\end{definition}

\begin{center}
\begin{tikzpicture}[scale=1.2]
    % Ejes
    \draw[->] (-1, 0) -- (4, 0) node[right] {$x$};
    \draw[->] (0, -2.5) -- (0, 2.5) node[above] {$y$};
    
    % Curva (aproximación de y² = x³ - 2x + 2)
    \draw[thick, blue, domain=0.3:3.2, samples=100] plot (\x, {sqrt(\x*\x*\x - 2*\x + 2)});
    \draw[thick, blue, domain=0.3:3.2, samples=100] plot (\x, {-sqrt(\x*\x*\x - 2*\x + 2)});
    
    % Puntos P, Q
    \fill[red] (0.7, 1.1) circle (2pt) node[above left] {$P$};
    \fill[red] (1.8, 1.7) circle (2pt) node[above] {$Q$};
    
    % Punto R' (intersección)
    \fill[orange] (2.8, -2.1) circle (2pt) node[below right] {$R'$};
    
    % Punto R = P + Q
    \fill[green!50!black] (2.8, 2.1) circle (2pt) node[above right] {$R = P + Q$};
    
    % Línea por P y Q
    \draw[dashed, thick] (0.3, 0.86) -- (3.2, 2.3);
    
    % Línea vertical de reflexión
    \draw[dotted, thick] (2.8, -2.1) -- (2.8, 2.1);
    
\end{tikzpicture}
\end{center}

%-------------------------------------------
\subsubsection{Fórmulas Algebraicas Explícitas}

\begin{definition}[Fórmulas de Suma]
Para $P = (x_1, y_1)$ y $Q = (x_2, y_2)$ en $E: y^2 = x^3 + ax + b$:

\textbf{Caso 1: $P \neq Q$ (suma de puntos distintos)}
\begin{align}
\lambda &= \frac{y_2 - y_1}{x_2 - x_1} \\
x_3 &= \lambda^2 - x_1 - x_2 \\
y_3 &= \lambda(x_1 - x_3) - y_1
\end{align}

\textbf{Caso 2: $P = Q$ (duplicación)}
\begin{align}
\lambda &= \frac{3x_1^2 + a}{2y_1} \\
x_3 &= \lambda^2 - 2x_1 \\
y_3 &= \lambda(x_1 - x_3) - y_1
\end{align}

\textbf{Casos especiales:}
\begin{itemize}
    \item Si $x_1 = x_2$ pero $y_1 = -y_2$: $P + Q = \mathcal{O}$
    \item Si $P = \mathcal{O}$: $P + Q = Q$
    \item Si $Q = \mathcal{O}$: $P + Q = P$
    \item Si $P = Q$ y $y_1 = 0$: $2P = \mathcal{O}$ (tangente vertical)
\end{itemize}
\end{definition}

\begin{mathnote}
El valor $\lambda$ es la \textbf{pendiente} de la línea:
\begin{itemize}
    \item En el caso $P \neq Q$: pendiente de la secante entre $P$ y $Q$
    \item En el caso $P = Q$: pendiente de la tangente en $P$, obtenida por diferenciación implícita de $y^2 = x^3 + ax + b$:
    \[
    2y \cdot \frac{dy}{dx} = 3x^2 + a \implies \frac{dy}{dx} = \frac{3x^2 + a}{2y}
    \]
\end{itemize}
\end{mathnote}

%-------------------------------------------
\subsubsection{Ejemplo Completo sobre Campo Finito}

\begin{example}[Suma de puntos en $E(\mathbb{F}_{17})$]
Consideremos la curva $E: y^2 = x^3 + 2x + 2$ sobre $\mathbb{F}_{17}$.

\textbf{Paso 1: Verificar que la curva es válida}
\begin{align}
\Delta &= 4(2)^3 + 27(2)^2 = 32 + 108 = 140 \equiv 4 \pmod{17} \neq 0 \quad \checkmark
\end{align}

\textbf{Paso 2: Encontrar puntos en la curva}

Probamos valores de $x$ y verificamos si $x^3 + 2x + 2$ es un cuadrado en $\mathbb{F}_{17}$:

\begin{center}
\begin{tabular}{c|c|c|c}
$x$ & $x^3 + 2x + 2$ & ¿Cuadrado? & Puntos $(x, y)$ \\
\hline
0 & 2 & $2 \equiv 6^2 = 36 \equiv 2$? No & --- \\
1 & $1 + 2 + 2 = 5$ & $5$: $6^2 = 2, 7^2 = 15$, no & --- \\
2 & $8 + 4 + 2 = 14$ & $14$: no & --- \\
3 & $27 + 6 + 2 = 35 \equiv 1$ & $1 = 1^2$ ✓ & $(3, 1), (3, 16)$ \\
5 & $125 + 10 + 2 = 137 \equiv 1$ & $1 = 1^2$ ✓ & $(5, 1), (5, 16)$ \\
6 & $216 + 12 + 2 = 230 \equiv 9$ & $9 = 3^2$ ✓ & $(6, 3), (6, 14)$ \\
\end{tabular}
\end{center}

\textbf{Paso 3: Sumar $P = (3, 1)$ y $Q = (6, 3)$}
\begin{align}
\lambda &= \frac{3 - 1}{6 - 3} = \frac{2}{3} \pmod{17}
\end{align}

Necesitamos $3^{-1} \mod 17$. Como $3 \cdot 6 = 18 \equiv 1 \pmod{17}$, tenemos $3^{-1} = 6$.
\begin{align}
\lambda &= 2 \cdot 6 = 12 \pmod{17}
\end{align}

Ahora calculamos las coordenadas:
\begin{align}
x_3 &= \lambda^2 - x_1 - x_2 = 12^2 - 3 - 6 = 144 - 9 = 135 \equiv 16 \pmod{17} \\
y_3 &= \lambda(x_1 - x_3) - y_1 = 12(3 - 16) - 1 = 12(-13) - 1 \\
    &= 12 \cdot 4 - 1 = 48 - 1 = 47 \equiv 13 \pmod{17}
\end{align}

\textbf{Resultado:} $P + Q = (3, 1) + (6, 3) = (16, 13)$

\textbf{Verificación:} ¿Está $(16, 13)$ en la curva?
\begin{align}
y^2 &= 13^2 = 169 \equiv 16 \pmod{17} \\
x^3 + 2x + 2 &= 16^3 + 2(16) + 2 = 4096 + 32 + 2 = 4130 \equiv 16 \pmod{17} \quad \checkmark
\end{align}
\end{example}

%-------------------------------------------
\subsubsection{Duplicación de Puntos}

\begin{example}[Duplicar $P = (3, 1)$ en $E: y^2 = x^3 + 2x + 2$ sobre $\mathbb{F}_{17}$]

Usamos la fórmula de duplicación con $a = 2$:
\begin{align}
\lambda &= \frac{3x_1^2 + a}{2y_1} = \frac{3(3)^2 + 2}{2(1)} = \frac{27 + 2}{2} = \frac{29}{2} \equiv \frac{12}{2} \pmod{17}
\end{align}

Necesitamos $2^{-1} \mod 17$. Como $2 \cdot 9 = 18 \equiv 1$, tenemos $2^{-1} = 9$.
\begin{align}
\lambda &= 12 \cdot 9 = 108 \equiv 6 \pmod{17}
\end{align}

Calculamos:
\begin{align}
x_3 &= \lambda^2 - 2x_1 = 6^2 - 2(3) = 36 - 6 = 30 \equiv 13 \pmod{17} \\
y_3 &= \lambda(x_1 - x_3) - y_1 = 6(3 - 13) - 1 = 6(-10) - 1 = -61 \equiv 7 \pmod{17}
\end{align}

\textbf{Resultado:} $2P = (13, 7)$

\textbf{Verificación:}
\begin{align}
y^2 &= 7^2 = 49 \equiv 15 \pmod{17} \\
x^3 + 2x + 2 &= 13^3 + 26 + 2 = 2197 + 28 = 2225 \equiv 15 \pmod{17} \quad \checkmark
\end{align}
\end{example}

%-------------------------------------------
\subsubsection{Propiedades del Grupo}

\begin{keyidea}
Las curvas elípticas forman un grupo abeliano. Esto significa que la operación de suma satisface:
\begin{enumerate}
    \item \textbf{Clausura}: Si $P, Q \in E$, entonces $P + Q \in E$
    \item \textbf{Asociatividad}: $(P + Q) + R = P + (Q + R)$
    \item \textbf{Elemento neutro}: Existe $\mathcal{O}$ tal que $P + \mathcal{O} = P$
    \item \textbf{Inversos}: Para cada $P$, existe $-P$ tal que $P + (-P) = \mathcal{O}$
    \item \textbf{Conmutatividad}: $P + Q = Q + P$
\end{enumerate}
\end{keyidea}

\begin{definition}[Inverso de un Punto]
Para $P = (x, y)$, su inverso es:
\[
-P = (x, -y)
\]

En un campo finito $\mathbb{F}_p$, esto es:
\[
-P = (x, p - y)
\]
\end{definition}

\begin{example}[Inversos en $\mathbb{F}_{17}$]
Si $P = (3, 1)$, entonces:
\[
-P = (3, 17 - 1) = (3, 16)
\]

Verificación: $P + (-P) = (3, 1) + (3, 16)$

Como $x_1 = x_2 = 3$ pero $y_1 = 1 \neq -y_2 = -16 \equiv 1$... 

Espera, $y_2 = 16 \equiv -1 \pmod{17}$, y $-y_2 = -16 \equiv 1 = y_1$. No, $y_1 \neq -y_2$ en este caso...

Realmente: $y_1 = 1$ y $-y_2 = -16 \equiv 1 \pmod{17}$, entonces $y_1 = -y_2$ ✗

Corrección: $-y_2 = 17 - 16 = 1 = y_1$, así que... los puntos SON opuestos cuando $y_2 = p - y_1$.

Por lo tanto: $(3, 1) + (3, 16) = \mathcal{O}$ ✓
\end{example}

%-------------------------------------------
\subsubsection{Multiplicación Escalar}

\begin{definition}[Multiplicación Escalar]
Para un punto $P$ y un entero $n$, definimos:
\[
nP = \underbrace{P + P + \cdots + P}_{n \text{ veces}}
\]

También:
\begin{itemize}
    \item $0 \cdot P = \mathcal{O}$
    \item $(-n) \cdot P = n \cdot (-P)$
\end{itemize}
\end{definition}

\begin{keyidea}
La multiplicación escalar es la operación fundamental en criptografía de curvas elípticas:
\begin{itemize}
    \item \textbf{Fácil}: Dado $n$ y $P$, calcular $Q = nP$ (usando double-and-add)
    \item \textbf{Difícil}: Dado $P$ y $Q$, encontrar $n$ tal que $Q = nP$ (ECDLP)
\end{itemize}

Esta asimetría es la base de la seguridad de ECC.
\end{keyidea}

\begin{definition}[Algoritmo Double-and-Add]
Para calcular $nP$ eficientemente:
\begin{verbatim}
function scalar_mult(n, P):
    if n == 0: return O
    if n < 0: return scalar_mult(-n, -P)
    
    R = O
    Q = P
    while n > 0:
        if n es impar:
            R = R + Q
        Q = 2*Q  # Duplicación
        n = n // 2
    return R
\end{verbatim}

Complejidad: $O(\log n)$ operaciones de grupo.
\end{definition}

\begin{example}[Calcular $7P$ usando double-and-add]
$7 = 111_2$ (en binario)

\begin{align}
7P &= 1 \cdot P + 1 \cdot 2P + 1 \cdot 4P \\
   &= P + 2P + 4P
\end{align}

Pasos:
\begin{enumerate}
    \item $2P$ = duplicar $P$
    \item $4P = 2(2P)$ = duplicar $2P$
    \item $P + 2P = 3P$
    \item $3P + 4P = 7P$
\end{enumerate}

Total: 2 duplicaciones + 2 sumas = 4 operaciones (vs. 6 sumas ingenuas)
\end{example}

%-------------------------------------------
\subsubsection{El Orden del Grupo}

\begin{definition}[Orden del Grupo]
El \textbf{orden} de $E(\mathbb{F}_p)$, denotado $\#E(\mathbb{F}_p)$ o $|E(\mathbb{F}_p)|$, es el número total de puntos en la curva, incluyendo $\mathcal{O}$.
\end{definition}

\begin{theorem}[Teorema de Hasse]
Para una curva elíptica $E$ sobre $\mathbb{F}_p$:
\[
|p + 1 - \#E(\mathbb{F}_p)| \leq 2\sqrt{p}
\]

Equivalentemente:
\[
p + 1 - 2\sqrt{p} \leq \#E(\mathbb{F}_p) \leq p + 1 + 2\sqrt{p}
\]
\end{theorem}

\begin{intuicion}
El teorema de Hasse dice que el número de puntos está ``cerca'' de $p + 1$. Para $p$ grande, el número de puntos es aproximadamente $p$, con una desviación de como mucho $2\sqrt{p}$.

Por ejemplo, para $p \approx 2^{256}$:
\begin{itemize}
    \item $\#E \approx 2^{256} \pm 2^{129}$
    \item El ``error'' es insignificante en términos relativos
\end{itemize}
\end{intuicion}

\begin{definition}[Orden de un Punto]
El \textbf{orden} de un punto $P$, denotado $\text{ord}(P)$, es el menor entero positivo $n$ tal que:
\[
nP = \mathcal{O}
\]

Por el teorema de Lagrange, $\text{ord}(P)$ divide a $\#E(\mathbb{F}_p)$.
\end{definition}

%-------------------------------------------
\subsubsection{Curvas Usadas en Criptografía}

\begin{center}
\begin{tabular}{llll}
\toprule
\textbf{Curva} & \textbf{Ecuación} & \textbf{Campo} & \textbf{Uso} \\
\midrule
secp256k1 & $y^2 = x^3 + 7$ & $\mathbb{F}_p$, $p = 2^{256} - 2^{32} - 977$ & Bitcoin, Ethereum \\
P-256 & $y^2 = x^3 - 3x + b$ & $\mathbb{F}_p$, $p = 2^{256} - 2^{224} + ...$ & TLS, NIST \\
Curve25519 & $y^2 = x^3 + 486662x^2 + x$ & $\mathbb{F}_p$, $p = 2^{255} - 19$ & Signal, SSH \\
BN254 & Barreto-Naehrig & $\mathbb{F}_p$, $p \approx 2^{254}$ & SNARKs \\
BLS12-381 & Barreto-Lynn-Scott & $\mathbb{F}_p$, $p \approx 2^{381}$ & Ethereum 2.0 \\
\bottomrule
\end{tabular}
\end{center}

\begin{remark}
\textbf{secp256k1} (usada en Bitcoin/Ethereum) tiene la forma simple $y^2 = x^3 + 7$ porque:
\begin{itemize}
    \item $a = 0$ simplifica las fórmulas (menos multiplicaciones)
    \item $b = 7$ es el menor valor que da una curva segura con $a = 0$
    \item Los parámetros son ``nothing-up-my-sleeve'': claramente no hay backdoor
\end{itemize}
\end{remark}

%-------------------------------------------
\subsubsection{Representación de Puntos}

\begin{definition}[Coordenadas Afines vs Proyectivas]
\textbf{Coordenadas afines:} $(x, y)$ donde $y^2 = x^3 + ax + b$

\textbf{Coordenadas proyectivas:} $(X : Y : Z)$ donde $Y^2 Z = X^3 + aXZ^2 + bZ^3$

Relación: $(x, y) \leftrightarrow (x : y : 1)$ y $(X : Y : Z) \leftrightarrow (X/Z, Y/Z)$ si $Z \neq 0$

El punto en el infinito: $\mathcal{O} = (0 : 1 : 0)$
\end{definition}

\begin{keyidea}
Las coordenadas proyectivas evitan divisiones (inversiones modulares), que son costosas. En su lugar, se trabaja con multiplicaciones y se hace una sola inversión al final para convertir a coordenadas afines.

En implementaciones optimizadas, esto puede dar speedups de 10x o más.
\end{keyidea}

\begin{definition}[Coordenadas Jacobianas]
Una variante popular: $(X : Y : Z)$ donde $(x, y) = (X/Z^2, Y/Z^3)$

La ecuación de la curva se transforma en:
\[
Y^2 = X^3 + aXZ^4 + bZ^6
\]

Las fórmulas de duplicación son especialmente eficientes en esta representación.
\end{definition}

%-------------------------------------------
\subsubsection{Compresión de Puntos}

\begin{definition}[Compresión de Puntos]
Un punto $(x, y)$ puede representarse de forma comprimida como $(x, \text{signo}(y))$, donde el signo indica cuál de las dos raíces cuadradas elegir.

\begin{itemize}
    \item \textbf{Sin comprimir}: $04 \,||\, x \,||\, y$ (65 bytes para curvas de 256 bits)
    \item \textbf{Comprimido}: $02/03 \,||\, x$ (33 bytes)
    \begin{itemize}
        \item $02$ si $y$ es par
        \item $03$ si $y$ es impar
    \end{itemize}
\end{itemize}
\end{definition}

\begin{example}[Compresión en secp256k1]
Clave pública sin comprimir (65 bytes):
\begin{verbatim}
04 79BE667EF9DCBBAC55A06295CE870B07029BFCDB2DCE28D959F2815B
   16F81798 483ADA7726A3C4655DA4FBFC0E1108A8FD17B448A6855419
   9C47D08FFB10D4B8
\end{verbatim}

Clave pública comprimida (33 bytes):
\begin{verbatim}
02 79BE667EF9DCBBAC55A06295CE870B07029BFCDB2DCE28D959F2815B
   16F81798
\end{verbatim}

El $02$ indica que $y$ es par. Para descomprimir, se calcula $y = \sqrt{x^3 + 7}$ y se elige la raíz par.
\end{example}

%-------------------------------------------
\subsubsection{Aplicación en Banco ZK}

\begin{bancozk}
Las curvas elípticas son fundamentales en Banco ZK:

\textbf{1. Firmas digitales (secp256k1)}
\begin{itemize}
    \item Compatible con Ethereum para verificación de transacciones
    \item Usuarios firman con sus claves existentes de Ethereum
\end{itemize}

\textbf{2. Curvas para SNARKs (BN254, BLS12-381)}
\begin{itemize}
    \item Verificación de pruebas ZK requiere operaciones de pairing
    \item BN254: ~190 bits de seguridad, muy eficiente
    \item BLS12-381: 128 bits de seguridad, usado en Ethereum 2.0
\end{itemize}

\textbf{3. Compromisos de Pedersen}
\[
C = v \cdot G + r \cdot H
\]
donde $G$ y $H$ son generadores de la curva. Usamos esto para ocultar valores de transacciones.

\textbf{4. Ciclos de curvas (Pallas/Vesta)}

Para pruebas recursivas eficientes, necesitamos ``ciclos'' de curvas donde el orden de una curva es el tamaño del campo de la otra:
\begin{itemize}
    \item Pallas: $E_1(\mathbb{F}_p)$ con $\#E_1 = q$
    \item Vesta: $E_2(\mathbb{F}_q)$ con $\#E_2 = p$
\end{itemize}

Esto permite verificar una prueba sobre $E_1$ dentro de un circuito sobre $E_2$ eficientemente.

\textbf{Decisión de diseño:}
\begin{center}
\begin{tabular}{lll}
\toprule
\textbf{Componente} & \textbf{Curva} & \textbf{Razón} \\
\midrule
Firmas de usuario & secp256k1 & Compatibilidad Ethereum \\
Verificación SNARK & BN254 & Eficiencia, soporte EVM \\
Pruebas recursivas & Pallas/Vesta & Ciclo nativo \\
Compromisos & Misma curva del SNARK & Eficiencia \\
\bottomrule
\end{tabular}
\end{center}
\end{bancozk}
% ALE FIN Grupo

% ALE Comienza secp256k1 
%===========================================
% SUBSECCIÓN: secp256k1 - VERSIÓN EXPANDIDA
% Usa tus entornos: definition, example, keyidea, bancozk, 
%                   mathnote, intuicion, advertencia, remark
%===========================================

\subsection{secp256k1}

Curva usada en Bitcoin y Ethereum.
\begin{itemize}
    \item Ecuación: $y^2 = x^3 + 7$ sobre $\Fp$
    \item $p = 2^{256} - 2^{32} - 977$
    \item Orden: $n \approx 2^{256}$
    \item Curva de Koblitz: parámetros especiales para eficiencia
\end{itemize}

%-------------------------------------------
\subsubsection{Origen y Nombre}

\begin{intuicion}
El nombre ``secp256k1'' sigue la convención de SECG (Standards for Efficient Cryptography Group):
\begin{itemize}
    \item \textbf{sec}: Standards for Efficient Cryptography
    \item \textbf{p}: Campo primo (vs. ``2'' para campos binarios)
    \item \textbf{256}: Tamaño del campo en bits
    \item \textbf{k}: Curva de Koblitz (parámetros especiales)
    \item \textbf{1}: Primera curva de este tipo en el estándar
\end{itemize}
\end{intuicion}

\begin{remark}
Antes de Bitcoin, secp256k1 era relativamente poco usada. La mayoría de aplicaciones preferían secp256r1 (P-256), recomendada por NIST. Satoshi Nakamoto eligió secp256k1 en 2008, posiblemente porque:
\begin{enumerate}
    \item Sus parámetros son ``nothing-up-my-sleeve'' (no hay sospecha de backdoors)
    \item La forma $y^2 = x^3 + 7$ (con $a = 0$) permite optimizaciones
    \item No fue generada por NIST/NSA (después de las revelaciones de Snowden, esto se valoró más)
\end{enumerate}
\end{remark}

%-------------------------------------------
\subsubsection{Parámetros Completos}

\begin{definition}[Parámetros de secp256k1]
La curva secp256k1 está definida por:

\textbf{Ecuación:}
\[ E: y^2 = x^3 + 7 \]

\textbf{Campo base} $\Fp$ con:
\begin{align}
p &= 2^{256} - 2^{32} - 2^9 - 2^8 - 2^7 - 2^6 - 2^4 - 1 \\
  &= 2^{256} - 2^{32} - 977 \\
  &= \texttt{FFFFFFFF FFFFFFFF FFFFFFFF FFFFFFFF} \\
  &\phantom{= }\texttt{FFFFFFFF FFFFFFFF FFFFFFFE FFFFFC2F}
\end{align}

\textbf{Orden del grupo} (número de puntos):
\begin{align}
n &= \texttt{FFFFFFFF FFFFFFFF FFFFFFFF FFFFFFFE} \\
  &\phantom{= }\texttt{BAAEDCE6 AF48A03B BFD25E8C D0364141}
\end{align}

\textbf{Cofactor:}
\[ h = 1 \]

\textbf{Generador} $G = (G_x, G_y)$:
\begin{align}
G_x &= \texttt{79BE667E F9DCBBAC 55A06295 CE870B07} \\
    &\phantom{= }\texttt{029BFCDB 2DCE28D9 59F2815B 16F81798} \\
G_y &= \texttt{483ADA77 26A3C465 5DA4FBFC 0E1108A8} \\
    &\phantom{= }\texttt{FD17B448 A6855419 9C47D08F FB10D4B8}
\end{align}
\end{definition}

\begin{mathnote}
Verificación del discriminante:
\begin{align}
a &= 0, \quad b = 7 \\
\Delta &= 4a^3 + 27b^2 = 0 + 27 \cdot 49 = 1323 \neq 0 \quad \checkmark
\end{align}

La curva es no singular, como se requiere.
\end{mathnote}

%-------------------------------------------
\subsubsection{¿Por Qué Esta Forma Especial?}

\begin{keyidea}
La ecuación $y^2 = x^3 + 7$ tiene $a = 0$, lo que elimina el término $ax$ y simplifica significativamente las operaciones:

\textbf{Duplicación estándar} (con $a$ genérico):
\[ \lambda = \frac{3x_1^2 + a}{2y_1} \]

\textbf{Duplicación en secp256k1} ($a = 0$):
\[ \lambda = \frac{3x_1^2}{2y_1} \]

Ahorro: una multiplicación por constante menos por cada duplicación. En millones de operaciones criptográficas, esto suma.
\end{keyidea}

\begin{mathnote}
Las fórmulas optimizadas para secp256k1 en coordenadas Jacobianas $(X, Y, Z)$ donde $(x, y) = (X/Z^2, Y/Z^3)$:

\textbf{Duplicación} de $P = (X_1, Y_1, Z_1)$:
\begin{align}
S &= 4 X_1 Y_1^2 \\
M &= 3 X_1^2 \quad \text{(no hay término } aZ_1^4 \text{ porque } a = 0\text{)} \\
X_3 &= M^2 - 2S \\
Y_3 &= M(S - X_3) - 8Y_1^4 \\
Z_3 &= 2Y_1 Z_1
\end{align}

Costo: 1S + 8M (1 cuadrado, 8 multiplicaciones), vs. 1S + 10M con $a \neq 0$.
\end{mathnote}

%-------------------------------------------
\subsubsection{El Primo $p$: Estructura Especial}

\begin{definition}[Primo de Mersenne Generalizado]
El primo de secp256k1:
\[ p = 2^{256} - 2^{32} - 977 \]

tiene una estructura que permite \textbf{reducción rápida} módulo $p$.
\end{definition}

\begin{intuicion}
La reducción modular normalmente requiere división, que es costosa. Pero con primos de forma especial, podemos usar trucos algebraicos.

Para $p = 2^{256} - c$ donde $c$ es pequeño:
\[ x \mod p = x \mod (2^{256} - c) \]

Si $x = x_{\text{high}} \cdot 2^{256} + x_{\text{low}}$, entonces:
\begin{align}
x &\equiv x_{\text{low}} + c \cdot x_{\text{high}} \pmod{p}
\end{align}

Como $c = 2^{32} + 977$ es pequeño (33 bits), la multiplicación $c \cdot x_{\text{high}}$ es mucho más barata que una división completa.
\end{intuicion}

\begin{example}[Reducción rápida]
Supongamos que tenemos un producto de 512 bits (resultado de multiplicar dos números de 256 bits):
\[ x = x_1 \cdot 2^{256} + x_0 \]

donde $x_1, x_0 < 2^{256}$.

\textbf{Paso 1:}
\[ x \equiv x_0 + (2^{32} + 977) \cdot x_1 \pmod{p} \]

\textbf{Paso 2:} El resultado tiene como mucho $256 + 33 = 289$ bits. Repetimos si es necesario.

\textbf{Paso 3:} Después de 2-3 iteraciones, el resultado está en $[0, p)$.

Esto es $\sim$10x más rápido que división general.
\end{example}

%-------------------------------------------
\subsubsection{El Orden $n$ y Seguridad}

\begin{definition}[Orden del Grupo]
\[ n = \texttt{FFFFFFFFFFFFFFFFFFFFFFFFFFFFFFFEBAAEDCE6AF48A03BBFD25E8CD0364141} \]

En decimal:
\[ n = 115792089237316195423570985008687907852837564279074904382605163141518161494337 \]

Esto es aproximadamente $2^{256} - 4.32 \times 10^{38}$.
\end{definition}

\begin{keyidea}
El orden $n$ es \textbf{primo}. Esto es crucial porque:
\begin{enumerate}
    \item El grupo generado por $G$ es cíclico de orden $n$
    \item Todo punto $\neq \mathcal{O}$ tiene orden $n$ (no hay subgrupos pequeños)
    \item El cofactor es $h = \#E(\Fp) / n = 1$
\end{enumerate}

Si hubiera subgrupos pequeños, ciertos ataques (como el ataque de Pohlig-Hellman) serían posibles.
\end{keyidea}

\begin{mathnote}
\textbf{Nivel de seguridad:}

El mejor ataque conocido contra ECDLP es el algoritmo rho de Pollard, que requiere $O(\sqrt{n})$ operaciones:
\[ \sqrt{n} \approx \sqrt{2^{256}} = 2^{128} \]

Por lo tanto, secp256k1 ofrece \textbf{128 bits de seguridad}, equivalente a:
\begin{itemize}
    \item AES-128
    \item RSA de $\sim$3072 bits
    \item Hash de 256 bits (para colisiones)
\end{itemize}
\end{mathnote}

%-------------------------------------------
\subsubsection{El Generador $G$}

\begin{definition}[Punto Generador]
El generador $G$ es un punto específico de la curva tal que:
\[ \{G, 2G, 3G, \ldots, nG\} = E(\Fp) \setminus \{\mathcal{O}\} \]

Es decir, $G$ genera todo el grupo (excepto el punto en el infinito).
\end{definition}

\begin{remark}
¿Cómo se eligió $G$? El estándar SECG especifica que $G$ se obtiene hasheando una semilla conocida. Esto es un procedimiento ``nothing-up-my-sleeve'': cualquiera puede verificar que $G$ no fue elegido maliciosamente.

Para secp256k1, el proceso fue:
\begin{enumerate}
    \item Elegir la coordenada $x$ como el menor entero tal que $x^3 + 7$ es un cuadrado en $\Fp$
    \item La coordenada $y$ es la raíz cuadrada con cierta paridad
\end{enumerate}
\end{remark}

\begin{example}[Verificar que $G$ está en la curva]
\begin{align}
G_x &= \texttt{0x79BE667E...16F81798} \\
G_y &= \texttt{0x483ADA77...FB10D4B8}
\end{align}

Verificamos: $G_y^2 \stackrel{?}{=} G_x^3 + 7 \pmod{p}$

\begin{verbatim}
>>> Gx = 0x79BE667EF9DCBBAC55A06295CE870B07029BFCDB2DCE28D959F2815B16F81798
>>> Gy = 0x483ADA7726A3C4655DA4FBFC0E1108A8FD17B448A68554199C47D08FFB10D4B8
>>> p = 0xFFFFFFFFFFFFFFFFFFFFFFFFFFFFFFFFFFFFFFFFFFFFFFFFFFFFFFFEFFFFFC2F
>>> (Gy**2 - Gx**3 - 7) % p
0
\end{verbatim}

$G$ está en la curva ✓
\end{example}

%-------------------------------------------
\subsubsection{Claves Públicas y Privadas}

\begin{definition}[Par de Claves en secp256k1]
\begin{itemize}
    \item \textbf{Clave privada}: Un entero aleatorio $d \in [1, n-1]$
    \item \textbf{Clave pública}: El punto $Q = dG$
\end{itemize}

La seguridad se basa en que conocer $Q$ y $G$ no permite recuperar $d$ (problema ECDLP).
\end{definition}

\begin{example}[Generación de claves en Bitcoin/Ethereum]
\textbf{Paso 1:} Generar 256 bits aleatorios seguros
\begin{verbatim}
private_key = 0xE8F32E723DECF4051AEFAC8E2C93C9C5B214313817CDB01A1494B917C8436B35
\end{verbatim}

\textbf{Paso 2:} Verificar que está en rango válido
\begin{verbatim}
assert 1 <= private_key < n
\end{verbatim}

\textbf{Paso 3:} Calcular clave pública
\begin{verbatim}
public_key = private_key * G  # Multiplicación escalar
# Resultado: (Qx, Qy)
\end{verbatim}

\textbf{Paso 4:} Serializar
\begin{verbatim}
# Sin comprimir (65 bytes): 04 || Qx || Qy
# Comprimido (33 bytes): 02/03 || Qx (según paridad de Qy)
\end{verbatim}
\end{example}

\begin{advertencia}
\textbf{Generación de claves privadas:}

La clave privada DEBE ser:
\begin{itemize}
    \item Generada con un CSPRNG (generador criptográficamente seguro)
    \item Verificada que esté en $[1, n-1]$
    \item Nunca reutilizada entre sistemas diferentes
    \item Almacenada de forma segura
\end{itemize}

Fallos históricos:
\begin{itemize}
    \item Android (2013): Bug en SecureRandom causó reutilización de nonces, perdiendo fondos
    \item Blockchain.info (2014): Generador débil predecible
    \item Numerosas wallets con entropía insuficiente
\end{itemize}
\end{advertencia}

%-------------------------------------------
\subsubsection{Direcciones de Bitcoin}

\begin{definition}[Derivación de Dirección Bitcoin]
Desde la clave pública, se deriva la dirección así:

\begin{enumerate}
    \item $\text{pubkey\_hash} = \text{RIPEMD160}(\text{SHA256}(\text{pubkey}))$ (20 bytes)
    \item $\text{versioned} = \text{version\_byte} \,||\, \text{pubkey\_hash}$ (21 bytes)
    \item $\text{checksum} = \text{SHA256}(\text{SHA256}(\text{versioned}))[0:4]$ (4 bytes)
    \item $\text{address} = \text{Base58}(\text{versioned} \,||\, \text{checksum})$
\end{enumerate}

Version bytes:
\begin{itemize}
    \item \texttt{0x00}: Mainnet P2PKH (empieza con ``1'')
    \item \texttt{0x05}: Mainnet P2SH (empieza con ``3'')
    \item \texttt{0x6F}: Testnet P2PKH (empieza con ``m'' o ``n'')
\end{itemize}
\end{definition}

\begin{example}[Derivación completa]
\begin{verbatim}
1. Clave pública (comprimida):
   02 50863AD64A87AE8A2FE83C1AF1A8403CB53F53E486D8511DAD8A04887E5B2352

2. SHA256:
   0B7C28C9B7290C98D7438E70B3D3F7C848FBD7D1DC194FF83F4F7CC9B1378E98

3. RIPEMD160:
   F54A5851E9372B87810A8E60CDD2E7CFD80B6E31

4. Añadir version byte (0x00 para mainnet):
   00 F54A5851E9372B87810A8E60CDD2E7CFD80B6E31

5. Doble SHA256 y tomar primeros 4 bytes (checksum):
   C7F18FE8

6. Concatenar:
   00 F54A5851E9372B87810A8E60CDD2E7CFD80B6E31 C7F18FE8

7. Base58 encode:
   1PMycacnJaSqwwJqjawXBErnLsZ7RkXUAs
\end{verbatim}
\end{example}

%-------------------------------------------
\subsubsection{Direcciones de Ethereum}

\begin{definition}[Derivación de Dirección Ethereum]
Ethereum usa un proceso más simple:
\begin{enumerate}
    \item Clave pública \textbf{sin comprimir, sin prefijo} (64 bytes, solo $x \,||\, y$)
    \item $\text{hash} = \text{Keccak256}(\text{pubkey})$ (32 bytes)
    \item $\text{address} = \text{hash}[12:32]$ (últimos 20 bytes)
    \item Representación: ``0x'' + hexadecimal (42 caracteres)
\end{enumerate}
\end{definition}

\begin{example}[Dirección Ethereum]
\begin{verbatim}
1. Clave pública (64 bytes, sin 04):
   50863AD64A87AE8A2FE83C1AF1A8403CB53F53E486D8511DAD8A04887E5B2352
   2CD470243453A299FA9E77237716103ABC11A1DF38855ED6F2EE187E9C582BA6

2. Keccak256:
   FC12AD814631BA689F7ADE2B6F9F4D17C8B6B6B4BF1D7E1E0C1C5D8F9A0B1C2D...

3. Últimos 20 bytes:
   C8B6B6B4BF1D7E1E0C1C5D8F9A0B1C2D3E4F5A6B

4. Dirección:
   0xC8B6B6B4BF1D7E1E0C1C5D8F9A0B1C2D3E4F5A6B
\end{verbatim}
\end{example}

\begin{remark}
Ethereum añade un checksum opcional mediante capitalización mixta (EIP-55):
\begin{verbatim}
Sin checksum: 0xc8b6b6b4bf1d7e1e0c1c5d8f9a0b1c2d3e4f5a6b
Con checksum: 0xC8b6B6b4bF1D7e1E0c1c5d8F9A0b1c2D3e4F5A6B
\end{verbatim}

Las letras se capitalizan según los bits del hash de la dirección.
\end{remark}

%-------------------------------------------
\subsubsection{Comparación con Otras Curvas}

\begin{center}
\begin{tabular}{lcccc}
\toprule
\textbf{Curva} & \textbf{$a$} & \textbf{$b$} & \textbf{Bits} & \textbf{Origen} \\
\midrule
secp256k1 & 0 & 7 & 256 & SECG \\
secp256r1 (P-256) & $-3$ & variable & 256 & NIST \\
Curve25519 & -- & -- & 255 & Bernstein \\
BN254 & 0 & 3 & 254 & Barreto-Naehrig \\
BLS12-381 & 0 & 4 & 381 & BLS \\
\bottomrule
\end{tabular}
\end{center}

\begin{keyidea}
\textbf{¿Por qué no P-256 (secp256r1)?}

P-256 es más popular en TLS, certificados, etc. Pero:
\begin{enumerate}
    \item Sus parámetros fueron generados por NIST/NSA con proceso opaco
    \item El coeficiente $b$ parece aleatorio, sin justificación clara
    \item Después de Snowden, hay sospechas de posibles ``backdoors''
\end{enumerate}

secp256k1 tiene parámetros transparentes: $a = 0$, $b = 7$ (el menor $b$ que da curva segura con $a = 0$).
\end{keyidea}

\begin{center}
\begin{tabular}{lcc}
\toprule
\textbf{Aspecto} & \textbf{secp256k1} & \textbf{P-256} \\
\midrule
Transparencia & Alta (``nothing-up-my-sleeve'') & Baja (proceso NIST opaco) \\
Velocidad ($a = 0$) & Más rápida & Estándar \\
Adopción cripto & Bitcoin, Ethereum & Rara \\
Adopción TLS/PKI & Rara & Dominante \\
Librerías & libsecp256k1 (muy optimizada) & Muchas \\
\bottomrule
\end{tabular}
\end{center}

%-------------------------------------------
\subsubsection{Vulnerabilidades y Consideraciones}

\begin{advertencia}
\textbf{Curvas potencialmente vulnerables:}

secp256k1 es segura contra ataques conocidos, pero algunas clases de curvas son vulnerables:

\begin{itemize}
    \item \textbf{Curvas anómalas}: $\#E(\Fp) = p$ → ECDLP se reduce a DLP en $\Fp$ (resuelto en tiempo lineal)
    \item \textbf{Curvas supersingulares}: Vulnerable al ataque MOV (embedding en campo de extensión)
    \item \textbf{Embedding degree bajo}: Similar al ataque MOV
\end{itemize}

secp256k1 no tiene ninguna de estas propiedades.
\end{advertencia}

\begin{remark}
\textbf{Computación cuántica:}

El algoritmo de Shor puede resolver ECDLP en tiempo polinomial con una computadora cuántica suficientemente grande. Para secp256k1:
\begin{itemize}
    \item Se necesitan $\sim$2500 qubits lógicos
    \item Con corrección de errores: millones de qubits físicos
    \item Estimación optimista: 10-20 años para amenaza real
\end{itemize}

Por esto, se investigan alternativas ``post-cuánticas'' (lattices, hashes, códigos).
\end{remark}

%-------------------------------------------
\subsubsection{Implementación: libsecp256k1}

\begin{intuicion}
La biblioteca \texttt{libsecp256k1}, desarrollada por Pieter Wuille para Bitcoin Core, es considerada la implementación de referencia. Sus características:
\begin{itemize}
    \item Escrita en C puro, sin dependencias
    \item Optimizaciones de ensamblador para x86\_64, ARM
    \item Tiempo constante (resistente a side-channels)
    \item Verificación formal parcial
    \item Usada por Bitcoin, Ethereum, y cientos de proyectos
\end{itemize}
\end{intuicion}

\begin{example}[Uso de libsecp256k1 en C]
\begin{verbatim}
#include <secp256k1.h>

// Crear contexto
secp256k1_context* ctx = secp256k1_context_create(
    SECP256K1_CONTEXT_SIGN | SECP256K1_CONTEXT_VERIFY);

// Generar par de claves
unsigned char seckey[32];
// ... llenar seckey con 32 bytes aleatorios seguros ...

secp256k1_pubkey pubkey;
int result = secp256k1_ec_pubkey_create(ctx, &pubkey, seckey);

// Serializar clave pública (comprimida)
unsigned char pubkey_ser[33];
size_t len = 33;
secp256k1_ec_pubkey_serialize(ctx, pubkey_ser, &len, &pubkey,
    SECP256K1_EC_COMPRESSED);

// Firmar mensaje
unsigned char msg_hash[32] = /* SHA256 del mensaje */;
secp256k1_ecdsa_signature sig;
secp256k1_ecdsa_sign(ctx, &sig, msg_hash, seckey, NULL, NULL);

// Verificar firma
int valid = secp256k1_ecdsa_verify(ctx, &sig, msg_hash, &pubkey);

secp256k1_context_destroy(ctx);
\end{verbatim}
\end{example}

%-------------------------------------------
\subsubsection{Optimizaciones Específicas}

\begin{definition}[Endomorfismo GLV]
secp256k1 admite el endomorfismo eficiente:
\[ \phi(x, y) = (\beta x, y) \]

donde $\beta$ es una raíz cúbica de la unidad en $\Fp$ (existe porque $p \equiv 1 \pmod{3}$).

Este endomorfismo satisface $\phi^2 + \phi + 1 = 0$ en el anillo de endomorfismos, correspondiente a un $\lambda$ tal que:
\[ \phi(P) = \lambda P \]
\end{definition}

\begin{keyidea}
\textbf{Speedup GLV:}

Para calcular $nP$, descomponemos:
\[ n = n_1 + n_2 \lambda \]

donde $n_1, n_2 \approx \sqrt{n}$ (128 bits cada uno).

Entonces:
\[ nP = n_1 P + n_2 \phi(P) = n_1 P + n_2 (\beta x_P, y_P) \]

Usando double-and-add simultáneo sobre ambos sumandos, reducimos el número de duplicaciones a la mitad. Speedup: $\sim$33\%.
\end{keyidea}

\begin{mathnote}
Los valores concretos para secp256k1:
\begin{align}
\beta &= \texttt{0x7AE96A2B657C07106E64479EAC3434E99CF0497512F58995C1396C28719501EE} \\
\lambda &= \texttt{0x5363AD4CC05C30E0A5261C028812645A122E22EA20816678DF02967C1B23BD72}
\end{align}

Verificación: $\beta^3 = 1 \pmod{p}$ y $\lambda^2 + \lambda + 1 = 0 \pmod{n}$.
\end{mathnote}

%-------------------------------------------
\subsubsection{Aplicación en Banco ZK}

\begin{bancozk}
secp256k1 es fundamental para la interoperabilidad de Banco ZK con Ethereum:

\textbf{1. Autenticación de usuarios}
\begin{itemize}
    \item Los usuarios conectan sus wallets Ethereum existentes
    \item Firman transacciones de Banco ZK con sus claves secp256k1
    \item No necesitan crear nuevas claves ni aprender nuevos procesos
\end{itemize}

\textbf{2. Verificación de firmas en circuitos ZK}

Verificar una firma ECDSA dentro de un circuito ZK es costoso pero necesario:
\begin{center}
\begin{tabular}{lc}
\toprule
\textbf{Operación} & \textbf{Constraints} \\
\midrule
Multiplicación escalar en secp256k1 & $\sim$500,000 \\
Verificación ECDSA completa & $\sim$600,000 \\
\bottomrule
\end{tabular}
\end{center}

\textbf{3. Estrategia de optimización}

En lugar de verificar firmas secp256k1 dentro del SNARK, usamos:
\begin{enumerate}
    \item Verificar firma en el contrato L1 (barato en gas: $\sim$3000 gas)
    \item Pasar la clave pública como input público al circuito
    \item El circuito verifica que el nullifier corresponde a esa clave
\end{enumerate}

Esto reduce de 600,000 a $\sim$1,000 constraints.

\textbf{4. Bridge de fondos}
\begin{itemize}
    \item Depósitos: Usuario firma con secp256k1 para autorizar
    \item Retiros: El contrato verifica firma antes de liberar fondos
    \item La misma clave protege fondos en Ethereum y en Banco ZK
\end{itemize}

\textbf{5. Compatibilidad futura}

Consideramos migrar a EdDSA (ed25519) para nuevas funcionalidades porque:
\begin{itemize}
    \item Firmas determinísticas (no hay riesgo de nonce)
    \item Más eficiente en circuitos ZK
    \item Batch verification nativa
\end{itemize}

Pero mantenemos secp256k1 para compatibilidad con el ecosistema Ethereum existente.
\end{bancozk}

%-------------------------------------------
\subsubsection{Constantes Útiles}

\begin{mathnote}
\textbf{Resumen de constantes secp256k1:}

\begin{verbatim}
# Campo base
p = 0xFFFFFFFFFFFFFFFFFFFFFFFFFFFFFFFFFFFFFFFFFFFFFFFFFFFFFFFEFFFFFC2F

# Orden del grupo
n = 0xFFFFFFFFFFFFFFFFFFFFFFFFFFFFFFFEBAAEDCE6AF48A03BBFD25E8CD0364141

# Coeficientes de la curva
a = 0
b = 7

# Generador G
Gx = 0x79BE667EF9DCBBAC55A06295CE870B07029BFCDB2DCE28D959F2815B16F81798
Gy = 0x483ADA7726A3C4655DA4FBFC0E1108A8FD17B448A68554199C47D08FFB10D4B8

# Para endomorfismo GLV
beta = 0x7AE96A2B657C07106E64479EAC3434E99CF0497512F58995C1396C28719501EE
lambda = 0x5363AD4CC05C30E0A5261C028812645A122E22EA20816678DF02967C1B23BD72

# Cofactor
h = 1
\end{verbatim}
\end{mathnote}

% ALE termina secp256k1
% ALE comienza bn128
%===========================================
% SUBSECCIONES: BN254 y BLS12-381 - VERSIÓN EXPANDIDA
% Usa tus entornos: definition, example, keyidea, bancozk, 
%                   mathnote, intuicion, advertencia, remark
%===========================================

\subsection{BN254 (alt\_bn128)}

Curva pairing-friendly usada en ZK-SNARKs de Ethereum.
\begin{itemize}
    \item Embedding degree: 12
    \item Seguridad: $\sim$100 bits (suficiente para SNARKs actuales)
    \item Precompilados en EVM: EIP-196, EIP-197
    \item Usada en: Groth16, PLONK en Ethereum
\end{itemize}

%-------------------------------------------
\subsubsection{¿Qué es una Curva Pairing-Friendly?}

\begin{intuicion}
Las curvas ``normales'' como secp256k1 permiten operaciones de grupo (suma, multiplicación escalar), pero nada más. Las curvas \textbf{pairing-friendly} tienen una propiedad mágica adicional: permiten ``multiplicar'' puntos de curvas elípticas de una forma que no es posible en curvas ordinarias.

Esta operación extra, llamada \textbf{pairing} o emparejamiento bilineal, es lo que hace posible los SNARKs eficientes.
\end{intuicion}

\begin{definition}[Pairing Bilineal]
Un pairing es una función:
\[ e: \G_1 \times \G_2 \to \G_T \]

donde $\G_1, \G_2$ son grupos de puntos de curva elíptica y $\G_T$ es un grupo multiplicativo (subgrupo del campo de extensión), que satisface:

\textbf{Bilinealidad:}
\[ e(aP, bQ) = e(P, Q)^{ab} \]

para todos los puntos $P \in \G_1$, $Q \in \G_2$ y escalares $a, b$.
\end{definition}

\begin{keyidea}
La bilinealidad permite ``mover'' escalares entre argumentos:
\begin{align}
e(aP, Q) &= e(P, Q)^a = e(P, aQ)
\end{align}

Esto es imposible en curvas normales y es la base de:
\begin{itemize}
    \item Verificación de SNARKs en tiempo constante
    \item Firmas BLS agregables
    \item Identity-Based Encryption (IBE)
    \item Attribute-Based Encryption (ABE)
\end{itemize}
\end{keyidea}

\begin{example}[Verificación de Groth16]
En Groth16, la prueba consiste en tres puntos $\pi = (A, B, C)$ donde $A, C \in \G_1$ y $B \in \G_2$.

La verificación comprueba:
\[ e(A, B) = e(\alpha, \beta) \cdot e(\sum_i x_i L_i, \gamma) \cdot e(C, \delta) \]

Sin pairings, esta ecuación no tendría sentido. Con pairings, verificar una prueba de cualquier computación compleja requiere solo 3-4 pairings, independientemente de la complejidad del circuito.
\end{example}

%-------------------------------------------
\subsubsection{Historia de BN254}

\begin{remark}
BN254 pertenece a la familia de curvas Barreto-Naehrig (BN), propuesta por Paulo Barreto y Michael Naehrig en 2005. Estas curvas se caracterizan por:
\begin{itemize}
    \item Construcción paramétrica simple
    \item Embedding degree $k = 12$
    \item Eficiencia en implementación
\end{itemize}

BN254 específicamente fue popularizada por la biblioteca \texttt{bn128} y adoptada por Ethereum en 2017 mediante los EIPs 196 y 197 para habilitar SNARKs on-chain.
\end{remark}

\begin{intuicion}
El nombre ``alt\_bn128'' que aparece en documentación de Ethereum es un nombre alternativo histórico. El ``128'' se refería originalmente a la seguridad estimada en bits, pero análisis posteriores mostraron que la seguridad real es $\sim$100 bits debido a avances en ataques a pairings (NFS modificado).
\end{intuicion}

%-------------------------------------------
\subsubsection{Parámetros de BN254}

\begin{definition}[Construcción Barreto-Naehrig]
Las curvas BN se parametrizan por un entero $u$. Para BN254:
\[ u = \texttt{0x44E992B44A6909F1} \]

A partir de $u$ se derivan todos los parámetros:
\begin{align}
p &= 36u^4 + 36u^3 + 24u^2 + 6u + 1 \quad \text{(primo del campo base)} \\
n &= 36u^4 + 36u^3 + 18u^2 + 6u + 1 \quad \text{(orden del grupo)}
\end{align}
\end{definition}

\begin{mathnote}
\textbf{Parámetros explícitos de BN254:}

\textbf{Campo base} $\Fp$:
\begin{verbatim}
p = 21888242871839275222246405745257275088696311157297823662689037894645226208583
\end{verbatim}
En hexadecimal:
\begin{verbatim}
p = 0x30644E72E131A029B85045B68181585D97816A916871CA8D3C208C16D87CFD47
\end{verbatim}

\textbf{Orden del grupo} $n$:
\begin{verbatim}
n = 21888242871839275222246405745257275088548364400416034343698204186575808495617
\end{verbatim}

\textbf{Ecuación de la curva} $E(\Fp)$:
\[ y^2 = x^3 + 3 \]

\textbf{Generador} $G_1$:
\begin{align}
G_{1x} &= 1 \\
G_{1y} &= 2
\end{align}
\end{mathnote}

\begin{remark}
La simplicidad de $G_1 = (1, 2)$ y $b = 3$ no es coincidencia. Las curvas BN se diseñan buscando parámetros pequeños cuando es posible, lo que facilita la implementación y verificación.
\end{remark}

%-------------------------------------------
\subsubsection{Los Tres Grupos: $\G_1$, $\G_2$, $\G_T$}

\begin{definition}[Grupos en BN254]
\textbf{$\G_1$}: Puntos de la curva sobre el campo base
\[ \G_1 \subset E(\Fp): y^2 = x^3 + 3 \]

\textbf{$\G_2$}: Puntos de la curva sobre una extensión cuadrática
\[ \G_2 \subset E(\Fp^2): y^2 = x^3 + \frac{3}{9+i} \]
donde $\Fp^2 = \Fp[i]/(i^2 + 1)$ (adjuntando $\sqrt{-1}$).

\textbf{$\G_T$}: Subgrupo multiplicativo del campo de extensión
\[ \G_T \subset \Fp^{12*} \]
\end{definition}

\begin{intuicion}
\begin{itemize}
    \item $\G_1$: Puntos ``normales'', coordenadas de 256 bits
    \item $\G_2$: Puntos ``extendidos'', coordenadas de 512 bits (par de elementos de $\Fp$)
    \item $\G_T$: Elemento del campo de extensión, 3072 bits ($12 \times 256$)
\end{itemize}

El pairing toma un punto ``pequeño'' y uno ``grande'' y produce un elemento ``muy grande''.
\end{intuicion}

%-------------------------------------------
\subsubsection{El Embedding Degree}

\begin{definition}[Embedding Degree]
El \textbf{embedding degree} $k$ es el menor entero positivo tal que:
\[ n \mid p^k - 1 \]

donde $n$ es el orden del grupo de la curva y $p$ es el primo del campo base.
\end{definition}

\begin{keyidea}
El embedding degree determina:
\begin{enumerate}
    \item \textbf{Seguridad}: El pairing mapea ECDLP a DLP en $\Fp^k$. Para seguridad, necesitamos que DLP en $\Fp^k$ sea difícil.
    \item \textbf{Eficiencia}: $k$ más pequeño = computación más rápida, pero también potencialmente menos seguridad.
\end{enumerate}

Para BN254: $k = 12$, lo que significa que el grupo target $\G_T$ vive en $\Fp^{12}$ ($\sim$3072 bits).
\end{keyidea}

\begin{mathnote}
\textbf{¿Por qué $k = 12$ da $\sim$100 bits de seguridad?}

La seguridad del sistema se determina por el eslabón más débil:
\begin{itemize}
    \item ECDLP en $\G_1$: $\sqrt{n} \approx 2^{127}$ operaciones (Pollard rho)
    \item DLP en $\Fp^{12}$: $\sim 2^{100}$ operaciones (NFS especial para pairings)
\end{itemize}

El ataque al campo de extensión resultó ser más eficiente de lo esperado inicialmente, reduciendo la seguridad a $\sim$100 bits.
\end{mathnote}

%-------------------------------------------
\subsubsection{Precompilados de Ethereum: EIP-196 y EIP-197}

\begin{definition}[Precompilados EVM para BN254]
Ethereum incluye contratos precompilados (código nativo, no EVM) para operaciones BN254:

\textbf{EIP-196} (direcciones 0x06, 0x07):
\begin{itemize}
    \item \texttt{0x06}: Suma de puntos en $\G_1$
    \item \texttt{0x07}: Multiplicación escalar en $\G_1$
\end{itemize}

\textbf{EIP-197} (dirección 0x08):
\begin{itemize}
    \item \texttt{0x08}: Verificación de producto de pairings
\end{itemize}
\end{definition}

\begin{example}[Costos de gas]
\begin{center}
\begin{tabular}{llc}
\toprule
\textbf{Operación} & \textbf{Dirección} & \textbf{Gas} \\
\midrule
Suma $\G_1$ & 0x06 & 150 \\
Multiplicación escalar $\G_1$ & 0x07 & 6,000 \\
Pairing (por par) & 0x08 & 34,000 + 45,000 × pares \\
\bottomrule
\end{tabular}
\end{center}

Para verificar un Groth16 típico (3 pairings):
\[ \text{Gas} \approx 34,000 + 3 \times 45,000 = 169,000 \text{ gas} \]
\end{example}

\begin{keyidea}
Sin precompilados, verificar un pairing en EVM puro costaría millones de gas (si fuera posible). Los precompilados hacen SNARKs on-chain económicamente viables.

Esto es por qué BN254 domina en Ethereum: es la única curva con soporte nativo en la EVM.
\end{keyidea}

\begin{example}[Llamada al precompilado de pairing en Solidity]
\begin{verbatim}
function verifyPairing(
    uint256[2] memory a,      // Punto A en G1
    uint256[2][2] memory b,   // Punto B en G2
    uint256[2] memory c,      // Punto C en G1
    uint256[2] memory alpha,  // Parámetro de setup
    uint256[2][2] memory beta // Parámetro de setup
) internal view returns (bool) {
    uint256[24] memory input;
    // Codificar los puntos...
    
    assembly {
        // Llamar al precompilado 0x08
        let success := staticcall(
            sub(gas(), 2000),
            0x08,  // Dirección del precompilado
            input,
            768,   // 24 * 32 bytes
            input,
            32     // Resultado: 0 o 1
        )
        // ...
    }
}
\end{verbatim}
\end{example}

%-------------------------------------------
\subsubsection{El Pairing Óptimo Ate}

\begin{definition}[Pairing Ate Óptimo]
BN254 usa el pairing ``Optimal Ate'', que es la versión más eficiente conocida para curvas BN:
\[ e: \G_1 \times \G_2 \to \G_T \]

El cálculo involucra:
\begin{enumerate}
    \item \textbf{Miller loop}: Construye una función racional evaluada en los puntos
    \item \textbf{Final exponentiation}: Eleva el resultado a $(p^{12}-1)/n$
\end{enumerate}
\end{definition}

\begin{mathnote}
La complejidad del pairing Ate en BN254:
\begin{itemize}
    \item Miller loop: $O(\log n)$ iteraciones, cada una con operaciones en $\Fp^{12}$
    \item Final exponentiation: Varias multiplicaciones y exponenciaciones en $\Fp^{12}$
\end{itemize}

Tiempo típico en CPU moderno: $\sim$1-2 ms por pairing.
\end{mathnote}

%-------------------------------------------
\subsubsection{Seguridad: ¿Son Suficientes 100 Bits?}

\begin{advertencia}
La seguridad de $\sim$100 bits de BN254 está por debajo del estándar moderno de 128 bits. Sin embargo:

\textbf{Argumentos a favor de seguir usando BN254:}
\begin{enumerate}
    \item Los ataques conocidos requieren $2^{100}$ operaciones, aún impracticables
    \item Las pruebas ZK típicamente tienen validez temporal corta
    \item El costo de migrar el ecosistema Ethereum es enorme
    \item Para aplicaciones no críticas, es suficiente
\end{enumerate}

\textbf{Argumentos para migrar a BLS12-381:}
\begin{enumerate}
    \item 128 bits es el estándar de la industria
    \item El hardware mejora; el margen de seguridad se reduce
    \item Nuevos protocolos deberían usar curvas más seguras
\end{enumerate}
\end{advertencia}

\begin{remark}
La comunidad de Ethereum está considerando añadir precompilados para BLS12-381 (EIP-2537), pero aún no está implementado en mainnet. Mientras tanto, verificar BLS12-381 en EVM es prohibitivamente caro.
\end{remark}

%-------------------------------------------
\subsubsection{Aplicación en Banco ZK}

\begin{bancozk}
BN254 es la curva principal de Banco ZK para verificación on-chain:

\textbf{1. Verificación de pruebas Groth16}
\begin{itemize}
    \item El verificador on-chain usa los precompilados 0x06, 0x07, 0x08
    \item Costo: $\sim$200,000 gas por verificación
    \item Independiente del tamaño del circuito
\end{itemize}

\textbf{2. Por qué elegimos BN254 sobre BLS12-381}
\begin{center}
\begin{tabular}{lcc}
\toprule
\textbf{Criterio} & \textbf{BN254} & \textbf{BLS12-381} \\
\midrule
Precompilados EVM & Sí & No \\
Gas para verificar & $\sim$200K & $\sim$3M+ (estimado) \\
Seguridad & $\sim$100 bits & 128 bits \\
\bottomrule
\end{tabular}
\end{center}

La diferencia de costo es 15x, haciendo BLS12-381 impracticable actualmente.

\textbf{3. Estructura de la prueba}
\begin{verbatim}
struct Proof {
    // Elementos en G1 (256 bits cada coordenada)
    uint256[2] a;      // Punto A
    uint256[2] c;      // Punto C
    
    // Elemento en G2 (512 bits cada coordenada)
    uint256[2][2] b;   // Punto B
}
\end{verbatim}

Tamaño total de prueba: $2 + 2 + 4 = 8$ elementos de campo = 256 bytes.

\textbf{4. Migración futura}

Cuando EIP-2537 se active:
\begin{enumerate}
    \item Desplegaremos verificadores BLS12-381 para nuevos circuitos
    \item Mantendremos verificadores BN254 para compatibilidad
    \item Los usuarios no notarán diferencia (abstracción en contrato)
\end{enumerate}
\end{bancozk}

%-------------------------------------------
\subsubsection{Librerías y Herramientas}

\begin{center}
\begin{tabular}{lll}
\toprule
\textbf{Lenguaje} & \textbf{Librería} & \textbf{Notas} \\
\midrule
Rust & arkworks (ark-bn254) & Referencia, muy completa \\
Rust & bellman & Usada por Zcash \\
Go & gnark & Muy rápida \\
JavaScript & snarkjs & Para desarrollo/testing \\
Solidity & Precompilados nativos & Producción \\
\bottomrule
\end{tabular}
\end{center}

%===========================================
% SECCIÓN: BLS12-381
%===========================================

\subsection{BLS12-381}

Curva de mayor seguridad para pairings.
\begin{itemize}
    \item Embedding degree: 12
    \item Seguridad: $\sim$128 bits
    \item Usada en: Ethereum 2.0 (firmas BLS), Zcash Sapling
    \item Más segura que BN254, pero sin precompilados EVM
\end{itemize}

%-------------------------------------------
\subsubsection{Origen y Motivación}

\begin{intuicion}
BLS12-381 fue diseñada específicamente para abordar las limitaciones de seguridad de curvas anteriores como BN254. El nombre proviene de:
\begin{itemize}
    \item \textbf{BLS}: Barreto-Lynn-Scott (los diseñadores de la familia de curvas)
    \item \textbf{12}: Embedding degree
    \item \textbf{381}: Tamaño del campo base en bits
\end{itemize}
\end{intuicion}

\begin{remark}
BLS12-381 fue propuesta por Sean Bowe (Zcash) en 2017 como parte del trabajo para Zcash Sapling. Los parámetros fueron elegidos cuidadosamente para:
\begin{enumerate}
    \item Alcanzar 128 bits de seguridad
    \item Maximizar eficiencia dentro de esa restricción
    \item Permitir implementaciones de tiempo constante
    \item Tener buenas propiedades para FFT (para SNARKs)
\end{enumerate}
\end{remark}

%-------------------------------------------
\subsubsection{Parámetros de BLS12-381}

\begin{definition}[Construcción BLS12]
Las curvas BLS12 se parametrizan por un entero $x$. Para BLS12-381:
\[ x = -\texttt{0xD201000000010000} \]

(Valor negativo, lo que da aritmética eficiente)

A partir de $x$:
\begin{align}
p &= (x-1)^2 (x^4 - x^2 + 1)/3 + x \\
n &= x^4 - x^2 + 1
\end{align}
\end{definition}

\begin{mathnote}
\textbf{Parámetros explícitos de BLS12-381:}

\textbf{Campo base} $\Fp$ (381 bits):
\begin{verbatim}
p = 0x1a0111ea397fe69a4b1ba7b6434bacd764774b84f38512bf6730d2a0f6b0f624
      1eabfffeb153ffffb9feffffffffaaab
\end{verbatim}

\textbf{Orden del subgrupo} $n$ (255 bits):
\begin{verbatim}
n = 0x73eda753299d7d483339d80809a1d80553bda402fffe5bfeffffffff00000001
\end{verbatim}

\textbf{Ecuación de la curva}:
\[ E(\Fp): y^2 = x^3 + 4 \]

\textbf{Twist para $\G_2$}:
\[ E'(\Fp^2): y^2 = x^3 + 4(1 + i) \]
donde $\Fp^2 = \Fp[i]/(i^2 + 1)$.

\textbf{Cofactores}:
\begin{align}
h_1 &= (x-1)^2/3 = \texttt{0x396c8c005555e1568c00aaab0000aaab} \\
h_2 &= (x^8 - 4x^7 + 5x^6 - 4x^4 + 6x^3 - 4x^2 - 4x + 13)/9
\end{align}
\end{mathnote}

\begin{advertencia}
A diferencia de BN254 donde el cofactor es 1, BLS12-381 tiene \textbf{cofactores grandes}. Esto significa que no todos los puntos de la curva están en el subgrupo de orden $n$. Las implementaciones deben:
\begin{enumerate}
    \item Multiplicar por el cofactor para proyectar al subgrupo correcto, o
    \item Verificar membresía en el subgrupo explícitamente
\end{enumerate}

Fallar en esto puede causar vulnerabilidades de seguridad.
\end{advertencia}

%-------------------------------------------
\subsubsection{Comparación de Tamaños}

\begin{center}
\begin{tabular}{lcc}
\toprule
\textbf{Elemento} & \textbf{BN254} & \textbf{BLS12-381} \\
\midrule
Campo base $p$ & 254 bits & 381 bits \\
Orden del grupo $n$ & 254 bits & 255 bits \\
Elemento de $\G_1$ & 64 bytes & 96 bytes \\
Elemento de $\G_2$ & 128 bytes & 192 bytes \\
Elemento de $\G_T$ & 384 bytes & 576 bytes \\
Campo de extensión & $\Fp^{12}$ (3072 bits) & $\Fp^{12}$ (4572 bits) \\
\bottomrule
\end{tabular}
\end{center}

\begin{intuicion}
BLS12-381 usa $\sim$50\% más espacio que BN254 para cada elemento. Esto impacta:
\begin{itemize}
    \item Tamaño de claves públicas y firmas
    \item Ancho de banda para transmitir pruebas
    \item Almacenamiento on-chain (más gas en calldata)
\end{itemize}

Pero la seguridad adicional (128 vs 100 bits) generalmente vale este costo.
\end{intuicion}

%-------------------------------------------
\subsubsection{Seguridad: Los 128 Bits}

\begin{keyidea}
BLS12-381 fue diseñada específicamente para 128 bits de seguridad considerando:
\begin{enumerate}
    \item \textbf{ECDLP en $\G_1, \G_2$}: $\sqrt{n} \approx 2^{127.5}$ operaciones
    \item \textbf{DLP en $\G_T \subset \Fp^{12}$}: NFS especial requiere $\sim 2^{128}$ operaciones
\end{enumerate}

Ambos ataques están balanceados, sin eslabón débil.
\end{keyidea}

\begin{mathnote}
El tamaño del campo (381 bits) fue elegido precisamente para que:
\[ \text{Seguridad de } \Fp^{12} \approx \text{Seguridad de ECDLP} \approx 128 \text{ bits} \]

Si el campo fuera más pequeño (como en BN254), el ataque al campo de extensión sería más fácil.
\end{mathnote}

%-------------------------------------------
\subsubsection{Firmas BLS}

\begin{definition}[Firma BLS]
El esquema de firmas BLS aprovecha los pairings:

\textbf{Generación de claves:}
\begin{itemize}
    \item Clave privada: $sk \in \mathbb{Z}_n$ aleatorio
    \item Clave pública: $pk = sk \cdot G_2 \in \G_2$
\end{itemize}

\textbf{Firma:}
\[ \sigma = sk \cdot H(m) \in \G_1 \]
donde $H: \{0,1\}^* \to \G_1$ es un hash-to-curve.

\textbf{Verificación:}
\[ e(\sigma, G_2) \stackrel{?}{=} e(H(m), pk) \]
\end{definition}

\begin{mathnote}
¿Por qué funciona?
\begin{align}
e(\sigma, G_2) &= e(sk \cdot H(m), G_2) \\
               &= e(H(m), G_2)^{sk} \quad \text{(bilinealidad)} \\
               &= e(H(m), sk \cdot G_2) \\
               &= e(H(m), pk) \quad \checkmark
\end{align}
\end{mathnote}

%-------------------------------------------
\subsubsection{Agregación de Firmas}

\begin{keyidea}
La propiedad más poderosa de BLS es la \textbf{agregación de firmas}:

\textbf{Mismo mensaje:}
Si $n$ firmantes firman el mismo mensaje $m$:
\[ \sigma_{\text{agg}} = \sigma_1 + \sigma_2 + \cdots + \sigma_n \]

Verificación:
\[ e(\sigma_{\text{agg}}, G_2) = e(H(m), pk_1 + pk_2 + \cdots + pk_n) \]

¡Una sola verificación de pairing para $n$ firmas!

\textbf{Mensajes diferentes} (más complejo):
\[ \sigma_{\text{agg}} = \sigma_1 + \sigma_2 + \cdots + \sigma_n \]
\[ e(\sigma_{\text{agg}}, G_2) = \prod_{i=1}^{n} e(H(m_i), pk_i) \]

Requiere $n$ pairings, pero solo una firma que transmitir.
\end{keyidea}

\begin{example}[Ethereum 2.0 Beacon Chain]
La Beacon Chain tiene $\sim$900,000 validadores. Sin agregación:
\begin{itemize}
    \item 900,000 firmas ECDSA × 64 bytes = 57.6 MB por bloque
    \item Verificar 900,000 firmas individualmente
\end{itemize}

Con agregación BLS:
\begin{itemize}
    \item Los validadores se agrupan en comités de 128
    \item Cada comité produce una firma agregada (48 bytes)
    \item $\sim$2,000 firmas agregadas por slot
    \item Reducción de $\sim$450x en datos y verificación
\end{itemize}
\end{example}

%-------------------------------------------
\subsubsection{Uso en Zcash Sapling}

\begin{remark}
Zcash migró de BN128 (usado en Sprout) a BLS12-381 (en Sapling) por razones de seguridad. El protocolo Sapling usa:
\begin{itemize}
    \item \textbf{Groth16} sobre BLS12-381 para pruebas de transacciones privadas
    \item \textbf{RedDSA} (variante de Schnorr) sobre Jubjub (curva embebida en BLS12-381)
\end{itemize}
\end{remark}

\begin{definition}[Curva Jubjub]
Jubjub es una curva de Edwards embebida en BLS12-381:
\[ -x^2 + y^2 = 1 + d \cdot x^2 y^2 \]

Su orden es $n_{\text{BLS12-381}}$, lo que permite operaciones de curva elíptica \textbf{dentro} de circuitos SNARK sobre BLS12-381 de forma eficiente.
\end{definition}

%-------------------------------------------
\subsubsection{EIP-2537: Precompilados BLS12-381}

\begin{definition}[EIP-2537]
Propuesta para añadir precompilados BLS12-381 a Ethereum:

\begin{center}
\begin{tabular}{clc}
\toprule
\textbf{Dirección} & \textbf{Operación} & \textbf{Gas propuesto} \\
\midrule
0x0a & Suma $\G_1$ & 500 \\
0x0b & Multiplicación escalar $\G_1$ & 12,000 \\
0x0c & Multi-scalar multiplication $\G_1$ & Variable \\
0x0d & Suma $\G_2$ & 800 \\
0x0e & Multiplicación escalar $\G_2$ & 45,000 \\
0x0f & Multi-scalar multiplication $\G_2$ & Variable \\
0x10 & Pairing & 43,000 + 32,600 × pares \\
0x11 & Map to $\G_1$ & 5,500 \\
0x12 & Map to $\G_2$ & 75,000 \\
\bottomrule
\end{tabular}
\end{center}
\end{definition}

\begin{remark}
EIP-2537 ha estado en desarrollo desde 2020 pero aún no se ha incluido en un hard fork. Las razones incluyen:
\begin{itemize}
    \item Complejidad de implementación y testing
    \item Preocupaciones sobre aumento del tamaño del cliente
    \item Debate sobre si los costos de gas son apropiados
    \item Competencia con otras mejoras prioritarias
\end{itemize}

Mientras tanto, verificar BLS12-381 en EVM requiere implementación en Solidity puro ($\sim$millones de gas).
\end{remark}

%-------------------------------------------
\subsubsection{Rendimiento Comparativo}

\begin{center}
\begin{tabular}{lccc}
\toprule
\textbf{Operación} & \textbf{BN254} & \textbf{BLS12-381} & \textbf{Ratio} \\
\midrule
Suma $\G_1$ & 0.5 µs & 0.8 µs & 1.6x \\
Mult. escalar $\G_1$ & 50 µs & 80 µs & 1.6x \\
Suma $\G_2$ & 1.5 µs & 2.5 µs & 1.7x \\
Mult. escalar $\G_2$ & 150 µs & 250 µs & 1.7x \\
Pairing & 1.5 ms & 2 ms & 1.3x \\
\bottomrule
\end{tabular}
\end{center}

\begin{intuicion}
BLS12-381 es aproximadamente 1.3-1.7x más lento que BN254 para la mayoría de operaciones. Este overhead es aceptable considerando el aumento de seguridad de 100 a 128 bits.
\end{intuicion}

%-------------------------------------------
\subsubsection{Comparación Final: BN254 vs BLS12-381}

\begin{center}
\begin{tabular}{lcc}
\toprule
\textbf{Característica} & \textbf{BN254} & \textbf{BLS12-381} \\
\midrule
Seguridad & $\sim$100 bits & $\sim$128 bits \\
Campo base & 254 bits & 381 bits \\
Precompilados EVM & Sí (desde 2017) & No (propuesto) \\
Tamaño de $\G_1$ punto & 64 bytes & 96 bytes \\
Tamaño de $\G_2$ punto & 128 bytes & 192 bytes \\
Velocidad & Más rápida & $\sim$1.5x más lenta \\
Ethereum usage & SNARKs on-chain & Beacon Chain (off-EVM) \\
Zcash & Sprout (deprecado) & Sapling (actual) \\
Cofactor $h_1$ & 1 & Grande \\
Madurez & Alta & Alta \\
\bottomrule
\end{tabular}
\end{center}

%-------------------------------------------
\subsubsection{Aplicación en Banco ZK}

\begin{bancozk}
\textbf{Estrategia dual de curvas en Banco ZK:}

\textbf{1. Actualmente (BN254)}
\begin{itemize}
    \item Todos los verificadores on-chain usan BN254
    \item Aprovechamos precompilados para costos bajos
    \item Aceptamos 100 bits de seguridad como compromiso práctico
\end{itemize}

\textbf{2. Infraestructura preparada para BLS12-381}
\begin{itemize}
    \item Circuitos compilados para ambas curvas
    \item Verificadores BLS12-381 listos (esperando EIP-2537)
    \item Testing continuo en redes de prueba
\end{itemize}

\textbf{3. Casos de uso por curva}

\begin{center}
\begin{tabular}{lcc}
\toprule
\textbf{Componente} & \textbf{Curva} & \textbf{Razón} \\
\midrule
Verificación SNARK (L1) & BN254 & Precompilados disponibles \\
Agregación de firmas (L2) & BLS12-381 & Mejor seguridad, off-chain \\
Firmas de validadores & BLS12-381 & Estándar Ethereum 2.0 \\
Pruebas recursivas & BN254/BLS12-381 & Depende del caso \\
\bottomrule
\end{tabular}
\end{center}

\textbf{4. Plan de migración post-EIP-2537}

Cuando los precompilados BLS12-381 estén disponibles:
\begin{enumerate}
    \item Desplegar verificador BLS12-381
    \item Período de transición con ambos verificadores activos
    \item Deprecar verificador BN254 para nuevas pruebas
    \item Mantener BN254 solo para verificación de pruebas históricas
\end{enumerate}

\textbf{5. Consideraciones de diseño}

Para maximizar flexibilidad, nuestro contrato verificador es una interfaz:
\begin{verbatim}
interface IVerifier {
    function verify(
        bytes calldata proof,
        uint256[] calldata publicInputs
    ) external view returns (bool);
}

contract BN254Verifier is IVerifier { ... }
contract BLS12381Verifier is IVerifier { ... }  // Futuro
\end{verbatim}

El contrato principal puede actualizar la implementación del verificador sin cambiar la lógica de negocio.
\end{bancozk}

%-------------------------------------------
\subsubsection{Librerías y Herramientas para BLS12-381}

\begin{center}
\begin{tabular}{lll}
\toprule
\textbf{Lenguaje} & \textbf{Librería} & \textbf{Notas} \\
\midrule
Rust & arkworks (ark-bls12-381) & Referencia \\
Rust & blst & Muy optimizada (usado en Ethereum 2.0) \\
Go & gnark-crypto & Producción \\
Go & bls12-381 (herumi) & Optimizada \\
C & blst & Ensamblador optimizado \\
JavaScript & noble-bls12-381 & Audited, pura JS \\
Python & py\_ecc & Educacional \\
\bottomrule
\end{tabular}
\end{center}

\begin{remark}
La librería \texttt{blst} (``blast'') desarrollada por Supranational es considerada la implementación más rápida y segura. Es usada por la mayoría de clientes de Ethereum 2.0 (Lighthouse, Prysm, Teku, etc.).
\end{remark}

% termina bn128 
%===========================================
% SUBSECCIÓN: PASTA (PALLAS Y VESTA) - VERSIÓN EXPANDIDA
% Usa tus entornos: definition, example, keyidea, bancozk, 
%                   mathnote, intuicion, advertencia, remark
%===========================================

\subsection{Pasta (Pallas y Vesta)}

Ciclo de curvas para recursión eficiente.

\begin{definition}[Ciclo de curvas]
Dos curvas $E_1, E_2$ forman un ciclo si:
\begin{itemize}
    \item El orden de $E_1$ es el campo base de $E_2$
    \item El orden de $E_2$ es el campo base de $E_1$
\end{itemize}
\end{definition}

\begin{itemize}
    \item \textbf{Pallas:} $y^2 = x^3 + 5$ sobre $\Fp$
    \item \textbf{Vesta:} $y^2 = x^3 + 5$ sobre $\mathbb{F}_q$
    \item Donde $p = |E_{\text{Vesta}}|$ y $q = |E_{\text{Pallas}}|$
    \item Usadas en: Halo2, Mina, recursión de pruebas
\end{itemize}

%-------------------------------------------
\subsubsection{El Problema de la Recursión}

\begin{intuicion}
Imagina que quieres verificar una prueba ZK \textit{dentro} de otra prueba ZK. Por ejemplo:
\begin{enumerate}
    \item Tienes una prueba $\pi_1$ que demuestra ``Conozco $x$ tal que $f(x) = y$''
    \item Quieres crear una prueba $\pi_2$ que demuestra ``Verifiqué correctamente $\pi_1$''
    \item Y luego $\pi_3$ que demuestra ``Verifiqué $\pi_2$''...
\end{enumerate}

Esto es \textbf{recursión de pruebas}, fundamental para:
\begin{itemize}
    \item Comprimir cadenas largas de verificaciones
    \item Rollups que acumulan miles de transacciones
    \item Blockchains comprimidas (como Mina)
\end{itemize}
\end{intuicion}

\begin{advertencia}
\textbf{El problema técnico:}

Si usas una curva como BN254, verificar una prueba requiere operaciones de curva elíptica (multiplicaciones escalares, pairings). Pero dentro de un circuito SNARK sobre BN254, las operaciones de campo son módulo $p$ (el orden del grupo), no módulo el campo base de BN254.

Resultado: verificar una prueba BN254 dentro de un circuito BN254 requiere \textbf{aritmética de campo no nativa}, que es extremadamente costosa ($\sim$1,000,000+ constraints para una verificación).
\end{advertencia}

%-------------------------------------------
\subsubsection{La Solución: Ciclos de Curvas}

\begin{definition}[Ciclo de Curvas (2-cycle)]
Dos curvas elípticas $E_1$ y $E_2$ forman un \textbf{2-ciclo} si:
\begin{align}
\#E_1(\mathbb{F}_p) &= q \\
\#E_2(\mathbb{F}_q) &= p
\end{align}

donde $p$ es el campo base de $E_1$ y $q$ es el campo base de $E_2$.

Dicho de otra forma:
\begin{itemize}
    \item El \textbf{orden del grupo} de $E_1$ es el \textbf{campo base} de $E_2$
    \item El \textbf{orden del grupo} de $E_2$ es el \textbf{campo base} de $E_1$
\end{itemize}
\end{definition}

\begin{center}
\begin{tikzpicture}[
    curve/.style={draw, ellipse, minimum width=3cm, minimum height=2cm, thick},
    arrow/.style={->, thick, >=stealth}
]
    % Curva E1 (Pallas)
    \node[curve, fill=blue!10] (e1) at (0, 0) {};
    \node at (0, 0) {$E_1$ (Pallas)};
    \node[below=0.1cm of e1, font=\small] {Campo: $\mathbb{F}_p$};
    \node[above=0.1cm of e1, font=\small] {Orden: $q$};
    
    % Curva E2 (Vesta)
    \node[curve, fill=red!10] (e2) at (6, 0) {};
    \node at (6, 0) {$E_2$ (Vesta)};
    \node[below=0.1cm of e2, font=\small] {Campo: $\mathbb{F}_q$};
    \node[above=0.1cm of e2, font=\small] {Orden: $p$};
    
    % Flechas
    \draw[arrow, blue] (e1.north east) to[bend left=30] node[above] {orden $\to$ campo} (e2.north west);
    \draw[arrow, red] (e2.south west) to[bend left=30] node[below] {orden $\to$ campo} (e1.south east);
\end{tikzpicture}
\end{center}

\begin{keyidea}
Con un ciclo de curvas:
\begin{enumerate}
    \item Creas una prueba $\pi_1$ sobre $E_1$ (Pallas)
    \item Verificas $\pi_1$ \textbf{dentro} de un circuito sobre $E_2$ (Vesta)
    \item La verificación usa aritmética de $\mathbb{F}_q$, que es el campo nativo de $E_2$
    \item Las operaciones de curva de $E_1$ (que tienen orden $q$) son \textbf{nativas} en $E_2$
    \item Y viceversa para verificar pruebas de $E_2$ en circuitos de $E_1$
\end{enumerate}

No hay aritmética de campo no nativa. La verificación es eficiente.
\end{keyidea}

%-------------------------------------------
\subsubsection{Historia: De MNT a Pasta}

\begin{remark}
Los ciclos de curvas no son nuevos, pero encontrar ciclos con buenas propiedades criptográficas es difícil:

\textbf{2001 - Curvas MNT:} Miyaji, Nakabayashi, y Takano descubrieron familias de curvas que forman ciclos. Usadas en los primeros sistemas recursivos.

\textbf{2014 - Ben-Sasson et al.:} Proponen usar ciclos MNT para SNARKs recursivos en el paper ``Scalable Zero Knowledge via Cycles of Elliptic Curves''.

\textbf{2019 - Bowe et al. (Zcash):} Diseñan las curvas \textbf{Pasta} (Pallas + Vesta) específicamente para Halo 2, optimizadas para eficiencia moderna.

\textbf{2020 - Mina Protocol:} Usa Pasta para crear una blockchain de tamaño constante ($\sim$22 KB) mediante recursión.
\end{remark}

%-------------------------------------------
\subsubsection{Las Curvas Pasta: Pallas y Vesta}

\begin{definition}[Curvas Pasta]
\textbf{Pallas} (la ``verde'' ��):
\begin{align}
E_{\text{Pallas}}: y^2 &= x^3 + 5 \quad \text{sobre } \mathbb{F}_p \\
p &= 2^{254} + 45560315531419706090280762371685220353
\end{align}

\textbf{Vesta} (la ``roja'' ��):
\begin{align}
E_{\text{Vesta}}: y^2 &= x^3 + 5 \quad \text{sobre } \mathbb{F}_q \\
q &= 2^{254} + 45560315531506369815346746415080538113
\end{align}

La propiedad de ciclo:
\begin{align}
\#E_{\text{Pallas}}(\mathbb{F}_p) &= q \\
\#E_{\text{Vesta}}(\mathbb{F}_q) &= p
\end{align}
\end{definition}

\begin{mathnote}
\textbf{Parámetros explícitos:}

\begin{verbatim}
# Pallas
p = 0x40000000000000000000000000000000224698fc094cf91b992d30ed00000001
# Orden de Pallas = q (campo base de Vesta)
q = 0x40000000000000000000000000000000224698fc0994a8dd8c46eb2100000001

# Ecuación: y² = x³ + 5
a = 0
b = 5

# Generador de Pallas
Gx = -1 (mod p)
Gy = 2
\end{verbatim}

Ambas curvas tienen la misma ecuación ($y^2 = x^3 + 5$) pero sobre campos diferentes.
\end{mathnote}

\begin{intuicion}
\textbf{¿Por qué el nombre ``Pasta''?}

Es un juego de palabras: \textbf{Pa}llas + Ve\textbf{sta} = Pasta ��

Los nombres individuales vienen de asteroides/diosas:
\begin{itemize}
    \item \textbf{Pallas}: Segundo asteroide descubierto (1802), nombrado por Palas Atenea
    \item \textbf{Vesta}: Cuarto asteroide descubierto (1807), nombrado por la diosa romana del hogar
\end{itemize}

La comunidad ZK tiene tradición de nombres creativos (Groth16, PLONK, Halo, etc.).
\end{intuicion}

%-------------------------------------------
\subsubsection{¿Por Qué Ambas Curvas son $y^2 = x^3 + 5$?}

\begin{keyidea}
Tener la \textbf{misma ecuación} para ambas curvas simplifica enormemente:
\begin{enumerate}
    \item \textbf{Código compartido}: Las mismas fórmulas de punto funcionan para ambas
    \item \textbf{Menos errores}: Una implementación, dos curvas
    \item \textbf{Optimizaciones compartidas}: Mejoras aplican a ambas
\end{enumerate}

La única diferencia es el campo sobre el que operan.
\end{keyidea}

\begin{example}[Suma de puntos - mismo código]
\begin{verbatim}
// Funciona para Pallas Y Vesta
fn point_add<F: PrimeField>(p1: Point<F>, p2: Point<F>) -> Point<F> {
    let lambda = (p2.y - p1.y) / (p2.x - p1.x);
    let x3 = lambda * lambda - p1.x - p2.x;
    let y3 = lambda * (p1.x - x3) - p1.y;
    Point { x: x3, y: y3 }
}

// Para Pallas: F = Fp
// Para Vesta: F = Fq
\end{verbatim}
\end{example}

%-------------------------------------------
\subsubsection{Propiedades de Seguridad}

\begin{center}
\begin{tabular}{lcc}
\toprule
\textbf{Propiedad} & \textbf{Pallas} & \textbf{Vesta} \\
\midrule
Tamaño del campo & 254 bits & 254 bits \\
Orden del grupo & 254 bits & 254 bits \\
Seguridad ECDLP & $\sim$127 bits & $\sim$127 bits \\
Cofactor & 1 & 1 \\
Embedding degree & Alto & Alto \\
Discriminante CM & Grande & Grande \\
\bottomrule
\end{tabular}
\end{center}

\begin{keyidea}
Ambas curvas tienen \textbf{cofactor 1}, lo que significa:
\begin{itemize}
    \item Todo punto (excepto $\mathcal{O}$) tiene orden $n$ (el orden del grupo)
    \item No hay subgrupos pequeños que explotar
    \item No se necesita multiplicación por cofactor para seguridad
\end{itemize}

Esto simplifica implementaciones y evita clases enteras de ataques.
\end{keyidea}

\begin{remark}
Las curvas Pasta fueron diseñadas para \textbf{no} ser pairing-friendly (embedding degree alto). Esto es intencional:
\begin{enumerate}
    \item Evita ataques de MOV/FR que explotan pairings
    \item Simplifica la implementación (no hay código de pairing)
    \item Para sistemas como Halo 2, los pairings no son necesarios
\end{enumerate}
\end{remark}

%-------------------------------------------
\subsubsection{Recursión en Acción}

\begin{example}[Cadena de verificación recursiva]
\textbf{Objetivo:} Demostrar que verificaste correctamente 1000 pruebas.

\textbf{Sin recursión:}
\begin{itemize}
    \item Verificar 1000 pruebas on-chain
    \item Costo: $1000 \times 200,000 = 200$ millones de gas
    \item Imposible en un solo bloque
\end{itemize}

\textbf{Con recursión usando Pasta:}

\textbf{Paso 1:} Crear prueba $\pi_1$ sobre Pallas que demuestra ``Verifiqué pruebas 1-10''

\textbf{Paso 2:} Crear prueba $\pi_2$ sobre \textbf{Vesta} que demuestra:
\begin{itemize}
    \item ``$\pi_1$ es válida'' (verificación nativa, eficiente)
    \item ``También verifiqué pruebas 11-20''
\end{itemize}

\textbf{Paso 3:} Crear prueba $\pi_3$ sobre \textbf{Pallas} que demuestra:
\begin{itemize}
    \item ``$\pi_2$ es válida''
    \item ``También verifiqué pruebas 21-30''
\end{itemize}

... continúa alternando ...

\textbf{Paso 100:} Prueba final $\pi_{100}$ que demuestra todo.

\textbf{Resultado:}
\begin{itemize}
    \item Una sola prueba de tamaño constante
    \item Verifica 1000 pruebas originales
    \item Costo on-chain: $\sim$200,000 gas (igual que una prueba)
\end{itemize}
\end{example}

\begin{center}
\begin{tikzpicture}[
    proof/.style={draw, rectangle, rounded corners, minimum width=1.5cm, minimum height=0.8cm},
    arrow/.style={->, thick}
]
    % Pruebas
    \node[proof, fill=blue!20] (p1) at (0, 0) {$\pi_1$};
    \node[below=0.1cm of p1, font=\tiny] {Pallas};
    
    \node[proof, fill=red!20] (p2) at (2.5, 0) {$\pi_2$};
    \node[below=0.1cm of p2, font=\tiny] {Vesta};
    
    \node[proof, fill=blue!20] (p3) at (5, 0) {$\pi_3$};
    \node[below=0.1cm of p3, font=\tiny] {Pallas};
    
    \node[proof, fill=red!20] (p4) at (7.5, 0) {$\pi_4$};
    \node[below=0.1cm of p4, font=\tiny] {Vesta};
    
    \node at (9.5, 0) {$\cdots$};
    
    \node[proof, fill=green!20] (pn) at (11.5, 0) {$\pi_n$};
    \node[below=0.1cm of pn, font=\tiny] {Final};
    
    % Flechas
    \draw[arrow] (p1) -- node[above, font=\tiny] {verifica} (p2);
    \draw[arrow] (p2) -- node[above, font=\tiny] {verifica} (p3);
    \draw[arrow] (p3) -- node[above, font=\tiny] {verifica} (p4);
    \draw[arrow] (p4) -- (9, 0);
    \draw[arrow] (10, 0) -- (pn);
\end{tikzpicture}
\end{center}

%-------------------------------------------
\subsubsection{Halo 2: El Sistema que Usa Pasta}

\begin{definition}[Halo 2]
Halo 2 es un sistema de pruebas desarrollado por Zcash que:
\begin{itemize}
    \item Usa las curvas Pasta para recursión
    \item No requiere trusted setup (es ``transparent'')
    \item Basado en el esquema de compromiso IPA (Inner Product Argument)
    \item Implementa un sistema de constraints llamado PLONKish
\end{itemize}
\end{definition}

\begin{keyidea}
\textbf{¿Por qué Halo 2 no necesita pairings?}

Los sistemas como Groth16 usan pairings para verificación:
\[ e(A, B) = e(C, D) \cdot e(E, F) \cdots \]

Halo 2 usa una técnica diferente basada en \textbf{acumuladores}:
\begin{enumerate}
    \item En lugar de verificar completamente cada prueba, ``acumula'' el trabajo
    \item Las verificaciones parciales se combinan en un ``acumulador''
    \item Solo al final se hace una verificación completa
    \item Esta verificación final es eficiente y no requiere pairings
\end{enumerate}

Esto evita la necesidad de curvas pairing-friendly, permitiendo usar curvas más simples como Pasta.
\end{keyidea}

\begin{remark}
El paper original ``Halo: Recursive Proof Composition without a Trusted Setup'' (2019) de Bowe, Grigg, y Hopwood fue revolucionario al demostrar que la recursión eficiente era posible sin pairings y sin trusted setup.
\end{remark}

%-------------------------------------------
\subsubsection{Mina Protocol: Blockchain de 22 KB}

\begin{intuicion}
Mina Protocol usa Pasta para crear una blockchain donde el estado completo siempre es de $\sim$22 KB, independientemente del historial.

\textbf{¿Cómo?}

Cada bloque contiene una prueba recursiva que dice:
``El bloque anterior era válido, y esta nueva transición de estado también es válida''

Verificar la cadena completa = verificar UNA prueba.
\end{intuicion}

\begin{example}[Comparación de tamaños de blockchain]
\begin{center}
\begin{tabular}{lcc}
\toprule
\textbf{Blockchain} & \textbf{Tamaño para verificar} & \textbf{Tiempo sync} \\
\midrule
Bitcoin & $\sim$500 GB & Horas/días \\
Ethereum & $\sim$1 TB+ & Días \\
Mina & $\sim$22 KB & Segundos \\
\bottomrule
\end{tabular}
\end{center}

Mina logra esto acumulando todas las verificaciones en una sola prueba recursiva usando Pasta.
\end{example}

%-------------------------------------------
\subsubsection{Comparación con Otros Ciclos}

\begin{center}
\begin{tabular}{lccc}
\toprule
\textbf{Ciclo} & \textbf{Tamaño (bits)} & \textbf{Seguridad} & \textbf{Uso} \\
\midrule
MNT4/MNT6 & 298/753 & $\sim$80/$\sim$128 bits & Celo, antiguo Zcash \\
Pasta (Pallas/Vesta) & 254/254 & $\sim$127 bits & Halo 2, Mina \\
Curve-5/Curve-6 & Variable & Variable & Experimental \\
\bottomrule
\end{tabular}
\end{center}

\begin{keyidea}
Ventajas de Pasta sobre MNT:
\begin{enumerate}
    \item \textbf{Tamaños balanceados}: Ambas curvas de 254 bits (MNT tiene una curva mucho más grande)
    \item \textbf{Misma ecuación}: Código compartido
    \item \textbf{Cofactor 1}: Sin complicaciones de subgrupos
    \item \textbf{Diseño moderno}: Optimizado para hardware actual
\end{enumerate}
\end{keyidea}

%-------------------------------------------
\subsubsection{Implementación: Coordinación de Campos}

\begin{mathnote}
\textbf{El truco de implementación:}

Cuando trabajas con Pasta, constantemente cambias entre dos contextos:
\begin{itemize}
    \item Aritmética en $\mathbb{F}_p$ (campo base de Pallas, campo escalar de Vesta)
    \item Aritmética en $\mathbb{F}_q$ (campo base de Vesta, campo escalar de Pallas)
\end{itemize}

Un elemento puede tener dos ``interpretaciones'':
\begin{itemize}
    \item Como coordenada de un punto Pallas (vive en $\mathbb{F}_p$)
    \item Como escalar para multiplicación en Vesta (vive en $\mathbb{F}_q = $ orden de Pallas)
\end{itemize}

Pero $p \neq q$ (difieren en $\sim 2^{126}$), así que hay que tener cuidado con las conversiones.
\end{mathnote}

\begin{example}[Estructura de datos en Rust (Halo 2)]
\begin{verbatim}
use pasta_curves::{Fp, Fq, pallas, vesta};

// Punto en Pallas - coordenadas en Fp
let pallas_point: pallas::Point = pallas::Point::generator();

// Escalar para Pallas - vive en Fq (el orden de Pallas)
let pallas_scalar: Fq = Fq::from(42u64);

// Multiplicación escalar en Pallas
let result = pallas_point * pallas_scalar;

// Punto en Vesta - coordenadas en Fq  
let vesta_point: vesta::Point = vesta::Point::generator();

// Escalar para Vesta - vive en Fp (el orden de Vesta)
let vesta_scalar: Fp = Fp::from(42u64);

// Aquí está la magia del ciclo:
// Fp es TANTO el campo base de Pallas COMO el campo escalar de Vesta
// Fq es TANTO el campo base de Vesta COMO el campo escalar de Pallas
\end{verbatim}
\end{example}

%-------------------------------------------
\subsubsection{Consideraciones de Seguridad}

\begin{advertencia}
\textbf{Riesgos específicos de ciclos de curvas:}

\textbf{1. Curvas relacionadas:}
Un ataque que debilite una curva podría afectar a la otra, ya que comparten estructura algebraica.

\textbf{2. Errores de implementación:}
Confundir $\mathbb{F}_p$ con $\mathbb{F}_q$ puede causar bugs sutiles que compilan pero son incorrectos.

\textbf{3. Side-channels amplificados:}
Si hay un leak de información en una curva, podría revelar información sobre la otra.

\textbf{Mitigaciones:}
\begin{itemize}
    \item Tipos de datos fuertes que distinguen $\mathbb{F}_p$ de $\mathbb{F}_q$
    \item Tests extensivos de ida y vuelta
    \item Auditorías de seguridad específicas para sistemas recursivos
\end{itemize}
\end{advertencia}

%-------------------------------------------
\subsubsection{Rendimiento}

\begin{center}
\begin{tabular}{lcc}
\toprule
\textbf{Operación} & \textbf{Pallas/Vesta} & \textbf{Comparación} \\
\midrule
Suma de puntos & $\sim$0.3 µs & Similar a secp256k1 \\
Multiplicación escalar & $\sim$40 µs & Similar a secp256k1 \\
MSM (1024 puntos) & $\sim$15 ms & Optimizado para Halo 2 \\
Verificación recursiva & $\sim$10-50 ms & Depende del circuito \\
\bottomrule
\end{tabular}
\end{center}

\begin{remark}
Las curvas Pasta están optimizadas para operaciones de \textbf{Multi-Scalar Multiplication (MSM)}, que dominan el tiempo de prueba en sistemas como Halo 2. Implementaciones como la de Zcash usan:
\begin{itemize}
    \item Algoritmo de Pippenger para MSM
    \item Paralelización en múltiples cores
    \item Representaciones de punto optimizadas (coordenadas extendidas)
\end{itemize}
\end{remark}

%-------------------------------------------
\subsubsection{Futuro: Más Allá de 2-Ciclos}

\begin{intuicion}
Los 2-ciclos (como Pasta) alternan entre dos curvas. Pero hay investigación en:
\begin{itemize}
    \item \textbf{1-ciclos}: Una sola curva que ``cicla'' consigo misma (teóricamente posible pero no encontrado aún con buena seguridad)
    \item \textbf{k-ciclos}: Cadenas más largas de curvas
    \item \textbf{Ciclos con pairings}: Combinar recursión eficiente con verificación por pairings
\end{itemize}
\end{intuicion}

\begin{remark}
Encontrar un 1-ciclo seguro sería un avance significativo: eliminaría la necesidad de alternar entre curvas y simplificaría aún más las implementaciones. Sin embargo, las restricciones matemáticas hacen esto muy difícil.
\end{remark}

%-------------------------------------------
\subsubsection{Aplicación en Banco ZK}

\begin{bancozk}
Las curvas Pasta son parte de nuestra estrategia de recursión:

\textbf{1. Agregación de pruebas}

En lugar de verificar cada prueba de transacción individualmente on-chain:
\begin{enumerate}
    \item Generamos pruebas individuales sobre Pallas
    \item Las agregamos recursivamente alternando Pallas/Vesta
    \item Enviamos una sola prueba agregada al L1
\end{enumerate}

\textbf{Beneficio:} 1000 transacciones $\to$ 1 verificación on-chain.

\textbf{2. Arquitectura del sistema}

\begin{center}
\begin{tabular}{lll}
\toprule
\textbf{Capa} & \textbf{Curva} & \textbf{Función} \\
\midrule
Transacción individual & Pallas & Prueba de validez \\
Agregación nivel 1 & Vesta & Combinar 10 pruebas \\
Agregación nivel 2 & Pallas & Combinar 10 agregados \\
$\vdots$ & Alterna & $\vdots$ \\
Prueba final & BN254 & Para verificar on-chain \\
\bottomrule
\end{tabular}
\end{center}

\textbf{3. Conversión Pasta $\to$ BN254}

La prueba final debe verificarse en Ethereum (que tiene precompilados BN254). Hacemos una ``traducción'' final:
\begin{enumerate}
    \item Prueba recursiva final está sobre Pallas o Vesta
    \item Creamos una prueba BN254 que dice ``Verifiqué la prueba Pasta''
    \item Esta verificación incluye emular aritmética de Pasta en BN254
    \item Costoso ($\sim$1M constraints), pero solo se hace una vez
\end{enumerate}

\textbf{4. Trade-offs considerados}

\begin{center}
\begin{tabular}{lcc}
\toprule
\textbf{Opción} & \textbf{Ventaja} & \textbf{Desventaja} \\
\midrule
Solo BN254 & Simple, directo a L1 & Recursión cara \\
Pasta + traducción & Recursión eficiente & Traducción costosa \\
Solo Pasta & Más eficiente & No hay precompilados L1 \\
\bottomrule
\end{tabular}
\end{center}

Elegimos Pasta + traducción porque la eficiencia de recursión compensa el costo de traducción cuando hay muchas transacciones.

\textbf{5. Código de ejemplo}

\begin{verbatim}
// Estructura de agregación recursiva
struct RecursiveProof {
    // Punto de acumulación en Pallas o Vesta
    accumulator: CurvePoint,
    // Prueba de que el acumulador es válido
    proof: Proof,
    // Metadata
    num_proofs_accumulated: u64,
}

fn aggregate(proofs: Vec<Proof>) -> RecursiveProof {
    let mut acc = RecursiveProof::empty();
    
    for (i, proof) in proofs.iter().enumerate() {
        if i % 2 == 0 {
            // Verificar en circuito Vesta
            acc = verify_and_accumulate_on_vesta(acc, proof);
        } else {
            // Verificar en circuito Pallas
            acc = verify_and_accumulate_on_pallas(acc, proof);
        }
    }
    
    acc
}
\end{verbatim}
\end{bancozk}

%-------------------------------------------
\subsubsection{Librerías y Herramientas}

\begin{center}
\begin{tabular}{lll}
\toprule
\textbf{Librería} & \textbf{Lenguaje} & \textbf{Notas} \\
\midrule
pasta\_curves & Rust & Implementación oficial de Zcash \\
halo2 & Rust & Sistema de pruebas completo \\
arkworks (ark-pallas) & Rust & Parte del ecosistema arkworks \\
o1js (SnarkyJS) & TypeScript & Para desarrollo en Mina \\
\bottomrule
\end{tabular}
\end{center}

\begin{example}[Uso básico de pasta\_curves]
\begin{verbatim}
use pasta_curves::{
    pallas, vesta,
    group::Group,
    arithmetic::CurveAffine,
};

fn main() {
    // Generadores
    let g_pallas = pallas::Point::generator();
    let g_vesta = vesta::Point::generator();
    
    // Escalares (notar los tipos cruzados)
    let scalar_for_pallas: pallas::Scalar = pallas::Scalar::from(42u64);
    let scalar_for_vesta: vesta::Scalar = vesta::Scalar::from(42u64);
    
    // Multiplicación escalar
    let p = g_pallas * scalar_for_pallas;
    let v = g_vesta * scalar_for_vesta;
    
    // Conversión a coordenadas afines
    let p_affine = p.to_affine();
    println!("Pallas point: ({:?}, {:?})", p_affine.x, p_affine.y);
}
\end{verbatim}
\end{example}

%-------------------------------------------
\subsubsection{Constantes de Referencia}

\begin{mathnote}
\textbf{Resumen de constantes Pasta:}

\begin{verbatim}
# Pallas
p_pallas = 0x40000000000000000000000000000000224698fc094cf91b992d30ed00000001
# = 2^254 + 45560315531419706090280762371685220353

# Orden de Pallas (= campo base de Vesta)
q_pallas = 0x40000000000000000000000000000000224698fc0994a8dd8c46eb2100000001
# = 2^254 + 45560315531506369815346746415080538113

# Ecuación de ambas curvas
a = 0
b = 5

# Generador de Pallas
Gx_pallas = p_pallas - 1  # = -1 mod p
Gy_pallas = 2

# Generador de Vesta  
Gx_vesta = q_pallas - 1   # = -1 mod q
Gy_vesta = 2

# Verificación del ciclo
assert orden_pallas == campo_base_vesta
assert orden_vesta == campo_base_pallas
\end{verbatim}
\end{mathnote}

% ----- FIRMAS DIGITALES -----
\section{Firmas Digitales}

%===========================================
% SUBSECCIÓN: ECDSA - VERSIÓN EXPANDIDA
% Usa tus entornos: definition, example, keyidea, bancozk, 
%                   mathnote, intuicion, advertencia, remark
%===========================================

\subsection{ECDSA}

\begin{definition}[ECDSA]
Elliptic Curve Digital Signature Algorithm.
\end{definition}

Algoritmo:
\begin{enumerate}
    \item \textbf{Keygen:} $sk \leftarrow \Z_n$, $pk = sk \cdot G$
    \item \textbf{Sign:}
    \begin{itemize}
        \item Elegir nonce $k \leftarrow \Z_n$
        \item $R = kG$, sea $r = x_R \mod n$
        \item $s = k^{-1}(H(m) + r \cdot sk) \mod n$
        \item Firma: $(r, s)$
    \end{itemize}
    \item \textbf{Verify:}
    \begin{itemize}
        \item $u_1 = H(m) \cdot s^{-1}$, $u_2 = r \cdot s^{-1}$
        \item Verificar $x_{u_1 G + u_2 \cdot pk} \equiv r \pmod{n}$
    \end{itemize}
\end{enumerate}

\begin{keyidea}
\textbf{Peligro crítico:} Reutilizar el nonce $k$ expone la clave privada.
Si dos firmas $(r, s_1)$ y $(r, s_2)$ usan el mismo $k$:
\[
sk = \frac{s_1 H(m_2) - s_2 H(m_1)}{r(s_2 - s_1)}
\]
Esto causó pérdidas millonarias en la historia de Bitcoin.
\end{keyidea}

%-------------------------------------------
\subsubsection{Contexto Histórico}

\begin{intuicion}
Las firmas digitales son el equivalente criptográfico de una firma manuscrita: demuestran que el poseedor de cierta clave privada autorizó un mensaje específico. ECDSA es la variante de curva elíptica del algoritmo DSA (Digital Signature Algorithm), estandarizado por NIST en 1991.
\end{intuicion}

\begin{remark}
\textbf{Línea temporal de las firmas digitales:}
\begin{center}
\begin{tabular}{ll}
\toprule
\textbf{Año} & \textbf{Desarrollo} \\
\midrule
1976 & Diffie-Hellman introduce firmas digitales (concepto) \\
1978 & RSA - primera implementación práctica \\
1985 & ElGamal propone firmas basadas en logaritmo discreto \\
1991 & DSA estandarizado por NIST (FIPS 186) \\
1992 & Vanstone propone ECDSA \\
1999 & ECDSA incluido en ANSI X9.62 \\
2005 & ECDSA en FIPS 186-3 \\
2009 & Bitcoin usa ECDSA sobre secp256k1 \\
2015 & Ethereum adopta el mismo esquema \\
\bottomrule
\end{tabular}
\end{center}
\end{remark}

%-------------------------------------------
\subsubsection{¿Por Qué ECDSA sobre DSA Clásico?}

\begin{keyidea}
ECDSA ofrece la misma seguridad que DSA con claves mucho más cortas:

\begin{center}
\begin{tabular}{ccc}
\toprule
\textbf{Seguridad} & \textbf{RSA/DSA} & \textbf{ECDSA} \\
\midrule
80 bits & 1024 bits & 160 bits \\
112 bits & 2048 bits & 224 bits \\
128 bits & 3072 bits & 256 bits \\
256 bits & 15360 bits & 512 bits \\
\bottomrule
\end{tabular}
\end{center}

Para 128 bits de seguridad:
\begin{itemize}
    \item Clave RSA: 3072 bits = 384 bytes
    \item Clave ECDSA: 256 bits = 32 bytes (12x más pequeña)
\end{itemize}
\end{keyidea}

%-------------------------------------------
\subsubsection{El Algoritmo en Detalle}

\paragraph{Parámetros del Sistema}

\begin{definition}[Parámetros de Dominio ECDSA]
Para usar ECDSA, se fijan:
\begin{itemize}
    \item Curva elíptica $E$ sobre $\Fp$
    \item Punto generador $G \in E(\Fp)$ de orden primo $n$
    \item Función hash $H: \{0,1\}^* \to \Z_n$
\end{itemize}

Para Bitcoin/Ethereum (secp256k1):
\begin{itemize}
    \item $E: y^2 = x^3 + 7$ sobre $\Fp$
    \item $p = 2^{256} - 2^{32} - 977$
    \item $n = $ \texttt{0xFFFFFFFFFFFFFFFFFFFFFFFFFFFFFFFE...}
    \item $H = $ SHA-256 (Bitcoin) o Keccak-256 (Ethereum)
\end{itemize}
\end{definition}

\paragraph{Generación de Claves}

\begin{definition}[KeyGen]
\begin{enumerate}
    \item Elegir $sk \xleftarrow{\$} [1, n-1]$ uniformemente aleatorio
    \item Calcular $pk = sk \cdot G$
    \item Publicar $pk$, mantener $sk$ secreto
\end{enumerate}
\end{definition}

\begin{intuicion}
La clave pública $pk$ es un punto en la curva. Aunque el atacante conoce $pk$ y $G$, no puede encontrar $sk$ porque esto requeriría resolver el problema del logaritmo discreto en curvas elípticas (ECDLP), que es computacionalmente intratable.
\end{intuicion}

\paragraph{Firma}

\begin{definition}[Sign$(sk, m)$]
Para firmar un mensaje $m$ con clave privada $sk$:

\begin{enumerate}
    \item Calcular el hash del mensaje: $e = H(m)$
    
    \item Elegir nonce aleatorio: $k \xleftarrow{\$} [1, n-1]$
    
    \item Calcular el punto efímero: $R = k \cdot G = (x_R, y_R)$
    
    \item Calcular $r = x_R \mod n$
    \begin{itemize}
        \item Si $r = 0$, volver al paso 2
    \end{itemize}
    
    \item Calcular $s = k^{-1}(e + r \cdot sk) \mod n$
    \begin{itemize}
        \item Si $s = 0$, volver al paso 2
    \end{itemize}
    
    \item Retornar firma $\sigma = (r, s)$
\end{enumerate}
\end{definition}

\begin{mathnote}
\textbf{Descomposición de la fórmula de $s$:}
\[ s = k^{-1}(e + r \cdot sk) \mod n \]

\begin{itemize}
    \item $e = H(m)$: el mensaje ``entra'' en la firma
    \item $r \cdot sk$: enlaza con la clave privada
    \item $k^{-1}$: ``oculta'' los componentes usando el nonce
\end{itemize}

La seguridad depende de que $k$ sea secreto y único para cada firma.
\end{mathnote}

\paragraph{Verificación}

\begin{definition}[Verify$(pk, m, \sigma)$]
Para verificar una firma $\sigma = (r, s)$ sobre mensaje $m$ con clave pública $pk$:

\begin{enumerate}
    \item Verificar que $r, s \in [1, n-1]$
    
    \item Calcular $e = H(m)$
    
    \item Calcular $w = s^{-1} \mod n$
    
    \item Calcular $u_1 = e \cdot w \mod n$
    
    \item Calcular $u_2 = r \cdot w \mod n$
    
    \item Calcular el punto $R' = u_1 \cdot G + u_2 \cdot pk$
    
    \item Verificar que $x_{R'} \equiv r \pmod{n}$
\end{enumerate}

Si la verificación pasa, la firma es válida.
\end{definition}

%-------------------------------------------
\subsubsection{¿Por Qué Funciona?}

\begin{mathnote}
\textbf{Demostración de correctitud:}

Queremos mostrar que $x_{R'} = r$ cuando la firma es válida.

\textbf{Partimos de:}
\[ R' = u_1 \cdot G + u_2 \cdot pk \]

\textbf{Sustituyendo:}
\begin{align}
R' &= (e \cdot w) \cdot G + (r \cdot w) \cdot pk \\
   &= w \cdot (e \cdot G + r \cdot pk) \\
   &= s^{-1} \cdot (e \cdot G + r \cdot sk \cdot G) \quad \text{(porque } pk = sk \cdot G\text{)} \\
   &= s^{-1} \cdot (e + r \cdot sk) \cdot G
\end{align}

\textbf{Recordando que:} $s = k^{-1}(e + r \cdot sk)$, entonces $s^{-1} = k \cdot (e + r \cdot sk)^{-1}$

\begin{align}
R' &= k \cdot (e + r \cdot sk)^{-1} \cdot (e + r \cdot sk) \cdot G \\
   &= k \cdot G \\
   &= R
\end{align}

Por lo tanto, $x_{R'} = x_R = r$. $\square$
\end{mathnote}

%-------------------------------------------
\subsubsection{Ejemplo Numérico Completo}

\begin{example}[ECDSA sobre una curva pequeña]
Usaremos una curva pequeña para ilustrar (NO segura para uso real).

\textbf{Parámetros:}
\begin{itemize}
    \item Curva: $E: y^2 = x^3 + 7$ sobre $\mathbb{F}_{17}$
    \item Generador: $G = (6, 11)$
    \item Orden: $n = 18$ (simplificado)
\end{itemize}

\textbf{1. Generación de claves:}
\begin{align}
sk &= 7 \quad \text{(elegido aleatoriamente)} \\
pk &= 7 \cdot G = 7 \cdot (6, 11) = (13, 10)
\end{align}

\textbf{2. Firma del mensaje $m$ donde $H(m) = 5$:}
\begin{align}
k &= 11 \quad \text{(nonce aleatorio)} \\
R &= 11 \cdot G = (7, 8) \\
r &= x_R = 7 \\
k^{-1} &= 11^{-1} \mod 18 = 5 \quad \text{(porque } 11 \cdot 5 = 55 \equiv 1 \mod 18\text{)} \\
s &= k^{-1}(e + r \cdot sk) = 5(5 + 7 \cdot 7) = 5(5 + 49) = 5 \cdot 54 \\
  &= 5 \cdot 54 \mod 18 = 270 \mod 18 = 0
\end{align}

¡$s = 0$! Esto es inválido, debemos elegir otro $k$.

\textbf{2b. Segundo intento con $k = 13$:}
\begin{align}
R &= 13 \cdot G = (10, 2) \\
r &= 10 \\
k^{-1} &= 13^{-1} \mod 18 = 7 \quad \text{(porque } 13 \cdot 7 = 91 \equiv 1 \mod 18\text{)} \\
s &= 7(5 + 10 \cdot 7) = 7(5 + 70) = 7 \cdot 75 = 7 \cdot 3 = 21 \equiv 3 \mod 18
\end{align}

\textbf{Firma:} $(r, s) = (10, 3)$

\textbf{3. Verificación:}
\begin{align}
w &= s^{-1} = 3^{-1} \mod 18 = 6 \quad \text{(porque } 3 \cdot 6 = 18 \equiv 0... \\
\end{align}

Hmm, $3 \cdot 6 = 18 \equiv 0 \mod 18$, así que $3^{-1}$ no existe... Esto es porque $\gcd(3, 18) = 3 \neq 1$.

\textbf{Nota:} Este ejemplo simplificado tiene problemas porque $n = 18$ no es primo. En la práctica, $n$ siempre es primo, garantizando que todos los elementos no nulos tienen inverso.
\end{example}

\begin{example}[ECDSA con parámetros realistas (valores truncados)]
\textbf{Usando secp256k1:}

\textbf{Clave privada:}
\begin{verbatim}
sk = 0xc9afa9d845ba75166b5c215767b1d6934e50c3db36e89b127b8a622b120f6721
\end{verbatim}

\textbf{Clave pública:} $pk = sk \cdot G$
\begin{verbatim}
pk.x = 0x60fed4ba255a9d31c961eb74c6356d68c049b8923b61fa6ce669622e60f29fb6
pk.y = 0x7903fe1008b8bc99a41ae9e95628bc64f2f1b20c2d7e9f5177a3c294d4462299
\end{verbatim}

\textbf{Mensaje:} ``Hello, ECDSA!''
\begin{verbatim}
e = SHA256("Hello, ECDSA!") 
  = 0x7846e3be8abd2e089ed812e9d8c5e23a5ad76d5af9d65e0f9d89e8c76a5f8b12
\end{verbatim}

\textbf{Firma (usando RFC 6979 para $k$ determinístico):}
\begin{verbatim}
r = 0x2b826c8c5f1e8e7d3a9b2c4d5e6f7a8b9c0d1e2f3a4b5c6d7e8f9a0b1c2d3e4f
s = 0x1a2b3c4d5e6f7a8b9c0d1e2f3a4b5c6d7e8f9a0b1c2d3e4f5a6b7c8d9e0f1a2b
\end{verbatim}

\textbf{Verificación:}
\begin{verbatim}
w = s^(-1) mod n
u1 = e * w mod n
u2 = r * w mod n
R' = u1*G + u2*pk
assert R'.x == r  # ✓ Firma válida
\end{verbatim}
\end{example}

%-------------------------------------------
\subsubsection{El Peligro Mortal: Reutilización del Nonce}

\begin{advertencia}
\textbf{LA REGLA MÁS IMPORTANTE DE ECDSA:}

\begin{center}
\textbf{NUNCA, JAMÁS, BAJO NINGUNA CIRCUNSTANCIA, REUTILIZAR EL NONCE $k$}
\end{center}

Si dos firmas usan el mismo $k$, la clave privada queda expuesta inmediatamente.
\end{advertencia}

\begin{mathnote}
\textbf{Demostración del ataque:}

Supongamos que firmamos dos mensajes $m_1, m_2$ con el mismo nonce $k$:
\begin{align}
s_1 &= k^{-1}(e_1 + r \cdot sk) \mod n \\
s_2 &= k^{-1}(e_2 + r \cdot sk) \mod n
\end{align}

Nota que $r$ es igual en ambas (porque $R = kG$ es igual).

\textbf{Restando:}
\begin{align}
s_1 - s_2 &= k^{-1}(e_1 - e_2) \mod n
\end{align}

\textbf{Despejando $k$:}
\begin{align}
k &= (e_1 - e_2)(s_1 - s_2)^{-1} \mod n
\end{align}

\textbf{Una vez conocido $k$, recuperamos $sk$:}
\begin{align}
sk &= (s_1 \cdot k - e_1) \cdot r^{-1} \mod n
\end{align}

\textbf{O equivalentemente:}
\[
sk = \frac{s_1 e_2 - s_2 e_1}{r(s_2 - s_1)} \mod n
\]

El atacante solo necesita ver dos firmas con el mismo $r$ para robar todos los fondos.
\end{mathnote}

\begin{example}[El hack de Sony PlayStation 3 (2010)]
Sony usaba ECDSA para firmar el software de PS3. Su error fatal:

\begin{verbatim}
// Código de Sony (simplificado)
k = 4;  // ¡CONSTANTE! El mismo para TODAS las firmas
\end{verbatim}

El grupo ``fail0verflow'' observó que todas las firmas tenían el mismo valor $r$.
Aplicando el ataque descrito, recuperaron la clave privada de Sony.

\textbf{Resultado:} Cualquiera podía firmar software como si fuera Sony, permitiendo piratería masiva y homebrew.

\textbf{Lección:} Sony perdió el control de su plataforma por un error de una línea de código.
\end{example}

\begin{example}[Pérdidas en Bitcoin por nonces repetidos]
\textbf{Android SecureRandom Bug (2013):}

El generador de números aleatorios de Android en ciertas versiones producía outputs predecibles. Múltiples wallets Android generaron nonces idénticos para diferentes transacciones.

\textbf{Resultado:}
\begin{itemize}
    \item Miles de claves privadas comprometidas
    \item 55 BTC robados en las primeras horas del descubrimiento
    \item Todas las wallets afectadas vaciadas eventualmente
\end{itemize}

\textbf{Blockchain.info Bug (2014):}

Un error en el código de generación de nonces causó colisiones.

\textbf{Resultado:}
\begin{itemize}
    \item Cientos de BTC robados
    \item La empresa tuvo que compensar a usuarios
\end{itemize}
\end{example}

%-------------------------------------------
\subsubsection{Solución: Nonces Determinísticos (RFC 6979)}

\begin{definition}[RFC 6979]
Para evitar la catastrófica reutilización de nonces, RFC 6979 define un método \textbf{determinístico} para generar $k$:

\[ k = \text{HMAC-DRBG}(sk, H(m)) \]

El nonce se deriva de:
\begin{itemize}
    \item La clave privada $sk$
    \item El hash del mensaje $H(m)$
\end{itemize}

\textbf{Propiedades:}
\begin{enumerate}
    \item Mismo mensaje + misma clave = misma firma (determinístico)
    \item Diferente mensaje = diferente $k$ (único)
    \item Sin necesidad de generador aleatorio (portable)
    \item Imposible para atacante predecir $k$ sin conocer $sk$
\end{enumerate}
\end{definition}

\begin{intuicion}
RFC 6979 elimina completamente el riesgo de nonce repetido:
\begin{itemize}
    \item Si firmas el mismo mensaje dos veces, obtienes la misma firma (no revela nada nuevo)
    \item Si firmas mensajes diferentes, los nonces son diferentes (seguro)
    \item No dependes de un generador aleatorio que podría fallar
\end{itemize}

Hoy en día, \textbf{todas} las implementaciones serias usan RFC 6979 o similar.
\end{intuicion}

\begin{example}[Pseudocódigo RFC 6979]
\begin{verbatim}
def generate_k_rfc6979(sk, message):
    # Paso 1: Hash del mensaje
    h1 = SHA256(message)
    
    # Paso 2: Inicializar HMAC-DRBG
    V = b'\x01' * 32
    K = b'\x00' * 32
    
    # Paso 3: Actualizar con clave privada y hash
    K = HMAC_SHA256(K, V + b'\x00' + sk + h1)
    V = HMAC_SHA256(K, V)
    K = HMAC_SHA256(K, V + b'\x01' + sk + h1)
    V = HMAC_SHA256(K, V)
    
    # Paso 4: Generar k
    while True:
        V = HMAC_SHA256(K, V)
        k = int.from_bytes(V, 'big')
        if 1 <= k < n:
            return k
        K = HMAC_SHA256(K, V + b'\x00')
        V = HMAC_SHA256(K, V)
\end{verbatim}
\end{example}

%-------------------------------------------
\subsubsection{Maleabilidad de Firmas}

\begin{advertencia}
ECDSA tiene un problema de \textbf{maleabilidad}: dada una firma válida $(r, s)$, cualquiera puede crear otra firma válida $(r, n - s)$ para el mismo mensaje.

\textbf{¿Por qué?}
\[ s' = n - s = -s \mod n \]

La verificación usa $s^{-1}$:
\[ (s')^{-1} = (-s)^{-1} = -s^{-1} \mod n \]

Los signos se cancelan y la verificación pasa.
\end{advertencia}

\begin{keyidea}
\textbf{Impacto de la maleabilidad:}

En Bitcoin, cada transacción se identifica por su hash (txid). Si alguien modifica la firma de una transacción no confirmada:
\begin{itemize}
    \item La transacción sigue siendo válida
    \item Pero el txid cambia
    \item Esto puede romper cadenas de transacciones dependientes
\end{itemize}

\textbf{Solución (BIP-62, BIP-146):}

Bitcoin requiere que $s \leq n/2$ (``low-S''). Esto elimina la maleabilidad porque solo una de $s$ o $n-s$ puede ser $\leq n/2$.
\end{keyidea}

\begin{example}[Normalización low-S]
\begin{verbatim}
def normalize_signature(r, s, n):
    """Asegurar que s <= n/2 para evitar maleabilidad"""
    if s > n // 2:
        s = n - s
    return (r, s)

# Ejemplo
n = 0xFFFFFFFFFFFFFFFFFFFFFFFFFFFFFFFEBAAEDCE6AF48A03BBFD25E8CD0364141
s = 0xDEADBEEFDEADBEEFDEADBEEFDEADBEEFDEADBEEFDEADBEEFDEADBEEFDEADBEEF

if s > n // 2:
    s_normalized = n - s
    print("Firma normalizada a low-S")
\end{verbatim}
\end{example}

%-------------------------------------------
\subsubsection{Formato de Firma}

\begin{definition}[Formatos de Firma ECDSA]
\textbf{DER (Distinguished Encoding Rules):}
\begin{verbatim}
30 <length>
   02 <r-length> <r>
   02 <s-length> <s>
\end{verbatim}
Tamaño variable: 70-72 bytes típicamente.

\textbf{Compacto (raw):}
\begin{verbatim}
<r: 32 bytes> <s: 32 bytes>
\end{verbatim}
Tamaño fijo: 64 bytes.

\textbf{Ethereum (con recovery id):}
\begin{verbatim}
<r: 32 bytes> <s: 32 bytes> <v: 1 byte>
\end{verbatim}
Tamaño: 65 bytes.
\end{definition}

\begin{remark}
El byte $v$ en Ethereum (27 o 28, o 0/1 post-EIP-155) permite \textbf{recuperar la clave pública} a partir de la firma. Esto es útil porque:
\begin{itemize}
    \item No necesitas incluir la clave pública en cada transacción
    \item El contrato puede derivar la dirección del firmante
    \item Ahorra espacio en la blockchain
\end{itemize}
\end{remark}

%-------------------------------------------
\subsubsection{Recuperación de Clave Pública (ecrecover)}

\begin{definition}[ECDSA Public Key Recovery]
Dada una firma $(r, s)$ y el hash del mensaje $e$, se puede recuperar la clave pública $pk$ (con hasta 4 candidatos, reducidos a 2 con el recovery id):

\begin{enumerate}
    \item Calcular $R = (r, y_R)$ donde $y_R$ se deriva de $r$ (dos opciones)
    \item Calcular $pk = r^{-1}(s \cdot R - e \cdot G)$
\end{enumerate}
\end{definition}

\begin{example}[ecrecover en Solidity]
\begin{verbatim}
// SPDX-License-Identifier: MIT
pragma solidity ^0.8.0;

contract ECDSAExample {
    function recoverSigner(
        bytes32 messageHash,
        uint8 v,
        bytes32 r,
        bytes32 s
    ) public pure returns (address) {
        // ecrecover es un precompilado de la EVM
        address signer = ecrecover(messageHash, v, r, s);
        require(signer != address(0), "Invalid signature");
        return signer;
    }
    
    function verifySignature(
        bytes32 messageHash,
        bytes memory signature,
        address expectedSigner
    ) public pure returns (bool) {
        bytes32 r;
        bytes32 s;
        uint8 v;
        
        assembly {
            r := mload(add(signature, 32))
            s := mload(add(signature, 64))
            v := byte(0, mload(add(signature, 96)))
        }
        
        return recoverSigner(messageHash, v, r, s) == expectedSigner;
    }
}
\end{verbatim}
\end{example}

%-------------------------------------------
\subsubsection{ECDSA vs Otros Esquemas de Firma}

\begin{center}
\begin{tabular}{lcccc}
\toprule
\textbf{Propiedad} & \textbf{ECDSA} & \textbf{Schnorr} & \textbf{EdDSA} & \textbf{BLS} \\
\midrule
Tamaño de firma & 64-72 B & 64 B & 64 B & 48 B (G1) \\
Determinístico & Opcional & Sí & Sí & Sí \\
Maleabilidad & Sí & No & No & No \\
Agregación nativa & No & Sí (MuSig) & Limitada & Sí \\
Batch verification & No & Sí & Sí & Sí \\
Prueba de seguridad & ROM & ROM/Standard & Standard & ROM \\
Adopción & Alta & Media & Alta & Media \\
\bottomrule
\end{tabular}
\end{center}

\begin{keyidea}
\textbf{¿Por qué Bitcoin/Ethereum eligieron ECDSA sobre Schnorr?}

En 2008-2015, las firmas Schnorr estaban bajo patente (expirada en 2008, pero la incertidumbre legal persistía). ECDSA era el estándar ``seguro'' establecido.

\textbf{Taproot (2021)} añadió Schnorr a Bitcoin (BIP-340), ofreciendo:
\begin{itemize}
    \item Firmas más pequeñas
    \item Agregación de firmas (multisig más eficiente)
    \item Sin maleabilidad
\end{itemize}

Ethereum planea adoptar BLS para Ethereum 2.0 (agregación de firmas de validadores).
\end{keyidea}

%-------------------------------------------
\subsubsection{Verificación en Circuitos ZK}

\begin{mathnote}
Verificar una firma ECDSA dentro de un circuito ZK es costoso porque involucra:
\begin{enumerate}
    \item Multiplicación escalar en curva elíptica: $u_1 \cdot G$
    \item Otra multiplicación escalar: $u_2 \cdot pk$
    \item Suma de puntos: $u_1 G + u_2 \cdot pk$
    \item Inversión modular: $s^{-1}$
    \item Hash del mensaje (SHA-256 o Keccak)
\end{enumerate}

\textbf{Costo aproximado en R1CS:}
\begin{center}
\begin{tabular}{lc}
\toprule
\textbf{Operación} & \textbf{Constraints} \\
\midrule
Multiplicación escalar (secp256k1) & $\sim$250,000 \\
Inversión modular & $\sim$10,000 \\
SHA-256 & $\sim$25,000 \\
\midrule
\textbf{Total ECDSA verify} & $\sim$600,000 \\
\bottomrule
\end{tabular}
\end{center}
\end{mathnote}

%-------------------------------------------
\subsubsection{Aplicación en Banco ZK}

\begin{bancozk}
ECDSA es la interfaz principal entre usuarios de Ethereum y Banco ZK:

\textbf{1. Autenticación de usuarios}
\begin{itemize}
    \item Los usuarios firman transacciones de Banco ZK con sus claves Ethereum existentes
    \item No necesitan crear nuevas claves ni aprender nuevos procesos
    \item Compatibilidad total con MetaMask, Ledger, etc.
\end{itemize}

\textbf{2. Verificación de firmas}

\textbf{Opción A: Fuera del circuito (lo que hacemos)}
\begin{enumerate}
    \item El contrato L1 verifica la firma ECDSA (barato: $\sim$3,000 gas)
    \item Pasa la dirección del firmante como input público al circuito
    \item El circuito verifica que los nullifiers corresponden a esa dirección
\end{enumerate}

\textbf{Opción B: Dentro del circuito (costoso)}
\begin{itemize}
    \item Verificar ECDSA en el circuito: $\sim$600,000 constraints
    \item Solo útil si necesitas privacidad sobre QUIÉN firmó
\end{itemize}

\textbf{3. Flujo de depósito}
\begin{verbatim}
1. Usuario firma: sign(sk, "deposit 1 ETH to commitment C")
2. Contrato verifica: ecrecover(hash, v, r, s) == msg.sender
3. Contrato ejecuta: deposit(commitment)
4. El circuito ZK opera sobre el commitment, no sobre la firma
\end{verbatim}

\textbf{4. Flujo de retiro}
\begin{verbatim}
1. Usuario genera prueba ZK de que conoce el preimagen de un nullifier
2. Usuario firma el retiro: sign(sk, "withdraw to address X")
3. Contrato verifica la firma ECDSA
4. Contrato verifica la prueba ZK
5. Contrato envía fondos a X
\end{verbatim}

\textbf{5. Por qué no usamos solo EdDSA o Schnorr internamente}

Aunque son más eficientes en circuitos ZK, mantener ECDSA tiene ventajas:
\begin{center}
\begin{tabular}{lcc}
\toprule
\textbf{Aspecto} & \textbf{ECDSA} & \textbf{EdDSA} \\
\midrule
Wallets compatibles & Todas & Pocas \\
Hardware wallets & Sí & Limitado \\
UX de usuario & Conocida & Nueva \\
Gas en L1 & 3,000 (precompilado) & $\sim$30,000 (sin precompilado) \\
\bottomrule
\end{tabular}
\end{center}

\textbf{6. Consideración de seguridad}

Nuestros contratos SIEMPRE verifican:
\begin{verbatim}
require(s <= secp256k1_n / 2, "Malleable signature");
\end{verbatim}

Esto previene ataques de maleabilidad que podrían alterar el txid.

\textbf{7. Migración futura}

Cuando Account Abstraction (EIP-4337) esté más maduro:
\begin{itemize}
    \item Usuarios podrán usar cualquier esquema de firma
    \item Banco ZK soportará Schnorr, EdDSA, BLS
    \item La interfaz permanecerá igual para el usuario
\end{itemize}
\end{bancozk}

%-------------------------------------------
\subsubsection{Implementación de Referencia}

\begin{example}[ECDSA en Python (educacional)]
\begin{verbatim}
import hashlib
import secrets
from dataclasses import dataclass

# Parámetros secp256k1 (simplificados)
p = 0xFFFFFFFFFFFFFFFFFFFFFFFFFFFFFFFFFFFFFFFFFFFFFFFFFFFFFFFEFFFFFC2F
n = 0xFFFFFFFFFFFFFFFFFFFFFFFFFFFFFFFEBAAEDCE6AF48A03BBFD25E8CD0364141
Gx = 0x79BE667EF9DCBBAC55A06295CE870B07029BFCDB2DCE28D959F2815B16F81798
Gy = 0x483ADA7726A3C4655DA4FBFC0E1108A8FD17B448A68554199C47D08FFB10D4B8

@dataclass
class Point:
    x: int
    y: int

def mod_inverse(a, m):
    return pow(a, m - 2, m)

def point_add(P, Q):
    if P is None: return Q
    if Q is None: return P
    if P.x == Q.x and P.y != Q.y: return None
    
    if P.x == Q.x:  # Duplicación
        lam = (3 * P.x * P.x) * mod_inverse(2 * P.y, p) % p
    else:  # Suma
        lam = (Q.y - P.y) * mod_inverse(Q.x - P.x, p) % p
    
    x = (lam * lam - P.x - Q.x) % p
    y = (lam * (P.x - x) - P.y) % p
    return Point(x, y)

def scalar_mult(k, P):
    result = None
    addend = P
    while k:
        if k & 1:
            result = point_add(result, addend)
        addend = point_add(addend, addend)
        k >>= 1
    return result

G = Point(Gx, Gy)

def keygen():
    sk = secrets.randbelow(n - 1) + 1
    pk = scalar_mult(sk, G)
    return sk, pk

def sign(sk, message):
    e = int(hashlib.sha256(message).hexdigest(), 16) % n
    
    while True:
        k = secrets.randbelow(n - 1) + 1
        R = scalar_mult(k, G)
        r = R.x % n
        if r == 0: continue
        
        s = mod_inverse(k, n) * (e + r * sk) % n
        if s == 0: continue
        
        # Normalizar a low-S
        if s > n // 2:
            s = n - s
        
        return r, s

def verify(pk, message, r, s):
    if not (1 <= r < n and 1 <= s < n):
        return False
    
    e = int(hashlib.sha256(message).hexdigest(), 16) % n
    w = mod_inverse(s, n)
    u1 = (e * w) % n
    u2 = (r * w) % n
    
    R = point_add(scalar_mult(u1, G), scalar_mult(u2, pk))
    if R is None:
        return False
    
    return R.x % n == r

# Uso
sk, pk = keygen()
message = b"Hello, ECDSA!"
r, s = sign(sk, message)
assert verify(pk, message, r, s), "Verification failed!"
print("Signature verified successfully!")
\end{verbatim}
\end{example}

%-------------------------------------------
\subsubsection{Resumen de Buenas Prácticas}

\begin{keyidea}
\textbf{Checklist de seguridad ECDSA:}

\begin{enumerate}
    \item[$\square$] \textbf{NUNCA} reutilizar nonces. Usar RFC 6979.
    \item[$\square$] Verificar que $r, s \in [1, n-1]$ antes de procesar.
    \item[$\square$] Normalizar firmas a low-S ($s \leq n/2$).
    \item[$\square$] Usar librerías auditadas (libsecp256k1, etc.).
    \item[$\square$] Generar claves con CSPRNG de calidad.
    \item[$\square$] Validar claves públicas (verificar que están en la curva).
    \item[$\square$] Proteger claves privadas en memoria (zeroing después de uso).
    \item[$\square$] Considerar timing attacks en implementaciones.
\end{enumerate}
\end{keyidea}

%===========================================
% SUBSECCIÓN: SCHNORR - VERSIÓN EXPANDIDA
% Usa tus entornos: definition, example, keyidea, bancozk, 
%                   mathnote, intuicion, advertencia, remark
%===========================================

\subsection{Schnorr}

\begin{definition}[Firma Schnorr]
Esquema de firma con prueba de conocimiento del logaritmo discreto.
\end{definition}

Algoritmo:
\begin{enumerate}
    \item Elegir nonce $k$, calcular $R = kG$
    \item Challenge: $e = H(R \| pk \| m)$
    \item Response: $s = k + e \cdot sk$
    \item Firma: $(R, s)$ o $(e, s)$
    \item Verificar: $sG = R + e \cdot pk$
\end{enumerate}

Ventajas:
\begin{itemize}
    \item Prueba de seguridad más simple
    \item Linealidad permite agregación: $\sum s_i G = \sum R_i + e \sum pk_i$
    \item MuSig: firmas multisig que parecen firmas normales
\end{itemize}

%-------------------------------------------
\subsubsection{Contexto Histórico}

\begin{intuicion}
Las firmas Schnorr son elegantes, eficientes y matemáticamente más simples que ECDSA. Fueron inventadas en 1989 por Claus-Peter Schnorr, un criptógrafo alemán. Entonces, ¿por qué Bitcoin no las usó desde el principio?

La respuesta: \textbf{patentes}.
\end{intuicion}

\begin{remark}
\textbf{La saga de las patentes Schnorr:}

\begin{center}
\begin{tabular}{ll}
\toprule
\textbf{Año} & \textbf{Evento} \\
\midrule
1989 & Claus Schnorr publica el esquema \\
1990 & Patente US 4,995,082 otorgada \\
1991 & NIST estandariza DSA (diseñado para evitar la patente) \\
2008 & Patente expira \\
2008 & Satoshi elige ECDSA (incertidumbre legal residual) \\
2018 & BIP-340 propuesto para Bitcoin \\
2021 & Taproot activa Schnorr en Bitcoin \\
\bottomrule
\end{tabular}
\end{center}

DSA (y por extensión ECDSA) fue diseñado específicamente por NIST para evitar la patente de Schnorr. Es por eso que ECDSA tiene una estructura más complicada.
\end{remark}

\begin{keyidea}
Ahora que las patentes expiraron, Schnorr está siendo adoptado ampliamente:
\begin{itemize}
    \item \textbf{Bitcoin}: Taproot (BIP-340, BIP-341) desde noviembre 2021
    \item \textbf{EdDSA}: Variante de Schnorr sobre curvas de Edwards (ed25519)
    \item \textbf{Monero}: Usa firmas Schnorr para ring signatures
    \item \textbf{Polkadot}: Sr25519 (Schnorr sobre Ristretto)
\end{itemize}
\end{keyidea}

%-------------------------------------------
\subsubsection{El Protocolo Schnorr: Identificación Interactiva}

\begin{intuicion}
Antes de entender las firmas, es útil ver el \textbf{protocolo de identificación} de Schnorr. Es un protocolo interactivo donde un prover demuestra que conoce un secreto sin revelarlo.
\end{intuicion}

\begin{definition}[Protocolo de Identificación Schnorr]
\textbf{Setup:} Prover tiene $sk$, Verifier conoce $pk = sk \cdot G$.

\textbf{Protocolo:}
\begin{enumerate}
    \item \textbf{Commitment:} Prover elige $k \xleftarrow{\$} \Z_n$, envía $R = k \cdot G$
    \item \textbf{Challenge:} Verifier elige $e \xleftarrow{\$} \Z_n$, envía $e$
    \item \textbf{Response:} Prover calcula $s = k + e \cdot sk$, envía $s$
    \item \textbf{Verificación:} Verifier acepta si $s \cdot G = R + e \cdot pk$
\end{enumerate}
\end{definition}

\begin{center}
\begin{tikzpicture}[
    entity/.style={draw, rectangle, minimum width=2cm, minimum height=1cm},
    arrow/.style={->, thick, >=stealth}
]
    % Entidades
    \node[entity] (prover) at (0, 0) {Prover};
    \node[entity] (verifier) at (8, 0) {Verifier};
    
    % Conocimiento
    \node[above=0.3cm of prover, font=\small] {Conoce: $sk$};
    \node[above=0.3cm of verifier, font=\small] {Conoce: $pk = sk \cdot G$};
    
    % Mensajes
    \draw[arrow] (1.5, -0.5) -- node[above] {$R = kG$} (6.5, -0.5);
    \draw[arrow] (6.5, -1.2) -- node[above] {$e$} (1.5, -1.2);
    \draw[arrow] (1.5, -1.9) -- node[above] {$s = k + e \cdot sk$} (6.5, -1.9);
    
    % Verificación
    \node at (8, -2.5) [font=\small] {Verifica: $sG \stackrel{?}{=} R + e \cdot pk$};
\end{tikzpicture}
\end{center}

\begin{mathnote}
\textbf{¿Por qué funciona?}
\begin{align}
sG &= (k + e \cdot sk) \cdot G \\
   &= k \cdot G + e \cdot sk \cdot G \\
   &= R + e \cdot pk \quad \checkmark
\end{align}

\textbf{¿Por qué es seguro?}

Un impostor que no conoce $sk$ tendría que:
\begin{itemize}
    \item Adivinar $e$ antes de enviar $R$ (probabilidad $1/n \approx 2^{-256}$), o
    \item Calcular $sk$ a partir de $pk$ (problema ECDLP)
\end{itemize}
\end{mathnote}

%-------------------------------------------
\subsubsection{De Identificación a Firma: Transformación Fiat-Shamir}

\begin{definition}[Transformación Fiat-Shamir]
Para convertir un protocolo interactivo en una firma no interactiva, reemplazamos el challenge aleatorio del verifier por el hash de la transcripción:
\[ e = H(R \,\|\, pk \,\|\, m) \]

donde $m$ es el mensaje a firmar.
\end{definition}

\begin{intuicion}
La idea es que el hash actúa como un ``verifier virtual'' que genera un challenge impredecible basado en el commitment $R$. El prover no puede ``hacer trampa'' eligiendo $R$ después de conocer $e$, porque cambiar $R$ cambiaría $e$.
\end{intuicion}

\begin{keyidea}
La firma Schnorr es esencialmente una \textbf{prueba de conocimiento cero no interactiva (NIZK)} de que el firmante conoce el logaritmo discreto de $pk$.

Esto conecta directamente con los SNARKs: las firmas Schnorr son el caso más simple de ``prueba de conocimiento''.
\end{keyidea}

%-------------------------------------------
\subsubsection{El Algoritmo de Firma Schnorr}

\begin{definition}[Firma Schnorr]
\textbf{Parámetros:} Curva $E$, generador $G$, orden $n$, hash $H$.

\textbf{KeyGen:}
\begin{enumerate}
    \item $sk \xleftarrow{\$} [1, n-1]$
    \item $pk = sk \cdot G$
\end{enumerate}

\textbf{Sign$(sk, m)$:}
\begin{enumerate}
    \item Elegir nonce: $k \xleftarrow{\$} [1, n-1]$ (o determinístico vía RFC 6979)
    \item Calcular commitment: $R = k \cdot G$
    \item Calcular challenge: $e = H(R \,\|\, pk \,\|\, m) \mod n$
    \item Calcular response: $s = k + e \cdot sk \mod n$
    \item Retornar $\sigma = (R, s)$ o $\sigma = (e, s)$
\end{enumerate}

\textbf{Verify$(pk, m, \sigma)$:}
\begin{enumerate}
    \item Parsear $\sigma = (R, s)$
    \item Calcular $e = H(R \,\|\, pk \,\|\, m) \mod n$
    \item Verificar: $s \cdot G \stackrel{?}{=} R + e \cdot pk$
\end{enumerate}
\end{definition}

\begin{remark}
Hay dos formatos comunes de firma Schnorr:
\begin{itemize}
    \item $(R, s)$: 64 bytes (32 para $x_R$, 32 para $s$). Usado en BIP-340.
    \item $(e, s)$: 64 bytes (32 para $e$, 32 para $s$). Más compacto pero requiere recalcular $R$ en verificación.
\end{itemize}

Bitcoin usa $(R, s)$ porque simplifica la verificación por lotes (batch verification).
\end{remark}

%-------------------------------------------
\subsubsection{Comparación con ECDSA}

\begin{center}
\begin{tabular}{lcc}
\toprule
\textbf{Aspecto} & \textbf{ECDSA} & \textbf{Schnorr} \\
\midrule
Firma & $s = k^{-1}(e + r \cdot sk)$ & $s = k + e \cdot sk$ \\
Verificación & $R' = s^{-1}(eG + r \cdot pk)$ & $sG = R + e \cdot pk$ \\
Inversiones & 2 (en firma y verificación) & 0 \\
Prueba de seguridad & Modelo del oráculo aleatorio & Más simple, más fuerte \\
Linealidad & No & Sí \\
Agregación nativa & No & Sí \\
Maleabilidad & Sí ($(r, s)$ y $(r, n-s)$) & No \\
Batch verification & No eficiente & Sí, nativo \\
\bottomrule
\end{tabular}
\end{center}

\begin{keyidea}
La diferencia clave es la \textbf{linealidad}:
\begin{itemize}
    \item ECDSA: $s = k^{-1}(e + r \cdot sk)$ — la inversión $k^{-1}$ rompe la linealidad
    \item Schnorr: $s = k + e \cdot sk$ — completamente lineal
\end{itemize}

Esta linealidad permite sumar firmas, lo cual es imposible con ECDSA.
\end{keyidea}

%-------------------------------------------
\subsubsection{Ejemplo Numérico}

\begin{example}[Firma Schnorr paso a paso]
\textbf{Parámetros simplificados:}
\begin{itemize}
    \item Campo: $\mathbb{F}_{97}$
    \item Curva: $y^2 = x^3 + 7$ (como secp256k1 pero pequeña)
    \item Generador: $G = (5, 32)$
    \item Orden: $n = 79$ (asumamos que es primo para el ejemplo)
\end{itemize}

\textbf{1. Generación de claves:}
\begin{align}
sk &= 23 \quad \text{(secreto)} \\
pk &= 23 \cdot G = 23 \cdot (5, 32) = (41, 67) \quad \text{(calculado)}
\end{align}

\textbf{2. Firma del mensaje $m$ donde $H(R \| pk \| m) = 17$:}
\begin{align}
k &= 47 \quad \text{(nonce aleatorio)} \\
R &= 47 \cdot G = (62, 15) \\
e &= H(R \| pk \| m) = 17 \quad \text{(simplificado)} \\
s &= k + e \cdot sk = 47 + 17 \cdot 23 = 47 + 391 = 438 \equiv 43 \pmod{79}
\end{align}

\textbf{Firma:} $\sigma = (R, s) = ((62, 15), 43)$

\textbf{3. Verificación:}

Calcular ambos lados de $sG = R + e \cdot pk$:

\textbf{Lado izquierdo:}
\[ sG = 43 \cdot G = 43 \cdot (5, 32) = (?, ?) \]

\textbf{Lado derecho:}
\[ R + e \cdot pk = (62, 15) + 17 \cdot (41, 67) = (62, 15) + (?, ?) = (?, ?) \]

Si ambos son iguales, la firma es válida. ✓
\end{example}

%-------------------------------------------
\subsubsection{Seguridad: El Nonce Sigue Siendo Crítico}

\begin{advertencia}
Al igual que ECDSA, \textbf{reutilizar el nonce en Schnorr expone la clave privada}.

Si dos firmas $(R, s_1)$ y $(R, s_2)$ usan el mismo $k$ (mismo $R$):
\begin{align}
s_1 &= k + e_1 \cdot sk \\
s_2 &= k + e_2 \cdot sk
\end{align}

Restando:
\[ s_1 - s_2 = (e_1 - e_2) \cdot sk \]

Despejando:
\[ sk = \frac{s_1 - s_2}{e_1 - e_2} \mod n \]

¡Incluso más simple que el ataque a ECDSA!
\end{advertencia}

\begin{keyidea}
\textbf{Solución:} Usar nonces determinísticos.

BIP-340 especifica:
\[ k = H_{\text{aux}}(sk \,\|\, \text{aux\_rand}) \oplus H_{\text{nonce}}(sk \,\|\, pk \,\|\, m) \]

Esto combina:
\begin{itemize}
    \item Determinismo (derivado de $sk$, $pk$, $m$)
    \item Aleatoriedad opcional (\texttt{aux\_rand}) para protección contra side-channels
\end{itemize}
\end{keyidea}

%-------------------------------------------
\subsubsection{La Propiedad Mágica: Linealidad}

\begin{keyidea}
La ecuación de verificación Schnorr es \textbf{lineal}:
\[ s \cdot G = R + e \cdot pk \]

Esto significa que podemos \textbf{sumar} firmas:
\[ (s_1 + s_2) \cdot G = (R_1 + R_2) + e \cdot (pk_1 + pk_2) \]

¡Una firma agregada que verifica contra una clave pública agregada!
\end{keyidea}

\begin{example}[Agregación simple (ingenua, insegura)]
Alice y Bob quieren firmar conjuntamente:
\begin{itemize}
    \item Alice: $sk_A$, $pk_A = sk_A \cdot G$
    \item Bob: $sk_B$, $pk_B = sk_B \cdot G$
\end{itemize}

\textbf{Clave pública conjunta:}
\[ pk_{AB} = pk_A + pk_B = (sk_A + sk_B) \cdot G \]

\textbf{Firma conjunta:}
\begin{enumerate}
    \item Alice elige $k_A$, envía $R_A = k_A G$
    \item Bob elige $k_B$, envía $R_B = k_B G$
    \item $R = R_A + R_B$
    \item $e = H(R \| pk_{AB} \| m)$
    \item Alice calcula $s_A = k_A + e \cdot sk_A$
    \item Bob calcula $s_B = k_B + e \cdot sk_B$
    \item $s = s_A + s_B$
\end{enumerate}

\textbf{Firma final:} $(R, s)$

\textbf{Verificación:}
\[ sG = (s_A + s_B)G = s_A G + s_B G = (R_A + e \cdot pk_A) + (R_B + e \cdot pk_B) = R + e \cdot pk_{AB} \quad \checkmark \]
\end{example}

\begin{advertencia}
\textbf{¡La agregación ingenua es INSEGURA!}

El ataque ``Rogue Key'': Un atacante puede elegir su ``clave pública'' como:
\[ pk_{\text{attack}} = pk_{\text{attack}}' - pk_{\text{victim}} \]

Entonces:
\[ pk_{\text{joint}} = pk_{\text{victim}} + pk_{\text{attack}} = pk_{\text{attack}}' \]

El atacante controla completamente la clave conjunta y puede firmar solo.

Solución: Protocolos seguros como MuSig.
\end{advertencia}

%-------------------------------------------
\subsubsection{MuSig: Agregación Segura de Firmas}

\begin{definition}[MuSig]
MuSig (Multi-Signature Schnorr) es un protocolo seguro para crear firmas $n$-de-$n$ que son indistinguibles de firmas individuales.

\textbf{Idea clave:} Cada clave pública se multiplica por un coeficiente derivado de todas las claves, previniendo ataques de clave rogue.
\end{definition}

\begin{example}[MuSig simplificado (2-de-2)]
\textbf{Setup:}
\begin{itemize}
    \item Alice: $sk_A$, $pk_A$
    \item Bob: $sk_B$, $pk_B$
\end{itemize}

\textbf{Paso 1: Calcular coeficientes}
\begin{align}
L &= H(pk_A \| pk_B) \quad \text{(lista ordenada de claves)} \\
a_A &= H(L \| pk_A) \\
a_B &= H(L \| pk_B)
\end{align}

\textbf{Paso 2: Clave pública agregada}
\[ \tilde{pk} = a_A \cdot pk_A + a_B \cdot pk_B \]

\textbf{Paso 3: Firma (3 rondas)}
\begin{enumerate}
    \item Cada parte elige nonce $k_i$, envía commitment $H(R_i)$
    \item Revelar $R_i$, calcular $R = R_A + R_B$
    \item Calcular $e = H(R \| \tilde{pk} \| m)$
    \item Cada parte calcula $s_i = k_i + e \cdot a_i \cdot sk_i$
    \item $s = s_A + s_B$
\end{enumerate}

\textbf{Firma final:} $(R, s)$ — ¡indistinguible de firma individual!
\end{example}

\begin{remark}
\textbf{Evolución de MuSig:}
\begin{itemize}
    \item \textbf{MuSig} (2018): 3 rondas de comunicación
    \item \textbf{MuSig2} (2020): 2 rondas, más eficiente
    \item \textbf{MuSig-DN} (2020): Nonces determinísticos
    \item \textbf{FROST} (2020): Threshold signatures ($t$-de-$n$)
\end{itemize}
\end{remark}

%-------------------------------------------
\subsubsection{MuSig2: Reduciendo Rondas}

\begin{definition}[MuSig2]
MuSig2 reduce las rondas de comunicación de 3 a 2 usando múltiples nonces por firmante.

\textbf{Idea:} Cada firmante envía 2 nonces $(R_{i,1}, R_{i,2})$ en la primera ronda. El nonce efectivo se combina con un coeficiente derivado del mensaje.
\end{definition}

\begin{mathnote}
\textbf{Por qué funcionan 2 rondas:}

En MuSig original, la primera ronda es solo para commitments de nonces. Esto previene que un atacante vea los nonces de otros y ajuste el suyo.

MuSig2 usa una técnica diferente:
\begin{enumerate}
    \item Enviar $(R_{i,1}, R_{i,2})$ directamente
    \item El nonce efectivo es $R_i = R_{i,1} + b \cdot R_{i,2}$ donde $b = H(\text{todos los } R_{j,k}, m)$
    \item El atacante no puede predecir $b$ al elegir sus nonces
\end{enumerate}
\end{mathnote}

%-------------------------------------------
\subsubsection{Batch Verification}

\begin{keyidea}
Otra ventaja de la linealidad: verificar $n$ firmas Schnorr es más rápido que verificarlas individualmente.

\textbf{Verificación individual:} $n$ verificaciones, cada una con 2 multiplicaciones escalares = $2n$ mult. escalares.

\textbf{Batch verification:} Combinar en una sola ecuación usando pesos aleatorios.
\end{keyidea}

\begin{definition}[Batch Verification]
Para verificar firmas $(R_1, s_1), \ldots, (R_n, s_n)$ sobre mensajes $m_1, \ldots, m_n$ con claves $pk_1, \ldots, pk_n$:

\begin{enumerate}
    \item Elegir pesos aleatorios $\alpha_1, \ldots, \alpha_n \xleftarrow{\$} \Z_n$
    \item Calcular challenges: $e_i = H(R_i \| pk_i \| m_i)$
    \item Verificar:
    \[ \left(\sum_{i=1}^n \alpha_i s_i\right) \cdot G = \sum_{i=1}^n \alpha_i R_i + \sum_{i=1}^n \alpha_i e_i \cdot pk_i \]
\end{enumerate}

Si la ecuación se cumple, todas las firmas son válidas (con alta probabilidad).
\end{definition}

\begin{mathnote}
\textbf{Análisis de eficiencia:}

\begin{center}
\begin{tabular}{lcc}
\toprule
& \textbf{Individual} & \textbf{Batch} \\
\midrule
Mult. escalares & $2n$ & $n + 1$ \\
Sumas de puntos & $n$ & $2n$ \\
\bottomrule
\end{tabular}
\end{center}

Las multiplicaciones escalares son mucho más costosas que las sumas. Para $n$ grande, batch verification es $\sim$2x más rápido.

Los pesos aleatorios $\alpha_i$ previenen que un atacante cree firmas inválidas que se ``cancelen'' en la suma.
\end{mathnote}

%-------------------------------------------
\subsubsection{BIP-340: Schnorr en Bitcoin}

\begin{definition}[BIP-340]
Bitcoin Improvement Proposal 340 especifica Schnorr para Bitcoin (activado con Taproot en noviembre 2021).

\textbf{Características:}
\begin{itemize}
    \item Curva: secp256k1 (misma que ECDSA)
    \item Tamaño de firma: 64 bytes (vs. 71-72 para ECDSA DER)
    \item Clave pública: 32 bytes (solo coordenada $x$)
    \item Hash: SHA256 con tags para domain separation
\end{itemize}
\end{definition}

\begin{mathnote}
\textbf{Innovaciones de BIP-340:}

\textbf{1. Solo coordenada $x$ para claves y $R$:}

En lugar de transmitir $(x, y)$, BIP-340 usa solo $x$ y asume que $y$ es par. Si $y$ es impar, se niega la clave.

Ahorro: 32 bytes por clave pública.

\textbf{2. Tagged hashes:}
\[ H_{\text{tag}}(m) = \text{SHA256}(\text{SHA256}(\text{tag}) \| \text{SHA256}(\text{tag}) \| m) \]

Tags usados:
\begin{itemize}
    \item ``BIP0340/aux'' para mezclar \texttt{aux\_rand}
    \item ``BIP0340/nonce'' para generar nonce
    \item ``BIP0340/challenge'' para el challenge $e$
\end{itemize}

Esto previene colisiones entre diferentes usos del hash.
\end{mathnote}

\begin{example}[Firma BIP-340]
\begin{verbatim}
# Inputs
secret_key: 32 bytes
message: arbitrary length
aux_rand: 32 bytes (optional randomness)

# Key generation
d = int(secret_key)
P = d * G
if P.y is odd:
    d = n - d  # Negate to make y even
pk = P.x  # Just x-coordinate, 32 bytes

# Nonce generation
t = xor(secret_key, tagged_hash("BIP0340/aux", aux_rand))
k' = int(tagged_hash("BIP0340/nonce", t || pk || message))
k' = k' mod n
R = k' * G
if R.y is odd:
    k = n - k'  # Negate to make y even
else:
    k = k'

# Challenge
e = int(tagged_hash("BIP0340/challenge", R.x || pk || message)) mod n

# Signature
s = (k + e * d) mod n
sig = R.x || s  # 64 bytes total
\end{verbatim}
\end{example}

%-------------------------------------------
\subsubsection{Taproot: Schnorr + MAST}

\begin{intuicion}
Taproot combina Schnorr con Merkelized Abstract Syntax Trees (MAST) para crear condiciones de gasto flexibles y privadas.

\textbf{Idea:} Una dirección Taproot puede gastarse de dos formas:
\begin{enumerate}
    \item \textbf{Key path:} Firma Schnorr simple (parece una transacción normal)
    \item \textbf{Script path:} Revelar un script oculto en un árbol Merkle
\end{enumerate}

Si todos los participantes cooperan, usan el key path y nadie sabe que había otras condiciones posibles.
\end{intuicion}

\begin{definition}[Dirección Taproot]
Una dirección Taproot codifica:
\[ P = Q + H(Q \| \text{script\_root}) \cdot G \]

donde:
\begin{itemize}
    \item $Q$ es la clave pública interna (para key path)
    \item $\text{script\_root}$ es la raíz del árbol Merkle de scripts
\end{itemize}

Para gastar por key path, se firma con $q + H(Q \| \text{script\_root})$ donde $q$ es la clave privada de $Q$.
\end{definition}

%-------------------------------------------
\subsubsection{Schnorr vs EdDSA}

\begin{remark}
EdDSA (Ed25519, Ed448) es técnicamente una variante de Schnorr, pero con diferencias importantes:

\begin{center}
\begin{tabular}{lcc}
\toprule
\textbf{Aspecto} & \textbf{Schnorr (BIP-340)} & \textbf{EdDSA (Ed25519)} \\
\midrule
Curva & secp256k1 & Curve25519 \\
Forma de curva & Weierstrass & Edwards \\
Nonce & Determinístico + aux & Determinístico puro \\
Challenge & $H(R \| pk \| m)$ & $H(R \| pk \| m)$ \\
Firma & $(R.x, s)$ & $(R, s)$ \\
Clave pública & 32 bytes (solo $x$) & 32 bytes (comprimido) \\
Cofactor & 1 & 8 \\
\bottomrule
\end{tabular}
\end{center}
\end{remark}

\begin{keyidea}
EdDSA es más popular fuera de Bitcoin porque:
\begin{enumerate}
    \item Curve25519 es más rápida que secp256k1
    \item Implementaciones de referencia bien auditadas (libsodium)
    \item Sin historial de patentes
    \item Usado en SSH, TLS 1.3, Signal, etc.
\end{enumerate}

Bitcoin usa Schnorr sobre secp256k1 para compatibilidad con infraestructura existente.
\end{keyidea}

%-------------------------------------------
\subsubsection{Verificación en Circuitos ZK}

\begin{mathnote}
Verificar Schnorr en un circuito ZK es significativamente más eficiente que ECDSA:

\begin{center}
\begin{tabular}{lcc}
\toprule
\textbf{Operación} & \textbf{ECDSA} & \textbf{Schnorr} \\
\midrule
Inversiones modulares & 2 & 0 \\
Multiplicaciones escalares & 2 & 2 \\
Estructura & No lineal & Lineal \\
\midrule
\textbf{Constraints totales} & $\sim$600,000 & $\sim$400,000 \\
\bottomrule
\end{tabular}
\end{center}

La ausencia de inversiones y la estructura lineal reducen los constraints en $\sim$30\%.
\end{mathnote}

%-------------------------------------------
\subsubsection{Adapter Signatures (Firmas Adaptadas)}

\begin{definition}[Adapter Signature]
Una firma adaptada es una firma Schnorr ``incompleta'' que se completa al revelar un secreto.

\textbf{Setup:}
\begin{itemize}
    \item Alice quiere crear una firma que Bob solo puede completar si revela un secreto $t$
    \item Alice conoce $T = t \cdot G$ (el ``adaptor point'')
\end{itemize}

\textbf{Protocolo:}
\begin{enumerate}
    \item Alice crea firma adaptada: $\tilde{\sigma} = (R, \tilde{s})$ donde $R$ incluye $T$
    \item Bob puede verificar que $\tilde{\sigma}$ se convertirá en firma válida con $t$
    \item Cuando Alice publica la firma completa $\sigma = (R, s)$
    \item Bob calcula $t = s - \tilde{s}$
\end{enumerate}
\end{definition}

\begin{intuicion}
Las adapter signatures permiten \textbf{atomic swaps sin HTLCs}:
\begin{enumerate}
    \item Alice crea firma adaptada para enviar BTC a Bob
    \item Bob crea firma adaptada para enviar ETH a Alice, usando el mismo adaptor $T$
    \item Cuando Alice reclama su ETH (revelando $t$), Bob puede extraer $t$ y reclamar su BTC
\end{enumerate}

Todo ocurre sin timelocks ni scripts complejos.
\end{intuicion}

%-------------------------------------------
\subsubsection{Aplicación en Banco ZK}

\begin{bancozk}
Las firmas Schnorr ofrecen varias ventajas para Banco ZK:

\textbf{1. Eficiencia en circuitos}

Verificar Schnorr en un SNARK cuesta $\sim$30\% menos que ECDSA:
\begin{center}
\begin{tabular}{lcc}
\toprule
\textbf{Esquema} & \textbf{Constraints} & \textbf{Tiempo de prueba} \\
\midrule
ECDSA & $\sim$600,000 & $\sim$10s \\
Schnorr & $\sim$400,000 & $\sim$7s \\
EdDSA (sobre curva embebida) & $\sim$50,000 & $\sim$1s \\
\bottomrule
\end{tabular}
\end{center}

\textbf{2. Agregación de firmas para batch deposits}

Múltiples usuarios pueden agregar sus firmas de depósito:
\begin{verbatim}
// 100 usuarios depositan
// Sin agregación: 100 verificaciones on-chain
// Con MuSig: 1 verificación on-chain

aggregate_pubkey = sum(coeff_i * pk_i)
aggregate_signature = (sum(R_i), sum(s_i))

// El contrato verifica UNA firma
verify(aggregate_pubkey, message, aggregate_signature)
\end{verbatim}

Ahorro: $100 \times 3000 = 300,000$ gas $\to$ $3,000$ gas

\textbf{3. Threshold signatures para governance}

El protocolo puede requerir $t$-de-$n$ firmas para cambios críticos:
\begin{itemize}
    \item Actualización de verificadores: 3-de-5 multisig
    \item Emergency pause: 2-de-3 multisig
    \item Treasury movements: 4-de-7 multisig
\end{itemize}

Con FROST (Schnorr threshold), estas parecen firmas normales on-chain.

\textbf{4. Por qué aún soportamos ECDSA}

A pesar de las ventajas de Schnorr:
\begin{itemize}
    \item No hay precompilado de Schnorr en EVM
    \item Wallets Ethereum usan ECDSA nativamente
    \item Cambiar requiere nuevo estándar de firma (EIP)
\end{itemize}

\textbf{Estrategia híbrida:}
\begin{center}
\begin{tabular}{ll}
\toprule
\textbf{Uso} & \textbf{Esquema} \\
\midrule
Firmas de usuario (L1) & ECDSA (compatibilidad) \\
Firmas internas (L2) & Schnorr/EdDSA \\
Multisig de governance & MuSig2 \\
Dentro de circuitos ZK & EdDSA sobre Jubjub \\
\bottomrule
\end{tabular}
\end{center}

\textbf{5. Adapter signatures para puentes}

Para puentes trustless entre Banco ZK y otras cadenas:
\begin{enumerate}
    \item Usuario crea adapter signature para lock en Banco ZK
    \item Relayer crea adapter signature para mint en destino
    \item Al completar una, la otra se puede completar atómicamente
\end{enumerate}

Esto elimina la necesidad de confiar en el relayer.

\textbf{6. Implementación actual}
\begin{verbatim}
// Usamos la librería schnorr-signatures de noble
import { schnorr } from '@noble/curves/secp256k1';

// Para firmas de governance (off-chain, luego verificadas)
const signature = schnorr.sign(messageHash, privateKey);
const isValid = schnorr.verify(signature, messageHash, publicKey);

// La verificación on-chain usa un contrato custom
// (no hay precompilado)
contract SchnorrVerifier {
    function verify(
        bytes32 px,
        bytes32 rx,
        bytes32 s,
        bytes32 e
    ) public pure returns (bool) {
        // Implementación en Solidity...
        // ~30,000 gas (vs 3,000 para ECDSA precompilado)
    }
}
\end{verbatim}
\end{bancozk}

%-------------------------------------------
\subsubsection{Resumen: ¿Por Qué Schnorr es Mejor?}

\begin{keyidea}
\textbf{Ventajas de Schnorr sobre ECDSA:}

\begin{enumerate}
    \item \textbf{Simplicidad matemática}: Sin inversiones, estructura lineal
    \item \textbf{Prueba de seguridad}: Reducción más directa al problema DL
    \item \textbf{Agregación}: Múltiples firmas $\to$ una firma
    \item \textbf{Batch verification}: Verificar $n$ firmas casi tan rápido como 1
    \item \textbf{Sin maleabilidad}: No hay firmas equivalentes
    \item \textbf{Firmas adaptadas}: Atomic swaps elegantes
    \item \textbf{Privacidad}: Multisig parece firma normal
    \item \textbf{Eficiencia en ZK}: Menos constraints
\end{enumerate}

\textbf{Desventajas:}
\begin{enumerate}
    \item Adopción más reciente (menos battle-tested)
    \item Sin precompilado EVM
    \item El nonce sigue siendo crítico (mismo riesgo que ECDSA)
\end{enumerate}

\textbf{Veredicto:} Para nuevos sistemas, Schnorr es generalmente preferible. Para sistemas existentes con ECDSA, la migración tiene costos.
\end{keyidea}

%-------------------------------------------
\subsubsection{Implementación de Referencia}

\begin{example}[Schnorr en Python (educacional)]
\begin{verbatim}
import hashlib
import secrets

# Parámetros (simplificados, usar librerías reales en producción)
p = 0xFFFFFFFFFFFFFFFFFFFFFFFFFFFFFFFFFFFFFFFFFFFFFFFFFFFFFFFEFFFFFC2F
n = 0xFFFFFFFFFFFFFFFFFFFFFFFFFFFFFFFEBAAEDCE6AF48A03BBFD25E8CD0364141
Gx = 0x79BE667EF9DCBBAC55A06295CE870B07029BFCDB2DCE28D959F2815B16F81798
Gy = 0x483ADA7726A3C4655DA4FBFC0E1108A8FD17B448A68554199C47D08FFB10D4B8

def tagged_hash(tag, data):
    tag_hash = hashlib.sha256(tag.encode()).digest()
    return hashlib.sha256(tag_hash + tag_hash + data).digest()

def schnorr_sign(sk, message):
    # Asegurar que pk tiene y par
    pk_point = scalar_mult(sk, G)
    if pk_point.y % 2 != 0:
        sk = n - sk
    pk = pk_point.x.to_bytes(32, 'big')
    
    # Generar nonce determinístico
    k = int.from_bytes(
        tagged_hash("BIP0340/nonce", sk.to_bytes(32, 'big') + pk + message),
        'big'
    ) % n
    
    R = scalar_mult(k, G)
    if R.y % 2 != 0:
        k = n - k
    
    # Challenge
    e = int.from_bytes(
        tagged_hash("BIP0340/challenge", R.x.to_bytes(32, 'big') + pk + message),
        'big'
    ) % n
    
    # Response
    s = (k + e * sk) % n
    
    return R.x.to_bytes(32, 'big') + s.to_bytes(32, 'big')

def schnorr_verify(pk_bytes, message, sig):
    rx = int.from_bytes(sig[:32], 'big')
    s = int.from_bytes(sig[32:], 'big')
    
    # Reconstruir R (asumiendo y par)
    R = lift_x(rx)
    if R is None:
        return False
    
    # Challenge
    e = int.from_bytes(
        tagged_hash("BIP0340/challenge", sig[:32] + pk_bytes + message),
        'big'
    ) % n
    
    # Verificar: sG = R + e*pk
    pk_point = lift_x(int.from_bytes(pk_bytes, 'big'))
    lhs = scalar_mult(s, G)
    rhs = point_add(R, scalar_mult(e, pk_point))
    
    return lhs.x == rhs.x and lhs.y == rhs.y
\end{verbatim}
\end{example}

%===========================================
% SUBSECCIONES: EdDSA (Ed25519) y BLS - VERSIÓN EXPANDIDA
% Usa tus entornos: definition, example, keyidea, bancozk, 
%                   mathnote, intuicion, advertencia, remark
%===========================================

\subsection{EdDSA (Ed25519)}

\begin{definition}[EdDSA]
Schnorr sobre curvas de Edwards con nonce determinístico.
\end{definition}

Características:
\begin{itemize}
    \item Curva Edwards: $-x^2 + y^2 = 1 + dx^2y^2$
    \item Nonce: $k = H(h_{sk} \| m)$ (determinístico, elimina riesgo)
    \item Ed25519: curva Curve25519, hash SHA-512
    \item Muy rápido, usado en: Signal, SSH, Solana
\end{itemize}

%-------------------------------------------
\subsubsection{Contexto y Motivación}

\begin{intuicion}
EdDSA (Edwards-curve Digital Signature Algorithm) fue diseñado por Daniel J. Bernstein y sus colaboradores con un objetivo claro: crear un esquema de firmas que sea \textbf{imposible de implementar mal}.

Los problemas históricos con ECDSA (nonces repetidos, side-channels, maleabilidad) motivaron un rediseño desde cero con seguridad como prioridad absoluta.
\end{intuicion}

\begin{remark}
\textbf{El equipo detrás de EdDSA:}
\begin{itemize}
    \item Daniel J. Bernstein (djb) - Creador de Curve25519, qmail, NaCl
    \item Niels Duif
    \item Tanja Lange
    \item Peter Schwabe
    \item Bo-Yin Yang
\end{itemize}

El paper ``High-speed high-security signatures'' (2011) introdujo Ed25519, que rápidamente se convirtió en el estándar de facto para firmas modernas.
\end{remark}

\begin{center}
\begin{tabular}{ll}
\toprule
\textbf{Año} & \textbf{Adopción} \\
\midrule
2011 & Publicación de Ed25519 \\
2013 & OpenSSH añade soporte \\
2014 & GnuPG añade soporte \\
2016 & Signal Protocol adopta Ed25519 \\
2017 & TLS 1.3 incluye Ed25519 \\
2018 & Solana usa Ed25519 nativamente \\
2020 & FIDO2/WebAuthn usa Ed25519 \\
\bottomrule
\end{tabular}
\end{center}

%-------------------------------------------
\subsubsection{Curvas de Edwards}

\begin{definition}[Curva de Edwards]
Una curva de Edwards sobre un campo $\mathbb{F}$ tiene la forma:
\[ E: x^2 + y^2 = 1 + d \cdot x^2 y^2 \]

donde $d \in \mathbb{F}$ es un no-cuadrado (para que la curva sea no singular).

La forma \textbf{twisted Edwards} generaliza esto:
\[ E: ax^2 + y^2 = 1 + d \cdot x^2 y^2 \]
\end{definition}

\begin{keyidea}
\textbf{¿Por qué curvas de Edwards sobre Weierstrass?}

Las curvas de Edwards tienen propiedades algebraicas superiores:
\begin{enumerate}
    \item \textbf{Ley de grupo unificada}: Una sola fórmula para suma y duplicación
    \item \textbf{Fórmulas completas}: Funcionan para TODOS los puntos, sin casos especiales
    \item \textbf{Más rápidas}: Menos operaciones de campo por operación de grupo
    \item \textbf{Resistentes a side-channels}: Ejecución en tiempo constante natural
\end{enumerate}
\end{keyidea}

\begin{mathnote}
\textbf{Ley de grupo en curvas de Edwards:}

Para $P_1 = (x_1, y_1)$ y $P_2 = (x_2, y_2)$:
\[ P_1 + P_2 = \left( \frac{x_1 y_2 + x_2 y_1}{1 + d x_1 x_2 y_1 y_2}, \frac{y_1 y_2 - a x_1 x_2}{1 - d x_1 x_2 y_1 y_2} \right) \]

Esta fórmula:
\begin{itemize}
    \item Funciona para suma de puntos distintos
    \item Funciona para duplicación ($P_1 = P_2$)
    \item Funciona cuando uno de los puntos es el neutro $(0, 1)$
    \item Nunca tiene divisiones por cero (si $d$ es no-cuadrado)
\end{itemize}

Compárese con Weierstrass, que tiene 3+ casos especiales.
\end{mathnote}

%-------------------------------------------
\subsubsection{Curve25519 y Ed25519}

\begin{definition}[Curve25519]
Curve25519 es una curva de Montgomery:
\[ y^2 = x^3 + 486662x^2 + x \]

sobre $\mathbb{F}_p$ donde $p = 2^{255} - 19$.

Su forma twisted Edwards equivalente (Edwards25519) es:
\[ -x^2 + y^2 = 1 - \frac{121665}{121666} x^2 y^2 \]

El número 25519 viene del primo $p = 2^{255} - 19$.
\end{definition}

\begin{mathnote}
\textbf{Parámetros de Ed25519:}

\begin{verbatim}
# Campo
p = 2^255 - 19
  = 0x7fffffffffffffffffffffffffffffffffffffffffffffffffffffffffffffed

# Curva (twisted Edwards)
a = -1
d = -121665/121666 mod p
  = 0x52036cee2b6ffe738cc740797779e89800700a4d4141d8ab75eb4dca135978a3

# Orden del subgrupo
n = 2^252 + 27742317777372353535851937790883648493
  = 0x1000000000000000000000000000000014def9dea2f79cd65812631a5cf5d3ed

# Cofactor
h = 8

# Punto base B
B_x = 0x216936d3cd6e53fec0a4e231fdd6dc5c692cc7609525a7b2c9562d608f25d51a
B_y = 0x6666666666666666666666666666666666666666666666666666666666666658
\end{verbatim}
\end{mathnote}

\begin{remark}
\textbf{¿Por qué $p = 2^{255} - 19$?}

\begin{enumerate}
    \item Es primo
    \item Permite aritmética muy rápida (reducción módulo $2^{255} - 19$ es eficiente)
    \item 255 bits caben perfectamente en 32 bytes con 1 bit de sobra
    \item 19 es el menor entero positivo que hace $2^{255} - 19$ primo
\end{enumerate}
\end{remark}

%-------------------------------------------
\subsubsection{El Algoritmo EdDSA}

\begin{definition}[EdDSA - Algoritmo Completo]
\textbf{Parámetros:} Curva Edwards con punto base $B$ de orden $n$, hash $H$.

\textbf{KeyGen:}
\begin{enumerate}
    \item Generar semilla aleatoria: $seed \xleftarrow{\$} \{0,1\}^{256}$
    \item Expandir con hash: $(h_L, h_R) = H(seed)$ (64 bytes)
    \item Procesar $h_L$: ajustar bits para clamping
    \item Clave privada efectiva: $s = \text{clamp}(h_L)$
    \item Clave pública: $A = s \cdot B$
\end{enumerate}

\textbf{Sign$(seed, A, m)$:}
\begin{enumerate}
    \item Calcular nonce determinístico: $r = H(h_R \| m) \mod n$
    \item Calcular commitment: $R = r \cdot B$
    \item Calcular challenge: $k = H(R \| A \| m) \mod n$
    \item Calcular response: $S = (r + k \cdot s) \mod n$
    \item Retornar firma: $\sigma = (R, S)$
\end{enumerate}

\textbf{Verify$(A, m, \sigma)$:}
\begin{enumerate}
    \item Parsear $\sigma = (R, S)$
    \item Calcular $k = H(R \| A \| m) \mod n$
    \item Verificar: $8 \cdot S \cdot B \stackrel{?}{=} 8 \cdot R + 8 \cdot k \cdot A$
\end{enumerate}
\end{definition}

\begin{keyidea}
\textbf{El nonce determinístico es la innovación clave:}
\[ r = H(h_R \| m) \]

Donde $h_R$ es la mitad derecha del hash de la semilla (nunca revelada).

Propiedades:
\begin{itemize}
    \item \textbf{Determinístico}: Mismo mensaje = mismo nonce = misma firma
    \item \textbf{Impredecible}: Sin conocer $h_R$, no se puede calcular $r$
    \item \textbf{Sin riesgo de repetición}: Diferentes mensajes $\to$ diferentes nonces
    \item \textbf{Sin PRNG}: No depende de la calidad del generador aleatorio
\end{itemize}
\end{keyidea}

\begin{advertencia}
\textbf{El cofactor 8:}

Ed25519 tiene cofactor $h = 8$, lo que significa que hay subgrupos pequeños. La verificación multiplica por 8 para proyectar al subgrupo principal:
\[ 8 \cdot S \cdot B = 8 \cdot R + 8 \cdot k \cdot A \]

Sin esto, un atacante podría crear firmas ``válidas'' usando puntos de orden pequeño.

Variante: \textbf{Ed25519ctx} y \textbf{Ed25519ph} añaden domain separation.
\end{advertencia}

%-------------------------------------------
\subsubsection{Clamping: ¿Qué y Por Qué?}

\begin{definition}[Clamping de clave privada]
El ``clamping'' ajusta los bits de $h_L$ antes de usarlo como escalar:

\begin{verbatim}
def clamp(h):
    h[0] &= 248    # Limpiar los 3 bits inferiores
    h[31] &= 127   # Limpiar el bit más alto
    h[31] |= 64    # Establecer el segundo bit más alto
    return h
\end{verbatim}

Resultado: El escalar $s$ siempre tiene:
\begin{itemize}
    \item Bits 0, 1, 2 = 0 (divisible por 8)
    \item Bit 254 = 1
    \item Bit 255 = 0
\end{itemize}
\end{definition}

\begin{mathnote}
\textbf{¿Por qué clamping?}

\textbf{1. Divisible por cofactor (8):}
Garantiza que $s \cdot B$ está en el subgrupo principal, evitando ataques de subgrupo pequeño.

\textbf{2. Bit 254 = 1:}
Garantiza que la multiplicación escalar siempre toma el mismo número de pasos (tiempo constante).

\textbf{3. Bit 255 = 0:}
El escalar cabe en 255 bits, evitando overflow.
\end{mathnote}

%-------------------------------------------
\subsubsection{Comparación de Rendimiento}

\begin{center}
\begin{tabular}{lccc}
\toprule
\textbf{Operación} & \textbf{Ed25519} & \textbf{ECDSA P-256} & \textbf{RSA-2048} \\
\midrule
KeyGen & 30 µs & 80 µs & 300 ms \\
Sign & 50 µs & 100 µs & 10 ms \\
Verify & 120 µs & 200 µs & 0.3 ms \\
Firma tamaño & 64 B & 64 B & 256 B \\
Clave pública & 32 B & 64 B & 256 B \\
Clave privada & 32 B & 32 B & 256+ B \\
\bottomrule
\end{tabular}
\end{center}

\begin{intuicion}
Ed25519 es $\sim$2x más rápido que ECDSA y órdenes de magnitud más rápido que RSA. Además, las claves son más pequeñas, lo que reduce ancho de banda y almacenamiento.
\end{intuicion}

%-------------------------------------------
\subsubsection{Ed25519 en la Práctica}

\begin{example}[Ed25519 en diferentes lenguajes]
\textbf{Python (usando PyNaCl):}
\begin{verbatim}
from nacl.signing import SigningKey, VerifyKey

# KeyGen
sk = SigningKey.generate()
pk = sk.verify_key

# Sign
message = b"Hello, Ed25519!"
signed = sk.sign(message)
signature = signed.signature  # 64 bytes

# Verify
try:
    pk.verify(signed)
    print("Valid signature!")
except:
    print("Invalid signature!")
\end{verbatim}

\textbf{JavaScript (usando @noble/ed25519):}
\begin{verbatim}
import * as ed from '@noble/ed25519';

// KeyGen
const privateKey = ed.utils.randomPrivateKey();
const publicKey = await ed.getPublicKey(privateKey);

// Sign
const message = new TextEncoder().encode("Hello, Ed25519!");
const signature = await ed.sign(message, privateKey);

// Verify
const isValid = await ed.verify(signature, message, publicKey);
\end{verbatim}

\textbf{Rust (usando ed25519-dalek):}
\begin{verbatim}
use ed25519_dalek::{Keypair, Signature, Signer, Verifier};
use rand::rngs::OsRng;

// KeyGen
let mut csprng = OsRng{};
let keypair = Keypair::generate(&mut csprng);

// Sign
let message = b"Hello, Ed25519!";
let signature = keypair.sign(message);

// Verify
assert!(keypair.public.verify(message, &signature).is_ok());
\end{verbatim}
\end{example}

%-------------------------------------------
\subsubsection{Variantes de EdDSA}

\begin{center}
\begin{tabular}{lccl}
\toprule
\textbf{Variante} & \textbf{Curva} & \textbf{Hash} & \textbf{Uso} \\
\midrule
Ed25519 & Curve25519 & SHA-512 & General, más común \\
Ed25519ctx & Curve25519 & SHA-512 + context & Cuando se necesita domain separation \\
Ed25519ph & Curve25519 & SHA-512(SHA-512(m)) & Para mensajes muy grandes \\
Ed448 & Curve448 & SHAKE256 & Mayor seguridad (224 bits) \\
\bottomrule
\end{tabular}
\end{center}

\begin{remark}
\textbf{Ed448 (Goldilocks):}

Para aplicaciones que necesitan más de 128 bits de seguridad, Ed448 ofrece $\sim$224 bits usando la curva Curve448 (también diseñada por djb):
\[ p = 2^{448} - 2^{224} - 1 \]

Es más lento que Ed25519 pero sigue siendo muy rápido comparado con alternativas.
\end{remark}

%-------------------------------------------
\subsubsection{Ristretto: Resolviendo el Problema del Cofactor}

\begin{definition}[Ristretto255]
Ristretto es una técnica para construir un grupo primo a partir de una curva con cofactor.

Para Curve25519 (cofactor 8), Ristretto255 proporciona:
\begin{itemize}
    \item Un grupo de orden primo $n$ (sin subgrupos pequeños)
    \item Codificación canónica de puntos
    \item Compatibilidad con implementaciones existentes de Curve25519
\end{itemize}
\end{definition}

\begin{intuicion}
El problema del cofactor causa bugs sutiles:
\begin{itemize}
    \item Diferentes codificaciones del mismo punto
    \item Ataques de subgrupo pequeño si no se verifica membresía
    \item Complejidad en protocolos multi-party
\end{itemize}

Ristretto ``quotients out'' los puntos de torsión, dando una abstracción limpia de grupo primo.
\end{intuicion}

\begin{remark}
\textbf{Sr25519 (Schnorrkel):}

Polkadot usa Sr25519, que es Schnorr sobre Ristretto255. Combina:
\begin{itemize}
    \item La velocidad de Curve25519
    \item La limpieza de un grupo primo (Ristretto)
    \item Las ventajas de Schnorr (agregación, etc.)
\end{itemize}
\end{remark}

%-------------------------------------------
\subsubsection{Ed25519 en Circuitos ZK}

\begin{mathnote}
Ed25519 es particularmente eficiente para verificación en circuitos ZK, especialmente cuando se usa una curva ``embebida''.

\textbf{Curva embebida:} Una curva cuyo orden es el campo escalar del sistema SNARK.

Para SNARKs sobre BLS12-381, se usa \textbf{Jubjub}:
\begin{itemize}
    \item Curva twisted Edwards sobre $\mathbb{F}_q$ (campo escalar de BLS12-381)
    \item Orden del subgrupo primo $\approx 2^{251}$
    \item Operaciones de Jubjub son ``nativas'' en el circuito
\end{itemize}

\textbf{Costo comparativo:}
\begin{center}
\begin{tabular}{lc}
\toprule
\textbf{Verificación} & \textbf{Constraints} \\
\midrule
ECDSA (secp256k1) en BN254 & $\sim$600,000 \\
EdDSA (Ed25519) en BN254 & $\sim$400,000 \\
EdDSA (Jubjub) en BLS12-381 & $\sim$15,000 \\
\bottomrule
\end{tabular}
\end{center}

¡40x más eficiente con curva embebida!
\end{mathnote}

%-------------------------------------------
\subsubsection{Aplicación en Banco ZK}

\begin{bancozk}
EdDSA es el esquema de firma preferido para operaciones internas de Banco ZK:

\textbf{1. Firmas dentro de circuitos}

Para verificar firmas de usuario dentro del SNARK, usamos EdDSA sobre una curva embebida:
\begin{center}
\begin{tabular}{lcc}
\toprule
\textbf{SNARK base} & \textbf{Curva de firma} & \textbf{Constraints} \\
\midrule
BN254 (Groth16) & Baby Jubjub & $\sim$20,000 \\
BLS12-381 (Halo2) & Jubjub & $\sim$15,000 \\
\bottomrule
\end{tabular}
\end{center}

Comparado con verificar ECDSA ($\sim$600,000), esto es 30-40x más eficiente.

\textbf{2. Por qué no usamos Ed25519 para todo}
\begin{center}
\begin{tabular}{lcc}
\toprule
\textbf{Aspecto} & \textbf{Ed25519} & \textbf{ECDSA secp256k1} \\
\midrule
Wallets compatibles & Pocas & Todas (Ethereum) \\
Precompilado EVM & No & Sí (ecrecover) \\
Verificación on-chain & $\sim$50,000 gas & $\sim$3,000 gas \\
En circuito ZK & Eficiente & Muy costoso \\
\bottomrule
\end{tabular}
\end{center}

\textbf{3. Arquitectura híbrida}

\begin{verbatim}
Usuario (L1) ──ECDSA──> Contrato ──┐
                                   │
                          [verificación L1]
                                   │
                                   v
Circuito ZK <───EdDSA/Jubjub───> Estado interno
\end{verbatim}

El usuario firma con ECDSA (compatible con MetaMask), pero las operaciones internas del protocolo usan EdDSA sobre curva embebida.

\textbf{4. Migración futura}

Con Account Abstraction (EIP-4337), los usuarios podrán usar wallets Ed25519:
\begin{itemize}
    \item Verificación delegada a contratos custom
    \item Sin dependencia de precompilados
    \item Mejor UX móvil (Ed25519 más rápido en ARM)
\end{itemize}
\end{bancozk}

%===========================================
% SECCIÓN: BLS
%===========================================

\subsection{BLS}

\begin{definition}[Firma BLS]
Firma basada en pairings con agregación nativa.
\end{definition}

Sea $e: \G_1 \times \G_2 \to \G_T$ un pairing.
\begin{enumerate}
    \item \textbf{Keygen:} $sk \in \Z_q$, $pk = sk \cdot g_2 \in \G_2$
    \item \textbf{Sign:} $\sigma = sk \cdot H(m) \in \G_1$
    \item \textbf{Verify:} $e(\sigma, g_2) = e(H(m), pk)$
\end{enumerate}

\begin{keyidea}
\textbf{Agregación:} Múltiples firmas se combinan en una:
\[
\sigma_{\text{agg}} = \sum_i \sigma_i, \quad \text{verifica con } \sum_i pk_i
\]
Ethereum 2.0 agrega miles de firmas de validadores en una sola.
\end{keyidea}

%-------------------------------------------
\subsubsection{Contexto Histórico}

\begin{intuicion}
Las firmas BLS (Boneh-Lynn-Shacham) fueron propuestas en 2001 y representan un paradigma diferente: usan \textbf{pairings} en curvas elípticas, una operación matemática más poderosa que permite cosas imposibles con firmas tradicionales.
\end{intuicion}

\begin{remark}
\textbf{Los autores:}
\begin{itemize}
    \item \textbf{Dan Boneh}: Profesor en Stanford, pionero en criptografía de pairings
    \item \textbf{Ben Lynn}: PhD de Stanford, ahora en Google
    \item \textbf{Hovav Shacham}: Profesor en UT Austin
\end{itemize}

El paper original ``Short Signatures from the Weil Pairing'' (2001) mostró que los pairings permitían firmas más cortas que cualquier esquema conocido hasta entonces.
\end{remark}

\begin{center}
\begin{tabular}{ll}
\toprule
\textbf{Año} & \textbf{Desarrollo} \\
\midrule
2001 & Paper original de BLS \\
2003 & Agregación de firmas propuesta \\
2017 & Ethereum 2.0 elige BLS para consenso \\
2018 & Zcash Sapling usa BLS12-381 (para SNARKs) \\
2020 & Beacon Chain lanzada con firmas BLS \\
2022 & The Merge: BLS en producción masiva \\
\bottomrule
\end{tabular}
\end{center}

%-------------------------------------------
\subsubsection{Pairings: El Ingrediente Mágico}

\begin{definition}[Pairing Bilineal (Repaso)]
Un pairing es una función:
\[ e: \G_1 \times \G_2 \to \G_T \]

que satisface \textbf{bilinealidad}:
\[ e(aP, bQ) = e(P, Q)^{ab} \]

Propiedades adicionales:
\begin{itemize}
    \item \textbf{No degenerado}: $e(G_1, G_2) \neq 1$ para generadores
    \item \textbf{Eficientemente computable}: $O(\log n)$ operaciones de campo
\end{itemize}
\end{definition}

\begin{keyidea}
Los pairings permiten una forma de ``multiplicación'' en curvas elípticas que no existe de otra forma:
\begin{itemize}
    \item Sin pairing: podemos sumar puntos y multiplicar por escalares
    \item Con pairing: podemos ``combinar'' información de dos puntos de manera bilineal
\end{itemize}

Esto permite construcciones criptográficas imposibles sin pairings, como:
\begin{itemize}
    \item Firmas cortas (BLS)
    \item Identity-Based Encryption (IBE)
    \item Verificación de SNARKs en tiempo constante
\end{itemize}
\end{keyidea}

%-------------------------------------------
\subsubsection{El Algoritmo BLS en Detalle}

\begin{definition}[BLS Signature Scheme]
\textbf{Setup:} Curva pairing-friendly con grupos $\G_1, \G_2, \G_T$ de orden $q$, generadores $g_1 \in \G_1$, $g_2 \in \G_2$, y hash-to-curve $H: \{0,1\}^* \to \G_1$.

\textbf{KeyGen:}
\begin{enumerate}
    \item $sk \xleftarrow{\$} \Z_q$
    \item $pk = sk \cdot g_2 \in \G_2$
\end{enumerate}

\textbf{Sign$(sk, m)$:}
\begin{enumerate}
    \item Hashear mensaje a punto de curva: $h = H(m) \in \G_1$
    \item Firma: $\sigma = sk \cdot h \in \G_1$
\end{enumerate}

\textbf{Verify$(pk, m, \sigma)$:}
\begin{enumerate}
    \item Calcular $h = H(m)$
    \item Verificar: $e(\sigma, g_2) \stackrel{?}{=} e(h, pk)$
\end{enumerate}
\end{definition}

\begin{mathnote}
\textbf{¿Por qué funciona?}

\begin{align}
e(\sigma, g_2) &= e(sk \cdot H(m), g_2) \\
               &= e(H(m), g_2)^{sk} \quad \text{(bilinealidad)} \\
               &= e(H(m), sk \cdot g_2) \\
               &= e(H(m), pk) \quad \checkmark
\end{align}
\end{mathnote}

\begin{intuicion}
La firma BLS es increíblemente simple:
\begin{itemize}
    \item KeyGen: multiplicar generador por escalar
    \item Sign: multiplicar hash por escalar
    \item Verify: comparar dos pairings
\end{itemize}

¡No hay nonces! Esto elimina completamente el riesgo de nonce repetido que afecta a ECDSA y Schnorr.
\end{intuicion}

%-------------------------------------------
\subsubsection{Hash-to-Curve}

\begin{advertencia}
BLS requiere hashear mensajes a puntos de curva: $H: \{0,1\}^* \to \G_1$.

Esto es más complejo que hashear a enteros:
\begin{enumerate}
    \item Hashear a un valor de campo: $h' = H'(m) \in \Fp$
    \item Mapear a la curva usando un algoritmo como Simplified SWU
    \item Multiplicar por cofactor si es necesario
\end{enumerate}

Un hash-to-curve incorrecto puede causar vulnerabilidades graves.
\end{advertencia}

\begin{remark}
El estándar \textbf{RFC 9380} (``Hashing to Elliptic Curves'') especifica métodos seguros para hash-to-curve. Para BLS12-381:
\begin{itemize}
    \item $\G_1$: hash\_to\_field + map\_to\_curve (Simplified SWU) + clear\_cofactor
    \item $\G_2$: Similar pero sobre la extensión cuadrática
\end{itemize}
\end{remark}

%-------------------------------------------
\subsubsection{Agregación de Firmas: El Superpoder de BLS}

\begin{keyidea}
La característica más importante de BLS es la \textbf{agregación nativa}:

\textbf{Múltiples firmas sobre el mismo mensaje:}
\begin{align}
\sigma_{\text{agg}} &= \sigma_1 + \sigma_2 + \cdots + \sigma_n = \sum_i sk_i \cdot H(m) \\
pk_{\text{agg}} &= pk_1 + pk_2 + \cdots + pk_n = \sum_i pk_i
\end{align}

\textbf{Verificación:}
\[ e(\sigma_{\text{agg}}, g_2) = e(H(m), pk_{\text{agg}}) \]

¡Una sola verificación de pairing para $n$ firmas!
\end{keyidea}

\begin{mathnote}
\textbf{¿Por qué funciona?}
\begin{align}
e(\sigma_{\text{agg}}, g_2) &= e\left(\sum_i sk_i \cdot H(m), g_2\right) \\
                            &= e\left(H(m), g_2\right)^{\sum_i sk_i} \\
                            &= e\left(H(m), \sum_i sk_i \cdot g_2\right) \\
                            &= e(H(m), pk_{\text{agg}}) \quad \checkmark
\end{align}
\end{mathnote}

\begin{example}[Agregación en Ethereum 2.0]
\textbf{Contexto:} La Beacon Chain tiene $\sim$900,000 validadores. Cada slot (12 segundos), miles de validadores deben atestiguar (firmar) un bloque.

\textbf{Sin agregación:}
\begin{itemize}
    \item $\sim$100,000 firmas por slot
    \item 48 bytes × 100,000 = 4.8 MB por slot
    \item 100,000 verificaciones de pairing por slot
\end{itemize}

\textbf{Con agregación BLS:}
\begin{itemize}
    \item Validadores se agrupan en comités de $\sim$128
    \item Cada comité produce UNA firma agregada
    \item $\sim$800 firmas por slot
    \item 48 bytes × 800 = 38.4 KB por slot
    \item 800 verificaciones por slot
\end{itemize}

\textbf{Mejora:} 125x menos datos, 125x menos verificaciones.
\end{example}

%-------------------------------------------
\subsubsection{Agregación de Mensajes Diferentes}

\begin{definition}[Agregación General]
Cuando cada firmante firma un mensaje diferente:
\[ \sigma_{\text{agg}} = \sum_i \sigma_i = \sum_i sk_i \cdot H(m_i) \]

\textbf{Verificación (n pairings):}
\[ e(\sigma_{\text{agg}}, g_2) \stackrel{?}{=} \prod_i e(H(m_i), pk_i) \]

Requiere $n$ pairings pero solo transmite UNA firma.
\end{definition}

\begin{mathnote}
\textbf{Verificación batch:}

Con técnicas de multi-pairing, verificar $n$ pairings es más rápido que $n$ verificaciones individuales:
\[ \text{Costo de } n \text{ pairings} \approx 1 + 0.1 \times n \text{ pairings individuales} \]

Esto hace que la agregación de mensajes diferentes siga siendo muy eficiente.
\end{mathnote}

%-------------------------------------------
\subsubsection{El Ataque Rogue Key y Proof of Possession}

\begin{advertencia}
\textbf{Ataque Rogue Key en BLS:}

Similar a Schnorr, un atacante puede elegir su clave pública maliciosamente:
\[ pk_{\text{attack}} = sk_{\text{attack}} \cdot g_2 - pk_{\text{victim}} \]

Entonces:
\[ pk_{\text{agg}} = pk_{\text{victim}} + pk_{\text{attack}} = sk_{\text{attack}} \cdot g_2 \]

¡El atacante puede firmar por la ``víctima'' sin su consentimiento!
\end{advertencia}

\begin{definition}[Proof of Possession (PoP)]
Para prevenir ataques rogue key, cada participante debe demostrar que conoce la clave privada correspondiente a su clave pública:

\[ \text{PoP} = sk \cdot H_{\text{pop}}(pk) \]

Esto se verifica una vez al registrar la clave:
\[ e(\text{PoP}, g_2) = e(H_{\text{pop}}(pk), pk) \]

Después de esto, la clave puede usarse para agregación segura.
\end{definition}

\begin{remark}
Ethereum 2.0 requiere que cada validador registre un PoP al hacer el depósito de 32 ETH. Sin el PoP, el depósito es rechazado.
\end{remark}

%-------------------------------------------
\subsubsection{BLS vs Otros Esquemas}

\begin{center}
\begin{tabular}{lcccc}
\toprule
\textbf{Propiedad} & \textbf{ECDSA} & \textbf{Schnorr} & \textbf{EdDSA} & \textbf{BLS} \\
\midrule
Tamaño de firma & 64-72 B & 64 B & 64 B & 48 B \\
Sin nonces & No & No & Sí & Sí \\
Agregación nativa & No & Requiere MuSig & No & Sí \\
Determinístico & Opcional & Opcional & Sí & Sí \\
Requiere pairing & No & No & No & Sí \\
Velocidad sign & Rápido & Rápido & Muy rápido & Más lento \\
Velocidad verify & Rápido & Rápido & Rápido & Lento \\
Verify agregada & N/A & Batch & Batch & 1 pairing \\
\bottomrule
\end{tabular}
\end{center}

\begin{keyidea}
\textbf{Cuándo usar BLS:}
\begin{itemize}
    \item Muchos firmantes necesitan agregarse (consenso, multisig masivo)
    \item El tamaño de firma es crítico
    \item La velocidad de verificación individual no es el cuello de botella
\end{itemize}

\textbf{Cuándo NO usar BLS:}
\begin{itemize}
    \item Firmas individuales sin agregación
    \item Dispositivos con recursos limitados (pairings son costosos)
    \item Cuando no hay curvas pairing-friendly disponibles
\end{itemize}
\end{keyidea}

%-------------------------------------------
\subsubsection{Curvas para BLS}

\begin{center}
\begin{tabular}{lccc}
\toprule
\textbf{Curva} & \textbf{Seguridad} & \textbf{Firma} & \textbf{Clave pública} \\
\midrule
BN254 & $\sim$100 bits & 32 B ($\G_1$) & 64 B ($\G_2$) \\
BLS12-381 & $\sim$128 bits & 48 B ($\G_1$) & 96 B ($\G_2$) \\
BLS12-377 & $\sim$128 bits & 48 B & 96 B \\
\bottomrule
\end{tabular}
\end{center}

\begin{remark}
\textbf{¿Por qué firma en $\G_1$ y clave en $\G_2$?}

En BLS12-381:
\begin{itemize}
    \item $\G_1$: puntos de 48 bytes (comprimidos)
    \item $\G_2$: puntos de 96 bytes (comprimidos)
\end{itemize}

Como transmitimos muchas más firmas que claves públicas, ponemos firmas en el grupo más pequeño ($\G_1$).

Ethereum 2.0 usa esta configuración. Alternativamente, se pueden intercambiar si hay más verificaciones que firmas.
\end{remark}

%-------------------------------------------
\subsubsection{Threshold BLS}

\begin{definition}[Threshold BLS]
Firma $t$-de-$n$ donde cualquier $t$ participantes pueden crear una firma válida, pero menos de $t$ no pueden.

\textbf{Setup:} Shamir Secret Sharing distribuye $sk$ en $n$ shares $sk_1, \ldots, sk_n$.

\textbf{Firma (t participantes):}
\begin{enumerate}
    \item Cada participante $i$ calcula $\sigma_i = sk_i \cdot H(m)$
    \item Se combinan $t$ firmas parciales usando interpolación de Lagrange:
    \[ \sigma = \sum_{i \in S} \lambda_i \sigma_i \]
    donde $S$ es el conjunto de $t$ participantes y $\lambda_i$ son coeficientes de Lagrange.
\end{enumerate}

La firma resultante es indistinguible de una firma BLS normal.
\end{definition}

\begin{intuicion}
Threshold BLS es más simple que threshold Schnorr porque:
\begin{enumerate}
    \item No hay nonces que coordinar
    \item Cada firma parcial es determinística
    \item La combinación es simplemente suma ponderada
\end{enumerate}

Esto lo hace ideal para protocolos distribuidos.
\end{intuicion}

%-------------------------------------------
\subsubsection{BLS en Ethereum 2.0}

\begin{mathnote}
\textbf{Especificación de Ethereum 2.0:}

\textbf{Curva:} BLS12-381
\textbf{Hash:} SHA-256 + hash\_to\_G1
\textbf{Domain separation:} Diferentes dominios para diferentes tipos de mensajes

\textbf{Tipos de firmas:}
\begin{itemize}
    \item \textbf{Attestations}: Validadores atestiguan bloques (agregadas por comité)
    \item \textbf{Sync committee}: 512 validadores rotativos firman headers
    \item \textbf{Block proposals}: Firma individual del proposer
    \item \textbf{Voluntary exits}: Validador sale del protocolo
    \item \textbf{Slashings}: Evidencia de comportamiento malicioso
\end{itemize}

\textbf{Flujo de attestation:}
\begin{verbatim}
1. Validador recibe bloque B en slot S
2. Calcula: msg = (slot, block_root, source, target)
3. Firma: σ = sk · H(domain || msg)
4. Broadcast: (attestation_data, σ, validator_index)
5. Aggregador colecta σ₁, σ₂, ... de su comité
6. Produce: σ_agg = Σ σᵢ, bitfield de participantes
7. Beacon chain verifica: e(σ_agg, g₂) = e(H(msg), Σ pkᵢ)
\end{verbatim}
\end{mathnote}

%-------------------------------------------
\subsubsection{Rendimiento de BLS}

\begin{center}
\begin{tabular}{lcc}
\toprule
\textbf{Operación (BLS12-381)} & \textbf{Tiempo} & \textbf{Comparación} \\
\midrule
Hash to $\G_1$ & $\sim$0.3 ms & --- \\
Multiplicación escalar $\G_1$ & $\sim$0.5 ms & 5x más lento que Ed25519 \\
Multiplicación escalar $\G_2$ & $\sim$1.5 ms & --- \\
Pairing & $\sim$2 ms & --- \\
Verify 1 firma & $\sim$2.5 ms & 20x más lento que Ed25519 \\
Verify 100 firmas (agregadas) & $\sim$3 ms & 80x más rápido que individual \\
\bottomrule
\end{tabular}
\end{center}

\begin{keyidea}
BLS es lento para firmas individuales pero brilla con agregación:
\begin{itemize}
    \item 1 firma: $\sim$2.5 ms (20x más lento que Ed25519)
    \item 100 firmas agregadas: $\sim$3 ms (prácticamente igual que 1)
    \item 100 firmas Ed25519: $\sim$12 ms (4x más lento)
\end{itemize}

El crossover está alrededor de 4-5 firmas: con más de eso, BLS agregado gana.
\end{keyidea}

%-------------------------------------------
\subsubsection{Aplicación en Banco ZK}

\begin{bancozk}
BLS juega varios roles importantes en Banco ZK:

\textbf{1. Firmas de validadores L2}

Nuestro sequencer tiene un comité de validadores que firman batches:
\begin{verbatim}
// Comité de 21 validadores
// Requiere 14/21 (2/3+1) para finalizar un batch

batch_hash = keccak256(transactions, state_root, batch_number)
signatures = [validator.sign(batch_hash) for validator in committee]
agg_sig = aggregate(signatures)

// On-chain: verificar UNA firma en lugar de 14
verify(agg_sig, batch_hash, agg_pubkey)
\end{verbatim}

\textbf{Ahorro:}
\begin{itemize}
    \item Sin BLS: 14 × ecrecover = 42,000 gas
    \item Con BLS: 1 verificación = $\sim$100,000 gas (más caro individualmente)
\end{itemize}

¿Por qué BLS es más caro? No hay precompilado BLS en EVM. Pero el ahorro viene en:
\begin{itemize}
    \item Calldata: 14 × 65 bytes = 910 bytes → 48 bytes
    \item Con EIP-2537: $\sim$50,000 gas (mejor que ECDSA)
\end{itemize}

\textbf{2. Agregación de depósitos}

Múltiples usuarios pueden agregar sus firmas de depósito off-chain:
\begin{verbatim}
// 100 usuarios quieren depositar
// Cada uno firma: "deposit(amount, commitment)"

off_chain_aggregate:
    agg_sig = sum(user_signatures)
    agg_pk = sum(user_pubkeys)
    
on_chain:
    verify_aggregate(agg_sig, deposit_data, agg_pk)
    process_deposits(deposit_list)
\end{verbatim}

\textbf{3. Sync committee para bridges}

Para verificar estado de Ethereum en otras cadenas (o viceversa):
\begin{itemize}
    \item 512 validadores firman headers (sync committee)
    \item La firma agregada permite light client verification
    \item Banco ZK puede verificar estado de Ethereum L1 eficientemente
\end{itemize}

\textbf{4. Por qué no usamos BLS para todo}

\begin{center}
\begin{tabular}{lcc}
\toprule
\textbf{Uso} & \textbf{Esquema elegido} & \textbf{Razón} \\
\midrule
Usuarios (L1) & ECDSA & Wallets existentes \\
Validadores L2 & BLS & Agregación \\
Dentro de SNARK & EdDSA & Eficiencia en circuito \\
Multisig governance & BLS threshold & Sin coordinación de nonces \\
\bottomrule
\end{tabular}
\end{center}

\textbf{5. Futuro: BLS para usuarios}

Con Account Abstraction, los usuarios podrán optar por BLS:
\begin{itemize}
    \item Wallet contract verifica firma BLS
    \item Múltiples operaciones del mismo usuario se agregan
    \item Ahorro significativo en gas para usuarios activos
\end{itemize}
\end{bancozk}

%-------------------------------------------
\subsubsection{Librerías BLS}

\begin{center}
\begin{tabular}{lll}
\toprule
\textbf{Lenguaje} & \textbf{Librería} & \textbf{Notas} \\
\midrule
C & blst (Supranational) & La más rápida, usada en Ethereum \\
Rust & blst-rs & Bindings de blst \\
Rust & arkworks bls12-381 & Más flexible \\
Go & gnark-crypto/bls12-381 & Producción \\
JavaScript & @noble/curves & Pura JS, auditada \\
Python & py\_ecc & Educacional \\
\bottomrule
\end{tabular}
\end{center}

\begin{example}[BLS con @noble/curves]
\begin{verbatim}
import { bls12_381 as bls } from '@noble/curves/bls12-381';

// KeyGen
const privateKey = bls.utils.randomPrivateKey();
const publicKey = bls.getPublicKey(privateKey);

// Sign
const message = new TextEncoder().encode("Hello, BLS!");
const signature = bls.sign(message, privateKey);

// Verify
const isValid = bls.verify(signature, message, publicKey);

// Aggregation (same message)
const sk1 = bls.utils.randomPrivateKey();
const sk2 = bls.utils.randomPrivateKey();
const pk1 = bls.getPublicKey(sk1);
const pk2 = bls.getPublicKey(sk2);

const sig1 = bls.sign(message, sk1);
const sig2 = bls.sign(message, sk2);

const aggSig = bls.aggregateSignatures([sig1, sig2]);
const aggPk = bls.aggregatePublicKeys([pk1, pk2]);

const isAggValid = bls.verify(aggSig, message, aggPk);
console.log("Aggregated signature valid:", isAggValid);
\end{verbatim}
\end{example}

%-------------------------------------------
\subsubsection{Resumen: El Ecosistema de Firmas}

\begin{keyidea}
\textbf{Guía de selección de esquema de firmas:}

\begin{center}
\begin{tabular}{ll}
\toprule
\textbf{Caso de uso} & \textbf{Esquema recomendado} \\
\midrule
Compatibilidad Ethereum/Bitcoin & ECDSA \\
Nuevos proyectos, velocidad máxima & EdDSA (Ed25519) \\
Agregación de muchas firmas & BLS \\
Multisig que parece normal & Schnorr (MuSig2) \\
Threshold sin coordinación & BLS threshold \\
Dentro de circuitos ZK & EdDSA sobre curva embebida \\
Post-cuántico & Dilithium, SPHINCS+ \\
\bottomrule
\end{tabular}
\end{center}

En la práctica, los sistemas modernos como Banco ZK usan \textbf{múltiples esquemas} optimizando para cada contexto.
\end{keyidea}

%===========================================
% SECCIÓN: MERKLE TREES - VERSIÓN EXPANDIDA
% Usa tus entornos: definition, example, keyidea, bancozk, 
%                   mathnote, intuicion, advertencia, remark
%===========================================

\section{Merkle Trees}

\begin{definition}[Merkle Tree]
Árbol binario donde cada nodo interno es el hash de sus hijos.
\end{definition}

Propiedades:
\begin{itemize}
    \item La raíz compromete a todos los datos
    \item Prueba de inclusión: $O(\log n)$ hashes
    \item Actualización: $O(\log n)$
    \item Usado en: Bitcoin (tx), Ethereum (state), ZK (commitments)
\end{itemize}

%-------------------------------------------
\subsection{Contexto Histórico}

\begin{intuicion}
Los Merkle trees son una de las estructuras de datos más importantes en criptografía aplicada. Permiten \textbf{comprometer} grandes cantidades de datos en un solo hash (la raíz) y después \textbf{probar} la inclusión de cualquier elemento sin revelar todo el conjunto.
\end{intuicion}

\begin{remark}
\textbf{Ralph Merkle y el nacimiento del hashing:}

Ralph Merkle propuso esta estructura en 1979 como parte de su tesis doctoral en Stanford, en el mismo período que inventó el concepto de criptografía de clave pública (independientemente de Diffie-Hellman).

\begin{center}
\begin{tabular}{ll}
\toprule
\textbf{Año} & \textbf{Desarrollo} \\
\midrule
1979 & Merkle propone el árbol hash \\
1980 & Patente US 4,309,569 \\
1991 & Uso en timestamps digitales (Haber-Stornetta) \\
2008 & Bitcoin usa Merkle trees para transacciones \\
2015 & Ethereum usa Patricia Merkle Tries para estado \\
2016 & Zcash usa Merkle trees para notas privadas \\
\bottomrule
\end{tabular}
\end{center}

Curiosamente, los papers de Haber y Stornetta sobre timestamping (que usaban Merkle trees) están citados en el whitepaper de Bitcoin.
\end{remark}

%-------------------------------------------
\subsection{Construcción Básica}

\begin{definition}[Construcción de Merkle Tree]
Dado un conjunto de $n$ elementos (hojas) $L_0, L_1, \ldots, L_{n-1}$ y una función hash $H$:

\textbf{Nivel 0 (hojas):}
\[ h_i^{(0)} = H(L_i) \]

\textbf{Niveles superiores:}
\[ h_i^{(k+1)} = H(h_{2i}^{(k)} \| h_{2i+1}^{(k)}) \]

\textbf{Raíz:}
\[ \text{root} = h_0^{(\lceil \log_2 n \rceil)} \]
\end{definition}

\begin{center}
\begin{tikzpicture}[
    level distance=1.5cm,
    sibling distance=3cm,
    every node/.style={draw, circle, minimum size=0.8cm, font=\small},
    edge from parent/.style={draw, -}
]
    \node {root}
        child {node {$h_{01}$}
            child {node {$h_0$}
                child[sibling distance=1.5cm] {node[rectangle] {$L_0$}}
            }
            child {node {$h_1$}
                child[sibling distance=1.5cm] {node[rectangle] {$L_1$}}
            }
        }
        child {node {$h_{23}$}
            child {node {$h_2$}
                child[sibling distance=1.5cm] {node[rectangle] {$L_2$}}
            }
            child {node {$h_3$}
                child[sibling distance=1.5cm] {node[rectangle] {$L_3$}}
            }
        };
\end{tikzpicture}
\end{center}

\begin{example}[Construcción paso a paso]
\textbf{Datos:} 4 transacciones $\text{tx}_0, \text{tx}_1, \text{tx}_2, \text{tx}_3$

\textbf{Nivel 0 (hojas):}
\begin{align}
h_0 &= H(\text{tx}_0) = \texttt{a1b2...} \\
h_1 &= H(\text{tx}_1) = \texttt{c3d4...} \\
h_2 &= H(\text{tx}_2) = \texttt{e5f6...} \\
h_3 &= H(\text{tx}_3) = \texttt{7890...}
\end{align}

\textbf{Nivel 1:}
\begin{align}
h_{01} &= H(h_0 \| h_1) = H(\texttt{a1b2...} \| \texttt{c3d4...}) = \texttt{1234...} \\
h_{23} &= H(h_2 \| h_3) = H(\texttt{e5f6...} \| \texttt{7890...}) = \texttt{5678...}
\end{align}

\textbf{Nivel 2 (raíz):}
\[ \text{root} = H(h_{01} \| h_{23}) = H(\texttt{1234...} \| \texttt{5678...}) = \texttt{abcd...} \]

La raíz \texttt{abcd...} compromete a las 4 transacciones.
\end{example}

%-------------------------------------------
\subsection{Merkle Proof (Prueba de Inclusión)}

Para probar que $L_i$ está en el árbol:
\begin{enumerate}
    \item Proporcionar los ``hermanos'' en el camino a la raíz
    \item Verificador recalcula hashes hasta la raíz
    \item Compara con raíz conocida
\end{enumerate}

Tamaño de prueba: $\log_2(n)$ hashes.

\begin{definition}[Merkle Proof]
Una prueba de inclusión para el elemento $L_i$ en un árbol con raíz $r$ consiste en:
\begin{enumerate}
    \item El elemento $L_i$ (o su hash $h_i$)
    \item El \textbf{authentication path}: secuencia de hashes hermanos desde la hoja hasta la raíz
    \item Los \textbf{direction bits}: indica si cada hermano va a la izquierda o derecha
\end{enumerate}

\[ \text{proof} = (L_i, [(s_0, d_0), (s_1, d_1), \ldots, (s_{k-1}, d_{k-1})]) \]

donde $s_j$ es el hash hermano y $d_j \in \{L, R\}$ es su posición.
\end{definition}

\begin{example}[Merkle proof para $L_1$]
Queremos probar que $\text{tx}_1$ está en el árbol.

\textbf{Authentication path:}
\begin{enumerate}
    \item Hermano de $h_1$: $h_0$ (a la izquierda)
    \item Hermano de $h_{01}$: $h_{23}$ (a la derecha)
\end{enumerate}

\textbf{Proof:} $(\text{tx}_1, [(h_0, L), (h_{23}, R)])$

\textbf{Verificación:}
\begin{enumerate}
    \item $h_1 = H(\text{tx}_1)$
    \item $h_{01} = H(h_0 \| h_1)$ (hermano $h_0$ a la izquierda)
    \item $\text{root}' = H(h_{01} \| h_{23})$ (hermano $h_{23}$ a la derecha)
    \item Verificar: $\text{root}' \stackrel{?}{=} \text{root}_{\text{conocida}}$
\end{enumerate}
\end{example}

\begin{center}
\begin{tikzpicture}[
    level distance=1.5cm,
    sibling distance=3cm,
    every node/.style={draw, circle, minimum size=0.8cm, font=\small},
    edge from parent/.style={draw, -},
    highlight/.style={fill=green!30},
    sibling/.style={fill=blue!20}
]
    \node[highlight] {root}
        child {node[highlight] {$h_{01}$}
            child {node[sibling] {$h_0$}}
            child {node[highlight] {$h_1$}
                child[sibling distance=1.5cm] {node[rectangle, highlight] {$L_1$}}
            }
        }
        child {node[sibling] {$h_{23}$}
            child {node {$h_2$}}
            child {node {$h_3$}}
        };
    
    % Leyenda
    \node[draw=none, fill=none] at (5, 0) {\textcolor{green!50!black}{Verde: Calculado}};
    \node[draw=none, fill=none] at (5, -0.5) {\textcolor{blue!50!black}{Azul: Proporcionado}};
\end{tikzpicture}
\end{center}

%-------------------------------------------
\subsection{Propiedades y Complejidad}

\begin{keyidea}
\textbf{Complejidad de operaciones en Merkle Trees:}

\begin{center}
\begin{tabular}{lcc}
\toprule
\textbf{Operación} & \textbf{Tiempo} & \textbf{Espacio} \\
\midrule
Construir árbol & $O(n)$ & $O(n)$ \\
Calcular raíz & $O(n)$ & $O(\log n)$ (streaming) \\
Generar proof & $O(\log n)$ & $O(\log n)$ \\
Verificar proof & $O(\log n)$ & $O(\log n)$ \\
Actualizar hoja & $O(\log n)$ & $O(\log n)$ \\
\bottomrule
\end{tabular}
\end{center}

Para $n = 2^{20} \approx 1$ millón de elementos, el proof tiene solo 20 hashes.
\end{keyidea}

\begin{mathnote}
\textbf{Propiedades de seguridad:}

Sea $H$ una función hash resistente a colisiones.

\textbf{1. Binding:} Es computacionalmente imposible encontrar dos conjuntos de hojas diferentes que produzcan la misma raíz. (Se reduce a resistencia a colisiones de $H$.)

\textbf{2. Soundness de proofs:} Un proof válido para un elemento $L$ con raíz $r$ implica que $L$ estaba en el conjunto original (salvo encontrar colisiones en $H$).

\textbf{3. Hiding (opcional):} El árbol por sí solo no oculta información. Para privacidad, se pueden añadir hojas aleatorias o usar hashes con salt.
\end{mathnote}

%-------------------------------------------
\subsection{Árboles con Número Impar de Hojas}

\begin{advertencia}
¿Qué pasa si $n$ no es potencia de 2?

\textbf{Opción 1: Duplicar último elemento}
\[ h_{n} = h_{n-1} \]
Problema: Dos conjuntos diferentes pueden tener la misma raíz.

\textbf{Opción 2: Padding con hash especial}
\[ h_{n} = H(\texttt{0x00}) \]
Problema: Inconsistencia si el conjunto crece.

\textbf{Opción 3: Árbol desbalanceado}
Permitir nodos con un solo hijo, propagando el hash directamente.

\textbf{Opción 4: Fijar profundidad} (común en ZK)
Usar un árbol de profundidad fija $d$, con hojas vacías = hash de cero.
\end{advertencia}

\begin{example}[Árbol de profundidad fija]
Para un árbol de profundidad 3 (capacidad 8) con 5 elementos:
\begin{verbatim}
Hojas: [L0, L1, L2, L3, L4, EMPTY, EMPTY, EMPTY]
donde EMPTY = H(0)
\end{verbatim}

Esto es común en protocolos ZK donde la profundidad del árbol es un parámetro del circuito.
\end{example}

%-------------------------------------------
\subsection{Merkle Trees en Bitcoin}

\begin{definition}[Merkle Root de Bitcoin]
Cada bloque de Bitcoin contiene un Merkle root que compromete todas las transacciones del bloque.

\begin{verbatim}
Block Header:
  - version
  - prev_block_hash
  - merkle_root        <-- Hash de todas las transacciones
  - timestamp
  - difficulty_target
  - nonce
\end{verbatim}
\end{definition}

\begin{keyidea}
\textbf{¿Por qué Merkle trees en Bitcoin?}

\textbf{1. Verificación SPV (Simplified Payment Verification):}

Un nodo ligero puede verificar que una transacción está en un bloque sin descargar todas las transacciones:
\begin{itemize}
    \item Descarga solo headers ($\sim$80 bytes/bloque)
    \item Pide Merkle proof de la transacción específica
    \item Verifica proof contra merkle\_root del header
\end{itemize}

\textbf{2. Eficiencia:}
\begin{itemize}
    \item Header: 80 bytes (fijo)
    \item Bloque completo: $\sim$1-2 MB
    \item Proof para 1 tx en bloque de 4000 tx: $\sim$12 hashes = 384 bytes
\end{itemize}
\end{keyidea}

\begin{example}[Verificación SPV]
Alice quiere verificar que su pago está confirmado.

\textbf{Sin Merkle tree:}
\begin{enumerate}
    \item Descargar bloque completo ($\sim$1 MB)
    \item Buscar transacción
    \item Verificar
\end{enumerate}

\textbf{Con Merkle tree:}
\begin{enumerate}
    \item Pedir Merkle proof ($\sim$400 bytes)
    \item Verificar proof contra header ($\sim$80 bytes)
\end{enumerate}

Reducción: 1 MB $\to$ 500 bytes ($\sim$2000x menos datos).
\end{example}

\begin{remark}
Bitcoin usa SHA-256 doble: $H(x) = \text{SHA256}(\text{SHA256}(x))$

Esto fue una decisión de Satoshi, posiblemente para prevenir ataques de extensión de longitud (aunque SHA-256 truncado hubiera sido suficiente).
\end{remark}

%-------------------------------------------
\subsection{Merkle Patricia Trie en Ethereum}

\begin{definition}[Patricia Trie]
Un Patricia Trie (Practical Algorithm To Retrieve Information Coded In Alphanumeric) es un trie comprimido donde nodos con un solo hijo se fusionan con su padre.
\end{definition}

\begin{definition}[Merkle Patricia Trie]
Ethereum combina Patricia Tries con Merkle Trees para crear una estructura que permite:
\begin{itemize}
    \item Búsqueda eficiente por clave
    \item Pruebas de inclusión Y exclusión
    \item Compromiso a todo el estado en una raíz
\end{itemize}

Tres tipos de nodos:
\begin{enumerate}
    \item \textbf{Branch node:} 17 elementos [16 hijos + valor]
    \item \textbf{Extension node:} [prefijo compartido, siguiente nodo]
    \item \textbf{Leaf node:} [sufijo de clave, valor]
\end{enumerate}
\end{definition}

\begin{intuicion}
Mientras Bitcoin solo necesita probar ``esta transacción está en el bloque'', Ethereum necesita probar cosas más complejas:
\begin{itemize}
    \item ``La cuenta 0x1234... tiene balance X''
    \item ``El storage slot Y del contrato Z tiene valor W''
    \item ``La cuenta 0xABCD... NO existe'' (prueba de exclusión)
\end{itemize}

El Merkle Patricia Trie permite todo esto con pruebas logarítmicas.
\end{intuicion}

\begin{remark}
Ethereum tiene múltiples MPTs:
\begin{itemize}
    \item \textbf{State trie:} Mapea direcciones a cuentas (balance, nonce, codeHash, storageRoot)
    \item \textbf{Storage trie:} Para cada contrato, mapea slots a valores
    \item \textbf{Transaction trie:} Transacciones del bloque indexadas por posición
    \item \textbf{Receipt trie:} Receipts indexados por posición
\end{itemize}

El header del bloque contiene las raíces de todos estos tries.
\end{remark}

%-------------------------------------------
\subsection{Sparse Merkle Trees}

\begin{definition}[Sparse Merkle Tree (SMT)]
Un SMT es un Merkle tree de profundidad fija $d$ donde la mayoría de las hojas están vacías.

Propiedades:
\begin{itemize}
    \item $2^d$ hojas posibles, pero solo $n \ll 2^d$ tienen valores no vacíos
    \item Permite pruebas de NO inclusión
    \item Se pueden precomputar los hashes de subárboles vacíos
\end{itemize}
\end{definition}

\begin{keyidea}
\textbf{Truco de eficiencia para SMTs:}

Los hashes de subárboles vacíos son determinísticos:
\begin{align}
\text{empty}_0 &= H(\texttt{0x00}) \quad \text{(hoja vacía)} \\
\text{empty}_1 &= H(\text{empty}_0 \| \text{empty}_0) \\
\text{empty}_2 &= H(\text{empty}_1 \| \text{empty}_1) \\
&\vdots \\
\text{empty}_d &= H(\text{empty}_{d-1} \| \text{empty}_{d-1})
\end{align}

Estos valores se precomputan una vez. Un subárbol completamente vacío se representa con su hash precomputado, sin almacenar los nodos individuales.
\end{keyidea}

\begin{example}[SMT para un conjunto de 3 elementos]
Profundidad $d = 32$ (soporta $2^{32} \approx 4$ mil millones de claves).

Elementos: clave 5 $\to$ valor A, clave 1000 $\to$ valor B, clave 7 $\to$ valor C

\begin{verbatim}
Almacenamiento real:
- Camino a clave 5: ~32 nodos
- Camino a clave 1000: ~32 nodos  
- Camino a clave 7: ~32 nodos
- Nodos compartidos entre caminos

Total: << 2^32 nodos
\end{verbatim}

Los subárboles sin elementos almacenan solo una referencia al hash vacío precomputado.
\end{example}

\begin{mathnote}
\textbf{Prueba de no inclusión en SMT:}

Para probar que la clave $k$ NO está en el árbol:
\begin{enumerate}
    \item Proporcionar el camino desde la raíz hasta la posición de $k$
    \item Mostrar que la hoja en esa posición está vacía (= hash vacío)
\end{enumerate}

Esto es imposible en un Merkle tree normal sin profundidad fija.
\end{mathnote}

%-------------------------------------------
\subsection{Indexed Merkle Trees}

\begin{definition}[Indexed Merkle Tree]
Variante donde cada hoja contiene no solo un valor, sino también información sobre el siguiente elemento en orden:

\[ \text{leaf} = (k, v, \text{next\_key}, \text{next\_idx}) \]

Esto permite pruebas de no inclusión más eficientes que SMTs completos.
\end{definition}

\begin{remark}
Los Indexed Merkle Trees son usados en Aztec Protocol para su árbol de nullifiers, combinando las ventajas de:
\begin{itemize}
    \item Merkle trees normales: tamaño proporcional a elementos insertados
    \item SMTs: pruebas de no inclusión
\end{itemize}
\end{remark}

%-------------------------------------------
\subsection{Append-Only Merkle Trees}

\begin{definition}[Append-Only Merkle Tree]
Un Merkle tree donde solo se permiten inserciones al final, nunca modificaciones o eliminaciones.

Propiedades adicionales:
\begin{itemize}
    \item \textbf{Consistency proofs:} Dado root antiguo $r_1$ y nuevo $r_2$, probar que el árbol viejo es prefijo del nuevo
    \item \textbf{Tamper evidence:} Cualquier modificación histórica es detectable
\end{itemize}
\end{definition}

\begin{intuicion}
Certificate Transparency (CT) usa append-only Merkle trees para crear un log público de certificados TLS. Si una CA emite un certificado malicioso, queda registrado públicamente y es detectable.
\end{intuicion}

%-------------------------------------------
\subsection{Hash Functions para Merkle Trees}

\begin{center}
\begin{tabular}{lcccc}
\toprule
\textbf{Hash} & \textbf{Bitcoin} & \textbf{Ethereum} & \textbf{ZK SNARKs} & \textbf{Costo ZK} \\
\midrule
SHA-256 & ✓ & & & $\sim$25,000 \\
Keccak-256 & & ✓ & & $\sim$150,000 \\
Poseidon & & & ✓ & $\sim$250 \\
Pedersen & & & ✓ & $\sim$1,000 \\
MiMC & & & ✓ & $\sim$700 \\
\bottomrule
\end{tabular}
\end{center}

\begin{keyidea}
\textbf{La elección del hash importa enormemente en ZK:}

Para un Merkle tree de profundidad 20 (1M hojas):
\begin{itemize}
    \item Con SHA-256: $20 \times 25,000 = 500,000$ constraints
    \item Con Keccak: $20 \times 150,000 = 3,000,000$ constraints
    \item Con Poseidon: $20 \times 250 = 5,000$ constraints
\end{itemize}

¡El mismo árbol puede costar 100-600x más con un hash inadecuado!
\end{keyidea}

%-------------------------------------------
\subsection{Merkle Trees en Sistemas ZK}

\begin{mathnote}
Los Merkle trees son fundamentales en casi todos los protocolos ZK:

\textbf{Tornado Cash:}
\begin{itemize}
    \item Árbol de profundidad 20 con compromisos de depósitos
    \item Probar: ``Conozco el secreto de algún commitment en el árbol''
    \item Sin revelar CUÁL commitment
\end{itemize}

\textbf{Zcash:}
\begin{itemize}
    \item Árbol de notas (notes = UTXOs privados)
    \item Probar: ``Esta nota existe y puedo gastarla''
    \item Nullifier previene doble gasto sin revelar la nota
\end{itemize}

\textbf{ZK Rollups:}
\begin{itemize}
    \item Árbol de estado (cuentas y balances)
    \item Probar: ``Las transacciones actualizaron el estado correctamente''
    \item Raíz antigua $\to$ raíz nueva
\end{itemize}
\end{mathnote}

\begin{example}[Merkle proof en circuito ZK]
\textbf{Inputs públicos:}
\begin{itemize}
    \item Merkle root: $r$
\end{itemize}

\textbf{Inputs privados (witness):}
\begin{itemize}
    \item Hoja: $L$
    \item Índice: $i$
    \item Authentication path: $[s_0, s_1, \ldots, s_{d-1}]$
\end{itemize}

\textbf{Circuito:}
\begin{verbatim}
h = Hash(L)
for j in 0..d:
    bit = (i >> j) & 1
    if bit == 0:
        h = Hash(h, siblings[j])
    else:
        h = Hash(siblings[j], h)

assert h == root
\end{verbatim}

El circuito verifica el path sin revelar $L$, $i$, ni los siblings.
\end{example}

%-------------------------------------------
\subsection{Optimizaciones}

\begin{keyidea}
\textbf{1. Merkle Mountain Ranges (MMR):}

Para datos append-only, MMR permite agregar elementos en $O(1)$ amortizado (vs $O(\log n)$ para árboles balanceados).

\textbf{2. Verkle Trees:}

Usan compromisos polinomiales (KZG) en lugar de hashes. Los proofs son de tamaño constante independiente de la profundidad. Ethereum planea migrar de MPT a Verkle trees.

\textbf{3. Árboles de aridad alta:}

En lugar de árboles binarios, usar árboles con 4, 8 o 16 hijos por nodo. Reduce la profundidad pero aumenta el tamaño de cada nodo en el proof.
\end{keyidea}

\begin{mathnote}
\textbf{Trade-off de aridad:}

Para $n$ hojas con árbol de aridad $a$:
\begin{itemize}
    \item Profundidad: $\log_a(n)$
    \item Tamaño de proof: $\log_a(n) \times (a-1)$ hashes
\end{itemize}

\begin{center}
\begin{tabular}{ccc}
\toprule
\textbf{Aridad} & \textbf{Profundidad} & \textbf{Hashes en proof} \\
\midrule
2 & $\log_2(n) = 20$ & 20 \\
4 & $\log_4(n) = 10$ & 30 \\
8 & $\log_8(n) \approx 7$ & 49 \\
16 & $\log_{16}(n) = 5$ & 75 \\
\bottomrule
\end{tabular}
\end{center}

Para $n = 2^{20}$. El árbol binario es óptimo en tamaño de proof.
\end{mathnote}

%-------------------------------------------
\subsection{Ataques y Consideraciones de Seguridad}

\begin{advertencia}
\textbf{1. Second preimage attacks:}

Si el hash no distingue hojas de nodos internos, un atacante podría presentar un nodo interno como hoja.

\textbf{Solución:} Prefijos diferentes para hojas y nodos:
\begin{align}
H_{\text{leaf}}(x) &= H(\texttt{0x00} \| x) \\
H_{\text{node}}(l, r) &= H(\texttt{0x01} \| l \| r)
\end{align}

\textbf{2. Length extension en Merkle roots:}

Si se usa un hash vulnerable a length extension (SHA-256 raw), cuidado con cómo se usa la raíz.

\textbf{3. Índice fuera de rango:}

Verificar que el índice proporcionado está en rango válido.
\end{advertencia}

%-------------------------------------------
\subsection{Implementación}

\begin{example}[Merkle Tree en Python]
\begin{verbatim}
import hashlib

def hash_node(left, right):
    """Hash two nodes together"""
    return hashlib.sha256(b'\x01' + left + right).digest()

def hash_leaf(data):
    """Hash a leaf node"""
    return hashlib.sha256(b'\x00' + data).digest()

class MerkleTree:
    def __init__(self, leaves):
        self.leaves = [hash_leaf(leaf) for leaf in leaves]
        self.layers = [self.leaves]
        self._build()
    
    def _build(self):
        """Build tree from leaves up"""
        current = self.leaves
        while len(current) > 1:
            next_layer = []
            for i in range(0, len(current), 2):
                left = current[i]
                right = current[i+1] if i+1 < len(current) else left
                next_layer.append(hash_node(left, right))
            self.layers.append(next_layer)
            current = next_layer
    
    @property
    def root(self):
        return self.layers[-1][0]
    
    def get_proof(self, index):
        """Get authentication path for leaf at index"""
        proof = []
        for layer in self.layers[:-1]:
            sibling_idx = index ^ 1  # Flip last bit
            if sibling_idx < len(layer):
                sibling = layer[sibling_idx]
            else:
                sibling = layer[index]
            direction = 'L' if index % 2 == 1 else 'R'
            proof.append((sibling, direction))
            index //= 2
        return proof

def verify_proof(leaf, proof, root):
    """Verify a Merkle proof"""
    current = hash_leaf(leaf)
    for sibling, direction in proof:
        if direction == 'L':
            current = hash_node(sibling, current)
        else:
            current = hash_node(current, sibling)
    return current == root

# Uso
leaves = [b'tx0', b'tx1', b'tx2', b'tx3']
tree = MerkleTree(leaves)
proof = tree.get_proof(1)  # Proof para tx1
assert verify_proof(b'tx1', proof, tree.root)
\end{verbatim}
\end{example}

\begin{example}[Merkle Tree en Solidity]
\begin{verbatim}
// SPDX-License-Identifier: MIT
pragma solidity ^0.8.0;

library MerkleProof {
    function verify(
        bytes32[] memory proof,
        bytes32 root,
        bytes32 leaf,
        uint256 index
    ) internal pure returns (bool) {
        bytes32 computedHash = leaf;
        
        for (uint256 i = 0; i < proof.length; i++) {
            bytes32 proofElement = proof[i];
            
            if (index % 2 == 0) {
                // Sibling is on the right
                computedHash = keccak256(
                    abi.encodePacked(computedHash, proofElement)
                );
            } else {
                // Sibling is on the left
                computedHash = keccak256(
                    abi.encodePacked(proofElement, computedHash)
                );
            }
            
            index = index / 2;
        }
        
        return computedHash == root;
    }
}

contract MerkleAirdrop {
    bytes32 public merkleRoot;
    mapping(address => bool) public claimed;
    
    constructor(bytes32 _merkleRoot) {
        merkleRoot = _merkleRoot;
    }
    
    function claim(
        uint256 index,
        address account,
        uint256 amount,
        bytes32[] calldata proof
    ) external {
        require(!claimed[account], "Already claimed");
        
        bytes32 leaf = keccak256(
            abi.encodePacked(index, account, amount)
        );
        require(
            MerkleProof.verify(proof, merkleRoot, leaf, index),
            "Invalid proof"
        );
        
        claimed[account] = true;
        // Transfer tokens...
    }
}
\end{verbatim}
\end{example}

%-------------------------------------------
\subsection{Aplicación en Banco ZK}

\begin{bancozk}
Los Merkle trees son el corazón de la arquitectura de Banco ZK:

\textbf{1. Árbol de Notas (Note Tree)}

Cada depósito crea un ``commitment'' que se inserta en el árbol:
\[ \text{commitment} = \text{Poseidon}(\text{value}, \text{pubkey}, \text{randomness}) \]

\begin{center}
\begin{tabular}{ll}
\toprule
\textbf{Parámetro} & \textbf{Valor} \\
\midrule
Profundidad & 20 \\
Capacidad & $2^{20} \approx 1$ millón de notas \\
Hash & Poseidon \\
Tipo & Append-only \\
\bottomrule
\end{tabular}
\end{center}

\textbf{2. Árbol de Nullifiers (Nullifier Tree)}

Para prevenir doble gasto sin revelar qué nota se gastó:
\[ \text{nullifier} = \text{Poseidon}(\text{commitment}, \text{spending\_key}) \]

\begin{center}
\begin{tabular}{ll}
\toprule
\textbf{Parámetro} & \textbf{Valor} \\
\midrule
Profundidad & 32 \\
Tipo & Sparse (indexed) \\
Función & Detectar doble gasto \\
\bottomrule
\end{tabular}
\end{center}

\textbf{3. Flujo de transacción privada}

\begin{verbatim}
DEPOSITAR:
1. Usuario genera: commitment = Poseidon(value, pk, r)
2. Contrato inserta commitment en Note Tree
3. Usuario guarda (value, pk, r) localmente

TRANSFERIR (privadamente):
1. Usuario prueba en ZK:
   - "Conozco una nota en el Note Tree" (Merkle proof)
   - "El nullifier de esta nota no está usado"
   - "Estoy creando nueva(s) nota(s) válida(s)"
2. Publica: nullifier, nuevos commitments, prueba ZK
3. Contrato:
   - Verifica prueba
   - Añade nullifier al Nullifier Tree
   - Añade nuevos commitments al Note Tree

RETIRAR:
1. Similar a transferir, pero el output es público
2. Contrato envía fondos a la dirección especificada
\end{verbatim}

\textbf{4. Costo en circuitos}

\begin{center}
\begin{tabular}{lcc}
\toprule
\textbf{Componente} & \textbf{Hash} & \textbf{Constraints} \\
\midrule
Note tree proof (depth 20) & Poseidon & $20 \times 250 = 5,000$ \\
Nullifier computation & Poseidon & 250 \\
Commitment computation & Poseidon & 250 \\
Output commitments (×2) & Poseidon & 500 \\
\midrule
\textbf{Total hashing} & & $\sim$6,000 \\
\bottomrule
\end{tabular}
\end{center}

Comparado con usar SHA-256 ($\sim$500,000$), Poseidon reduce el costo 80x.

\textbf{5. On-chain storage}

El contrato almacena:
\begin{itemize}
    \item Note Tree root: 32 bytes (actualizado por batch)
    \item Nullifier Tree root: 32 bytes (actualizado por batch)
    \item Historical roots: Array para verificar pruebas antiguas
\end{itemize}

\begin{verbatim}
contract BancoZK {
    bytes32 public noteTreeRoot;
    bytes32 public nullifierTreeRoot;
    
    // Permitir pruebas contra roots recientes
    bytes32[100] public historicalRoots;
    uint256 public rootIndex;
    
    function isKnownRoot(bytes32 root) public view returns (bool) {
        if (root == noteTreeRoot) return true;
        for (uint i = 0; i < 100; i++) {
            if (historicalRoots[i] == root) return true;
        }
        return false;
    }
}
\end{verbatim}

\textbf{6. Actualización de árboles}

Las actualizaciones se procesan en batches:
\begin{enumerate}
    \item Sequencer colecta transacciones
    \item Genera prueba ZK de transición de estado válida
    \item Publica: nueva raíz + prueba
    \item Contrato verifica prueba y actualiza raíz
\end{enumerate}

\textbf{7. Consideraciones de privacidad}

Para máxima privacidad:
\begin{itemize}
    \item El árbol tiene hojas aleatorias (no todas son depósitos reales)
    \item Las inserciones son en posiciones semi-aleatorias
    \item El timing de operaciones se ofusca donde es posible
\end{itemize}
\end{bancozk}

%-------------------------------------------
\subsection{Resumen}

\begin{keyidea}
\textbf{Merkle Trees: Puntos clave}

\begin{enumerate}
    \item \textbf{Compromiso eficiente:} Una raíz de 32 bytes compromete millones de elementos
    
    \item \textbf{Proofs logarítmicos:} $O(\log n)$ hashes para probar inclusión
    
    \item \textbf{Variantes para diferentes usos:}
    \begin{itemize}
        \item Normal: inclusión
        \item Sparse: inclusión y exclusión
        \item Indexed: eficiencia intermedia
        \item Append-only: logs y auditoría
    \end{itemize}
    
    \item \textbf{Hash matters:} En ZK, elegir Poseidon sobre SHA-256 puede ahorrar 100x en constraints
    
    \item \textbf{Ubicuidad:} Bitcoin, Ethereum, Zcash, rollups, airdrops, certificate transparency...
\end{enumerate}
\end{keyidea}
% ----- MERKLE PATRICIA TRIES -----
%===========================================
% SECCIÓN: MERKLE PATRICIA TRIES - VERSIÓN EXPANDIDA
% Usa tus entornos: definition, example, keyidea, bancozk, 
%                   mathnote, intuicion, advertencia, remark
%===========================================

\section{Merkle Patricia Tries}

\begin{definition}[Patricia Trie]
Trie comprimido donde caminos sin bifurcación se colapsan.
\end{definition}

\begin{definition}[Merkle Patricia Trie]
Patricia Trie donde cada nodo está hasheado (Merkleizado).
\end{definition}

Usado en Ethereum para:
\begin{itemize}
    \item State trie: cuentas y storage
    \item Transaction trie: transacciones por bloque
    \item Receipt trie: logs y resultados
\end{itemize}

Tipos de nodos:
\begin{itemize}
    \item \textbf{Branch:} 16 hijos (hex) + valor opcional
    \item \textbf{Extension:} prefijo compartido + siguiente nodo
    \item \textbf{Leaf:} resto del path + valor
\end{itemize}

Ventajas:
\begin{itemize}
    \item Pruebas de inclusión Y exclusión
    \item Actualizaciones eficientes
    \item State root en header de bloque
\end{itemize}

%-------------------------------------------
\subsection{Contexto: ¿Por Qué No un Merkle Tree Simple?}

\begin{intuicion}
Bitcoin usa Merkle trees simples para transacciones: funciona porque solo necesita probar ``esta transacción está en el bloque''. Pero Ethereum necesita algo más poderoso.

Ethereum es una \textbf{máquina de estados}: cada bloque transiciona el estado global (balances, storage de contratos, código). Necesitamos:
\begin{enumerate}
    \item Probar que una cuenta tiene cierto balance
    \item Probar que una variable de contrato tiene cierto valor
    \item Probar que una cuenta \textbf{NO existe} (prueba de exclusión)
    \item Actualizar eficientemente cuando el estado cambia
\end{enumerate}

Un Merkle tree simple no puede hacer el punto 3 (exclusión) ni el 4 eficientemente.
\end{intuicion}

\begin{keyidea}
El Merkle Patricia Trie combina tres conceptos:
\begin{enumerate}
    \item \textbf{Trie}: Estructura de árbol indexada por claves (permite búsqueda)
    \item \textbf{Patricia}: Compresión de caminos sin bifurcación (eficiencia espacial)
    \item \textbf{Merkle}: Cada nodo está hasheado (compromiso criptográfico)
\end{enumerate}

Resultado: Una estructura que soporta búsquedas, actualizaciones, pruebas de inclusión Y exclusión, todo con compromiso criptográfico.
\end{keyidea}

%-------------------------------------------
\subsection{Fundamentos: El Trie Básico}

\begin{definition}[Trie (Prefix Tree)]
Un trie es un árbol donde:
\begin{itemize}
    \item Cada arista está etiquetada con un carácter
    \item La clave de un nodo es la concatenación de etiquetas desde la raíz
    \item Los valores se almacenan en los nodos
\end{itemize}

Para claves que son strings sobre alfabeto $\Sigma$, cada nodo tiene hasta $|\Sigma|$ hijos.
\end{definition}

\begin{example}[Trie simple con palabras]
Almacenando: ``cat'' $\to$ 1, ``car'' $\to$ 2, ``card'' $\to$ 3, ``dog'' $\to$ 4

\begin{verbatim}
         (root)
        /      \
       c        d
       |        |
       a        o
      / \       |
     t   r      g
     |   |      |
    [1]  |     [4]
         |
        [2]
         |
         d
         |
        [3]
\end{verbatim}

Para buscar ``car'': seguir c $\to$ a $\to$ r, encontrar valor 2.
\end{example}

\begin{mathnote}
\textbf{Complejidad del trie básico:}
\begin{itemize}
    \item Búsqueda: $O(|k|)$ donde $|k|$ es la longitud de la clave
    \item Inserción: $O(|k|)$
    \item Espacio: $O(n \cdot |\Sigma| \cdot |k|)$ en el peor caso
\end{itemize}

El problema: con alfabeto grande ($|\Sigma| = 256$ para bytes), cada nodo usa mucha memoria aunque tenga pocos hijos.
\end{mathnote}

%-------------------------------------------
\subsection{Patricia Trie: Compresión de Caminos}

\begin{definition}[Patricia Trie]
PATRICIA = Practical Algorithm To Retrieve Information Coded In Alphanumeric

Un Patricia trie comprime cadenas de nodos con un solo hijo en un solo nodo con una etiqueta multi-carácter.
\end{definition}

\begin{example}[Compresión Patricia]
El trie anterior comprimido:

\begin{verbatim}
Antes (trie normal):          Después (Patricia):
      (root)                        (root)
     /      \                      /      \
    c        d                   "ca"    "dog"
    |        |                   / \       |
    a        o                 "t" "r"    [4]
   / \       |                  |   |
  t   r      g                 [1] [2]
  |   |      |                      |
 [1] [2]    [4]                    "d"
      |                             |
      d                            [3]
      |
     [3]
\end{verbatim}

Nodos reducidos de 8 a 6. Con claves más largas, la compresión es mayor.
\end{example}

\begin{keyidea}
La compresión Patricia es especialmente efectiva cuando:
\begin{itemize}
    \item Las claves son largas (como hashes de 32 bytes)
    \item Las claves comparten prefijos largos
    \item El conjunto de claves es disperso (sparse)
\end{itemize}

En Ethereum, las claves son hashes Keccak de 32 bytes. Sin compresión, el árbol tendría profundidad 64 (nibbles hexadecimales). Con Patricia, la profundidad promedio es mucho menor.
\end{keyidea}

%-------------------------------------------
\subsection{Merkle Patricia Trie: Añadiendo Hashes}

\begin{definition}[Merkleización]
Para convertir un Patricia trie en un Merkle Patricia Trie:
\begin{enumerate}
    \item Serializar cada nodo usando RLP (Recursive Length Prefix)
    \item Si el RLP del nodo tiene $\geq 32$ bytes, almacenar su hash Keccak-256
    \item Si tiene $< 32$ bytes, almacenar el nodo directamente (inline)
    \item La raíz del trie es el hash del nodo raíz
\end{enumerate}
\end{definition}

\begin{intuicion}
El hash de la raíz ``compromete'' a todo el contenido del trie. Cualquier cambio en cualquier valor propagará cambios de hash hasta la raíz.

Esto permite incluir un solo hash de 32 bytes en el header del bloque que compromete todo el estado de Ethereum ($\sim$100 GB de datos).
\end{intuicion}

%-------------------------------------------
\subsection{Tipos de Nodos en Detalle}

\begin{definition}[Nodo Branch]
Un nodo branch tiene 17 elementos:
\[ [\text{child}_0, \text{child}_1, \ldots, \text{child}_{15}, \text{value}] \]

\begin{itemize}
    \item $\text{child}_i$: referencia al hijo para el nibble $i$ (0-9, a-f)
    \item $\text{value}$: valor almacenado si alguna clave termina en este nodo
\end{itemize}

Cada $\text{child}_i$ es:
\begin{itemize}
    \item Vacío (null) si no hay hijo
    \item Hash del hijo si el hijo serializado $\geq 32$ bytes
    \item El hijo serializado directamente (inline) si $< 32$ bytes
\end{itemize}
\end{definition}

\begin{definition}[Nodo Extension]
Un nodo extension tiene 2 elementos:
\[ [\text{encoded\_path}, \text{next\_node}] \]

\begin{itemize}
    \item $\text{encoded\_path}$: prefijo compartido comprimido (con hex-prefix encoding)
    \item $\text{next\_node}$: referencia al siguiente nodo (hash o inline)
\end{itemize}

Se usa cuando múltiples claves comparten un prefijo largo.
\end{definition}

\begin{definition}[Nodo Leaf]
Un nodo leaf tiene 2 elementos:
\[ [\text{encoded\_path}, \text{value}] \]

\begin{itemize}
    \item $\text{encoded\_path}$: resto del path desde el nodo padre (con hex-prefix)
    \item $\text{value}$: el valor almacenado
\end{itemize}
\end{definition}

\begin{center}
\begin{tikzpicture}[
    node distance=1.5cm,
    box/.style={draw, rectangle, minimum width=2cm, minimum height=0.8cm}
]
    % Branch node
    \node[box, fill=blue!20] (branch) at (0, 0) {Branch};
    \node[below=0.2cm of branch, font=\small] {[c₀, c₁, ..., c₁₅, val]};
    
    % Extension node
    \node[box, fill=green!20] (ext) at (4, 0) {Extension};
    \node[below=0.2cm of ext, font=\small] {[path, next]};
    
    % Leaf node
    \node[box, fill=yellow!20] (leaf) at (8, 0) {Leaf};
    \node[below=0.2cm of leaf, font=\small] {[path, value]};
    
    % Descriptions
    \node[below=1cm of branch, font=\footnotesize, align=center] {16 hijos\\+ valor opcional};
    \node[below=1cm of ext, font=\footnotesize, align=center] {Prefijo compartido\\comprimido};
    \node[below=1cm of leaf, font=\footnotesize, align=center] {Resto de clave\\+ valor final};
\end{tikzpicture}
\end{center}

%-------------------------------------------
\subsection{Hex-Prefix Encoding}

\begin{definition}[Hex-Prefix Encoding]
Para distinguir nodos Extension de Leaf y manejar paths de longitud impar, se usa un prefijo especial:

\begin{center}
\begin{tabular}{ccc}
\toprule
\textbf{Tipo de nodo} & \textbf{Longitud path} & \textbf{Prefijo} \\
\midrule
Extension & Par & 0x00 \\
Extension & Impar & 0x1 \\
Leaf & Par & 0x20 \\
Leaf & Impar & 0x3 \\
\bottomrule
\end{tabular}
\end{center}

Si la longitud es impar, el primer nibble del path se combina con el prefijo.
\end{definition}

\begin{example}[Hex-Prefix encoding]
\textbf{Path = [1, 2, 3, 4] (par), tipo Leaf:}
\begin{verbatim}
Prefijo: 0x20
Encoded: [0x20, 0x12, 0x34]
\end{verbatim}

\textbf{Path = [1, 2, 3] (impar), tipo Leaf:}
\begin{verbatim}
Prefijo: 0x3, combinado con primer nibble
Encoded: [0x31, 0x23]
\end{verbatim}

\textbf{Path = [a, b, c, d] (par), tipo Extension:}
\begin{verbatim}
Prefijo: 0x00
Encoded: [0x00, 0xab, 0xcd]
\end{verbatim}
\end{example}

%-------------------------------------------
\subsection{RLP Encoding}

\begin{definition}[Recursive Length Prefix (RLP)]
RLP es el formato de serialización usado en Ethereum para codificar datos estructurados.

\textbf{Reglas básicas:}
\begin{itemize}
    \item Byte simple $[0x00, 0x7f]$: se codifica como sí mismo
    \item String corto ($\leq 55$ bytes): $[0x80 + \text{len}] + \text{data}$
    \item String largo ($> 55$ bytes): $[0xb7 + \text{len\_of\_len}] + \text{len} + \text{data}$
    \item Lista corta: $[0xc0 + \text{len}] + \text{items}$
    \item Lista larga: $[0xf7 + \text{len\_of\_len}] + \text{len} + \text{items}$
\end{itemize}
\end{definition}

\begin{example}[RLP de nodo Leaf]
\begin{verbatim}
Leaf node:
  path = [0x20, 0x12, 0x34]  (3 bytes)
  value = "hello"           (5 bytes)

RLP([path, value]):
  path_rlp = [0x83, 0x20, 0x12, 0x34]  (0x83 = 0x80 + 3)
  value_rlp = [0x85] + "hello"         (0x85 = 0x80 + 5)
  
  list_content = path_rlp + value_rlp = 12 bytes
  list_prefix = 0xc0 + 12 = 0xcc
  
Final RLP: [0xcc, 0x83, 0x20, 0x12, 0x34, 0x85, 'h', 'e', 'l', 'l', 'o']
\end{verbatim}

Como tiene 12 bytes (< 32), este nodo puede ser inline en su padre.
\end{example}

%-------------------------------------------
\subsection{Ejemplo Completo de MPT}

\begin{example}[Construcción de MPT]
Almacenemos estos pares clave-valor (claves en hex):

\begin{center}
\begin{tabular}{ll}
\toprule
\textbf{Clave} & \textbf{Valor} \\
\midrule
\texttt{0x1234} & ``Alice'' \\
\texttt{0x1235} & ``Bob'' \\
\texttt{0x5678} & ``Carol'' \\
\bottomrule
\end{tabular}
\end{center}

\textbf{Paso 1: Convertir claves a nibbles}
\begin{verbatim}
0x1234 → [1, 2, 3, 4]
0x1235 → [1, 2, 3, 5]
0x5678 → [5, 6, 7, 8]
\end{verbatim}

\textbf{Paso 2: Construir trie}

\begin{verbatim}
                    [Branch]
                   /        \
              [1]            [5]
               |              |
          [Extension]     [Leaf]
          path: [2,3]    path: [6,7,8]
               |         value: "Carol"
            [Branch]
           /        \
         [4]        [5]
          |          |
       [Leaf]     [Leaf]
       path: []   path: []
       val:"Alice" val:"Bob"
\end{verbatim}

\textbf{Paso 3: Merkleizar (bottom-up)}

\begin{verbatim}
1. Leaf "Alice": RLP([hp([]), "Alice"]) → hash₁
2. Leaf "Bob":   RLP([hp([]), "Bob"])   → hash₂
3. Branch en [1,2,3]:
   children[4] = hash₁
   children[5] = hash₂
   RLP → hash₃
4. Extension [2,3]:
   RLP([hp([2,3], ext), hash₃]) → hash₄
5. Leaf "Carol":
   RLP([hp([6,7,8], leaf), "Carol"]) → hash₅
6. Root Branch:
   children[1] = hash₄
   children[5] = hash₅
   RLP → ROOT
\end{verbatim}
\end{example}

%-------------------------------------------
\subsection{Operaciones}

\subsubsection{Búsqueda (Get)}

\begin{definition}[Algoritmo de Búsqueda]
\begin{verbatim}
function get(node, path):
    if path is empty:
        if node is Leaf or Branch with value:
            return node.value
        else:
            return NOT_FOUND
    
    if node is Branch:
        child = node.children[path[0]]
        if child is null:
            return NOT_FOUND
        return get(resolve(child), path[1:])
    
    if node is Extension:
        if path starts with node.path:
            return get(resolve(node.next), path[len(node.path):])
        else:
            return NOT_FOUND
    
    if node is Leaf:
        if path == node.path:
            return node.value
        else:
            return NOT_FOUND
\end{verbatim}
\end{definition}

\subsubsection{Inserción (Put)}

\begin{mathnote}
La inserción puede requerir:
\begin{enumerate}
    \item Crear nuevo Leaf (clave no existe, sin conflictos)
    \item Dividir un Extension (nuevo prefijo compartido más corto)
    \item Dividir un Leaf (claves comparten prefijo parcial)
    \item Actualizar valor existente (clave ya existe)
\end{enumerate}

Cada caso requiere recalcular hashes hasta la raíz.
\end{mathnote}

\subsubsection{Eliminación (Delete)}

\begin{advertencia}
La eliminación es la operación más compleja porque puede requerir:
\begin{enumerate}
    \item Eliminar Leaf y posiblemente colapsar Branch padre
    \item Fusionar Extension con hijo si queda un solo camino
    \item Actualizar hashes hasta la raíz
\end{enumerate}

La complejidad de mantener el trie ``comprimido'' después de deletes hace que la implementación sea propensa a bugs.
\end{advertencia}

%-------------------------------------------
\subsection{Pruebas de Inclusión y Exclusión}

\begin{keyidea}
A diferencia de los Merkle trees simples, los MPT permiten \textbf{pruebas de exclusión}: demostrar que una clave NO está en el trie.
\end{keyidea}

\begin{definition}[Prueba de Inclusión]
Para probar que la clave $k$ tiene valor $v$:
\begin{enumerate}
    \item Proporcionar todos los nodos en el camino desde la raíz hasta el Leaf de $k$
    \item El verificador reconstruye los hashes y verifica que la raíz coincide
\end{enumerate}
\end{definition}

\begin{definition}[Prueba de Exclusión]
Para probar que la clave $k$ NO está en el trie:
\begin{enumerate}
    \item Proporcionar el camino hasta donde la búsqueda ``falla''
    \item Mostrar que:
    \begin{itemize}
        \item Un Branch no tiene hijo para el siguiente nibble, o
        \item Un Extension tiene un prefijo que no coincide, o
        \item Un Leaf tiene un path diferente
    \end{itemize}
\end{enumerate}
\end{definition}

\begin{example}[Prueba de exclusión]
En el MPT anterior, probar que \texttt{0x1236} NO existe:

\begin{verbatim}
Camino de búsqueda:
1. Root Branch → child[1] existe → seguir
2. Extension [2,3] → coincide → seguir
3. Branch → child[6] = null → ¡NO EXISTE!

Proof:
- Root Branch (mostrando child[1] = hash₄)
- Extension con path [2,3]
- Branch (mostrando child[6] = null)

Verificador:
1. Verifica que los hashes son consistentes con la raíz conocida
2. Observa que child[6] está vacío
3. Concluye: 0x1236 no puede existir
\end{verbatim}
\end{example}

%-------------------------------------------
\subsection{Ethereum: Los Cuatro Tries}

\begin{definition}[State Trie]
Mapea direcciones (20 bytes) a cuentas (Account objects).

\textbf{Clave:} $\text{Keccak256}(\text{address})$
\textbf{Valor:} $\text{RLP}([\text{nonce}, \text{balance}, \text{storageRoot}, \text{codeHash}])$

\begin{verbatim}
Account:
  nonce: número de transacciones enviadas
  balance: wei poseídos
  storageRoot: raíz del Storage Trie de la cuenta
  codeHash: Keccak256(código) para contratos, 
            Keccak256("") para EOAs
\end{verbatim}
\end{definition}

\begin{definition}[Storage Trie]
Cada contrato tiene su propio Storage Trie.

\textbf{Clave:} $\text{Keccak256}(\text{slot})$ donde slot es de 32 bytes
\textbf{Valor:} RLP del valor (32 bytes)

Los slots se asignan según las reglas del compilador Solidity:
\begin{itemize}
    \item Variables simples: slots consecutivos 0, 1, 2, ...
    \item Mappings: $\text{Keccak256}(\text{key} \| \text{slot})$
    \item Arrays dinámicos: elementos en $\text{Keccak256}(\text{slot}) + i$
\end{itemize}
\end{definition}

\begin{definition}[Transaction Trie]
Contiene las transacciones de un bloque.

\textbf{Clave:} RLP(índice de transacción): 0, 1, 2, ...
\textbf{Valor:} RLP(transacción)

Raíz incluida en block header como \texttt{transactionsRoot}.
\end{definition}

\begin{definition}[Receipt Trie]
Contiene los receipts de transacciones.

\textbf{Clave:} RLP(índice de transacción)
\textbf{Valor:} RLP(receipt) con status, gasUsed, logs, etc.

Raíz incluida en block header como \texttt{receiptsRoot}.
\end{definition}

\begin{center}
\begin{tikzpicture}[
    box/.style={draw, rectangle, minimum width=3cm, minimum height=1cm, align=center},
    trie/.style={draw, ellipse, minimum width=2cm}
]
    % Block header
    \node[box, fill=gray!20] (header) at (0, 3) {Block Header};
    
    % Roots
    \node[below=0.5cm of header, font=\small] (roots) {
        stateRoot, transactionsRoot, receiptsRoot
    };
    
    % Tries
    \node[trie, fill=blue!20] (state) at (-4, 0) {State Trie};
    \node[trie, fill=green!20] (tx) at (0, 0) {Tx Trie};
    \node[trie, fill=yellow!20] (receipt) at (4, 0) {Receipt Trie};
    
    % Storage tries
    \node[trie, fill=blue!10] (storage1) at (-5, -2) {Storage};
    \node[trie, fill=blue!10] (storage2) at (-3, -2) {Storage};
    
    % Connections
    \draw[->] (header) -- (-4, 2);
    \draw[->] (header) -- (0, 2);
    \draw[->] (header) -- (4, 2);
    \draw[->] (state) -- (storage1);
    \draw[->] (state) -- (storage2);
    
    % Labels
    \node[below=0.1cm of storage1, font=\tiny] {Contrato A};
    \node[below=0.1cm of storage2, font=\tiny] {Contrato B};
\end{tikzpicture}
\end{center}

%-------------------------------------------
\subsection{Tamaño del Estado de Ethereum}

\begin{mathnote}
\textbf{Estado de Ethereum (circa 2024):}

\begin{center}
\begin{tabular}{lr}
\toprule
\textbf{Métrica} & \textbf{Valor aproximado} \\
\midrule
Cuentas únicas & $\sim$250 millones \\
Contratos & $\sim$50 millones \\
Tamaño del state trie & $\sim$100 GB \\
Nodos en el trie & $\sim$500 millones \\
Profundidad promedio & $\sim$8 niveles \\
Profundidad máxima & 64 nibbles (32 bytes) \\
\bottomrule
\end{tabular}
\end{center}

El state trie es la principal razón por la que los nodos completos de Ethereum requieren $>$ 1 TB de almacenamiento.
\end{mathnote}

%-------------------------------------------
\subsection{Problemas y Limitaciones}

\begin{advertencia}
\textbf{1. State bloat:}

El estado de Ethereum crece constantemente. Cada cuenta nueva, cada variable de storage, añade al trie permanentemente. No hay ``garbage collection'' automático.

\textbf{2. Proof size:}

Una prueba de inclusión/exclusión puede requerir varios KB de datos (múltiples nodos del camino).

\textbf{3. Witness size para statelessness:}

Para ``stateless clients'' (clientes que no almacenan el estado), cada transacción necesita incluir pruebas de todos los accesos al estado. Esto puede ser 100+ KB por transacción.

\textbf{4. Complejidad de implementación:}

La combinación de Patricia compression + hex-prefix + RLP + Merkleization hace que las implementaciones sean complejas y propensas a bugs.
\end{advertencia}

%-------------------------------------------
\subsection{Verkle Trees: El Futuro}

\begin{definition}[Verkle Tree]
Verkle trees (Vector commitment Merkle trees) reemplazan hashes con \textbf{compromisos polinomiales} (usualmente KZG).

Ventajas sobre MPT:
\begin{itemize}
    \item Proofs de tamaño constante ($\sim$150 bytes vs $\sim$4 KB)
    \item Agregación de proofs
    \item Mejor para stateless clients
\end{itemize}
\end{definition}

\begin{remark}
Ethereum tiene planeado migrar de MPT a Verkle trees en una actualización futura (probablemente post-2025). Esto:
\begin{enumerate}
    \item Reducirá drásticamente el tamaño de los witnesses
    \item Habilitará stateless clients prácticos
    \item Requerirá una migración masiva del estado existente
\end{enumerate}
\end{remark}

%-------------------------------------------
\subsection{Implementación}

\begin{example}[MPT simplificado en Python]
\begin{verbatim}
from typing import Optional, List, Union
import hashlib

def keccak256(data: bytes) -> bytes:
    return hashlib.sha3_256(data).digest()

class MPT:
    def __init__(self):
        self.db = {}  # hash -> node
        self.root = None
    
    def _hash_node(self, node: bytes) -> bytes:
        if len(node) < 32:
            return node  # Inline
        h = keccak256(node)
        self.db[h] = node
        return h
    
    def _nibbles(self, key: bytes) -> List[int]:
        """Convert bytes to nibbles (hex digits)"""
        result = []
        for byte in key:
            result.append(byte >> 4)
            result.append(byte & 0x0f)
        return result
    
    def get(self, key: bytes) -> Optional[bytes]:
        """Get value for key, or None if not found"""
        if self.root is None:
            return None
        nibbles = self._nibbles(keccak256(key))
        return self._get(self.root, nibbles)
    
    def _get(self, node_ref, path: List[int]) -> Optional[bytes]:
        node = self._resolve(node_ref)
        if node is None:
            return None
        
        node_type, *data = self._decode_node(node)
        
        if node_type == 'branch':
            if not path:
                return data[16]  # Value at branch
            child = data[path[0]]
            if child is None:
                return None
            return self._get(child, path[1:])
        
        elif node_type == 'extension':
            ext_path, next_node = data
            if path[:len(ext_path)] != ext_path:
                return None
            return self._get(next_node, path[len(ext_path):])
        
        elif node_type == 'leaf':
            leaf_path, value = data
            if path == leaf_path:
                return value
            return None
    
    def _resolve(self, ref):
        """Resolve hash reference to node data"""
        if ref is None:
            return None
        if len(ref) < 32:
            return ref  # Inline node
        return self.db.get(ref)
    
    # put() and delete() omitted for brevity - they're complex!
\end{verbatim}
\end{example}

\begin{example}[Verificación de prueba en Solidity]
\begin{verbatim}
// SPDX-License-Identifier: MIT
pragma solidity ^0.8.0;

library MPTVerifier {
    // Verifica una prueba de inclusión en un MPT
    function verifyInclusionProof(
        bytes32 root,
        bytes memory key,
        bytes memory value,
        bytes[] memory proof
    ) internal pure returns (bool) {
        bytes32 keyHash = keccak256(key);
        bytes memory path = _toNibbles(keyHash);
        
        bytes32 currentHash = root;
        uint256 pathIndex = 0;
        
        for (uint256 i = 0; i < proof.length; i++) {
            bytes memory node = proof[i];
            
            // Verify current node hash
            if (keccak256(node) != currentHash && 
                node.length >= 32) {
                return false;
            }
            
            // Decode and traverse node
            // (simplified - real implementation needs RLP decoding)
            // ...
        }
        
        return true;
    }
    
    function _toNibbles(bytes32 data) 
        internal pure returns (bytes memory) 
    {
        bytes memory result = new bytes(64);
        for (uint256 i = 0; i < 32; i++) {
            result[i*2] = bytes1(uint8(data[i]) >> 4);
            result[i*2+1] = bytes1(uint8(data[i]) & 0x0f);
        }
        return result;
    }
}
\end{verbatim}
\end{example}

%-------------------------------------------
\subsection{Aplicación en Banco ZK}

\begin{bancozk}
Banco ZK interactúa con MPT de Ethereum para verificar estado L1:

\textbf{1. Verificación de depósitos}

Cuando un usuario deposita en el contrato L1:
\begin{enumerate}
    \item El evento Deposit se emite y va al Receipt Trie
    \item El sequencer de L2 verifica el receipt proof
    \item Una vez verificado, acredita fondos en L2
\end{enumerate}

\begin{verbatim}
// Verificar que un depósito ocurrió
function verifyDeposit(
    bytes32 receiptRoot,    // Del block header
    bytes[] calldata proof, // MPT proof
    uint256 txIndex,
    address depositor,
    uint256 amount
) internal returns (bool) {
    // Decodificar receipt y verificar evento Deposit
    // ...
}
\end{verbatim}

\textbf{2. Verificación de estado para bridges}

Para verificar balances o storage de contratos L1:
\begin{verbatim}
// Probar que una cuenta tiene cierto balance
struct AccountProof {
    bytes[] accountProof;  // Proof en state trie
    bytes32 stateRoot;     // Del block header
}

function verifyBalance(
    AccountProof calldata proof,
    address account,
    uint256 expectedBalance
) internal pure returns (bool) {
    // Verificar MPT proof
    // Decodificar Account struct
    // Comparar balance
}
\end{verbatim}

\textbf{3. Por qué NO usamos MPT internamente}

Para el estado interno de L2, usamos Sparse Merkle Trees con Poseidon en lugar de MPT:

\begin{center}
\begin{tabular}{lcc}
\toprule
\textbf{Aspecto} & \textbf{MPT (Ethereum)} & \textbf{SMT (Banco ZK)} \\
\midrule
Hash & Keccak-256 & Poseidon \\
Costo en ZK & $\sim$150K constraints/hash & $\sim$250 constraints/hash \\
Prueba típica & $\sim$4 KB & $\sim$1 KB \\
Complejidad & Alta (3 tipos nodos) & Baja (solo branch/leaf) \\
\bottomrule
\end{tabular}
\end{center}

\textbf{4. Interoperabilidad L1 $\leftrightarrow$ L2}

\begin{verbatim}
L1 (Ethereum)                    L2 (Banco ZK)
┌──────────────┐                ┌──────────────┐
│  MPT State   │                │  SMT State   │
│  - Keccak    │                │  - Poseidon  │
│  - RLP       │                │  - Simple    │
└──────────────┘                └──────────────┘
       │                               │
       └──────── Bridge ───────────────┘
                  │
         MPT proofs verificados
         en circuito ZK (costoso)
\end{verbatim}

Verificar MPT proofs en ZK es costoso ($\sim$150K constraints por hash Keccak), pero solo lo hacemos para:
\begin{itemize}
    \item Verificar depósitos desde L1
    \item Sincronizar con estado L1 cuando es necesario
\end{itemize}

El estado interno de L2 usa estructuras ZK-friendly.

\textbf{5. Light client verification}

Para que usuarios verifiquen el estado de L2 sin confiar en el sequencer:
\begin{enumerate}
    \item Publicamos state root de L2 en contrato L1
    \item Usuario puede pedir proof de su balance en L2
    \item Proof se verifica contra state root publicado
\end{enumerate}

Esto permite ``exits'' de emergencia si el sequencer se vuelve malicioso.
\end{bancozk}

%-------------------------------------------
\subsection{Resumen}

\begin{keyidea}
\textbf{Merkle Patricia Tries: Puntos clave}

\begin{enumerate}
    \item \textbf{Combina tres conceptos:} Trie (búsqueda) + Patricia (compresión) + Merkle (hashing)
    
    \item \textbf{Tres tipos de nodos:} Branch (16 hijos), Extension (prefijo compartido), Leaf (valor final)
    
    \item \textbf{Cuatro tries en Ethereum:} State, Storage (por contrato), Transaction, Receipt
    
    \item \textbf{Ventaja única:} Pruebas de inclusión Y exclusión
    
    \item \textbf{Desventajas:} Proofs grandes, estado creciente, implementación compleja
    
    \item \textbf{Futuro:} Verkle trees para proofs más pequeños y stateless clients
    
    \item \textbf{En ZK:} Costoso verificar (Keccak), preferimos SMT con Poseidon para estado interno
\end{enumerate}
\end{keyidea}

% ----- COMMITMENT SCHEMES -----
%===========================================
% SECCIÓN: COMMITMENT SCHEMES - VERSIÓN EXPANDIDA
% Usa tus entornos: definition, example, keyidea, bancozk, 
%                   mathnote, intuicion, advertencia, remark
%===========================================

\section{Commitment Schemes}

\subsection{Qué son y para qué sirven}

\begin{definition}[Esquema de compromiso]
Protocolo de dos fases:
\begin{enumerate}
    \item \textbf{Commit:} $c = \text{Commit}(m, r)$
    \item \textbf{Open:} Revelar $(m, r)$, verificar $c$
\end{enumerate}
\end{definition}

Propiedades:
\begin{itemize}
    \item \textbf{Hiding:} $c$ no revela información sobre $m$
    \item \textbf{Binding:} No se puede abrir a otro mensaje
\end{itemize}

%-------------------------------------------
\subsubsection{Intuición: El Sobre Sellado}

\begin{intuicion}
Un commitment scheme es como poner un mensaje en un sobre sellado:
\begin{enumerate}
    \item \textbf{Commit:} Escribes tu predicción, la metes en un sobre, lo sellas y lo entregas
    \item \textbf{Open:} Más tarde, abres el sobre para revelar tu predicción
\end{enumerate}

\textbf{Propiedades del sobre:}
\begin{itemize}
    \item \textbf{Hiding:} Nadie puede ver a través del sobre (tu predicción está oculta)
    \item \textbf{Binding:} No puedes cambiar el contenido después de sellarlo (sin romperlo)
\end{itemize}

Los commitments criptográficos logran esto matemáticamente, sin sobres físicos.
\end{intuicion}

\begin{example}[Juego de ``piedra, papel, tijera'' remoto]
Alice y Bob quieren jugar piedra-papel-tijera por internet. Sin commitments:
\begin{itemize}
    \item Si Alice envía primero, Bob puede ver y elegir para ganar
    \item Si envían simultáneamente, ¿cómo garantizarlo?
\end{itemize}

\textbf{Con commitments:}
\begin{enumerate}
    \item Alice: $c_A = \text{Commit}(\text{``piedra''}, r_A)$, envía $c_A$
    \item Bob: $c_B = \text{Commit}(\text{``papel''}, r_B)$, envía $c_B$
    \item Alice revela: $(\text{``piedra''}, r_A)$
    \item Bob revela: $(\text{``papel''}, r_B)$
    \item Ambos verifican los commitments
    \item Bob gana (papel vence piedra)
\end{enumerate}

Ninguno puede hacer trampa: no pueden cambiar su elección después de ver el commitment del otro.
\end{example}

%-------------------------------------------
\subsubsection{Definición Formal}

\begin{definition}[Commitment Scheme]
Un esquema de compromiso consiste en tres algoritmos:

\textbf{Setup}$() \to \text{params}$:
Genera parámetros públicos (si los hay).

\textbf{Commit}$(\text{params}, m) \to (c, d)$:
\begin{itemize}
    \item Input: mensaje $m$ del espacio de mensajes $\mathcal{M}$
    \item Output: commitment $c$ y decommitment (opening) $d$
\end{itemize}

\textbf{Verify}$(\text{params}, c, m, d) \to \{0, 1\}$:
\begin{itemize}
    \item Input: commitment $c$, mensaje $m$, decommitment $d$
    \item Output: 1 (aceptar) o 0 (rechazar)
\end{itemize}
\end{definition}

\begin{definition}[Propiedades de Seguridad]
\textbf{Hiding (Ocultamiento):}

Un commitment es \textit{hiding} si el commitment $c$ no revela información sobre $m$.

\begin{itemize}
    \item \textbf{Perfectly hiding:} Incluso con poder computacional ilimitado, $c$ no revela nada sobre $m$. Formalmente: $\forall m_0, m_1$, las distribuciones de $\text{Commit}(m_0)$ y $\text{Commit}(m_1)$ son idénticas.
    
    \item \textbf{Computationally hiding:} Un adversario PPT no puede distinguir $\text{Commit}(m_0)$ de $\text{Commit}(m_1)$ con ventaja no negligible.
\end{itemize}

\textbf{Binding (Vinculante):}

Un commitment es \textit{binding} si no se puede abrir a dos mensajes diferentes.

\begin{itemize}
    \item \textbf{Perfectly binding:} Incluso con poder ilimitado, no existe $(m', d')$ con $m' \neq m$ tal que $\text{Verify}(c, m', d') = 1$.
    
    \item \textbf{Computationally binding:} Un adversario PPT no puede encontrar tal $(m', d')$ con probabilidad no negligible.
\end{itemize}
\end{definition}

\begin{mathnote}
\textbf{Teorema de imposibilidad:}

Ningún esquema de compromiso puede ser simultáneamente \textbf{perfectly hiding} Y \textbf{perfectly binding}.

\textbf{Demostración informal:}
\begin{itemize}
    \item Si es perfectly hiding, entonces para cada $c$, múltiples mensajes $m$ son ``compatibles'' con $c$
    \item Pero entonces un adversario con poder ilimitado podría encontrar aperturas para diferentes mensajes
    \item Contradiciendo perfect binding
\end{itemize}

Por lo tanto, debemos elegir cuál propiedad es perfecta y cuál es computacional.
\end{mathnote}

%-------------------------------------------
\subsubsection{El Commitment más Simple: Hash}

\begin{definition}[Hash Commitment]
Usando una función hash $H$:
\[ c = H(m \| r) \]

donde $r$ es aleatoriedad (randomness) elegida uniformemente.
\end{definition}

\begin{keyidea}
\textbf{Propiedades del hash commitment:}
\begin{itemize}
    \item \textbf{Computationally hiding:} Asumiendo $H$ es un oráculo aleatorio, $c$ no revela $m$ (por $r$)
    \item \textbf{Computationally binding:} Encontrar $m' \neq m$ y $r'$ con $H(m' \| r') = H(m \| r)$ requiere romper resistencia a colisiones
\end{itemize}

Ni perfectly hiding ni perfectly binding, pero ambas propiedades computacionalmente fuertes.
\end{keyidea}

\begin{example}[Hash commitment en la práctica]
\begin{verbatim}
import hashlib
import secrets

def commit(message: bytes) -> tuple[bytes, bytes]:
    r = secrets.token_bytes(32)  # 256 bits de aleatoriedad
    c = hashlib.sha256(message + r).digest()
    return (c, r)  # commitment y decommitment

def verify(c: bytes, message: bytes, r: bytes) -> bool:
    return hashlib.sha256(message + r).digest() == c

# Uso
msg = b"Mi prediccion: Bitcoin a $100,000"
commitment, opening = commit(msg)

print(f"Commitment: {commitment.hex()[:16]}...")
# Más tarde...
assert verify(commitment, msg, opening)
\end{verbatim}
\end{example}

%-------------------------------------------
\subsubsection{Aplicaciones de Commitments}

\begin{keyidea}
Los commitments son un \textbf{building block} fundamental en criptografía:

\begin{center}
\begin{tabular}{ll}
\toprule
\textbf{Aplicación} & \textbf{Uso del commitment} \\
\midrule
Coin flipping & Comprometer a bit antes de revelar \\
Zero-knowledge proofs & Prover compromete a witness \\
Auction protocols & Comprometer a bid sellado \\
Voting & Comprometer a voto antes de contar \\
Timestamping & Comprometer a documento en blockchain \\
Confidential transactions & Comprometer a monto sin revelarlo \\
\bottomrule
\end{tabular}
\end{center}
\end{keyidea}

%-------------------------------------------
\subsection{Pedersen Commitments}

\begin{definition}[Pedersen]
Con generadores $g, h$ donde $\log_g h$ es desconocido:
\[
C = g^m \cdot h^r
\]
\end{definition}

\begin{itemize}
    \item \textbf{Perfectly hiding:} Para cualquier $c$ y $m'$, existe $r'$
    \item \textbf{Computationally binding:} Encontrar dos aperturas = DLP
    \item \textbf{Homomórfico:} $C_1 \cdot C_2 = \text{Commit}(m_1 + m_2, r_1 + r_2)$
\end{itemize}

Usado en: Confidential Transactions, Bulletproofs, ZK.

%-------------------------------------------
\subsubsection{Contexto Histórico}

\begin{remark}
Torben Pedersen propuso este esquema en 1991 en el paper ``Non-Interactive and Information-Theoretic Secure Verifiable Secret Sharing''.

El esquema se basa en el trabajo previo de:
\begin{itemize}
    \item Diffie-Hellman (1976): Problema del logaritmo discreto
    \item ElGamal (1985): Cifrado basado en DLP
    \item Feldman (1987): VSS verificable
\end{itemize}

La innovación de Pedersen fue combinar la estructura multiplicativa con dos generadores independientes para lograr perfect hiding.
\end{remark}

%-------------------------------------------
\subsubsection{Construcción Detallada}

\begin{definition}[Pedersen Commitment - Formal]
\textbf{Setup:}
\begin{enumerate}
    \item Elegir grupo cíclico $\mathbb{G}$ de orden primo $q$
    \item Elegir dos generadores $g, h \in \mathbb{G}$ tales que $\log_g h$ es desconocido
    \item Publicar $(q, \mathbb{G}, g, h)$
\end{enumerate}

\textbf{Commit}$(m, r)$ donde $m \in \mathbb{Z}_q$ y $r \xleftarrow{\$} \mathbb{Z}_q$:
\[ C = g^m \cdot h^r \]

\textbf{Open:} Revelar $(m, r)$

\textbf{Verify}$(C, m, r)$: Aceptar si $C = g^m \cdot h^r$
\end{definition}

\begin{center}
\begin{tikzpicture}[
    gen/.style={draw, circle, fill=blue!20, minimum size=1cm},
    commit/.style={draw, rectangle, fill=green!20, minimum width=1.5cm, minimum height=0.8cm}
]
    % Generadores
    \node[gen] (g) at (0, 0) {$g$};
    \node[gen] (h) at (4, 0) {$h$};
    
    % Exponentes
    \node at (0, 1.5) {$m$ (mensaje)};
    \node at (4, 1.5) {$r$ (randomness)};
    
    % Flechas
    \draw[->, thick] (0, 1.2) -- (g);
    \draw[->, thick] (4, 1.2) -- (h);
    
    % Resultado
    \node[commit] (c) at (2, -1.5) {$C = g^m \cdot h^r$};
    
    \draw[->, thick] (g) -- (c);
    \draw[->, thick] (h) -- (c);
\end{tikzpicture}
\end{center}

%-------------------------------------------
\subsubsection{Análisis de Seguridad}

\begin{mathnote}
\textbf{Perfect Hiding:}

Para cualquier commitment $C$ y cualquier mensaje objetivo $m'$, existe un $r'$ tal que $C = g^{m'} h^{r'}$.

\textbf{Demostración:}
\begin{align}
C &= g^m h^r \\
g^{m'} h^{r'} &= C \\
h^{r'} &= C \cdot g^{-m'} = g^m h^r \cdot g^{-m'} = g^{m-m'} h^r \\
r' &= r + (m - m') \cdot \log_g h
\end{align}

Para cada $m'$, existe exactamente un $r'$ que satisface la ecuación. Por lo tanto, viendo solo $C$, todos los mensajes son igualmente probables.

\textbf{Nota:} Esto asume que no conocemos $\log_g h$. Si lo conociéramos, podríamos calcular $r'$, pero aún no podríamos determinar el $m$ original solo viendo $C$.
\end{mathnote}

\begin{mathnote}
\textbf{Computational Binding:}

Supongamos que un adversario encuentra $(m_1, r_1)$ y $(m_2, r_2)$ con $m_1 \neq m_2$ tales que:
\[ g^{m_1} h^{r_1} = g^{m_2} h^{r_2} \]

Entonces:
\begin{align}
g^{m_1 - m_2} &= h^{r_2 - r_1} \\
g^{(m_1 - m_2) / (r_2 - r_1)} &= h
\end{align}

Por lo tanto: $\log_g h = (m_1 - m_2) / (r_2 - r_1)$

¡El adversario ha resuelto el DLP entre $g$ y $h$! Si el DLP es difícil en $\mathbb{G}$, esto es computacionalmente imposible.
\end{mathnote}

%-------------------------------------------
\subsubsection{La Propiedad Homomórfica}

\begin{keyidea}
\textbf{El superpoder de Pedersen:} Los commitments son \textbf{aditivamente homomórficos}.

\[ C_1 \cdot C_2 = (g^{m_1} h^{r_1}) \cdot (g^{m_2} h^{r_2}) = g^{m_1 + m_2} h^{r_1 + r_2} = \text{Commit}(m_1 + m_2, r_1 + r_2) \]

Podemos sumar commitments sin conocer los mensajes originales.
\end{keyidea}

\begin{example}[Verificación de sumas sin revelar valores]
\textbf{Escenario:} Alice tiene 3 valores secretos que suman 100. Quiere probar esto sin revelar los valores individuales.

\begin{verbatim}
Valores secretos: a=30, b=25, c=45
Randomness: r_a, r_b, r_c (elegidos aleatoriamente)

Commitments:
C_a = g^30 · h^(r_a)
C_b = g^25 · h^(r_b)
C_c = g^45 · h^(r_c)

Alice publica: C_a, C_b, C_c

Verificador calcula:
C_total = C_a · C_b · C_c = g^(30+25+45) · h^(r_a+r_b+r_c)
        = g^100 · h^(r_total)

Alice revela: r_total = r_a + r_b + r_c

Verificador comprueba:
C_total ?= g^100 · h^(r_total)

Si coincide, la suma es 100 (sin saber los valores individuales).
\end{verbatim}
\end{example}

%-------------------------------------------
\subsubsection{Pedersen en Notación Aditiva (Curvas Elípticas)}

\begin{definition}[Pedersen sobre Curvas Elípticas]
En una curva elíptica con puntos base $G, H$ donde $\log_G H$ es desconocido:
\[ C = m \cdot G + r \cdot H \]

Es la misma construcción, pero en notación aditiva.
\end{definition}

\begin{remark}
En la práctica, casi siempre usamos curvas elípticas para Pedersen porque:
\begin{enumerate}
    \item Claves más cortas para mismo nivel de seguridad
    \item Operaciones más rápidas
    \item Compatibilidad con otros protocolos (ECDSA, Schnorr)
\end{enumerate}
\end{remark}

%-------------------------------------------
\subsubsection{Generación de $G$ y $H$}

\begin{advertencia}
\textbf{Crítico:} Nadie debe conocer $\log_G H$.

Si alguien conoce $x$ tal que $H = xG$:
\begin{align}
C &= mG + rH = mG + r(xG) = (m + rx)G
\end{align}

Puede abrir el commitment a cualquier $m'$ usando:
\[ r' = (m - m' + rx) / x \]
\[ C = m'G + r'H = m'G + r'(xG) = (m' + r'x)G = (m + rx)G \quad \checkmark \]

¡El commitment ya no es binding!
\end{advertencia}

\begin{keyidea}
\textbf{Métodos para generar $H$ de forma segura:}

\textbf{1. Hash-to-curve (Nothing-up-my-sleeve):}
\[ H = \text{HashToCurve}(\text{``Pedersen generator H''}) \]

Nadie puede conocer $\log_G H$ porque tendría que invertir el hash.

\textbf{2. Trusted setup con múltiples partes:}

Cada parte $i$ elige $x_i$ secreto:
\[ H = \sum_i x_i \cdot G \]

Si al menos una parte es honesta y descarta su $x_i$, nadie conoce $\log_G H$.

\textbf{3. Usar punto estándar:}

Curvas como secp256k1 tienen un segundo generador estándar derivado de forma verificable.
\end{keyidea}

%-------------------------------------------
\subsubsection{Ejemplo Numérico}

\begin{example}[Pedersen sobre curva pequeña]
\textbf{Parámetros (toy example):}
\begin{itemize}
    \item Grupo: $\mathbb{Z}_{23}^*$ (multiplicativo módulo 23)
    \item $g = 5$ (generador)
    \item $h = 7$ (segundo generador, $\log_5 7$ desconocido)
    \item Orden: $q = 22$
\end{itemize}

\textbf{Commit al mensaje $m = 10$ con $r = 15$:}
\begin{align}
C &= g^m \cdot h^r \mod 23 \\
  &= 5^{10} \cdot 7^{15} \mod 23 \\
  &= 9765625 \cdot 4747561509943 \mod 23 \\
  &\equiv 9 \cdot 12 \mod 23 \\
  &= 108 \mod 23 \\
  &= 16
\end{align}

\textbf{Verificación:}
\[ 5^{10} \cdot 7^{15} \mod 23 = 16 \quad \checkmark \]

\textbf{Perfect hiding:} Para $m' = 5$, ¿existe $r'$?
\begin{align}
16 &= 5^5 \cdot 7^{r'} \mod 23 \\
16 &= 3125 \cdot 7^{r'} \mod 23 \\
16 &= 20 \cdot 7^{r'} \mod 23 \\
7^{r'} &= 16 \cdot 20^{-1} = 16 \cdot 15 = 240 \equiv 10 \mod 23
\end{align}

Buscamos $r'$ tal que $7^{r'} \equiv 10 \pmod{23}$. Por prueba: $7^{12} = 10 \mod 23$.

Por lo tanto, $(m', r') = (5, 12)$ también abre al commitment 16.

¡Esto demuestra perfect hiding! Pero encontrar este $r'$ requirió resolver DLP (computationally hard en grupos grandes).
\end{example}

%-------------------------------------------
\subsubsection{Vector Pedersen Commitments}

\begin{definition}[Vector Pedersen]
Para comprometer a un vector $\vec{m} = (m_1, \ldots, m_n)$:

\[ C = \sum_{i=1}^n m_i \cdot G_i + r \cdot H \]

donde $G_1, \ldots, G_n, H$ son generadores independientes.
\end{definition}

\begin{keyidea}
Los vector Pedersen commitments son fundamentales para:
\begin{itemize}
    \item \textbf{Bulletproofs}: Prueban que un valor está en un rango $[0, 2^n)$
    \item \textbf{Inner Product Arguments}: Prueban $\langle \vec{a}, \vec{b} \rangle = c$
    \item \textbf{Polynomial commitments}: Comprometer a coeficientes de polinomio
\end{itemize}
\end{keyidea}

%-------------------------------------------
\subsubsection{Aplicaciones}

\begin{remark}
\textbf{Confidential Transactions (Monero, Mimblewimble):}

En lugar de montos públicos, cada output tiene un Pedersen commitment al monto:
\[ C = m \cdot G + r \cdot H \]

\textbf{Propiedad clave:} Se puede verificar que inputs = outputs sin conocer montos:
\[ \sum C_{\text{inputs}} - \sum C_{\text{outputs}} = 0 \cdot G + r_{\text{fee}} \cdot H \]

Solo funciona si la suma de montos es correcta (homomorfismo).

\textbf{Problema:} Sin restricciones adicionales, alguien podría crear monto negativo.

\textbf{Solución:} Range proofs (Bulletproofs) prueban $m \in [0, 2^{64})$.
\end{remark}

%-------------------------------------------
\subsection{KZG Commitments}

\begin{definition}[KZG]
Compromiso a polinomio con evaluación verificable.
\end{definition}

Setup: $\{g, g^\tau, g^{\tau^2}, \ldots, g^{\tau^d}\}$ donde $\tau$ es secreto (toxic waste).

\begin{itemize}
    \item Commit: $C = g^{p(\tau)} = \prod_i (g^{\tau^i})^{a_i}$
    \item Evaluación: Probar $p(z) = y$ con prueba $\pi = g^{q(\tau)}$
    \item Donde $q(x) = \frac{p(x) - y}{x - z}$
    \item Verificar: $e(C - g^y, g) = e(\pi, g^\tau - g^z)$
\end{itemize}

Tamaño de prueba: 1 elemento de grupo (48 bytes en BLS12-381).

%-------------------------------------------
\subsubsection{Contexto y Motivación}

\begin{intuicion}
KZG permite hacer algo mágico: comprometer a un polinomio entero en un solo punto de curva, y después probar evaluaciones individuales sin revelar el polinomio completo.

Es como tener una ``caja negra'' que contiene un polinomio $p(x)$:
\begin{itemize}
    \item La caja es pequeña (48 bytes)
    \item Puedes probar ``mi caja contiene un polinomio donde $p(5) = 17$''
    \item La prueba también es pequeña (48 bytes)
    \item Sin revelar ningún otro punto del polinomio
\end{itemize}
\end{intuicion}

\begin{remark}
\textbf{Kate-Zaverucha-Goldberg (2010):}

Aniket Kate, Gregory Zaverucha, y Ian Goldberg propusieron este esquema en ``Constant-Size Commitments to Polynomials and Their Applications''.

El paper original lo llamaba ``polynomial commitment scheme'', pero la comunidad lo nombró KZG por los autores.

\textbf{Adopción:}
\begin{center}
\begin{tabular}{ll}
\toprule
\textbf{Protocolo} & \textbf{Uso de KZG} \\
\midrule
PLONK & Commitment a polinomios del circuito \\
Groth16 & Verificación de polinomios \\
Ethereum (Danksharding) & Data availability sampling \\
Scroll, Polygon zkEVM & zkRollups basados en PLONK \\
\bottomrule
\end{tabular}
\end{center}
\end{remark}

%-------------------------------------------
\subsubsection{El Setup: Structured Reference String (SRS)}

\begin{definition}[Trusted Setup KZG]
Para polinomios de grado hasta $d$, el setup genera:

\textbf{En $\mathbb{G}_1$:}
\[ \{g_1, g_1^\tau, g_1^{\tau^2}, \ldots, g_1^{\tau^d}\} \]

\textbf{En $\mathbb{G}_2$:}
\[ \{g_2, g_2^\tau\} \]

donde $\tau \xleftarrow{\$} \mathbb{F}_p$ es el ``toxic waste'' que debe destruirse.
\end{definition}

\begin{advertencia}
\textbf{El Toxic Waste:}

Si alguien conoce $\tau$, puede:
\begin{enumerate}
    \item Crear commitments falsos a cualquier polinomio
    \item Crear pruebas de evaluación falsas
    \item Romper completamente la seguridad del sistema
\end{enumerate}

Por esto se llama ``toxic waste'': debe destruirse inmediatamente después del setup.
\end{advertencia}

\begin{keyidea}
\textbf{Ceremonia de Powers of Tau:}

Para generar el SRS sin que nadie conozca $\tau$:
\begin{enumerate}
    \item Participante 1 elige $\tau_1$, computa $(g^{\tau_1}, g^{\tau_1^2}, \ldots)$, destruye $\tau_1$
    \item Participante 2 recibe, elige $\tau_2$, computa $(g^{\tau_1 \tau_2}, g^{(\tau_1\tau_2)^2}, \ldots)$, destruye $\tau_2$
    \item ...continúa con $n$ participantes
\end{enumerate}

El $\tau$ final es $\tau = \tau_1 \cdot \tau_2 \cdots \tau_n$.

\textbf{Seguridad:} Si AL MENOS UN participante es honesto y destruye su $\tau_i$, nadie conoce $\tau$ completo.

\textbf{Ceremonias famosas:}
\begin{itemize}
    \item Zcash Powers of Tau (2018): 87 participantes
    \item Ethereum KZG Ceremony (2023): 141,000+ participantes
\end{itemize}
\end{keyidea}

%-------------------------------------------
\subsubsection{El Algoritmo KZG}

\begin{definition}[KZG Commitment Scheme - Completo]
\textbf{Commit}$(p(x))$:

Dado polinomio $p(x) = \sum_{i=0}^d a_i x^i$:
\[ C = g_1^{p(\tau)} = \prod_{i=0}^d (g_1^{\tau^i})^{a_i} \]

Nota: Calculamos $g_1^{p(\tau)}$ sin conocer $\tau$, usando el SRS.

\textbf{CreateProof}$(p(x), z)$ para probar $p(z) = y$:

\begin{enumerate}
    \item Calcular $y = p(z)$
    \item Calcular polinomio cociente: $q(x) = \frac{p(x) - y}{x - z}$
    
    (Esto es exactamente divisible porque $p(z) = y$ implica $(x-z) | (p(x) - y)$)
    
    \item Computar prueba: $\pi = g_1^{q(\tau)}$
\end{enumerate}

\textbf{Verify}$(C, z, y, \pi)$:

Verificar la ecuación de pairing:
\[ e(C - g_1^y, g_2) = e(\pi, g_2^\tau - g_2^z) \]

O equivalentemente:
\[ e(C - g_1^y, g_2) = e(\pi, g_2^\tau \cdot g_2^{-z}) \]
\end{definition}

%-------------------------------------------
\subsubsection{¿Por Qué Funciona?}

\begin{mathnote}
\textbf{Verificación de correctitud:}

Queremos verificar que $p(z) = y$ usando el pairing.

\textbf{Lado izquierdo:}
\[ e(C - g_1^y, g_2) = e(g_1^{p(\tau)} \cdot g_1^{-y}, g_2) = e(g_1^{p(\tau) - y}, g_2) = e(g_1, g_2)^{p(\tau) - y} \]

\textbf{Lado derecho:}
\begin{align}
e(\pi, g_2^\tau - g_2^z) &= e(g_1^{q(\tau)}, g_2^{\tau - z}) \\
                         &= e(g_1, g_2)^{q(\tau) \cdot (\tau - z)}
\end{align}

\textbf{La clave:} Si $p(z) = y$, entonces:
\[ p(x) - y = q(x) \cdot (x - z) \]

Evaluando en $x = \tau$:
\[ p(\tau) - y = q(\tau) \cdot (\tau - z) \]

Por lo tanto, ambos lados del pairing son iguales:
\[ e(g_1, g_2)^{p(\tau) - y} = e(g_1, g_2)^{q(\tau)(\tau - z)} \quad \checkmark \]
\end{mathnote}

%-------------------------------------------
\subsubsection{Ejemplo Numérico}

\begin{example}[KZG simplificado]
\textbf{Polinomio:} $p(x) = 2x^2 + 3x + 1$

\textbf{Queremos probar:} $p(2) = 2(4) + 3(2) + 1 = 8 + 6 + 1 = 15$

\textbf{Paso 1: Commitment}

Usando SRS $\{g, g^\tau, g^{\tau^2}\}$:
\[ C = g^{p(\tau)} = g^{2\tau^2 + 3\tau + 1} = (g^{\tau^2})^2 \cdot (g^\tau)^3 \cdot g^1 \]

\textbf{Paso 2: Polinomio cociente}
\begin{align}
q(x) &= \frac{p(x) - 15}{x - 2} \\
     &= \frac{2x^2 + 3x + 1 - 15}{x - 2} \\
     &= \frac{2x^2 + 3x - 14}{x - 2}
\end{align}

División polinómica:
\[ 2x^2 + 3x - 14 = (x - 2)(2x + 7) \]

Verificación: $(x-2)(2x+7) = 2x^2 + 7x - 4x - 14 = 2x^2 + 3x - 14$ ✓

Por lo tanto: $q(x) = 2x + 7$

\textbf{Paso 3: Prueba}
\[ \pi = g^{q(\tau)} = g^{2\tau + 7} = (g^\tau)^2 \cdot g^7 \]

\textbf{Paso 4: Verificación}

El verificador comprueba:
\[ e(C \cdot g^{-15}, g) \stackrel{?}{=} e(\pi, g^\tau \cdot g^{-2}) \]

Ambos lados evalúan a $e(g,g)^{q(\tau)(\tau-2)} = e(g,g)^{p(\tau)-15}$.
\end{example}

%-------------------------------------------
\subsubsection{Propiedades de KZG}

\begin{center}
\begin{tabular}{lc}
\toprule
\textbf{Propiedad} & \textbf{KZG} \\
\midrule
Tamaño de commitment & 1 elemento $\mathbb{G}_1$ (48 bytes) \\
Tamaño de prueba (1 evaluación) & 1 elemento $\mathbb{G}_1$ (48 bytes) \\
Tiempo de verificación & 2 pairings \\
Hiding & Computationally hiding \\
Binding & Computationally binding (bajo q-SDH) \\
Homomórfico & Sí (aditivo) \\
Trusted setup & Sí (universal, actualizable) \\
\bottomrule
\end{tabular}
\end{center}

\begin{keyidea}
\textbf{Homomorfismo de KZG:}

Si $C_1 = \text{Commit}(p_1)$ y $C_2 = \text{Commit}(p_2)$:
\[ C_1 \cdot C_2 = g^{p_1(\tau)} \cdot g^{p_2(\tau)} = g^{(p_1 + p_2)(\tau)} = \text{Commit}(p_1 + p_2) \]

Podemos combinar commitments a polinomios sin conocer los polinomios.
\end{keyidea}

%-------------------------------------------
\subsubsection{Batch Opening}

\begin{definition}[Batch Opening]
Para probar múltiples evaluaciones $(z_1, y_1), \ldots, (z_k, y_k)$ del mismo polinomio:

\textbf{Opción 1: Pruebas individuales}
$k$ pruebas de 48 bytes cada una = $48k$ bytes

\textbf{Opción 2: Batch proof}
Una sola prueba de 48 bytes (con técnicas de agregación)
\end{definition}

\begin{remark}
El batch opening es crucial para eficiencia en SNARKs como PLONK, donde se verifican muchas evaluaciones simultáneamente.
\end{remark}

%-------------------------------------------
\subsubsection{KZG vs Otros Polynomial Commitments}

\begin{center}
\begin{tabular}{lccc}
\toprule
\textbf{Propiedad} & \textbf{KZG} & \textbf{FRI} & \textbf{Bulletproofs} \\
\midrule
Tamaño prueba & $O(1)$ & $O(\log^2 d)$ & $O(\log d)$ \\
Tiempo verificación & $O(1)$ pairings & $O(\log^2 d)$ hashes & $O(d)$ \\
Trusted setup & Sí & No & No \\
Quantum-safe & No & Sí & No \\
Asunción & q-SDH & Hash collision & DLP \\
\bottomrule
\end{tabular}
\end{center}

\begin{keyidea}
\textbf{Cuándo usar cada uno:}
\begin{itemize}
    \item \textbf{KZG}: Cuando se pueden tolerar trusted setups y se necesita máxima eficiencia (PLONK, Danksharding)
    \item \textbf{FRI}: Cuando se necesita transparencia (sin trusted setup) y post-quantum es importante (STARKs)
    \item \textbf{Bulletproofs}: Para range proofs y cuando el tamaño de prueba importa más que el tiempo de verificación
\end{itemize}
\end{keyidea}

%-------------------------------------------
\subsubsection{KZG en Ethereum (Danksharding)}

\begin{definition}[Proto-Danksharding (EIP-4844)]
Ethereum usa KZG para ``blob transactions'':
\begin{enumerate}
    \item Rollups publican datos como ``blobs'' ($\sim$128 KB cada uno)
    \item Cada blob se compromete con KZG
    \item El commitment va en el header del bloque
    \item Los datos completos se ``podan'' después de $\sim$2 semanas
\end{enumerate}
\end{definition}

\begin{remark}
\textbf{Data Availability Sampling (futuro):}

Con KZG, los light clients podrán verificar que los datos están disponibles sin descargarlos:
\begin{enumerate}
    \item Solicitar evaluaciones aleatorias del polinomio blob
    \item Verificar proofs KZG de esas evaluaciones
    \item Con suficientes muestras válidas, hay alta confianza de que los datos existen
\end{enumerate}

Esto permite escalar Ethereum sin que cada nodo descargue todos los datos.
\end{remark}

%-------------------------------------------
\subsubsection{Implementación}

\begin{example}[KZG con arkworks (Rust)]
\begin{verbatim}
use ark_bls12_381::{Bls12_381, Fr, G1Projective};
use ark_poly::{univariate::DensePolynomial, UVPolynomial};
use ark_poly_commit::kzg10::{KZG10, Powers};

fn main() {
    // Setup (en producción, usar ceremonia)
    let rng = &mut ark_std::test_rng();
    let degree = 10;
    let params = KZG10::<Bls12_381, DensePolynomial<Fr>>::setup(
        degree, false, rng
    ).unwrap();
    
    // Polinomio: p(x) = 2x^2 + 3x + 1
    let poly = DensePolynomial::from_coefficients_vec(vec![
        Fr::from(1u64),  // constante
        Fr::from(3u64),  // coef de x
        Fr::from(2u64),  // coef de x^2
    ]);
    
    // Commit
    let (commitment, randomness) = KZG10::commit(&params, &poly, None, None)
        .unwrap();
    
    // Evaluar y probar en z = 2
    let z = Fr::from(2u64);
    let y = poly.evaluate(&z);  // p(2) = 15
    
    let proof = KZG10::open(&params, &poly, z, &randomness).unwrap();
    
    // Verificar
    let valid = KZG10::check(&params.vk, &commitment, z, y, &proof).unwrap();
    assert!(valid);
    
    println!("p(2) = {} verificado!", y);
}
\end{verbatim}
\end{example}

%-------------------------------------------
\subsection{Comparación de Commitment Schemes}

\begin{center}
\begin{tabular}{lcccc}
\toprule
\textbf{Propiedad} & \textbf{Hash} & \textbf{Pedersen} & \textbf{KZG} \\
\midrule
Tamaño commitment & 32 B & 32-48 B & 48 B \\
Hiding & Comput. & Perfect & Comput. \\
Binding & Comput. & Comput. & Comput. \\
Homomórfico & No & Sí (aditivo) & Sí (aditivo) \\
Trusted setup & No & No & Sí \\
Evaluación verificable & No & No & Sí \\
Velocidad commit & Muy rápido & Rápido & Medio \\
\bottomrule
\end{tabular}
\end{center}

%-------------------------------------------
\subsection{Aplicación en Banco ZK}

\begin{bancozk}
Los commitment schemes son fundamentales en toda la arquitectura de Banco ZK:

\textbf{1. Note Commitments (Pedersen)}

Cada ``nota'' (output de transacción) se compromete:
\[ C_{\text{note}} = \text{Pedersen}(\text{value}, \text{pubkey}, \text{salt}) \]

\begin{verbatim}
struct Note {
    value: u64,      // Monto en wei
    pubkey: Point,   // Clave pública del dueño
    salt: Field,     // Aleatoriedad
}

commitment = value * G_value + pubkey.x * G_pk + salt * H
\end{verbatim}

\textbf{¿Por qué Pedersen?}
\begin{itemize}
    \item Perfect hiding: Nadie puede saber el monto viendo $C$
    \item Homomórfico: Verificar suma de inputs = suma de outputs
    \item Eficiente en circuitos ZK
\end{itemize}

\textbf{2. Nullifier Derivation}

Para gastar una nota sin revelar cuál:
\[ \text{nullifier} = \text{Poseidon}(C_{\text{note}}, \text{spending\_key}) \]

El nullifier es público (para prevenir doble gasto), pero no revela qué nota se gastó.

\textbf{3. KZG para PLONK Proofs}

Nuestro sistema SNARK usa PLONK, que internamente usa KZG:
\begin{verbatim}
Circuito de transacción:
1. Verificar Merkle proof de nota
2. Verificar nullifier correcto
3. Verificar firma
4. Verificar conservación de valor

PLONK convierte esto en:
- Polinomios de constraints
- KZG commitments a esos polinomios
- Prueba final: ~700 bytes
\end{verbatim}

\textbf{4. Confidential Transactions}

Los montos permanecen ocultos usando Pedersen:
\begin{verbatim}
Input 1:  C_1 = 50 * G + r_1 * H
Input 2:  C_2 = 30 * G + r_2 * H
Output 1: C_3 = 70 * G + r_3 * H
Output 2: C_4 = 10 * G + r_4 * H  (fee)

Verificación (homomorfismo):
C_1 + C_2 - C_3 - C_4 
  = (50+30-70-10)*G + (r_1+r_2-r_3-r_4)*H
  = 0*G + r_diff*H

Prover revela r_diff para verificar suma = 0.
\end{verbatim}

\textbf{5. Range Proofs}

Para evitar montos negativos (que romperían la conservación), incluimos range proofs:
\begin{verbatim}
Para cada output con commitment C:
1. Probar que value ∈ [0, 2^64)
2. Usamos Bulletproofs sobre el Pedersen commitment
3. O verificación dentro del circuito SNARK
\end{verbatim}

\textbf{6. Data Blobs (futuro)}

Con EIP-4844, publicaremos datos de batch como blobs:
\begin{verbatim}
Batch de 1000 transacciones:
1. Serializar transacciones → ~128 KB
2. Codificar como polinomio
3. KZG commitment publicado on-chain
4. Datos disponibles por 2 semanas para verificadores
5. Después se prunan (solo commitment permanece)
\end{verbatim}

\textbf{7. Resumen de uso}

\begin{center}
\begin{tabular}{lll}
\toprule
\textbf{Componente} & \textbf{Commitment} & \textbf{Razón} \\
\midrule
Note values & Pedersen & Perfect hiding, homomórfico \\
Merkle tree nodes & Poseidon hash & ZK-friendly \\
SNARK proofs & KZG (PLONK) & Tamaño constante \\
Data availability & KZG (blobs) & DAS eficiente \\
Addresses & Hash commitment & Simple, binding \\
\bottomrule
\end{tabular}
\end{center}
\end{bancozk}

%-------------------------------------------
\subsection{Resumen}

\begin{keyidea}
\textbf{Commitment Schemes: Puntos clave}

\begin{enumerate}
    \item \textbf{Dos propiedades:} Hiding (oculta mensaje) y Binding (no se puede cambiar)
    
    \item \textbf{Trade-off:} No pueden ser perfectamente ambas simultáneamente
    
    \item \textbf{Hash commitments:} Simples, pero no homomórficos
    
    \item \textbf{Pedersen:} Perfect hiding, homomórfico, ideal para valores confidenciales
    
    \item \textbf{KZG:} Compromete polinomios, pruebas de evaluación en tamaño constante, requiere trusted setup
    
    \item \textbf{Aplicaciones:} ZK proofs, confidential transactions, data availability, voting, auctions
\end{enumerate}
\end{keyidea}

%===========================================
% SUBSECCIÓN: POLYNOMIAL COMMITMENTS - VERSIÓN EXPANDIDA
% Usa tus entornos: definition, example, keyidea, bancozk, 
%                   mathnote, intuicion, advertencia, remark
%===========================================

\subsection{Polynomial Commitments}

Generalización de compromisos a polinomios:

\begin{table}[htbp]
\centering
\begin{tabular}{lccc}
\toprule
\textbf{Esquema} & \textbf{Prueba} & \textbf{Verificación} & \textbf{Setup} \\
\midrule
KZG & $O(1)$ & $O(1)$ pairing & Trusted \\
IPA & $O(\log n)$ & $O(n)$ & Ninguno \\
FRI & $O(\log^2 n)$ & $O(\log^2 n)$ & Ninguno \\
\bottomrule
\end{tabular}
\end{table}

%-------------------------------------------
\subsubsection{¿Por Qué Polinomios?}

\begin{intuicion}
Los polinomios son el ``lenguaje universal'' de los SNARKs modernos. Casi todos los sistemas de pruebas convierten el problema a verificar en ecuaciones polinomiales.

\textbf{¿Por qué?}
\begin{enumerate}
    \item \textbf{Expresividad}: Cualquier circuito se puede representar como restricciones polinomiales
    \item \textbf{Schwartz-Zippel}: Si dos polinomios diferentes de grado $d$ coinciden en un punto aleatorio, la probabilidad es $\leq d/|F|$ (negligible)
    \item \textbf{Compresión}: Un polinomio de grado $n$ se puede ``resumir'' evaluándolo en un punto
    \item \textbf{Homomorfismo}: Muchos esquemas permiten operar sobre commitments
\end{enumerate}

El polynomial commitment es la herramienta que permite verificar propiedades de polinomios sin revelarlos.
\end{intuicion}

\begin{keyidea}
\textbf{El paradigma moderno de SNARKs:}

\begin{enumerate}
    \item Convertir el problema en restricciones polinomiales
    \item El prover compromete a los polinomios (usando PCS)
    \item El verificador pide evaluaciones en puntos aleatorios
    \item El prover proporciona evaluaciones con pruebas
    \item Por Schwartz-Zippel, si las evaluaciones son correctas, los polinomios probablemente son correctos
\end{enumerate}

El \textbf{Polynomial Commitment Scheme (PCS)} es el componente que hace posible los pasos 2-4.
\end{keyidea}

%-------------------------------------------
\subsubsection{Definición Formal}

\begin{definition}[Polynomial Commitment Scheme]
Un PCS consiste en los siguientes algoritmos:

\textbf{Setup}$(1^\lambda, d) \to \text{pp}$:
\begin{itemize}
    \item Input: parámetro de seguridad $\lambda$, grado máximo $d$
    \item Output: parámetros públicos pp
\end{itemize}

\textbf{Commit}$(\text{pp}, p(x)) \to (C, \text{aux})$:
\begin{itemize}
    \item Input: polinomio $p(x)$ de grado $\leq d$
    \item Output: commitment $C$ y datos auxiliares (opcionales)
\end{itemize}

\textbf{Open}$(\text{pp}, p(x), z, \text{aux}) \to \pi$:
\begin{itemize}
    \item Input: polinomio, punto de evaluación $z$
    \item Output: prueba $\pi$ de que $p(z) = y$
\end{itemize}

\textbf{Verify}$(\text{pp}, C, z, y, \pi) \to \{0, 1\}$:
\begin{itemize}
    \item Input: commitment, punto, valor declarado, prueba
    \item Output: aceptar o rechazar
\end{itemize}
\end{definition}

\begin{definition}[Propiedades de Seguridad]
\textbf{Correctitud:} Si $p(z) = y$, entonces Verify acepta.

\textbf{Binding:} Es computacionalmente difícil encontrar $C, z, y, y', \pi, \pi'$ con $y \neq y'$ tales que Verify acepta ambos $(C, z, y, \pi)$ y $(C, z, y', \pi')$.

\textbf{Hiding (opcional):} El commitment $C$ no revela información sobre $p(x)$ más allá de lo que revelan las evaluaciones abiertas.
\end{definition}

%-------------------------------------------
\subsubsection{KZG (Kate-Zaverucha-Goldberg)}

\begin{remark}
KZG fue cubierto en detalle en la sección anterior. Aquí resumimos sus características clave como PCS.
\end{remark}

\begin{definition}[KZG como PCS]
\textbf{Setup:} Trusted setup genera $\{g^{\tau^i}\}_{i=0}^d$ (Powers of Tau)

\textbf{Commit:} $C = g^{p(\tau)} = \prod_i (g^{\tau^i})^{a_i}$ para $p(x) = \sum_i a_i x^i$

\textbf{Open en $z$:} Calcular $q(x) = \frac{p(x) - p(z)}{x - z}$, retornar $\pi = g^{q(\tau)}$

\textbf{Verify:} Comprobar $e(C - g^{p(z)}, g) = e(\pi, g^\tau - g^z)$
\end{definition}

\begin{center}
\begin{tabular}{ll}
\toprule
\textbf{Aspecto} & \textbf{KZG} \\
\midrule
Tamaño commitment & 48 bytes (1 elem. $\mathbb{G}_1$) \\
Tamaño prueba & 48 bytes \\
Tiempo commit & $O(n)$ multiplicaciones \\
Tiempo prove & $O(n)$ multiplicaciones \\
Tiempo verify & $O(1)$ (2 pairings) \\
Setup & Trusted (ceremony) \\
Asunción & q-SDH \\
Post-quantum & No \\
\bottomrule
\end{tabular}
\end{center}

\begin{keyidea}
\textbf{Fortalezas de KZG:}
\begin{itemize}
    \item Pruebas de tamaño constante (¡48 bytes para cualquier grado!)
    \item Verificación en tiempo constante
    \item Batch opening eficiente
    \item Homomórfico (se pueden combinar commitments)
\end{itemize}

\textbf{Debilidades:}
\begin{itemize}
    \item Requiere trusted setup
    \item No es post-quantum (vulnerable a computadores cuánticos)
    \item Setup es específico para grado máximo
\end{itemize}
\end{keyidea}

%-------------------------------------------
\subsubsection{IPA (Inner Product Argument)}

\begin{definition}[Inner Product Argument]
IPA es un polynomial commitment scheme basado en el ``inner product argument'' de Bulletproofs. No requiere trusted setup ni pairings.

\textbf{Idea central:} Reducir recursivamente un inner product de tamaño $n$ a uno de tamaño $n/2$, hasta llegar a tamaño 1.
\end{definition}

\begin{intuicion}
Imagina que quieres probar que conoces vectores $\vec{a}, \vec{b}$ tales que $\langle \vec{a}, \vec{b} \rangle = c$.

En lugar de enviar los vectores completos ($O(n)$ datos), IPA permite:
\begin{enumerate}
    \item Comprimir a vectores de tamaño $n/2$
    \item Repetir recursivamente $\log n$ veces
    \item Prueba final de tamaño $O(\log n)$
\end{enumerate}

Para polynomial commitments: representamos $p(x)$ como inner product de coeficientes con potencias de $x$.
\end{intuicion}

\begin{mathnote}
\textbf{El protocolo IPA (simplificado):}

\textbf{Setup:} Generadores $\vec{G} = (G_1, \ldots, G_n)$, $\vec{H} = (H_1, \ldots, H_n)$, y $Q$ (para binding).

\textbf{Commitment a polinomio $p(x) = \sum_i a_i x^i$:}
\[ C = \langle \vec{a}, \vec{G} \rangle = \sum_i a_i G_i \]

\textbf{Para probar evaluación $p(z) = y$:}

Sea $\vec{b} = (1, z, z^2, \ldots, z^{n-1})$. Entonces $p(z) = \langle \vec{a}, \vec{b} \rangle = y$.

\textbf{Reducción recursiva:}

Dividir $\vec{a} = (\vec{a}_L, \vec{a}_R)$ y $\vec{b} = (\vec{b}_L, \vec{b}_R)$ en mitades.

Calcular cross-terms:
\begin{align}
L &= \langle \vec{a}_L, \vec{G}_R \rangle + \langle \vec{a}_L, \vec{b}_R \rangle \cdot Q \\
R &= \langle \vec{a}_R, \vec{G}_L \rangle + \langle \vec{a}_R, \vec{b}_L \rangle \cdot Q
\end{align}

Verificador envía challenge aleatorio $x$.

Nuevos vectores:
\begin{align}
\vec{a}' &= \vec{a}_L + x \cdot \vec{a}_R \\
\vec{b}' &= \vec{b}_L + x^{-1} \cdot \vec{b}_R \\
\vec{G}' &= \vec{G}_L + x^{-1} \cdot \vec{G}_R
\end{align}

Repetir hasta $n = 1$.

\textbf{Prueba final:} $(L_1, R_1, L_2, R_2, \ldots, L_{\log n}, R_{\log n}, a_{\text{final}})$
\end{mathnote}

\begin{center}
\begin{tabular}{ll}
\toprule
\textbf{Aspecto} & \textbf{IPA} \\
\midrule
Tamaño commitment & 32 bytes (1 elem. grupo) \\
Tamaño prueba & $O(\log n)$ elementos \\
Tiempo commit & $O(n)$ \\
Tiempo prove & $O(n)$ \\
Tiempo verify & $O(n)$ (!!) \\
Setup & Ninguno (transparent) \\
Asunción & DLP \\
Post-quantum & No \\
\bottomrule
\end{tabular}
\end{center}

\begin{advertencia}
\textbf{El problema del tiempo de verificación:}

IPA tiene verificación en tiempo $O(n)$, no $O(\log n)$. Aunque la prueba es logarítmica, el verificador debe reconstruir los generadores $\vec{G}'$ en cada ronda.

\textbf{Soluciones:}
\begin{enumerate}
    \item Usar ``inner product argument'' para verificar múltiples evaluaciones en batch
    \item Halo-style accumulation: diferir la verificación costosa
    \item Precomputar tablas de generadores
\end{enumerate}
\end{advertencia}

\begin{keyidea}
\textbf{Fortalezas de IPA:}
\begin{itemize}
    \item Sin trusted setup (completamente transparent)
    \item Pruebas razonablemente pequeñas ($O(\log n)$)
    \item Simple de implementar
    \item Base para Bulletproofs y Halo
\end{itemize}

\textbf{Debilidades:}
\begin{itemize}
    \item Verificación $O(n)$ (costosa)
    \item No es post-quantum
    \item Menos eficiente que KZG para verificación individual
\end{itemize}
\end{keyidea}

\begin{example}[IPA en Halo 2]
Halo 2 (usado por Zcash y protocolos como Scroll) usa IPA con una técnica de ``accumulation'':

\begin{enumerate}
    \item En lugar de verificar completamente cada prueba, se ``acumula'' el trabajo
    \item Múltiples pruebas se combinan en un ``acumulador''
    \item Solo al final se hace la verificación costosa una vez
\end{enumerate}

Esto amortiza el costo $O(n)$ sobre muchas pruebas.
\end{example}

%-------------------------------------------
\subsubsection{FRI (Fast Reed-Solomon IOP)}

\begin{definition}[FRI]
Fast Reed-Solomon Interactive Oracle Proof (FRI) es un protocolo para probar que un commitment corresponde a la evaluación de un polinomio de bajo grado.

\textbf{Diferencia clave:} FRI trabaja con evaluaciones del polinomio sobre un dominio, no con coeficientes.
\end{definition}

\begin{intuicion}
FRI opera de forma diferente a KZG e IPA:

\textbf{KZG/IPA:} ``He aquí un commitment a los \textit{coeficientes} de mi polinomio''

\textbf{FRI:} ``He aquí las \textit{evaluaciones} de mi polinomio en muchos puntos. Puedo probar que son consistentes con un polinomio de grado $< d$''

FRI usa la propiedad de que los códigos Reed-Solomon tienen máxima distancia: un polinomio de grado $d$ evaluado en $n \gg d$ puntos tiene una ``firma'' única.
\end{intuicion}

\begin{mathnote}
\textbf{El protocolo FRI (simplificado):}

\textbf{Setup:} Campo $\mathbb{F}$, dominio de evaluación $D$ de tamaño $n = 2^k$, grado objetivo $d < n$.

\textbf{Commit a polinomio $p(x)$ de grado $< d$:}
\begin{enumerate}
    \item Evaluar $p$ en todo $D$: $(p(\omega^0), p(\omega^1), \ldots, p(\omega^{n-1}))$ donde $\omega$ es raíz de unidad
    \item Construir Merkle tree de las evaluaciones
    \item Commitment = raíz del Merkle tree
\end{enumerate}

\textbf{Protocolo de reducción (query phase):}

\textbf{Ronda 1:}
\begin{enumerate}
    \item Verificador envía challenge aleatorio $\alpha_1$
    \item Prover construye $p_1(x)$ de grado $< d/2$ tal que:
    \[ p(x) = p_1(x^2) + x \cdot p_1^{\text{odd}}(x^2) \]
    donde ``fold'' combina coeficientes pares e impares usando $\alpha_1$
    \item Prover envía commitment a evaluaciones de $p_1$
\end{enumerate}

\textbf{Rondas 2 a $\log(d)$:} Repetir, reduciendo grado a la mitad cada vez.

\textbf{Ronda final:} El polinomio tiene grado constante; prover lo envía directamente.

\textbf{Verificación (query phase):}
\begin{enumerate}
    \item Verificador elige puntos aleatorios $z_1, \ldots, z_\lambda$
    \item Para cada $z_i$, pide evaluaciones y Merkle proofs
    \item Verifica consistencia de las reducciones
\end{enumerate}
\end{mathnote}

\begin{center}
\begin{tikzpicture}[
    level/.style={draw, rectangle, minimum width=2.5cm, minimum height=0.8cm},
    arrow/.style={->, thick}
]
    % Niveles
    \node[level, fill=blue!20] (p0) at (0, 0) {$p(x)$, grado $< d$};
    \node[level, fill=blue!15] (p1) at (0, -1.5) {$p_1(x)$, grado $< d/2$};
    \node[level, fill=blue!10] (p2) at (0, -3) {$p_2(x)$, grado $< d/4$};
    \node at (0, -4) {$\vdots$};
    \node[level, fill=green!20] (pf) at (0, -5) {$p_{\log d}$, constante};
    
    % Flechas
    \draw[arrow] (p0) -- node[right] {fold con $\alpha_1$} (p1);
    \draw[arrow] (p1) -- node[right] {fold con $\alpha_2$} (p2);
    \draw[arrow] (0, -3.5) -- (pf);
    
    % Queries
    \node at (4, -2.5) {Queries en puntos aleatorios};
    \draw[dashed, ->] (2.5, -2.5) -- (1.5, 0);
    \draw[dashed, ->] (2.5, -2.5) -- (1.5, -1.5);
    \draw[dashed, ->] (2.5, -2.5) -- (1.5, -3);
\end{tikzpicture}
\end{center}

\begin{center}
\begin{tabular}{ll}
\toprule
\textbf{Aspecto} & \textbf{FRI} \\
\midrule
Tamaño commitment & $O(1)$ (raíz Merkle) \\
Tamaño prueba & $O(\lambda \cdot \log^2 n)$ \\
Tiempo commit & $O(n \log n)$ FFT \\
Tiempo prove & $O(n \log n)$ \\
Tiempo verify & $O(\lambda \cdot \log^2 n)$ \\
Setup & Ninguno (transparent) \\
Asunción & Hash collision resistance \\
Post-quantum & Sí (!) \\
\bottomrule
\end{tabular}
\end{center}

\begin{keyidea}
\textbf{Fortalezas de FRI:}
\begin{itemize}
    \item \textbf{Post-quantum}: Solo usa hashes, sobrevive a computadores cuánticos
    \item Sin trusted setup (completamente transparent)
    \item Verificación polylogarítmica
    \item Base para STARKs
\end{itemize}

\textbf{Debilidades:}
\begin{itemize}
    \item Pruebas más grandes que KZG ($\sim$100 KB vs 48 bytes)
    \item Más costoso computacionalmente
    \item Complejidad de implementación
\end{itemize}
\end{keyidea}

\begin{remark}
\textbf{STARKs = FRI + AIR:}

Los STARKs (Scalable Transparent ARguments of Knowledge) combinan:
\begin{itemize}
    \item \textbf{AIR} (Algebraic Intermediate Representation): El sistema de constraints
    \item \textbf{FRI}: El polynomial commitment scheme
\end{itemize}

Sistemas como StarkWare, Cairo, y Polygon Miden usan STARKs basados en FRI.
\end{remark}

%-------------------------------------------
\subsubsection{Comparación Detallada}

\begin{center}
\begin{tabular}{lccc}
\toprule
\textbf{Propiedad} & \textbf{KZG} & \textbf{IPA} & \textbf{FRI} \\
\midrule
Tamaño prueba & $O(1)$ & $O(\log n)$ & $O(\log^2 n)$ \\
Verificación & $O(1)$ & $O(n)$ & $O(\log^2 n)$ \\
Trusted setup & Sí & No & No \\
Post-quantum & No & No & Sí \\
Basado en & Pairings & DLP & Hashes \\
Homomórfico & Sí & Sí & No \\
Batch friendly & Muy sí & Sí & Sí \\
\bottomrule
\end{tabular}
\end{center}

\begin{mathnote}
\textbf{Tamaños concretos para $n = 2^{20}$ (1 millón de grado):}

\begin{center}
\begin{tabular}{lccc}
\toprule
& \textbf{KZG} & \textbf{IPA} & \textbf{FRI} \\
\midrule
Commitment & 48 B & 32 B & 32 B \\
Prueba (1 eval) & 48 B & 640 B & $\sim$50 KB \\
Tiempo verify & $\sim$2 ms & $\sim$500 ms & $\sim$50 ms \\
\bottomrule
\end{tabular}
\end{center}

\textbf{Nota:} Los tiempos son aproximados y dependen de la implementación y hardware.
\end{mathnote}

%-------------------------------------------
\subsubsection{¿Cuál Elegir?}

\begin{keyidea}
\textbf{Guía de selección:}

\textbf{Usar KZG cuando:}
\begin{itemize}
    \item El tamaño de prueba es crítico (ej: verificación on-chain)
    \item Se puede tolerar trusted setup (ceremonias disponibles)
    \item No se necesita seguridad post-quantum
    \item Ejemplo: PLONK, Ethereum data availability
\end{itemize}

\textbf{Usar IPA cuando:}
\begin{itemize}
    \item No se puede tolerar trusted setup
    \item Se usará accumulation/recursión (ej: Halo 2)
    \item Verificación individual no es el cuello de botella
    \item Ejemplo: Zcash Orchard, Halo 2
\end{itemize}

\textbf{Usar FRI cuando:}
\begin{itemize}
    \item Se necesita seguridad post-quantum
    \item No hay trusted setup disponible
    \item El tamaño de prueba no es crítico ($\sim$100 KB aceptable)
    \item Ejemplo: STARKs, zkVMs como Cairo
\end{itemize}
\end{keyidea}

%-------------------------------------------
\subsubsection{Otros Polynomial Commitment Schemes}

\begin{remark}
Además de KZG, IPA, y FRI, existen otros PCS:

\textbf{DARK (2019):}
\begin{itemize}
    \item Basado en grupos de clase (class groups)
    \item Sin trusted setup
    \item Teóricamente interesante, poco usado en práctica
\end{itemize}

\textbf{Dory (2020):}
\begin{itemize}
    \item Combina IPA con pairings
    \item Verificación $O(\log n)$ (mejor que IPA puro)
    \item Sin trusted setup
\end{itemize}

\textbf{Brakedown (2021):}
\begin{itemize}
    \item Basado en códigos de corrección de errores
    \item Sin trusted setup, eficiente para el prover
    \item Usado en algunos zkVMs
\end{itemize}

\textbf{Gemini (2022):}
\begin{itemize}
    \item Split accumulation scheme
    \item Combina ventajas de varios enfoques
\end{itemize}
\end{remark}

%-------------------------------------------
\subsubsection{Multilinear Polynomial Commitments}

\begin{definition}[Polinomios Multilineales]
Un polinomio multilineal en $n$ variables tiene la forma:
\[ f(x_1, \ldots, x_n) = \sum_{b \in \{0,1\}^n} c_b \cdot x_1^{b_1} \cdots x_n^{b_n} \]

donde cada variable aparece con grado $\leq 1$.
\end{definition}

\begin{intuicion}
Los polinomios multilineales son naturales para circuitos:
\begin{itemize}
    \item Un circuito de $n$ wires se representa con polinomio de $\log n$ variables
    \item Cada ``hypercube evaluation'' corresponde a un wire
    \item La estructura permite verificación eficiente vía sumcheck
\end{itemize}
\end{intuicion}

\begin{remark}
\textbf{PCS para multilineales:}
\begin{itemize}
    \item \textbf{Zeromorph}: Adapta KZG para multilineales
    \item \textbf{Binius}: Basado en campos binarios, muy eficiente
    \item \textbf{HyperKZG}: KZG para hypercube
\end{itemize}

Sistemas como Lasso y Jolt usan multilinear PCS para lookup arguments eficientes.
\end{remark}

%-------------------------------------------
\subsubsection{Aplicación en Banco ZK}

\begin{bancozk}
El PCS es un componente fundamental de nuestro sistema de pruebas:

\textbf{1. Arquitectura actual: PLONK con KZG}

Usamos PLONK con KZG como PCS porque:
\begin{itemize}
    \item Pruebas pequeñas ($\sim$700 bytes totales)
    \item Verificación on-chain eficiente ($\sim$200K gas)
    \item Ceremony de Ethereum disponible
\end{itemize}

\begin{verbatim}
Circuito de transacción → PLONK → Polinomios → KZG commits
                                                   ↓
                                    Prueba final: ~700 bytes
                                    - 7 elementos G1 (KZG commits)
                                    - 2 elementos G1 (KZG opening proofs)
                                    - Evaluaciones de campo
\end{verbatim}

\textbf{2. Trade-offs considerados}

\begin{center}
\begin{tabular}{lccc}
\toprule
\textbf{Opción} & \textbf{Prueba} & \textbf{Gas} & \textbf{Setup} \\
\midrule
PLONK + KZG & $\sim$700 B & $\sim$200K & Trusted \\
Halo 2 + IPA & $\sim$5 KB & $\sim$2M & None \\
STARK + FRI & $\sim$100 KB & $\sim$5M & None \\
\bottomrule
\end{tabular}
\end{center}

Para L2 con verificación on-chain, el gas domina. KZG gana claramente.

\textbf{3. Recursión con diferentes PCS}

Nuestro sistema de recursión combina:
\begin{verbatim}
Nivel 0: Pruebas individuales (KZG, rápido de generar)
    ↓
Nivel 1: Agregación (IPA/Halo style, sin trusted setup interno)
    ↓
Nivel 2: Prueba final (KZG para verificación on-chain)
\end{verbatim}

\textbf{4. Consideraciones de seguridad}

\textbf{Trusted setup:}
\begin{itemize}
    \item Usamos el SRS de la Ethereum KZG Ceremony (141,000+ participantes)
    \item Si 1 participante fue honesto, el setup es seguro
    \item Verificable públicamente
\end{itemize}

\textbf{Post-quantum:}
\begin{itemize}
    \item KZG NO es post-quantum seguro
    \item Plan de migración: cuando amenaza cuántica sea realista, migrar a FRI
    \item Las pruebas tienen ``temporal validity'': solo importa que sean seguras hoy
\end{itemize}

\textbf{5. Optimizaciones específicas}

\begin{verbatim}
// Batch opening en KZG para múltiples evaluaciones
// En lugar de k proofs separados, uno combinado

fn batch_open(
    commitments: [C_1, ..., C_k],
    points: [z_1, ..., z_k], 
    values: [y_1, ..., y_k]
) -> BatchProof {
    // 1. Combinar polinomios con challenges aleatorios
    // 2. Single KZG opening proof
    // Tamaño: O(1) en lugar de O(k)
}
\end{verbatim}

\textbf{6. Métricas de producción}

\begin{center}
\begin{tabular}{lc}
\toprule
\textbf{Métrica} & \textbf{Valor} \\
\midrule
Grado del polinomio (circuito típico) & $\sim 2^{18}$ \\
KZG commits por prueba & 7 \\
Tiempo de commit & $\sim$500 ms \\
Tiempo de prove total & $\sim$3 s \\
Tamaño de prueba & 704 bytes \\
Gas de verificación & 213,000 \\
\bottomrule
\end{tabular}
\end{center}

\textbf{7. Roadmap futuro}

\begin{itemize}
    \item \textbf{Corto plazo:} Optimizar MSM para commits más rápidos
    \item \textbf{Medio plazo:} Explorar Gemini/Nova para mejor recursión
    \item \textbf{Largo plazo:} Preparar migración a FRI cuando post-quantum sea necesario
\end{itemize}
\end{bancozk}

%-------------------------------------------
\subsubsection{Resumen}

\begin{keyidea}
\textbf{Polynomial Commitments: Puntos clave}

\begin{enumerate}
    \item \textbf{El por qué:} SNARKs modernos reducen todo a polinomios. PCS permite verificar sin revelar.
    
    \item \textbf{KZG:} Pruebas $O(1)$, verificación $O(1)$, pero trusted setup. Ideal para on-chain.
    
    \item \textbf{IPA:} Sin trusted setup, pruebas $O(\log n)$, pero verificación $O(n)$. Bueno para recursión.
    
    \item \textbf{FRI:} Post-quantum, sin setup, pero pruebas grandes. Base para STARKs.
    
    \item \textbf{Trade-off fundamental:} Tamaño de prueba vs trusted setup vs post-quantum.
    
    \item \textbf{En la práctica:} Muchos sistemas usan KZG para eficiencia, con plan de migración a FRI.
\end{enumerate}
\end{keyidea}

\cleardoublepage

% ============================================
% BLOQUE 2: BLOCKCHAIN BASE (L1) - VERSIÓN EXPANDIDA
% Parte 1: Qué es L1 + Consenso
% ============================================
\chapter{Blockchain Base (L1)}
\label{ch:bloque2}

\begin{center}
\textit{Bloque 2 de 20}
\end{center}

%===========================================
\section{Qué es L1}
%===========================================

\begin{definition}[Layer 1]
Blockchain base que proporciona:
\begin{itemize}
    \item Consenso sobre el estado
    \item Disponibilidad de datos
    \item Ejecución de transacciones
    \item Settlement final
\end{itemize}
\end{definition}

Ejemplos: Bitcoin, Ethereum, Solana, Avalanche, Polkadot.

%-------------------------------------------
\subsection{Contexto Histórico}

\begin{intuicion}
Layer 1 es la ``capa base'' de una blockchain: el fundamento sobre el cual todo lo demás se construye. Es como el sistema operativo de una computadora---todo lo demás (aplicaciones, protocolos, tokens) depende de que L1 funcione correctamente.

Cuando hablamos de ``heredar seguridad de L1'', significa que un protocolo puede confiar en las garantías que L1 proporciona sin tener que reimplementarlas.
\end{intuicion}

\begin{remark}
\textbf{La evolución de las blockchains L1:}

\begin{center}
\begin{tabular}{lll}
\toprule
\textbf{Generación} & \textbf{Años} & \textbf{Características} \\
\midrule
1ª Gen & 2009-2013 & Solo transferencias (Bitcoin) \\
2ª Gen & 2015-2017 & Smart contracts (Ethereum) \\
3ª Gen & 2018-2020 & Escalabilidad nativa (Solana, Avalanche) \\
4ª Gen & 2021+ & Modular, ZK-native (Celestia, Mina) \\
\bottomrule
\end{tabular}
\end{center}

Cada generación intentó resolver las limitaciones de la anterior, pero ninguna ha logrado satisfacer completamente el ``trilema'' de escalabilidad, descentralización y seguridad.
\end{remark}

%-------------------------------------------
\subsection{Las Cuatro Funciones Fundamentales}

\begin{definition}[Las Cuatro Funciones de L1]
\textbf{1. Consenso sobre el estado:}
Todos los nodos acuerdan cuál es el estado actual del sistema. En Ethereum, esto incluye todos los balances, código de contratos, y storage.

\textbf{2. Disponibilidad de datos (DA):}
Los datos necesarios para verificar transacciones están disponibles para cualquiera que los necesite. Sin DA, no se puede verificar que el estado es correcto.

\textbf{3. Ejecución de transacciones:}
Procesar transacciones y calcular el nuevo estado. En Ethereum, esto significa ejecutar bytecode EVM.

\textbf{4. Settlement final:}
Una vez que una transacción está ``finalizada'', no puede ser revertida. Diferentes L1s tienen diferentes tiempos y garantías de finalidad.
\end{definition}

\begin{mathnote}
\textbf{Formalización del estado:}

El estado de una blockchain en el bloque $n$ se puede expresar como:
\[ S_n = f(S_{n-1}, B_n) \]

donde:
\begin{itemize}
    \item $S_n$ es el estado después del bloque $n$
    \item $S_{n-1}$ es el estado anterior
    \item $B_n$ es el bloque $n$ (con sus transacciones)
    \item $f$ es la función de transición de estado (STF)
\end{itemize}

El estado génesis $S_0$ está definido en el protocolo.

\textbf{En Ethereum:}
\[ \sigma: \text{Address} \to (\text{nonce}, \text{balance}, \text{storageRoot}, \text{codeHash}) \]

La función de transición ejecuta cada transacción secuencialmente:
\[ S_n = f(f(f(S_{n-1}, \text{tx}_1), \text{tx}_2), \ldots, \text{tx}_k) \]

donde $B_n = \{\text{tx}_1, \text{tx}_2, \ldots, \text{tx}_k\}$.
\end{mathnote}

\begin{keyidea}
\textbf{Monolítico vs Modular:}

\textbf{L1 Monolítico} (Bitcoin, Ethereum, Solana):
\begin{itemize}
    \item Una sola cadena hace todo: consenso, DA, ejecución, settlement
    \item Más simple pero menos escalable
    \item Trade-offs fijos
\end{itemize}

\textbf{L1 Modular} (Celestia + rollups):
\begin{itemize}
    \item Diferentes capas especializadas
    \item Celestia: solo DA y consenso
    \item Rollups: ejecución
    \item Ethereum: settlement
    \item Más complejo pero más escalable y flexible
\end{itemize}

\begin{center}
\begin{tikzpicture}[
    box/.style={draw, rectangle, minimum width=2.5cm, minimum height=0.8cm, align=center},
    arrow/.style={->, thick}
]
    % Monolítico
    \node[box, fill=blue!30] (mono) at (0, 0) {L1 Monolítico\\(Ethereum)};
    \node[below=0.3cm of mono, font=\small] {Consenso + DA +\\Ejecución + Settlement};
    
    % Modular
    \node[box, fill=green!20] (da) at (6, 1) {Celestia\\(DA + Consenso)};
    \node[box, fill=yellow!20] (exec) at (6, 0) {Rollup\\(Ejecución)};
    \node[box, fill=red!20] (settle) at (6, -1) {Ethereum\\(Settlement)};
    
    \draw[arrow] (da) -- (exec);
    \draw[arrow] (exec) -- (settle);
    
    \node at (3, -2) {vs};
\end{tikzpicture}
\end{center}
\end{keyidea}

%-------------------------------------------
\subsection{Comparación de L1s Principales}

\begin{center}
\begin{tabular}{lccccc}
\toprule
\textbf{L1} & \textbf{Consenso} & \textbf{TPS} & \textbf{Finalidad} & \textbf{VM} & \textbf{Modelo} \\
\midrule
Bitcoin & PoW & $\sim$7 & $\sim$60 min & Script & UTXO \\
Ethereum & PoS & $\sim$15 & $\sim$15 min & EVM & Account \\
Solana & PoH+PoS & $\sim$4000 & $\sim$13 s & SVM & Account \\
Avalanche & Avalanche & $\sim$4500 & $\sim$2 s & EVM & Account \\
Polkadot & NPoS & $\sim$1000 & $\sim$60 s & WASM & Account \\
Cardano & Ouroboros & $\sim$250 & $\sim$20 min & Plutus & eUTXO \\
Near & Nightshade & $\sim$1000 & $\sim$2 s & WASM & Account \\
Aptos & AptosBFT & $\sim$10000 & $\sim$1 s & Move VM & Account \\
\bottomrule
\end{tabular}
\end{center}

\begin{advertencia}
\textbf{TPS (Transactions Per Second) es una métrica engañosa:}

\begin{itemize}
    \item Diferentes transacciones tienen diferentes complejidades
    \item Una transferencia simple $\neq$ un swap en DEX $\neq$ deploy de contrato
    \item Condiciones de prueba vs producción real difieren enormemente
    \item No captura la descentralización ni la seguridad
\end{itemize}

\textbf{Métricas más útiles:}
\begin{itemize}
    \item Gas/compute units por segundo
    \item Costo por transacción
    \item Número de validadores independientes
    \item Nakamoto coefficient (cuántas entidades controlan 51\%)
    \item Tiempo a finalidad económica
\end{itemize}
\end{advertencia}

\begin{example}[Nakamoto Coefficient]
El coeficiente de Nakamoto mide cuántas entidades independientes necesitan coludirse para atacar la red:

\begin{center}
\begin{tabular}{lc}
\toprule
\textbf{Red} & \textbf{Nakamoto Coef.} \\
\midrule
Bitcoin & $\sim$4 (mining pools) \\
Ethereum & $\sim$2-3 (Lido + exchanges) \\
Solana & $\sim$19 \\
Avalanche & $\sim$26 \\
Cosmos Hub & $\sim$7 \\
\bottomrule
\end{tabular}
\end{center}

Un coeficiente más alto indica mayor descentralización.
\end{example}

%===========================================
\section{Consenso}
%===========================================

\begin{intuicion}
El consenso es el mecanismo por el cual miles de computadoras independientes, que no confían entre sí, acuerdan una única verdad. Es el corazón de cualquier blockchain.

Sin consenso, cada nodo tendría su propia versión de la ``verdad'' y no podríamos tener un sistema coherente.

El problema fundamental: en un sistema distribuido abierto, ¿cómo decidimos qué transacciones son válidas y en qué orden ocurrieron?
\end{intuicion}

\begin{mathnote}
\textbf{El Problema de los Generales Bizantinos (1982):}

Lamport, Shostak y Pease formalizaron el problema:
\begin{itemize}
    \item $n$ generales deben acordar atacar o retirarse
    \item Algunos generales pueden ser traidores
    \item Comunicación solo por mensajeros (pueden ser interceptados)
\end{itemize}

\textbf{Resultado:} Con $f$ traidores, se necesitan $n \geq 3f + 1$ generales para consenso.

En blockchain:
\begin{itemize}
    \item Generales = nodos validadores
    \item Traidores = nodos maliciosos o caídos
    \item Mensaje = bloques y votos
\end{itemize}
\end{mathnote}

%-------------------------------------------
\subsection{Proof of Work (PoW)}

\begin{definition}[PoW]
Para crear bloque, encontrar nonce tal que:
\[
H(\text{header} \| \text{nonce}) < \text{target}
\]
\end{definition}

\begin{itemize}
    \item Dificultad ajustada para tiempo promedio constante
    \item Seguridad: requiere 51\% del hashpower para atacar
    \item Consumo energético proporcional a seguridad
    \item Usado en: Bitcoin, Litecoin, (Ethereum pre-merge)
\end{itemize}

\subsubsection{Historia y Orígenes}

\begin{remark}
\textbf{La genealogía de Proof of Work:}

\begin{center}
\begin{tabular}{lll}
\toprule
\textbf{Año} & \textbf{Autor} & \textbf{Contribución} \\
\midrule
1992 & Dwork \& Naor & PoW para combatir spam de email \\
1997 & Adam Back & Hashcash: PoW práctico \\
2004 & Hal Finney & RPOW: Reusable Proofs of Work \\
2005 & Nick Szabo & Bit Gold: proto-Bitcoin \\
2008 & Satoshi Nakamoto & Bitcoin: PoW + incentivos + blockchain \\
2009 & --- & Bitcoin mainnet launch \\
2011 & --- & Primeras pools de minería \\
2013 & --- & Primeros ASICs para Bitcoin \\
2015 & Ethereum & PoW con Ethash (memory-hard) \\
2022 & Ethereum & The Merge: fin de PoW en Ethereum \\
\bottomrule
\end{tabular}
\end{center}

La innovación de Satoshi no fue PoW en sí, sino combinarlo con:
\begin{enumerate}
    \item Incentivos económicos (block rewards)
    \item Una cadena de bloques con referencia al anterior
    \item Ajuste dinámico de dificultad
    \item Un mecanismo para resolver forks (cadena más larga)
\end{enumerate}
\end{remark}

\subsubsection{El Algoritmo de Minería}

\begin{definition}[Algoritmo de Minería PoW]
\begin{verbatim}
def mine_block(prev_hash, transactions, difficulty_target):
    """
    Minar un bloque válido.
    
    Args:
        prev_hash: Hash del bloque anterior
        transactions: Lista de transacciones a incluir
        difficulty_target: Número máximo que el hash puede ser
    
    Returns:
        (block_header, nonce, block_hash)
    """
    # Construir header
    merkle_root = compute_merkle_root(transactions)
    timestamp = current_time()
    
    header = {
        'version': CURRENT_VERSION,
        'prev_hash': prev_hash,
        'merkle_root': merkle_root,
        'timestamp': timestamp,
        'bits': encode_target(difficulty_target),
        'nonce': 0
    }
    
    # Buscar nonce válido
    nonce = 0
    extra_nonce = 0  # En coinbase transaction
    
    while True:
        # Doble SHA-256 (Bitcoin)
        header['nonce'] = nonce
        header_bytes = serialize(header)
        block_hash = sha256(sha256(header_bytes))
        
        # ¿Es válido?
        if int(block_hash, 16) < difficulty_target:
            return header, nonce, block_hash
        
        nonce += 1
        
        # Nonce es solo 32 bits, puede agotarse
        if nonce > 0xFFFFFFFF:
            # Modificar extra_nonce en coinbase y recalcular merkle
            extra_nonce += 1
            transactions[0] = update_coinbase(transactions[0], extra_nonce)
            merkle_root = compute_merkle_root(transactions)
            header['merkle_root'] = merkle_root
            nonce = 0
\end{verbatim}
\end{definition}

\begin{mathnote}
\textbf{Probabilidad y Dificultad:}

El espacio de hashes SHA-256 es $\{0, 1, \ldots, 2^{256}-1\}$.

El target determina cuántos hashes son ``válidos'':
\[ P(\text{hash válido}) = \frac{\text{target}}{2^{256}} \]

\textbf{Ejemplo para Bitcoin (2024):}

Con dificultad $D \approx 80 \times 10^{12}$:
\[ \text{target} = \frac{\text{target}_{\text{max}}}{D} \approx \frac{2^{224}}{80 \times 10^{12}} \approx 2^{177} \]

\[ P(\text{válido}) = \frac{2^{177}}{2^{256}} = 2^{-79} \approx 1.6 \times 10^{-24} \]

Número esperado de hashes para encontrar bloque:
\[ E[\text{hashes}] = \frac{1}{P} = 2^{79} \approx 6 \times 10^{23} \]

Con hashrate global de $\sim$500 EH/s = $5 \times 10^{20}$ hashes/segundo:
\[ \text{Tiempo esperado} = \frac{6 \times 10^{23}}{5 \times 10^{20}} = 1200 \text{ segundos} \approx 20 \text{ minutos} \]

Pero como hay muchos mineros compitiendo, en promedio se encuentra un bloque cada $\sim$10 minutos.
\end{mathnote}

\subsubsection{Ajuste de Dificultad}

\begin{definition}[Ajuste de Dificultad de Bitcoin]
Cada 2016 bloques ($\sim$2 semanas si cada bloque toma 10 min):
\[ \text{nuevo\_target} = \text{viejo\_target} \times \frac{\text{tiempo\_real}}{2016 \times 10 \text{ min}} \]

Con límites:
\begin{itemize}
    \item Máximo incremento: $\times 4$ (si tiempo fue muy corto)
    \item Máximo decremento: $\div 4$ (si tiempo fue muy largo)
\end{itemize}
\end{definition}

\begin{example}[Cálculo de ajuste]
Supongamos que los últimos 2016 bloques tomaron 12 días en lugar de 14:

\[ \text{factor} = \frac{12 \text{ días}}{14 \text{ días}} = \frac{12}{14} \approx 0.857 \]

El nuevo target será:
\[ \text{nuevo\_target} = \text{viejo\_target} \times 0.857 \]

Como el target es más pequeño, la dificultad aumenta $\sim$17\%. Los bloques estaban llegando muy rápido porque el hashrate aumentó.
\end{example}

\begin{center}
\begin{tikzpicture}
    \begin{axis}[
        width=12cm,
        height=6cm,
        xlabel={Año},
        ylabel={Dificultad (log scale)},
        ymode=log,
        grid=major,
        legend pos=north west
    ]
    \addplot[blue, thick, mark=none] coordinates {
        (2009, 1)
        (2011, 1e6)
        (2013, 1e9)
        (2015, 5e10)
        (2017, 1e12)
        (2019, 1e13)
        (2021, 2e13)
        (2023, 5e13)
        (2024, 8e13)
    };
    \addlegendentry{Bitcoin Difficulty}
    \end{axis}
\end{tikzpicture}
\end{center}

\subsubsection{Seguridad de PoW}

\begin{keyidea}
\textbf{El Ataque del 51\%:}

Un atacante con más del 50\% del hashpower puede:
\begin{enumerate}
    \item Minar una cadena privada en secreto
    \item Gastar monedas en la cadena pública
    \item Cuando la cadena privada es más larga, publicarla
    \item La red adopta la cadena más larga (Nakamoto consensus)
    \item Las transacciones en la cadena abandonada desaparecen
    \item El atacante tiene sus monedas de vuelta (doble gasto)
\end{enumerate}
\end{keyidea}

\begin{mathnote}
\textbf{Probabilidad de éxito del ataque:}

Sea $q$ la fracción de hashpower del atacante y $p = 1 - q$ la fracción honesta.

Si el atacante está $z$ bloques detrás, la probabilidad de alcanzar a la cadena honesta es:
\[
P(\text{catch up}) = \begin{cases}
1 & \text{si } q > 0.5 \\
\left(\frac{q}{p}\right)^z & \text{si } q \leq 0.5
\end{cases}
\]

\textbf{Ejemplo:} Con $q = 0.3$ (30\% del hashpower) y $z = 6$ confirmaciones:
\[ P = \left(\frac{0.3}{0.7}\right)^6 = 0.43^6 \approx 0.006 = 0.6\% \]

Por eso Bitcoin recomienda esperar 6 confirmaciones para transacciones grandes.
\end{mathnote}

\begin{advertencia}
\textbf{Costo de un ataque del 51\% a Bitcoin (2024):}

\begin{center}
\begin{tabular}{lr}
\toprule
\textbf{Componente} & \textbf{Costo estimado} \\
\midrule
Hardware (>50\% hashrate) & $>\$10$ mil millones \\
Electricidad por hora & $\sim$\$2 millones \\
Infraestructura & $>\$100$ millones \\
Tiempo de adquisición & Meses a años \\
\bottomrule
\end{tabular}
\end{center}

Además, un ataque exitoso destruiría la confianza en Bitcoin, colapsando el precio y haciendo que el atacante pierda el valor de su inversión en hardware.

\textbf{Ataques exitosos históricos:}

\begin{center}
\begin{tabular}{llrr}
\toprule
\textbf{Moneda} & \textbf{Fecha} & \textbf{Pérdida} & \textbf{Confirmaciones} \\
\midrule
Ethereum Classic & Ene 2019 & \$1.1M & 100+ \\
Bitcoin Gold & May 2018 & \$18M & 22 \\
Verge & Abr 2018 & \$1.7M & -- \\
Bitcoin SV & Ago 2021 & -- & 14 \\
\bottomrule
\end{tabular}
\end{center}

Las monedas con bajo hashrate son vulnerables. Se puede ``alquilar'' hashpower en NiceHash.
\end{advertencia}

\subsubsection{Selfish Mining}

\begin{definition}[Selfish Mining]
Estrategia donde un minero con $>$25-33\% del hashpower puede obtener más recompensas que su proporción justa:
\begin{enumerate}
    \item Minero encuentra bloque pero NO lo publica
    \item Sigue minando en secreto
    \item Cuando la red pública está cerca de alcanzarlo, publica su cadena
    \item Red adopta cadena más larga, desperdiciando trabajo de otros
\end{enumerate}
\end{definition}

\begin{mathnote}
El paper de Eyal \& Sirer (2014) demostró que con $\alpha$ = fracción de hashpower y $\gamma$ = fracción de la red que el atacante alcanza primero con su bloque:

El selfish mining es rentable cuando:
\[ \alpha > \frac{1 - \gamma}{3 - 2\gamma} \]

Con $\gamma = 0$: rentable si $\alpha > 0.333$ (33\%)
Con $\gamma = 0.5$: rentable si $\alpha > 0.25$ (25\%)
\end{mathnote}

\subsubsection{Algoritmos de Hash para PoW}

\begin{center}
\begin{tabular}{llll}
\toprule
\textbf{Algoritmo} & \textbf{Usado en} & \textbf{Tipo} & \textbf{Hardware dominante} \\
\midrule
SHA-256d & Bitcoin, BCH & CPU-friendly & ASICs \\
Scrypt & Litecoin, Doge & Memory-hard & ASICs (ahora) \\
Ethash & Ethereum (era) & Memory-hard DAG & GPUs \\
RandomX & Monero & CPU-optimizado & CPUs \\
Equihash & Zcash & Memory-hard & ASICs (moderado) \\
X11 & Dash & 11 algoritmos & ASICs \\
CryptoNight & Monero (old) & Memory-hard & CPUs/GPUs \\
\bottomrule
\end{tabular}
\end{center}

\begin{remark}
\textbf{La carrera armamentista del hardware de minería:}

\begin{enumerate}
    \item \textbf{2009-2010:} Minería con CPU (cualquier computadora)
    \item \textbf{2010-2012:} GPUs (10-100x más eficientes)
    \item \textbf{2013+:} ASICs (Application-Specific Integrated Circuits)
    \item ASICs son 100,000x+ más eficientes que CPUs para SHA-256
\end{enumerate}

\textbf{Trade-off:}
\begin{itemize}
    \item ASICs: máxima eficiencia, pero centraliza minería
    \item Memory-hard: resiste ASICs, más descentralizado
    \item CPU-only: máxima descentralización, pero botnets son problema
\end{itemize}
\end{remark}

\subsubsection{Pools de Minería}

\begin{definition}[Mining Pool]
Grupo de mineros que combinan su hashpower y comparten recompensas proporcionalmente.
\end{definition}

\begin{intuicion}
Un minero pequeño con 0.001\% del hashrate encontraría un bloque cada $\sim$1000 días en promedio. Con alta varianza, podría ser 5 años o 1 mes.

En un pool, recibe pequeños pagos constantes proporcionales a su contribución. Menos varianza, mismo valor esperado (menos fee del pool).
\end{intuicion}

\textbf{Métodos de pago:}

\begin{center}
\begin{tabular}{ll}
\toprule
\textbf{Método} & \textbf{Descripción} \\
\midrule
PPS (Pay Per Share) & Pago fijo por cada share válido \\
PPLNS & Pago proporcional a shares en últimos N shares \\
FPPS & PPS + estimación de fees de transacción \\
SOLO & Solo ayuda a encontrar bloque, toda la recompensa \\
\bottomrule
\end{tabular}
\end{center}

\begin{advertencia}
\textbf{Distribución de hashpower Bitcoin (2024):}

\begin{center}
\begin{tabular}{lr}
\toprule
\textbf{Pool} & \textbf{\% Hashrate} \\
\midrule
Foundry USA & $\sim$27\% \\
AntPool & $\sim$16\% \\
F2Pool & $\sim$13\% \\
ViaBTC & $\sim$12\% \\
Binance Pool & $\sim$8\% \\
Otros & $\sim$24\% \\
\bottomrule
\end{tabular}
\end{center}

Los top 4 pools controlan $>$65\% del hashrate. Esto es preocupante para la descentralización, aunque los mineros individuales dentro de los pools pueden cambiar fácilmente.
\end{advertencia}

%-------------------------------------------
\subsection{Proof of Stake (PoS)}

\begin{definition}[PoS]
Validadores depositan stake como colateral. Seleccionados para proponer bloques según su stake.
\end{definition}

Mecanismos:
\begin{itemize}
    \item \textbf{Selección:} Aleatoria ponderada por stake
    \item \textbf{Slashing:} Perder stake por mal comportamiento
    \item \textbf{Rewards:} Ganar por participación honesta
\end{itemize}

Variantes:
\begin{itemize}
    \item \textbf{Delegated PoS:} Usuarios delegan a validadores
    \item \textbf{Nominated PoS:} Polkadot, nominadores eligen
    \item \textbf{Liquid staking:} stETH, stake tokenizado
\end{itemize}

\subsubsection{Historia y Motivación}

\begin{remark}
\textbf{¿Por qué PoS?}

\begin{enumerate}
    \item \textbf{Eficiencia energética:} PoW consume tanta electricidad como países enteros
    \item \textbf{Descentralización de hardware:} No requiere ASICs especializados
    \item \textbf{Seguridad económica:} Atacante pierde su propio stake
    \item \textbf{Finalidad más rápida:} BFT puede dar finalidad en segundos
    \item \textbf{Menor barrera de entrada:} Cualquiera con stake puede participar
\end{enumerate}

\textbf{Timeline de PoS:}
\begin{center}
\begin{tabular}{ll}
\toprule
\textbf{Año} & \textbf{Evento} \\
\midrule
2012 & Peercoin: primer PoS híbrido \\
2014 & Primeras propuestas formales de PoS \\
2017 & Cardano Ouroboros (primer PoS con prueba de seguridad) \\
2019 & Cosmos Hub launch (Tendermint BFT) \\
2020 & Ethereum Beacon Chain launch (PoS) \\
2020 & Solana mainnet (PoH + PoS) \\
2022 & Ethereum ``The Merge'': transición completa a PoS \\
\bottomrule
\end{tabular}
\end{center}
\end{remark}

\subsubsection{El Problema ``Nothing at Stake''}

\begin{advertencia}
\textbf{Nothing at Stake:}

En PoW, minar en múltiples forks divide tu hashpower (costo real).

En PoS naive, votar en múltiples forks no tiene costo---puedes votar en todos ``por si acaso''.

\textbf{Problema:} Si todos votan en todos los forks, nunca se alcanza consenso.

\textbf{Solución: Slashing}
\begin{itemize}
    \item Si se detecta que votaste en múltiples forks conflictivos
    \item Pierdes parte o todo tu stake
    \item Ahora hay un costo real por comportamiento deshonesto
\end{itemize}
\end{advertencia}

\subsubsection{Ethereum PoS (Gasper)}

\begin{definition}[Gasper]
El protocolo de consenso de Ethereum combina:
\begin{itemize}
    \item \textbf{Casper FFG:} Mecanismo de finalidad
    \item \textbf{LMD-GHOST:} Regla de elección de fork
\end{itemize}
\end{definition}

\begin{mathnote}
\textbf{Estructura de Ethereum PoS:}

\textbf{Slots y Epochs:}
\begin{itemize}
    \item 1 slot = 12 segundos
    \item 1 epoch = 32 slots = 6.4 minutos
    \item Cada slot tiene un proposer designado
    \item Validadores divididos en comités por slot
\end{itemize}

\textbf{Roles:}
\begin{itemize}
    \item \textbf{Proposer:} Propone bloque en su slot
    \item \textbf{Attester:} Vota por el bloque que considera correcto
    \item \textbf{Sync Committee:} 512 validadores que firman headers (para light clients)
\end{itemize}

\textbf{Requisitos:}
\begin{itemize}
    \item Mínimo stake: 32 ETH
    \item Tiempo de activación: $\sim$1 día (cola de entrada)
    \item Tiempo de salida: $\sim$1 día + período de retiro
\end{itemize}
\end{mathnote}

\begin{definition}[Slashing Conditions en Ethereum]
Un validador es slashed si:
\begin{enumerate}
    \item \textbf{Double voting:} Votar por dos bloques diferentes en el mismo slot
    \item \textbf{Surround voting:} Hacer attestation que ``rodea'' otra attestation propia
\end{enumerate}

\textbf{Penalidades:}
\begin{itemize}
    \item Mínimo: 1/32 del stake ($\sim$1 ETH)
    \item Correlación: Si muchos son slashed juntos, penalidad aumenta hasta 100\%
    \item Exclusión forzosa de la red
\end{itemize}
\end{definition}

\begin{example}[Cálculo de recompensas de validador]
\textbf{Parámetros (circa 2024):}
\begin{itemize}
    \item Total staked: $\sim$30 millones ETH
    \item Base reward por epoch: $\sim$0.000003 ETH
    \item Epochs por año: $\sim$82,125
\end{itemize}

\textbf{APR aproximado:}
\[ \text{APR} \approx \frac{C}{\sqrt{\text{total\_staked}}} \]

Con 30M ETH staked: APR $\approx$ 4-5\%

Para un validador con 32 ETH:
\[ \text{Ganancia anual} \approx 32 \times 0.045 = 1.44 \text{ ETH} \]

Más fees de transacción y MEV (variable).
\end{example}

\subsubsection{Casper FFG: Finalidad}

\begin{definition}[Finalidad en Casper FFG]
Un bloque está \textbf{finalizado} cuando:
\begin{enumerate}
    \item Es \textbf{justified}: $>$2/3 del stake votó por él
    \item El siguiente checkpoint también es justified, con él como source
\end{enumerate}

Un bloque finalizado NUNCA puede ser revertido (salvo que $>$1/3 del stake sea slashed).
\end{definition}

\begin{center}
\begin{tikzpicture}[
    block/.style={draw, rectangle, minimum width=1.5cm, minimum height=0.8cm},
    checkpoint/.style={draw, rectangle, minimum width=1.5cm, minimum height=0.8cm, fill=blue!20},
    finalized/.style={draw, rectangle, minimum width=1.5cm, minimum height=0.8cm, fill=green!30}
]
    % Bloques
    \node[finalized] (b0) at (0, 0) {Epoch 0};
    \node[finalized] (b1) at (2.5, 0) {Epoch 1};
    \node[checkpoint] (b2) at (5, 0) {Epoch 2};
    \node[block] (b3) at (7.5, 0) {Epoch 3};
    
    % Flechas
    \draw[->] (b0) -- (b1);
    \draw[->] (b1) -- (b2);
    \draw[->] (b2) -- (b3);
    
    % Labels
    \node[below=0.3cm of b0, font=\small] {Finalizado};
    \node[below=0.3cm of b1, font=\small] {Finalizado};
    \node[below=0.3cm of b2, font=\small] {Justified};
    \node[below=0.3cm of b3, font=\small] {Propuesto};
    
    % Votos
    \draw[->, dashed, blue] (b1) to[bend left=30] node[above, font=\tiny] {$>$2/3 votos} (b2);
\end{tikzpicture}
\end{center}

\subsubsection{LMD-GHOST: Regla de Fork Choice}

\begin{definition}[LMD-GHOST]
Latest Message Driven Greedy Heaviest Observed SubTree.

Para elegir la cadena canónica:
\begin{enumerate}
    \item Empezar desde el último bloque finalizado
    \item En cada fork, elegir la rama con más votos recientes
    \item ``Reciente'' = último voto de cada validador
\end{enumerate}
\end{definition}

\begin{intuicion}
LMD-GHOST es como una votación continua donde cada validador tiene un voto, y en cada bifurcación, la rama con más votos gana. Pero solo cuenta el voto más reciente de cada validador.

Esto permite:
\begin{itemize}
    \item Adaptarse rápido a nuevos bloques
    \item Resistir censura (un bloque censurado puede ser votado después)
    \item Mantener liveness incluso con particiones de red temporales
\end{itemize}
\end{intuicion}

\subsubsection{Liquid Staking}

\begin{definition}[Liquid Staking]
Mecanismo donde el stake bloqueado se representa con un token líquido.
\end{definition}

\begin{example}[stETH de Lido]
\begin{enumerate}
    \item Usuario deposita ETH en Lido
    \item Lido stakea el ETH con validadores
    \item Usuario recibe stETH (1:1 inicialmente)
    \item stETH acumula recompensas (rebasing)
    \item Usuario puede usar stETH en DeFi
    \item Cuando quiere salir: intercambia stETH por ETH
\end{enumerate}

\textbf{Ventajas:}
\begin{itemize}
    \item Liquidez: no hay lock-up
    \item Composabilidad: usar como colateral
    \item Sin mínimo de 32 ETH
\end{itemize}

\textbf{Riesgos:}
\begin{itemize}
    \item Centralización: Lido tiene $\sim$30\% del stake
    \item Smart contract risk
    \item Depeg en mercados secundarios
\end{itemize}
\end{example}

\begin{advertencia}
\textbf{Riesgo de centralización en Ethereum PoS:}

\begin{center}
\begin{tabular}{lr}
\toprule
\textbf{Entidad} & \textbf{\% del stake} \\
\midrule
Lido & $\sim$32\% \\
Coinbase & $\sim$13\% \\
Kraken & $\sim$7\% \\
Binance & $\sim$5\% \\
Rocket Pool & $\sim$3\% \\
Otros & $\sim$40\% \\
\bottomrule
\end{tabular}
\end{center}

Si Lido + Coinbase + Kraken cooperaran, tendrían $>$50\% del stake. Esto es un riesgo de censura, aunque hay mecanismos sociales y técnicos para mitigarlo.
\end{advertencia}

\subsubsection{Comparación PoW vs PoS}

\begin{center}
\begin{tabular}{lll}
\toprule
\textbf{Aspecto} & \textbf{PoW} & \textbf{PoS} \\
\midrule
Recurso escaso & Energía + hardware & Capital \\
Costo de ataque & Comprar/alquilar hashpower & Comprar stake \\
Después de ataque & Hardware sigue valiendo & Stake es slashed \\
Barrera de entrada & Hardware especializado & Capital + técnico \\
Consumo energético & Muy alto & Bajo \\
Finalidad & Probabilística & Determinística (con BFT) \\
Tiempo a finalidad & $\sim$60 min (6 conf) & $\sim$15 min \\
Descentralización & Por hashpower & Por stake \\
\bottomrule
\end{tabular}
\end{center}

%-------------------------------------------
\subsection{Proof of Authority (PoA)}

\begin{definition}[PoA]
Validadores son entidades conocidas y autorizadas.
\end{definition}

\begin{itemize}
    \item Muy rápido (no hay competencia)
    \item Menos descentralizado
    \item Usado en: testnets, redes empresariales, sidechains
\end{itemize}

\subsubsection{Cuándo Usar PoA}

\begin{keyidea}
\textbf{PoA es apropiado cuando:}
\begin{itemize}
    \item Los validadores son conocidos y tienen reputación que perder
    \item La descentralización completa no es el objetivo principal
    \item Se necesita alta velocidad y bajo costo
    \item Hay una entidad central que puede autorizar validadores
\end{itemize}

\textbf{Casos de uso:}
\begin{itemize}
    \item \textbf{Testnets:} Ethereum Goerli, Sepolia
    \item \textbf{Sidechains:} Polygon PoS (originalmente), Ronin
    \item \textbf{Enterprise:} Hyperledger Besu, Quorum
    \item \textbf{Gaming:} Ronin (Axie Infinity)
\end{itemize}
\end{keyidea}

\begin{example}[Ronin Network]
Ronin es una sidechain para Axie Infinity que usaba PoA:
\begin{itemize}
    \item 9 validadores autorizados por Sky Mavis
    \item 5/9 requeridos para firmar transacciones
    \item Muy rápido y barato para el juego
\end{itemize}

\textbf{El problema:} En marzo 2022, un atacante comprometió 5 de los 9 validadores (incluyendo 4 de Sky Mavis) y robó \$625M del bridge.

\textbf{Lección:} PoA con pocos validadores es un punto único de fallo.
\end{example}

\subsubsection{Algoritmos de PoA}

\begin{center}
\begin{tabular}{lll}
\toprule
\textbf{Algoritmo} & \textbf{Usado en} & \textbf{Características} \\
\midrule
Clique & Ethereum testnets & Simple, round-robin \\
Aura & Parity/OpenEthereum & Round-robin con slots \\
IBFT 2.0 & Hyperledger Besu & BFT, finalidad inmediata \\
QBFT & ConsenSys Quorum & BFT optimizado \\
\bottomrule
\end{tabular}
\end{center}

%-------------------------------------------
\subsection{Otros Mecanismos de Consenso}

\begin{remark}
\textbf{Variantes y mecanismos alternativos:}

\textbf{Proof of History (PoH) - Solana:}
\begin{itemize}
    \item No es consenso, es un reloj criptográfico
    \item SHA-256 secuencial crea timestamps verificables
    \item Combinado con PoS para consenso real
    \item Permite ordenar transacciones sin comunicación
\end{itemize}

\textbf{Avalanche Consensus:}
\begin{itemize}
    \item Muestreo aleatorio repetido
    \item Pregunta a $k$ nodos aleatorios
    \item Si $>\alpha k$ coinciden, adopta su opinión
    \item Repite hasta convergencia
    \item Finalidad sub-segundo
\end{itemize}

\textbf{Tendermint BFT (Cosmos):}
\begin{itemize}
    \item BFT clásico con proposer rotativo
    \item Tolera hasta 1/3 maliciosos
    \item Finalidad inmediata (1 bloque)
    \item Usado en Cosmos Hub, Binance Chain
\end{itemize}

\textbf{Ouroboros (Cardano):}
\begin{itemize}
    \item PoS con prueba formal de seguridad
    \item Divide tiempo en slots y epochs
    \item Líder elegido por VRF ponderado por stake
\end{itemize}

\textbf{Nominated PoS (Polkadot):}
\begin{itemize}
    \item Nominadores eligen validadores
    \item Algoritmo Phragmén para distribución equitativa
    \item Slashing compartido entre validador y nominadores
\end{itemize}
\end{remark}

\begin{bancozk}
\textbf{Consenso en Banco ZK:}

Banco ZK como L2 no necesita su propio consenso para seguridad---hereda la seguridad de Ethereum L1.

Sin embargo, necesitamos ordenar transacciones:
\begin{itemize}
    \item \textbf{Fase actual:} Sequencer centralizado (eficiente pero requiere confianza)
    \item \textbf{Futuro:} Sequencer descentralizado o based sequencing (L1 ordena)
\end{itemize}

La privacidad de Banco ZK ayuda con consenso:
\begin{itemize}
    \item Transacciones encriptadas previenen MEV del sequencer
    \item El sequencer no puede hacer frontrunning si no ve el contenido
\end{itemize}
\end{bancozk}

% ============================================
% BLOQUE 2: BLOCKCHAIN BASE (L1) - VERSIÓN EXPANDIDA
% Parte 2: Bloques, Transacciones, Gas y Fees
% ============================================

%===========================================
\section{Bloques y Transacciones}
%===========================================

\begin{intuicion}
Un bloque es como una página en un libro de contabilidad: contiene un conjunto de transacciones que ocurrieron en un período de tiempo. Cada página hace referencia a la anterior, creando una cadena ininterrumpida desde el inicio.

La estructura exacta del bloque varía entre blockchains, pero todas comparten el concepto de:
\begin{itemize}
    \item Header: metadatos sobre el bloque
    \item Body: las transacciones en sí
\end{itemize}
\end{intuicion}

%-------------------------------------------
\subsection{Estructura de Bloque}

\subsubsection{Bloque de Bitcoin}

\begin{definition}[Bloque de Bitcoin]
\textbf{Header (80 bytes):}
\begin{itemize}
    \item \texttt{version} (4 bytes): Versión del protocolo
    \item \texttt{prev\_block\_hash} (32 bytes): Hash del bloque anterior
    \item \texttt{merkle\_root} (32 bytes): Raíz del árbol de transacciones
    \item \texttt{timestamp} (4 bytes): Tiempo Unix
    \item \texttt{bits} (4 bytes): Target de dificultad codificado
    \item \texttt{nonce} (4 bytes): Valor para PoW
\end{itemize}

\textbf{Body:}
\begin{itemize}
    \item \texttt{tx\_count}: Número de transacciones (varint)
    \item \texttt{transactions}: Lista de transacciones
\end{itemize}
\end{definition}

\begin{center}
\begin{tikzpicture}[
    header/.style={draw, rectangle, fill=blue!20, minimum width=8cm, minimum height=0.6cm},
    field/.style={draw, rectangle, fill=blue!10, minimum width=2.5cm, minimum height=0.5cm, font=\small},
    body/.style={draw, rectangle, fill=green!20, minimum width=8cm, minimum height=2cm}
]
    % Header
    \node[header] (h) at (0, 3) {\textbf{Block Header (80 bytes)}};
    
    % Header fields
    \node[field] at (-2.5, 2.2) {version};
    \node[field] at (0, 2.2) {prev\_hash};
    \node[field] at (2.5, 2.2) {merkle\_root};
    \node[field] at (-2.5, 1.5) {timestamp};
    \node[field] at (0, 1.5) {bits};
    \node[field] at (2.5, 1.5) {nonce};
    
    % Body
    \node[body] (b) at (0, -0.5) {};
    \node at (0, 0.3) {\textbf{Block Body}};
    \node at (0, -0.3) {Coinbase TX + Transacciones};
    \node at (0, -0.9) {(hasta $\sim$1 MB)};
\end{tikzpicture}
\end{center}

\begin{example}[Bloque de Bitcoin real]
\textbf{Bloque \#800000 (Julio 2023):}
\begin{verbatim}
Header:
  version: 0x27ffe000
  prev_block: 00000000000000000002a7c4c1e48d76...
  merkle_root: b2f7c58b5c8f2e7d4a1b6c9d8e2f3a4b...
  timestamp: 1690243200 (2023-07-24 16:00:00 UTC)
  bits: 0x17053894
  nonce: 3127834567
  
Block hash: 00000000000000000002a7c4c1e48d76c9...
  (empieza con muchos ceros por PoW)

Body:
  tx_count: 3721
  transactions: [coinbase, tx1, tx2, ..., tx3720]
  size: 1,563,789 bytes
\end{verbatim}
\end{example}

\subsubsection{Bloque de Ethereum (Post-Merge)}

\begin{definition}[Bloque de Ethereum]
\textbf{Header:}
\begin{itemize}
    \item \texttt{parentHash}: Hash del bloque padre
    \item \texttt{stateRoot}: Raíz del Merkle Patricia Trie de estado
    \item \texttt{transactionsRoot}: Raíz del trie de transacciones
    \item \texttt{receiptsRoot}: Raíz del trie de receipts
    \item \texttt{logsBloom}: Bloom filter de logs (256 bytes)
    \item \texttt{number}: Número de bloque
    \item \texttt{gasLimit}: Límite de gas del bloque
    \item \texttt{gasUsed}: Gas usado en el bloque
    \item \texttt{timestamp}: Tiempo Unix
    \item \texttt{baseFeePerGas}: Base fee (EIP-1559)
    \item \texttt{withdrawalsRoot}: Raíz de withdrawals (post-Shanghai)
    \item \texttt{blobGasUsed}: Gas de blobs usado (EIP-4844)
    \item \texttt{excessBlobGas}: Exceso de blob gas
\end{itemize}

\textbf{Body:}
\begin{itemize}
    \item \texttt{transactions}: Lista de transacciones
    \item \texttt{withdrawals}: Retiros de stake
\end{itemize}
\end{definition}

\begin{mathnote}
\textbf{Los tres tries de Ethereum:}

Cada bloque tiene raíces de tres Merkle Patricia Tries:

\textbf{1. State Trie} (\texttt{stateRoot}):
\[ \text{Address} \to (\text{nonce}, \text{balance}, \text{storageRoot}, \text{codeHash}) \]
Representa TODO el estado de Ethereum.

\textbf{2. Transaction Trie} (\texttt{transactionsRoot}):
\[ \text{index} \to \text{RLP(transaction)} \]
Solo las transacciones de este bloque.

\textbf{3. Receipt Trie} (\texttt{receiptsRoot}):
\[ \text{index} \to \text{RLP(receipt)} \]
Resultado de cada transacción (status, gas usado, logs).
\end{mathnote}

%-------------------------------------------
\subsection{Transacciones}

\subsubsection{Modelo UTXO (Bitcoin)}

\begin{definition}[UTXO - Unspent Transaction Output]
En Bitcoin no hay ``cuentas'' con balances. En su lugar:
\begin{itemize}
    \item Cada transacción consume UTXOs existentes (inputs)
    \item Y crea nuevos UTXOs (outputs)
    \item Tu ``balance'' es la suma de UTXOs que puedes gastar
\end{itemize}
\end{definition}

\begin{center}
\begin{tikzpicture}[
    utxo/.style={draw, rectangle, fill=yellow!30, minimum width=2cm, minimum height=0.8cm},
    tx/.style={draw, rectangle, fill=blue!20, minimum width=2cm, minimum height=1cm},
    arrow/.style={->, thick}
]
    % UTXOs de entrada
    \node[utxo] (u1) at (-3, 1) {UTXO: 0.5 BTC};
    \node[utxo] (u2) at (-3, -1) {UTXO: 0.3 BTC};
    
    % Transacción
    \node[tx] (tx) at (0, 0) {TX};
    
    % UTXOs de salida
    \node[utxo] (o1) at (3, 1) {UTXO: 0.7 BTC};
    \node[utxo] (o2) at (3, -1) {UTXO: 0.09 BTC};
    
    % Flechas
    \draw[arrow] (u1) -- (tx);
    \draw[arrow] (u2) -- (tx);
    \draw[arrow] (tx) -- (o1);
    \draw[arrow] (tx) -- (o2);
    
    % Labels
    \node[below=0.3cm of tx, font=\small] {Fee: 0.01 BTC};
    \node[above=0.1cm of u1, font=\tiny] {De Alice};
    \node[above=0.1cm of u2, font=\tiny] {De Alice};
    \node[above=0.1cm of o1, font=\tiny] {A Bob};
    \node[above=0.1cm of o2, font=\tiny] {Cambio a Alice};
\end{tikzpicture}
\end{center}

\begin{definition}[Estructura de Transacción Bitcoin]
\begin{verbatim}
Transaction:
  version: 4 bytes
  
  inputs: [
    {
      prev_txid: 32 bytes    # Hash de tx anterior
      prev_vout: 4 bytes     # Índice del output
      scriptSig: variable    # Script de desbloqueo
      sequence: 4 bytes      # Para timelocks
    },
    ...
  ]
  
  outputs: [
    {
      value: 8 bytes         # Cantidad en satoshis
      scriptPubKey: variable # Script de bloqueo
    },
    ...
  ]
  
  locktime: 4 bytes          # Tiempo mínimo de validez
\end{verbatim}
\end{definition}

\begin{keyidea}
\textbf{Scripts de Bitcoin:}

Cada UTXO está ``bloqueado'' por un script que especifica las condiciones para gastarlo.

\textbf{P2PKH (Pay to Public Key Hash):}
\begin{verbatim}
scriptPubKey: OP_DUP OP_HASH160 <pubKeyHash> OP_EQUALVERIFY OP_CHECKSIG
scriptSig: <signature> <publicKey>
\end{verbatim}

\textbf{P2SH (Pay to Script Hash):}
\begin{verbatim}
scriptPubKey: OP_HASH160 <scriptHash> OP_EQUAL
scriptSig: <signatures...> <redeemScript>
\end{verbatim}

\textbf{P2WPKH (Native SegWit):}
\begin{verbatim}
scriptPubKey: OP_0 <20-byte-pubKeyHash>
witness: [<signature>, <publicKey>]
\end{verbatim}
\end{keyidea}

\begin{remark}
\textbf{Ventajas del modelo UTXO:}
\begin{itemize}
    \item Paralelismo natural: transacciones independientes no conflictan
    \item Privacidad: cada UTXO es distinto, dificulta tracking
    \item Validación simple: solo verificar que inputs existen y no gastados
    \item Sin estado global: no necesitas conocer todo el estado
\end{itemize}

\textbf{Desventajas:}
\begin{itemize}
    \item Smart contracts complejos son difíciles
    \item Gestión de UTXOs puede ser engorrosa
    \item ``Dust'' (UTXOs muy pequeños) es problema
\end{itemize}
\end{remark}

\subsubsection{Modelo Account (Ethereum)}

\begin{definition}[Transacción de Ethereum]
\textbf{Campos:}
\begin{itemize}
    \item \texttt{type}: Tipo de transacción (0, 1, 2, 3)
    \item \texttt{chainId}: ID de la red (1 = mainnet)
    \item \texttt{nonce}: Contador de transacciones del sender
    \item \texttt{to}: Dirección destino (null para deploy)
    \item \texttt{value}: Wei a transferir
    \item \texttt{data}: Calldata (para contratos)
    \item \texttt{gasLimit}: Máximo gas a usar
    \item \texttt{maxFeePerGas}: Máximo fee por gas (EIP-1559)
    \item \texttt{maxPriorityFeePerGas}: Máximo tip (EIP-1559)
    \item \texttt{accessList}: Direcciones/slots a pre-calentar (EIP-2930)
    \item \texttt{v, r, s}: Firma ECDSA
\end{itemize}
\end{definition}

\begin{example}[Transacción Ethereum decodificada]
\begin{verbatim}
{
  "type": 2,                    // EIP-1559
  "chainId": 1,                 // Mainnet
  "nonce": 42,                  // 43ª tx de esta cuenta
  "to": "0x7a250d5630B4cF539739dF2C5dAcb4c659F2488D",  // Uniswap
  "value": "0",                 // No envía ETH
  "data": "0x38ed1739...",      // swap() function call
  "gasLimit": 250000,
  "maxFeePerGas": "30000000000",      // 30 gwei
  "maxPriorityFeePerGas": "1000000000", // 1 gwei tip
  "accessList": [],
  
  // Firma
  "v": 1,
  "r": "0x3a4b5c...",
  "s": "0x8d9e0f..."
}

// Decodificando data (ABI):
// Function: swapExactTokensForTokens(
//   uint256 amountIn: 1000000000,        // 1000 USDC
//   uint256 amountOutMin: 999000000000,  // Min 0.999 ETH
//   address[] path: [USDC, WETH],
//   address to: 0xMyAddress,
//   uint256 deadline: 1699999999
// )
\end{verbatim}
\end{example}

\begin{mathnote}
\textbf{Tipos de transacción en Ethereum:}

\begin{center}
\begin{tabular}{clll}
\toprule
\textbf{Type} & \textbf{Nombre} & \textbf{EIP} & \textbf{Características} \\
\midrule
0 & Legacy & Pre-EIP-2718 & gasPrice fijo \\
1 & Access List & EIP-2930 & + accessList \\
2 & EIP-1559 & EIP-1559 & maxFee + priorityFee \\
3 & Blob & EIP-4844 & + blobVersionedHashes \\
\bottomrule
\end{tabular}
\end{center}
\end{mathnote}

\subsubsection{El Nonce}

\begin{definition}[Nonce de Cuenta]
Contador que indica cuántas transacciones ha enviado una cuenta.

\textbf{Reglas:}
\begin{itemize}
    \item Transacción con nonce $n$ solo válida si la cuenta tiene nonce $n$
    \item Después de ejecutar, nonce incrementa a $n+1$
    \item No se pueden saltar nonces
    \item Se pueden enviar múltiples txs con nonces consecutivos
\end{itemize}
\end{definition}

\begin{advertencia}
\textbf{Problemas comunes con nonces:}

\textbf{1. Nonce gap:}
\begin{itemize}
    \item Enviaste tx con nonce 5, pero nonce 4 falló
    \item La tx 5 queda pendiente para siempre
    \item Solución: reenviar tx 4 o cancelarla
\end{itemize}

\textbf{2. Nonce reuse:}
\begin{itemize}
    \item Enviar dos txs con mismo nonce
    \item Solo una puede ser incluida
    \item Útil para ``cancelar'': enviar tx vacía con mismo nonce y más gas
\end{itemize}

\textbf{3. Pending nonce:}
\begin{itemize}
    \item Si tx está pending, siguiente debe usar nonce+1
    \item Algunos RPCs dan ``pending nonce'', otros no
\end{itemize}
\end{advertencia}

\subsubsection{Comparación UTXO vs Account}

\begin{center}
\begin{tabular}{lll}
\toprule
\textbf{Aspecto} & \textbf{UTXO (Bitcoin)} & \textbf{Account (Ethereum)} \\
\midrule
Estado & Set de UTXOs & Mapping address $\to$ account \\
Balance & Suma de UTXOs & Campo en account \\
Transacción & Consume UTXOs, crea nuevos & Modifica estado de cuentas \\
Paralelismo & Natural (UTXOs independientes) & Requiere cuidado (mismo nonce) \\
Privacidad & Mejor (direcciones únicas) & Peor (misma dirección) \\
Smart contracts & Limitados (Script) & Completos (EVM) \\
Complejidad & Gestión de UTXOs & Gestión de nonces \\
Espacio & Crece con UTXOs & Crece con cuentas \\
\bottomrule
\end{tabular}
\end{center}

%===========================================
\section{Gas y Fees}
%===========================================

\begin{definition}[Gas]
Unidad de medida del costo computacional en Ethereum.
\end{definition}

\begin{intuicion}
El gas es como la gasolina de un auto:
\begin{itemize}
    \item Cada operación consume cierta cantidad de gas
    \item Más compleja la operación = más gas
    \item Pagas por el gas que usas, no más
    \item Si te quedas sin gas, la transacción falla (pero pagas igual)
\end{itemize}

\textbf{¿Por qué gas y no ETH directamente?}

Separar la unidad de computación (gas) del precio (gwei) permite:
\begin{itemize}
    \item El costo de operaciones es estable en gas
    \item El precio en ETH fluctúa según demanda
    \item Protege contra ataques de spam
\end{itemize}
\end{intuicion}

%-------------------------------------------
\subsection{Costos de Gas por Operación}

\begin{definition}[Costos de Gas (EVM)]
\begin{center}
\begin{tabular}{llr}
\toprule
\textbf{Categoría} & \textbf{Operación} & \textbf{Gas} \\
\midrule
\multirow{4}{*}{Aritmética} 
    & ADD, SUB & 3 \\
    & MUL, DIV & 5 \\
    & MOD & 5 \\
    & EXP & $10 + 50 \times \text{bytes}$ \\
\midrule
\multirow{3}{*}{Comparación}
    & LT, GT, EQ & 3 \\
    & ISZERO & 3 \\
    & AND, OR, XOR & 3 \\
\midrule
\multirow{4}{*}{Storage}
    & SLOAD & 2,100 (cold), 100 (warm) \\
    & SSTORE (0 $\to$ non-0) & 20,000 \\
    & SSTORE (non-0 $\to$ non-0) & 2,900 \\
    & SSTORE (non-0 $\to$ 0) & 2,900 + refund \\
\midrule
\multirow{3}{*}{Memoria}
    & MLOAD, MSTORE & 3 \\
    & Memory expansion & $3 \times \text{words} + \text{words}^2 / 512$ \\
    & MSTORE8 & 3 \\
\midrule
\multirow{3}{*}{Control}
    & JUMP & 8 \\
    & JUMPI & 10 \\
    & CALL (cold) & 2,600 \\
    & CALL (warm) & 100 \\
\midrule
\multirow{3}{*}{Crypto}
    & SHA3/KECCAK256 & $30 + 6 \times \text{words}$ \\
    & ECRECOVER (precompile) & 3,000 \\
    & ECADD (precompile) & 150 \\
\midrule
\multirow{3}{*}{Transacción}
    & Base cost & 21,000 \\
    & Zero byte in data & 4 \\
    & Non-zero byte in data & 16 \\
\bottomrule
\end{tabular}
\end{center}
\end{definition}

\begin{example}[Cálculo de gas para transferencia simple]
\textbf{Transferencia de ETH (sin contrato):}
\begin{verbatim}
Base transaction cost:     21,000 gas
Data (vacío):                   0 gas
----------------------------------------
Total:                     21,000 gas
\end{verbatim}

Con gas price de 30 gwei:
\[ \text{Costo} = 21,000 \times 30 \text{ gwei} = 630,000 \text{ gwei} = 0.00063 \text{ ETH} \]

A \$2500/ETH: \$1.58
\end{example}

\begin{example}[Cálculo de gas para swap en Uniswap]
\textbf{Swap típico:}
\begin{verbatim}
Base transaction:           21,000
Calldata (228 bytes):       ~3,200  (mix de zeros y non-zeros)
Contract execution:
  - SLOAD (prices): 2x      4,200  (cold)
  - SLOAD (balances): 4x    8,400  (cold)
  - SSTORE (balances): 4x  11,600  (warm, modify)
  - Math operations:        ~1,000
  - Memory ops:               ~500
  - Event emission:         ~2,000
  - Other:                  ~5,000
----------------------------------------
Total:                    ~150,000 gas
\end{verbatim}

Con gas price de 30 gwei:
\[ \text{Costo} = 150,000 \times 30 \text{ gwei} = 0.0045 \text{ ETH} \approx \$11 \]
\end{example}

%-------------------------------------------
\subsection{EIP-1559: El Nuevo Modelo de Fees}

\begin{definition}[EIP-1559]
Modelo de fees introducido en el London hard fork (Agosto 2021):

\textbf{Componentes:}
\begin{itemize}
    \item \textbf{Base fee:} Fee mínimo por gas, ajustado algorítmicamente, \textbf{quemado}
    \item \textbf{Priority fee (tip):} Propina opcional al validador
    \item \textbf{Max fee:} Máximo que el usuario está dispuesto a pagar
\end{itemize}

\textbf{Fee efectivo por gas:}
\[ \text{effectiveFee} = \text{baseFee} + \min(\text{priorityFee}, \text{maxFee} - \text{baseFee}) \]
\end{definition}

\begin{mathnote}
\textbf{Ajuste del Base Fee:}

El base fee se ajusta para mantener bloques al 50\% de capacidad:
\begin{itemize}
    \item Target: 15M gas (50\% de 30M max)
    \item Si bloque $>$ 15M gas: base fee sube (hasta +12.5\%)
    \item Si bloque $<$ 15M gas: base fee baja (hasta -12.5\%)
\end{itemize}

\textbf{Fórmula:}
\[ \text{baseFee}_{n+1} = \text{baseFee}_n \times \left(1 + \frac{1}{8} \times \frac{\text{gasUsed} - \text{target}}{\text{target}}\right) \]

\textbf{Ejemplo:}
\begin{itemize}
    \item Base fee actual: 20 gwei
    \item Bloque usó: 25M gas (target: 15M)
    \item Exceso: $(25M - 15M) / 15M = 0.667$
    \item Cambio: $0.667 / 8 = 0.0833 = 8.33\%$
    \item Nuevo base fee: $20 \times 1.0833 = 21.67$ gwei
\end{itemize}
\end{mathnote}

\begin{center}
\begin{tikzpicture}
    \begin{axis}[
        width=12cm,
        height=6cm,
        xlabel={Gas usado en bloque (M)},
        ylabel={Cambio en base fee (\%)},
        xmin=0, xmax=30,
        ymin=-15, ymax=15,
        grid=major,
        legend pos=north west
    ]
    
    % Línea de cambio
    \addplot[blue, thick, domain=0:30] {((x-15)/15)/8 * 100};
    
    % Líneas de referencia
    \draw[dashed, red] (axis cs:15, -15) -- (axis cs:15, 15);
    \draw[dashed, gray] (axis cs:0, 0) -- (axis cs:30, 0);
    
    \node at (axis cs:15, -13) {Target: 15M};
    \node at (axis cs:7, 10) {Base fee baja};
    \node at (axis cs:23, 10) {Base fee sube};
    
    \end{axis}
\end{tikzpicture}
\end{center}

\begin{keyidea}
\textbf{¿Por qué quemar el base fee?}

\textbf{Antes de EIP-1559:}
\begin{itemize}
    \item Todo el fee iba al minero
    \item Mineros tenían incentivo de incluir txs propias con fee alto para manipular precio
    \item Users no sabían cuánto pagar
\end{itemize}

\textbf{Después de EIP-1559:}
\begin{itemize}
    \item Base fee se quema: nadie lo recibe
    \item Solo el tip va al validador
    \item Validadores no pueden manipular el base fee
    \item ETH se vuelve potencialmente deflacionario
\end{itemize}

\textbf{Impacto deflacionario:}

Si se quema más ETH del que se emite como rewards:
\[ \text{ETH quemado} > \text{ETH emitido} \implies \text{Supply decrece} \]

En períodos de alta actividad, Ethereum se ha vuelto deflacionario.
\end{keyidea}

\begin{example}[Transacción EIP-1559]
\begin{verbatim}
Usuario configura:
  maxFeePerGas: 50 gwei
  maxPriorityFeePerGas: 2 gwei

Estado actual:
  baseFee: 30 gwei

Cálculo:
  effectivePriorityFee = min(2, 50 - 30) = min(2, 20) = 2 gwei
  effectiveFee = 30 + 2 = 32 gwei

Distribución por gas:
  - Quemado: 30 gwei
  - Al validador: 2 gwei
  - Refund al usuario: 50 - 32 = 18 gwei (por unidad de gas limit)

Para 100,000 gas usado:
  - Total pagado: 3,200,000 gwei = 0.0032 ETH
  - Quemado: 3,000,000 gwei = 0.003 ETH
  - Validador: 200,000 gwei = 0.0002 ETH
\end{verbatim}
\end{example}

%-------------------------------------------
\subsection{Gas Limit y Gas Used}

\begin{definition}[Gas Limit]
\textbf{Por transacción:}
\begin{itemize}
    \item Máximo gas que la transacción puede consumir
    \item Usuario lo especifica
    \item Si se excede, transacción falla pero gas no se devuelve
\end{itemize}

\textbf{Por bloque:}
\begin{itemize}
    \item Máximo gas total en un bloque
    \item Actualmente: 30M gas (target: 15M)
    \item Validadores pueden votar para ajustarlo
\end{itemize}
\end{definition}

\begin{advertencia}
\textbf{Out of Gas:}

Si tu transacción se queda sin gas:
\begin{enumerate}
    \item La transacción se revierte completamente
    \item PERO pagas por todo el gas que especificaste
    \item El estado no cambia (excepto tu balance)
\end{enumerate}

\textbf{Estrategia:} Siempre poner un poco más de gas limit del necesario. El gas no usado se devuelve.
\end{advertencia}

%-------------------------------------------
\subsection{EIP-4844: Blob Gas}

\begin{definition}[Blob Gas]
EIP-4844 introduce un mercado de gas separado para blobs:
\begin{itemize}
    \item \textbf{Blob base fee:} Independiente del base fee normal
    \item \textbf{Target:} 3 blobs por bloque
    \item \textbf{Max:} 6 blobs por bloque
    \item Mismo mecanismo EIP-1559 para ajuste
\end{itemize}
\end{definition}

\begin{mathnote}
\textbf{Fórmula del blob base fee:}

\[ \text{blobBaseFee}_{n+1} = \text{blobBaseFee}_n \times e^{\frac{\text{excessBlobGas}}{3338477}} \]

Donde:
\begin{itemize}
    \item Cada blob consume 131,072 gas de blob
    \item Target por bloque: 393,216 blob gas (3 blobs)
    \item Max por bloque: 786,432 blob gas (6 blobs)
\end{itemize}

El blob base fee puede variar mucho más que el base fee normal (sin límite de 12.5\%).
\end{mathnote}

%-------------------------------------------
\subsection{Optimización de Gas}

\begin{keyidea}
\textbf{Técnicas para reducir gas en smart contracts:}

\textbf{1. Storage es caro, memoria es barata:}
\begin{verbatim}
// Malo: múltiples SLOADs
function bad() external {
    total += values[0];  // SLOAD
    total += values[1];  // SLOAD
    total += values[2];  // SLOAD
}

// Bueno: cargar una vez a memoria
function good() external {
    uint256 _total = total;      // SLOAD una vez
    uint256[] memory _values = values; // SLOAD una vez
    for (uint i = 0; i < 3; i++) {
        _total += _values[i];    // MLOAD (barato)
    }
    total = _total;              // SSTORE una vez
}
\end{verbatim}

\textbf{2. Packing de variables:}
\begin{verbatim}
// Malo: 3 slots de storage
struct Bad {
    uint256 a;  // slot 0
    uint8 b;    // slot 1 (desperdicia 31 bytes)
    uint256 c;  // slot 2
}

// Bueno: 2 slots
struct Good {
    uint256 a;  // slot 0
    uint256 c;  // slot 1
    uint8 b;    // slot 1 (comparte con c)
}
\end{verbatim}

\textbf{3. Usar eventos en vez de storage para datos históricos:}
\begin{verbatim}
// Malo: guardar historial en storage
mapping(uint => Trade) public tradeHistory;

// Bueno: emitir eventos (logs son más baratos)
event TradeExecuted(address indexed user, uint amount, uint price);
\end{verbatim}

\textbf{4. Short-circuit evaluation:}
\begin{verbatim}
// Poner condición más probable de fallar primero
require(amount > 0 && balances[msg.sender] >= amount);
\end{verbatim}

\textbf{5. Usar calldata en vez de memory para inputs:}
\begin{verbatim}
// Malo
function process(uint256[] memory data) external { }

// Bueno (data es solo lectura)
function process(uint256[] calldata data) external { }
\end{verbatim}
\end{keyidea}

\begin{bancozk}
\textbf{Gas en Banco ZK:}

Como L2, Banco ZK tiene sus propios costos de gas internos, más el costo de publicar datos en L1.

\textbf{Costos L2 (dentro de Banco ZK):}
\begin{itemize}
    \item Transferencia privada: $\sim$500 gas interno
    \item Verificación de prueba: incluida en batch
    \item Mucho más barato que L1
\end{itemize}

\textbf{Costos L1 (data availability):}
\begin{itemize}
    \item Publicar batch: $\sim$100K-500K gas L1
    \item Dividido entre todas las txs del batch
    \item Con 1000 txs: $\sim$100-500 gas L1 por tx
\end{itemize}

\textbf{Optimizaciones específicas:}
\begin{itemize}
    \item Comprimir datos al máximo antes de publicar
    \item Usar blobs (EIP-4844) cuando disponible
    \item Batching agresivo para amortizar costos
    \item Pruebas ZK permiten no publicar datos completos
\end{itemize}
\end{bancozk}

% ============================================
% BLOQUE 2: BLOCKCHAIN BASE (L1) - VERSIÓN EXPANDIDA
% Parte 3: EVM, Estado Global, Nodos
% ============================================

%===========================================
\section{EVM (Ethereum Virtual Machine)}
%===========================================

\begin{definition}[EVM]
Máquina virtual stack-based que ejecuta bytecode de smart contracts.
\end{definition}

\begin{intuicion}
La EVM es como una computadora virtual que existe en todos los nodos de Ethereum simultáneamente. Cuando ejecutas un smart contract:
\begin{enumerate}
    \item Tu transacción llega a los validadores
    \item Cada validador ejecuta el mismo código en su EVM
    \item Todos llegan al mismo resultado (determinismo)
    \item El nuevo estado se guarda
\end{enumerate}

Es ``la computadora mundial'': una máquina que nadie controla pero todos pueden usar.
\end{intuicion}

%-------------------------------------------
\subsection{Arquitectura de la EVM}

\begin{definition}[Componentes de la EVM]
\textbf{1. Stack:}
\begin{itemize}
    \item Máximo 1024 elementos
    \item Cada elemento: 256 bits (32 bytes)
    \item Operaciones: PUSH, POP, DUP, SWAP
    \item LIFO (Last In, First Out)
\end{itemize}

\textbf{2. Memory:}
\begin{itemize}
    \item Array de bytes, expandible
    \item Direccionada por byte (no por word)
    \item Volátil: se borra después de la ejecución
    \item Costo: 3 gas base + expansión cuadrática
\end{itemize}

\textbf{3. Storage:}
\begin{itemize}
    \item Mapping persistente: uint256 $\to$ uint256
    \item Cada contrato tiene su propio storage
    \item Persistente entre transacciones
    \item MUY caro: 20,000 gas para escribir nuevo valor
\end{itemize}

\textbf{4. Calldata:}
\begin{itemize}
    \item Datos de entrada de la transacción
    \item Read-only
    \item Más barato que memory
\end{itemize}

\textbf{5. Code:}
\begin{itemize}
    \item Bytecode del contrato
    \item Inmutable después del deploy
    \item Accesible con CODECOPY
\end{itemize}
\end{definition}

\begin{center}
\begin{tikzpicture}[
    component/.style={draw, rectangle, minimum width=3cm, minimum height=1.5cm, align=center},
    arrow/.style={->, thick}
]
    % Componentes
    \node[component, fill=blue!20] (stack) at (0, 0) {Stack\\1024 × 256 bits\\Volátil};
    \node[component, fill=green!20] (mem) at (4, 0) {Memory\\Bytes array\\Volátil};
    \node[component, fill=red!20] (storage) at (8, 0) {Storage\\uint256 $\to$ uint256\\Persistente};
    
    \node[component, fill=yellow!20] (calldata) at (2, -3) {Calldata\\Input data\\Read-only};
    \node[component, fill=purple!20] (code) at (6, -3) {Code\\Bytecode\\Inmutable};
    
    % EVM
    \node[draw, rectangle, minimum width=10cm, minimum height=1cm, fill=gray!20] (evm) at (4, 2.5) {EVM Execution Context};
    
    % Flechas
    \draw[arrow] (evm) -- (stack);
    \draw[arrow] (evm) -- (mem);
    \draw[arrow] (evm) -- (storage);
    \draw[arrow] (calldata) -- (evm);
    \draw[arrow] (code) -- (evm);
\end{tikzpicture}
\end{center}

%-------------------------------------------
\subsection{Opcodes de la EVM}

\begin{definition}[Categorías de Opcodes]
La EVM tiene $\sim$140 opcodes organizados en categorías:
\end{definition}

\subsubsection{Operaciones Aritméticas}

\begin{center}
\begin{tabular}{llcl}
\toprule
\textbf{Opcode} & \textbf{Hex} & \textbf{Gas} & \textbf{Descripción} \\
\midrule
ADD & 0x01 & 3 & $a + b$ \\
MUL & 0x02 & 5 & $a \times b$ \\
SUB & 0x03 & 3 & $a - b$ \\
DIV & 0x04 & 5 & $a \div b$ (entero) \\
SDIV & 0x05 & 5 & División con signo \\
MOD & 0x06 & 5 & $a \mod b$ \\
SMOD & 0x07 & 5 & Módulo con signo \\
ADDMOD & 0x08 & 8 & $(a + b) \mod N$ \\
MULMOD & 0x09 & 8 & $(a \times b) \mod N$ \\
EXP & 0x0a & $10 + 50 \times \text{bytes}$ & $a^b$ \\
SIGNEXTEND & 0x0b & 5 & Extensión de signo \\
\bottomrule
\end{tabular}
\end{center}

\subsubsection{Operaciones de Comparación y Lógica}

\begin{center}
\begin{tabular}{llcl}
\toprule
\textbf{Opcode} & \textbf{Hex} & \textbf{Gas} & \textbf{Descripción} \\
\midrule
LT & 0x10 & 3 & $a < b$ \\
GT & 0x11 & 3 & $a > b$ \\
SLT & 0x12 & 3 & $<$ con signo \\
SGT & 0x13 & 3 & $>$ con signo \\
EQ & 0x14 & 3 & $a = b$ \\
ISZERO & 0x15 & 3 & $a = 0$ \\
AND & 0x16 & 3 & Bitwise AND \\
OR & 0x17 & 3 & Bitwise OR \\
XOR & 0x18 & 3 & Bitwise XOR \\
NOT & 0x19 & 3 & Bitwise NOT \\
SHL & 0x1b & 3 & Shift left \\
SHR & 0x1c & 3 & Shift right (logical) \\
SAR & 0x1d & 3 & Shift right (arithmetic) \\
\bottomrule
\end{tabular}
\end{center}

\subsubsection{Operaciones de Storage y Memory}

\begin{center}
\begin{tabular}{llcl}
\toprule
\textbf{Opcode} & \textbf{Hex} & \textbf{Gas} & \textbf{Descripción} \\
\midrule
SLOAD & 0x54 & 2100/100 & Leer de storage \\
SSTORE & 0x55 & 20000/2900 & Escribir a storage \\
MLOAD & 0x51 & 3 & Leer 32 bytes de memory \\
MSTORE & 0x52 & 3 & Escribir 32 bytes a memory \\
MSTORE8 & 0x53 & 3 & Escribir 1 byte a memory \\
CALLDATALOAD & 0x35 & 3 & Leer de calldata \\
CALLDATACOPY & 0x37 & 3+ & Copiar calldata a memory \\
CODECOPY & 0x39 & 3+ & Copiar code a memory \\
\bottomrule
\end{tabular}
\end{center}

\subsubsection{Operaciones de Control de Flujo}

\begin{center}
\begin{tabular}{llcl}
\toprule
\textbf{Opcode} & \textbf{Hex} & \textbf{Gas} & \textbf{Descripción} \\
\midrule
STOP & 0x00 & 0 & Terminar ejecución \\
JUMP & 0x56 & 8 & Salto incondicional \\
JUMPI & 0x57 & 10 & Salto condicional \\
JUMPDEST & 0x5b & 1 & Marca destino válido \\
PC & 0x58 & 2 & Program counter \\
RETURN & 0xf3 & 0 & Retornar datos \\
REVERT & 0xfd & 0 & Revertir con datos \\
INVALID & 0xfe & all & Instrucción inválida \\
SELFDESTRUCT & 0xff & 5000+ & Destruir contrato \\
\bottomrule
\end{tabular}
\end{center}

\subsubsection{Operaciones de Llamada}

\begin{center}
\begin{tabular}{llcl}
\toprule
\textbf{Opcode} & \textbf{Hex} & \textbf{Gas} & \textbf{Descripción} \\
\midrule
CALL & 0xf1 & 100-2600 & Llamar otro contrato \\
CALLCODE & 0xf2 & 100-2600 & Llamar con contexto propio \\
DELEGATECALL & 0xf4 & 100-2600 & Llamar preservando context \\
STATICCALL & 0xfa & 100-2600 & Llamada read-only \\
CREATE & 0xf0 & 32000 & Crear nuevo contrato \\
CREATE2 & 0xf5 & 32000 & Crear con salt determinístico \\
\bottomrule
\end{tabular}
\end{center}

\begin{mathnote}
\textbf{Diferencia entre CALL, DELEGATECALL, STATICCALL:}

\textbf{CALL:}
\begin{itemize}
    \item Ejecuta código del contrato destino
    \item msg.sender = contrato que hace CALL
    \item Storage del contrato destino
\end{itemize}

\textbf{DELEGATECALL:}
\begin{itemize}
    \item Ejecuta código del contrato destino
    \item msg.sender = sender original
    \item Storage del contrato que hace DELEGATECALL
    \item Usado para proxies y librerías
\end{itemize}

\textbf{STATICCALL:}
\begin{itemize}
    \item Como CALL pero read-only
    \item Revierte si intenta modificar estado
    \item Usado para views y pure functions
\end{itemize}
\end{mathnote}

\subsubsection{Precompilados (Precompiled Contracts)}

\begin{definition}[Precompilados]
Contratos especiales con direcciones fijas que implementan operaciones costosas de forma nativa.
\end{definition}

\begin{center}
\begin{tabular}{clcl}
\toprule
\textbf{Address} & \textbf{Nombre} & \textbf{Gas} & \textbf{Función} \\
\midrule
0x01 & ECRECOVER & 3000 & Recuperar pubkey de firma \\
0x02 & SHA256 & 60+ & Hash SHA-256 \\
0x03 & RIPEMD160 & 600+ & Hash RIPEMD-160 \\
0x04 & IDENTITY & 15+ & Copiar datos (usado por DATACOPY) \\
0x05 & MODEXP & variable & Exponenciación modular \\
0x06 & ECADD & 150 & Suma en BN254 \\
0x07 & ECMUL & 6000 & Multiplicación escalar BN254 \\
0x08 & ECPAIRING & 34000+ & Pairing check BN254 \\
0x09 & BLAKE2F & 1+ & Compresión BLAKE2 \\
0x0a & KZG\_POINT\_EVALUATION & 50000 & Verificar evaluación KZG \\
\bottomrule
\end{tabular}
\end{center}

\begin{keyidea}
\textbf{¿Por qué precompilados?}

Algunas operaciones son demasiado costosas para implementar en bytecode EVM:
\begin{itemize}
    \item Verificación de firmas ECDSA
    \item Operaciones de curvas elípticas para ZK
    \item Hashes criptográficos
\end{itemize}

Los precompilados las implementan en código nativo, haciendo prácticas operaciones que serían prohibitivamente caras.

\textbf{Ejemplo:} Verificar una prueba Groth16 on-chain usa el precompilado de pairing (0x08).
\end{keyidea}

%-------------------------------------------
\subsection{Ejemplo de Ejecución EVM}

\begin{example}[Ejecución paso a paso]
\textbf{Código Solidity:}
\begin{verbatim}
function add(uint256 a, uint256 b) public pure returns (uint256) {
    return a + b;
}
\end{verbatim}

\textbf{Bytecode (simplificado):}
\begin{verbatim}
PUSH1 0x00      // Posición en calldata para a
CALLDATALOAD    // Cargar a
PUSH1 0x20      // Posición para b
CALLDATALOAD    // Cargar b
ADD             // a + b
PUSH1 0x00      // Posición en memory
MSTORE          // Guardar resultado
PUSH1 0x20      // Tamaño del resultado
PUSH1 0x00      // Posición del resultado
RETURN          // Retornar
\end{verbatim}

\textbf{Ejecución con a=5, b=3:}

\begin{center}
\begin{tabular}{clcl}
\toprule
\textbf{Step} & \textbf{Opcode} & \textbf{Stack (top first)} & \textbf{Gas} \\
\midrule
0 & PUSH1 0x00 & [0x00] & 3 \\
1 & CALLDATALOAD & [0x05] & 3 \\
2 & PUSH1 0x20 & [0x20, 0x05] & 3 \\
3 & CALLDATALOAD & [0x03, 0x05] & 3 \\
4 & ADD & [0x08] & 3 \\
5 & PUSH1 0x00 & [0x00, 0x08] & 3 \\
6 & MSTORE & [] & 3 \\
7 & PUSH1 0x20 & [0x20] & 3 \\
8 & PUSH1 0x00 & [0x00, 0x20] & 3 \\
9 & RETURN & -- & 0 \\
\midrule
& Total & & 27 \\
\bottomrule
\end{tabular}
\end{center}

Memory después de MSTORE:
\begin{verbatim}
0x00: 0x0000000000000000000000000000000000000000000000000000000000000008
\end{verbatim}

RETURN devuelve los bytes [0x00, 0x20), que es el número 8.
\end{example}

%-------------------------------------------
\subsection{Contexto de Ejecución}

\begin{definition}[Contexto de Ejecución]
Variables disponibles durante la ejecución:

\textbf{Información del bloque:}
\begin{itemize}
    \item \texttt{BLOCKHASH}: Hash de bloque reciente
    \item \texttt{COINBASE}: Address del validador
    \item \texttt{TIMESTAMP}: Tiempo del bloque
    \item \texttt{NUMBER}: Número de bloque
    \item \texttt{DIFFICULTY/PREVRANDAO}: Aleatoriedad (post-merge)
    \item \texttt{GASLIMIT}: Gas limit del bloque
    \item \texttt{CHAINID}: ID de la cadena
    \item \texttt{BASEFEE}: Base fee actual
\end{itemize}

\textbf{Información de la transacción:}
\begin{itemize}
    \item \texttt{ORIGIN}: EOA que inició la transacción
    \item \texttt{GASPRICE}: Precio del gas
    \item \texttt{BLOBHASH}: Hashes de blobs (EIP-4844)
    \item \texttt{BLOBBASEFEE}: Base fee de blobs
\end{itemize}

\textbf{Información del call:}
\begin{itemize}
    \item \texttt{ADDRESS}: Dirección del contrato actual
    \item \texttt{CALLER}: Quien llamó (msg.sender)
    \item \texttt{CALLVALUE}: ETH enviado (msg.value)
    \item \texttt{CALLDATASIZE}: Tamaño de calldata
\end{itemize}
\end{definition}

\begin{advertencia}
\textbf{tx.origin vs msg.sender:}

\begin{verbatim}
Usuario (EOA) -> Contrato A -> Contrato B

En Contrato B:
  tx.origin = Usuario (EOA)     // ORIGIN opcode
  msg.sender = Contrato A       // CALLER opcode
\end{verbatim}

\textbf{Nunca usar tx.origin para autenticación:}
\begin{itemize}
    \item Un contrato malicioso puede engañar al usuario para llamarlo
    \item Luego llamar a tu contrato
    \item tx.origin será el usuario, pero msg.sender será el contrato malicioso
\end{itemize}
\end{advertencia}

%===========================================
\section{Estado Global (State)}
%===========================================

\begin{definition}[Estado]
Mapeo de direcciones a cuentas:
\[
\sigma: \text{Address} \to (\text{nonce}, \text{balance}, \text{storageRoot}, \text{codeHash})
\]
\end{definition}

\begin{itemize}
    \item \textbf{EOA:} Externally Owned Account (usuario)
    \item \textbf{Contract:} Tiene código y storage
    \item Estado almacenado en Merkle Patricia Trie
    \item State root en cada block header
\end{itemize}

%-------------------------------------------
\subsection{Tipos de Cuentas}

\begin{definition}[Externally Owned Account (EOA)]
Cuenta controlada por una clave privada.

\textbf{Características:}
\begin{itemize}
    \item Creada generando un par de claves
    \item Address = últimos 20 bytes de Keccak256(pubkey)
    \item Puede iniciar transacciones
    \item No tiene código
    \item storageRoot y codeHash son valores especiales vacíos
\end{itemize}

\textbf{Campos de estado:}
\begin{itemize}
    \item \texttt{nonce}: Número de transacciones enviadas
    \item \texttt{balance}: ETH poseído (en wei)
    \item \texttt{storageRoot}: EMPTY\_ROOT
    \item \texttt{codeHash}: KECCAK256\_NULL
\end{itemize}
\end{definition}

\begin{definition}[Contract Account]
Cuenta con código ejecutable.

\textbf{Características:}
\begin{itemize}
    \item Creada por una transacción de deploy
    \item Address determinada por creator + nonce (CREATE) o salt (CREATE2)
    \item No puede iniciar transacciones (solo responder)
    \item Tiene código inmutable
    \item Tiene storage persistente
\end{itemize}

\textbf{Campos de estado:}
\begin{itemize}
    \item \texttt{nonce}: Número de contratos creados por este contrato
    \item \texttt{balance}: ETH poseído
    \item \texttt{storageRoot}: Raíz del MPT de storage
    \item \texttt{codeHash}: Keccak256 del bytecode
\end{itemize}
\end{definition}

\begin{mathnote}
\textbf{Cálculo de direcciones:}

\textbf{Para EOA:}
\[ \text{address} = \text{Keccak256}(\text{pubkey})[12:32] \]

\textbf{Para CREATE (nonce-based):}
\[ \text{address} = \text{Keccak256}(\text{RLP}([\text{sender}, \text{nonce}]))[12:32] \]

\textbf{Para CREATE2 (determinístico):}
\[ \text{address} = \text{Keccak256}(\texttt{0xff} \| \text{sender} \| \text{salt} \| \text{Keccak256}(\text{initCode}))[12:32] \]

CREATE2 permite predecir la dirección antes del deploy, útil para:
\begin{itemize}
    \item Counterfactual instantiation
    \item Factory patterns
    \item Upgradeable proxies
\end{itemize}
\end{mathnote}

%-------------------------------------------
\subsection{Storage Layout}

\begin{definition}[Storage Layout]
El storage de un contrato es un mapping de 256-bit a 256-bit:
\[ \text{storage}: \mathbb{Z}_{2^{256}} \to \mathbb{Z}_{2^{256}} \]

Inicialmente, todos los slots son 0.
\end{definition}

\begin{keyidea}
\textbf{Layout de Solidity:}

\textbf{Variables estáticas:}
\begin{verbatim}
contract Example {
    uint256 a;      // slot 0
    uint256 b;      // slot 1
    uint128 c;      // slot 2 (primeros 16 bytes)
    uint128 d;      // slot 2 (siguientes 16 bytes, packed)
    uint256 e;      // slot 3
}
\end{verbatim}

\textbf{Arrays dinámicos:}
\begin{verbatim}
uint256[] arr;  // slot n guarda length
// arr[i] está en: keccak256(n) + i
\end{verbatim}

\textbf{Mappings:}
\begin{verbatim}
mapping(address => uint256) balances;  // slot m
// balances[addr] está en: keccak256(addr . m)
// donde . es concatenación
\end{verbatim}

\textbf{Mappings anidados:}
\begin{verbatim}
mapping(address => mapping(address => uint256)) allowances;  // slot k
// allowances[a][b] está en: keccak256(b . keccak256(a . k))
\end{verbatim}
\end{keyidea}

\begin{example}[Encontrar slot de storage]
\begin{verbatim}
contract Token {
    string public name;           // slot 0
    string public symbol;         // slot 1  
    uint8 public decimals;        // slot 2
    uint256 public totalSupply;   // slot 3
    mapping(address => uint256) public balanceOf;  // slot 4
    mapping(address => mapping(address => uint256)) public allowance; // slot 5
}

// Para encontrar balanceOf[0x123...]:
slot = keccak256(abi.encode(0x123..., 4))

// Para encontrar allowance[0x123...][0x456...]:
inner_slot = keccak256(abi.encode(0x123..., 5))
slot = keccak256(abi.encode(0x456..., inner_slot))
\end{verbatim}

En Python:
\begin{verbatim}
from web3 import Web3

def get_balance_slot(address, mapping_slot=4):
    # Pad address to 32 bytes
    address_padded = address.lower()[2:].zfill(64)
    slot_padded = hex(mapping_slot)[2:].zfill(64)
    
    # Concatenate and hash
    data = address_padded + slot_padded
    slot = Web3.keccak(hexstr=data).hex()
    return slot

# Ejemplo
slot = get_balance_slot("0x123...")
# Ahora puedes usar eth_getStorageAt(contract, slot)
\end{verbatim}
\end{example}

%-------------------------------------------
\subsection{State Trie}

\begin{mathnote}
\textbf{Estructura del State Trie:}

El estado completo de Ethereum se almacena en un Merkle Patricia Trie:
\begin{itemize}
    \item \textbf{Key}: Keccak256(address) - 32 bytes
    \item \textbf{Value}: RLP([nonce, balance, storageRoot, codeHash])
\end{itemize}

Cada contrato tiene su propio Storage Trie:
\begin{itemize}
    \item \textbf{Key}: Keccak256(slot) - 32 bytes
    \item \textbf{Value}: RLP(valor) - hasta 32 bytes
\end{itemize}

\textbf{Tamaño del estado (2024):}
\begin{itemize}
    \item $\sim$250 millones de cuentas
    \item $\sim$100+ GB de datos de estado
    \item Creciendo constantemente
\end{itemize}
\end{mathnote}

%===========================================
\section{Nodos}
%===========================================

\begin{intuicion}
Un nodo es una computadora que participa en la red de Ethereum. Diferentes tipos de nodos tienen diferentes responsabilidades y requisitos de recursos.

La descentralización de la red depende de que mucha gente pueda correr nodos.
\end{intuicion}

%-------------------------------------------
\subsection{Full Node}

\begin{definition}[Full Node]
\begin{itemize}
    \item Almacena estado actual completo
    \item Valida todas las transacciones y bloques
    \item Puede no tener historial completo (pruned)
    \item Participa en la red P2P
\end{itemize}
\end{definition}

\begin{mathnote}
\textbf{Requisitos de Full Node (2024):}

\begin{center}
\begin{tabular}{ll}
\toprule
\textbf{Recurso} & \textbf{Requisito} \\
\midrule
CPU & 4+ cores \\
RAM & 16+ GB \\
Storage (pruned) & 1+ TB SSD \\
Bandwidth & 25+ Mbps \\
Tiempo de sync & 1-7 días \\
\bottomrule
\end{tabular}
\end{center}

\textbf{Clientes populares:}
\begin{itemize}
    \item \textbf{Geth}: Go, el más usado
    \item \textbf{Nethermind}: C\#, rápido sync
    \item \textbf{Besu}: Java, enterprise-friendly
    \item \textbf{Erigon}: Go, optimizado para storage
    \item \textbf{Reth}: Rust, nuevo y rápido
\end{itemize}
\end{mathnote}

%-------------------------------------------
\subsection{Light Node}

\begin{definition}[Light Node]
\begin{itemize}
    \item Solo block headers
    \item Verifica con Merkle proofs
    \item Confía en full nodes para datos
    \item Usado en wallets móviles
\end{itemize}
\end{definition}

\begin{keyidea}
\textbf{Cómo funciona un Light Node:}

\begin{enumerate}
    \item Descarga solo headers ($\sim$500 bytes cada uno)
    \item Para verificar un dato (ej: balance), pide Merkle proof a un full node
    \item Verifica que el proof es consistente con el state root en el header
    \item Confía en que el header es válido (verificado por consenso)
\end{enumerate}

\textbf{Trade-off:}
\begin{itemize}
    \item Pro: Muy ligero ($\sim$100 MB vs $\sim$1 TB)
    \item Con: Depende de full nodes para datos
    \item Con: Puede ser engañado si todos los peers son maliciosos
\end{itemize}
\end{keyidea}

%-------------------------------------------
\subsection{Archive Node}

\begin{definition}[Archive Node]
\begin{itemize}
    \item Todo el historial de estados
    \item Puede consultar estado en cualquier bloque
    \item Requiere varios TB de almacenamiento
    \item Usado por exploradores, analytics
\end{itemize}
\end{definition}

\begin{mathnote}
\textbf{Requisitos de Archive Node (2024):}

\begin{center}
\begin{tabular}{ll}
\toprule
\textbf{Recurso} & \textbf{Requisito} \\
\midrule
Storage & 15+ TB SSD \\
RAM & 32+ GB \\
Tiempo de sync & Semanas a meses \\
Costo de hardware & \$5,000+ \\
\bottomrule
\end{tabular}
\end{center}

\textbf{Uso:}
\begin{itemize}
    \item Etherscan y otros exploradores
    \item Servicios de analytics (Dune, Nansen)
    \item Debugging de transacciones históricas
    \item Investigación
\end{itemize}
\end{mathnote}

%-------------------------------------------
\subsection{Validator Node (Post-Merge)}

\begin{definition}[Validator Node]
Full node que además:
\begin{itemize}
    \item Tiene 32 ETH depositados
    \item Corre Consensus Layer client (Lighthouse, Prysm, etc.)
    \item Propone y attesta bloques
    \item Gana recompensas por participación
\end{itemize}
\end{definition}

\begin{center}
\begin{tikzpicture}[
    client/.style={draw, rectangle, minimum width=3cm, minimum height=1.2cm, align=center},
    arrow/.style={<->, thick}
]
    % Execution Layer
    \node[client, fill=blue!20] (el) at (0, 0) {Execution Layer\\(Geth, Nethermind)};
    
    % Consensus Layer
    \node[client, fill=green!20] (cl) at (5, 0) {Consensus Layer\\(Lighthouse, Prysm)};
    
    % Conexión
    \draw[arrow] (el) -- node[above] {Engine API} (cl);
    
    % Labels
    \node[below=0.5cm of el, font=\small, align=center] {Ejecuta txs\\Mantiene estado};
    \node[below=0.5cm of cl, font=\small, align=center] {Consenso PoS\\Propone bloques};
\end{tikzpicture}
\end{center}

%-------------------------------------------
\subsection{RPC Nodes y Providers}

\begin{definition}[RPC Node]
Nodo que expone una API JSON-RPC para que aplicaciones interactúen con Ethereum.
\end{definition}

\textbf{Métodos RPC comunes:}
\begin{verbatim}
eth_blockNumber          // Último bloque
eth_getBalance           // Balance de cuenta
eth_getTransactionCount  // Nonce de cuenta
eth_call                 // Simular llamada (view)
eth_sendRawTransaction   // Enviar tx firmada
eth_getTransactionReceipt // Resultado de tx
eth_getLogs              // Eventos filtrados
eth_getStorageAt         // Leer storage slot
\end{verbatim}

\textbf{Providers populares:}
\begin{center}
\begin{tabular}{lll}
\toprule
\textbf{Provider} & \textbf{Tier gratuito} & \textbf{Notas} \\
\midrule
Infura & 100k req/día & El más antiguo \\
Alchemy & 300M CU/mes & Buen dashboard \\
QuickNode & Limitado & Multi-chain \\
Ankr & 1M req/mes & Descentralizado \\
\bottomrule
\end{tabular}
\end{center}

\begin{advertencia}
\textbf{Dependencia de providers centralizados:}

La mayoría de dApps usan Infura o Alchemy en lugar de nodos propios.

\textbf{Riesgos:}
\begin{itemize}
    \item Censura: provider puede bloquear ciertas txs
    \item Privacidad: provider ve todas tus queries
    \item Single point of failure: si cae, tu app no funciona
\end{itemize}

\textbf{Solución:} Correr tu propio nodo o usar providers descentralizados (Pocket Network).
\end{advertencia}

\begin{bancozk}
\textbf{Nodos en Banco ZK:}

\textbf{L1 (Ethereum):}
\begin{itemize}
    \item Necesitamos nodo completo para verificar estado
    \item El sequencer monitorea eventos del contrato puente
    \item Archive node útil para debugging
\end{itemize}

\textbf{L2 (Banco ZK):}
\begin{itemize}
    \item Nuestros propios nodos procesan txs privadas
    \item Prover nodes generan las pruebas ZK
    \item Users pueden correr light nodes para verificar
\end{itemize}

\textbf{Descentralización futura:}
\begin{itemize}
    \item Múltiples sequencers en rotación
    \item Red de provers descentralizada
    \item Light clients verifican proofs sin confiar en nadie
\end{itemize}
\end{bancozk}

% ============================================
% BLOQUE 2: BLOCKCHAIN BASE (L1) - VERSIÓN EXPANDIDA
% Parte 4: Mempool y MEV
% ============================================

%===========================================
\section{Mempool}
%===========================================

\begin{definition}[Mempool]
Pool de transacciones pendientes esperando inclusión en un bloque.
\end{definition}

\begin{itemize}
    \item Transacciones broadcast a la red P2P
    \item Cada nodo mantiene su propia mempool
    \item Validadores/mineros seleccionan de su mempool
    \item Típicamente ordenado por fee (mayor fee primero)
    \item Público por defecto: cualquiera puede ver txs pendientes
\end{itemize}

%-------------------------------------------
\subsection{Cómo Funciona la Mempool}

\begin{intuicion}
La mempool es como una sala de espera para transacciones. Cuando envías una tx:
\begin{enumerate}
    \item Tu wallet la firma y la envía a un nodo (vía RPC)
    \item El nodo la valida (firma correcta, balance suficiente, etc.)
    \item Si es válida, la añade a su mempool
    \item La propaga a otros nodos via P2P (gossip)
    \item Un validador la incluye en su próximo bloque
    \item La tx sale de la mempool de todos los nodos
\end{enumerate}

Como la mempool es pública, cualquiera puede ver qué transacciones están esperando. Esto tiene implicaciones importantes para MEV.
\end{intuicion}

\begin{definition}[Ciclo de Vida de una Transacción]
\begin{center}
\begin{tikzpicture}[
    state/.style={draw, rectangle, rounded corners, minimum width=2.5cm, minimum height=0.8cm, align=center},
    arrow/.style={->, thick}
]
    % Estados
    \node[state, fill=blue!20] (created) at (0, 0) {Creada\\y firmada};
    \node[state, fill=yellow!20] (pending) at (4, 0) {Pending\\(mempool)};
    \node[state, fill=green!20] (included) at (8, 0) {Incluida\\(bloque)};
    \node[state, fill=green!40] (confirmed) at (12, 0) {Confirmada\\(finalizada)};
    
    % Alternativas
    \node[state, fill=red!20] (dropped) at (4, -2) {Dropped};
    \node[state, fill=orange!20] (replaced) at (4, -3.5) {Reemplazada};
    
    % Flechas
    \draw[arrow] (created) -- node[above, font=\small] {broadcast} (pending);
    \draw[arrow] (pending) -- node[above, font=\small] {minado} (included);
    \draw[arrow] (included) -- node[above, font=\small] {finality} (confirmed);
    
    \draw[arrow] (pending) -- node[right, font=\small] {timeout/low fee} (dropped);
    \draw[arrow] (pending) -- node[right, font=\small] {mismo nonce, más fee} (replaced);
\end{tikzpicture}
\end{center}
\end{definition}

%-------------------------------------------
\subsection{Validación en Mempool}

\begin{definition}[Checks de Validación]
Antes de aceptar una tx en mempool, el nodo verifica:

\textbf{Checks estáticos:}
\begin{itemize}
    \item Firma ECDSA válida
    \item Formato de tx correcto
    \item Gas limit $\geq$ gas intrínseco (21000 + data)
    \item Gas limit $\leq$ block gas limit
\end{itemize}

\textbf{Checks de estado:}
\begin{itemize}
    \item Nonce = nonce actual de la cuenta
    \item Balance $\geq$ value + gasLimit $\times$ gasPrice
    \item Sender no está en blacklist (si aplica)
\end{itemize}

\textbf{Checks de mempool:}
\begin{itemize}
    \item No duplicada (mismo hash)
    \item Si mismo nonce existe: nuevo fee debe ser $>$ 10\% mayor
\end{itemize}
\end{definition}

%-------------------------------------------
\subsection{Políticas de Mempool}

\begin{mathnote}
\textbf{Límites y políticas típicas:}

\begin{center}
\begin{tabular}{ll}
\toprule
\textbf{Parámetro} & \textbf{Valor típico (Geth)} \\
\midrule
Tamaño máximo & 5,000 txs \\
Memoria máxima & 1 GB \\
Tiempo máximo & 3 horas \\
Gap máximo de nonce & 64 \\
Fee mínimo & 1 gwei \\
Reemplazo mínimo & +10\% fee \\
\bottomrule
\end{tabular}
\end{center}

Cuando la mempool está llena:
\begin{itemize}
    \item Txs con menor fee son evictadas primero
    \item Txs muy viejas son descartadas
    \item Nonce gaps excesivos son rechazados
\end{itemize}
\end{mathnote}

%-------------------------------------------
\subsection{Propagación P2P}

\begin{definition}[Gossip Protocol]
Las transacciones se propagan por la red usando gossip:
\begin{enumerate}
    \item Nodo A recibe tx nueva
    \item A envía announcement a sus peers
    \item Peers que no tienen la tx la solicitan
    \item A envía la tx completa
    \item Peers repiten el proceso
\end{enumerate}
\end{definition}

\begin{mathnote}
\textbf{Optimizaciones de propagación:}

\textbf{Transaction announcements (EIP-2464):}
\begin{itemize}
    \item En lugar de enviar tx completa, enviar solo hash
    \item Receptor pide txs que no tiene
    \item Reduce bandwidth en $\sim$50\%
\end{itemize}

\textbf{Transaction pooling:}
\begin{itemize}
    \item Agrupar múltiples announcements
    \item Reducir overhead de mensajes
\end{itemize}

\textbf{Latencia típica de propagación:}
\begin{itemize}
    \item A mayoría de nodos: $<$ 1 segundo
    \item A 95\% de la red: $<$ 3 segundos
    \item Depende de la conectividad
\end{itemize}
\end{mathnote}

%-------------------------------------------
\subsection{Mempool Privadas y PBS}

\begin{definition}[Private Mempool]
Mempool no pública donde las transacciones no se propagan por gossip.

\textbf{Ejemplos:}
\begin{itemize}
    \item Flashbots Protect: txs van directo a builders
    \item MEV Blocker: coalition de builders éticos
    \item Private RPC endpoints de builders
\end{itemize}
\end{definition}

\begin{keyidea}
\textbf{Proposer-Builder Separation (PBS):}

En Ethereum post-merge con MEV-Boost:

\begin{center}
\begin{tikzpicture}[
    actor/.style={draw, rectangle, minimum width=2cm, minimum height=0.8cm, align=center},
    arrow/.style={->, thick}
]
    % Actores
    \node[actor, fill=blue!20] (user) at (0, 0) {Usuario};
    \node[actor, fill=yellow!20] (builder) at (4, 0) {Builder};
    \node[actor, fill=green!20] (relay) at (8, 0) {Relay};
    \node[actor, fill=red!20] (proposer) at (12, 0) {Proposer\\(Validador)};
    
    % Flechas
    \draw[arrow] (user) -- node[above, font=\small] {tx} (builder);
    \draw[arrow] (builder) -- node[above, font=\small] {bloque} (relay);
    \draw[arrow] (relay) -- node[above, font=\small] {header} (proposer);
    \draw[arrow] (proposer) to[bend right=30] node[below, font=\small] {firma} (relay);
    \draw[arrow] (relay) to[bend right=30] node[below, font=\small] {bloque completo} (proposer);
\end{tikzpicture}
\end{center}

\textbf{Flujo:}
\begin{enumerate}
    \item Users envían txs a builders (privado o público)
    \item Builders construyen bloques optimizados
    \item Builders envían bloques a relays
    \item Relays muestran solo headers a proposers
    \item Proposer firma el header más rentable
    \item Relay revela el bloque completo
    \item Proposer lo propone a la red
\end{enumerate}

\textbf{Beneficios:}
\begin{itemize}
    \item Proposers no necesitan construir bloques
    \item Builders compiten por eficiencia
    \item MEV se redistribuye más equitativamente
\end{itemize}
\end{keyidea}

%===========================================
\section{MEV (Maximal Extractable Value)}
%===========================================

\begin{definition}[MEV]
Valor que productores de bloques pueden extraer reordenando, incluyendo o excluyendo transacciones.
\end{definition}

\begin{intuicion}
Imagina que eres el dueño de un restaurante y puedes decidir el orden en que atiendes a los clientes. Si un cliente quiere comprar el último pastel, y ves que otro cliente también lo quiere, podrías:
\begin{enumerate}
    \item Comprarlo tú primero
    \item Vendérselo al que pague más
    \item O simplemente cobrar una ``cuota de prioridad''
\end{enumerate}

En Ethereum, el ``dueño del restaurante'' es el builder/validator, y las ``órdenes'' son transacciones. MEV es el valor que pueden extraer por controlar el orden.
\end{intuicion}

%-------------------------------------------
\subsection{Tipos de MEV}

\begin{definition}[Tipos de MEV]
\textbf{1. Frontrunning:}
Ejecutar ANTES de una tx víctima.
\begin{verbatim}
Víctima: comprar 100 ETH de TOKEN a precio X
Attacker: compra TOKEN antes → sube precio
Víctima: compra a precio X + Δ (peor precio)
Attacker: vende TOKEN con profit
\end{verbatim}

\textbf{2. Backrunning:}
Ejecutar INMEDIATAMENTE DESPUÉS de una tx.
\begin{verbatim}
Tx legítima: gran compra que mueve precio
Backrunner: arbitraje entre DEXs aprovechando el movimiento
\end{verbatim}

\textbf{3. Sandwich Attack:}
Frontrun + Backrun combinados.
\begin{verbatim}
1. Attacker: compra TOKEN (frontrun)
2. Víctima: compra TOKEN (precio más alto)
3. Attacker: vende TOKEN (backrun, profit)
\end{verbatim}

\textbf{4. Liquidations:}
Competir por liquidar posiciones en lending protocols.
\begin{verbatim}
Posición underwater detectada
Múltiples bots compiten por liquidarla
El más rápido/mejor posicionado gana la recompensa
\end{verbatim}

\textbf{5. Arbitraje:}
Explotar diferencias de precio entre mercados.
\begin{verbatim}
ETH/USDC en Uniswap: 2000
ETH/USDC en Sushiswap: 2010
Arbitrageur: compra en Uni, vende en Sushi
Profit: $10 por ETH (menos fees)
\end{verbatim}
\end{definition}

%-------------------------------------------
\subsection{Anatomía de un Sandwich Attack}

\begin{example}[Sandwich Attack detallado]
\textbf{Escenario:}
\begin{itemize}
    \item Alice quiere comprar TOKEN con 10 ETH en Uniswap
    \item Slippage tolerance: 1\%
    \item Pool: 1000 ETH / 1,000,000 TOKEN (precio: 1000 TOKEN/ETH)
\end{itemize}

\textbf{Sin ataque:}
\begin{verbatim}
Alice compra con 10 ETH
Pool después: 1010 ETH / 990,099 TOKEN
Alice recibe: ~9,901 TOKEN
Precio efectivo: ~990 TOKEN/ETH (0.99% slippage)
\end{verbatim}

\textbf{Con sandwich:}
\begin{verbatim}
1. Frontrun - Attacker compra con 5 ETH:
   Pool: 1005 ETH / 995,025 TOKEN
   Attacker recibe: ~4,975 TOKEN
   
2. Víctima (Alice) compra con 10 ETH:
   Pool: 1015 ETH / 985,197 TOKEN
   Alice recibe: ~9,828 TOKEN (peor que sin ataque)
   
3. Backrun - Attacker vende ~4,975 TOKEN:
   Pool: 1009.95 ETH / 990,172 TOKEN
   Attacker recibe: ~5.05 ETH
   
Resultado:
- Alice: recibió 9,828 en vez de 9,901 TOKEN (73 TOKEN menos)
- Attacker: gastó 5 ETH, recibió 5.05 ETH = 0.05 ETH profit
- Attacker también pagó gas, pero con bundles eficientes es rentable
\end{verbatim}
\end{example}

%-------------------------------------------
\subsection{Ecosistema MEV}

\begin{definition}[Actores del Ecosistema MEV]
\textbf{1. Searchers:}
\begin{itemize}
    \item Encuentran oportunidades de MEV
    \item Escriben bots y algoritmos
    \item Compiten enviando bundles a builders
    \item Pagan ``tips'' a builders por inclusión
\end{itemize}

\textbf{2. Builders:}
\begin{itemize}
    \item Construyen bloques óptimos
    \item Reciben txs de searchers y mempool pública
    \item Ordenan para maximizar valor
    \item Compiten por que proposers elijan su bloque
\end{itemize}

\textbf{3. Relays:}
\begin{itemize}
    \item Intermediarios entre builders y proposers
    \item Verifican bloques de builders
    \item Protegen a proposers de bloques inválidos
    \item Ejemplos: Flashbots, BloXroute, Ultrasound
\end{itemize}

\textbf{4. Proposers (Validators):}
\begin{itemize}
    \item Proponen bloques a la red
    \item Eligen el bloque más rentable de los relays
    \item Reciben pago del builder
    \item No necesitan expertise en MEV
\end{itemize}
\end{definition}

\begin{center}
\begin{tikzpicture}[
    actor/.style={draw, ellipse, minimum width=2cm, minimum height=1cm, align=center},
    arrow/.style={->, thick},
    label/.style={font=\small}
]
    % Usuarios
    \node[actor, fill=blue!20] (users) at (0, 3) {Users};
    
    % Searchers
    \node[actor, fill=yellow!20] (searchers) at (-3, 1) {Searchers};
    
    % Builders
    \node[actor, fill=orange!20] (builders) at (3, 1) {Builders};
    
    % Relays
    \node[actor, fill=green!20] (relays) at (0, -1) {Relays};
    
    % Proposers
    \node[actor, fill=red!20] (proposers) at (0, -3.5) {Proposers};
    
    % Blockchain
    \node[draw, rectangle, fill=gray!20, minimum width=3cm] (chain) at (0, -5.5) {Ethereum};
    
    % Flechas
    \draw[arrow] (users) -- node[label, right] {txs} (builders);
    \draw[arrow] (users) -- node[label, left] {txs} (searchers);
    \draw[arrow] (searchers) -- node[label, above] {bundles + tips} (builders);
    \draw[arrow] (builders) -- node[label, right] {bloques} (relays);
    \draw[arrow] (relays) -- node[label, right] {headers + bids} (proposers);
    \draw[arrow] (proposers) -- node[label, right] {bloque firmado} (chain);
\end{tikzpicture}
\end{center}

%-------------------------------------------
\subsection{Flashbots}

\begin{definition}[Flashbots]
Organización e infraestructura para MEV ``fair''.

\textbf{Productos:}
\begin{itemize}
    \item \textbf{Flashbots Auction}: Mercado de bundles
    \item \textbf{Flashbots Protect}: RPC que protege de sandwich
    \item \textbf{MEV-Boost}: Software para PBS
    \item \textbf{MEV-Share}: Redistribución de MEV a users
    \item \textbf{SUAVE}: Futuro sistema de encrypted mempools
\end{itemize}
\end{definition}

\begin{keyidea}
\textbf{Flashbots Protect - Cómo funciona:}

\begin{enumerate}
    \item User envía tx a Flashbots RPC (no mempool pública)
    \item Tx va directo a builders de Flashbots
    \item Builders incluyen sin posibilidad de sandwich
    \item Si tx no es incluida en X bloques, se puede enviar a mempool pública
\end{enumerate}

\textbf{Beneficios:}
\begin{itemize}
    \item No aparece en mempool pública = no puede ser sandwiched
    \item Si tx falla, no pagas gas (no hay revert público)
    \item Gratis para usuarios
\end{itemize}

\textbf{Trade-offs:}
\begin{itemize}
    \item Confías en Flashbots para no hacer MEV ellos mismos
    \item Puede ser más lento (menos builders)
    \item Centralización hacia Flashbots
\end{itemize}
\end{keyidea}

%-------------------------------------------
\subsection{MEV por los Números}

\begin{mathnote}
\textbf{Estadísticas de MEV (2023-2024):}

\begin{center}
\begin{tabular}{lr}
\toprule
\textbf{Métrica} & \textbf{Valor} \\
\midrule
MEV total extraído (histórico) & $>\$700M$ \\
MEV diario promedio & $\sim\$1-5M$ \\
\% bloques via MEV-Boost & $\sim90\%$ \\
Builders activos & $\sim10-20$ \\
Top builder market share & $\sim40-60\%$ \\
Sandwich attacks diarios & $\sim1000-5000$ \\
\bottomrule
\end{tabular}
\end{center}

\textbf{Distribución típica del MEV:}
\begin{itemize}
    \item Searchers: $\sim$10-30\% (compiten fuertemente)
    \item Builders: $\sim$5-15\%
    \item Proposers/Validators: $\sim$60-80\%
    \item La competencia ha comprimido márgenes
\end{itemize}
\end{mathnote}

%-------------------------------------------
\subsection{Impacto Negativo del MEV}

\begin{advertencia}
\textbf{Problemas causados por MEV:}

\textbf{1. Peor ejecución para usuarios:}
\begin{itemize}
    \item Sandwich attacks roban valor de swaps
    \item Usuarios pagan más slippage del necesario
    \item Estimación: usuarios pierden \$1-10 por swap afectado
\end{itemize}

\textbf{2. Centralización:}
\begin{itemize}
    \item Economías de escala favorecen grandes operadores
    \item Top 3 builders dominan $>$80\% del mercado
    \item Proposers pequeños no pueden competir
\end{itemize}

\textbf{3. Congestión de red:}
\begin{itemize}
    \item Bots compiten con txs de spam
    \item Gas prices artificialmente altos durante oportunidades
    \item Failed txs desperdician blockspace
\end{itemize}

\textbf{4. Inestabilidad del consenso:}
\begin{itemize}
    \item Incentivo para reorgs si MEV $>$ block reward
    \item ``Time-bandit attacks'' teóricamente posibles
    \item Reduce seguridad de finalidad
\end{itemize}
\end{advertencia}

%-------------------------------------------
\subsection{Soluciones y Mitigaciones}

\begin{keyidea}
\textbf{Estrategias contra MEV:}

\textbf{1. Private mempools:}
\begin{itemize}
    \item Flashbots Protect, MEV Blocker
    \item Txs no visibles hasta inclusión
    \item Elimina frontrunning
\end{itemize}

\textbf{2. Commit-reveal schemes:}
\begin{itemize}
    \item Usuario envía hash de tx primero
    \item Revela contenido después de inclusión
    \item Ejemplo: Submarine sends
\end{itemize}

\textbf{3. Batch auctions:}
\begin{itemize}
    \item Agrupar órdenes y ejecutar al mismo precio
    \item CoW Protocol, 1inch Fusion
    \item Elimina ventaja de orden
\end{itemize}

\textbf{4. Encrypted mempools:}
\begin{itemize}
    \item Txs encriptadas hasta threshold de validadores
    \item Shutter Network, SUAVE (futuro)
    \item Nadie puede ver contenido antes de ordenar
\end{itemize}

\textbf{5. MEV redistribution:}
\begin{itemize}
    \item MEV-Share: devolver parte a usuarios
    \item MEV taxes: protocolo captura MEV
    \item Reducir daño aunque no eliminar
\end{itemize}
\end{keyidea}

%-------------------------------------------
\subsection{MEV-Share}

\begin{definition}[MEV-Share]
Protocolo de Flashbots para redistribuir MEV a usuarios.

\textbf{Funcionamiento:}
\begin{enumerate}
    \item Usuario envía tx a Flashbots con hints
    \item Hints revelan info parcial (ej: ``swap en Uniswap'')
    \item Searchers pueden backrun basándose en hints
    \item MEV generado se divide entre user y searcher
\end{enumerate}

\textbf{Ejemplo:}
\begin{verbatim}
User swap genera $100 de oportunidad de arbitraje
Sin MEV-Share: Searcher toma $100, user recibe $0
Con MEV-Share: Searcher toma $20, user recibe $80
\end{verbatim}
\end{definition}

%-------------------------------------------
\subsection{Aplicación en Banco ZK}

\begin{bancozk}
\textbf{La privacidad de Banco ZK mitiga MEV:}

\textbf{1. Montos ocultos:}
\begin{itemize}
    \item Transacciones no revelan cantidades
    \item No se puede calcular profit de sandwich
    \item Frontrunning es inútil sin saber el impacto
\end{itemize}

\textbf{2. Intenciones ocultas:}
\begin{itemize}
    \item No se sabe qué quiere hacer el usuario
    \item ¿Es un swap? ¿Una transferencia? ¿Un depósito?
    \item Sin esta info, no hay MEV que extraer
\end{itemize}

\textbf{3. Encrypted mempools nativos:}
\begin{itemize}
    \item Las txs están encriptadas por diseño
    \item Sequencer no puede ver contenido
    \item Ordenamiento ``fair'' por defecto
\end{itemize}

\textbf{4. Batch processing:}
\begin{itemize}
    \item Múltiples txs se procesan juntas
    \item Orden dentro del batch puede ser aleatorio
    \item Reduce ventaja de timing
\end{itemize}

\textbf{Trade-offs:}
\begin{itemize}
    \item Privacidad tiene costo computacional (ZK proofs)
    \item Algunas operaciones no pueden ser privadas (bridges)
    \item MEV del sequencer aún posible en teoría
\end{itemize}

\textbf{Diseño del sequencer:}
\begin{verbatim}
// El sequencer de Banco ZK:
1. Recibe txs encriptadas de usuarios
2. No puede ver contenido (solo metadata mínima)
3. Ordena por timestamp de llegada (FIFO)
4. Genera prueba ZK del batch
5. Publica batch + proof en L1

// Incluso si quisiera hacer MEV:
- No sabe qué hacen las txs
- No puede reordenar basándose en contenido
- Solo puede censurar (detectable y punishable)
\end{verbatim}

\textbf{Mitigación de censura:}
\begin{itemize}
    \item Múltiples sequencers en rotación (futuro)
    \item Forced inclusion via L1 si sequencer censura
    \item Incentivos alineados: sequencer gana fees, no MEV
\end{itemize}
\end{bancozk}

%===========================================
\section{Resumen del Bloque 2}
%===========================================

\begin{keyidea}
\textbf{Puntos clave de Blockchain Base (L1):}

\textbf{1. L1 proporciona cuatro funciones:}
\begin{itemize}
    \item Consenso sobre el estado
    \item Disponibilidad de datos
    \item Ejecución de transacciones
    \item Settlement final
\end{itemize}

\textbf{2. Consenso:}
\begin{itemize}
    \item PoW: seguridad por energía gastada, 51\% attack
    \item PoS: seguridad por capital en stake, slashing
    \item PoA: validadores autorizados, centralizado pero rápido
\end{itemize}

\textbf{3. Modelos de transacciones:}
\begin{itemize}
    \item UTXO (Bitcoin): outputs gastables, privacy natural
    \item Account (Ethereum): balances persistentes, smart contracts
\end{itemize}

\textbf{4. Gas y fees:}
\begin{itemize}
    \item Gas mide costo computacional
    \item EIP-1559: base fee (quemado) + priority fee
    \item Storage es muy caro, optimizar es crucial
\end{itemize}

\textbf{5. EVM:}
\begin{itemize}
    \item Máquina virtual stack-based
    \item $\sim$140 opcodes
    \item Precompilados para operaciones costosas
    \item Determinística: todos los nodos llegan al mismo resultado
\end{itemize}

\textbf{6. Estado:}
\begin{itemize}
    \item Merkle Patricia Trie
    \item EOAs y Contract accounts
    \item Storage layout importa para gas
\end{itemize}

\textbf{7. Nodos:}
\begin{itemize}
    \item Full: estado actual, valida todo
    \item Light: solo headers, usa proofs
    \item Archive: todo el historial
    \item Validator: participa en consenso
\end{itemize}

\textbf{8. Mempool:}
\begin{itemize}
    \item Pool de txs pendientes
    \item Pública por defecto (problema para MEV)
    \item PBS separa construcción de proposición
\end{itemize}

\textbf{9. MEV:}
\begin{itemize}
    \item Valor extraíble por reordenamiento
    \item Tipos: frontrunning, sandwich, arbitraje, liquidaciones
    \item Soluciones: private mempools, batch auctions, encryption
    \item Privacidad (como en Banco ZK) mitiga MEV naturalmente
\end{itemize}
\end{keyidea}

\cleardoublepage

% ============================================
% BLOQUE 3: ESCALABILIDAD (L2/L3) - VERSIÓN EXPANDIDA
% Parte 1: Trilema, L2, Sidechains vs Rollups
% ============================================
\chapter{Escalabilidad (L2/L3)}
\label{ch:bloque3}

\begin{center}
\textit{Bloque 3 de 20}
\end{center}

%===========================================
\section{El Problema de Escalabilidad (Trilema)}
%===========================================

\begin{definition}[Trilema de Escalabilidad]
Es difícil optimizar simultáneamente:
\begin{itemize}
    \item Escalabilidad (throughput)
    \item Descentralización (número de nodos)
    \item Seguridad (costo de ataque)
\end{itemize}
\end{definition}

Throughput actual:
\begin{itemize}
    \item Bitcoin: $\sim$7 TPS
    \item Ethereum: $\sim$15 TPS
    \item Visa: $\sim$65,000 TPS
\end{itemize}

%-------------------------------------------
\subsection{Historia del Trilema}

\begin{remark}
El término ``trilema de escalabilidad'' fue popularizado por Vitalik Buterin alrededor de 2017, aunque el concepto existía antes en teoría de sistemas distribuidos.

\textbf{La intuición:}
\begin{itemize}
    \item Para más throughput: procesar más txs por segundo
    \item Esto requiere: nodos más potentes O menos validación
    \item Nodos más potentes = menos gente puede participar = centralización
    \item Menos validación = más fácil de atacar = menos seguridad
\end{itemize}

No es un teorema matemático estricto, pero captura un trade-off fundamental en el diseño de blockchains.
\end{remark}

\begin{center}
\begin{tikzpicture}[
    vertex/.style={draw, circle, minimum size=2.5cm, align=center, font=\small},
    edge/.style={thick}
]
    % Triángulo
    \node[vertex, fill=blue!20] (scale) at (90:3) {Escalabilidad\\(TPS)};
    \node[vertex, fill=green!20] (decent) at (210:3) {Descentralización\\(\# nodos)};
    \node[vertex, fill=red!20] (secure) at (330:3) {Seguridad\\(costo ataque)};
    
    % Aristas
    \draw[edge] (scale) -- (decent);
    \draw[edge] (decent) -- (secure);
    \draw[edge] (secure) -- (scale);
    
    % Centro
    \node at (0, 0) {\textbf{Trilema}};
    
    % Ejemplos
    \node[font=\tiny, right] at (1.5, 2) {Solana: alto TPS};
    \node[font=\tiny, right] at (1.5, 1.5) {pero centralizado};
    
    \node[font=\tiny, left] at (-1.5, -1) {Bitcoin: muy seguro};
    \node[font=\tiny, left] at (-1.5, -1.5) {y descentralizado};
    \node[font=\tiny, left] at (-1.5, -2) {pero 7 TPS};
\end{tikzpicture}
\end{center}

%-------------------------------------------
\subsection{Análisis de las Tres Propiedades}

\begin{definition}[Escalabilidad]
Capacidad de procesar un número creciente de transacciones.

\textbf{Métricas:}
\begin{itemize}
    \item TPS (Transactions Per Second): transacciones por segundo
    \item Throughput: datos procesados por segundo
    \item Latency: tiempo desde envío hasta confirmación
    \item Finality: tiempo hasta que tx es irreversible
\end{itemize}
\end{definition}

\begin{mathnote}
\textbf{¿Por qué Ethereum solo hace $\sim$15 TPS?}

\textbf{Límites actuales:}
\begin{itemize}
    \item Block gas limit: 30M gas
    \item Block time: 12 segundos
    \item Tx simple (transfer): 21,000 gas
\end{itemize}

\textbf{Cálculo:}
\[ \text{TPS}_{\text{max}} = \frac{30,000,000 \text{ gas}}{12 \text{ s} \times 21,000 \text{ gas/tx}} \approx 119 \text{ TPS} \]

Pero en la práctica:
\begin{itemize}
    \item Muchas txs son contratos complejos ($>$100K gas)
    \item El promedio real es $\sim$15-30 TPS
    \item Con más gas por tx = menos TPS
\end{itemize}

\textbf{¿Por qué no simplemente aumentar el gas limit?}
\begin{itemize}
    \item Bloques más grandes = más tiempo para propagar
    \item Más validación = nodos más caros
    \item Estado crece más rápido = storage más caro
    \item Centralización hacia operadores grandes
\end{itemize}
\end{mathnote}

\begin{definition}[Descentralización]
Distribución del poder de decisión y validación entre múltiples entidades independientes.

\textbf{Métricas:}
\begin{itemize}
    \item Número de nodos full
    \item Nakamoto coefficient (mínimo para 51\%)
    \item Distribución geográfica
    \item Diversidad de clientes (software)
    \item Requisitos de hardware para participar
\end{itemize}
\end{definition}

\begin{example}[Requisitos de hardware y descentralización]
\begin{center}
\begin{tabular}{lcccc}
\toprule
\textbf{Red} & \textbf{RAM} & \textbf{Storage} & \textbf{Bandwidth} & \textbf{Nodos} \\
\midrule
Bitcoin & 2 GB & 500 GB & 5 Mbps & $\sim$15,000 \\
Ethereum & 16 GB & 1 TB & 25 Mbps & $\sim$6,000 \\
Solana & 128 GB & 2 TB & 1 Gbps & $\sim$1,900 \\
BSC & 32 GB & 2 TB & 100 Mbps & $\sim$21 (validators) \\
\bottomrule
\end{tabular}
\end{center}

Solana tiene alto TPS pero pocos pueden correr nodos. Bitcoin tiene bajo TPS pero cualquiera puede validar con hardware modesto.
\end{example}

\begin{definition}[Seguridad]
Resistencia a ataques que intenten modificar el estado o censurar transacciones.

\textbf{Métricas:}
\begin{itemize}
    \item Costo de ataque 51\%
    \item Tiempo para finalidad
    \item Resistencia a censura
    \item Historial de ataques exitosos
\end{itemize}
\end{definition}

%-------------------------------------------
\subsection{Enfoques para Escalar}

\begin{keyidea}
\textbf{Dos filosofías de escalabilidad:}

\textbf{1. Escalar L1 (on-chain):}
\begin{itemize}
    \item Aumentar block size/gas limit
    \item Reducir block time
    \item Optimizar ejecución
    \item Sharding del estado
    \item Trade-off: sacrificar algo del trilema
\end{itemize}

\textbf{2. Escalar off-chain (L2):}
\begin{itemize}
    \item Procesar txs fuera de L1
    \item Solo publicar resultados/proofs en L1
    \item Heredar seguridad de L1
    \item Trade-off: complejidad adicional
\end{itemize}

\textbf{Ethereum eligió el camino L2:}

``The Rollup-Centric Roadmap'' (2020):
\begin{itemize}
    \item L1 optimizado para seguridad y DA
    \item L2s manejan la ejecución masiva
    \item Mejor de ambos mundos
\end{itemize}
\end{keyidea}

\begin{remark}
\textbf{Historia de los intentos de escalar Ethereum:}

\begin{center}
\begin{tabular}{lll}
\toprule
\textbf{Año} & \textbf{Propuesta} & \textbf{Estado} \\
\midrule
2017 & State channels (Raiden) & Abandonado mayormente \\
2017 & Plasma & Abandonado (data availability) \\
2018 & Sharding original & Reemplazado por rollup-centric \\
2019 & Optimistic Rollups & Activo (Arbitrum, Optimism) \\
2019 & ZK Rollups & Activo (zkSync, StarkNet) \\
2020 & Rollup-centric roadmap & Adoptado oficialmente \\
2024 & EIP-4844 (blobs) & Implementado \\
202X & Full Danksharding & En desarrollo \\
\bottomrule
\end{tabular}
\end{center}
\end{remark}

%-------------------------------------------
\subsection{El Throughput en Perspectiva}

\begin{mathnote}
\textbf{¿Cuánto necesitamos realmente?}

\textbf{Demanda actual:}
\begin{itemize}
    \item Visa procesa $\sim$65,000 TPS pico, $\sim$1,700 TPS promedio
    \item Ethereum: $\sim$1.2M txs/día $\approx$ 14 TPS promedio
    \item En picos de demanda: fees de \$50-500 por tx
\end{itemize}

\textbf{Demanda futura (si crypto escala):}
\begin{itemize}
    \item 8 mil millones de personas
    \item Si cada una hace 10 txs/día
    \item = 80 mil millones de txs/día
    \item = 925,000 TPS
\end{itemize}

\textbf{¿Pueden los L2s lograrlo?}

\begin{center}
\begin{tabular}{lccc}
\toprule
\textbf{Escenario} & \textbf{TPS L2} & \textbf{\# L2s} & \textbf{TPS Total} \\
\midrule
Hoy & 100-1000 & 10 & $\sim$10,000 \\
2025 & 1000-5000 & 100 & $\sim$500,000 \\
2030 (proyección) & 10,000+ & 1000+ & $\sim$10,000,000 \\
\bottomrule
\end{tabular}
\end{center}

Con múltiples L2s, L3s, y optimizaciones, el throughput agregado puede escalar dramáticamente.
\end{mathnote}

%===========================================
\section{Qué es L2}
%===========================================

\begin{definition}[Layer 2]
Protocolo que procesa transacciones off-chain pero hereda seguridad de L1.
\end{definition}

Tipos:
\begin{itemize}
    \item State channels
    \item Plasma
    \item Rollups (Optimistic y ZK)
\end{itemize}

%-------------------------------------------
\subsection{La Definición Precisa de L2}

\begin{intuicion}
Un L2 es como una sucursal de un banco:
\begin{itemize}
    \item Puedes hacer transacciones en la sucursal (rápido, barato)
    \item La sucursal reporta al banco central (L1)
    \item Si la sucursal hace trampa, el banco central puede corregirlo
    \item Tu dinero está seguro porque está respaldado por el banco central
\end{itemize}

La clave es que \textbf{heredas la seguridad de L1}. No necesitas confiar en la sucursal si puedes verificar con el banco central.
\end{intuicion}

\begin{definition}[Criterios para ser ``verdadero L2'']
Según la comunidad de Ethereum, un L2 debe:

\textbf{1. Heredar seguridad de L1:}
\begin{itemize}
    \item Un usuario puede forzar retiro solo interactuando con L1
    \item No requiere cooperación del operador L2
    \item L1 puede verificar correctitud del estado L2
\end{itemize}

\textbf{2. Publicar datos en L1:}
\begin{itemize}
    \item Datos necesarios para reconstruir estado
    \item O prueba de que datos están disponibles
    \item Permite validación independiente
\end{itemize}

\textbf{3. No tener custodio confiable:}
\begin{itemize}
    \item Fondos no controlados por multisig
    \item O escape hatch si multisig es malicioso
\end{itemize}
\end{definition}

\begin{advertencia}
\textbf{¿Es realmente un L2?}

Muchos proyectos se llaman ``L2'' pero no cumplen todos los criterios:

\begin{center}
\begin{tabular}{lccc}
\toprule
\textbf{Proyecto} & \textbf{Hereda seg.} & \textbf{DA en L1} & \textbf{Sin custodio} \\
\midrule
Arbitrum & Sí & Sí & Parcial* \\
Optimism & Sí & Sí & Parcial* \\
zkSync Era & Sí & Sí & Parcial* \\
Polygon PoS & No (propia) & No & No \\
Ronin & No (propia) & No & No \\
BSC & No (L1 separada) & N/A & N/A \\
\bottomrule
\end{tabular}
\end{center}

*Los rollups actuales aún tienen ``training wheels'': multisigs que pueden pausar o actualizar. Ver L2Beat para estado actual de descentralización.
\end{advertencia}

%-------------------------------------------
\subsection{Tipos de Soluciones L2}

\subsubsection{State Channels}

\begin{definition}[State Channel]
Canal bilateral o multilateral donde partes hacen transacciones off-chain, solo tocando L1 para abrir y cerrar el canal.
\end{definition}

\begin{example}[Lightning Network (Bitcoin)]
\begin{verbatim}
1. Alice y Bob abren canal depositando 1 BTC cada uno (tx L1)
   Estado inicial: Alice=1, Bob=1

2. Off-chain, intercambian estados firmados:
   Estado 1: Alice=0.9, Bob=1.1 (Alice paga 0.1 a Bob)
   Estado 2: Alice=0.8, Bob=1.2 (Alice paga 0.1 más)
   ...
   Estado 100: Alice=0.3, Bob=1.7

3. Pueden hacer miles de txs sin tocar L1

4. Para cerrar: publican último estado en L1 (1 tx)
   Alice recibe 0.3 BTC, Bob recibe 1.7 BTC

5. Si alguien intenta publicar estado viejo:
   - Periodo de disputa permite al otro probar fraude
   - El fraudulento pierde todo su balance
\end{verbatim}
\end{example}

\begin{keyidea}
\textbf{Ventajas de State Channels:}
\begin{itemize}
    \item Latencia instantánea (sin esperar bloques)
    \item Costo por tx: prácticamente cero
    \item Privacidad: txs intermedias no se publican
    \item Finalidad: inmediata entre las partes
\end{itemize}

\textbf{Desventajas:}
\begin{itemize}
    \item Requiere capital bloqueado
    \item Solo funciona entre partes del canal
    \item Routing complejo para pagos multi-hop
    \item Requiere estar online para disputar
    \item No soporta smart contracts generales
\end{itemize}
\end{keyidea}

\subsubsection{Plasma}

\begin{definition}[Plasma]
Cadena hija que publica merkle roots de su estado en L1, con mecanismo de exits en caso de fraude.
\end{definition}

\begin{remark}
Plasma fue propuesto en 2017 por Vitalik Buterin y Joseph Poon (co-autor de Lightning).

\textbf{El problema fatal: Data Availability}

En Plasma, los datos completos NO van a L1 (solo merkle roots). Esto crea el problema de ``data withholding'':
\begin{enumerate}
    \item Operador publica root válido
    \item Pero NO publica los datos
    \item Usuarios no pueden probar qué poseen
    \item No pueden hacer exit correctamente
\end{enumerate}

\textbf{Solución intentada: Mass exits}

Si operador no da datos, todos hacen exit simultáneo. Pero esto:
\begin{itemize}
    \item Congestiona L1
    \item Es costosísimo
    \item Crea pánico
\end{itemize}

Por estas razones, Plasma fue mayormente abandonado en favor de Rollups.
\end{remark}

\subsubsection{Rollups}

\begin{definition}[Rollup]
L2 que ``enrolla'' múltiples transacciones en un batch, publica datos en L1, y usa pruebas para verificar correctitud.
\end{definition}

\begin{keyidea}
\textbf{¿Por qué Rollups funcionan mejor que Plasma?}

\textbf{Rollups publican datos en L1:}
\begin{itemize}
    \item Cualquiera puede reconstruir el estado
    \item No hay problema de data withholding
    \item Exits siempre son posibles
\end{itemize}

\textbf{Compresión hace que sea económico:}
\begin{itemize}
    \item Una tx Ethereum: $\sim$110 bytes
    \item Una tx en rollup (compressed): $\sim$12 bytes
    \item 10x menos espacio = 10x más txs por bloque
\end{itemize}

\textbf{Verificación es barata:}
\begin{itemize}
    \item Optimistic: solo verificar si hay disputa
    \item ZK: verificar proof es $O(1)$
\end{itemize}
\end{keyidea}

%-------------------------------------------
\subsection{Comparación de Tipos de L2}

\begin{center}
\begin{tabular}{lcccc}
\toprule
\textbf{Propiedad} & \textbf{Channels} & \textbf{Plasma} & \textbf{Opt. Rollup} & \textbf{ZK Rollup} \\
\midrule
DA & Off-chain & Off-chain & On-chain & On-chain \\
Finalidad L1 & Inmediata & $\sim$7 días & $\sim$7 días & Minutos \\
Smart contracts & Limitados & Limitados & Completos & Completos* \\
Capital locked & Sí & No & No & No \\
Online required & Sí & Sí & No & No \\
Complejidad & Media & Alta & Media & Alta \\
Madurez & Alta & Abandonado & Alta & Media \\
\bottomrule
\end{tabular}
\end{center}

*ZK Rollups para smart contracts generales (zkEVM) son relativamente nuevos.

%===========================================
\section{Sidechains vs Rollups}
%===========================================

\begin{table}[htbp]
\centering
\begin{tabular}{lll}
\toprule
\textbf{Aspecto} & \textbf{Sidechain} & \textbf{Rollup} \\
\midrule
Seguridad & Propia (validadores) & Hereda de L1 \\
Datos & En la sidechain & En L1 \\
Trust & Confiar en validadores & Trustless (cripto) \\
Bridge & Requiere confianza & Native, seguro \\
\bottomrule
\end{tabular}
\end{table}

Sidechains: Polygon PoS, Ronin, Gnosis Chain.

%-------------------------------------------
\subsection{Qué es una Sidechain}

\begin{definition}[Sidechain]
Blockchain independiente con su propio consenso, conectada a otra blockchain mediante un bridge.
\end{definition}

\begin{intuicion}
Una sidechain es como un país extranjero:
\begin{itemize}
    \item Tiene sus propias leyes (consenso)
    \item Su propia policía (validadores)
    \item Su propia moneda puede estar ``pegged'' a otra
    \item Puedes ir y venir (bridge), pero si el país colapsa, pierdes todo
\end{itemize}

Un rollup es más como una provincia autónoma:
\begin{itemize}
    \item Tiene cierta autonomía (procesa sus txs)
    \item Pero está bajo la jurisdicción del país central (L1)
    \item Si hay problemas, el gobierno central interviene
    \item Tu propiedad está protegida por ley nacional (L1)
\end{itemize}
\end{intuicion}

%-------------------------------------------
\subsection{Análisis de Seguridad}

\begin{mathnote}
\textbf{Modelo de seguridad de Sidechain:}

Sea $n$ el número de validadores y $f$ los maliciosos.

Para BFT típico (Tendermint, etc.):
\begin{itemize}
    \item Requiere $f < n/3$ para safety
    \item Si $f \geq n/3$: pueden detener la cadena
    \item Si $f \geq 2n/3$: pueden robar fondos
\end{itemize}

\textbf{Ejemplo: Ronin}
\begin{itemize}
    \item 9 validadores, threshold 5/9
    \item 4 controlados por Sky Mavis
    \item 1 más comprometido por phishing
    \item Resultado: \$625M robados
\end{itemize}

\textbf{Modelo de seguridad de Rollup:}

\begin{itemize}
    \item 1 participante honesto puede forzar correctitud
    \item Operador malicioso: usuarios hacen exit via L1
    \item Seguridad = seguridad de L1 (miles de validadores)
\end{itemize}
\end{mathnote}

\begin{advertencia}
\textbf{El problema de los bridges de sidechains:}

\begin{enumerate}
    \item Usuario deposita ETH en bridge de sidechain
    \item Bridge mintea ``wrapped ETH'' en sidechain
    \item Usuario usa wrapped ETH normalmente
    \item Para salir: quemar wrapped, bridge libera ETH real
\end{enumerate}

\textbf{¿Quién controla el ETH real?}

Típicamente un multisig de validadores de la sidechain.

Si el multisig es comprometido:
\begin{itemize}
    \item Pueden robar todo el ETH del bridge
    \item Los wrapped tokens en sidechain se vuelven sin respaldo
    \item Usuarios pierden todo
\end{itemize}

\textbf{Hacks de bridges de sidechains:}

\begin{center}
\begin{tabular}{llr}
\toprule
\textbf{Bridge} & \textbf{Causa} & \textbf{Pérdida} \\
\midrule
Ronin & Multisig comprometido & \$625M \\
Harmony & 2/5 multisig comprometido & \$100M \\
\bottomrule
\end{tabular}
\end{center}
\end{advertencia}

%-------------------------------------------
\subsection{Bridges de Rollups vs Sidechains}

\begin{keyidea}
\textbf{Bridge nativo de Rollup (ejemplo: Arbitrum)}

\textbf{Depósito (L1 → L2):}
\begin{enumerate}
    \item Usuario envía ETH a contrato bridge en L1
    \item Contrato bloquea ETH y emite evento
    \item Sequencer ve evento y acredita ETH en L2
    \item Si sequencer no lo hace: después del challenge period, usuario puede reclamar en L2 forzadamente
\end{enumerate}

\textbf{Retiro (L2 → L1):}
\begin{enumerate}
    \item Usuario inicia retiro en L2
    \item Transacción incluida en batch publicado en L1
    \item Esperar challenge period ($\sim$7 días para Optimistic)
    \item Usuario ejecuta retiro en L1
    \item Contrato libera ETH
\end{enumerate}

\textbf{Seguridad:}
\begin{itemize}
    \item No hay multisig que pueda robar fondos
    \item Contrato en L1 solo libera con proof válido
    \item Usuario siempre puede hacer exit si tiene los datos
\end{itemize}
\end{keyidea}

\begin{center}
\begin{tikzpicture}[
    layer/.style={draw, rectangle, minimum width=10cm, minimum height=2cm, align=center},
    arrow/.style={->, thick},
    label/.style={font=\small}
]
    % L1
    \node[layer, fill=blue!20] (l1) at (0, 0) {Ethereum L1\\(Bridge Contract)};
    
    % Sidechain
    \node[layer, fill=red!20] (side) at (-5, 3.5) {Sidechain\\(Validadores propios)};
    
    % Rollup
    \node[layer, fill=green!20] (rollup) at (5, 3.5) {Rollup\\(Hereda seguridad L1)};
    
    % Bridges
    \draw[arrow, red, very thick] (l1.north west) -- node[left, label] {Multisig\\(trust)} (side.south);
    \draw[arrow, red] (side.south) -- (l1.north west);
    
    \draw[arrow, green!50!black, very thick] (l1.north east) -- node[right, label] {Crypto proof\\(trustless)} (rollup.south);
    \draw[arrow, green!50!black] (rollup.south) -- (l1.north east);
    
    % Labels
    \node[below=0.5cm of l1, font=\small] {Seguridad: ETH stake + miles de validadores};
    \node[above=0.3cm of side, font=\small, red] {Seguridad: N validadores sidechain};
    \node[above=0.3cm of rollup, font=\small, green!50!black] {Seguridad: = L1};
\end{tikzpicture}
\end{center}

%-------------------------------------------
\subsection{¿Cuándo Usar Cada Uno?}

\begin{keyidea}
\textbf{Sidechains son apropiadas cuando:}
\begin{itemize}
    \item Los validadores son confiables (empresa, consorcio)
    \item El valor en juego es limitado
    \item La velocidad es más importante que la seguridad
    \item Se necesita compatibilidad EVM inmediata
    \item Ejemplo: testnet, juegos con items de bajo valor
\end{itemize}

\textbf{Rollups son apropiados cuando:}
\begin{itemize}
    \item Se maneja valor significativo
    \item La seguridad es prioritaria
    \item Se acepta algo de latencia para finalidad
    \item Se quiere heredar seguridad de Ethereum
    \item Ejemplo: DeFi, pagos, cualquier aplicación seria
\end{itemize}
\end{keyidea}

\begin{remark}
\textbf{La evolución de Polygon:}

\begin{center}
\begin{tabular}{lll}
\toprule
\textbf{Producto} & \textbf{Tipo} & \textbf{Seguridad} \\
\midrule
Polygon PoS & Sidechain & Validadores propios \\
Polygon zkEVM & ZK Rollup & Hereda de Ethereum \\
Polygon CDK & Framework & Configurable \\
\bottomrule
\end{tabular}
\end{center}

Polygon comenzó como sidechain pero está migrando hacia ZK rollups para mejor seguridad. Planean ``upgrade'' Polygon PoS a un validium/rollup.
\end{remark}

\begin{bancozk}
\textbf{Banco ZK como Rollup:}

Banco ZK es un \textbf{ZK Rollup}, no una sidechain. Esto significa:

\begin{itemize}
    \item \textbf{Seguridad de Ethereum:} Aunque el sequencer sea malicioso, tus fondos están seguros
    \item \textbf{Datos en L1:} Publicamos datos comprimidos en Ethereum
    \item \textbf{Bridge trustless:} El bridge está asegurado por proofs ZK
    \item \textbf{Forced exits:} Usuarios pueden salir sin cooperación del sequencer
\end{itemize}

La privacidad es ortogonal a esto:
\begin{itemize}
    \item Podríamos ser sidechain privada (menos seguro)
    \item O rollup privado (lo que somos, más seguro)
    \item La privacidad viene de ZK proofs sobre estado encriptado
    \item La seguridad viene de ser rollup con DA en L1
\end{itemize}
\end{bancozk}

% ============================================
% BLOQUE 3: ESCALABILIDAD (L2/L3) - VERSIÓN EXPANDIDA
% Parte 2: Optimistic Rollups
% ============================================

%===========================================
\section{Optimistic Rollups}
%===========================================

\begin{definition}[Optimistic Rollup]
Asume validez de transacciones, permite disputas.
\end{definition}

\begin{intuicion}
``Optimistic'' porque asume que todo es correcto hasta que alguien demuestre lo contrario. Es como un sistema legal:
\begin{itemize}
    \item Las personas son inocentes hasta que se pruebe su culpabilidad
    \item Cualquiera puede acusar (disputar)
    \item El acusador debe probar el fraude
    \item Si hay fraude probado, se castiga al culpable
\end{itemize}

Esto es eficiente porque la mayoría del tiempo no hay fraude, y solo pagamos el costo de verificación cuando es necesario.
\end{intuicion}

%-------------------------------------------
\subsection{Cómo Funcionan}

\begin{enumerate}
    \item Sequencer ejecuta txs off-chain
    \item Publica batch + state root en L1
    \item Período de disputa ($\sim$7 días)
    \item Cualquiera puede disputar con fraud proof
    \item Si no hay disputa, estado es final
\end{enumerate}

\subsubsection{Flujo Detallado}

\begin{definition}[Flujo de Optimistic Rollup]
\textbf{1. Recepción de transacciones:}
\begin{verbatim}
Usuario envía tx al sequencer (via RPC)
    ↓
Sequencer valida tx (firma, balance, nonce)
    ↓
Sequencer añade tx a su mempool local
    ↓
Sequencer ordena txs y las ejecuta
\end{verbatim}

\textbf{2. Ejecución y batching:}
\begin{verbatim}
Sequencer ejecuta txs secuencialmente
    ↓
Calcula nuevo state root después de cada tx
    ↓
Acumula txs en un batch (ej: cada 1-10 minutos)
    ↓
Comprime datos del batch
\end{verbatim}

\textbf{3. Publicación en L1:}
\begin{verbatim}
Sequencer envía tx a L1 con:
  - Compressed batch data (calldata o blobs)
  - Previous state root
  - New state root (claimed)
  - Batch metadata
    ↓
Contrato en L1 registra el batch
\end{verbatim}

\textbf{4. Challenge period:}
\begin{verbatim}
Durante ~7 días:
  - Cualquiera puede verificar el batch
  - Si encuentra error: inicia disputa
  - Si no hay disputa: batch se considera final
    ↓
Después del periodo:
  - State root se marca como finalizado
  - Retiros de ese batch pueden ejecutarse
\end{verbatim}
\end{definition}

\begin{center}
\begin{tikzpicture}[
    state/.style={draw, rectangle, rounded corners, minimum width=2cm, minimum height=0.8cm, align=center, font=\small},
    arrow/.style={->, thick}
]
    % Timeline
    \draw[thick, ->] (0, 0) -- (14, 0) node[right] {Tiempo};
    
    % Estados
    \node[state, fill=blue!20] (submit) at (1, 1.5) {Batch\\enviado};
    \node[state, fill=yellow!20] (challenge) at (5, 1.5) {Challenge\\period};
    \node[state, fill=orange!20] (dispute) at (5, -1.5) {Disputa\\(si hay)};
    \node[state, fill=green!20] (final) at (10, 1.5) {Finalizado};
    \node[state, fill=red!20] (reverted) at (10, -1.5) {Revertido\\(si fraude)};
    
    % Flechas
    \draw[arrow] (submit) -- (challenge);
    \draw[arrow] (challenge) -- node[above] {7 días, sin disputa} (final);
    \draw[arrow, dashed] (challenge) -- node[right] {disputa} (dispute);
    \draw[arrow, dashed] (dispute) -- node[below] {fraude probado} (reverted);
    \draw[arrow, dashed] (dispute) -- node[right, font=\tiny] {disputa\\fallida} (final);
    
    % Marcas de tiempo
    \draw (1, -0.1) -- (1, 0.1) node[below=0.2cm] {T};
    \draw (10, -0.1) -- (10, 0.1) node[below=0.2cm] {T+7d};
\end{tikzpicture}
\end{center}

%-------------------------------------------
\subsection{Fraud Proofs}

\begin{definition}[Fraud Proof]
Prueba de que una transacción fue mal ejecutada.
\end{definition}

Tipos:
\begin{itemize}
    \item \textbf{Non-interactive:} Prueba completa en una tx
    \item \textbf{Interactive:} Bisection game para encontrar paso inválido
\end{itemize}

\subsubsection{Non-Interactive Fraud Proofs}

\begin{definition}[Non-Interactive Fraud Proof]
Una sola transacción en L1 que demuestra el fraude completo.

\textbf{Contenido:}
\begin{itemize}
    \item Estado antes de la tx disputada
    \item La transacción en cuestión
    \item Merkle proofs de todos los datos accedidos
    \item Ejecución que muestra resultado diferente al claimed
\end{itemize}
\end{definition}

\begin{advertencia}
\textbf{Problema con non-interactive proofs:}

Una transacción puede acceder mucho estado:
\begin{itemize}
    \item Loops que leen muchos storage slots
    \item Llamadas a múltiples contratos
    \item Cada acceso requiere Merkle proof
\end{itemize}

En el peor caso, el fraud proof podría ser más grande que un bloque L1. Por esto, la mayoría de optimistic rollups usan proofs interactivos.
\end{advertencia}

\subsubsection{Interactive Fraud Proofs (Bisection Game)}

\begin{definition}[Bisection Game]
Protocolo interactivo para encontrar el primer paso de ejecución donde el operador diverge del resultado correcto.

\textbf{Participantes:}
\begin{itemize}
    \item Asserter: quien publicó el state root (defiende)
    \item Challenger: quien disputa (ataca)
\end{itemize}
\end{definition}

\begin{mathnote}
\textbf{Cómo funciona el bisection:}

Sea $N$ el número de pasos de ejecución en el batch disputado.

\textbf{Ronda 1:}
\begin{itemize}
    \item Asserter divide ejecución en 2 mitades
    \item Publica hash del estado en el punto medio ($N/2$)
    \item Challenger indica en qué mitad está el desacuerdo
\end{itemize}

\textbf{Ronda 2:}
\begin{itemize}
    \item Asserter divide la mitad disputada en 2
    \item Publica hash del estado en nuevo punto medio ($N/4$ o $3N/4$)
    \item Challenger indica en qué cuarto está el desacuerdo
\end{itemize}

\textbf{Continúa hasta:}
\begin{itemize}
    \item Llegar a un solo paso de ejecución
    \item Este único paso se ejecuta on-chain
    \item L1 determina quién tiene razón
\end{itemize}

\textbf{Complejidad:}
\begin{itemize}
    \item Número de rondas: $O(\log N)$
    \item Si $N = 2^{30}$ pasos: solo 30 rondas
    \item Cada ronda es una tx en L1
    \item Final: ejecutar 1 paso on-chain (factible)
\end{itemize}
\end{mathnote}

\begin{center}
\begin{tikzpicture}[
    step/.style={draw, rectangle, minimum width=0.3cm, minimum height=0.5cm},
    brace/.style={decorate, decoration={brace, amplitude=5pt}}
]
    % Pasos iniciales
    \foreach \x in {0, 0.4, 0.8, 1.2, 1.6, 2, 2.4, 2.8, 3.2, 3.6, 4, 4.4, 4.8, 5.2, 5.6, 6} {
        \node[step, fill=gray!30] at (\x, 0) {};
    }
    
    % Ronda 1: dividir en 2
    \draw[brace] (0, 0.6) -- node[above=0.2cm] {Primera mitad} (2.8, 0.6);
    \draw[brace] (3.2, 0.6) -- node[above=0.2cm] {Segunda mitad} (6, 0.6);
    \draw[red, thick] (3, 0.8) -- (3, -0.3);
    \node[below, font=\tiny] at (3, -0.5) {Punto medio};
    
    % Ronda 2: disputar primera mitad
    \draw[brace, blue] (0, -0.8) -- node[below=0.2cm, blue] {Cuarto 1} (1.2, -0.8);
    \draw[brace, blue] (1.6, -0.8) -- node[below=0.2cm, blue] {Cuarto 2} (2.8, -0.8);
    
    % Indicar disputa
    \node[red] at (0.6, -1.5) {$\leftarrow$ Disputa aquí};
    
    % Final
    \node[step, fill=red!50] at (0.4, -2.5) {};
    \node[right, font=\small] at (0.7, -2.5) {$\leftarrow$ Un solo paso ejecutado on-chain};
\end{tikzpicture}
\end{center}

\begin{example}[Bisection en Arbitrum]
\begin{verbatim}
Batch con 1 billón de pasos de ejecución (N = 10^9)

Rondas necesarias: log2(10^9) ≈ 30 rondas

Cada ronda:
  - Asserter: publica commitment (~100K gas)
  - Challenger: responde con elección (~50K gas)
  
Costo total (sin contar bond):
  ~30 × 150K gas = 4.5M gas ≈ $20-100 (depende de gas price)

Tiempo total:
  - Cada ronda permite ~1 día para responder
  - En la práctica: 1-2 semanas para resolver
  
Final:
  - Un solo paso de WASM/EVM se ejecuta en L1
  - Contrato determina quién tenía razón
  - Perdedor pierde su bond (~ETH)
\end{verbatim}
\end{example}

%-------------------------------------------
\subsection{Arquitectura de Arbitrum}

\begin{definition}[Arbitrum One]
El optimistic rollup más grande por TVL, desarrollado por Offchain Labs.
\end{definition}

\begin{keyidea}
\textbf{Componentes de Arbitrum:}

\textbf{1. Sequencer:}
\begin{itemize}
    \item Recibe txs de usuarios
    \item Ordena y ejecuta off-chain
    \item Publica batches en L1
    \item Actualmente centralizado (Offchain Labs)
\end{itemize}

\textbf{2. Validators/Asserters:}
\begin{itemize}
    \item Publican state roots (assertions)
    \item Deben depositar bond
    \item Cualquiera puede ser validator (permissionless)
\end{itemize}

\textbf{3. Challengers:}
\begin{itemize}
    \item Monitorean assertions
    \item Disputan si encuentran error
    \item También necesitan bond
\end{itemize}

\textbf{4. Contracts en L1:}
\begin{itemize}
    \item Rollup contract: gestiona assertions y challenges
    \item Bridge: maneja depósitos y retiros
    \item Inbox: recibe batches del sequencer
    \item Outbox: procesa retiros
\end{itemize}
\end{keyidea}

\begin{remark}
\textbf{Arbitrum vs Optimism (diferencias técnicas):}

\begin{center}
\begin{tabular}{lll}
\toprule
\textbf{Aspecto} & \textbf{Arbitrum} & \textbf{Optimism} \\
\midrule
VM & AVM (custom) / WASM & EVM equivalente \\
Fraud proof & Interactivo (bisection) & Single-round* \\
Challenge period & $\sim$7 días & $\sim$7 días \\
EVM compat & Nitro (muy alta) & Muy alta \\
Governance & ARB token & OP token \\
\bottomrule
\end{tabular}
\end{center}

*Optimism está migrando a fraud proofs basados en Cannon (MIPS).
\end{remark}

%-------------------------------------------
\subsection{El Problema de los 7 Días}

\begin{advertencia}
\textbf{El delay de retiro:}

Para retirar de un optimistic rollup a L1:
\begin{enumerate}
    \item Iniciar retiro en L2
    \item Esperar que batch sea publicado en L1
    \item Esperar challenge period completo ($\sim$7 días)
    \item Ejecutar retiro en L1
\end{enumerate}

\textbf{¿Por qué 7 días?}

\begin{itemize}
    \item Tiempo suficiente para que alguien detecte fraude
    \item Tiempo para ejecutar el bisection game completo
    \item Balance entre seguridad y UX
\end{itemize}

\textbf{Soluciones para usuarios impacientes:}

\begin{enumerate}
    \item \textbf{Liquidity providers:} Servicios como Hop, Across
    \begin{itemize}
        \item LP te da fondos en L1 inmediatamente
        \item LP espera los 7 días y cobra fee
        \item Fee típico: 0.1-0.5\%
    \end{itemize}
    
    \item \textbf{CEX:} Algunos exchanges aceptan depósitos directos desde L2
    
    \item \textbf{Cross-L2:} Bridges entre L2s sin pasar por L1
\end{enumerate}
\end{advertencia}

%-------------------------------------------
\subsection{Compresión de Datos}

\begin{mathnote}
\textbf{Cómo Arbitrum comprime transacciones:}

\textbf{Tx original Ethereum:}
\begin{verbatim}
{
  nonce: 4 bytes
  gasPrice: 8 bytes  
  gasLimit: 4 bytes
  to: 20 bytes
  value: 8 bytes
  data: variable
  v, r, s: 65 bytes
  ---
  Total: ~110+ bytes (sin data)
}
\end{verbatim}

\textbf{Tx comprimida en Arbitrum batch:}
\begin{verbatim}
{
  to: 4 bytes (índice, no dirección completa)
  value: 4 bytes (comprimido)
  nonce: 0-4 bytes (delta encoding)
  signature: 64 bytes (sin recovery byte)
  data: comprimido
  ---
  Total: ~12-20 bytes (sin data)
}
\end{verbatim}

\textbf{Técnicas de compresión:}
\begin{itemize}
    \item Direcciones indexadas: mapa de direcciones frecuentes
    \item Delta encoding: solo diferencia del nonce anterior
    \item Valores pequeños: encoding variable-length
    \item Firmas batch: BLS aggregation (futuro)
    \item Compresión general: zlib sobre todo el batch
\end{itemize}

\textbf{Resultado:}
\begin{itemize}
    \item Compresión 5-10x típica
    \item Una transferencia: $\sim$12 bytes
    \item Costo L1 por tx: $\sim$\$0.01-0.10
\end{itemize}
\end{mathnote}

%-------------------------------------------
\subsection{Seguridad y Trust Assumptions}

\begin{keyidea}
\textbf{¿En qué confías al usar un optimistic rollup?}

\textbf{Asunciones de seguridad:}
\begin{enumerate}
    \item \textbf{Al menos 1 validator honesto:}
    \begin{itemize}
        \item Alguien debe disputar si hay fraude
        \item Si todos los validators son maliciosos, fraude pasa
        \item En la práctica: múltiples validators independientes
    \end{itemize}
    
    \item \textbf{L1 está disponible:}
    \begin{itemize}
        \item Disputas ocurren en L1
        \item Si L1 está censurado, no se puede disputar
    \end{itemize}
    
    \item \textbf{Datos disponibles:}
    \begin{itemize}
        \item Batch data publicado en L1 (calldata/blobs)
        \item Necesario para verificar y disputar
    \end{itemize}
\end{enumerate}

\textbf{Riesgos actuales (training wheels):}

La mayoría de optimistic rollups tienen:
\begin{itemize}
    \item Upgrade keys: multisig puede actualizar contratos
    \item Pausability: pueden pausar el rollup
    \item Sequencer centralizado: puede censurar
    \item Whitelist de validators: no cualquiera puede disputar
\end{itemize}

\textbf{Roadmap a trustlessness:}
\begin{itemize}
    \item Remover upgrade keys (time delays primero)
    \item Validators permissionless
    \item Sequencer descentralizado
    \item Emergency exit mechanisms
\end{itemize}
\end{keyidea}

%-------------------------------------------
\subsection{Ejemplos: Arbitrum, Optimism, Base}

\begin{remark}
\textbf{Estado del ecosistema Optimistic Rollup (2024):}

\begin{center}
\begin{tabular}{lcccr}
\toprule
\textbf{Rollup} & \textbf{TVL} & \textbf{TPS} & \textbf{Fraud Proofs} & \textbf{Sequencer} \\
\midrule
Arbitrum One & $\sim$\$10B & $\sim$40 & Live (permissioned) & Centralizado \\
Optimism & $\sim$\$5B & $\sim$30 & En desarrollo & Centralizado \\
Base & $\sim$\$3B & $\sim$25 & En desarrollo & Centralizado \\
Mantle & $\sim$\$1B & $\sim$20 & En desarrollo & Centralizado \\
\bottomrule
\end{tabular}
\end{center}

\textbf{OP Stack (Superchain):}

Optimism creó el OP Stack, un framework para lanzar optimistic rollups:
\begin{itemize}
    \item Base (Coinbase)
    \item Zora
    \item Mode
    \item Worldcoin
    \item Muchos más
\end{itemize}

La visión ``Superchain'' es que todos estos L2s compartan:
\begin{itemize}
    \item Sequencing (eventualmente)
    \item Bridge nativo entre ellos
    \item Standards comunes
\end{itemize}
\end{remark}

%-------------------------------------------
\subsection{Comparación con ZK Rollups}

\begin{center}
\begin{tabular}{lll}
\toprule
\textbf{Aspecto} & \textbf{Optimistic} & \textbf{ZK} \\
\midrule
Validez & Asumida, disputable & Probada matemáticamente \\
Finalidad L1 & $\sim$7 días & Minutos-horas \\
Costo L2 tx & Muy bajo & Bajo (más overhead) \\
Costo L1 (proof) & 0 si no hay disputa & Siempre (verify) \\
EVM compat & Nativa & Requiere zkEVM \\
Madurez & Alta & Media \\
Complejidad & Media & Alta \\
\bottomrule
\end{tabular}
\end{center}

\begin{keyidea}
\textbf{¿Cuándo elegir Optimistic vs ZK?}

\textbf{Optimistic es mejor para:}
\begin{itemize}
    \item Máxima compatibilidad EVM hoy
    \item Costos más bajos (especialmente con poco volumen)
    \item Aplicaciones donde 7 días de delay es aceptable
    \item Casos donde ya hay liquidez en Optimistic L2s
\end{itemize}

\textbf{ZK es mejor para:}
\begin{itemize}
    \item Finalidad rápida es crítica
    \item Privacidad (como Banco ZK)
    \item Bridges seguros (verificación instantánea)
    \item Aplicaciones financieras de alto valor
\end{itemize}

\textbf{Tendencia:} A largo plazo, muchos creen que ZK dominará porque:
\begin{itemize}
    \item Verificación instantánea > esperar 7 días
    \item Hardware ZK mejora rápidamente
    \item zkEVMs están madurando
\end{itemize}
\end{keyidea}

% ============================================
% BLOQUE 3: ESCALABILIDAD (L2/L3) - VERSIÓN EXPANDIDA
% Parte 3: ZK Rollups
% ============================================

%===========================================
\section{ZK Rollups}
%===========================================

\begin{definition}[ZK Rollup]
Genera prueba criptográfica de ejecución correcta.
\end{definition}

\begin{intuicion}
Un ZK Rollup es como un examen con respuestas verificadas:
\begin{itemize}
    \item El estudiante (prover) resuelve muchos problemas (ejecuta txs)
    \item En lugar de mostrar todo su trabajo, genera un ``certificado de corrección''
    \item El profesor (L1) verifica el certificado en segundos
    \item Si el certificado es válido, las respuestas son correctas garantizado
    \item No hay necesidad de esperar quejas de otros estudiantes
\end{itemize}

La magia es que verificar el certificado (proof) es muchísimo más barato que re-ejecutar todos los cálculos.
\end{intuicion}

%-------------------------------------------
\subsection{Cómo Funcionan}

\begin{enumerate}
    \item Sequencer ejecuta batch off-chain
    \item Prover genera ZK proof de validez
    \item Publica datos + proof en L1
    \item L1 verifica proof (barato) y acepta estado
    \item Finalidad inmediata (sin período de disputa)
\end{enumerate}

\subsubsection{Flujo Detallado}

\begin{definition}[Flujo de ZK Rollup]
\textbf{1. Recepción y ejecución:}
\begin{verbatim}
Usuario envía tx al sequencer
    ↓
Sequencer valida y ejecuta tx
    ↓
Actualiza estado interno (off-chain)
    ↓
Acumula txs en batch
\end{verbatim}

\textbf{2. Generación de prueba:}
\begin{verbatim}
Prover recibe batch de txs
    ↓
Construye circuito ZK que representa:
  - Verificación de firmas
  - Transiciones de estado
  - Actualización de Merkle trees
  - Cualquier lógica del protocolo
    ↓
Ejecuta el circuito con witnesses (datos privados)
    ↓
Genera proof (puede tomar minutos a horas)
\end{verbatim}

\textbf{3. Publicación en L1:}
\begin{verbatim}
Sequencer/Prover envía a L1:
  - Compressed batch data (para DA)
  - Previous state root
  - New state root
  - ZK proof
    ↓
Contrato en L1:
  - Verifica proof (precompilado)
  - Si válido: actualiza state root
  - Emite evento
\end{verbatim}

\textbf{4. Finalización:}
\begin{verbatim}
Inmediatamente después de verificación:
  - Nuevo estado es final
  - Retiros pueden procesarse
  - No hay challenge period
\end{verbatim}
\end{definition}

\begin{center}
\begin{tikzpicture}[
    box/.style={draw, rectangle, rounded corners, minimum width=2.5cm, minimum height=1cm, align=center},
    arrow/.style={->, thick}
]
    % Off-chain
    \node[box, fill=blue!20] (seq) at (0, 2) {Sequencer\\ejecuta txs};
    \node[box, fill=yellow!20] (prover) at (4, 2) {Prover\\genera proof};
    
    % On-chain
    \node[box, fill=green!20] (verify) at (8, 2) {L1 Verifier\\verifica proof};
    \node[box, fill=green!40] (final) at (12, 2) {Estado\\finalizado};
    
    % Flechas
    \draw[arrow] (seq) -- node[above, font=\small] {batch} (prover);
    \draw[arrow] (prover) -- node[above, font=\small] {proof} (verify);
    \draw[arrow] (verify) -- node[above, font=\small] {válido} (final);
    
    % Tiempo
    \draw[dashed] (6, 0) -- (6, 3.5);
    \node[below] at (3, 0) {Off-chain (minutos-horas)};
    \node[below] at (10, 0) {On-chain (segundos)};
\end{tikzpicture}
\end{center}

%-------------------------------------------
\subsection{Validity Proofs}

La prueba ZK garantiza:
\begin{itemize}
    \item Todas las firmas son válidas
    \item Transiciones de estado son correctas
    \item Balances no son negativos
    \item Reglas del protocolo se cumplen
\end{itemize}

\begin{mathnote}
\textbf{¿Qué exactamente prueba el ZK proof?}

El circuito ZK codifica la función de transición de estado:
\[ \text{proof}: \pi \text{ tal que } \text{Verify}(\pi, S_{\text{old}}, S_{\text{new}}, \text{batch}) = 1 \]

\textbf{La prueba garantiza:}
\begin{enumerate}
    \item Para cada tx en el batch:
    \begin{itemize}
        \item La firma ECDSA/EdDSA es válida
        \item El nonce es correcto
        \item El sender tiene balance suficiente
        \item La ejecución sigue las reglas
    \end{itemize}
    
    \item El state root $S_{\text{new}}$ es el resultado correcto de aplicar todas las txs a $S_{\text{old}}$
    
    \item Los datos publicados coinciden con las txs ejecutadas
\end{enumerate}

\textbf{Propiedades:}
\begin{itemize}
    \item \textbf{Soundness:} Es computacionalmente imposible crear un proof falso
    \item \textbf{Zero-knowledge (opcional):} El proof no revela información adicional
    \item \textbf{Succinctness:} El proof es pequeño y rápido de verificar
\end{itemize}
\end{mathnote}

%-------------------------------------------
\subsection{Tipos de ZK Rollups}

\begin{definition}[Tipos según sistema de pruebas]
\textbf{1. SNARK-based:}
\begin{itemize}
    \item zkSync Era, Scroll, Polygon zkEVM, Linea
    \item Pruebas pequeñas ($\sim$200-500 bytes)
    \item Verificación barata (200-500K gas)
    \item Típicamente requieren trusted setup
\end{itemize}

\textbf{2. STARK-based:}
\begin{itemize}
    \item StarkNet, StarkEx
    \item Pruebas grandes ($\sim$50-200 KB)
    \item No trusted setup (transparent)
    \item Post-quantum seguro
\end{itemize}
\end{definition}

\begin{definition}[Tipos según compatibilidad EVM]
\textbf{Tipo 1 - Ethereum equivalent:}
\begin{itemize}
    \item Idéntico a Ethereum, incluyendo bugs
    \item Máxima compatibilidad
    \item Más difícil de probar en ZK
    \item Ejemplo: Taiko (en desarrollo)
\end{itemize}

\textbf{Tipo 2 - EVM equivalent:}
\begin{itemize}
    \item Compatible a nivel EVM
    \item Pequeñas diferencias en edge cases
    \item Ejemplo: Scroll, Polygon zkEVM
\end{itemize}

\textbf{Tipo 3 - EVM compatible:}
\begin{itemize}
    \item La mayoría de código funciona
    \item Algunas modificaciones necesarias
    \item Ejemplo: Scroll (en transición a tipo 2)
\end{itemize}

\textbf{Tipo 4 - Language compatible:}
\begin{itemize}
    \item Lenguaje propio que compila similar a EVM
    \item Solidity/Vyper compilan a VM diferente
    \item Ejemplo: zkSync Era, StarkNet (Cairo)
\end{itemize}
\end{definition}

\begin{center}
\begin{tikzpicture}
    % Espectro
    \draw[thick, ->] (0, 0) -- (12, 0);
    \node[below] at (0, -0.3) {Más compatible};
    \node[below] at (12, -0.3) {Más eficiente};
    
    % Tipos
    \node[draw, fill=blue!20, minimum width=2cm] at (1.5, 1) {Tipo 1};
    \node[draw, fill=blue!30, minimum width=2cm] at (4.5, 1) {Tipo 2};
    \node[draw, fill=blue!40, minimum width=2cm] at (7.5, 1) {Tipo 3};
    \node[draw, fill=blue!50, minimum width=2cm] at (10.5, 1) {Tipo 4};
    
    % Ejemplos
    \node[font=\tiny] at (1.5, 0.4) {Taiko};
    \node[font=\tiny] at (4.5, 0.4) {Scroll};
    \node[font=\tiny] at (7.5, 0.4) {Polygon zkEVM};
    \node[font=\tiny] at (10.5, 0.4) {zkSync, StarkNet};
\end{tikzpicture}
\end{center}

%-------------------------------------------
\subsection{zkEVM: Probando la EVM}

\begin{keyidea}
\textbf{El desafío del zkEVM:}

La EVM no fue diseñada para ser probada en ZK:
\begin{itemize}
    \item Usa Keccak256 (no ZK-friendly)
    \item Aritmética de 256 bits (muy grande para circuitos)
    \item Stack machine (muchas operaciones por paso)
    \item Storage como Merkle Patricia Trie (complejo)
\end{itemize}

\textbf{Estrategias de los zkEVMs:}

\textbf{1. Probar EVM directamente (tipo 1-2):}
\begin{itemize}
    \item Circuito que simula cada opcode
    \item Muy grande pero compatible
    \item Scroll, Polygon zkEVM
\end{itemize}

\textbf{2. Compilar a VM ZK-friendly (tipo 4):}
\begin{itemize}
    \item Transformar Solidity a otra representación
    \item Más eficiente pero menos compatible
    \item zkSync Era, StarkNet
\end{itemize}
\end{keyidea}

\begin{mathnote}
\textbf{Costos de opcodes en zkEVM:}

\begin{center}
\begin{tabular}{lccc}
\toprule
\textbf{Operación} & \textbf{Gas EVM} & \textbf{Constraints ZK} & \textbf{Ratio} \\
\midrule
ADD & 3 & $\sim$100 & 33x \\
MUL & 5 & $\sim$500 & 100x \\
SHA3/Keccak & 30+6/word & $\sim$150,000 & 5000x \\
ECDSA verify & 3000 & $\sim$500,000 & 167x \\
SLOAD & 2100 & $\sim$10,000 & 5x \\
Merkle update & N/A & $\sim$50,000 & N/A \\
\bottomrule
\end{tabular}
\end{center}

\textbf{Nota:} Keccak es especialmente caro en ZK. Por eso algunos zkEVMs usan Poseidon para hashes internos.
\end{mathnote}

%-------------------------------------------
\subsection{Ejemplos: zkSync, StarkNet, Scroll, Polygon zkEVM}

\begin{remark}
\textbf{Comparación de ZK Rollups principales (2024):}

\begin{center}
\begin{tabular}{lcccc}
\toprule
\textbf{} & \textbf{zkSync Era} & \textbf{StarkNet} & \textbf{Scroll} & \textbf{Polygon zkEVM} \\
\midrule
Tipo zkEVM & 4 & 4 (Cairo) & 2 & 2-3 \\
Proof system & PLONK & STARK & PLONK & PLONK \\
DA & L1 (blobs) & L1 + validium & L1 & L1 \\
TVL & $\sim$\$800M & $\sim$\$300M & $\sim$\$200M & $\sim$\$100M \\
Mainnet & Mar 2023 & Nov 2022 & Oct 2023 & Mar 2023 \\
\bottomrule
\end{tabular}
\end{center}

\textbf{zkSync Era:}
\begin{itemize}
    \item Desarrollado por Matter Labs
    \item Usa su propia VM (zkEVM, tipo 4)
    \item Soporta account abstraction nativo
    \item Compilador llvm-based para Solidity
\end{itemize}

\textbf{StarkNet:}
\begin{itemize}
    \item Desarrollado por StarkWare
    \item Usa Cairo (lenguaje propio, no EVM)
    \item STARKs (no trusted setup, post-quantum)
    \item Pruebas más grandes pero más seguras a largo plazo
\end{itemize}

\textbf{Scroll:}
\begin{itemize}
    \item Colaboración con Ethereum Foundation
    \item zkEVM tipo 2 (muy compatible)
    \item Usa Halo2 (sin trusted setup)
    \item Enfocado en equivalencia bytecode
\end{itemize}

\textbf{Polygon zkEVM:}
\begin{itemize}
    \item Parte del ecosistema Polygon
    \item zkEVM tipo 2-3
    \item Proof recursion para eficiencia
    \item Integración con Polygon ecosystem
\end{itemize}
\end{remark}

%-------------------------------------------
\subsection{Prover: El Cuello de Botella}

\begin{advertencia}
\textbf{El costo del proving:}

Generar una prueba ZK es computacionalmente intensivo:

\begin{center}
\begin{tabular}{lcc}
\toprule
\textbf{Batch size} & \textbf{Tiempo} & \textbf{Hardware} \\
\midrule
100 txs simples & $\sim$1 min & GPU potente \\
1000 txs & $\sim$10 min & Cluster GPUs \\
10000 txs & $\sim$1 hora & Datacenter \\
\bottomrule
\end{tabular}
\end{center}

\textbf{Por qué es caro:}
\begin{itemize}
    \item Operaciones de campo grande (FFT, MSM)
    \item Circuitos con millones de constraints
    \item Mucha memoria (100+ GB para batches grandes)
\end{itemize}

\textbf{Soluciones:}
\begin{itemize}
    \item Hardware especializado (FPGAs, ASICs)
    \item Paralelización masiva (GPUs)
    \item Proof recursion (probar proofs)
    \item Optimización de circuitos
    \item Mercados de proving descentralizados
\end{itemize}
\end{advertencia}

\begin{mathnote}
\textbf{Proof Recursion:}

En lugar de una prueba gigante para todo el batch:

\begin{verbatim}
Batch de 10,000 txs:
    ↓
Dividir en 100 mini-batches de 100 txs
    ↓
Generar 100 pruebas pequeñas (paralelo)
    ↓
Generar prueba que verifica las 100 pruebas
    ↓
Una sola prueba final para L1
\end{verbatim}

\textbf{Beneficios:}
\begin{itemize}
    \item Paralelización: 100 provers trabajan simultáneamente
    \item Latencia reducida: no esperar por mega-batch
    \item Flexibilidad: agregar más provers linealmente
\end{itemize}

\textbf{Overhead:}
\begin{itemize}
    \item Prueba recursiva tiene overhead
    \item Pero la paralelización lo compensa
\end{itemize}
\end{mathnote}

%-------------------------------------------
\subsection{Data Availability en ZK Rollups}

\begin{keyidea}
\textbf{¿Por qué publicar datos si ya tenemos proof?}

El proof garantiza que la transición de estado es correcta, pero:
\begin{itemize}
    \item Usuarios necesitan conocer su estado para interactuar
    \item Sin datos, no pueden construir txs futuras
    \item Sin datos, no pueden hacer exit si operador desaparece
\end{itemize}

\textbf{Opciones de DA para ZK Rollups:}

\textbf{1. ZK Rollup (datos on-chain):}
\begin{itemize}
    \item Publica batch data en L1 (calldata/blobs)
    \item Máxima seguridad
    \item Más caro
\end{itemize}

\textbf{2. Validium (datos off-chain):}
\begin{itemize}
    \item Solo proof en L1
    \item Datos en comité (DAC) o DA layer
    \item Más barato pero menos seguro
\end{itemize}

\textbf{3. Volition (híbrido):}
\begin{itemize}
    \item Usuario elige por tx
    \item Txs importantes: on-chain
    \item Txs menos críticas: off-chain
\end{itemize}
\end{keyidea}

%-------------------------------------------
\subsection{Ventajas y Desventajas de ZK Rollups}

\begin{center}
\begin{tabular}{ll}
\toprule
\textbf{Ventajas} & \textbf{Desventajas} \\
\midrule
Finalidad rápida (minutos) & Proving es caro/lento \\
Sin challenge period & Compatibilidad EVM limitada* \\
Seguridad matemática & Complejidad técnica alta \\
Bridges más seguros & Circuitos pueden tener bugs \\
Privacidad posible & Menos maduros que Optimistic \\
\bottomrule
\end{tabular}
\end{center}

*Mejorando rápidamente con zkEVMs.

%-------------------------------------------
\subsection{ZK Rollups con Privacidad}

\begin{bancozk}
\textbf{Banco ZK: ZK Rollup + Privacidad}

Banco ZK combina las ventajas de ZK Rollups con privacidad transaccional:

\textbf{Como ZK Rollup estándar:}
\begin{itemize}
    \item Publica datos en L1 (para DA)
    \item Pruebas de validez verificadas en L1
    \item Hereda seguridad de Ethereum
    \item Finalidad rápida
\end{itemize}

\textbf{Capa de privacidad adicional:}
\begin{itemize}
    \item Txs son commitments encriptados
    \item Montos ocultos con Pedersen
    \item Identidades ocultas con nullifiers
    \item Prueba ZK incluye prueba de privacidad
\end{itemize}

\textbf{Arquitectura:}
\begin{verbatim}
Usuario genera tx privada:
  - Note commitment (oculta monto y owner)
  - Nullifier (prueba gasto sin revelar cuál note)
  - ZK proof de validez
    ↓
Sequencer incluye en batch:
  - No puede ver contenido
  - Solo verifica que proof es válido
    ↓
Prover genera proof de batch:
  - Incluye verificación de cada tx
  - Transiciones de estado privadas
    ↓
L1 verifica:
  - Proof de todo el batch
  - Nullifiers no duplicados
  - State root correcto
\end{verbatim}

\textbf{Diferencia con zkSync/StarkNet/Scroll:}
\begin{itemize}
    \item Ellos: txs públicas, solo escalabilidad
    \item Banco ZK: txs privadas, escalabilidad + privacidad
    \item El costo de privacidad: circuitos más grandes
\end{itemize}

\textbf{Similar a:}
\begin{itemize}
    \item Aztec Protocol (ZK rollup privado)
    \item Aleo (L1 ZK nativo)
    \item Zcash (pero como rollup sobre Ethereum)
\end{itemize}
\end{bancozk}

% ============================================
% BLOQUE 3: ESCALABILIDAD (L2/L3) - VERSIÓN EXPANDIDA
% Parte 4: Sequencers, Data Availability, Blobs
% ============================================

%===========================================
\section{Sequencers}
%===========================================

\begin{definition}[Sequencer]
Entidad que ordena y ejecuta transacciones en un rollup.
\end{definition}

Estado actual:
\begin{itemize}
    \item Mayoría de rollups tienen sequencer centralizado
    \item Más eficiente pero punto único de falla
    \item Riesgos: censura, MEV extraction, downtime
\end{itemize}

%-------------------------------------------
\subsection{Rol del Sequencer}

\begin{intuicion}
El sequencer es como el director de orquesta del rollup:
\begin{itemize}
    \item Recibe todas las ``notas'' (transacciones) de los músicos (usuarios)
    \item Decide el orden en que se tocan
    \item Coordina la ejecución
    \item Presenta el resultado final
\end{itemize}

El problema es que un director corrupto podría:
\begin{itemize}
    \item Ignorar a ciertos músicos (censura)
    \item Reordenar para su beneficio (MEV)
    \item Simplemente no presentarse (downtime)
\end{itemize}
\end{intuicion}

\begin{definition}[Funciones del Sequencer]
\textbf{1. Recepción:}
\begin{itemize}
    \item Acepta txs de usuarios vía RPC
    \item Valida formato y firmas
    \item Mantiene mempool local
\end{itemize}

\textbf{2. Ordenamiento:}
\begin{itemize}
    \item Decide el orden de las txs
    \item Puede ser FIFO, por fee, o arbitrario
    \item Este orden es crítico para MEV
\end{itemize}

\textbf{3. Ejecución:}
\begin{itemize}
    \item Ejecuta txs contra el estado
    \item Calcula nuevos state roots
    \item Proporciona ``soft confirmations'' a usuarios
\end{itemize}

\textbf{4. Batching:}
\begin{itemize}
    \item Agrupa txs en batches
    \item Comprime datos
    \item Publica en L1
\end{itemize}
\end{definition}

%-------------------------------------------
\subsection{Centralización vs Descentralización}

\begin{advertencia}
\textbf{Riesgos de sequencer centralizado:}

\textbf{1. Censura:}
\begin{itemize}
    \item Sequencer puede ignorar ciertas txs
    \item Usuarios de ciertos países/wallets bloqueados
    \item Mitigación: forced inclusion via L1
\end{itemize}

\textbf{2. MEV extraction:}
\begin{itemize}
    \item Sequencer ve todas las txs pendientes
    \item Puede hacer frontrunning/sandwiching
    \item Captura valor que debería ir a usuarios
\end{itemize}

\textbf{3. Downtime:}
\begin{itemize}
    \item Si sequencer cae, L2 se detiene
    \item Usuarios no pueden hacer txs
    \item Dependen de que vuelva a funcionar
\end{itemize}

\textbf{4. Jurisdicción:}
\begin{itemize}
    \item Sequencer es una empresa en algún país
    \item Puede ser obligado a censurar
    \item Single point of regulatory attack
\end{itemize}

\textbf{Estado actual (2024):}

\begin{center}
\begin{tabular}{lcc}
\toprule
\textbf{Rollup} & \textbf{Sequencer} & \textbf{Forced inclusion} \\
\midrule
Arbitrum & Offchain Labs & Sí (delayed) \\
Optimism & Optimism Foundation & Sí (delayed) \\
zkSync & Matter Labs & En desarrollo \\
StarkNet & StarkWare & Sí \\
Base & Coinbase & Sí (delayed) \\
\bottomrule
\end{tabular}
\end{center}
\end{advertencia}

%-------------------------------------------
\subsection{Soluciones para Descentralización}

\begin{keyidea}
\textbf{Enfoques para descentralizar sequencers:}

\textbf{1. Rotación de sequencers:}
\begin{itemize}
    \item Múltiples sequencers que rotan
    \item Selección por stake, lottery, o round-robin
    \item Si uno censura, el siguiente incluye
    \item Ejemplo: Diseño futuro de Arbitrum
\end{itemize}

\textbf{2. Based rollups (L1 sequencing):}
\begin{itemize}
    \item L1 validators ordenan txs del rollup
    \item Hereda descentralización de L1
    \item No hay sequencer L2 separado
    \item Trade-off: latencia de L1 (12s)
    \item Ejemplo: Taiko
\end{itemize}

\textbf{3. Shared sequencing:}
\begin{itemize}
    \item Múltiples rollups comparten un sequencer descentralizado
    \item Beneficios de atomic cross-rollup txs
    \item Economías de escala
    \item Ejemplos: Espresso, Astria
\end{itemize}

\textbf{4. Encrypted mempools:}
\begin{itemize}
    \item Txs encriptadas hasta inclusión
    \item Sequencer ordena sin ver contenido
    \item Elimina MEV del sequencer
    \item Ejemplo: Shutter Network, SUAVE
\end{itemize}
\end{keyidea}

\begin{mathnote}
\textbf{Based Rollups - Cómo funcionan:}

\begin{verbatim}
Flujo tradicional:
  User → Sequencer L2 → L1 (datos)
                ↑
          Ordering aquí (centralizado)

Flujo based rollup:
  User → L1 Mempool → L1 Block → L2 deriva estado
                ↑
          Ordering por L1 validators (descentralizado)
\end{verbatim}

\textbf{Ventajas:}
\begin{itemize}
    \item Descentralización heredada de L1
    \item Sin operador L2 que pueda censurar
    \item Atomicidad con L1 (composability)
\end{itemize}

\textbf{Desventajas:}
\begin{itemize}
    \item Latencia mínima = block time L1 (12s)
    \item No hay ``soft confirmations'' instantáneas
    \item Menos optimización posible
\end{itemize}
\end{mathnote}

\begin{remark}
\textbf{Shared Sequencing - Espresso Systems:}

Espresso propone un sequencer compartido entre múltiples rollups:

\begin{center}
\begin{tikzpicture}[
    rollup/.style={draw, rectangle, minimum width=1.5cm, minimum height=0.8cm},
    seq/.style={draw, ellipse, minimum width=3cm, minimum height=1.5cm, fill=yellow!30}
]
    % Sequencer compartido
    \node[seq] (esp) at (5, 2) {Espresso\\(Shared Sequencer)};
    
    % Rollups
    \node[rollup, fill=blue!20] (r1) at (1, 0) {Rollup A};
    \node[rollup, fill=green!20] (r2) at (5, 0) {Rollup B};
    \node[rollup, fill=red!20] (r3) at (9, 0) {Rollup C};
    
    % Flechas
    \draw[->] (r1) -- (esp);
    \draw[->] (r2) -- (esp);
    \draw[->] (r3) -- (esp);
    
    % L1
    \node[draw, rectangle, fill=gray!20, minimum width=8cm] (l1) at (5, -1.5) {Ethereum L1};
    
    \draw[->] (esp) -- (l1);
    \draw[->] (r1) -- (l1);
    \draw[->] (r2) -- (l1);
    \draw[->] (r3) -- (l1);
\end{tikzpicture}
\end{center}

\textbf{Beneficios:}
\begin{itemize}
    \item Atomic bundles cross-rollup
    \item Compartir infraestructura
    \item Interoperabilidad mejorada
\end{itemize}
\end{remark}

%===========================================
\section{Data Availability (DA)}
%===========================================

\begin{definition}[Data Availability]
Garantía de que datos para reconstruir estado están disponibles.
\end{definition}

%-------------------------------------------
\subsection{Por Qué Importa}

Sin DA, operador podría:
\begin{itemize}
    \item Publicar estado válido pero ocultar txs
    \item Usuarios no pueden calcular sus balances
    \item No pueden generar proofs para retiros
\end{itemize}

\begin{intuicion}
Imagina que el banco te dice ``tu balance es \$1000'' pero no te muestra los movimientos:
\begin{itemize}
    \item No puedes verificar si es correcto
    \item No sabes de dónde vino el dinero
    \item Si el banco desaparece, no puedes probar nada
\end{itemize}

Data Availability es tener acceso a los ``movimientos'' para poder verificar y actuar independientemente.
\end{intuicion}

\begin{mathnote}
\textbf{El problema de Data Withholding:}

Escenario malicioso:
\begin{enumerate}
    \item Operador ejecuta batch correctamente
    \item Publica state root válido en L1
    \item PERO no publica los datos del batch
    \item Estado en L1 es ``correcto'' pero nadie sabe qué contiene
\end{enumerate}

\textbf{Consecuencias:}
\begin{itemize}
    \item Usuarios no saben sus balances actuales
    \item No pueden construir txs (necesitan Merkle proofs)
    \item No pueden hacer exit (necesitan probar ownership)
    \item Operador puede robar fondos en el siguiente batch
\end{itemize}

\textbf{Por qué no solo pedir los datos?}
\begin{itemize}
    \item Operador malicioso simplemente no responde
    \item ¿Cómo probar que NO te dieron los datos?
    \item Problema fundamentalmente asimétrico
\end{itemize}
\end{mathnote}

%-------------------------------------------
\subsection{Opciones de DA}

\begin{table}[htbp]
\centering
\begin{tabular}{lll}
\toprule
\textbf{Opción} & \textbf{Seguridad} & \textbf{Costo} \\
\midrule
Calldata L1 & Máxima & Alto \\
Blobs (EIP-4844) & Alta & Medio \\
DAC (comité) & Media & Bajo \\
DA layers & Media-Alta & Bajo \\
\bottomrule
\end{tabular}
\end{table}

DA layers: Celestia, EigenDA, Avail, NearDA.

\subsubsection{Calldata en L1}

\begin{definition}[Calldata]
Datos incluidos en transacciones L1 que son permanentemente almacenados por todos los nodos.
\end{definition}

\begin{itemize}
    \item \textbf{Seguridad:} Máxima (misma que L1)
    \item \textbf{Costo:} 16 gas por byte non-zero, 4 por zero
    \item \textbf{Disponibilidad:} Para siempre
    \item \textbf{Acceso:} No accesible desde EVM, pero almacenado
\end{itemize}

\begin{mathnote}
\textbf{Cálculo de costo de calldata:}

Para un batch de 1000 txs comprimidas ($\sim$15 KB):
\begin{verbatim}
15,000 bytes × ~12 gas/byte (promedio) = 180,000 gas
A 30 gwei: 180,000 × 30 = 5.4M gwei = 0.0054 ETH
A $2500/ETH: ~$13.50

Por tx: $13.50 / 1000 = $0.0135
\end{verbatim}

Este era el costo antes de EIP-4844. Aún usado como fallback.
\end{mathnote}

\subsubsection{Data Availability Committees (DAC)}

\begin{definition}[DAC]
Grupo de entidades que prometen mantener datos disponibles.
\end{definition}

\begin{itemize}
    \item Rollup envía datos a comité (no a L1)
    \item Comité firma attestation de que tienen los datos
    \item Solo la attestation va a L1 (muy pequeña)
    \item Si usuario pide datos, comité debe proporcionar
\end{itemize}

\begin{advertencia}
\textbf{Riesgos de DAC:}

\begin{itemize}
    \item Si comité colude: pueden withhold datos
    \item Trust assumption: mayoría honesta del comité
    \item Típicamente N-of-M: si N miembros honestos, datos disponibles
\end{itemize}

\textbf{Ejemplos:}
\begin{itemize}
    \item StarkEx (validium mode): DAC de entidades conocidas
    \item Immutable X: DAC para gaming
\end{itemize}
\end{advertencia}

\subsubsection{DA Layers Externos}

\begin{definition}[DA Layer]
Blockchain especializada en almacenar datos con garantías criptoeconómicas.
\end{definition}

\begin{remark}
\textbf{Principales DA Layers:}

\textbf{Celestia:}
\begin{itemize}
    \item Primera DA layer dedicada (mainnet Oct 2023)
    \item Usa Data Availability Sampling (DAS)
    \item Light nodes pueden verificar DA sin descargar todo
    \item Costo: \$$\sim$0.001 por KB
\end{itemize}

\textbf{EigenDA:}
\begin{itemize}
    \item DA sobre EigenLayer (restaking)
    \item Seguridad heredada de ETH restaked
    \item Integración nativa con Ethereum validators
\end{itemize}

\textbf{Avail:}
\begin{itemize}
    \item Fork de Polygon, enfocado en DA
    \item KZG commitments
    \item Validity proofs para DA
\end{itemize}

\textbf{NearDA:}
\begin{itemize}
    \item DA usando Near Protocol
    \item Sharding nativo de Near
    \item Bajo costo
\end{itemize}
\end{remark}

\begin{mathnote}
\textbf{Data Availability Sampling (DAS):}

DAS permite verificar que datos están disponibles sin descargarlos completos:

\begin{enumerate}
    \item Datos se codifican con Reed-Solomon (redundancia)
    \item Se divide en muchos chunks pequeños
    \item Light node solicita chunks aleatorios
    \item Si recibe suficientes chunks, alta confianza de que todo está disponible
\end{enumerate}

\textbf{Matemática:}
\begin{itemize}
    \item Con 50\% de datos disponibles, se puede reconstruir 100\%
    \item Si un atacante oculta $>$50\%, será detectado
    \item Probabilidad de no detectar con $k$ samples: $(0.5)^k$
    \item Con $k = 30$ samples: probabilidad de error $< 10^{-9}$
\end{itemize}

Esto permite que nodes con poco bandwidth verifiquen DA.
\end{mathnote}

%===========================================
\section{Blobs (EIP-4844)}
%===========================================

\begin{definition}[Blob]
Nuevo tipo de dato en Ethereum optimizado para rollups.
\end{definition}

Características:
\begin{itemize}
    \item 128 KB por blob, hasta 6 por bloque
    \item No accesible desde EVM (solo por hash)
    \item Se elimina después de $\sim$18 días
    \item Mercado de fees separado
    \item Compromiso KZG para verificación
\end{itemize}

%-------------------------------------------
\subsection{Contexto y Motivación}

\begin{remark}
\textbf{El camino a EIP-4844:}

\textbf{Problema:}
\begin{itemize}
    \item Rollups publicaban datos como calldata
    \item Calldata es permanente y cara
    \item Rollups competían con txs normales por blockspace
    \item Fees de L2 dependían de congestión L1
\end{itemize}

\textbf{Solución - Danksharding (visión completa):}
\begin{itemize}
    \item Nuevos ``data shards'' dedicados para rollups
    \item Separar datos de ejecución
    \item Mucho más espacio para datos
\end{itemize}

\textbf{Proto-Danksharding (EIP-4844):}
\begin{itemize}
    \item Paso intermedio hacia full Danksharding
    \item Introduce blobs sin sharding completo
    \item Mercado de fees separado
    \item Implementado en Dencun (Marzo 2024)
\end{itemize}
\end{remark}

%-------------------------------------------
\subsection{Anatomía de un Blob}

\begin{definition}[Estructura de Blob]
\textbf{Tamaño:} 128 KB (131,072 bytes) = 4096 field elements de 32 bytes

\textbf{Representación:}
\begin{itemize}
    \item Polinomio de grado 4095 sobre BLS12-381
    \item Cada ``field element'' es un coeficiente
    \item KZG commitment: un punto de curva (48 bytes)
\end{itemize}

\textbf{En el bloque:}
\begin{itemize}
    \item Header contiene hashes de blobs
    \item Blobs se propagan separadamente
    \item Solo el hash va al execution payload
\end{itemize}
\end{definition}

\begin{mathnote}
\textbf{KZG Commitment para Blobs:}

Cada blob tiene un KZG commitment que permite:
\begin{enumerate}
    \item Verificar que ciertos datos están en el blob
    \item Sin descargar el blob completo
    \item Proofs de evaluación en puntos específicos
\end{enumerate}

\textbf{Proceso:}
\begin{verbatim}
Blob data (128 KB)
    ↓
Interpretar como polinomio p(x)
    ↓
Commitment: C = [p(τ)]_1 (punto en G1)
    ↓
Para probar p(z) = y:
    - Quotient: q(x) = (p(x) - y) / (x - z)
    - Proof: π = [q(τ)]_1
    - Verify: e(C - [y]_1, [1]_2) = e(π, [τ - z]_2)
\end{verbatim}

Esto es la base para DAS en Full Danksharding.
\end{mathnote}

%-------------------------------------------
\subsection{Blob Transactions}

\begin{definition}[Blob Transaction (Type 3)]
Nueva transacción que incluye blobs:

\begin{verbatim}
{
  type: 3,
  chainId: ...,
  nonce: ...,
  maxFeePerGas: ...,
  maxPriorityFeePerGas: ...,
  gasLimit: ...,
  to: ...,
  value: ...,
  data: ...,
  accessList: [...],
  
  // Nuevos campos para blobs:
  maxFeePerBlobGas: ...,
  blobVersionedHashes: [hash1, hash2, ...],
  
  // Blobs van en "sidecar" separado
}
\end{verbatim}
\end{definition}

\begin{keyidea}
\textbf{Mercado de fees separado para blobs:}

\textbf{Execution gas (existente):}
\begin{itemize}
    \item Paga por computación y calldata
    \item Base fee + priority fee
    \item Target: 15M gas/bloque
\end{itemize}

\textbf{Blob gas (nuevo):}
\begin{itemize}
    \item Paga por espacio de blob
    \item Su propio base fee (empieza muy bajo)
    \item Target: 3 blobs/bloque (384 KB)
    \item Max: 6 blobs/bloque (768 KB)
\end{itemize}

\textbf{Ajuste del blob base fee:}
\begin{itemize}
    \item Mismo mecanismo EIP-1559
    \item Pero con parámetros diferentes
    \item Puede subir/bajar más rápido que el normal
\end{itemize}
\end{keyidea}

%-------------------------------------------
\subsection{Proto-Danksharding vs Full Danksharding}

\begin{itemize}
    \item \textbf{EIP-4844 (actual):} 6 blobs/bloque, $\sim$1 MB/bloque
    \item \textbf{Full Danksharding:} 64+ blobs, Data Availability Sampling
    \item DAS permite verificar DA sin descargar todo
\end{itemize}

\begin{mathnote}
\textbf{Comparación:}

\begin{center}
\begin{tabular}{lcc}
\toprule
\textbf{Aspecto} & \textbf{Proto-Dank (4844)} & \textbf{Full Danksharding} \\
\midrule
Blobs/bloque & 6 & 64-256 \\
DA por bloque & $\sim$750 KB & $\sim$8-32 MB \\
DAS & No & Sí \\
Verificación & Full nodes descargan todo & Light nodes muestrean \\
Timeline & Live (Mar 2024) & 2025-2026+ \\
\bottomrule
\end{tabular}
\end{center}

\textbf{Por qué DAS es necesario para más blobs:}
\begin{itemize}
    \item Con 6 blobs: full nodes pueden descargar todo
    \item Con 256 blobs: 32 MB/bloque = 160 Mbps solo para DA
    \item DAS permite verificar sin descargar todo
    \item Mantiene descentralización con más datos
\end{itemize}
\end{mathnote}

%-------------------------------------------
\subsection{Impacto de EIP-4844}

\begin{remark}
\textbf{Reducción de costos post-Dencun:}

\begin{center}
\begin{tabular}{lccc}
\toprule
\textbf{Rollup} & \textbf{Pre-4844} & \textbf{Post-4844} & \textbf{Reducción} \\
\midrule
Arbitrum & \$0.10-0.50 & \$0.001-0.01 & 90-99\% \\
Optimism & \$0.10-0.50 & \$0.001-0.01 & 90-99\% \\
Base & \$0.05-0.20 & \$0.001-0.005 & 95-99\% \\
zkSync Era & \$0.20-1.00 & \$0.01-0.05 & 90-95\% \\
\bottomrule
\end{tabular}
\end{center}

\textbf{Nota:} Los costos varían según congestión. El blob base fee puede subir si la demanda excede el target.

\textbf{Impacto en adopción:}
\begin{itemize}
    \item Txs de \$0.01 hacen viable micropagos
    \item Gaming y social apps ahora factibles
    \item Más usuarios pueden permitirse usar L2s
\end{itemize}
\end{remark}

\begin{bancozk}
\textbf{Data Availability en Banco ZK:}

Banco ZK usa una estrategia híbrida de DA:

\textbf{Datos on-chain (blobs):}
\begin{itemize}
    \item Commitments de notas (encriptados)
    \item Nullifiers (públicos, para prevenir doble gasto)
    \item State root updates
\end{itemize}

\textbf{¿Por qué no validium?}

Aunque más barato, validium requiere confiar en el comité:
\begin{itemize}
    \item Si comité colude con operador, pueden robar fondos
    \item Para un sistema privado, la confianza es especialmente problemática
    \item Preferimos seguridad de L1 aunque cueste más
\end{itemize}

\textbf{Optimizaciones de DA:}
\begin{itemize}
    \item Solo publicamos commitments, no datos completos
    \item Compresión agresiva de nullifiers
    \item Batch aggregation para amortizar costos
\end{itemize}

\textbf{Costo estimado por tx:}
\begin{itemize}
    \item Pre-4844: \$0.10-0.30
    \item Post-4844: \$0.01-0.05
    \item Target futuro: $<$\$0.01
\end{itemize}
\end{bancozk}

% ============================================
% BLOQUE 3: ESCALABILIDAD (L2/L3) - VERSIÓN EXPANDIDA
% Parte 5: L3, App-chains, Bridges, Resumen
% ============================================

%===========================================
\section{L3 y App-chains}
%===========================================

\begin{definition}[L3]
Layer 3: rollup sobre un L2.
\end{definition}

Motivaciones:
\begin{itemize}
    \item \textbf{Escalabilidad:} Más compresión
    \item \textbf{Personalización:} Parámetros específicos
    \item \textbf{Privacidad:} Features especializadas
    \item \textbf{Aislamiento:} Separar aplicaciones
\end{itemize}

%-------------------------------------------
\subsection{¿Por Qué L3?}

\begin{intuicion}
Si L2 escala L1, ¿por qué no L3 escala L2?

La respuesta no es tan simple:
\begin{itemize}
    \item Cada capa añade latencia
    \item Cada capa tiene su propio sequencer/prover
    \item Los beneficios de escalabilidad tienen rendimientos decrecientes
\end{itemize}

\textbf{Entonces, ¿cuál es el punto de L3?}

No es solo escalabilidad, sino \textbf{especialización}:
\begin{itemize}
    \item L2 general-purpose (Arbitrum, zkSync)
    \item L3 específico para una app o caso de uso
    \item Parámetros optimizados para ese caso
\end{itemize}
\end{intuicion}

\begin{keyidea}
\textbf{Casos de uso para L3:}

\textbf{1. Privacidad:}
\begin{itemize}
    \item L2 público para la mayoría de txs
    \item L3 privado para txs confidenciales
    \item Banco ZK podría ser L3 privado sobre un L2
\end{itemize}

\textbf{2. Gaming:}
\begin{itemize}
    \item Alto volumen de txs pequeñas
    \item Latencia ultra-baja necesaria
    \item Puede permitirse DA más barata (validium)
    \item Items de bajo valor = menor necesidad de seguridad
\end{itemize}

\textbf{3. Enterprise:}
\begin{itemize}
    \item Permisos y compliance específicos
    \item Aislamiento de otras apps
    \item Gobernanza controlada
\end{itemize}

\textbf{4. Experimentación:}
\begin{itemize}
    \item Probar nuevos mecanismos
    \item Fallos no afectan al L2 principal
    \item Iteración rápida
\end{itemize}
\end{keyidea}

\begin{mathnote}
\textbf{Vitalik sobre L2 vs L3:}

En su post ``What kind of layer 3s make sense?'' (2022):

\textbf{L3 para escalabilidad pura:} No tiene mucho sentido
\begin{itemize}
    \item Si L2 da 100x, L3 daría 10000x
    \item Pero L2 puede simplemente usar más DA
    \item No necesitas otra capa
\end{itemize}

\textbf{L3 para customización:} Tiene sentido
\begin{itemize}
    \item L2: plataforma general (EVM)
    \item L3: app-specific (gaming VM, privacy VM)
    \item Diferentes trade-offs para diferentes usos
\end{itemize}

\textbf{Arquitectura recomendada:}
\begin{verbatim}
L1 (Ethereum): Seguridad + DA
    ↑
L2 (Arbitrum/zkSync): Ejecución general
    ↑
L3 (App-chain): Customización específica
\end{verbatim}
\end{mathnote}

%-------------------------------------------
\subsection{App-chains}

\begin{definition}[App-chain]
Blockchain/rollup dedicado a una sola aplicación.
\end{definition}

Frameworks: Arbitrum Orbit, OP Stack, Polygon CDK, zkSync Hyperchains.

\begin{remark}
\textbf{Frameworks para App-chains:}

\textbf{Arbitrum Orbit:}
\begin{itemize}
    \item Fork de Arbitrum One/Nova
    \item Puede ser L2 o L3
    \item Customizable: DA, gas token, permisos
    \item Licencia específica de Arbitrum
\end{itemize}

\textbf{OP Stack:}
\begin{itemize}
    \item Framework modular de Optimism
    \item Base, Zora, Mode lo usan
    \item Open source (MIT license)
    \item ``Superchain'' vision de interoperabilidad
\end{itemize}

\textbf{Polygon CDK:}
\begin{itemize}
    \item Chain Development Kit de Polygon
    \item ZK-based (zkEVM)
    \item Puede conectar al ``AggLayer''
    \item Validium o rollup configurable
\end{itemize}

\textbf{zkSync Hyperchains:}
\begin{itemize}
    \item Chains conectadas a zkSync Era
    \item Comparten prover infrastructure
    \item Interoperabilidad nativa
    \item ZK-based
\end{itemize}
\end{remark}

\begin{example}[App-chains en producción]
\begin{center}
\begin{tabular}{llll}
\toprule
\textbf{App-chain} & \textbf{Tipo} & \textbf{Framework} & \textbf{Caso de uso} \\
\midrule
Xai & L3 & Arbitrum Orbit & Gaming (Arbitrum) \\
Degen Chain & L3 & Arbitrum Orbit & Social/Memes \\
Zora & L2 & OP Stack & NFTs \\
Worldcoin & L2 & OP Stack & Identity \\
Immutable X & L2 & StarkEx & Gaming NFTs \\
dYdX v4 & L1 & Cosmos SDK & Derivatives \\
\bottomrule
\end{tabular}
\end{center}

\textbf{Nota:} dYdX migró de StarkEx L2 a su propia L1 en Cosmos, mostrando que la elección óptima depende del caso de uso.
\end{example}

%-------------------------------------------
\subsection{Interoperabilidad entre L2s/L3s}

\begin{advertencia}
\textbf{El problema de fragmentación:}

Con docenas de L2s y cientos de L3s:
\begin{itemize}
    \item Liquidez fragmentada entre chains
    \item UX complicada (bridging constante)
    \item Composability perdida (no puedes llamar contratos cross-chain atómicamente)
\end{itemize}

\textbf{Soluciones en desarrollo:}

\textbf{1. Shared sequencing:}
\begin{itemize}
    \item Mismo sequencer ordena txs de múltiples chains
    \item Permite atomic bundles cross-chain
\end{itemize}

\textbf{2. Proof aggregation:}
\begin{itemize}
    \item Múltiples L3s comparten un proof en L2/L1
    \item Reduce costo de verificación
\end{itemize}

\textbf{3. Messaging protocols:}
\begin{itemize}
    \item LayerZero, Hyperlane, Axelar
    \item Mensajes verificables cross-chain
    \item No atomic, pero funciona
\end{itemize}

\textbf{4. Superchain/AggLayer:}
\begin{itemize}
    \item Visión de chains interconectadas
    \item Standards comunes
    \item Interoperabilidad nativa
\end{itemize}
\end{advertencia}

%===========================================
\section{Bridges}
%===========================================

\begin{definition}[Bridge]
Protocolo para transferir activos entre blockchains.
\end{definition}

%-------------------------------------------
\subsection{Cómo Funcionan}

\begin{enumerate}
    \item Usuario deposita en contrato de origen
    \item Bridge detecta depósito
    \item Emite tokens equivalentes en destino
    \item Para retirar, proceso inverso
\end{enumerate}

\begin{center}
\begin{tikzpicture}[
    chain/.style={draw, rectangle, minimum width=3cm, minimum height=1.5cm, align=center},
    bridge/.style={draw, ellipse, minimum width=2cm, fill=yellow!30},
    arrow/.style={->, thick}
]
    % Chains
    \node[chain, fill=blue!20] (origin) at (0, 0) {Chain A\\(Origen)};
    \node[chain, fill=green!20] (dest) at (8, 0) {Chain B\\(Destino)};
    
    % Bridge
    \node[bridge] (bridge) at (4, 0) {Bridge};
    
    % Flujo de depósito
    \draw[arrow, blue] (0, 1.5) -- node[above, font=\small] {1. Lock ETH} (0, 0.75);
    \draw[arrow, blue] (origin) -- node[above, font=\small] {2. Detecta} (bridge);
    \draw[arrow, blue] (bridge) -- node[above, font=\small] {3. Mint wETH} (dest);
    \draw[arrow, blue] (8, 0.75) -- node[above, font=\small] {4. Recibe} (8, 1.5);
    
    % Flujo de retiro
    \draw[arrow, red, dashed] (8, -0.75) -- (8, -1.5);
    \draw[arrow, red, dashed] (dest) -- node[below, font=\small] {Burn} (bridge);
    \draw[arrow, red, dashed] (bridge) -- node[below, font=\small] {Unlock} (origin);
    \draw[arrow, red, dashed] (0, -1.5) -- (0, -0.75);
\end{tikzpicture}
\end{center}

%-------------------------------------------
\subsection{Tipos de Bridges}

\begin{definition}[Tipos de Bridges]
\textbf{1. Native (Canonical):}
\begin{itemize}
    \item Integrado en el rollup
    \item Usa el mecanismo de seguridad del rollup
    \item El más seguro para ese rollup
    \item Ejemplo: Bridge nativo de Arbitrum/Optimism
\end{itemize}

\textbf{2. Lock-and-mint:}
\begin{itemize}
    \item Bloquear token en chain A
    \item Mintear wrapped token en chain B
    \item Wrapped token representa el original
    \item Ejemplo: WBTC (Bitcoin wrapped en Ethereum)
\end{itemize}

\textbf{3. Burn-and-mint:}
\begin{itemize}
    \item Quemar token en chain A
    \item Mintear token ``nativo'' en chain B
    \item Token no es wrapped, es nuevo
    \item Requiere coordinación entre chains
    \item Ejemplo: USDC (Circle mint/burn)
\end{itemize}

\textbf{4. Liquidity Networks:}
\begin{itemize}
    \item Pools de liquidez en ambas chains
    \item No hay lock/mint, solo swaps
    \item Más rápido pero limitado por liquidez
    \item Ejemplo: Hop, Across, Stargate
\end{itemize}
\end{definition}

\begin{mathnote}
\textbf{Comparación de tipos:}

\begin{center}
\begin{tabular}{lcccc}
\toprule
\textbf{Tipo} & \textbf{Velocidad} & \textbf{Liquidez} & \textbf{Seguridad} & \textbf{Costo} \\
\midrule
Native & Lento (7d opt) & Ilimitada & Máxima & Medio \\
Lock-mint & Medio & Ilimitada & Variable & Medio \\
Burn-mint & Medio & Ilimitada & Alta & Medio \\
Liquidity & Rápido & Limitada & Alta & Alto (slippage) \\
\bottomrule
\end{tabular}
\end{center}
\end{mathnote}

%-------------------------------------------
\subsection{Mecanismos de Verificación}

\begin{keyidea}
\textbf{¿Cómo sabe Chain B que el depósito en Chain A es real?}

\textbf{1. Multisig/MPC:}
\begin{itemize}
    \item Grupo de firmantes atestiguan el depósito
    \item Si M-de-N firman, se considera válido
    \item Trust assumption: mayoría honesta
    \item Problema: punto único de fallo
\end{itemize}

\textbf{2. Optimistic (con challenge):}
\begin{itemize}
    \item Cualquiera puede atestiguar
    \item Período de disputa para cuestionar
    \item Fraud proof si attestation es falsa
    \item Lento pero trustless
\end{itemize}

\textbf{3. Light client:}
\begin{itemize}
    \item Verificar headers de la otra chain
    \item Verificar Merkle proof de la tx
    \item Trustless si light client es correcto
    \item Requiere implementar consenso de la otra chain
\end{itemize}

\textbf{4. ZK:}
\begin{itemize}
    \item Generar proof de que tx existe en chain A
    \item Verificar proof en chain B
    \item Trustless y rápido
    \item Costoso de generar proofs
\end{itemize}
\end{keyidea}

%-------------------------------------------
\subsection{Riesgos de Bridges}

Hacks históricos:
\begin{itemize}
    \item Ronin: \$625M (multisig comprometido)
    \item Wormhole: \$320M (bug en verificación)
    \item Nomad: \$190M (bug permite mint arbitrario)
\end{itemize}

\begin{advertencia}
\textbf{Anatomía de los hacks de bridges:}

\textbf{Ronin (Marzo 2022):}
\begin{itemize}
    \item Bridge controlado por 9 validadores, threshold 5
    \item 4 de Sky Mavis + 1 de Axie DAO comprometidos
    \item Atacante obtuvo 5 claves privadas
    \item Firmó retiros fraudulentos
    \item \$625M robados, no detectado por 6 días
\end{itemize}

\textbf{Wormhole (Febrero 2022):}
\begin{itemize}
    \item Bug en verificación de firmas
    \item Atacante bypassed signature check
    \item Minteó 120,000 wETH sin depósito real
    \item Jump Crypto (backer) cubrió las pérdidas
\end{itemize}

\textbf{Nomad (Agosto 2022):}
\begin{itemize}
    \item Bug en inicialización de contrato
    \item Cualquier mensaje se consideraba válido
    \item Cientos de atacantes copiaron el exploit
    \item ``Decentralized robbery''
\end{itemize}

\textbf{Lecciones:}
\begin{itemize}
    \item Multisigs son puntos únicos de fallo
    \item Bugs en verificación son catastróficos
    \item Bridges son los targets más valiosos
    \item ZK bridges son más seguros pero más complejos
\end{itemize}
\end{advertencia}

%-------------------------------------------
\subsection{ZK Bridges}

Soluciones:
\begin{itemize}
    \item ZK bridges: verificar consenso en ZK
    \item Light client bridges: verificar headers
    \item Optimistic bridges: período de disputa
\end{itemize}

\begin{definition}[ZK Bridge]
Bridge que usa pruebas ZK para verificar eventos en la chain de origen.
\end{definition}

\begin{mathnote}
\textbf{Cómo funciona un ZK Bridge:}

\begin{verbatim}
Chain A (origen):
1. Usuario deposita en contrato
2. Tx incluida en bloque
3. Bloque tiene header firmado por validators

Prover (off-chain):
4. Obtiene header y signatures de Chain A
5. Genera ZK proof de:
   - Header es válido (cumple reglas de consenso)
   - Signatures son válidas (>2/3 de validators)
   - Tx de depósito está en el bloque (Merkle proof)
6. Proof es pequeño (~200 bytes)

Chain B (destino):
7. Verifica ZK proof (barato, ~200K gas)
8. Si válido: mintea tokens
9. No necesita confiar en nadie
\end{verbatim}

\textbf{Ventajas:}
\begin{itemize}
    \item Trustless: solo confías en matemáticas
    \item Rápido: verificación instantánea
    \item Barato: proof pequeño
\end{itemize}

\textbf{Desafíos:}
\begin{itemize}
    \item Probar consenso de otra chain es complejo
    \item Cada chain tiene reglas diferentes
    \item Generar proofs toma tiempo
\end{itemize}
\end{mathnote}

\begin{remark}
\textbf{Proyectos de ZK Bridges:}

\begin{center}
\begin{tabular}{lll}
\toprule
\textbf{Proyecto} & \textbf{Enfoque} & \textbf{Estado} \\
\midrule
Succinct & ZK light clients genéricos & Producción \\
Polyhedra (zkBridge) & Cross-chain messaging & Producción \\
Lagrange & State proofs & Desarrollo \\
Electron Labs & Cosmos $\leftrightarrow$ Ethereum & Desarrollo \\
\bottomrule
\end{tabular}
\end{center}
\end{remark}

\begin{bancozk}
\textbf{Bridge de Banco ZK:}

El bridge de Banco ZK conecta L1 (Ethereum) con nuestro L2 privado:

\textbf{Depósitos (L1 → L2):}
\begin{enumerate}
    \item Usuario deposita ETH/tokens en contrato L1
    \item Genera commitment secreto off-chain
    \item Contrato registra el depósito
    \item Sequencer detecta evento
    \item Crea nota privada equivalente en L2
    \item Usuario tiene acceso con su secreto
\end{enumerate}

\textbf{Retiros (L2 → L1):}
\begin{enumerate}
    \item Usuario genera nullifier y ZK proof
    \item Proof demuestra: tiene nota válida sin revelar cuál
    \item Sequencer incluye en batch
    \item Batch publicado en L1 con ZK proof agregado
    \item Contrato L1 verifica proof
    \item Fondos liberados a dirección de retiro
\end{enumerate}

\textbf{Seguridad:}
\begin{itemize}
    \item Bridge nativo: misma seguridad que el rollup
    \item No hay multisig que pueda robar fondos
    \item Si sequencer es malicioso: forced exit via L1
    \item Privacidad: nadie puede vincular depósito con retiro
\end{itemize}

\textbf{Diferencia con bridges externos:}
\begin{itemize}
    \item Bridge nativo: trustless, lento (esperar proof)
    \item LP bridges (Hop/Across): rápido, pero no privados
    \item Para Banco ZK usamos principalmente el nativo
\end{itemize}
\end{bancozk}

%===========================================
\section{Resumen del Bloque 3}
%===========================================

\begin{keyidea}
\textbf{Puntos clave de Escalabilidad (L2/L3):}

\textbf{1. El Trilema:}
\begin{itemize}
    \item Escalabilidad vs Descentralización vs Seguridad
    \item No es un teorema, pero captura trade-off real
    \item L2s intentan ``escapar'' el trilema
\end{itemize}

\textbf{2. L2 = Off-chain execution + L1 security:}
\begin{itemize}
    \item Procesar txs fuera de L1
    \item Publicar resultados/proofs en L1
    \item Heredar seguridad de L1
\end{itemize}

\textbf{3. Sidechains $\neq$ Rollups:}
\begin{itemize}
    \item Sidechains: seguridad propia (validadores)
    \item Rollups: seguridad de L1 (proofs)
    \item Rollups son ``verdaderos L2s''
\end{itemize}

\textbf{4. Optimistic Rollups:}
\begin{itemize}
    \item Asume validez, permite disputas
    \item Fraud proofs (interactive bisection)
    \item 7 días de challenge period
    \item Ejemplos: Arbitrum, Optimism, Base
\end{itemize}

\textbf{5. ZK Rollups:}
\begin{itemize}
    \item Prueba criptográfica de validez
    \item Finalidad rápida (minutos)
    \item zkEVM para compatibilidad
    \item Ejemplos: zkSync, StarkNet, Scroll
\end{itemize}

\textbf{6. Sequencers:}
\begin{itemize}
    \item Actualmente centralizados
    \item Riesgos: censura, MEV, downtime
    \item Futuro: descentralización, shared sequencing, based rollups
\end{itemize}

\textbf{7. Data Availability:}
\begin{itemize}
    \item Crítico para seguridad
    \item Opciones: calldata, blobs, DAC, DA layers
    \item EIP-4844 (blobs) redujo costos 90\%+
\end{itemize}

\textbf{8. L3 y App-chains:}
\begin{itemize}
    \item L3: rollup sobre L2
    \item Motivación: especialización, no solo escalabilidad
    \item Frameworks: Orbit, OP Stack, Polygon CDK
\end{itemize}

\textbf{9. Bridges:}
\begin{itemize}
    \item Conectan diferentes chains
    \item Tipos: native, lock-mint, liquidity
    \item Riesgos: hacks multimillonarios históricos
    \item ZK bridges: futuro trustless
\end{itemize}

\textbf{10. Visión general:}
\begin{itemize}
    \item L1 (Ethereum): seguridad y DA
    \item L2s: ejecución escalable
    \item L3s/App-chains: especialización
    \item Bridges: interoperabilidad
    \item El ecosistema está convergiendo hacia esta arquitectura modular
\end{itemize}
\end{keyidea}

\cleardoublepage
% ============================================
% BLOQUE 4: ZERO-KNOWLEDGE PROOFS (CORE) - VERSIÓN EXPANDIDA
% Parte 1: Definición y Propiedades Fundamentales
% ============================================
\chapter{Zero-Knowledge Proofs (Core)}
\label{ch:bloque4}

\begin{center}
\textit{Bloque 4 de 20}
\end{center}

%===========================================
\section{Qué es una ZK Proof}
%===========================================

\begin{definition}[Prueba de Conocimiento Cero]
Protocolo donde un Prover convence a un Verifier de que una afirmación es verdadera, sin revelar información adicional.
\end{definition}

\begin{intuicion}
Imagina que quieres probar que conoces la contraseña de una caja fuerte sin revelar la contraseña. Una ZK proof permite exactamente esto: convencer a alguien de que \textit{sabes algo} sin revelar \textit{qué es} ese algo.

Esto parece mágico, pero es matemáticamente riguroso y tiene aplicaciones profundas en criptografía, privacidad, y escalabilidad de blockchains.
\end{intuicion}

%-------------------------------------------
\subsection{Historia y Contexto}

\begin{remark}
\textbf{La historia de las Zero-Knowledge Proofs:}

\begin{center}
\begin{tabular}{lll}
\toprule
\textbf{Año} & \textbf{Autores} & \textbf{Contribución} \\
\midrule
1985 & Goldwasser, Micali, Rackoff & Definición formal de ZK \\
1986 & Fiat, Shamir & Transformación a no-interactivo \\
1988 & Goldreich, Micali, Wigderson & ZK para todo NP \\
1991 & Schnorr & Protocolo eficiente de identificación \\
2010 & Groth & Pruebas succinctas eficientes \\
2012 & Bitansky et al. & SNARKs \\
2013 & Gennaro et al. & Pinocchio (primer SNARK práctico) \\
2016 & Groth & Groth16 (el más eficiente) \\
2018 & Ben-Sasson et al. & STARKs (sin trusted setup) \\
2019 & Gabizon et al. & PLONK (universal setup) \\
2019 & Bünz et al. & Bulletproofs recursivos \\
2020 & Zcash team & Halo (recursión sin trusted setup) \\
\bottomrule
\end{tabular}
\end{center}

El paper original de 1985 ``The Knowledge Complexity of Interactive Proof Systems'' ganó el premio Gödel en 1993 y sentó las bases de toda la criptografía moderna basada en ZK.
\end{remark}

%-------------------------------------------
\subsection{La Intuición: Cueva de Ali Baba}

\begin{example}[Cueva de Ali Baba]
Probar que conoces la palabra secreta de una puerta sin revelarla:
\begin{enumerate}
    \item Victor espera afuera, Peggy entra y elige camino A o B
    \item Victor entra y pide ``sal por A'' o ``sal por B''
    \item Peggy sale por el camino pedido (atravesando puerta si necesario)
    \item Repetir 20 veces: probabilidad de engaño $< 0.0001\%$
\end{enumerate}
\end{example}

\begin{center}
\begin{tikzpicture}[scale=0.8]
    % Cueva
    \draw[thick] (0, 0) -- (0, 4);
    \draw[thick] (0, 4) -- (-3, 6);
    \draw[thick] (0, 4) -- (3, 6);
    \draw[thick] (-3, 6) arc (180:0:3);
    
    % Puerta mágica
    \draw[red, very thick] (0, 8) -- (0, 9);
    \node[red] at (0.5, 8.5) {Puerta};
    
    % Caminos
    \node at (-2, 5) {A};
    \node at (2, 5) {B};
    
    % Personas
    \node[draw, circle, fill=blue!30] (peggy) at (-2, 7) {P};
    \node[draw, circle, fill=green!30] (victor) at (0, 2) {V};
    
    % Entrada
    \node at (0, -0.5) {Entrada};
    
    % Descripción
    \node[text width=6cm, align=left] at (8, 5) {
        \textbf{Protocolo:}\\
        1. Peggy entra, elige A o B\\
        2. Victor entra, dice ``sal por X''\\
        3. Peggy sale por X\\
        4. Si conoce la palabra: siempre puede\\
        5. Si no conoce: 50\% de fallar\\
        6. Después de 20 rondas: $2^{-20}$ de engaño
    };
\end{tikzpicture}
\end{center}

\begin{mathnote}
\textbf{Análisis del protocolo:}

Sea $p$ la probabilidad de que Peggy NO conozca la palabra secreta pero convenza a Victor.

\textbf{Una ronda:}
\begin{itemize}
    \item Peggy elige camino (digamos A) sin saber qué pedirá Victor
    \item Victor pide A o B con probabilidad 1/2 cada uno
    \item Si pide A: Peggy sale (éxito sin importar si conoce la palabra)
    \item Si pide B: Peggy necesita la palabra para cruzar
\end{itemize}

$P(\text{éxito sin palabra}) = 1/2$

\textbf{$n$ rondas:}
$P(\text{éxito en todas sin palabra}) = (1/2)^n$

Con $n = 20$: $P = 2^{-20} \approx 10^{-6}$
Con $n = 40$: $P = 2^{-40} \approx 10^{-12}$

\textbf{¿Por qué es Zero-Knowledge?}

Victor no aprende nada sobre la palabra porque:
\begin{itemize}
    \item Solo ve a Peggy salir por el camino que pidió
    \item Esto ocurriría 100\% si conoce la palabra
    \item No puede distinguir si cruzó la puerta o no
    \item La interacción no revela la palabra
\end{itemize}
\end{mathnote}

%-------------------------------------------
\subsection{Definición Formal}

\begin{definition}[Sistema de Prueba Interactivo]
Un sistema de prueba interactivo para un lenguaje $L$ consiste en:
\begin{itemize}
    \item \textbf{Prover $\mathcal{P}$:} Algoritmo (posiblemente sin límite de tiempo)
    \item \textbf{Verifier $\mathcal{V}$:} Algoritmo de tiempo polinomial
    \item \textbf{Interacción:} Secuencia de mensajes $\langle \mathcal{P}, \mathcal{V} \rangle(x)$
    \item \textbf{Resultado:} $\mathcal{V}$ acepta o rechaza
\end{itemize}
\end{definition}

\begin{definition}[Prueba de Conocimiento Cero - Formal]
Un sistema de prueba interactivo $(\mathcal{P}, \mathcal{V})$ para $L$ es \textbf{zero-knowledge} si satisface:

\textbf{1. Completeness:}
\[ \forall x \in L, \forall w \text{ válido}: \Pr[\langle \mathcal{P}(w), \mathcal{V} \rangle(x) = 1] = 1 \]

\textbf{2. Soundness:}
\[ \forall x \notin L, \forall \mathcal{P}^*: \Pr[\langle \mathcal{P}^*, \mathcal{V} \rangle(x) = 1] \leq \epsilon \]

\textbf{3. Zero-Knowledge:}
\[ \exists \text{ Simulador } \mathcal{S}: \forall \mathcal{V}^*, \text{View}_{\mathcal{V}^*}[\langle \mathcal{P}, \mathcal{V}^* \rangle(x)] \approx \mathcal{S}(x) \]
\end{definition}

\begin{keyidea}
\textbf{Lo que significa cada propiedad intuitivamente:}

\textbf{Completeness:} ``Si digo la verdad, siempre me creerán''
\begin{itemize}
    \item Un prover honesto con un witness válido siempre convence
    \item No hay falsos negativos
\end{itemize}

\textbf{Soundness:} ``Si miento, me descubrirán (casi seguro)''
\begin{itemize}
    \item Un prover malicioso no puede convencer si la afirmación es falsa
    \item Excepto con probabilidad negligible $\epsilon$
\end{itemize}

\textbf{Zero-Knowledge:} ``No revelo más de lo necesario''
\begin{itemize}
    \item El verificador solo aprende que $x \in L$
    \item Nada sobre el witness $w$
    \item Nada que no pudiera calcular solo
\end{itemize}
\end{keyidea}

%===========================================
\section{Propiedades: Completeness, Soundness, Zero-Knowledge}
%===========================================

\subsection{Completeness (Completitud)}

Si la afirmación es verdadera, un prover honesto siempre convence al verificador.
\[
x \in L \implies \Pr[\text{Verifier acepta}] = 1
\]

\begin{definition}[Completeness - Tipos]
\textbf{Perfect Completeness:}
\[ \Pr[\langle \mathcal{P}(w), \mathcal{V} \rangle(x) = 1] = 1 \]

El verificador acepta con probabilidad exactamente 1.

\textbf{Statistical Completeness:}
\[ \Pr[\langle \mathcal{P}(w), \mathcal{V} \rangle(x) = 1] \geq 1 - \text{negl}(\lambda) \]

El verificador acepta con probabilidad abrumadora.
\end{definition}

\begin{example}[Completeness en protocolo Schnorr]
Para probar conocimiento de $x$ tal que $Y = g^x$:

\begin{verbatim}
Prover (conoce x):
    1. Elige r aleatorio
    2. Envía R = g^r
    
Verifier:
    3. Envía challenge c aleatorio
    
Prover:
    4. Calcula z = r + c*x
    5. Envía z
    
Verifier:
    6. Acepta si g^z = R * Y^c
\end{verbatim}

\textbf{Verificación de completeness:}
\[ g^z = g^{r + cx} = g^r \cdot g^{cx} = R \cdot (g^x)^c = R \cdot Y^c \checkmark \]

Si el prover es honesto, la verificación siempre pasa.
\end{example}

%-------------------------------------------
\subsection{Soundness (Solidez)}

Si la afirmación es falsa, ningún prover puede convencer al verificador (excepto con probabilidad negligible).
\[
x \notin L \implies \Pr[\text{Verifier acepta}] \leq \epsilon
\]

\begin{definition}[Tipos de Soundness]
\textbf{Statistical Soundness:}
\[ \Pr[\langle \mathcal{P}^*, \mathcal{V} \rangle(x) = 1] \leq 2^{-\lambda} \]

Seguro contra provers computacionalmente ilimitados.

\textbf{Computational Soundness (Argument):}
\[ \Pr[\langle \mathcal{P}^*, \mathcal{V} \rangle(x) = 1] \leq \text{negl}(\lambda) \]

Solo seguro contra provers de tiempo polinomial.
\end{definition}

\begin{mathnote}
\textbf{Proof vs Argument:}

\textbf{Proof} (prueba): Soundness estadístico
\begin{itemize}
    \item Seguro incluso contra adversarios con poder computacional infinito
    \item Basado en teoría de la información
    \item Típicamente pruebas más grandes
\end{itemize}

\textbf{Argument} (argumento): Soundness computacional
\begin{itemize}
    \item Seguro solo contra adversarios eficientes
    \item Basado en asunciones criptográficas
    \item Permite pruebas más pequeñas (SNARKs)
\end{itemize}

SNARKs son \textit{arguments}, no \textit{proofs} en el sentido técnico.
\end{mathnote}

\textbf{Knowledge soundness:} Además, si el verifier acepta, existe un ``extractor'' que puede obtener el witness del prover.

\begin{definition}[Knowledge Soundness]
Un sistema tiene knowledge soundness si existe un extractor $\mathcal{E}$ tal que:

Para todo prover $\mathcal{P}^*$ que hace que el verificador acepte con probabilidad $\epsilon$, $\mathcal{E}$ puede extraer un witness válido $w$ con probabilidad $\epsilon - \text{negl}(\lambda)$.

\textbf{Formalmente:}
\[ \Pr[\langle \mathcal{P}^*, \mathcal{V} \rangle(x) = 1] = \epsilon \implies \Pr[\mathcal{E}^{\mathcal{P}^*}(x) = w : R(x,w) = 1] \geq \epsilon - \text{negl}(\lambda) \]
\end{definition}

\begin{keyidea}
\textbf{¿Por qué Knowledge Soundness importa?}

Soundness normal solo dice: ``no puedes convencer si $x \notin L$''

Knowledge soundness dice: ``si convences, realmente CONOCES un witness''

Esto es crucial para aplicaciones como:
\begin{itemize}
    \item Probar que conoces una clave privada
    \item Probar que conoces la preimagen de un hash
    \item Probar que ejecutaste un cálculo correctamente
\end{itemize}

Sin knowledge soundness, podrías ``adivinar'' la prueba correcta sin realmente conocer el secreto (aunque con probabilidad exponencialmente pequeña).
\end{keyidea}

\begin{example}[Extracción en Schnorr]
Para el protocolo Schnorr, el extractor funciona así:

\begin{verbatim}
Extractor E:
    1. Ejecuta P* hasta que envía R
    2. Envía challenge c1, obtiene z1
    3. "Rebobina" P* al estado después de R
    4. Envía challenge c2 ≠ c1, obtiene z2
    
    5. Si ambos verifican:
       g^z1 = R * Y^c1
       g^z2 = R * Y^c2
       
    6. Entonces:
       g^(z1-z2) = Y^(c1-c2)
       g^((z1-z2)/(c1-c2)) = Y
       x = (z1-z2)/(c1-c2)  ← ¡Witness extraído!
\end{verbatim}

Esto demuestra que si $\mathcal{P}^*$ puede responder a dos challenges diferentes para el mismo $R$, el extractor puede calcular $x$.
\end{example}

%-------------------------------------------
\subsection{Zero-Knowledge}

El verificador no aprende nada más allá de ``$x \in L$''.

\begin{definition}[Simulador]
Existe un algoritmo (simulador) que puede generar transcripciones indistinguibles de interacciones reales, sin conocer el witness.
\end{definition}

\begin{mathnote}
\textbf{La definición formal de Zero-Knowledge:}

Sea $\text{View}_{\mathcal{V}^*}[\langle \mathcal{P}(w), \mathcal{V}^* \rangle(x)]$ la ``vista'' del verificador:
\begin{itemize}
    \item Todos los mensajes que recibió
    \item Su aleatoriedad interna
    \item Su estado final
\end{itemize}

Un protocolo es ZK si existe un simulador $\mathcal{S}$ que, dado solo $x$ (sin $w$), puede producir una vista indistinguible:
\[ \text{View}_{\mathcal{V}^*}[\langle \mathcal{P}(w), \mathcal{V}^* \rangle(x)] \approx \mathcal{S}(x) \]

\textbf{Intuición:} Si el simulador puede producir la misma ``experiencia'' sin el witness, entonces el verificador no pudo haber aprendido nada del witness durante la interacción real.
\end{mathnote}

Tipos:
\begin{itemize}
    \item \textbf{Perfect ZK:} Distribuciones idénticas
    \item \textbf{Statistical ZK:} Distribuciones estadísticamente cercanas
    \item \textbf{Computational ZK:} Indistinguibles para adversarios eficientes
\end{itemize}

\begin{definition}[Tipos de Zero-Knowledge]
\textbf{Perfect ZK:}
\[ \text{View}[\text{real}] \equiv \mathcal{S}(x) \]
Distribuciones exactamente iguales.

\textbf{Statistical ZK:}
\[ \Delta(\text{View}[\text{real}], \mathcal{S}(x)) \leq \text{negl}(\lambda) \]
Distancia estadística negligible.

\textbf{Computational ZK:}
\[ \text{View}[\text{real}] \approx_c \mathcal{S}(x) \]
Computacionalmente indistinguibles.
\end{definition}

\begin{example}[Simulador para Schnorr]
\textbf{Protocolo real:}
\begin{verbatim}
P: R = g^r (r aleatorio)
V: c aleatorio
P: z = r + c*x
Transcript: (R, c, z)
Verifica: g^z = R * Y^c
\end{verbatim}

\textbf{Simulador (sin conocer x):}
\begin{verbatim}
S: 
    1. Elige z y c aleatorios
    2. Calcula R = g^z / Y^c
    3. Output: (R, c, z)
\end{verbatim}

\textbf{Verificación de que simula correctamente:}
\[ g^z = R \cdot Y^c \iff R = g^z / Y^c \checkmark \]

El transcript simulado tiene la misma distribución que el real:
\begin{itemize}
    \item En ambos casos, $z$ es uniforme
    \item $c$ es uniforme en ambos
    \item $R$ está determinado por $z$ y $c$
\end{itemize}

¡El simulador produce transcripts indistinguibles sin conocer $x$!
\end{example}

\begin{keyidea}
\textbf{¿Por qué el simulador demuestra zero-knowledge?}

El argumento es por contradicción:
\begin{enumerate}
    \item Supón que el verificador $\mathcal{V}^*$ aprende algo sobre $w$
    \item Entonces podría distinguir entre transcripts reales y simulados
    \item Pero el simulador produce transcripts indistinguibles
    \item Por lo tanto, $\mathcal{V}^*$ no puede haber aprendido nada
\end{enumerate}

Cualquier cosa que $\mathcal{V}^*$ puede computar del transcript real, también la puede computar del transcript simulado (que no usó $w$).
\end{keyidea}

%-------------------------------------------
\subsection{Trade-offs entre Propiedades}

\begin{remark}
\textbf{No puedes tener todo:}

\begin{center}
\begin{tabular}{lccc}
\toprule
\textbf{Sistema} & \textbf{Completeness} & \textbf{Soundness} & \textbf{ZK} \\
\midrule
Statistical Proof & Perfect & Statistical & Computational \\
SNARK & Perfect & Computational & Computational \\
STARK & Perfect & Statistical & Computational \\
\bottomrule
\end{tabular}
\end{center}

\textbf{Teorema (informal):} No puede existir un sistema con soundness perfecto, ZK perfecto, y completeness perfecto para lenguajes no triviales.

Por eso la mayoría de sistemas prácticos tienen:
\begin{itemize}
    \item Perfect completeness (lo más fácil)
    \item Computational soundness (para pruebas cortas)
    \item Computational ZK (suficiente en práctica)
\end{itemize}
\end{remark}

\begin{bancozk}
\textbf{Propiedades ZK en Banco ZK:}

Las pruebas de Banco ZK garantizan:

\textbf{Completeness:}
\begin{itemize}
    \item Si tu transacción es válida, siempre se genera una prueba válida
    \item Nunca se rechaza una tx honesta
\end{itemize}

\textbf{Soundness:}
\begin{itemize}
    \item No puedes crear pruebas para txs inválidas (doble gasto, monto negativo)
    \item Basado en la dureza del problema DLP/asunciones de pairings
\end{itemize}

\textbf{Zero-Knowledge:}
\begin{itemize}
    \item La prueba no revela: quién envía, quién recibe, cuánto
    \item Solo revela: ``la tx es válida según las reglas del protocolo''
    \item El verificador (L1) no aprende nada sobre los participantes
\end{itemize}
\end{bancozk}

% ============================================
% BLOQUE 4: ZERO-KNOWLEDGE PROOFS (CORE) - VERSIÓN EXPANDIDA
% Parte 2: Interactive vs NIZK, Fiat-Shamir, Witness
% ============================================

%===========================================
\section{Interactive vs Non-interactive Proofs}
%===========================================

\subsection{Pruebas Interactivas}

Múltiples rondas de comunicación:
\[
\mathcal{P} \leftrightarrow \mathcal{V}: \quad \text{commit} \to \text{challenge} \to \text{response}
\]

\begin{intuicion}
En una prueba interactiva, prover y verifier tienen un ``diálogo'':
\begin{enumerate}
    \item Prover hace un commitment (se ``compromete'' a algo)
    \item Verifier envía un challenge aleatorio
    \item Prover responde al challenge
    \item Verifier decide si acepta
\end{enumerate}

La interacción es crucial porque el challenge debe ser:
\begin{itemize}
    \item Impredecible para el prover (generado después del commitment)
    \item Aleatorio (para que no pueda prepararse)
\end{itemize}

Si el prover pudiera predecir el challenge, podría hacer trampa.
\end{intuicion}

%-------------------------------------------
\subsubsection{Protocolo Sigma ($\Sigma$-protocol)}

Protocolo $\Sigma$ (3 rondas):
\begin{enumerate}
    \item Prover envía commitment $a$
    \item Verifier envía challenge aleatorio $e$
    \item Prover envía response $z$
    \item Verifier comprueba relación
\end{enumerate}

\begin{definition}[$\Sigma$-Protocol]
Un protocolo de 3 rondas $(a, e, z)$ es un $\Sigma$-protocol si satisface:

\textbf{1. Completeness:} Prover honesto siempre convence.

\textbf{2. Special Soundness:} Dados dos transcripts aceptados $(a, e, z)$ y $(a, e', z')$ con $e \neq e'$, se puede extraer el witness.

\textbf{3. Special Honest-Verifier ZK (SHVZK):} Existe simulador para verificador honesto.
\end{definition}

\begin{center}
\begin{tikzpicture}[
    actor/.style={draw, rectangle, minimum width=2cm, minimum height=1cm, align=center},
    arrow/.style={->, thick}
]
    % Actores
    \node[actor, fill=blue!20] (P) at (0, 0) {Prover\\(conoce $w$)};
    \node[actor, fill=green!20] (V) at (8, 0) {Verifier};
    
    % Mensajes
    \draw[arrow] (P.east) -- ++(1.5, 0) |- node[above, pos=0.25] {$a$ (commitment)} ++(2, 1) -- (V.west |- 0, 1);
    \draw[arrow] (V.west) -- ++(-1.5, 0) |- node[above, pos=0.25] {$e$ (challenge)} ++(-2, 0) -- (P.east |- 0, 0);
    \draw[arrow] (P.east) -- ++(1.5, 0) |- node[above, pos=0.25] {$z$ (response)} ++(2, -1) -- (V.west |- 0, -1);
    
    % Timeline
    \node at (0, 2) {Ronda 1};
    \node at (4, 2) {Ronda 2};
    \node at (8, 2) {Ronda 3};
    
    % Decisión
    \node[below=0.5cm of V] {¿Acepta?};
\end{tikzpicture}
\end{center}

\begin{example}[Schnorr como $\Sigma$-protocol]
Para relación $R = \{(Y, x) : Y = g^x\}$:

\begin{center}
\begin{tabular}{ll}
\toprule
\textbf{Prover (conoce $x$)} & \textbf{Verifier} \\
\midrule
$r \xleftarrow{\$} \mathbb{Z}_q$ & \\
$a = g^r$ & \\
& $\xleftarrow{a}$ \\
& $e \xleftarrow{\$} \mathbb{Z}_q$ \\
$\xrightarrow{e}$ & \\
$z = r + e \cdot x \mod q$ & \\
& $\xleftarrow{z}$ \\
& Verificar: $g^z \stackrel{?}{=} a \cdot Y^e$ \\
\bottomrule
\end{tabular}
\end{center}

\textbf{Special Soundness:}

Si tenemos $(a, e, z)$ y $(a, e', z')$ con $e \neq e'$:
\begin{align*}
g^z &= a \cdot Y^e \\
g^{z'} &= a \cdot Y^{e'} \\
\implies g^{z-z'} &= Y^{e-e'} \\
\implies Y &= g^{(z-z')/(e-e')}
\end{align*}

Por lo tanto $x = (z-z')/(e-e')$.
\end{example}

%-------------------------------------------
\subsubsection{Múltiples Rondas}

\begin{mathnote}
\textbf{¿Por qué 3 rondas son comunes?}

\textbf{1 ronda:} El prover envía todo de una vez
\begin{itemize}
    \item No puede ser ZK (verifier no contribuye aleatoriedad)
    \item Excepto en el modelo de random oracle (NIZK)
\end{itemize}

\textbf{3 rondas ($\Sigma$-protocol):}
\begin{itemize}
    \item Mínimo para ZK con verifier honesto
    \item Commit $\to$ Challenge $\to$ Response
    \item Muy eficiente, bien estudiado
\end{itemize}

\textbf{Más rondas:}
\begin{itemize}
    \item Necesario para ZK contra verifier malicioso
    \item O para soundness más fuerte
    \item Generalmente se evita por eficiencia
\end{itemize}

\textbf{Teorema (Goldreich-Kahan):} 4 rondas son suficientes para ZK contra cualquier verifier.
\end{mathnote}

%-------------------------------------------
\subsection{NIZK (Non-Interactive ZK)}

Un solo mensaje del prover al verificador.

Requiere:
\begin{itemize}
    \item \textbf{CRS:} Common Reference String
    \item \textbf{Random Oracle:} Hash modelado como función aleatoria
\end{itemize}

\begin{definition}[NIZK]
Un NIZK para lenguaje $L$ consiste en:

\textbf{Setup:} $\text{crs} \leftarrow \text{Setup}(1^\lambda)$

\textbf{Prove:} $\pi \leftarrow \text{Prove}(\text{crs}, x, w)$

\textbf{Verify:} $b \leftarrow \text{Verify}(\text{crs}, x, \pi)$

Con las propiedades usuales de completeness, soundness, y zero-knowledge.
\end{definition}

\begin{keyidea}
\textbf{¿Por qué necesitamos CRS o Random Oracle?}

En pruebas interactivas, la aleatoriedad del verifier es crucial:
\begin{itemize}
    \item El prover no puede predecirla
    \item Por eso no puede hacer trampa
\end{itemize}

En NIZK, no hay verifier enviando challenges. Necesitamos otra fuente de aleatoriedad impredecible:

\textbf{Opción 1 - CRS (Common Reference String):}
\begin{itemize}
    \item String aleatorio generado una vez
    \item Todos confían en que es aleatorio
    \item El prover no lo eligió, así que no puede explotarlo
\end{itemize}

\textbf{Opción 2 - Random Oracle (Fiat-Shamir):}
\begin{itemize}
    \item Hash modelado como función aleatoria
    \item Challenge = Hash(commitment)
    \item El prover no puede predecir el output del hash
\end{itemize}
\end{keyidea}

%-------------------------------------------
\subsection{Transformación Fiat-Shamir}

Convierte prueba interactiva en no interactiva:
\begin{enumerate}
    \item En protocolo original, verifier envía challenge $e$
    \item En Fiat-Shamir: $e = H(\text{transcript})$
    \item Hash ``simula'' aleatoriedad del verifier
\end{enumerate}

\begin{definition}[Transformación Fiat-Shamir]
Dado un $\Sigma$-protocol $(a, e, z)$, la versión Fiat-Shamir es:

\textbf{Prove:}
\begin{enumerate}
    \item Prover genera commitment $a$
    \item Calcula challenge: $e = H(x \| a)$
    \item Calcula response $z$
    \item Output: $\pi = (a, z)$ (nota: $e$ no se incluye, es derivable)
\end{enumerate}

\textbf{Verify:}
\begin{enumerate}
    \item Recalcula $e = H(x \| a)$
    \item Verifica la relación con $(a, e, z)$
\end{enumerate}
\end{definition}

\begin{example}[Schnorr con Fiat-Shamir]
\textbf{Interactivo:}
\begin{verbatim}
P: a = g^r
V: e ← random
P: z = r + e*x
V: check g^z = a * Y^e
\end{verbatim}

\textbf{Non-interactivo (Fiat-Shamir):}
\begin{verbatim}
P: 
    a = g^r
    e = H(Y || a)           // Challenge derivado del hash
    z = r + e*x
    output π = (a, z)

V:
    e = H(Y || a)           // Recalcula el mismo challenge
    check g^z = a * Y^e
\end{verbatim}
\end{example}

\begin{mathnote}
\textbf{Seguridad de Fiat-Shamir:}

\textbf{En el Random Oracle Model (ROM):}

Si el hash $H$ es un ``random oracle'' (función verdaderamente aleatoria), entonces Fiat-Shamir es seguro:
\begin{itemize}
    \item El prover no puede predecir $H(x \| a)$ antes de elegir $a$
    \item Es como si el verifier eligiera $e$ después de ver $a$
\end{itemize}

\textbf{En la práctica (Standard Model):}

Hashes reales (SHA-256, Keccak) no son random oracles:
\begin{itemize}
    \item Hay ataques teóricos en algunos casos
    \item Pero en la práctica, funciona bien para esquemas diseñados cuidadosamente
    \item Se debe incluir suficiente contexto en el hash: $H(\text{dominio} \| x \| a)$
\end{itemize}

\textbf{Errores comunes:}
\begin{itemize}
    \item No incluir $x$ en el hash $\to$ maleable
    \item No incluir información de dominio $\to$ ataques cross-protocol
\end{itemize}
\end{mathnote}

\begin{advertencia}
\textbf{Weak Fiat-Shamir:}

Un error común es calcular $e = H(a)$ en lugar de $e = H(x \| a)$.

\textbf{Ataque:}
\begin{enumerate}
    \item Atacante quiere probar $x'$ (statement diferente)
    \item Reutiliza la misma prueba $(a, z)$
    \item Challenge $e = H(a)$ es el mismo
    \item Si la verificación solo usa $e$ y no $x$ directamente, acepta
\end{enumerate}

\textbf{Regla:} Siempre incluir el statement completo en el hash.

\begin{verbatim}
// INCORRECTO
e = H(a)

// CORRECTO
e = H(context || statement || a || public_inputs)
\end{verbatim}
\end{advertencia}

%-------------------------------------------
\subsection{Comparación Interactive vs NIZK}

\begin{center}
\begin{tabular}{lcc}
\toprule
\textbf{Aspecto} & \textbf{Interactive} & \textbf{NIZK} \\
\midrule
Comunicación & Múltiples rondas & 1 mensaje \\
Verificador online & Sí & No \\
Reusable & No (fresh challenges) & Sí \\
Publicable & No & Sí (on-chain) \\
Asunción & Ninguna adicional & ROM o CRS \\
Tamaño prueba & Similar & Similar \\
\bottomrule
\end{tabular}
\end{center}

\begin{keyidea}
\textbf{¿Por qué blockchains usan NIZK?}

Las pruebas en blockchain deben ser:
\begin{itemize}
    \item \textbf{No interactivas:} No hay ``verifier'' online
    \item \textbf{Verificables por cualquiera:} Cualquier nodo debe poder verificar
    \item \textbf{Publicables:} Van en la cadena
    \item \textbf{Reusables:} La misma prueba vale para todos
\end{itemize}

Esto hace que NIZK sea esencial para ZK rollups, Zcash, etc.
\end{keyidea}

%===========================================
\section{Witness (Testigo)}
%===========================================

\begin{definition}[Witness]
Información secreta que satisface la relación y permite generar la prueba.
\end{definition}

Para lenguaje NP con relación $R$:
\[
x \in L \iff \exists w : R(x, w) = 1
\]

\begin{itemize}
    \item $x$: statement (público)
    \item $w$: witness (privado)
\end{itemize}

%-------------------------------------------
\subsection{Lenguajes NP y Relaciones}

\begin{mathnote}
\textbf{Conexión con complejidad computacional:}

\textbf{Clase NP:}
Un lenguaje $L$ está en NP si existe:
\begin{itemize}
    \item Relación $R$ verificable en tiempo polinomial
    \item Para cada $x \in L$, existe witness $w$ de tamaño polinomial
    \item Tal que $R(x, w) = 1$
\end{itemize}

\textbf{ZK Proofs para NP:}

El teorema fundamental (Goldreich, Micali, Wigderson, 1991):
\begin{quote}
Todo lenguaje en NP tiene una prueba zero-knowledge.
\end{quote}

Esto significa que cualquier problema cuya solución se pueda verificar eficientemente, también se puede probar en zero-knowledge.
\end{mathnote}

\begin{example}[Witness para diferentes problemas]
\textbf{1. Conocimiento de preimagen:}
\begin{itemize}
    \item Statement: $y$ (el hash)
    \item Witness: $x$ tal que $H(x) = y$
    \item Relación: verificar $H(x) = y$
\end{itemize}

\textbf{2. Conocimiento de clave privada:}
\begin{itemize}
    \item Statement: $Y$ (clave pública)
    \item Witness: $x$ tal que $Y = g^x$
    \item Relación: verificar $g^x = Y$
\end{itemize}

\textbf{3. Transacción válida (simplificado):}
\begin{itemize}
    \item Statement: (old\_root, new\_root, nullifier)
    \item Witness: (secret\_key, amount, merkle\_path, randomness)
    \item Relación: verificar Merkle proof, firma, balance $\geq 0$
\end{itemize}

\textbf{4. Solución de Sudoku:}
\begin{itemize}
    \item Statement: puzzle con algunas celdas llenas
    \item Witness: valores de todas las celdas
    \item Relación: verificar reglas de Sudoku
\end{itemize}

\textbf{5. Grafo 3-coloreable:}
\begin{itemize}
    \item Statement: grafo $G = (V, E)$
    \item Witness: coloración $c: V \to \{1, 2, 3\}$
    \item Relación: $\forall (u,v) \in E: c(u) \neq c(v)$
\end{itemize}
\end{example}

%-------------------------------------------
\subsection{Witness en Circuitos ZK}

\begin{definition}[Witness en Circuitos Aritméticos]
En sistemas como SNARKs, el witness consiste en los valores de todos los ``wires'' internos del circuito.

\textbf{Public inputs:} Valores conocidos por el verifier
\textbf{Private inputs (witness):} Valores conocidos solo por el prover
\textbf{Intermediate wires:} Valores calculados durante la ejecución
\end{definition}

\begin{example}[Witness para $x^3 + x + 5 = y$]
\begin{verbatim}
Circuito:
    w1 = x        // Input privado
    w2 = w1 * w1  // x^2
    w3 = w2 * w1  // x^3
    w4 = w3 + w1  // x^3 + x
    w5 = w4 + 5   // x^3 + x + 5
    
    assert(w5 == y)  // y es público
    
Para x = 3, y = 35:
    Witness completo: (w1=3, w2=9, w3=27, w4=30, w5=35)
    
Public: y = 35
Private: x = 3
Intermediate: w2, w3, w4
\end{verbatim}
\end{example}

%-------------------------------------------
\subsection{Witness Generation}

\begin{keyidea}
\textbf{El proceso de generar el witness:}

\begin{verbatim}
Entrada: 
    - Circuito C
    - Inputs públicos (statement)
    - Inputs privados (secretos del prover)

Proceso:
    1. Asignar inputs a los wires correspondientes
    2. Evaluar el circuito gate por gate
    3. Calcular valor de cada wire intermedio
    4. Verificar que constraints se satisfacen
    
Salida:
    - Witness completo (todos los valores de wires)
    - O error si los inputs no satisfacen el circuito
\end{verbatim}

\textbf{Complejidad:} $O(|C|)$ donde $|C|$ es el tamaño del circuito.

El witness generation es generalmente la parte más rápida. El proving (generar la prueba dado el witness) es lo costoso.
\end{keyidea}

\begin{example}[Witness generation en Circom]
\begin{verbatim}
// Circuito Circom
template Multiplier() {
    signal input a;
    signal input b;
    signal output c;
    
    c <== a * b;
}

// Witness generation (JavaScript)
const input = {
    a: "3",
    b: "5"
};

// snarkjs calcula:
// c = 3 * 5 = 15
// witness = [1, 15, 3, 5]  // [one, c, a, b]
\end{verbatim}
\end{example}

%===========================================
\section{Prover vs Verifier}
%===========================================

\begin{table}[htbp]
\centering
\begin{tabular}{lll}
\toprule
\textbf{Aspecto} & \textbf{Prover} & \textbf{Verifier} \\
\midrule
Input & Statement + Witness & Statement + Proof \\
Cómputo & Costoso ($O(n)$ a $O(n \log n)$) & Barato ($O(1)$ a $O(\log n)$) \\
Memoria & Alta & Baja \\
Rol & Convencer & Verificar \\
\bottomrule
\end{tabular}
\end{table}

%-------------------------------------------
\subsection{El Rol del Prover}

\begin{definition}[Prover]
El prover tiene acceso a:
\begin{itemize}
    \item El statement $x$ (público)
    \item El witness $w$ (secreto)
    \item Los parámetros del sistema (CRS si aplica)
\end{itemize}

Su objetivo: generar una prueba $\pi$ que convenza al verifier de que $x \in L$.
\end{definition}

\begin{mathnote}
\textbf{Complejidad del Prover:}

Para diferentes sistemas:

\begin{center}
\begin{tabular}{lcc}
\toprule
\textbf{Sistema} & \textbf{Tiempo} & \textbf{Memoria} \\
\midrule
Groth16 & $O(n \log n)$ & $O(n)$ \\
PLONK & $O(n \log n)$ & $O(n)$ \\
STARK & $O(n \log^2 n)$ & $O(n)$ \\
Bulletproofs & $O(n)$ & $O(n)$ \\
\bottomrule
\end{tabular}
\end{center}

Donde $n$ es el tamaño del circuito/statement.

\textbf{Operaciones dominantes:}
\begin{itemize}
    \item FFT (Fast Fourier Transform): $O(n \log n)$
    \item MSM (Multi-Scalar Multiplication): $O(n)$
    \item Evaluación de polinomios: $O(n)$
\end{itemize}

El proving es inherentemente costoso porque ``comprime'' un cálculo de tamaño $n$ en una prueba pequeña.
\end{mathnote}

%-------------------------------------------
\subsection{El Rol del Verifier}

\begin{definition}[Verifier]
El verifier tiene acceso a:
\begin{itemize}
    \item El statement $x$ (público)
    \item La prueba $\pi$
    \item Los parámetros del sistema
\end{itemize}

Su objetivo: decidir si acepta o rechaza, idealmente en tiempo sub-lineal en el tamaño del statement original.
\end{definition}

\begin{mathnote}
\textbf{Complejidad del Verifier:}

\begin{center}
\begin{tabular}{lcc}
\toprule
\textbf{Sistema} & \textbf{Tiempo} & \textbf{Tamaño prueba} \\
\midrule
Groth16 & $O(1)$ (3 pairings) & 128 bytes \\
PLONK & $O(1)$ (1-2 pairings) & $\sim$400 bytes \\
STARK & $O(\log^2 n)$ & $\sim$50-200 KB \\
Bulletproofs & $O(n)$ & $O(\log n)$ \\
\bottomrule
\end{tabular}
\end{center}

\textbf{La magia de la succinctness:}

Para un circuito de 1 millón de gates:
\begin{itemize}
    \item Ejecutar el circuito directamente: $\sim$1M operaciones
    \item Verificar Groth16: $\sim$3 operaciones de pairing
    \item Ratio: $\sim$300,000x más rápido
\end{itemize}

Esta asimetría es lo que hace útiles los ZK proofs para blockchain.
\end{mathnote}

%-------------------------------------------
\subsection{Asimetría Prover-Verifier}

\begin{keyidea}
\textbf{¿Por qué la asimetría es buena?}

En blockchain:
\begin{itemize}
    \item \textbf{Provers:} Pocos, tienen recursos (servidores, GPUs)
    \item \textbf{Verifiers:} Muchos (todos los nodos), deben ser eficientes
\end{itemize}

\textbf{Ejemplo ZK Rollup:}
\begin{itemize}
    \item 1 prover genera prueba para 10,000 txs: 5 minutos, \$10 de compute
    \item 10,000 nodos verifican: 10ms cada uno, prácticamente gratis
\end{itemize}

Si la verificación fuera $O(n)$, los nodos tendrían que re-ejecutar todo, perdiendo el beneficio de escalabilidad.

\textbf{Regla general:}
\begin{itemize}
    \item Está bien que proving sea caro (pocos provers)
    \item Verificación DEBE ser barata (muchos verificadores)
    \item El costo de proving se amortiza entre todos los usuarios
\end{itemize}
\end{keyidea}

\begin{center}
\begin{tikzpicture}
    % Prover
    \node[draw, rectangle, minimum width=3cm, minimum height=2cm, fill=blue!20] (prover) at (0, 0) {
        \begin{tabular}{c}
        \textbf{Prover} \\
        CPU/GPU potente \\
        Minutos-horas \\
        Alto costo
        \end{tabular}
    };
    
    % Proof
    \node[draw, ellipse, fill=yellow!30] (proof) at (5, 0) {Prueba pequeña};
    
    % Verifiers
    \node[draw, rectangle, minimum width=3cm, minimum height=2cm, fill=green!20] (verifier) at (10, 0) {
        \begin{tabular}{c}
        \textbf{Verifiers} \\
        Cualquier nodo \\
        Milisegundos \\
        Casi gratis
        \end{tabular}
    };
    
    % Flechas
    \draw[->, thick] (prover) -- node[above] {genera} (proof);
    \draw[->, thick] (proof) -- node[above] {verifica} (verifier);
    
    % Witness
    \node[above=0.5cm of prover] {+ witness (secreto)};
    
    % Statement
    \node[below=0.5cm of prover] {statement (público)};
    \node[below=0.5cm of verifier] {statement (público)};
\end{tikzpicture}
\end{center}

\begin{bancozk}
\textbf{Prover y Verifier en Banco ZK:}

\textbf{Arquitectura de proving:}
\begin{verbatim}
Usuario genera tx privada
    ↓
Usuario genera witness local
(clave, monto, Merkle path, randomness)
    ↓
Prover (puede ser usuario o servidor dedicado)
genera ZK proof de la tx individual
    ↓
Sequencer agrega txs en batch
    ↓
Prover central genera proof del batch completo
(puede usar recursión para combinar proofs)
    ↓
Proof publicado en L1
    ↓
Contrato L1 verifica (3 pairings, ~200K gas)
\end{verbatim}

\textbf{Trade-offs:}
\begin{itemize}
    \item Proving en cliente: máxima privacidad, puede ser lento
    \item Proving en servidor: más rápido, pero servidor ve el witness
    \item Híbrido: proof parcial en cliente, agregación en servidor
\end{itemize}
\end{bancozk}

% ============================================
% BLOQUE 4: ZERO-KNOWLEDGE PROOFS (CORE) - VERSIÓN EXPANDIDA
% Parte 3: Trusted Setup vs Transparent, CRS
% ============================================

%===========================================
\section{Trusted Setup vs Transparent Setup}
%===========================================

\subsection{Trusted Setup}

\begin{definition}[Trusted Setup]
Ceremonia para generar parámetros públicos. Se crea ``toxic waste'' que debe destruirse.
\end{definition}

\begin{itemize}
    \item Si alguien tiene toxic waste: puede crear pruebas falsas
    \item Mitigación: MPC ceremony (basta 1 participante honesto)
    \item Powers of Tau: ceremonia reutilizable
\end{itemize}

Sistemas: Groth16, PLONK (KZG).

%-------------------------------------------
\subsubsection{¿Qué es el Toxic Waste?}

\begin{intuicion}
El trusted setup genera parámetros públicos que ``codifican'' un secreto $\tau$ (tau). Este secreto es el ``toxic waste''.

Imagina que los parámetros son una caja fuerte:
\begin{itemize}
    \item La caja fuerte está diseñada para que solo quien conoce un código específico pueda abrirla
    \item El código se destruye después de crear la caja
    \item Si alguien conservó el código, puede abrir cualquier caja de ese tipo
\end{itemize}

En ZK proofs:
\begin{itemize}
    \item Los parámetros permiten crear pruebas válidas
    \item Pero quien conoce $\tau$ puede crear pruebas de statements FALSOS
    \item Por eso $\tau$ debe ser destruido irrecuperablemente
\end{itemize}
\end{intuicion}

\begin{mathnote}
\textbf{El toxic waste en KZG/Groth16:}

\textbf{Setup:}

Se elige $\tau \xleftarrow{\$} \mathbb{Z}_p$ aleatorio (el toxic waste).

Se generan los parámetros públicos (SRS - Structured Reference String):
\[ \text{SRS} = \{g, g^\tau, g^{\tau^2}, \ldots, g^{\tau^d}\} \]

Donde $g$ es generador del grupo.

\textbf{Uso legítimo:}

Para commitir a un polinomio $p(x) = \sum_{i=0}^{d} a_i x^i$:
\[ C = \prod_{i=0}^{d} (g^{\tau^i})^{a_i} = g^{\sum a_i \tau^i} = g^{p(\tau)} \]

El prover no conoce $\tau$, solo puede evaluar $g^{p(\tau)}$ usando el SRS.

\textbf{Ataque si conoces $\tau$:}

Si conoces $\tau$, puedes:
\begin{enumerate}
    \item Crear commitment a cualquier polinomio
    \item Crear prueba de evaluación para cualquier punto
    \item Incluso para relaciones que no son verdad
\end{enumerate}

El sistema de prueba asume que nadie conoce $\tau$.
\end{mathnote}

%-------------------------------------------
\subsubsection{MPC Ceremonies (Multi-Party Computation)}

\begin{definition}[Ceremonia MPC]
Protocolo donde múltiples participantes contribuyen aleatoriedad al setup, de forma que:
\begin{itemize}
    \item Nadie conoce el secreto completo
    \item Basta que 1 participante sea honesto para que el setup sea seguro
\end{itemize}
\end{definition}

\begin{keyidea}
\textbf{Cómo funciona una ceremonia MPC:}

\textbf{Idea básica:}

En lugar de un $\tau$ único, se construye como producto de contribuciones:
\[ \tau = \tau_1 \cdot \tau_2 \cdot \tau_3 \cdots \tau_n \]

Cada participante $i$:
\begin{enumerate}
    \item Recibe los parámetros del participante anterior
    \item Elige su $\tau_i$ aleatorio
    \item ``Multiplica'' los parámetros por $\tau_i$
    \item Destruye $\tau_i$
    \item Pasa los nuevos parámetros al siguiente
\end{enumerate}

\textbf{Seguridad:}

Para conocer $\tau$ completo, necesitas TODOS los $\tau_i$.

Si al menos un participante:
\begin{itemize}
    \item Generó su $\tau_i$ aleatoriamente
    \item Lo destruyó correctamente
\end{itemize}

Entonces $\tau$ es desconocido para todos, incluyendo atacantes.

\textbf{Trust assumption:} ``1-of-n honest'' (al menos 1 de n es honesto)
\end{keyidea}

\begin{example}[Ceremonia simplificada]
\begin{verbatim}
Inicio: params_0 = (g, g^1) = (g, g)  // "tau = 1" inicial

Participante 1:
    tau_1 = random()
    params_1 = (g, g^tau_1)
    delete tau_1

Participante 2:
    tau_2 = random()
    params_2 = (g, (g^tau_1)^tau_2) = (g, g^(tau_1*tau_2))
    delete tau_2

Participante 3:
    tau_3 = random()
    params_3 = (g, g^(tau_1*tau_2*tau_3))
    delete tau_3

Final:
    tau_efectivo = tau_1 * tau_2 * tau_3
    
    Para conocer tau_efectivo necesitas tau_1 AND tau_2 AND tau_3
    Si cualquiera es desconocido, tau_efectivo es seguro
\end{verbatim}
\end{example}

%-------------------------------------------
\subsubsection{Ceremonias Históricas}

\begin{remark}
\textbf{Ceremonias famosas:}

\textbf{Zcash Sprout (2016):}
\begin{itemize}
    \item 6 participantes
    \item Protocolo muy seguro pero complejo
    \item Participantes destruyeron hardware físicamente
    \item Críticas: pocos participantes
\end{itemize}

\textbf{Zcash Powers of Tau (2017-2018):}
\begin{itemize}
    \item 87 participantes públicos
    \item Cualquiera podía participar
    \item ``Perpetual Powers of Tau'' - sigue abierta
\end{itemize}

\textbf{Ethereum KZG Ceremony (2023):}
\begin{itemize}
    \item Para EIP-4844 (blobs)
    \item 141,000+ participantes
    \item Récord de participación
    \item Interfaz web fácil de usar
    \item Probabilidad de que todos sean deshonestos: $\approx 0$
\end{itemize}

\begin{center}
\begin{tabular}{lccc}
\toprule
\textbf{Ceremonia} & \textbf{Participantes} & \textbf{Año} & \textbf{Propósito} \\
\midrule
Zcash Sprout & 6 & 2016 & Zcash shielded \\
Powers of Tau & 87+ & 2017 & Universal \\
Aztec Ignition & 200+ & 2020 & PLONK \\
Hermez & 100+ & 2021 & zkRollup \\
Ethereum KZG & 141,000+ & 2023 & EIP-4844 \\
\bottomrule
\end{tabular}
\end{center}
\end{remark}

%-------------------------------------------
\subsubsection{Per-Circuit vs Universal Setup}

\begin{definition}[Tipos de Trusted Setup]
\textbf{Per-circuit (Groth16):}
\begin{itemize}
    \item Setup específico para cada circuito
    \item Si cambias el circuito, nueva ceremonia
    \item Pruebas más pequeñas
\end{itemize}

\textbf{Universal (PLONK, Marlin):}
\begin{itemize}
    \item Setup para un tamaño máximo de circuito
    \item Cualquier circuito $\leq$ tamaño max funciona
    \item Una ceremonia para muchos circuitos
\end{itemize}

\textbf{Updatable (PLONK):}
\begin{itemize}
    \item Cualquiera puede ``actualizar'' el setup
    \item Añade su propia aleatoriedad
    \item Setup se vuelve más seguro con el tiempo
\end{itemize}
\end{definition}

%-------------------------------------------
\subsection{Transparent Setup}

\begin{definition}[Transparent Setup]
Parámetros derivados de valores públicos, sin secretos.
\end{definition}

\begin{itemize}
    \item ``Nothing up my sleeve'': semillas públicas
    \item Sin confianza en ceremonia
    \item Pruebas típicamente más grandes
\end{itemize}

Sistemas: STARKs, Bulletproofs, Halo (IPA).

%-------------------------------------------
\subsubsection{¿Cómo Funciona Transparent Setup?}

\begin{keyidea}
\textbf{Generación de parámetros sin secretos:}

\textbf{Método 1: Hash-to-curve (``nothing up my sleeve''):}
\begin{verbatim}
// Generar puntos de curva sin conocer logaritmo discreto
seed = "Banco ZK Genesis 2024"
for i in 0..n:
    G_i = hash_to_curve(seed || i)
    
// Nadie conoce x tal que G_i = x * G_0
// Los puntos son "aleatorios" pero determinísticos
\end{verbatim}

\textbf{Método 2: Usar constantes públicas:}
\begin{itemize}
    \item Dígitos de $\pi$: ``3.14159...''
    \item Dígitos de $e$: ``2.71828...''
    \item Hash de algo público (headline de periódico)
\end{itemize}

\textbf{Verificabilidad:}

Cualquiera puede regenerar los parámetros:
\begin{itemize}
    \item Ejecutar el mismo algoritmo
    \item Verificar que coincide con los parámetros publicados
    \item No hay ``ceremonia'' de la que depender
\end{itemize}
\end{keyidea}

\begin{example}[Generadores en Bulletproofs]
\begin{verbatim}
// Bulletproofs necesita generadores G_i, H_i

def generate_generators(n):
    generators_g = []
    generators_h = []
    
    for i in range(n):
        // Hash to curve con dominio separado
        g_i = hash_to_curve(b"bulletproofs_g" + i.to_bytes())
        h_i = hash_to_curve(b"bulletproofs_h" + i.to_bytes())
        
        generators_g.append(g_i)
        generators_h.append(h_i)
    
    return generators_g, generators_h

// Propiedades:
// - Determinístico: siempre genera los mismos puntos
// - Verificable: cualquiera puede regenerar
// - Sin trapdoor: nadie conoce log_g(g_i) o log_g(h_i)
\end{verbatim}
\end{example}

%-------------------------------------------
\subsubsection{Trade-offs: Trusted vs Transparent}

\begin{center}
\begin{tabular}{lcc}
\toprule
\textbf{Aspecto} & \textbf{Trusted Setup} & \textbf{Transparent} \\
\midrule
Confianza & 1-of-n honesto en ceremonia & Ninguna \\
Tamaño prueba & Pequeño ($\sim$100-500 B) & Grande ($\sim$10-200 KB) \\
Verificación & Muy rápida (pairings) & Más lenta \\
Post-quantum & No* & Puede ser sí \\
Complejidad setup & Ceremonia MPC & Determinístico \\
Actualizar circuito & Nueva ceremonia* & Sin cambios \\
\bottomrule
\end{tabular}
\end{center}

*Universal setup mitiga estos problemas.

\begin{advertencia}
\textbf{¿Es tan malo el trusted setup?}

\textbf{Argumentos en contra:}
\begin{itemize}
    \item Si la ceremonia falla, todo el sistema es inseguro
    \item Difícil de auditar
    \item ``Confianza'' parece contrario al ethos de crypto
\end{itemize}

\textbf{Argumentos a favor:}
\begin{itemize}
    \item Con 141,000+ participantes, la probabilidad de fallo es $\approx 0$
    \item Las pruebas son mucho más pequeñas y baratas de verificar
    \item En la práctica, el riesgo es mínimo
    \item El setup se hace una vez; los beneficios son permanentes
\end{itemize}

\textbf{Tendencia actual:}

La comunidad se mueve hacia transparent setups cuando es práctico, pero trusted setup con muchos participantes sigue siendo aceptable para producción.
\end{advertencia}

%===========================================
\section{CRS (Common Reference String)}
%===========================================

\begin{definition}[CRS]
Parámetros públicos compartidos entre prover y verifier.
\end{definition}

Tipos:
\begin{itemize}
    \item \textbf{Structured:} Codifica información sobre $\tau$ (KZG)
    \item \textbf{Uniform:} Aleatorio puro
    \item \textbf{Per-circuit:} Específico para cada circuito (Groth16)
    \item \textbf{Universal:} Funciona para cualquier circuito hasta cierto tamaño (PLONK)
    \item \textbf{Updatable:} Cualquiera puede contribuir
\end{itemize}

%-------------------------------------------
\subsection{Estructura del CRS}

\begin{mathnote}
\textbf{CRS para diferentes sistemas:}

\textbf{Groth16 (per-circuit):}
\begin{verbatim}
CRS = (
    // Proving key (para prover)
    pk = {
        [α]_1, [β]_1, [β]_2, [δ]_1, [δ]_2,
        [τ^i]_1 para i = 0..n,
        [α·A_j(τ) + β·B_j(τ)]_1 para cada wire j,
        ...
    },
    
    // Verification key (para verifier)
    vk = {
        [α]_1, [β]_2, [γ]_2, [δ]_2,
        [γ^(-1)·(β·A_i(τ) + α·B_i(τ) + C_i(τ))]_1
    }
)
\end{verbatim}

El CRS codifica el circuito específico. Cambiar el circuito = nuevo CRS.

\textbf{PLONK (universal):}
\begin{verbatim}
CRS = {
    // Powers of tau en G1
    [1]_1, [τ]_1, [τ^2]_1, ..., [τ^n]_1,
    
    // Powers of tau en G2 (menos)
    [1]_2, [τ]_2
}
\end{verbatim}

El CRS solo depende del tamaño máximo $n$. Cualquier circuito de tamaño $\leq n$ puede usar este CRS.

\textbf{STARKs (transparent):}
\begin{verbatim}
// No hay CRS estructurado
// Solo parámetros públicos:
params = {
    field: F_p,
    hash: Poseidon/Keccak,
    expansion_factor: 8,
    num_queries: 30
}
\end{verbatim}

Cualquiera puede verificar que estos son los parámetros correctos.
\end{mathnote}

%-------------------------------------------
\subsection{Proving Key vs Verification Key}

\begin{definition}[Keys en ZK Systems]
\textbf{Proving Key (pk):}
\begin{itemize}
    \item Usado por el prover
    \item Típicamente grande (MB a GB)
    \item Contiene información para generar pruebas
    \item Secreto del prover? No, es público
\end{itemize}

\textbf{Verification Key (vk):}
\begin{itemize}
    \item Usado por el verifier
    \item Pequeño (cientos de bytes a KB)
    \item Contiene información para verificar pruebas
    \item Público, puede ir on-chain
\end{itemize}
\end{definition}

\begin{example}[Tamaños típicos]
Para un circuito de 1 millón de constraints:

\begin{center}
\begin{tabular}{lcc}
\toprule
\textbf{Sistema} & \textbf{Proving Key} & \textbf{Verification Key} \\
\midrule
Groth16 & $\sim$1 GB & $\sim$1 KB \\
PLONK & $\sim$500 MB & $\sim$2 KB \\
STARK & N/A (transparent) & $\sim$100 bytes \\
\bottomrule
\end{tabular}
\end{center}

La proving key grande es un trade-off por pruebas pequeñas.
\end{example}

%-------------------------------------------
\subsection{Derivación de CRS Específico de Circuito}

\begin{keyidea}
\textbf{De CRS universal a CRS de circuito:}

Con PLONK universal:
\begin{verbatim}
// CRS universal (de la ceremonia)
srs = {[τ^i]_1 para i = 0..N_max}

// Circuito específico
circuit = compile("mi_programa.circom")

// Derivar proving/verification keys
(pk, vk) = derive_keys(srs, circuit)

// pk y vk son específicos del circuito
// pero srs es reutilizable
\end{verbatim}

\textbf{Proceso de derivación:}
\begin{enumerate}
    \item Tomar el SRS universal
    \item Computar polinomios que describen el circuito
    \item ``Evaluar'' esos polinomios en el SRS
    \item Resultado: keys específicas del circuito
\end{enumerate}

No se necesita nueva ceremonia, solo cómputo local.
\end{keyidea}

\begin{bancozk}
\textbf{Setup en Banco ZK:}

\textbf{Nuestra elección: PLONK con KZG}

Razones:
\begin{itemize}
    \item Setup universal: un setup para todos nuestros circuitos
    \item Pruebas pequeñas ($\sim$500 bytes)
    \item Verificación barata ($\sim$200K gas en Ethereum)
    \item Podemos actualizar circuitos sin nueva ceremonia
\end{itemize}

\textbf{SRS utilizado:}

Usamos el SRS de la Ethereum KZG Ceremony:
\begin{itemize}
    \item 141,000+ participantes
    \item Públicamente verificable
    \item Soporta circuitos hasta $2^{28}$ constraints
\end{itemize}

\textbf{Keys derivadas:}

\begin{verbatim}
// Circuitos de Banco ZK
circuits = {
    "transfer": circuito para transferencias privadas,
    "deposit": circuito para depósitos desde L1,
    "withdraw": circuito para retiros a L1,
    "batch": circuito para agregar múltiples txs
}

// Derivar keys para cada circuito
for circuit in circuits:
    (pk, vk) = derive_keys(ethereum_srs, circuit)
    
// vk va al contrato verificador en L1
// pk se distribuye a los provers
\end{verbatim}

\textbf{Consideración post-quantum:}

KZG/PLONK no es post-quantum. Nuestro roadmap incluye:
\begin{itemize}
    \item Corto plazo: PLONK+KZG (actual)
    \item Medio plazo: Evaluar Plonky3/STARK híbrido
    \item Largo plazo: Migrar a STARKs si amenaza cuántica se materializa
\end{itemize}
\end{bancozk}

% ============================================
% BLOQUE 4: ZERO-KNOWLEDGE PROOFS (CORE) - VERSIÓN EXPANDIDA
% Parte 4: Sistemas de Prueba
% ============================================

%===========================================
\section{Sistemas de Prueba}
%===========================================

\begin{intuicion}
Un ``sistema de prueba'' es un conjunto de algoritmos que implementan las primitivas ZK:
\begin{itemize}
    \item \textbf{Setup:} Generar parámetros
    \item \textbf{Prove:} Crear prueba dado witness
    \item \textbf{Verify:} Verificar prueba
\end{itemize}

Diferentes sistemas hacen diferentes trade-offs entre:
\begin{itemize}
    \item Tamaño de prueba
    \item Tiempo de proving
    \item Tiempo de verificación
    \item Necesidad de trusted setup
    \item Seguridad post-quantum
    \item Asunciones criptográficas
\end{itemize}
\end{intuicion}

%-------------------------------------------
\subsection{SNARKs}

\begin{definition}[SNARK]
Succinct Non-interactive ARgument of Knowledge.
\end{definition}

\begin{itemize}
    \item \textbf{Succinct:} Prueba pequeña ($O(1)$), verificación rápida
    \item \textbf{Non-interactive:} Un mensaje
    \item \textbf{ARgument:} Seguridad computacional
    \item \textbf{of Knowledge:} Extractor existe
\end{itemize}

\begin{mathnote}
\textbf{Desglose del acrónimo:}

\textbf{S - Succinct (Sucinto):}
\begin{itemize}
    \item Tamaño de prueba: $O(1)$ o $O(\log n)$
    \item Tiempo de verificación: $O(1)$ o $O(\log n)$
    \item Independiente (o logarítmico) del tamaño del statement
\end{itemize}

\textbf{N - Non-interactive (No interactivo):}
\begin{itemize}
    \item Una sola comunicación: prover $\to$ verifier
    \item No hay rondas de challenge-response
    \item Publicable, verificable offline
\end{itemize}

\textbf{AR - ARgument (Argumento):}
\begin{itemize}
    \item Soundness computacional (no estadístico)
    \item Seguro contra adversarios de tiempo polinomial
    \item Permite pruebas más pequeñas que ``proofs'' estadísticos
\end{itemize}

\textbf{K - of Knowledge (de Conocimiento):}
\begin{itemize}
    \item Si el prover convence, ``conoce'' un witness
    \item Existe extractor que puede obtener el witness
    \item Más fuerte que soundness básico
\end{itemize}
\end{mathnote}

\begin{remark}
\textbf{Historia de los SNARKs:}

\begin{center}
\begin{tabular}{lll}
\toprule
\textbf{Año} & \textbf{Sistema} & \textbf{Contribución} \\
\midrule
2010 & Groth & Primeros SNARKs eficientes \\
2013 & Pinocchio & Primer SNARK práctico \\
2013 & GGPR & Base teórica para muchos SNARKs \\
2016 & Groth16 & El más eficiente (aún usado) \\
2018 & Sonic & Primer universal SNARK \\
2019 & PLONK & Universal + custom gates \\
2019 & Marlin & Universal con mejoras \\
2020 & Halo & Sin trusted setup (IPA) \\
2022 & Plonky2 & PLONK + FRI, muy rápido \\
\bottomrule
\end{tabular}
\end{center}
\end{remark}

%-------------------------------------------
\subsection{STARKs}

\begin{definition}[STARK]
Scalable Transparent ARgument of Knowledge.
\end{definition}

\begin{itemize}
    \item \textbf{Scalable:} Prover quasi-lineal, verifier logarítmico
    \item \textbf{Transparent:} Sin trusted setup
    \item Basado en hashes (post-quantum)
    \item Pruebas más grandes ($O(\log^2 n)$)
\end{itemize}

\begin{mathnote}
\textbf{Diferencias clave SNARK vs STARK:}

\begin{center}
\begin{tabular}{lcc}
\toprule
\textbf{Propiedad} & \textbf{SNARK} & \textbf{STARK} \\
\midrule
Basado en & Curvas elípticas & Hashes \\
Trusted setup & Generalmente sí & No \\
Post-quantum & No & Sí \\
Tamaño prueba & $O(1)$ $\sim$100 B & $O(\log^2 n)$ $\sim$100 KB \\
Verify time & $O(1)$ ms & $O(\log^2 n)$ ms \\
Prove time & $O(n \log n)$ & $O(n \log^2 n)$ \\
Asunciones & DLP, Pairings & Collision-resistant hash \\
\bottomrule
\end{tabular}
\end{center}
\end{mathnote}

\begin{keyidea}
\textbf{¿Cómo funcionan los STARKs?}

\textbf{Idea central: AIR (Algebraic Intermediate Representation)}

\begin{enumerate}
    \item Representar cómputo como tabla de ``trace''
    \item Cada fila es un estado
    \item Constraints relacionan filas consecutivas
    \item Interpolar trace como polinomios
    \item Probar que polinomios satisfacen constraints
    \item Usar FRI para probar bajo grado
\end{enumerate}

\textbf{Ejemplo simplificado:}

Probar que calculé Fibonacci correctamente:
\begin{verbatim}
Trace (10 pasos):
| step | a  | b  |
|------|----|----|
| 0    | 1  | 1  |
| 1    | 1  | 2  |
| 2    | 2  | 3  |
| 3    | 3  | 5  |
| ...  | ...| ...|
| 9    | 34 | 55 |

Constraints:
- a_{i+1} = b_i
- b_{i+1} = a_i + b_i
- a_0 = 1, b_0 = 1

Estos constraints se convierten en ecuaciones polinomiales
que deben ser cero en todos los puntos del trace.
\end{verbatim}

\textbf{FRI (Fast Reed-Solomon IOP):}
\begin{itemize}
    \item Protocolo para probar que un polinomio tiene grado bajo
    \item Basado en códigos Reed-Solomon
    \item Solo usa hashes (no curvas)
    \item Pruebas logarítmicas
\end{itemize}
\end{keyidea}

\begin{remark}
\textbf{¿Quién usa STARKs?}

\begin{itemize}
    \item \textbf{StarkNet:} L2 de StarkWare
    \item \textbf{StarkEx:} Engine para dYdX, Immutable X
    \item \textbf{Polygon Miden:} VM basada en STARKs
    \item \textbf{Winterfell:} Librería de Facebook/Meta
\end{itemize}

Cairo es el lenguaje de StarkWare para escribir programas STARK-provable.
\end{remark}

%-------------------------------------------
\subsection{Bulletproofs}

\begin{itemize}
    \item Sin trusted setup
    \item Prueba logarítmica: $O(\log n)$
    \item Verificación lineal: $O(n)$
    \item Ideal para range proofs
    \item Usado en Monero
\end{itemize}

\begin{definition}[Bulletproofs]
Sistema de prueba basado en Inner Product Arguments (IPA).

\textbf{Componentes:}
\begin{enumerate}
    \item Pedersen commitments para comprometer valores
    \item Inner product argument para probar relaciones
    \item Range proofs como aplicación principal
\end{enumerate}
\end{definition}

\begin{mathnote}
\textbf{Inner Product Argument:}

Probar que $\langle \mathbf{a}, \mathbf{b} \rangle = c$ dado:
\begin{itemize}
    \item Commitment a $\mathbf{a}$: $A = \langle \mathbf{a}, \mathbf{G} \rangle$
    \item Commitment a $\mathbf{b}$: $B = \langle \mathbf{b}, \mathbf{H} \rangle$
    \item El producto interno $c$
\end{itemize}

\textbf{Protocolo (simplificado):}
\begin{enumerate}
    \item Split $\mathbf{a}, \mathbf{b}$ en mitades: $\mathbf{a}_L, \mathbf{a}_R, \mathbf{b}_L, \mathbf{b}_R$
    \item Calcular cross-terms $L, R$
    \item Verifier envía challenge $x$
    \item Fold: $\mathbf{a}' = \mathbf{a}_L + x \cdot \mathbf{a}_R$
    \item Repetir hasta $n = 1$
\end{enumerate}

\textbf{Resultado:}
\begin{itemize}
    \item Prueba de tamaño $O(\log n)$: $2 \log n$ puntos de curva
    \item Pero verificación es $O(n)$: necesita recomputar los folds
\end{itemize}
\end{mathnote}

\begin{example}[Range Proof con Bulletproofs]
Probar que $v \in [0, 2^{64})$ sin revelar $v$:

\begin{verbatim}
// Commitment al valor
V = g^v * h^r  // Pedersen commitment

// Descomponer v en bits
v = sum_{i=0}^{63} v_i * 2^i

// Probar:
// 1. Cada v_i es 0 o 1 (constraints cuadráticos)
// 2. Los bits reconstruyen v
// 3. Todo en ZK

Tamaño de prueba: 64 bytes base + 2*log2(64) puntos
                = 64 + 2*6*32 = ~448 bytes

(vs ~2KB para prueba naive bit por bit)
\end{verbatim}
\end{example}

\begin{remark}
\textbf{Bulletproofs en Monero:}

Antes de Bulletproofs, Monero usaba range proofs que escalaban linealmente. Una transacción con múltiples outputs era muy grande.

Con Bulletproofs:
\begin{itemize}
    \item Tamaño de tx reducido $\sim$80\%
    \item Fees reducidos proporcionalmente
    \item Verificación sigue siendo $O(n)$ pero más eficiente
\end{itemize}

Bulletproofs+ (mejora): verificación batched más eficiente.
\end{remark}

%-------------------------------------------
\subsection{PLONK}

\begin{itemize}
    \item Setup universal (una ceremonia para todos los circuitos)
    \item Prueba: $\sim$400 bytes
    \item Custom gates y lookups
    \item Variantes: TurboPLONK, UltraPLONK, HyperPLONK
\end{itemize}

\begin{definition}[PLONK]
Permutations over Lagrange-bases for Oecumenical Noninteractive arguments of Knowledge.

\textbf{Innovaciones:}
\begin{enumerate}
    \item \textbf{Universal setup:} Un SRS para cualquier circuito
    \item \textbf{Custom gates:} Más que solo multiplicación
    \item \textbf{Permutation argument:} Para copy constraints
    \item \textbf{Grand product:} Verificación eficiente
\end{enumerate}
\end{definition}

\begin{mathnote}
\textbf{Arithmetización PLONKish:}

PLONK representa circuitos como sistema de ecuaciones sobre columnas:

\textbf{Columnas básicas:}
\begin{itemize}
    \item $a, b, c$: valores de los wires
    \item $q_L, q_R, q_O, q_M, q_C$: selectores (constantes)
\end{itemize}

\textbf{Gate equation:}
\[ q_L \cdot a + q_R \cdot b + q_O \cdot c + q_M \cdot a \cdot b + q_C = 0 \]

\textbf{Ejemplos:}

\textit{Suma:} $a + b = c$
\[ q_L=1, q_R=1, q_O=-1, q_M=0, q_C=0 \]

\textit{Multiplicación:} $a \cdot b = c$
\[ q_L=0, q_R=0, q_O=-1, q_M=1, q_C=0 \]

\textit{Constante:} $a = 5$
\[ q_L=1, q_R=0, q_O=0, q_M=0, q_C=-5 \]

\textbf{Copy constraints:}

Para conectar wires entre gates (ej: output de gate 1 = input de gate 2), PLONK usa un \textit{permutation argument}.
\end{mathnote}

\begin{keyidea}
\textbf{Variantes de PLONK:}

\textbf{Vanilla PLONK:}
\begin{itemize}
    \item Gate equation básica
    \item 5 selectores
    \item Suficiente para circuitos generales
\end{itemize}

\textbf{TurboPLONK:}
\begin{itemize}
    \item Custom gates adicionales
    \item Más columnas de wires
    \item Más eficiente para operaciones específicas
\end{itemize}

\textbf{UltraPLONK:}
\begin{itemize}
    \item + Lookup tables (Plookup)
    \item Muy útil para operaciones no-algebraicas (XOR, ranges)
    \item Usado en Aztec
\end{itemize}

\textbf{HyperPLONK:}
\begin{itemize}
    \item Basado en polinomios multilineales
    \item Sumcheck protocol
    \item Potencialmente más eficiente
\end{itemize}
\end{keyidea}

%-------------------------------------------
\subsection{Groth16}

\begin{itemize}
    \item Prueba más pequeña: 128 bytes
    \item Verificación más rápida: 3 pairings
    \item Trusted setup por circuito
    \item Estándar en Ethereum (precompilados)
\end{itemize}

\begin{definition}[Groth16]
El sistema de prueba más eficiente en términos de tamaño de prueba y verificación.

\textbf{Estructura de la prueba:}
\[ \pi = ([A]_1, [B]_2, [C]_1) \]

Solo 3 elementos de grupo: 2 en $\mathbb{G}_1$, 1 en $\mathbb{G}_2$.

\textbf{Verificación:}
\[ e(A, B) = e(\alpha, \beta) \cdot e(\sum_{i \in \text{pub}} x_i \cdot L_i, \gamma) \cdot e(C, \delta) \]

Solo 3 pairings (más uno precalculado).
\end{definition}

\begin{mathnote}
\textbf{Tamaños exactos en Groth16:}

Para curva BN254:
\begin{itemize}
    \item $[A]_1$: 64 bytes (punto en $\mathbb{G}_1$)
    \item $[B]_2$: 128 bytes (punto en $\mathbb{G}_2$)
    \item $[C]_1$: 64 bytes (punto en $\mathbb{G}_1$)
    \item Total: 256 bytes sin comprimir
    \item Comprimido: 128 bytes (solo coordenada $x$ + signo)
\end{itemize}

Para BLS12-381:
\begin{itemize}
    \item $[A]_1$: 48 bytes comprimido
    \item $[B]_2$: 96 bytes comprimido
    \item $[C]_1$: 48 bytes comprimido
    \item Total: 192 bytes comprimido
\end{itemize}
\end{mathnote}

\begin{example}[Verificación Groth16 en Solidity]
\begin{verbatim}
// Contrato verificador Groth16
contract Groth16Verifier {
    // Puntos del verification key
    uint256[2] alpha1;
    uint256[2][2] beta2;
    uint256[2][2] gamma2;
    uint256[2][2] delta2;
    uint256[2][] IC;  // Para public inputs
    
    function verify(
        uint256[2] memory a,    // A en G1
        uint256[2][2] memory b, // B en G2
        uint256[2] memory c,    // C en G1
        uint256[] memory input  // Public inputs
    ) public view returns (bool) {
        
        // Calcular vk_x = IC[0] + sum(input[i] * IC[i+1])
        uint256[2] memory vk_x = IC[0];
        for (uint i = 0; i < input.length; i++) {
            vk_x = addition(vk_x, scalar_mul(IC[i+1], input[i]));
        }
        
        // Verificar pairing equation
        // e(A, B) == e(alpha, beta) * e(vk_x, gamma) * e(C, delta)
        return pairing(
            negate(a), b,           // -A, B
            alpha1, beta2,          // alpha, beta
            vk_x, gamma2,           // vk_x, gamma
            c, delta2               // C, delta
        );
    }
}
\end{verbatim}

Costo en gas: $\sim$200-250K gas.
\end{example}

%-------------------------------------------
\subsection{Halo2}

\begin{itemize}
    \item PLONKish arithmetization
    \item Sin trusted setup (usando IPA)
    \item Custom gates, lookups
    \item Recursión eficiente
    \item Usado en Zcash Orchard, Scroll
\end{itemize}

\begin{definition}[Halo2]
Sistema desarrollado por Electric Coin Company (Zcash) que combina:
\begin{enumerate}
    \item Arithmetización PLONK
    \item Polynomial commitment con IPA (sin trusted setup)
    \item Técnica de ``accumulation'' para recursión
\end{enumerate}
\end{definition}

\begin{keyidea}
\textbf{Innovación principal: Recursión sin trusted setup}

\textbf{El problema:}

Para recursión (probar una prueba dentro de otra), necesitas verificar un SNARK dentro de un circuito.

Verificar pairings en un circuito es muy caro ($\sim$millones de constraints).

\textbf{Solución Halo: Accumulation}

En lugar de verificar completamente:
\begin{enumerate}
    \item Acumula las verificaciones en un ``acumulador''
    \item Al final, una sola verificación para todas
    \item Costo amortizado: casi gratis por prueba recursiva
\end{enumerate}

\textbf{Trade-off:}
\begin{itemize}
    \item Pro: recursión muy eficiente
    \item Pro: sin trusted setup
    \item Con: verificación final es $O(n)$ en lugar de $O(1)$
    \item Con: pruebas más grandes que Groth16
\end{itemize}
\end{keyidea}

\begin{remark}
\textbf{Halo2 en producción:}

\textbf{Zcash Orchard:}
\begin{itemize}
    \item Nuevo pool de shielded transactions
    \item Reemplaza a Sprout (Groth16)
    \item Sin trusted setup = más confianza
\end{itemize}

\textbf{Scroll:}
\begin{itemize}
    \item zkEVM usando Halo2
    \item Custom gates para opcodes EVM
    \item Lookups para tablas de operaciones
\end{itemize}
\end{remark}

%-------------------------------------------
\subsection{Plonky2/Plonky3}

\begin{itemize}
    \item PLONK + FRI (de STARKs)
    \item Campos pequeños (Goldilocks: 64 bits)
    \item Muy rápido para recursión
    \item Usado en Polygon zkEVM
\end{itemize}

\begin{definition}[Plonky2]
Sistema que combina lo mejor de SNARKs y STARKs:
\begin{itemize}
    \item Arithmetización PLONKish (flexible)
    \item Polynomial commitment con FRI (transparent)
    \item Campo pequeño (64 bits) para eficiencia
\end{itemize}
\end{definition}

\begin{mathnote}
\textbf{¿Por qué campos pequeños?}

\textbf{Campo Goldilocks:}
\[ p = 2^{64} - 2^{32} + 1 \]

\textbf{Ventajas:}
\begin{itemize}
    \item Operaciones en 64 bits: nativas en CPU
    \item Multiplicación: una instrucción
    \item NTT (FFT en campo finito): muy rápido
\end{itemize}

\textbf{Comparación de velocidad:}
\begin{center}
\begin{tabular}{lcc}
\toprule
\textbf{Campo} & \textbf{Bits} & \textbf{Velocidad relativa} \\
\midrule
BN254 & 254 & 1x \\
BLS12-381 & 381 & 0.7x \\
Goldilocks & 64 & 10-20x \\
Baby Bear & 31 & 30-50x \\
\bottomrule
\end{tabular}
\end{center}

Plonky2 usa Goldilocks, Plonky3 explora campos aún más pequeños.
\end{mathnote}

\begin{keyidea}
\textbf{FRI en lugar de KZG:}

\textbf{KZG (usado en PLONK tradicional):}
\begin{itemize}
    \item Pruebas pequeñas (48 bytes)
    \item Verificación rápida (pairings)
    \item Requiere trusted setup
    \item No post-quantum
\end{itemize}

\textbf{FRI (usado en Plonky2):}
\begin{itemize}
    \item Pruebas más grandes ($\sim$KB)
    \item Verificación hash-based
    \item Sin trusted setup (transparent)
    \item Post-quantum seguro
\end{itemize}

\textbf{Plonky2 = PLONK + FRI:}

Obtiene:
\begin{itemize}
    \item Flexibilidad de PLONK (custom gates, lookups)
    \item Transparencia de STARKs (no trusted setup)
    \item Velocidad de campos pequeños
\end{itemize}
\end{keyidea}

%-------------------------------------------
\subsection{Tabla Comparativa de Sistemas}

\begin{table}[htbp]
\centering
\begin{tabular}{lcccc}
\toprule
\textbf{Sistema} & \textbf{Prueba} & \textbf{Verify} & \textbf{Setup} & \textbf{Post-Q} \\
\midrule
Groth16 & 128 B & 3 pairings & Trusted/circ & No \\
PLONK & 400 B & 1-2 pairings & Universal & No \\
STARK & 50-200 KB & $O(\log^2 n)$ & Ninguno & Sí \\
Bulletproof & $O(\log n)$ & $O(n)$ & Ninguno & No \\
Halo2 & $\sim$5 KB & $O(n)$ & Ninguno & No \\
Plonky2 & $\sim$50 KB & $O(\log^2 n)$ & Ninguno & Sí \\
\bottomrule
\end{tabular}
\end{table}

\begin{keyidea}
\textbf{Guía para elegir sistema:}

\textbf{Elige Groth16 si:}
\begin{itemize}
    \item Necesitas pruebas mínimas
    \item Verificación on-chain es crítica
    \item Trusted setup es aceptable
    \item El circuito es estable (no cambia)
\end{itemize}

\textbf{Elige PLONK si:}
\begin{itemize}
    \item Quieres flexibilidad (custom gates)
    \item Universal setup es importante
    \item Balancear tamaño y flexibilidad
\end{itemize}

\textbf{Elige STARK si:}
\begin{itemize}
    \item No quieres trusted setup
    \item Necesitas post-quantum security
    \item El tamaño de prueba no es crítico
    \item Tienes mucho cómputo (escalabilidad)
\end{itemize}

\textbf{Elige Bulletproofs si:}
\begin{itemize}
    \item Solo necesitas range proofs
    \item Verificación off-chain
    \item Sin trusted setup es prioritario
\end{itemize}

\textbf{Elige Halo2/Plonky2 si:}
\begin{itemize}
    \item Necesitas recursión eficiente
    \item Sin trusted setup
    \item Dispuesto a trade-off en tamaño
\end{itemize}
\end{keyidea}

\begin{bancozk}
\textbf{Elección de sistema en Banco ZK:}

\textbf{Sistema principal: PLONK con KZG}

Razones:
\begin{itemize}
    \item Universal setup: podemos iterar circuitos rápido
    \item Custom gates: optimizar operaciones de curva elíptica
    \item Lookups: eficiente para range checks
    \item Pruebas razonablemente pequeñas ($\sim$500 bytes)
    \item Verificación on-chain barata ($\sim$200K gas)
\end{itemize}

\textbf{Para recursión: Considerando Plonky2}
\begin{itemize}
    \item Agregar múltiples proofs de usuario
    \item Proof final KZG-based para L1
    \item Trade-off: prover más complejo
\end{itemize}

\textbf{Futuro: Migración a STARKs/Plonky3}
\begin{itemize}
    \item Cuando amenaza cuántica sea realista
    \item Hardware STARK-proving mejorando
    \item L1 agregando precompilados STARK
\end{itemize}
\end{bancozk}

%===========================================
\section{Resumen del Bloque 4}
%===========================================

\begin{keyidea}
\textbf{Puntos clave de Zero-Knowledge Proofs:}

\textbf{1. Definición:}
\begin{itemize}
    \item Probar conocimiento sin revelar el conocimiento
    \item Prover convence a Verifier de que $x \in L$
    \item Sin revelar el witness $w$
\end{itemize}

\textbf{2. Tres propiedades fundamentales:}
\begin{itemize}
    \item Completeness: honesto siempre convence
    \item Soundness: deshonesto no puede convencer
    \item Zero-Knowledge: verifier no aprende nada extra
\end{itemize}

\textbf{3. Interactivo vs NIZK:}
\begin{itemize}
    \item Interactivo: múltiples rondas de comunicación
    \item NIZK: un solo mensaje, requiere CRS o Random Oracle
    \item Fiat-Shamir: transforma interactivo a no-interactivo
\end{itemize}

\textbf{4. Witness:}
\begin{itemize}
    \item Información secreta que satisface la relación
    \item Lo que el prover conoce pero no revela
    \item Witness generation: calcular todos los valores del circuito
\end{itemize}

\textbf{5. Prover vs Verifier:}
\begin{itemize}
    \item Prover: costoso (tiempo y memoria)
    \item Verifier: barato (idealmente $O(1)$)
    \item Asimetría es característica clave de SNARKs
\end{itemize}

\textbf{6. Setup:}
\begin{itemize}
    \item Trusted: ceremonia MPC, toxic waste
    \item Transparent: parámetros públicos derivables
    \item Trade-off: tamaño vs confianza
\end{itemize}

\textbf{7. Sistemas de prueba:}
\begin{itemize}
    \item SNARKs: sucintos, trusted setup (Groth16, PLONK)
    \item STARKs: transparentes, post-quantum, pruebas grandes
    \item Bulletproofs: logarítmicos, verificación lineal
    \item Halo2/Plonky2: recursión eficiente, híbridos
\end{itemize}

\textbf{8. Tendencias:}
\begin{itemize}
    \item Hacia setups universales y transparentes
    \item Recursión para escalabilidad
    \item Hardware acelerado (GPUs, FPGAs, ASICs)
    \item Campos pequeños para velocidad
\end{itemize}
\end{keyidea}

\cleardoublepage

% ============================================
% BLOQUE 5: CIRCUITOS ZK - VERSIÓN EXPANDIDA
% Parte 1: Circuitos Aritméticos y R1CS
% ============================================
\chapter{Circuitos ZK}
\label{ch:bloque5}

\begin{center}
\textit{Bloque 5 de 20}
\end{center}

%===========================================
\section{Qué es un Circuito Aritmético}
%===========================================

\begin{definition}[Circuito Aritmético]
DAG (grafo dirigido acíclico) donde:
\begin{itemize}
    \item Nodos de entrada: variables o constantes
    \item Nodos internos: compuertas de suma o multiplicación
    \item Nodos de salida: resultados a verificar
\end{itemize}
\end{definition}

Operaciones sobre campo finito $\mathbb{F}_p$.

%-------------------------------------------
\subsection{Contexto e Intuición}

\begin{intuicion}
Un circuito aritmético es como una receta de cocina matemática:
\begin{itemize}
    \item \textbf{Ingredientes} (inputs): los valores que entran
    \item \textbf{Pasos} (gates): operaciones de suma y multiplicación
    \item \textbf{Resultado} (output): lo que sale al final
\end{itemize}

La diferencia con un programa normal es que todo se expresa como sumas y multiplicaciones sobre un campo finito. No hay if/else, loops, ni operaciones de bits directas---todo debe traducirse a aritmética.

\textbf{¿Por qué campos finitos?}

Los sistemas ZK (SNARKs, STARKs) están diseñados para trabajar sobre campos finitos porque:
\begin{itemize}
    \item Los polinomios sobre campos finitos tienen propiedades algebraicas útiles
    \item Permiten codificar constraints de forma eficiente
    \item Los compromisos criptográficos (KZG, Pedersen) trabajan sobre grupos finitos
\end{itemize}
\end{intuicion}

\begin{definition}[Campo Finito $\mathbb{F}_p$]
Un campo finito de orden primo $p$ es el conjunto $\{0, 1, 2, \ldots, p-1\}$ con:
\begin{itemize}
    \item Suma: $(a + b) \mod p$
    \item Multiplicación: $(a \cdot b) \mod p$
    \item Inverso aditivo: $-a = p - a$
    \item Inverso multiplicativo: $a^{-1}$ tal que $a \cdot a^{-1} \equiv 1 \pmod{p}$
\end{itemize}

\textbf{Campos comunes en ZK:}
\begin{center}
\begin{tabular}{lcc}
\toprule
\textbf{Curva} & \textbf{Campo escalar} & \textbf{Bits} \\
\midrule
BN254 & $p = 21888...$ & 254 \\
BLS12-381 & $p = 52435...$ & 255 \\
Pasta (Pallas) & $p = 28948...$ & 255 \\
Goldilocks & $p = 2^{64} - 2^{32} + 1$ & 64 \\
\bottomrule
\end{tabular}
\end{center}
\end{definition}

%-------------------------------------------
\subsection{Estructura de un Circuito}

\begin{mathnote}
\textbf{Componentes formales:}

Un circuito aritmético $C$ sobre campo $\mathbb{F}$ consiste en:
\begin{enumerate}
    \item \textbf{Nodos de entrada:} $x_1, \ldots, x_n$ (inputs públicos y privados)
    \item \textbf{Constantes:} $c_1, \ldots, c_k \in \mathbb{F}$
    \item \textbf{Gates:} Cada gate $g$ tiene:
    \begin{itemize}
        \item Tipo: ADD o MUL
        \item Dos inputs: wires de entrada
        \item Un output: wire de salida
    \end{itemize}
    \item \textbf{Wires:} Conexiones entre gates
    \item \textbf{Output:} Wire(s) designado(s) como salida
\end{enumerate}

\textbf{Evaluación:}

El circuito define una función $f: \mathbb{F}^n \to \mathbb{F}^m$:
\begin{itemize}
    \item Propagar valores desde inputs hacia outputs
    \item En cada gate, aplicar la operación correspondiente
    \item El output es el resultado de los wires de salida
\end{itemize}
\end{mathnote}

\begin{center}
\begin{tikzpicture}[
    input/.style={draw, circle, minimum size=0.8cm, fill=blue!20},
    gate/.style={draw, rectangle, minimum size=0.8cm, fill=yellow!30},
    output/.style={draw, circle, minimum size=0.8cm, fill=green!20},
    wire/.style={->, thick}
]
    % Inputs
    \node[input] (x1) at (0, 2) {$x$};
    \node[input] (x2) at (0, 0) {$y$};
    \node[input] (c) at (0, -2) {$5$};
    
    % Gates
    \node[gate] (mul1) at (3, 1) {$\times$};
    \node[gate] (mul2) at (3, -1) {$\times$};
    \node[gate] (add) at (6, 0) {$+$};
    
    % Output
    \node[output] (out) at (9, 0) {out};
    
    % Wires
    \draw[wire] (x1) -- (mul1);
    \draw[wire] (x1) to[bend left=20] (mul1);
    \draw[wire] (x2) -- (mul2);
    \draw[wire] (c) -- (mul2);
    \draw[wire] (mul1) -- node[above] {$x^2$} (add);
    \draw[wire] (mul2) -- node[below] {$5y$} (add);
    \draw[wire] (add) -- node[above] {$x^2+5y$} (out);
    
    % Labels
    \node[below=0.3cm of x1] {\small input};
    \node[below=0.3cm of x2] {\small input};
    \node[below=0.3cm of c] {\small const};
\end{tikzpicture}
\end{center}

\begin{example}[Circuito para $f(x, y) = x^2 + 5y$]
\textbf{Inputs:}
\begin{itemize}
    \item $x$: variable de entrada
    \item $y$: variable de entrada
    \item $5$: constante
\end{itemize}

\textbf{Gates:}
\begin{enumerate}
    \item $w_1 = x \times x$ (multiplicación: $x^2$)
    \item $w_2 = 5 \times y$ (multiplicación: $5y$)
    \item $w_3 = w_1 + w_2$ (suma: $x^2 + 5y$)
\end{enumerate}

\textbf{Output:} $w_3$

\textbf{Número de gates:} 2 multiplicaciones + 1 suma = 3 gates

En R1CS, solo las multiplicaciones generan constraints, así que este circuito tiene 2 constraints.
\end{example}

%-------------------------------------------
\subsection{De Programas a Circuitos}

\begin{keyidea}
\textbf{El desafío de la compilación a circuitos:}

Los programas normales tienen:
\begin{itemize}
    \item Condicionales (if/else)
    \item Loops (for, while)
    \item Operaciones de bits (AND, OR, XOR, shifts)
    \item Comparaciones ($<$, $>$, $=$)
    \item División y módulo
\end{itemize}

Pero los circuitos aritméticos solo tienen:
\begin{itemize}
    \item Suma
    \item Multiplicación
\end{itemize}

\textbf{Solución:} Expresar todo como aritmética de campo.
\end{keyidea}

\begin{example}[Transformaciones comunes]
\textbf{1. Booleano (0 o 1):}

Forzar que $b \in \{0, 1\}$:
\[ b \times (1 - b) = 0 \]

Si $b = 0$: $0 \times 1 = 0$ ✓
Si $b = 1$: $1 \times 0 = 0$ ✓
Si $b = 2$: $2 \times (-1) \neq 0$ ✗

\textbf{2. Selección (if/else):}

$\text{out} = b ? x : y$ donde $b \in \{0, 1\}$:
\[ \text{out} = b \cdot x + (1 - b) \cdot y \]

Si $b = 1$: $\text{out} = 1 \cdot x + 0 \cdot y = x$
Si $b = 0$: $\text{out} = 0 \cdot x + 1 \cdot y = y$

\textbf{3. Igualdad ($x = y$):}

Probar que $x = y$:
\[ x - y = 0 \]

\textbf{4. No-igualdad ($x \neq y$):}

Probar que $x \neq y$ requiere un witness auxiliar $inv$:
\[ inv = (x - y)^{-1} \]
\[ (x - y) \times inv = 1 \]

Si $x = y$, no existe $inv$ que satisfaga esto.

\textbf{5. Descomposición en bits:}

Para probar que $x$ tiene representación binaria $b_0, b_1, \ldots, b_{n-1}$:
\[ x = \sum_{i=0}^{n-1} b_i \cdot 2^i \]
\[ \forall i: b_i \times (1 - b_i) = 0 \]

Esto cuesta $n$ constraints (uno por cada bit).
\end{example}

\begin{advertencia}
\textbf{Operaciones costosas en circuitos:}

\begin{center}
\begin{tabular}{lcc}
\toprule
\textbf{Operación} & \textbf{Constraints} & \textbf{Razón} \\
\midrule
Suma & 0 & Combinación lineal \\
Multiplicación & 1 & Gate básico \\
División & 1 & Verificar $a \cdot b^{-1} = c$ \\
Comparación ($<$) & $\sim n$ & Descomposición en bits \\
Hash SHA-256 & $\sim 25,000$ & Muchas ops de bits \\
Hash Poseidon & $\sim 250$ & Diseñado para circuitos \\
ECDSA verify & $\sim 500,000$ & Aritmética de curva \\
\bottomrule
\end{tabular}
\end{center}

Los hashes como SHA-256 son muy caros porque usan XOR, rotaciones, y operaciones de bits que no son nativas en aritmética de campo.
\end{advertencia}

%===========================================
\section{R1CS (Rank-1 Constraint System)}
%===========================================

\begin{definition}[R1CS]
Sistema de $m$ ecuaciones sobre $n$ variables:
\[
(A_i \cdot \vec{w}) \times (B_i \cdot \vec{w}) = C_i \cdot \vec{w} \quad \forall i \in [m]
\]
donde $A_i, B_i, C_i$ son vectores de coeficientes y $\vec{w}$ es el witness.
\end{definition}

Forma matricial:
\[
(A \cdot \vec{w}) \circ (B \cdot \vec{w}) = C \cdot \vec{w}
\]
donde $\circ$ es producto Hadamard (elemento a elemento).

%-------------------------------------------
\subsection{Entendiendo R1CS}

\begin{intuicion}
R1CS es como un sistema de ecuaciones especial:
\begin{itemize}
    \item Cada ecuación tiene la forma: (algo) $\times$ (algo) $=$ (algo)
    \item Los ``algos'' son combinaciones lineales de las variables
    \item Una multiplicación por ecuación (por eso ``Rank-1'')
\end{itemize}

Es la representación intermedia más común entre un programa y un SNARK.

\textbf{Flujo típico:}
\begin{verbatim}
Programa → Circuito → R1CS → QAP → SNARK
     (Circom)        (compilador)   (Groth16/PLONK)
\end{verbatim}
\end{intuicion}

\begin{mathnote}
\textbf{Estructura detallada:}

\textbf{Witness vector:}
\[ \vec{w} = (1, x_1, x_2, \ldots, x_n, w_1, w_2, \ldots, w_k) \]

Donde:
\begin{itemize}
    \item $1$: constante (siempre el primer elemento)
    \item $x_i$: inputs públicos y privados
    \item $w_j$: variables intermedias (wires internos)
\end{itemize}

\textbf{Matrices $A$, $B$, $C$:}

Cada matriz tiene $m$ filas (una por constraint) y $n+k+1$ columnas (una por variable).

\textbf{Constraint $i$:}
\[ \underbrace{(a_{i,0} + a_{i,1}x_1 + \ldots + a_{i,n}w_k)}_{\text{lado izquierdo}} \times \underbrace{(b_{i,0} + b_{i,1}x_1 + \ldots)}_{\text{lado derecho}} = \underbrace{(c_{i,0} + c_{i,1}x_1 + \ldots)}_{\text{resultado}} \]
\end{mathnote}

%-------------------------------------------
\subsection{Ejemplo Completo de R1CS}

\begin{example}[R1CS para $x^3 + x + 5 = 35$]
\textbf{Paso 1: Descomponer en multiplicaciones}

La expresión $x^3 + x + 5$ requiere calcular $x^3$:
\begin{align*}
w_1 &= x \times x \quad \text{(compute } x^2\text{)} \\
w_2 &= w_1 \times x \quad \text{(compute } x^3\text{)} \\
w_3 &= w_2 + x + 5 \quad \text{(compute } x^3 + x + 5\text{)} \\
\end{align*}

Y verificar que $w_3 = 35$.

\textbf{Paso 2: Convertir a constraints R1CS}

Witness vector: $\vec{w} = (1, x, w_1, w_2, w_3, \text{out})$

Con índices: $w_0 = 1$, $w_1 = x$, $w_2 = x^2$, $w_3 = x^3$, $w_4 = x^3+x+5$, $w_5 = 35$

\textbf{Constraint 1:} $x \times x = w_1$
\begin{align*}
A_1 &= (0, 1, 0, 0, 0, 0) \quad \text{(selecciona } x\text{)} \\
B_1 &= (0, 1, 0, 0, 0, 0) \quad \text{(selecciona } x\text{)} \\
C_1 &= (0, 0, 1, 0, 0, 0) \quad \text{(selecciona } w_1\text{)}
\end{align*}
Verificación: $(0 + 1 \cdot x) \times (0 + 1 \cdot x) = (0 + 0 + 1 \cdot w_1)$
$\Rightarrow x \times x = w_1$ ✓

\textbf{Constraint 2:} $w_1 \times x = w_2$
\begin{align*}
A_2 &= (0, 0, 1, 0, 0, 0) \\
B_2 &= (0, 1, 0, 0, 0, 0) \\
C_2 &= (0, 0, 0, 1, 0, 0)
\end{align*}

\textbf{Constraint 3:} $1 \times (w_2 + x + 5) = w_3$
\begin{align*}
A_3 &= (1, 0, 0, 0, 0, 0) \quad \text{(selecciona 1)} \\
B_3 &= (5, 1, 0, 1, 0, 0) \quad \text{(selecciona } 5 + x + w_2\text{)} \\
C_3 &= (0, 0, 0, 0, 1, 0)
\end{align*}

\textbf{Constraint 4:} $1 \times w_3 = 35$
\begin{align*}
A_4 &= (1, 0, 0, 0, 0, 0) \\
B_4 &= (0, 0, 0, 0, 1, 0) \\
C_4 &= (35, 0, 0, 0, 0, 0)
\end{align*}

\textbf{Paso 3: Verificar con $x = 3$}

$\vec{w} = (1, 3, 9, 27, 35, 35)$

\begin{itemize}
    \item Constraint 1: $3 \times 3 = 9$ ✓
    \item Constraint 2: $9 \times 3 = 27$ ✓
    \item Constraint 3: $1 \times (27 + 3 + 5) = 35$ ✓
    \item Constraint 4: $1 \times 35 = 35$ ✓
\end{itemize}
\end{example}

%-------------------------------------------
\subsection{Matrices R1CS}

\begin{mathnote}
\textbf{Visualización matricial del ejemplo:}

\[
A = \begin{pmatrix}
0 & 1 & 0 & 0 & 0 & 0 \\
0 & 0 & 1 & 0 & 0 & 0 \\
1 & 0 & 0 & 0 & 0 & 0 \\
1 & 0 & 0 & 0 & 0 & 0
\end{pmatrix}
\quad
B = \begin{pmatrix}
0 & 1 & 0 & 0 & 0 & 0 \\
0 & 1 & 0 & 0 & 0 & 0 \\
5 & 1 & 0 & 1 & 0 & 0 \\
0 & 0 & 0 & 0 & 1 & 0
\end{pmatrix}
\quad
C = \begin{pmatrix}
0 & 0 & 1 & 0 & 0 & 0 \\
0 & 0 & 0 & 1 & 0 & 0 \\
0 & 0 & 0 & 0 & 1 & 0 \\
35 & 0 & 0 & 0 & 0 & 0
\end{pmatrix}
\]

\textbf{Verificación:}
\[
(A \cdot \vec{w}) \circ (B \cdot \vec{w}) = \begin{pmatrix} 3 \\ 9 \\ 1 \\ 1 \end{pmatrix} \circ \begin{pmatrix} 3 \\ 3 \\ 35 \\ 35 \end{pmatrix} = \begin{pmatrix} 9 \\ 27 \\ 35 \\ 35 \end{pmatrix} = C \cdot \vec{w}
\]
\end{mathnote}

%===========================================
\section{Constraints}
%===========================================

\begin{definition}[Constraint]
Ecuación que el witness debe satisfacer.
\end{definition}

Forma general R1CS:
\[
(\sum_j a_j w_j) \times (\sum_k b_k w_k) = \sum_l c_l w_l
\]

\begin{keyidea}
En R1CS:
\begin{itemize}
    \item Sumas son ``gratis'' (combinaciones lineales)
    \item Solo multiplicaciones generan constraints
    \item Número de constraints = tamaño del circuito
\end{itemize}
\end{keyidea}

%-------------------------------------------
\subsection{Costos de Operaciones}

\begin{table}[htbp]
\centering
\begin{tabular}{lr}
\toprule
\textbf{Operación} & \textbf{Constraints} \\
\midrule
Suma/resta & 0 \\
Multiplicación & 1 \\
División & 1 (verificar $a \times b^{-1} = c$) \\
Comparación ($<$) & $\sim n$ bits \\
Binario check & 1 ($b(1-b)=0$) \\
Hash Poseidon & $\sim 250$ \\
Hash SHA-256 & $\sim 25,000$ \\
ECDSA verify & $\sim 500,000$ \\
\bottomrule
\end{tabular}
\end{table}

%-------------------------------------------
\subsection{Constraints Comunes en Detalle}

\begin{example}[IsZero: verificar si $x = 0$]
Queremos $out = 1$ si $x = 0$, y $out = 0$ si $x \neq 0$.

\textbf{Estrategia:} Usar inverso como witness auxiliar.

\textbf{Constraints:}
\begin{align*}
\text{inv} &\leftarrow \begin{cases} 0 & \text{si } x = 0 \\ x^{-1} & \text{si } x \neq 0 \end{cases} \\
\text{out} &= 1 - x \cdot \text{inv} \\
x \cdot \text{out} &= 0
\end{align*}

\textbf{Análisis:}
\begin{itemize}
    \item Si $x = 0$: $\text{out} = 1 - 0 = 1$, y $0 \cdot 1 = 0$ ✓
    \item Si $x \neq 0$: $\text{inv} = x^{-1}$, $\text{out} = 1 - 1 = 0$, y $x \cdot 0 = 0$ ✓
    \item Si $x \neq 0$ pero intenta $\text{out} = 1$: $x \cdot 1 = x \neq 0$ ✗
\end{itemize}

\textbf{Costo:} 2 constraints.
\end{example}

\begin{example}[LessThan: verificar $a < b$]
Para comparar en un campo finito, necesitamos descomponer en bits.

\textbf{Estrategia:} 
\begin{enumerate}
    \item Calcular $\text{diff} = b - a + 2^n$ (asegurar positivo si $a < b$)
    \item Descomponer $\text{diff}$ en $n+1$ bits
    \item El bit más significativo indica si $a < b$
\end{enumerate}

\textbf{Para $n$-bit numbers:}
\begin{verbatim}
// Pseudo-código
diff = b - a + 2^n
bits = decompose(diff, n+1)
// Si a < b: diff está en [2^n, 2^(n+1)), bit[n] = 1
// Si a >= b: diff está en [0, 2^n), bit[n] = 0

// Constraints:
for i in 0..n+1:
    bits[i] * (1 - bits[i]) = 0  // Es binario
    
sum(bits[i] * 2^i) = diff        // Reconstrucción

out = bits[n]                     // MSB indica resultado
\end{verbatim}

\textbf{Costo:} $n + 2$ constraints (uno por bit + reconstrucción).
\end{example}

\begin{example}[Range check: $x \in [0, 2^n)$]
Verificar que $x$ está en un rango.

\textbf{Método 1: Descomposición en bits}
\begin{verbatim}
bits = decompose(x, n)
for i in 0..n:
    bits[i] * (1 - bits[i]) = 0  // Cada bit es 0 o 1
sum(bits[i] * 2^i) = x           // Reconstrucción
\end{verbatim}

Costo: $n$ constraints.

\textbf{Método 2: Lookup tables (PLONKish)}

En sistemas como PLONK con Plookup:
\begin{itemize}
    \item Tabla precalculada: $\{0, 1, 2, \ldots, 2^{16}-1\}$
    \item Verificar: $x$ está en la tabla
    \item Costo: $\sim$1 constraint (amortizado)
\end{itemize}

Mucho más eficiente para rangos grandes.
\end{example}

%-------------------------------------------
\subsection{Optimización de Constraints}

\begin{keyidea}
\textbf{Técnicas para reducir constraints:}

\textbf{1. Agrupar operaciones lineales:}
\begin{verbatim}
// Malo: 2 multiplicaciones
w1 = a * b
w2 = w1 * c

// Equivalente pero unavoidable
// Solo podemos multiplicar 2 cosas a la vez
\end{verbatim}

\textbf{2. Usar constraints de ``igualdad gratuita'':}
\begin{verbatim}
// Malo: copiar valor con constraint
w2 = w1 * 1  // 1 constraint

// Bueno: usar el mismo wire
// No hay constraint, solo referencia
\end{verbatim}

\textbf{3. Reutilizar cálculos:}
\begin{verbatim}
// Malo: calcular x^2 dos veces
y = x * x
z = x * x + 1

// Bueno: calcular una vez
w = x * x    // 1 constraint
y = w
z = w + 1   // 0 constraints (suma)
\end{verbatim}

\textbf{4. Elegir representaciones eficientes:}
\begin{verbatim}
// Para booleanos, usar 0/1 en lugar de true/false
// Para enums, usar índices numéricos
// Para hashes, usar Poseidon en lugar de SHA-256
\end{verbatim}
\end{keyidea}

\begin{bancozk}
\textbf{Constraints en Banco ZK:}

\textbf{Circuito de transferencia privada:}

\begin{center}
\begin{tabular}{lr}
\toprule
\textbf{Componente} & \textbf{Constraints aprox.} \\
\midrule
Verificar Merkle path (depth 20) & 5,000 \\
Hash del nullifier (Poseidon) & 250 \\
Hash del nuevo commitment (Poseidon) & 250 \\
Verificar firma (EdDSA) & 3,000 \\
Range check del monto & 64 \\
Verificar balance suficiente & 100 \\
\midrule
\textbf{Total aproximado} & $\sim$9,000 \\
\bottomrule
\end{tabular}
\end{center}

Con un prover optimizado, generar prueba para 9,000 constraints toma $\sim$1-2 segundos.

\textbf{Por qué usamos Poseidon:}
\begin{itemize}
    \item SHA-256: 25,000 constraints por hash
    \item Poseidon: 250 constraints por hash
    \item Ahorro: 100x menos constraints
    \item Un Merkle tree de depth 20 con SHA-256 sería $\sim$500,000 constraints
\end{itemize}
\end{bancozk}

% ============================================
% BLOQUE 5: CIRCUITOS ZK - VERSIÓN EXPANDIDA
% Parte 2: Señales, Witness Generation, Circom
% ============================================

%===========================================
\section{Señales (Signals)}
%===========================================

\begin{definition}[Signal]
Variable en un circuito, contiene elemento de $\mathbb{F}_p$.
\end{definition}

Tipos:
\begin{itemize}
    \item \textbf{input public:} Visible al verifier
    \item \textbf{input private:} Solo conocido por prover
    \item \textbf{output:} Resultado del circuito
    \item \textbf{intermediate:} Variables internas
\end{itemize}

%-------------------------------------------
\subsection{Anatomía de las Señales}

\begin{intuicion}
Las señales son los ``cables'' del circuito:
\begin{itemize}
    \item Transportan valores entre gates
    \item Cada señal tiene un valor en $\mathbb{F}_p$
    \item Las señales conectan inputs, gates, y outputs
\end{itemize}

En el contexto de ZK:
\begin{itemize}
    \item \textbf{Públicas:} El verifier las conoce (parte del statement)
    \item \textbf{Privadas:} Solo el prover las conoce (parte del witness)
\end{itemize}
\end{intuicion}

\begin{mathnote}
\textbf{Clasificación formal de señales:}

\textbf{1. Public Inputs ($x$):}
\begin{itemize}
    \item Conocidos por prover Y verifier
    \item Parte del ``statement'' que se prueba
    \item Ejemplo: el hash resultante, el root del Merkle tree
\end{itemize}

\textbf{2. Private Inputs (witness $w$):}
\begin{itemize}
    \item Conocidos SOLO por el prover
    \item Lo que queremos mantener secreto
    \item Ejemplo: la preimagen del hash, la clave privada
\end{itemize}

\textbf{3. Intermediate Signals:}
\begin{itemize}
    \item Calculados durante la ejecución
    \item Parte del witness extendido
    \item Ejemplo: resultados de multiplicaciones intermedias
\end{itemize}

\textbf{4. Output Signals:}
\begin{itemize}
    \item Resultados finales del circuito
    \item Pueden ser públicos o privados según el diseño
\end{itemize}

\textbf{Witness completo:}
\[ \vec{w} = (1, x_1, \ldots, x_l, w_1, \ldots, w_m) \]
donde $x_i$ son public inputs y $w_j$ son todas las señales privadas + intermedias.
\end{mathnote}

\begin{example}[Señales en un circuito de hash]
\begin{verbatim}
Circuito: Probar conocimiento de preimagen

Public inputs:
    - hash_output: H(x) = y (conocido públicamente)

Private inputs:
    - preimage: x (secreto del prover)

Intermediate signals:
    - Todos los estados internos del hash
    - Variables temporales de las rondas
    
Verificación:
    - El circuito calcula H(preimage)
    - Constraint: H(preimage) === hash_output
\end{verbatim}

El verifier solo ve \texttt{hash\_output}. No aprende nada sobre \texttt{preimage}.
\end{example}

%-------------------------------------------
\subsection{Señales en Diferentes Frameworks}

\begin{center}
\begin{tabular}{lccc}
\toprule
\textbf{Concepto} & \textbf{Circom} & \textbf{Halo2} & \textbf{Noir} \\
\midrule
Input público & \texttt{signal input} & \texttt{Instance} & \texttt{pub param} \\
Input privado & \texttt{signal private input} & \texttt{Advice} & \texttt{param} \\
Output & \texttt{signal output} & \texttt{Advice} & return value \\
Intermedio & \texttt{signal} & \texttt{Advice} & \texttt{let} \\
Constante & literal/\texttt{var} & \texttt{Fixed} & \texttt{const} \\
\bottomrule
\end{tabular}
\end{center}

%===========================================
\section{Witness Generation}
%===========================================

\begin{definition}[Witness]
Asignación de valores a todas las señales que satisface todos los constraints.
\end{definition}

Proceso:
\begin{enumerate}
    \item Input: valores de entrada (JSON)
    \item Ejecutar lógica del circuito
    \item Calcular todas las señales intermedias
    \item Output: archivo .wtns con todos los valores
\end{enumerate}

%-------------------------------------------
\subsection{El Proceso de Witness Generation}

\begin{keyidea}
\textbf{Witness generation es ``ejecutar'' el circuito:}

\begin{verbatim}
Input (JSON):
{
    "a": "3",
    "b": "5"
}

Circuito:
    signal input a;
    signal input b;
    signal c;
    signal output d;
    
    c <== a * b;      // c = 15
    d <== c + a;      // d = 18

Output (witness):
    w = [1, 3, 5, 15, 18]
    //   1  a  b  c   d
\end{verbatim}

El witness generator:
\begin{enumerate}
    \item Lee los inputs
    \item Propaga valores a través del circuito
    \item Evalúa cada constraint/asignación
    \item Produce el vector completo de señales
\end{enumerate}
\end{keyidea}

%-------------------------------------------
\subsection{Witness vs Constraints}

\begin{mathnote}
\textbf{Dos fases separadas:}

\textbf{1. Witness Generation (ejecución):}
\begin{itemize}
    \item Calcula valores de todas las señales
    \item Usa código imperativo (asignaciones \texttt{<--})
    \item Puede usar if/else, loops, funciones auxiliares
    \item Output: vector de valores
\end{itemize}

\textbf{2. Constraint Checking (verificación):}
\begin{itemize}
    \item Verifica que el witness satisface todos los constraints
    \item Solo usa ecuaciones algebraicas (\texttt{===})
    \item Sin control de flujo
    \item Output: válido/inválido
\end{itemize}

\textbf{Importante:} El prover debe generar un witness que satisfaga TODOS los constraints. El verifier no ejecuta el witness generation---solo verifica la prueba.
\end{mathnote}

\begin{advertencia}
\textbf{Error común: witness generation vs constraints}

\begin{verbatim}
// INCORRECTO: constraint sin asignación
signal x;
x === a * b;  // Error: x no tiene valor asignado

// CORRECTO: asignación + constraint
signal x;
x <== a * b;  // Asigna Y crea constraint

// TAMBIÉN CORRECTO: separado
signal x;
x <-- a * b;  // Solo asignación
x === a * b;  // Solo constraint
\end{verbatim}

El operador \texttt{<==} es azúcar sintáctica para \texttt{<--} seguido de \texttt{===}.
\end{advertencia}

%-------------------------------------------
\subsection{Formato del Witness}

\begin{example}[Archivo witness (.wtns)]
El formato .wtns es binario, pero conceptualmente:

\begin{verbatim}
// Header
version: 2
sections: [
    {type: "header", field: "bn128", n_vars: 5},
    {type: "data", values: [...]}
]

// Data section (little-endian field elements)
w[0] = 1                    // Siempre 1
w[1] = 3                    // Primera variable
w[2] = 5                    // Segunda variable
w[3] = 15                   // Señal intermedia
w[4] = 18                   // Output

// Cada valor es un elemento de F_p (32 bytes para BN254)
\end{verbatim}
\end{example}

%===========================================
\section{Circom}
%===========================================

\begin{definition}[Circom]
Lenguaje de dominio específico (DSL) para escribir circuitos ZK.
\end{definition}

%-------------------------------------------
\subsection{Historia y Contexto}

\begin{remark}
\textbf{Historia de Circom:}

\begin{center}
\begin{tabular}{lll}
\toprule
\textbf{Año} & \textbf{Versión} & \textbf{Desarrollo} \\
\midrule
2018 & Circom 1.0 & Creado por iden3/Jordi Baylina \\
2020 & Circom 1.x & Mejoras de sintaxis \\
2021 & Circom 2.0 & Reescrito en Rust, más rápido \\
2022 & Circom 2.1 & Custom templates, mejoras \\
2023 & Actual & Estándar de facto para ZK \\
\bottomrule
\end{tabular}
\end{center}

Circom es el lenguaje más usado para circuitos ZK en el ecosistema Ethereum. Proyectos que usan Circom:
\begin{itemize}
    \item Tornado Cash (mixer privado)
    \item Semaphore (identidad anónima)
    \item Hermez/Polygon zkEVM
    \item Dark Forest (juego ZK)
    \item WorldCoin
\end{itemize}
\end{remark}

%-------------------------------------------
\subsection{Sintaxis Básica}

\begin{verbatim}
pragma circom 2.0.0;

template Multiplicador() {
    signal input a;
    signal input b;
    signal output c;
    
    c <== a * b;
}

component main = Multiplicador();
\end{verbatim}

\begin{mathnote}
\textbf{Estructura de un archivo Circom:}

\begin{verbatim}
// 1. Pragma (versión del compilador)
pragma circom 2.0.0;

// 2. Includes (bibliotecas)
include "circomlib/poseidon.circom";

// 3. Templates (definiciones de circuitos)
template MiCircuito(param1, param2) {
    // Declaraciones de señales
    signal input x;
    signal output y;
    
    // Lógica y constraints
    y <== x * x;
}

// 4. Componente main
component main {public [x]} = MiCircuito(10, 20);
\end{verbatim}

\textbf{El main especifica:}
\begin{itemize}
    \item Qué template es el punto de entrada
    \item Qué inputs son públicos (\texttt{public [...]})
    \item Parámetros del template
\end{itemize}
\end{mathnote}

%-------------------------------------------
\subsection{Operadores}

\begin{itemize}
    \item \texttt{<--}: Asignación sin constraint
    \item \texttt{-->}: Asignación inversa sin constraint
    \item \texttt{<==}: Asignación CON constraint (lo más usado)
    \item \texttt{==>}: Asignación inversa con constraint
    \item \texttt{===}: Solo constraint, sin asignación
\end{itemize}

\begin{keyidea}
\textbf{Entendiendo los operadores de Circom:}

\textbf{\texttt{<--} (asignación sin constraint):}
\begin{verbatim}
// Solo calcula el valor, NO verifica nada
x <-- y / z;  // División (no es operación de campo nativa)
x <-- cond ? a : b;  // Condicional
\end{verbatim}
Útil cuando la operación no es directamente expresable como constraint.

\textbf{\texttt{<==} (asignación + constraint):}
\begin{verbatim}
// Calcula Y verifica
x <== a * b;  // Asigna x = a*b, añade constraint x === a*b
x <== a + b;  // Asigna x = a+b, añade constraint x === a+b
\end{verbatim}
Usar siempre que sea posible---es más seguro.

\textbf{\texttt{===} (solo constraint):}
\begin{verbatim}
// Solo verifica, no asigna
a * b === c;  // Verifica que a*b = c
x === 5;      // Verifica que x = 5
\end{verbatim}
Usar después de \texttt{<--} para completar la verificación.

\textbf{Patrón común:}
\begin{verbatim}
// Para división: z = x / y
z <-- x / y;      // Calcular (no es constraint válido)
z * y === x;      // Verificar que z * y = x
\end{verbatim}
\end{keyidea}

\begin{advertencia}
\textbf{Peligro: \texttt{<--} sin constraint}

\begin{verbatim}
// VULNERABLE: el prover puede poner cualquier valor
signal output result;
result <-- computeSecretValue();  // Sin constraint!

// Un prover malicioso puede:
result <-- 999999;  // Valor arbitrario, nadie lo verifica

// CORRECTO: siempre añadir constraint
result <-- computeSecretValue();
result === expectedValue;  // O algún constraint que lo ate
\end{verbatim}

Regla: Todo \texttt{<--} debe tener un \texttt{===} correspondiente que lo verifique.
\end{advertencia}

%-------------------------------------------
\subsection{Templates}

\begin{definition}[Template]
Función que genera porción de circuito, reutilizable y parametrizable.
\end{definition}

\begin{verbatim}
template IsZero() {
    signal input in;
    signal output out;
    
    signal inv;
    inv <-- in != 0 ? 1/in : 0;
    out <== -in * inv + 1;
    in * out === 0;
}
\end{verbatim}

\begin{mathnote}
\textbf{Análisis del template IsZero:}

\textbf{Objetivo:} $out = 1$ si $in = 0$, $out = 0$ si $in \neq 0$.

\textbf{Línea por línea:}
\begin{verbatim}
inv <-- in != 0 ? 1/in : 0;
\end{verbatim}
Calcula el inverso si existe. Si $in = 0$, pone $inv = 0$.

\begin{verbatim}
out <== -in * inv + 1;
\end{verbatim}
Si $in \neq 0$: $out = -in \cdot (1/in) + 1 = -1 + 1 = 0$
Si $in = 0$: $out = 0 \cdot 0 + 1 = 1$

\begin{verbatim}
in * out === 0;
\end{verbatim}
Garantiza que no puedes tener $in \neq 0$ Y $out \neq 0$.

\textbf{Constraints:} 2 (uno de \texttt{<==}, uno de \texttt{===})
\end{mathnote}

\begin{example}[Templates con parámetros]
\begin{verbatim}
// Template parametrizado para N bits
template Bits2Num(n) {
    signal input in[n];
    signal output out;
    
    var lc = 0;
    var e2 = 1;
    
    for (var i = 0; i < n; i++) {
        lc += in[i] * e2;
        e2 = e2 * 2;
    }
    
    out <== lc;
}

// Uso: convertir 8 bits a número
component b2n = Bits2Num(8);
b2n.in <== [1, 0, 1, 0, 0, 0, 0, 0];  // 5 en binario
// b2n.out será 5
\end{verbatim}
\end{example}

%-------------------------------------------
\subsection{Components}

\begin{definition}[Component]
Instancia de un template dentro de otro circuito.
\end{definition}

\begin{verbatim}
template Main() {
    signal input x;
    signal output isZ;
    
    component iz = IsZero();
    iz.in <== x;
    isZ <== iz.out;
}
\end{verbatim}

\begin{keyidea}
\textbf{Composición de circuitos:}

Los components permiten construir circuitos complejos a partir de bloques simples:

\begin{verbatim}
template ComplexCircuit() {
    signal input a, b, c;
    signal output result;
    
    // Usar múltiples components
    component mul = Multiplier();
    component add = Adder();
    component isz = IsZero();
    
    // Conectar
    mul.a <== a;
    mul.b <== b;
    
    add.a <== mul.out;
    add.b <== c;
    
    isz.in <== add.out;
    
    result <== isz.out;
}
\end{verbatim}

Cada component añade sus propios constraints al circuito total.
\end{keyidea}

%-------------------------------------------
\subsection{Arrays y Loops}

\begin{example}[Trabajando con arrays]
\begin{verbatim}
template SumArray(n) {
    signal input in[n];
    signal output out;
    
    // Variables auxiliares (no son señales)
    var sum = 0;
    
    // Loops se desenrollan en compile-time
    for (var i = 0; i < n; i++) {
        sum += in[i];
    }
    
    out <== sum;
}

template DotProduct(n) {
    signal input a[n];
    signal input b[n];
    signal output out;
    
    signal products[n];
    var sum = 0;
    
    for (var i = 0; i < n; i++) {
        products[i] <== a[i] * b[i];  // n constraints
        sum += products[i];
    }
    
    out <== sum;  // 1 constraint
}
\end{verbatim}
\end{example}

%-------------------------------------------
\subsection{Circomlib: Biblioteca Estándar}

\begin{remark}
\textbf{Circomlib es la biblioteca estándar de circuitos:}

\begin{center}
\begin{tabular}{ll}
\toprule
\textbf{Categoría} & \textbf{Templates} \\
\midrule
Comparadores & IsZero, IsEqual, LessThan, GreaterThan \\
Aritmética & Multiplier, Sum, Division \\
Bits & Num2Bits, Bits2Num, AND, OR, XOR \\
Hashes & Poseidon, MiMC, Pedersen \\
Merkle trees & MerkleTreeChecker, MerkleTreeInclusionProof \\
Firmas & EdDSA, ECDSA \\
Otros & Mux, Decoder, Switcher \\
\bottomrule
\end{tabular}
\end{center}

\textbf{Instalación:}
\begin{verbatim}
npm install circomlib
\end{verbatim}

\textbf{Uso:}
\begin{verbatim}
include "circomlib/circuits/poseidon.circom";
include "circomlib/circuits/comparators.circom";
\end{verbatim}
\end{remark}

\begin{example}[Verificación de Merkle proof con Circomlib]
\begin{verbatim}
pragma circom 2.0.0;

include "circomlib/circuits/poseidon.circom";
include "circomlib/circuits/mux1.circom";

template MerkleProof(levels) {
    signal input leaf;
    signal input pathElements[levels];
    signal input pathIndices[levels];  // 0 = izq, 1 = der
    signal output root;
    
    signal hashes[levels + 1];
    hashes[0] <== leaf;
    
    component hashers[levels];
    component mux[levels];
    
    for (var i = 0; i < levels; i++) {
        // Seleccionar orden según pathIndex
        mux[i] = MultiMux1(2);
        mux[i].c[0][0] <== hashes[i];
        mux[i].c[0][1] <== pathElements[i];
        mux[i].c[1][0] <== pathElements[i];
        mux[i].c[1][1] <== hashes[i];
        mux[i].s <== pathIndices[i];
        
        // Hash de los dos hijos
        hashers[i] = Poseidon(2);
        hashers[i].inputs[0] <== mux[i].out[0];
        hashers[i].inputs[1] <== mux[i].out[1];
        
        hashes[i + 1] <== hashers[i].out;
    }
    
    root <== hashes[levels];
}
\end{verbatim}
\end{example}

\begin{bancozk}
\textbf{Circom en Banco ZK:}

\textbf{Estructura de nuestros circuitos:}

\begin{verbatim}
banco-zk-circuits/
├── src/
│   ├── transfer.circom      // Transferencia privada
│   ├── deposit.circom       // Depósito desde L1
│   ├── withdraw.circom      // Retiro a L1
│   └── batch.circom         // Agregación de proofs
├── lib/
│   ├── merkle.circom        // Verificación Merkle
│   ├── nullifier.circom     // Generación de nullifiers
│   └── commitment.circom    // Creación de commitments
├── test/
│   └── *.test.js
└── circomlib/               // Dependencia
\end{verbatim}

\textbf{Template principal de transferencia:}
\begin{verbatim}
template Transfer(levels) {
    // Inputs públicos
    signal input root;              // Merkle root actual
    signal input nullifierHash;     // Nullifier de input
    signal input newCommitment;     // Commitment de output
    
    // Inputs privados (witness)
    signal private input secret;    // Clave del usuario
    signal private input amount;    // Monto
    signal private input pathElements[levels];
    signal private input pathIndices[levels];
    
    // Verificaciones
    component merkle = MerkleProof(levels);
    component nullGen = NullifierGenerator();
    component commGen = CommitmentGenerator();
    
    // ... constraints ...
}
\end{verbatim}
\end{bancozk}

% ============================================
% BLOQUE 5: CIRCUITOS ZK - VERSIÓN EXPANDIDA
% Parte 3: snarkjs - Flujo Completo
% ============================================

%===========================================
\section{snarkjs}
%===========================================

Herramienta JavaScript para el flujo completo:

\begin{definition}[snarkjs]
Biblioteca JavaScript/WASM que implementa el flujo completo de SNARKs:
\begin{itemize}
    \item Setup de ceremonias (Powers of Tau)
    \item Generación de claves (proving key, verification key)
    \item Generación de pruebas
    \item Verificación de pruebas
    \item Exportación a Solidity
\end{itemize}
\end{definition}

%-------------------------------------------
\subsection{Visión General del Flujo}

\begin{center}
\begin{tikzpicture}[
    box/.style={draw, rectangle, minimum width=2.5cm, minimum height=1cm, align=center, fill=blue!10},
    file/.style={draw, rectangle, minimum width=2cm, minimum height=0.7cm, align=center, fill=green!10},
    arrow/.style={->, thick}
]
    % Fase 1: Compilación
    \node[box] (circom) at (0, 6) {circuit.circom};
    \node[box] (compile) at (0, 4.5) {Compilador\\Circom};
    \node[file] (r1cs) at (-2, 3) {.r1cs};
    \node[file] (wasm) at (0, 3) {.wasm};
    \node[file] (sym) at (2, 3) {.sym};
    
    \draw[arrow] (circom) -- (compile);
    \draw[arrow] (compile) -- (r1cs);
    \draw[arrow] (compile) -- (wasm);
    \draw[arrow] (compile) -- (sym);
    
    % Fase 2: Setup
    \node[box] (ptau) at (-4, 1.5) {Powers of Tau\\Ceremony};
    \node[file] (ptauf) at (-4, 0) {.ptau};
    \node[box] (setup) at (0, 0) {Circuit Setup\\(Phase 2)};
    \node[file] (zkey) at (0, -1.5) {.zkey};
    \node[file] (vkey) at (2, -1.5) {vkey.json};
    
    \draw[arrow] (ptau) -- (ptauf);
    \draw[arrow] (ptauf) -- (setup);
    \draw[arrow] (r1cs) -- (setup);
    \draw[arrow] (setup) -- (zkey);
    \draw[arrow] (zkey) -- (vkey);
    
    % Fase 3: Prove
    \node[box] (input) at (5, 3) {input.json};
    \node[box] (wtns) at (5, 1.5) {Witness\\Generator};
    \node[file] (wtnsf) at (5, 0) {.wtns};
    \node[box] (prove) at (5, -1.5) {Prover};
    \node[file] (proof) at (5, -3) {proof.json};
    
    \draw[arrow] (input) -- (wtns);
    \draw[arrow] (wasm) -- (wtns);
    \draw[arrow] (wtns) -- (wtnsf);
    \draw[arrow] (wtnsf) -- (prove);
    \draw[arrow] (zkey) -- (prove);
    \draw[arrow] (prove) -- (proof);
    
    % Fase 4: Verify
    \node[box] (verify) at (5, -4.5) {Verifier};
    \node[box] (result) at (5, -6) {✓ / ✗};
    
    \draw[arrow] (proof) -- (verify);
    \draw[arrow] (vkey) |- (verify);
    \draw[arrow] (verify) -- (result);
    
    % Labels
    \node[left] at (-5, 4.5) {\textbf{Fase 1:} Compilación};
    \node[left] at (-5, 0.5) {\textbf{Fase 2:} Setup};
    \node[left] at (-5, -3) {\textbf{Fase 3:} Prove};
    \node[left] at (-5, -5) {\textbf{Fase 4:} Verify};
\end{tikzpicture}
\end{center}

%-------------------------------------------
\subsection{Setup}

\begin{verbatim}
# Compilar circuito
circom circuit.circom --r1cs --wasm --sym

# Powers of Tau (fase 1)
snarkjs powersoftau new bn128 12 pot12_0.ptau
snarkjs powersoftau contribute pot12_0.ptau pot12_1.ptau
snarkjs powersoftau prepare phase2 pot12_1.ptau pot12_final.ptau

# Setup del circuito (fase 2)
snarkjs groth16 setup circuit.r1cs pot12_final.ptau circuit.zkey
\end{verbatim}

%-------------------------------------------
\subsubsection{Paso 1: Compilar el Circuito}

\begin{mathnote}
\textbf{Comando de compilación:}

\begin{verbatim}
circom circuit.circom --r1cs --wasm --sym --c --O2
\end{verbatim}

\textbf{Flags:}
\begin{itemize}
    \item \texttt{--r1cs}: Genera el sistema R1CS (.r1cs)
    \item \texttt{--wasm}: Genera el witness generator en WebAssembly
    \item \texttt{--sym}: Genera archivo de símbolos (debug)
    \item \texttt{--c}: También genera código C++ (más rápido)
    \item \texttt{--O0/O1/O2}: Nivel de optimización
\end{itemize}

\textbf{Archivos generados:}
\begin{verbatim}
circuit.r1cs          # Sistema de constraints
circuit_js/
├── circuit.wasm      # Witness generator (WASM)
├── generate_witness.js
└── witness_calculator.js
circuit.sym           # Mapeo de símbolos
circuit_cpp/          # (si usaste --c)
└── ...
\end{verbatim}
\end{mathnote}

\begin{example}[Información del R1CS]
\begin{verbatim}
$ snarkjs r1cs info circuit.r1cs

[INFO]  snarkjs: Curve: bn-128
[INFO]  snarkjs: # of Wires: 1234
[INFO]  snarkjs: # of Constraints: 1000
[INFO]  snarkjs: # of Private Inputs: 5
[INFO]  snarkjs: # of Public Inputs: 3
[INFO]  snarkjs: # of Labels: 1234
[INFO]  snarkjs: # of Outputs: 2
\end{verbatim}

\textbf{Interpretación:}
\begin{itemize}
    \item \textbf{Wires:} Total de señales en el circuito
    \item \textbf{Constraints:} Número de ecuaciones R1CS (determina tiempo de proving)
    \item \textbf{Private Inputs:} Señales del witness
    \item \textbf{Public Inputs:} Señales conocidas por el verifier
\end{itemize}
\end{example}

%-------------------------------------------
\subsubsection{Paso 2: Powers of Tau (Fase 1)}

\begin{keyidea}
\textbf{¿Qué es Powers of Tau?}

Es la ceremonia de trusted setup universal:
\begin{itemize}
    \item Genera parámetros reutilizables para cualquier circuito
    \item Múltiples participantes contribuyen aleatoriedad
    \item El archivo .ptau es el resultado
\end{itemize}

\textbf{Tamaño del ceremony ($2^n$):}
\begin{center}
\begin{tabular}{ccc}
\toprule
\textbf{Power} & \textbf{Max constraints} & \textbf{Tamaño .ptau} \\
\midrule
12 & 4,096 & $\sim$4 MB \\
14 & 16,384 & $\sim$16 MB \\
16 & 65,536 & $\sim$64 MB \\
18 & 262,144 & $\sim$256 MB \\
20 & 1,048,576 & $\sim$1 GB \\
22 & 4,194,304 & $\sim$4 GB \\
\bottomrule
\end{tabular}
\end{center}

Usar el power mínimo necesario para tu circuito.
\end{keyidea}

\begin{verbatim}
# Crear ceremonia nueva (solo para testing)
snarkjs powersoftau new bn128 14 pot14_0.ptau -v

# Contribuir aleatoriedad
snarkjs powersoftau contribute pot14_0.ptau pot14_1.ptau \
    --name="First contribution" -v

# Más contribuciones (de diferentes personas)
snarkjs powersoftau contribute pot14_1.ptau pot14_2.ptau \
    --name="Second contribution" -v

# Aplicar random beacon (opcional, para finalizar)
snarkjs powersoftau beacon pot14_2.ptau pot14_beacon.ptau \
    0102030405060708 10 -n="Final Beacon"

# Preparar para fase 2
snarkjs powersoftau prepare phase2 pot14_beacon.ptau pot14_final.ptau -v
\end{verbatim}

\begin{remark}
\textbf{En producción, usa ceremonias existentes:}

Para no hacer tu propia ceremonia (que sería menos segura):

\begin{verbatim}
# Descargar .ptau de Hermez (trusted)
wget https://hermez.s3-eu-west-1.amazonaws.com/\
    powersOfTau28_hez_final_14.ptau
\end{verbatim}

Hermez realizó una ceremonia con 1000+ participantes.
\end{remark}

%-------------------------------------------
\subsubsection{Paso 3: Setup del Circuito (Fase 2)}

\begin{verbatim}
# Groth16 setup (per-circuit)
snarkjs groth16 setup circuit.r1cs pot14_final.ptau circuit_0.zkey

# Contribuir a la ceremonia de fase 2 (opcional pero recomendado)
snarkjs zkey contribute circuit_0.zkey circuit_1.zkey \
    --name="First phase2 contribution" -v

# Exportar verification key
snarkjs zkey export verificationkey circuit_1.zkey verification_key.json
\end{verbatim}

\begin{mathnote}
\textbf{Fase 1 vs Fase 2:}

\textbf{Fase 1 (Powers of Tau):}
\begin{itemize}
    \item Universal para todos los circuitos
    \item Solo depende del tamaño máximo
    \item Se hace una vez, se reutiliza
\end{itemize}

\textbf{Fase 2 (Circuit-specific):}
\begin{itemize}
    \item Específica para cada circuito
    \item Incorpora la estructura del R1CS
    \item Genera las proving/verification keys finales
\end{itemize}

\textbf{Archivos generados:}
\begin{itemize}
    \item \texttt{circuit.zkey}: Proving key (grande, para el prover)
    \item \texttt{verification\_key.json}: Verification key (pequeña, para el verifier)
\end{itemize}
\end{mathnote}

%-------------------------------------------
\subsection{Prove}

\begin{verbatim}
# Generar witness
node circuit_js/generate_witness.js circuit.wasm input.json witness.wtns

# Generar prueba
snarkjs groth16 prove circuit.zkey witness.wtns proof.json public.json
\end{verbatim}

%-------------------------------------------
\subsubsection{Paso 1: Generar el Witness}

\begin{example}[Input y witness generation]
\textbf{Archivo input.json:}
\begin{verbatim}
{
    "a": "3",
    "b": "5",
    "c": "15"
}
\end{verbatim}

\textbf{Generar witness:}
\begin{verbatim}
# Usando el WASM generator
node circuit_js/generate_witness.js \
    circuit_js/circuit.wasm \
    input.json \
    witness.wtns

# Alternativamente, usando C++ (más rápido)
./circuit_cpp/circuit input.json witness.wtns
\end{verbatim}

\textbf{Ver el witness:}
\begin{verbatim}
$ snarkjs wtns export json witness.wtns witness.json
$ cat witness.json
[
 "1",
 "15",   // output c (público)
 "3",    // input a
 "5"     // input b
]
\end{verbatim}
\end{example}

%-------------------------------------------
\subsubsection{Paso 2: Generar la Prueba}

\begin{verbatim}
snarkjs groth16 prove circuit.zkey witness.wtns proof.json public.json
\end{verbatim}

\begin{example}[Archivos de salida]
\textbf{proof.json:}
\begin{verbatim}
{
 "pi_a": [
  "12345678901234567890...",
  "98765432109876543210...",
  "1"
 ],
 "pi_b": [
  ["12345...", "67890..."],
  ["11111...", "22222..."],
  ["1", "0"]
 ],
 "pi_c": [
  "55555555555555555555...",
  "44444444444444444444...",
  "1"
 ],
 "protocol": "groth16",
 "curve": "bn128"
}
\end{verbatim}

\textbf{public.json:} (public inputs + outputs)
\begin{verbatim}
[
 "15"  // El valor público c
]
\end{verbatim}
\end{example}

\begin{mathnote}
\textbf{Estructura de la prueba Groth16:}

\begin{itemize}
    \item \texttt{pi\_a}: Punto en $\mathbb{G}_1$ (2 coordenadas + 1)
    \item \texttt{pi\_b}: Punto en $\mathbb{G}_2$ (2x2 coordenadas + 1)
    \item \texttt{pi\_c}: Punto en $\mathbb{G}_1$ (2 coordenadas + 1)
\end{itemize}

Tamaño total: $\sim$256 bytes sin comprimir, $\sim$128 bytes comprimido.

\textbf{Tiempo de proving:}
\begin{itemize}
    \item 1,000 constraints: $\sim$1 segundo
    \item 10,000 constraints: $\sim$5 segundos
    \item 100,000 constraints: $\sim$30 segundos
    \item 1,000,000 constraints: $\sim$5 minutos
\end{itemize}

(Tiempos aproximados en laptop moderna)
\end{mathnote}

%-------------------------------------------
\subsection{Verify}

\begin{verbatim}
# Exportar verification key
snarkjs zkey export verificationkey circuit.zkey vkey.json

# Verificar
snarkjs groth16 verify vkey.json public.json proof.json
\end{verbatim}

\begin{example}[Verificación]
\begin{verbatim}
$ snarkjs groth16 verify verification_key.json public.json proof.json
[INFO]  snarkjs: OK!

# Si la prueba es inválida:
$ snarkjs groth16 verify verification_key.json public.json bad_proof.json
[ERROR] snarkjs: Invalid proof
\end{verbatim}

\textbf{Tiempo de verificación:} $\sim$10-50 ms (independiente del tamaño del circuito).
\end{example}

%-------------------------------------------
\subsection{Exportar a Solidity}

\begin{verbatim}
snarkjs zkey export solidityverifier circuit.zkey Verifier.sol
\end{verbatim}

\begin{example}[Contrato verificador generado]
\begin{verbatim}
// SPDX-License-Identifier: GPL-3.0
pragma solidity >=0.7.0 <0.9.0;

contract Groth16Verifier {
    // Parámetros de la curva BN254
    uint256 constant PRIME_Q = 
        21888242871839275222246405745257275088696311157297823662689037894645226208583;
    
    // Verification key (hardcoded)
    uint256 constant alphax = ...;
    uint256 constant alphay = ...;
    // ... más constantes ...
    
    // IC (para public inputs)
    uint256[2][2] IC;
    
    function verifyProof(
        uint[2] calldata _pA,
        uint[2][2] calldata _pB,
        uint[2] calldata _pC,
        uint[1] calldata _pubSignals
    ) public view returns (bool) {
        // ... verificación con pairings ...
        return pairing(...);
    }
}
\end{verbatim}
\end{example}

\begin{keyidea}
\textbf{Usando el verificador en Solidity:}

\begin{verbatim}
// En tu contrato
contract MyApp {
    Groth16Verifier verifier;
    
    constructor(address _verifier) {
        verifier = Groth16Verifier(_verifier);
    }
    
    function doSomethingWithProof(
        uint[2] calldata a,
        uint[2][2] calldata b,
        uint[2] calldata c,
        uint[1] calldata publicInputs
    ) external {
        require(verifier.verifyProof(a, b, c, publicInputs), 
                "Invalid proof");
        
        // La prueba es válida, proceder con la lógica
        uint result = publicInputs[0];
        // ...
    }
}
\end{verbatim}

\textbf{Gas cost:} $\sim$200,000-300,000 gas por verificación.
\end{keyidea}

%-------------------------------------------
\subsection{Generar Calldata}

\begin{verbatim}
# Generar calldata para llamar al contrato
snarkjs zkey export soliditycalldata public.json proof.json
\end{verbatim}

\begin{example}[Calldata generado]
\begin{verbatim}
$ snarkjs zkey export soliditycalldata public.json proof.json

["0x12345...", "0x67890..."],
[["0x11111...", "0x22222..."], ["0x33333...", "0x44444..."]],
["0x55555...", "0x66666..."],
["0x0000000000000000000000000000000000000000000000000000000000000001"]
\end{verbatim}

Este output se puede usar directamente en Etherscan o en scripts de deployment.
\end{example}

%-------------------------------------------
\subsection{Script Completo de Ejemplo}

\begin{example}[Script de desarrollo end-to-end]
\begin{verbatim}
#!/bin/bash
set -e

CIRCUIT=multiplier

echo "=== 1. Compilar circuito ==="
circom ${CIRCUIT}.circom --r1cs --wasm --sym --c

echo "=== 2. Información del circuito ==="
snarkjs r1cs info ${CIRCUIT}.r1cs

echo "=== 3. Descargar Powers of Tau ==="
if [ ! -f pot14_final.ptau ]; then
    wget https://hermez.s3-eu-west-1.amazonaws.com/\
        powersOfTau28_hez_final_14.ptau \
        -O pot14_final.ptau
fi

echo "=== 4. Setup (Groth16) ==="
snarkjs groth16 setup ${CIRCUIT}.r1cs pot14_final.ptau ${CIRCUIT}_0.zkey

echo "=== 5. Contribuir a fase 2 ==="
snarkjs zkey contribute ${CIRCUIT}_0.zkey ${CIRCUIT}_final.zkey \
    --name="Dev contribution" -v -e="random entropy"

echo "=== 6. Exportar verification key ==="
snarkjs zkey export verificationkey ${CIRCUIT}_final.zkey vkey.json

echo "=== 7. Generar witness ==="
node ${CIRCUIT}_js/generate_witness.js \
    ${CIRCUIT}_js/${CIRCUIT}.wasm \
    input.json \
    witness.wtns

echo "=== 8. Generar prueba ==="
snarkjs groth16 prove ${CIRCUIT}_final.zkey witness.wtns proof.json public.json

echo "=== 9. Verificar prueba ==="
snarkjs groth16 verify vkey.json public.json proof.json

echo "=== 10. Exportar verificador Solidity ==="
snarkjs zkey export solidityverifier ${CIRCUIT}_final.zkey Verifier.sol

echo "=== 11. Generar calldata ==="
snarkjs zkey export soliditycalldata public.json proof.json > calldata.txt

echo "¡Completado!"
\end{verbatim}
\end{example}

%-------------------------------------------
\subsection{Otros Backends: PLONK y FFLONK}

\begin{remark}
\textbf{snarkjs soporta múltiples sistemas:}

\textbf{Groth16:}
\begin{verbatim}
snarkjs groth16 setup ...
snarkjs groth16 prove ...
snarkjs groth16 verify ...
\end{verbatim}

\textbf{PLONK:}
\begin{verbatim}
snarkjs plonk setup circuit.r1cs pot_final.ptau circuit.zkey
snarkjs plonk prove circuit.zkey witness.wtns proof.json public.json
snarkjs plonk verify vkey.json public.json proof.json
\end{verbatim}

\textbf{FFLONK (más nuevo, más eficiente):}
\begin{verbatim}
snarkjs fflonk setup circuit.r1cs pot_final.ptau circuit.zkey
snarkjs fflonk prove circuit.zkey witness.wtns proof.json public.json
snarkjs fflonk verify vkey.json public.json proof.json
\end{verbatim}

\begin{center}
\begin{tabular}{lccc}
\toprule
\textbf{Sistema} & \textbf{Setup} & \textbf{Prueba} & \textbf{Verify gas} \\
\midrule
Groth16 & Per-circuit & 128 B & $\sim$230K \\
PLONK & Universal & 768 B & $\sim$300K \\
FFLONK & Universal & 576 B & $\sim$250K \\
\bottomrule
\end{tabular}
\end{center}
\end{remark}

\begin{bancozk}
\textbf{Flujo de desarrollo en Banco ZK:}

\begin{verbatim}
# Nuestro Makefile
.PHONY: build prove verify deploy

PTAU=powersOfTau28_hez_final_20.ptau

build:
    circom src/transfer.circom -o build/ --r1cs --wasm --sym
    snarkjs groth16 setup build/transfer.r1cs $(PTAU) \
        build/transfer.zkey
    snarkjs zkey export verificationkey build/transfer.zkey \
        build/vkey.json
    snarkjs zkey export solidityverifier build/transfer.zkey \
        contracts/TransferVerifier.sol

test:
    npm run test:circuits

prove:
    node scripts/prove.js $(INPUT)

deploy:
    forge script script/Deploy.s.sol --broadcast
\end{verbatim}

\textbf{Optimizaciones en producción:}
\begin{itemize}
    \item Usamos rapidsnark (C++) en lugar de snarkjs para proving
    \item Witness generation con código C++ compilado
    \item Verificador Solidity optimizado para gas
    \item Batch proving para múltiples transacciones
\end{itemize}
\end{bancozk}

% ============================================
% BLOQUE 5: CIRCUITOS ZK - VERSIÓN EXPANDIDA
% Parte 4: Halo2, PLONKish, Custom Gates, Lookups
% ============================================

%===========================================
\section{Halo2}
%===========================================

\begin{definition}[Halo2]
Framework ZK basado en PLONKish arithmetization.
\end{definition}

%-------------------------------------------
\subsection{Historia y Contexto}

\begin{remark}
\textbf{Evolución de Halo a Halo2:}

\begin{center}
\begin{tabular}{lll}
\toprule
\textbf{Año} & \textbf{Versión} & \textbf{Innovación} \\
\midrule
2019 & Halo & Recursión sin trusted setup (ECC) \\
2020 & Halo 2 & PLONKish + lookups + mejor API \\
2021 & Orchard & Halo2 en Zcash mainnet \\
2022 & Scroll & Halo2 para zkEVM \\
2023 & PSE fork & Mejoras de la comunidad \\
\bottomrule
\end{tabular}
\end{center}

\textbf{Proyectos usando Halo2:}
\begin{itemize}
    \item Zcash (Orchard shielded pool)
    \item Scroll (zkEVM)
    \item Axiom (ZK coprocessor)
    \item Taiko (Type-1 zkEVM)
    \item zkEVM Community Edition
\end{itemize}
\end{remark}

%-------------------------------------------
\subsection{PLONKish Arithmetization}

En lugar de R1CS, usa tablas con columnas tipadas:
\begin{itemize}
    \item \textbf{Advice:} Witness privado
    \item \textbf{Instance:} Inputs públicos
    \item \textbf{Fixed:} Constantes del circuito
    \item \textbf{Selector:} Activan/desactivan gates
\end{itemize}

\begin{intuicion}
\textbf{De R1CS a PLONKish:}

\textbf{R1CS piensa en constraints individuales:}
\begin{verbatim}
Constraint 1: a * b = c
Constraint 2: c * d = e
Constraint 3: e + f = g
\end{verbatim}

\textbf{PLONKish piensa en una tabla:}
\begin{verbatim}
| row | a   | b   | c   | selector |
|-----|-----|-----|-----|----------|
| 0   | a   | b   | c   | mul      |
| 1   | c   | d   | e   | mul      |
| 2   | e   | f   | g   | add      |
\end{verbatim}

La tabla tiene:
\begin{itemize}
    \item \textbf{Columnas:} Variables del circuito (advice)
    \item \textbf{Filas:} Diferentes ``pasos'' o gates
    \item \textbf{Selectores:} Qué tipo de gate aplica en cada fila
\end{itemize}
\end{intuicion}

\begin{mathnote}
\textbf{Estructura de una tabla Halo2:}

\begin{center}
\begin{tabular}{|c|c|c|c|c|c|c|}
\hline
\textbf{row} & \textbf{a (advice)} & \textbf{b (advice)} & \textbf{c (advice)} & \textbf{s\_mul} & \textbf{s\_add} & \textbf{const (fixed)} \\
\hline
0 & 3 & 4 & 12 & 1 & 0 & 0 \\
1 & 12 & 5 & 60 & 1 & 0 & 0 \\
2 & 60 & 7 & 67 & 0 & 1 & 0 \\
3 & 67 & 0 & 67 & 0 & 0 & 0 \\
\hline
\end{tabular}
\end{center}

\textbf{Tipos de columnas:}
\begin{itemize}
    \item \textbf{Advice:} Valores del witness (privados)
    \item \textbf{Instance:} Inputs públicos (conocidos por verifier)
    \item \textbf{Fixed:} Constantes compiladas en el circuito
    \item \textbf{Selector:} Columnas fixed que activan gates
\end{itemize}

\textbf{Gates:}
\begin{itemize}
    \item Row 0: $s_{mul} \cdot (a \cdot b - c) = 0 \Rightarrow 3 \cdot 4 = 12$ ✓
    \item Row 1: $s_{mul} \cdot (a \cdot b - c) = 0 \Rightarrow 12 \cdot 5 = 60$ ✓
    \item Row 2: $s_{add} \cdot (a + b - c) = 0 \Rightarrow 60 + 7 = 67$ ✓
\end{itemize}
\end{mathnote}

%-------------------------------------------
\subsection{Custom Gates}

\begin{definition}[Custom Gate]
Expresión polinomial arbitraria sobre columnas.
\end{definition}

\begin{example}
Gate para Poseidon S-box ($x^5$):
\[
s(X) \cdot (a(X)^5 - b(X)) = 0
\]
Cuando selector $s$ activo, verifica $a^5 = b$.
\end{example}

\begin{keyidea}
\textbf{Ventaja de custom gates:}

\textbf{En R1CS:} Solo puedes hacer $a \cdot b = c$ (una multiplicación por constraint).

Para calcular $x^5$:
\begin{verbatim}
x2 = x * x      // constraint 1
x3 = x2 * x     // constraint 2
x4 = x3 * x     // constraint 3
x5 = x4 * x     // constraint 4
// Total: 4 constraints
\end{verbatim}

\textbf{En PLONKish:} Puedes definir gate de grado 5:
\begin{verbatim}
s_pow5 * (a^5 - b) = 0  // 1 constraint
// Total: 1 constraint (grado más alto)
\end{verbatim}

\textbf{Trade-off:}
\begin{itemize}
    \item Custom gates reducen número de rows
    \item Pero aumentan el grado del polinomio
    \item Grado más alto = prover más lento
    \item Balance óptimo depende del circuito
\end{itemize}
\end{keyidea}

\begin{example}[Custom gates para operaciones comunes]
\begin{verbatim}
// Gate de multiplicación estándar
s_mul * (a * b - c) = 0

// Gate de suma
s_add * (a + b - c) = 0

// Gate de multiplicación por constante
s_mul_const * (a * constant - b) = 0

// Gate bool (verifica que es 0 o 1)
s_bool * a * (1 - a) = 0

// Gate ternario (if-then-else)
// if cond then a else b = out
s_ternary * (cond * (a - b) + b - out) = 0

// Gate Poseidon full round
s_poseidon * (
    (state[0] + round_const)^5 - new_state[0] +
    matrix_multiply(...)
) = 0
\end{verbatim}
\end{example}

%-------------------------------------------
\subsection{Lookup Tables}

\begin{definition}[Lookup]
Verificar que valores están en una tabla predefinida.
\end{definition}

Usos:
\begin{itemize}
    \item Range checks (valor en $[0, 2^{16})$)
    \item Tablas XOR para operaciones de bits
    \item S-boxes para criptografía simétrica
\end{itemize}

\begin{intuicion}
\textbf{¿Por qué lookups?}

Algunas operaciones son difíciles de expresar algebraicamente:
\begin{itemize}
    \item XOR de bits
    \item Comparaciones
    \item Funciones no-algebraicas
\end{itemize}

\textbf{Solución:} Precalcular todos los resultados en una tabla y verificar que el valor está en la tabla.

\textbf{Ejemplo XOR de 4 bits:}
\begin{verbatim}
Tabla precalculada (256 entradas):
| a | b | a XOR b |
|---|---|---------|
| 0 | 0 |    0    |
| 0 | 1 |    1    |
| 1 | 0 |    1    |
| 1 | 1 |    0    |
| 0 | 2 |    2    |
| ...
\end{verbatim}

En lugar de 4 constraints por XOR (descomponer en bits), un lookup es $\sim$1 constraint amortizado.
\end{intuicion}

\begin{mathnote}
\textbf{Cómo funcionan los lookups (Plookup):}

\textbf{Setup:}
\begin{itemize}
    \item Tabla $T$ con $n$ filas de valores válidos
    \item La tabla es parte del circuito (fixed)
\end{itemize}

\textbf{Claim:}
\begin{itemize}
    \item El prover dice ``mis valores $(a, b, c)$ están todos en $T$''
\end{itemize}

\textbf{Verificación:}

Usando grand product argument:
\[ \prod_i \frac{(\gamma + f_i)}{(\gamma + t_i)} = 1 \]

Donde $f_i$ son los valores a verificar y $t_i$ son la tabla más valores repetidos.

\textbf{Costo:}
\begin{itemize}
    \item Cada lookup: $\sim$1 constraint adicional
    \item Overhead por tabla: polinomios adicionales
    \item Muy eficiente para muchos lookups a la misma tabla
\end{itemize}
\end{mathnote}

\begin{example}[Range check con lookup]
\begin{verbatim}
// Sin lookup (para 16 bits):
// Descomponer en bits: 16 constraints
// Verificar cada bit es 0/1: 16 constraints
// Total: 32 constraints

// Con lookup:
// Tabla: {0, 1, 2, ..., 65535}
// Lookup: 1 constraint
// Total: ~1 constraint

// En Halo2:
meta.lookup(|meta| {
    let value = meta.query_advice(a, Rotation::cur());
    let table = meta.query_fixed(range_table, Rotation::cur());
    
    vec![(value, table)]
});
\end{verbatim}
\end{example}

%-------------------------------------------
\subsection{API de Halo2 (Rust)}

\begin{example}[Circuito simple en Halo2]
\begin{verbatim}
use halo2_proofs::{
    circuit::{Layouter, SimpleFloorPlanner, Value},
    plonk::{Advice, Circuit, Column, ConstraintSystem, Error, Selector},
    poly::Rotation,
};

#[derive(Clone)]
struct MyConfig {
    advice: [Column<Advice>; 2],
    selector: Selector,
}

struct MyCircuit {
    a: Value<Fp>,
    b: Value<Fp>,
}

impl Circuit<Fp> for MyCircuit {
    type Config = MyConfig;
    type FloorPlanner = SimpleFloorPlanner;

    fn configure(meta: &mut ConstraintSystem<Fp>) -> Self::Config {
        let advice = [meta.advice_column(), meta.advice_column()];
        let selector = meta.selector();

        // Definir el gate de multiplicación
        meta.create_gate("mul", |meta| {
            let s = meta.query_selector(selector);
            let a = meta.query_advice(advice[0], Rotation::cur());
            let b = meta.query_advice(advice[1], Rotation::cur());
            let c = meta.query_advice(advice[0], Rotation::next());

            // s * (a * b - c) = 0
            vec![s * (a * b - c)]
        });

        MyConfig { advice, selector }
    }

    fn synthesize(
        &self,
        config: Self::Config,
        mut layouter: impl Layouter<Fp>,
    ) -> Result<(), Error> {
        layouter.assign_region(
            || "multiplication",
            |mut region| {
                config.selector.enable(&mut region, 0)?;
                
                region.assign_advice(
                    || "a",
                    config.advice[0],
                    0,
                    || self.a,
                )?;
                region.assign_advice(
                    || "b",
                    config.advice[1],
                    0,
                    || self.b,
                )?;
                region.assign_advice(
                    || "c = a * b",
                    config.advice[0],
                    1,
                    || self.a * self.b,
                )?;
                
                Ok(())
            },
        )
    }
}
\end{verbatim}
\end{example}

\begin{keyidea}
\textbf{Componentes de un circuito Halo2:}

\textbf{1. Config:} Define la estructura del circuito
\begin{itemize}
    \item Columnas (advice, fixed, instance)
    \item Selectores
    \item Gates y lookups
\end{itemize}

\textbf{2. configure():} Crea los gates
\begin{itemize}
    \item Define expresiones polinomiales
    \item Configura lookups
    \item Se ejecuta una vez al compilar
\end{itemize}

\textbf{3. synthesize():} Asigna valores
\begin{itemize}
    \item Llena las columnas con valores concretos
    \item Habilita selectores en filas específicas
    \item Se ejecuta para cada prueba
\end{itemize}

\textbf{4. FloorPlanner:} Organiza el layout
\begin{itemize}
    \item SimpleFloorPlanner: básico, secuencial
    \item V1: más optimizado
\end{itemize}
\end{keyidea}

%-------------------------------------------
\subsection{Chips y Regiones}

\begin{definition}[Chip]
Módulo reutilizable que implementa una funcionalidad específica.
\end{definition}

\begin{example}[Estructura de chips]
\begin{verbatim}
// Chip para operaciones de campo
struct FieldChip<F: FieldExt> {
    config: FieldConfig,
    _marker: PhantomData<F>,
}

impl<F: FieldExt> FieldChip<F> {
    fn mul(
        &self,
        layouter: &mut impl Layouter<F>,
        a: AssignedCell<F, F>,
        b: AssignedCell<F, F>,
    ) -> Result<AssignedCell<F, F>, Error> {
        layouter.assign_region(
            || "mul",
            |mut region| {
                self.config.selector.enable(&mut region, 0)?;
                a.copy_advice(|| "a", &mut region, self.config.advice[0], 0)?;
                b.copy_advice(|| "b", &mut region, self.config.advice[1], 0)?;
                
                let c = a.value().copied() * b.value();
                region.assign_advice(
                    || "a * b",
                    self.config.advice[2],
                    0,
                    || c,
                )
            },
        )
    }
}

// Uso:
let field_chip = FieldChip::construct(config);
let a = field_chip.load_private(layouter, self.a)?;
let b = field_chip.load_private(layouter, self.b)?;
let c = field_chip.mul(layouter, a, b)?;
\end{verbatim}
\end{example}

%-------------------------------------------
\subsection{Halo2 vs Circom}

\begin{center}
\begin{tabular}{lcc}
\toprule
\textbf{Aspecto} & \textbf{Circom} & \textbf{Halo2} \\
\midrule
Lenguaje & DSL propio & Rust \\
Arithmetization & R1CS & PLONKish \\
Custom gates & No & Sí \\
Lookups & No nativo & Sí \\
Curva learning & Baja & Alta \\
Flexibilidad & Media & Alta \\
Tooling & Maduro & En desarrollo \\
Performance & Buena & Muy buena \\
Trusted setup & Sí (Groth16) & No (IPA) \\
\bottomrule
\end{tabular}
\end{center}

\begin{remark}
\textbf{¿Cuándo usar cada uno?}

\textbf{Usar Circom si:}
\begin{itemize}
    \item Prototipando rápido
    \item Circuito relativamente simple
    \item Quieres ecosistema maduro (snarkjs, circomlib)
    \item Team no muy experimentado en ZK
\end{itemize}

\textbf{Usar Halo2 si:}
\begin{itemize}
    \item Circuito muy grande/complejo (zkEVM)
    \item Necesitas custom gates para eficiencia
    \item Necesitas lookups extensivamente
    \item No quieres trusted setup
    \item Team con experiencia en Rust/ZK
\end{itemize}
\end{remark}

\begin{bancozk}
\textbf{Elección en Banco ZK:}

\textbf{Actualmente: Circom + Groth16}
\begin{itemize}
    \item Desarrollo más rápido
    \item Circuitos de tamaño manejable ($\sim$10K constraints)
    \item Verificación en Ethereum es crítica (Groth16 más barato)
    \item Team familiar con el stack
\end{itemize}

\textbf{Roadmap: Migración gradual a Halo2}
\begin{itemize}
    \item Para circuitos nuevos más complejos
    \item Eliminar dependencia de trusted setup
    \item Mejor performance con custom gates para Poseidon
    \item Lookups para range checks eficientes
\end{itemize}

\textbf{Híbrido potencial:}
\begin{itemize}
    \item Circuitos de usuario: Circom (simplicidad)
    \item Agregador/batch prover: Halo2 (performance)
    \item Wrapper final: Groth16 (verificación barata en L1)
\end{itemize}
\end{bancozk}

% ============================================
% BLOQUE 5: CIRCUITOS ZK - VERSIÓN EXPANDIDA
% Parte 5: Arkworks, Noir, y Resumen
% ============================================

%===========================================
\section{Arkworks}
%===========================================

\begin{definition}[Arkworks]
Ecosistema Rust de bibliotecas para criptografía ZK.
\end{definition}

Componentes:
\begin{itemize}
    \item \textbf{ark-ff:} Aritmética de campos finitos
    \item \textbf{ark-ec:} Curvas elípticas
    \item \textbf{ark-poly:} Polinomios
    \item \textbf{ark-relations:} R1CS
    \item \textbf{ark-snark:} Interfaces de sistemas
    \item \textbf{ark-groth16:} Implementación Groth16
    \item \textbf{ark-marlin:} Sistema universal
\end{itemize}

Uso: investigación, implementaciones personalizadas.

%-------------------------------------------
\subsection{Arquitectura del Ecosistema}

\begin{center}
\begin{tikzpicture}[
    box/.style={draw, rectangle, minimum width=2.5cm, minimum height=0.8cm, align=center, fill=blue!10},
    layer/.style={draw, rectangle, minimum width=12cm, minimum height=1.2cm, align=center, fill=gray!10},
]
    % Layers
    \node[layer] (app) at (0, 4) {\textbf{Aplicaciones:} tu código, circuitos específicos};
    \node[layer] (snark) at (0, 2.5) {\textbf{Sistemas SNARK:} ark-groth16, ark-marlin, ark-gm17};
    \node[layer] (rel) at (0, 1) {\textbf{Relaciones:} ark-relations (R1CS), ark-r1cs-std};
    \node[layer] (crypto) at (0, -0.5) {\textbf{Cripto primitivas:} ark-ec (curvas), ark-poly (polinomios)};
    \node[layer] (math) at (0, -2) {\textbf{Matemáticas:} ark-ff (campos finitos), ark-serialize};
    
    % Arrows
    \draw[->, thick] (app) -- (snark);
    \draw[->, thick] (snark) -- (rel);
    \draw[->, thick] (rel) -- (crypto);
    \draw[->, thick] (crypto) -- (math);
\end{tikzpicture}
\end{center}

\begin{remark}
\textbf{Historia de Arkworks:}

\begin{itemize}
    \item Originalmente ``zexe'' de la Universidad de Berkeley
    \item Renombrado a Arkworks en 2021
    \item Mantenido por comunidad + academia
    \item Base de muchos proyectos de investigación
    \item Usado por Aleo, Mina, Penumbra, y otros
\end{itemize}

\textbf{Filosofía:} Bibliotecas modulares y componibles, no un framework monolítico.
\end{remark}

%-------------------------------------------
\subsection{ark-ff: Campos Finitos}

\begin{example}[Definiendo y usando campos]
\begin{verbatim}
use ark_ff::{Field, PrimeField, BigInteger};
use ark_bn254::Fr;  // Campo escalar de BN254

fn field_operations() {
    // Crear elementos
    let a = Fr::from(5u64);
    let b = Fr::from(7u64);
    
    // Operaciones básicas
    let sum = a + b;           // 12
    let product = a * b;       // 35
    let diff = a - b;          // p - 2 (wraps around)
    
    // Inverso
    let inv = a.inverse().unwrap();  // a^(-1) mod p
    assert_eq!(a * inv, Fr::one());
    
    // Exponenciación
    let pow = a.pow([3u64]);   // a^3 = 125
    
    // Raíz cuadrada (si existe)
    let sqrt = a.sqrt();
    
    // Serialización
    let bytes = a.into_repr().to_bytes_le();
}
\end{verbatim}
\end{example}

%-------------------------------------------
\subsection{ark-relations: R1CS}

\begin{example}[Definiendo un circuito R1CS]
\begin{verbatim}
use ark_relations::{
    lc,
    r1cs::{ConstraintSynthesizer, ConstraintSystemRef, SynthesisError},
};
use ark_ff::Field;

// Circuito: probar que conocemos x tal que x^3 + x + 5 = y
pub struct CubicCircuit<F: Field> {
    pub x: Option<F>,  // Witness (privado)
    pub y: F,          // Público
}

impl<F: Field> ConstraintSynthesizer<F> for CubicCircuit<F> {
    fn generate_constraints(
        self,
        cs: ConstraintSystemRef<F>,
    ) -> Result<(), SynthesisError> {
        // Allocate x (private witness)
        let x = cs.new_witness_variable(|| {
            self.x.ok_or(SynthesisError::AssignmentMissing)
        })?;
        
        // Allocate y (public input)
        let y = cs.new_input_variable(|| Ok(self.y))?;
        
        // x^2
        let x_squared = cs.new_witness_variable(|| {
            let x_val = self.x.ok_or(SynthesisError::AssignmentMissing)?;
            Ok(x_val * x_val)
        })?;
        
        // Constraint: x * x = x_squared
        cs.enforce_constraint(lc!() + x, lc!() + x, lc!() + x_squared)?;
        
        // x^3
        let x_cubed = cs.new_witness_variable(|| {
            let x_val = self.x.ok_or(SynthesisError::AssignmentMissing)?;
            Ok(x_val * x_val * x_val)
        })?;
        
        // Constraint: x_squared * x = x_cubed
        cs.enforce_constraint(lc!() + x_squared, lc!() + x, lc!() + x_cubed)?;
        
        // Constraint: x_cubed + x + 5 = y
        let five = F::from(5u64);
        cs.enforce_constraint(
            lc!() + x_cubed + x + (five, CS::one()),
            lc!() + CS::one(),
            lc!() + y,
        )?;
        
        Ok(())
    }
}
\end{verbatim}
\end{example}

%-------------------------------------------
\subsection{ark-groth16: Generando Pruebas}

\begin{example}[Flujo completo con Groth16]
\begin{verbatim}
use ark_groth16::{Groth16, ProvingKey, VerifyingKey};
use ark_bn254::{Bn254, Fr};
use ark_snark::SNARK;
use ark_std::rand::thread_rng;

fn prove_and_verify() {
    let mut rng = thread_rng();
    
    // 1. Crear circuito de ejemplo para setup
    let circuit = CubicCircuit::<Fr> {
        x: None,  // No necesitamos valores para setup
        y: Fr::from(35u64),
    };
    
    // 2. Setup (trusted setup)
    let (pk, vk) = Groth16::<Bn254>::circuit_specific_setup(
        circuit.clone(),
        &mut rng,
    ).unwrap();
    
    // 3. Crear circuito con witness real
    let circuit_with_witness = CubicCircuit {
        x: Some(Fr::from(3u64)),  // x = 3
        y: Fr::from(35u64),       // y = 3^3 + 3 + 5 = 35
    };
    
    // 4. Generar prueba
    let proof = Groth16::<Bn254>::prove(
        &pk,
        circuit_with_witness,
        &mut rng,
    ).unwrap();
    
    // 5. Preparar public inputs
    let public_inputs = vec![Fr::from(35u64)];
    
    // 6. Verificar
    let valid = Groth16::<Bn254>::verify(
        &vk,
        &public_inputs,
        &proof,
    ).unwrap();
    
    assert!(valid);
    println!("Proof verified successfully!");
}
\end{verbatim}
\end{example}

%-------------------------------------------
\subsection{ark-r1cs-std: Gadgets}

\begin{mathnote}
\textbf{Gadgets estándar disponibles:}

\begin{verbatim}
use ark_r1cs_std::{
    // Tipos básicos
    fields::fp::FpVar,
    bits::boolean::Boolean,
    bits::uint8::UInt8,
    
    // Comparaciones
    eq::EqGadget,
    cmp::CmpGadget,
    
    // Curvas elípticas
    groups::CurveVar,
    
    // Hashes
    // (en ark-crypto-primitives)
};

// Ejemplo: comparar dos valores
let a = FpVar::new_witness(cs.clone(), || Ok(Fr::from(5)))?;
let b = FpVar::new_witness(cs.clone(), || Ok(Fr::from(3)))?;

// a == b ?
let is_eq = a.is_eq(&b)?;
is_eq.enforce_equal(&Boolean::FALSE)?;

// a > b ?
let is_greater = a.is_gt(&b)?;
is_greater.enforce_equal(&Boolean::TRUE)?;
\end{verbatim}
\end{mathnote}

%-------------------------------------------
\subsection{Cuándo Usar Arkworks}

\begin{keyidea}
\textbf{Arkworks es ideal para:}

\begin{itemize}
    \item \textbf{Investigación:} Implementar nuevos esquemas, papers
    \item \textbf{Customización:} Cuando necesitas control total
    \item \textbf{Performance:} Optimizaciones específicas
    \item \textbf{Nuevas curvas:} Experimentar con diferentes curvas
    \item \textbf{Integración:} Combinar múltiples primitivas
\end{itemize}

\textbf{NO es ideal para:}

\begin{itemize}
    \item Prototipado rápido (usar Circom)
    \item Circuitos simples (overhead de Rust)
    \item Teams sin experiencia en Rust/crypto
\end{itemize}
\end{keyidea}

%===========================================
\section{Noir}
%===========================================

\begin{definition}[Noir]
Lenguaje de alto nivel de Aztec para circuitos ZK.
\end{definition}

\begin{verbatim}
fn main(x: Field, y: pub Field) {
    assert(x * x == y);
}
\end{verbatim}

Características:
\begin{itemize}
    \item Sintaxis similar a Rust
    \item Compila a ACIR (Abstract Circuit IR)
    \item Backend-agnostic
    \item Ideal para private smart contracts
\end{itemize}

%-------------------------------------------
\subsection{Historia y Filosofía}

\begin{remark}
\textbf{Evolución de Noir:}

\begin{center}
\begin{tabular}{lll}
\toprule
\textbf{Año} & \textbf{Versión} & \textbf{Desarrollo} \\
\midrule
2021 & Concepto & Aztec diseña lenguaje de alto nivel \\
2022 & 0.1-0.5 & Primera versión pública \\
2023 & 0.6-0.19 & Mejoras significativas \\
2024 & 0.20+ & Producción-ready \\
\bottomrule
\end{tabular}
\end{center}

\textbf{Filosofía:}
\begin{itemize}
    \item ``ZK para todos'': sintaxis familiar
    \item Abstracción del backend: escribir una vez, compilar a cualquier sistema
    \item Seguridad: el compilador previene errores comunes
    \item Performance: optimizaciones automáticas
\end{itemize}
\end{remark}

%-------------------------------------------
\subsection{Sintaxis Básica}

\begin{example}[Programa Noir completo]
\begin{verbatim}
// main.nr

// Función principal - define el circuito
fn main(
    // Inputs privados (witness)
    secret: Field,
    password: [u8; 32],
    
    // Inputs públicos (prefijo 'pub')
    pub expected_hash: Field,
    pub min_value: Field,
) {
    // Computar hash del secreto
    let computed_hash = std::hash::pedersen([secret]);
    
    // Verificar igualdad (constraint)
    assert(computed_hash == expected_hash);
    
    // Verificar rango
    assert(secret >= min_value);
    
    // Operaciones con arrays
    let mut sum: Field = 0;
    for i in 0..32 {
        sum += password[i] as Field;
    }
    assert(sum > 0);
}

// Función auxiliar (no genera constraints directamente)
fn helper_function(x: Field) -> Field {
    x * x + 1
}

// Tests
#[test]
fn test_main() {
    main(5, [1; 32], pedersen([5]), 3);
}
\end{verbatim}
\end{example}

%-------------------------------------------
\subsection{Tipos de Datos}

\begin{mathnote}
\textbf{Tipos disponibles en Noir:}

\begin{verbatim}
// Campos finitos
let a: Field = 5;
let b: Field = -1;  // p - 1 en el campo

// Enteros
let u: u32 = 100;
let i: i8 = -50;
// Soporta: u1, u8, u16, u32, u64, i8, i16, i32, i64

// Booleanos
let flag: bool = true;

// Arrays (tamaño fijo)
let arr: [Field; 3] = [1, 2, 3];
let nested: [[u8; 2]; 3] = [[1, 2], [3, 4], [5, 6]];

// Tuplas
let tuple: (Field, bool) = (5, true);

// Structs
struct Point {
    x: Field,
    y: Field,
}

// Strings (tamaño fijo)
let s: str<5> = "hello";
\end{verbatim}

\textbf{Nota:} No hay tipos dinámicos (no hay Vec, no hay String dinámico). Todo debe ser de tamaño conocido en tiempo de compilación.
\end{mathnote}

%-------------------------------------------
\subsection{Control de Flujo}

\begin{example}[Control de flujo en Noir]
\begin{verbatim}
fn control_examples(x: Field, condition: bool) -> Field {
    // If-else (ambas ramas se computan, resultado se selecciona)
    let result = if condition {
        x * 2
    } else {
        x + 1
    };
    
    // For loops (bounds deben ser constantes)
    let mut sum: Field = 0;
    for i in 0..10 {
        sum += i as Field;
    }
    
    // While no existe - usar for con break
    // (en versiones recientes)
    
    // Match
    let value = match x {
        0 => 100,
        1 => 200,
        _ => 0,
    };
    
    result + sum + value
}

// IMPORTANTE: Los loops deben tener bounds conocidos en compile-time
// Esto NO compila:
// for i in 0..n { ... }  // Error: n no es constante
\end{verbatim}
\end{example}

%-------------------------------------------
\subsection{Biblioteca Estándar}

\begin{remark}
\textbf{Funciones estándar de Noir:}

\begin{verbatim}
use dep::std;

// Hashes
let h1 = std::hash::pedersen([a, b, c]);
let h2 = std::hash::poseidon::bn254::hash_2([a, b]);
let h3 = std::hash::keccak256([...]);
let h4 = std::hash::sha256([...]);

// Merkle trees
let root = std::merkle::compute_merkle_root(leaf, index, path);
let valid = std::merkle::check_membership(root, leaf, index, path);

// ECDSA / Schnorr
let valid = std::ecdsa_secp256k1::verify_signature(
    pub_key_x, pub_key_y, signature, message_hash
);
let valid = std::schnorr::verify_signature(
    pub_key_x, pub_key_y, signature, message
);

// EdDSA (BabyJubJub)
let valid = std::eddsa::eddsa_poseidon_verify(
    pub_key_x, pub_key_y, signature_s, signature_r8_x, signature_r8_y, message
);
\end{verbatim}
\end{remark}

%-------------------------------------------
\subsection{Toolchain de Noir}

\begin{example}[Flujo de desarrollo con Nargo]
\begin{verbatim}
# Instalar Noir
curl -L https://noir-lang.org/install.sh | bash
noirup

# Crear proyecto
nargo new my_circuit
cd my_circuit

# Estructura del proyecto
my_circuit/
├── Nargo.toml        # Configuración
├── src/
│   └── main.nr       # Código del circuito
├── Prover.toml       # Inputs para prover
└── Verifier.toml     # Inputs públicos

# Compilar
nargo compile

# Ejecutar (witness generation + check)
nargo execute

# Generar prueba
nargo prove

# Verificar
nargo verify

# Generar verificador Solidity
nargo codegen-verifier
\end{verbatim}
\end{example}

\begin{example}[Archivos de configuración]
\textbf{Nargo.toml:}
\begin{verbatim}
[package]
name = "my_circuit"
type = "bin"
authors = [""]
compiler_version = ">=0.25.0"

[dependencies]
# Dependencias de otras librerías Noir
\end{verbatim}

\textbf{Prover.toml:}
\begin{verbatim}
# Inputs para generar la prueba
secret = "0x1234..."
password = [1, 2, 3, ...]
expected_hash = "0xabcd..."
min_value = "5"
\end{verbatim}
\end{example}

%-------------------------------------------
\subsection{ACIR: Representación Intermedia}

\begin{keyidea}
\textbf{ACIR (Abstract Circuit Intermediate Representation):}

Noir compila a ACIR, no directamente a R1CS o PLONKish.

\begin{verbatim}
Noir code → ACIR → Backend-specific format
                      ├─→ Barretenberg (default)
                      ├─→ Arkworks
                      └─→ Otros backends
\end{verbatim}

\textbf{Ventajas:}
\begin{itemize}
    \item Un lenguaje, múltiples backends
    \item Optimizaciones a nivel de ACIR
    \item Portabilidad del código
\end{itemize}

\textbf{Backends disponibles:}
\begin{itemize}
    \item \textbf{Barretenberg:} UltraPlonk, default de Aztec
    \item \textbf{arkworks:} Groth16, Marlin
    \item \textbf{halo2:} En desarrollo
\end{itemize}
\end{keyidea}

%-------------------------------------------
\subsection{Comparación de Lenguajes ZK}

\begin{center}
\begin{tabular}{lcccc}
\toprule
\textbf{Aspecto} & \textbf{Circom} & \textbf{Halo2} & \textbf{Noir} & \textbf{Arkworks} \\
\midrule
Tipo & DSL & Framework & Lenguaje & Biblioteca \\
Lenguaje base & Propio & Rust & Rust-like & Rust \\
Curva learning & Baja & Alta & Media & Alta \\
Abstracción & Media & Baja & Alta & Baja \\
Flexibilidad & Media & Alta & Media & Muy alta \\
Backends & Groth16 & IPA & Múltiples & Múltiples \\
Custom gates & No & Sí & Vía backend & Sí \\
Lookups & No* & Sí & Sí & Depende \\
Madurez & Alta & Media & Media & Alta \\
\bottomrule
\end{tabular}
\end{center}

*Circom tiene extensiones con lookups en desarrollo.

%===========================================
\section{Resumen del Bloque 5}
%===========================================

\begin{keyidea}
\textbf{Puntos clave de Circuitos ZK:}

\textbf{1. Circuitos Aritméticos:}
\begin{itemize}
    \item DAGs de operaciones sobre campos finitos
    \item Solo suma y multiplicación como operaciones base
    \item Todo debe traducirse a aritmética de campo
\end{itemize}

\textbf{2. R1CS:}
\begin{itemize}
    \item Sistema de constraints: $(A \cdot w) \times (B \cdot w) = C \cdot w$
    \item Una multiplicación por constraint
    \item Sumas son ``gratis''
    \item Número de constraints $\approx$ complejidad del circuito
\end{itemize}

\textbf{3. Señales y Witness:}
\begin{itemize}
    \item Public inputs: conocidos por verifier
    \item Private inputs: solo conocidos por prover
    \item Witness: asignación de valores que satisface constraints
\end{itemize}

\textbf{4. Circom:}
\begin{itemize}
    \item DSL más usado para ZK en Ethereum
    \item Templates reutilizables, components
    \item Operadores: \texttt{<--}, \texttt{<==}, \texttt{===}
    \item Circomlib: biblioteca estándar
\end{itemize}

\textbf{5. snarkjs:}
\begin{itemize}
    \item Flujo: Compilar → Setup → Prove → Verify
    \item Powers of Tau: ceremonia universal
    \item Exporta verificadores a Solidity
\end{itemize}

\textbf{6. Halo2/PLONKish:}
\begin{itemize}
    \item Tablas en lugar de constraints individuales
    \item Custom gates: expresiones polinomiales arbitrarias
    \item Lookups: verificar pertenencia a tablas
    \item Sin trusted setup (con IPA)
\end{itemize}

\textbf{7. Otros frameworks:}
\begin{itemize}
    \item Arkworks: bibliotecas Rust modulares
    \item Noir: lenguaje de alto nivel, backend-agnostic
\end{itemize}

\textbf{8. Elección de herramientas:}
\begin{itemize}
    \item Prototipo rápido: Circom
    \item Circuitos complejos: Halo2
    \item Abstracción: Noir
    \item Investigación: Arkworks
\end{itemize}
\end{keyidea}

\begin{bancozk}
\textbf{Stack de Banco ZK:}

\begin{verbatim}
Capa de usuario:
├── Wallet → genera witness local
└── SDK → helpers para inputs

Capa de circuitos:
├── Circom (actual)
│   ├── transfer.circom
│   ├── deposit.circom
│   └── withdraw.circom
└── Noir (futuro)
    └── Mismos circuitos, más legibles

Capa de proving:
├── snarkjs (desarrollo)
└── rapidsnark (producción)

Capa de verificación:
├── Verifier.sol (L1)
└── On-chain verification (~200K gas)
\end{verbatim}

\textbf{Métricas objetivo:}
\begin{itemize}
    \item Transfer circuit: $<$15K constraints
    \item Prove time: $<$5 segundos (cliente)
    \item Verify time: $<$10 ms (off-chain)
    \item Verify gas: $<$250K gas (on-chain)
\end{itemize}
\end{bancozk}
\cleardoublepage


% ============================================
% APÉNDICE: NOTACIÓN MATEMÁTICA
% ============================================
\chapter*{Apéndice: Notación Matemática}
\addcontentsline{toc}{chapter}{Apéndice: Notación}

\section*{Conjuntos y Estructuras}
\begin{tabular}{ll}
$\mathbb{Z}$ & Números enteros \\
$\mathbb{Z}_n$ & Enteros módulo $n$ \\
$\mathbb{Z}_n^*$ & Grupo multiplicativo de $\mathbb{Z}_n$ \\
$\mathbb{F}_p$ & Campo finito de orden primo $p$ \\
$\mathbb{F}_{p^k}$ & Extensión de grado $k$ de $\mathbb{F}_p$ \\
$\mathbb{G}$ & Grupo genérico \\
$\mathbb{G}_1, \mathbb{G}_2, \mathbb{G}_T$ & Grupos de pairing \\
$E(\mathbb{F}_p)$ & Puntos de curva elíptica sobre $\mathbb{F}_p$ \\
\end{tabular}

\section*{Operaciones}
\begin{tabular}{ll}
$a \mod n$ & Residuo de $a$ módulo $n$ \\
$a \equiv b \pmod{n}$ & $a$ congruente a $b$ módulo $n$ \\
$g^x$ & Exponenciación en grupo (notación multiplicativa) \\
$[x]G$ & Multiplicación escalar en curva (notación aditiva) \\
$e(P, Q)$ & Pairing bilineal de $P$ y $Q$ \\
$H(m)$ & Hash de mensaje $m$ \\
$\|$ & Concatenación \\
\end{tabular}

\section*{Criptografía}
\begin{tabular}{ll}
$\mathsf{pk}, \mathsf{sk}$ & Clave pública, clave secreta \\
$\mathsf{Enc}_k(m)$ & Cifrado de $m$ con clave $k$ \\
$\mathsf{Dec}_k(c)$ & Descifrado de $c$ con clave $k$ \\
$\mathsf{Sign}_{\mathsf{sk}}(m)$ & Firma de $m$ con clave secreta \\
$\mathsf{Verify}_{\mathsf{pk}}(m, \sigma)$ & Verificación de firma \\
$\mathsf{Commit}(v; r)$ & Commitment de valor $v$ con randomness $r$ \\
$\mathsf{Open}(c, v, r)$ & Apertura de commitment \\
\end{tabular}

\section*{Zero-Knowledge}
\begin{tabular}{ll}
$\pi$ & Prueba (proof) \\
$x$ & Statement público \\
$w$ & Witness (testigo privado) \\
$\mathcal{R}$ & Relación: $(x, w) \in \mathcal{R}$ si $w$ es witness válido para $x$ \\
$\mathsf{Setup}(\lambda) \to \mathsf{crs}$ & Generación de parámetros \\
$\mathsf{Prove}(\mathsf{crs}, x, w) \to \pi$ & Generación de prueba \\
$\mathsf{Verify}(\mathsf{crs}, x, \pi) \to \{0, 1\}$ & Verificación de prueba \\
\end{tabular}

\section*{Identidad y Credenciales}
\begin{tabular}{ll}
$\mathsf{Issue}(\mathsf{sk}_I, \mathsf{attrs}) \to \mathsf{cred}$ & Emisión de credencial \\
$\mathsf{Present}(\mathsf{cred}, \mathsf{disclosed}) \to \mathsf{pres}$ & Presentación selectiva \\
$\mathsf{VerifyPres}(\mathsf{pk}_I, \mathsf{pres}) \to \{0, 1\}$ & Verificación de presentación \\
$\mathsf{DID}$ & Decentralized Identifier \\
$\mathsf{VC}$ & Verifiable Credential \\
\end{tabular}

\section*{Complejidad}
\begin{tabular}{ll}
$\mathcal{O}(f(n))$ & Cota superior asintótica \\
$\lambda$ & Parámetro de seguridad (típicamente 128 o 256 bits) \\
$\mathsf{negl}(\lambda)$ & Función negligible en $\lambda$ \\
$\mathsf{poly}(\lambda)$ & Función polinomial en $\lambda$ \\
PPT & Probabilistic Polynomial Time \\
\end{tabular}

\cleardoublepage

% ============================================
% FIN DEL VOLUMEN 1
% ============================================
\thispagestyle{empty}
\begin{center}
\vspace*{3cm}

{\Huge\textbf{Fin del Volumen 1}}

\vspace{2cm}

{\Large Continúa en:}

\vspace{0.5cm}

{\LARGE\textbf{Banco ZK: Volumen 2}}

\vspace{0.3cm}

{\large Aplicaciones Avanzadas y Ecosistema}

\vspace{0.2cm}

{\large Bloques 6--13}

\vspace{3cm}

\begin{tikzpicture}
    \node[rectangle, draw, rounded corners, minimum width=10cm, minimum height=3cm, fill=blue!10] {
        \begin{minipage}{9cm}
        \centering
        \textbf{Volumen 2 incluye:}\\[0.3cm]
        Compromisos $\bullet$ Infraestructura ZK $\bullet$ Privacidad Avanzada\\[0.2cm]
        DeFi $\bullet$ Anonimización $\bullet$ Privacy Coins\\[0.2cm]
        Unlinkability $\bullet$ Investigación Avanzada
        \end{minipage}
    };
\end{tikzpicture}

\end{center}
\cleardoublepage


\end{document}